
% !TEX root = ../../../../../Esperimenti/.tex/MEF.tex
% LTeX: language=it

\infopageKeysI

\begin{tikzpicture}%
    %\begingroup Righe
        \begin{groupplot}[
            group style={
                group name=Row1,
                plotGroup,
                group size=2 by 1,
            },
            plotRow,
            row5Keys,
            xTickScaleLabelEmpty,
            plotLegendSE,
            %
            xmin=0,xmax=1e+05,
            ymin=100,ymax=1.4483e+05,
        ]
            \nextgroupplot[%
                % xmin=0,xmax=1e+05,
                % ymin=100,ymax=1.4349e+05,
            ] % This file was created by matlab2tikz.
%
\definecolor{mycolor1}{rgb}{0.00146,0.00047,0.01387}%
\definecolor{mycolor2}{rgb}{0.01166,0.00942,0.06346}%
\definecolor{mycolor3}{rgb}{0.02943,0.02150,0.11462}%
\definecolor{mycolor4}{rgb}{0.06134,0.03659,0.17764}%
\definecolor{mycolor5}{rgb}{0.08741,0.04456,0.22481}%
\definecolor{mycolor6}{rgb}{0.12291,0.04754,0.28162}%
\definecolor{mycolor7}{rgb}{0.14907,0.04547,0.31709}%
\definecolor{mycolor8}{rgb}{0.18343,0.04033,0.35497}%
\definecolor{mycolor9}{rgb}{0.21795,0.03662,0.38352}%
\definecolor{mycolor10}{rgb}{0.24497,0.03705,0.40001}%
\definecolor{mycolor11}{rgb}{0.27135,0.04092,0.41198}%
\definecolor{mycolor12}{rgb}{0.29076,0.04564,0.41864}%
\definecolor{mycolor13}{rgb}{0.31628,0.05349,0.42512}%
\definecolor{mycolor14}{rgb}{0.33522,0.06006,0.42852}%
\definecolor{mycolor15}{rgb}{0.36028,0.06925,0.43150}%
\definecolor{mycolor16}{rgb}{0.37900,0.07625,0.43272}%
\definecolor{mycolor17}{rgb}{0.39767,0.08326,0.43318}%
\definecolor{mycolor18}{rgb}{0.41633,0.09020,0.43294}%
\definecolor{mycolor19}{rgb}{0.42877,0.09479,0.43241}%
\definecolor{mycolor20}{rgb}{0.44743,0.10160,0.43108}%
\definecolor{mycolor21}{rgb}{0.46610,0.10832,0.42912}%
\definecolor{mycolor22}{rgb}{0.47856,0.11276,0.42747}%
\definecolor{mycolor23}{rgb}{0.49726,0.11938,0.42449}%
\definecolor{mycolor24}{rgb}{0.50973,0.12377,0.42216}%
\definecolor{mycolor25}{rgb}{0.52221,0.12815,0.41955}%
\definecolor{mycolor26}{rgb}{0.53468,0.13253,0.41667}%
\definecolor{mycolor27}{rgb}{0.55339,0.13913,0.41183}%
\definecolor{mycolor28}{rgb}{0.56585,0.14357,0.40826}%
\definecolor{mycolor29}{rgb}{0.57830,0.14804,0.40441}%
\definecolor{mycolor30}{rgb}{0.59073,0.15256,0.40029}%
\definecolor{mycolor31}{rgb}{0.60314,0.15715,0.39589}%
\definecolor{mycolor32}{rgb}{0.61551,0.16182,0.39122}%
\definecolor{mycolor33}{rgb}{0.62169,0.16418,0.38878}%
\definecolor{mycolor34}{rgb}{0.63400,0.16899,0.38370}%
\definecolor{mycolor35}{rgb}{0.64626,0.17391,0.37836}%
\definecolor{mycolor36}{rgb}{0.65846,0.17896,0.37275}%
\definecolor{mycolor37}{rgb}{0.66454,0.18154,0.36985}%
\definecolor{mycolor38}{rgb}{0.67664,0.18681,0.36385}%
\definecolor{mycolor39}{rgb}{0.68865,0.19224,0.35760}%
\definecolor{mycolor40}{rgb}{0.69463,0.19502,0.35439}%
\definecolor{mycolor41}{rgb}{0.70650,0.20073,0.34778}%
\definecolor{mycolor42}{rgb}{0.71240,0.20366,0.34438}%
\definecolor{mycolor43}{rgb}{0.72410,0.20967,0.33742}%
\definecolor{mycolor44}{rgb}{0.72991,0.21276,0.33386}%
\definecolor{mycolor45}{rgb}{0.74142,0.21911,0.32658}%
\definecolor{mycolor46}{rgb}{0.74713,0.22238,0.32286}%
\definecolor{mycolor47}{rgb}{0.75279,0.22571,0.31909}%
\definecolor{mycolor48}{rgb}{0.76401,0.23255,0.31140}%
\definecolor{mycolor49}{rgb}{0.76956,0.23608,0.30749}%
\definecolor{mycolor50}{rgb}{0.77506,0.23967,0.30353}%
\definecolor{mycolor51}{rgb}{0.79129,0.25086,0.29139}%
\definecolor{mycolor52}{rgb}{0.79661,0.25473,0.28726}%
\definecolor{mycolor53}{rgb}{0.80187,0.25867,0.28310}%
\definecolor{mycolor54}{rgb}{0.81224,0.26679,0.27466}%
\definecolor{mycolor55}{rgb}{0.81734,0.27095,0.27039}%
\definecolor{mycolor56}{rgb}{0.82737,0.27952,0.26175}%
\definecolor{mycolor57}{rgb}{0.83230,0.28391,0.25738}%
\definecolor{mycolor58}{rgb}{0.83717,0.28839,0.25299}%
\definecolor{mycolor59}{rgb}{0.84671,0.29756,0.24411}%
\definecolor{mycolor60}{rgb}{0.85138,0.30226,0.23964}%
\definecolor{mycolor61}{rgb}{0.85599,0.30704,0.23513}%
\definecolor{mycolor62}{rgb}{0.86053,0.31189,0.23061}%
\definecolor{mycolor63}{rgb}{0.87374,0.32691,0.21689}%
\definecolor{mycolor64}{rgb}{0.87800,0.33206,0.21227}%
\definecolor{mycolor65}{rgb}{0.89034,0.34796,0.19829}%
\definecolor{mycolor66}{rgb}{0.89819,0.35891,0.18886}%
\definecolor{mycolor67}{rgb}{0.90200,0.36449,0.18412}%
\definecolor{mycolor68}{rgb}{0.90573,0.37014,0.17935}%
\definecolor{mycolor69}{rgb}{0.90939,0.37586,0.17456}%
\definecolor{mycolor70}{rgb}{0.91297,0.38164,0.16975}%
\definecolor{mycolor71}{rgb}{0.91646,0.38748,0.16492}%
\definecolor{mycolor72}{rgb}{0.91988,0.39339,0.16007}%
\definecolor{mycolor73}{rgb}{0.92322,0.39936,0.15519}%
\definecolor{mycolor74}{rgb}{0.92647,0.40539,0.15029}%
\definecolor{mycolor75}{rgb}{0.92964,0.41148,0.14537}%
\definecolor{mycolor76}{rgb}{0.93274,0.41763,0.14042}%
\definecolor{mycolor77}{rgb}{0.93575,0.42383,0.13544}%
\definecolor{mycolor78}{rgb}{0.93868,0.43009,0.13044}%
\definecolor{mycolor79}{rgb}{0.94152,0.43640,0.12541}%
\definecolor{mycolor80}{rgb}{0.94429,0.44277,0.12035}%
\definecolor{mycolor81}{rgb}{0.94956,0.45566,0.11016}%
\definecolor{mycolor82}{rgb}{0.95208,0.46218,0.10503}%
\definecolor{mycolor83}{rgb}{0.96129,0.48872,0.08429}%
\definecolor{mycolor84}{rgb}{0.96540,0.50225,0.07386}%
\definecolor{mycolor85}{rgb}{0.96732,0.50908,0.06866}%
\definecolor{mycolor86}{rgb}{0.97259,0.52980,0.05332}%
\definecolor{mycolor87}{rgb}{0.97568,0.54380,0.04362}%
\definecolor{mycolor88}{rgb}{0.98082,0.57221,0.02851}%
\definecolor{mycolor89}{rgb}{0.98189,0.57939,0.02625}%
\definecolor{mycolor90}{rgb}{0.98378,0.59385,0.02377}%
\definecolor{mycolor91}{rgb}{0.98532,0.60842,0.02420}%
\definecolor{mycolor92}{rgb}{0.98696,0.63048,0.03091}%
\definecolor{mycolor93}{rgb}{0.98793,0.66025,0.05175}%
\definecolor{mycolor94}{rgb}{0.98771,0.68281,0.07249}%
\definecolor{mycolor95}{rgb}{0.97750,0.78226,0.18592}%
\definecolor{mycolor96}{rgb}{0.96624,0.83619,0.26153}%
\definecolor{mycolor97}{rgb}{0.96252,0.85148,0.28555}%
\definecolor{mycolor98}{rgb}{0.98836,0.99836,0.64492}%
%
\addplot [color=mycolor1]
  table[row sep=crcr]{%
0	4395.3\\
2000	16908\\
4000	36693\\
6000	52482\\
8000	65435\\
10000	76041\\
12000	84910\\
14000	92770\\
16000	99400\\
18000	1.0487e+05\\
20000	1.0976e+05\\
22000	1.1403e+05\\
24000	1.1768e+05\\
28000	1.238e+05\\
30000	1.262e+05\\
32000	1.281e+05\\
36000	1.3112e+05\\
40000	1.3436e+05\\
44000	1.3655e+05\\
48000	1.3817e+05\\
64000	1.4164e+05\\
80000	1.4256e+05\\
94000	1.4349e+05\\
98000	1.4335e+05\\
1e+05	1.4341e+05\\
};
\addlegendentry{Classe k=283}

\addplot [color=mycolor2, forget plot]
  table[row sep=crcr]{%
0	4395.3\\
2000	9728.2\\
4000	16754\\
6000	20791\\
8000	23191\\
10000	24588\\
12000	25379\\
14000	25776\\
18000	26041\\
20000	26119\\
22000	25967\\
28000	25921\\
30000	25764\\
36000	25789\\
40000	25483\\
42000	25430\\
44000	25473\\
46000	25224\\
50000	25044\\
52000	25139\\
58000	25091\\
60000	25032\\
62000	25233\\
74000	25298\\
78000	25068\\
82000	25305\\
84000	25345\\
86000	25065\\
88000	25221\\
90000	25167\\
92000	25007\\
94000	25052\\
96000	25246\\
98000	25228\\
1e+05	25076\\
};
\addplot [color=mycolor3, forget plot]
  table[row sep=crcr]{%
0	4395.3\\
2000	10931\\
4000	20153\\
6000	26219\\
8000	30318\\
10000	33199\\
12000	35166\\
14000	36202\\
16000	36896\\
18000	37074\\
20000	37149\\
28000	38160\\
32000	38097\\
34000	38305\\
38000	37969\\
44000	37594\\
46000	37398\\
50000	37441\\
52000	37680\\
56000	37644\\
60000	37877\\
66000	37756\\
68000	37884\\
76000	37663\\
78000	37792\\
82000	37622\\
86000	37963\\
88000	37968\\
94000	38411\\
1e+05	38225\\
};
\addplot [color=mycolor4, forget plot]
  table[row sep=crcr]{%
0	4395.3\\
2000	7749.6\\
4000	11740\\
6000	13738\\
8000	14689\\
10000	15169\\
12000	15279\\
18000	15056\\
22000	14817\\
26000	14561\\
28000	14627\\
30000	14426\\
32000	14332\\
34000	14375\\
36000	14373\\
38000	14287\\
40000	14408\\
42000	14355\\
44000	14129\\
46000	14010\\
50000	14117\\
52000	14028\\
54000	14171\\
56000	14154\\
58000	14075\\
60000	14123\\
62000	14014\\
64000	14022\\
68000	13888\\
70000	14032\\
72000	13944\\
74000	13978\\
76000	13909\\
78000	14069\\
80000	14144\\
82000	14039\\
84000	13997\\
88000	14046\\
90000	14007\\
92000	13920\\
94000	13782\\
96000	13973\\
98000	13994\\
1e+05	13954\\
};
\addplot [color=mycolor5, forget plot]
  table[row sep=crcr]{%
0	4395.3\\
2000	10267\\
4000	19800\\
6000	27442\\
8000	33546\\
10000	38473\\
12000	42692\\
14000	45961\\
16000	48667\\
18000	51031\\
20000	53107\\
22000	54758\\
24000	56042\\
28000	58029\\
34000	60289\\
38000	61544\\
42000	61863\\
46000	61910\\
52000	62397\\
56000	62564\\
58000	62365\\
64000	62519\\
68000	62280\\
70000	62017\\
84000	61528\\
88000	61269\\
92000	61190\\
94000	61433\\
1e+05	61633\\
};
\addplot [color=mycolor6, forget plot]
  table[row sep=crcr]{%
0	4395.3\\
2000	9511.3\\
4000	17544\\
6000	23829\\
8000	28641\\
10000	32522\\
12000	35637\\
14000	38312\\
16000	40653\\
18000	42495\\
20000	43940\\
22000	45305\\
24000	46442\\
28000	47801\\
34000	49610\\
40000	50316\\
46000	51384\\
48000	51578\\
50000	51981\\
52000	52092\\
54000	52011\\
66000	52559\\
70000	52464\\
80000	52904\\
84000	52584\\
88000	53124\\
92000	53132\\
1e+05	53254\\
};
\addplot [color=mycolor7, forget plot]
  table[row sep=crcr]{%
0	4395.3\\
2000	7968.5\\
4000	13166\\
6000	16426\\
8000	18757\\
10000	20544\\
12000	21814\\
14000	22739\\
16000	23369\\
18000	23927\\
20000	24292\\
30000	25393\\
34000	25737\\
38000	25634\\
42000	26110\\
44000	26182\\
46000	26381\\
48000	26331\\
56000	26762\\
58000	26427\\
60000	26562\\
62000	26467\\
66000	26799\\
68000	26976\\
74000	26969\\
80000	26717\\
82000	26977\\
84000	26820\\
86000	26956\\
88000	27173\\
94000	27075\\
96000	27105\\
98000	27242\\
1e+05	27210\\
};
\addplot [color=mycolor8, forget plot]
  table[row sep=crcr]{%
0	4395.3\\
2000	7767.9\\
4000	12561\\
6000	15752\\
8000	17761\\
10000	19098\\
12000	20002\\
14000	20693\\
16000	21115\\
20000	21400\\
22000	21471\\
26000	21391\\
28000	21522\\
30000	21475\\
32000	21347\\
34000	21382\\
36000	21228\\
42000	20950\\
44000	20850\\
46000	20825\\
48000	20895\\
50000	20879\\
54000	20569\\
58000	20505\\
66000	20374\\
72000	20348\\
74000	20223\\
76000	20182\\
78000	20312\\
84000	19959\\
92000	20389\\
96000	20225\\
1e+05	20086\\
};
\addplot [color=mycolor9, forget plot]
  table[row sep=crcr]{%
0	4395.3\\
2000	6531.4\\
4000	9323.9\\
6000	10845\\
8000	11686\\
10000	12225\\
12000	12608\\
14000	12886\\
16000	13061\\
18000	12966\\
22000	13017\\
26000	13184\\
28000	13118\\
30000	13149\\
34000	13452\\
36000	13475\\
38000	13337\\
42000	13363\\
44000	13410\\
48000	13315\\
50000	13331\\
52000	13309\\
54000	13379\\
58000	13368\\
60000	13432\\
62000	13550\\
68000	13462\\
70000	13594\\
72000	13618\\
74000	13574\\
76000	13581\\
80000	13463\\
82000	13500\\
84000	13695\\
86000	13682\\
90000	13750\\
92000	13688\\
96000	13715\\
98000	13685\\
1e+05	13604\\
};
\addplot [color=mycolor10, forget plot]
  table[row sep=crcr]{%
0	4395.3\\
2000	7018.2\\
4000	10740\\
6000	13224\\
8000	14952\\
10000	16272\\
12000	17164\\
14000	17746\\
16000	18259\\
18000	18664\\
20000	18974\\
22000	19235\\
30000	20117\\
34000	20312\\
38000	20559\\
40000	20537\\
42000	20588\\
46000	20972\\
48000	20961\\
52000	21115\\
54000	21010\\
60000	21080\\
64000	21119\\
66000	21129\\
68000	20972\\
70000	20948\\
72000	21025\\
74000	21173\\
76000	21155\\
78000	21204\\
80000	21169\\
82000	21336\\
84000	21391\\
86000	21367\\
88000	21212\\
90000	21250\\
96000	21162\\
98000	21177\\
1e+05	21289\\
};
\addplot [color=mycolor11, forget plot]
  table[row sep=crcr]{%
0	4395.3\\
2000	6303\\
4000	8658.4\\
6000	9804\\
8000	10253\\
10000	10310\\
12000	10236\\
16000	9927.4\\
18000	9739.8\\
22000	9550.5\\
24000	9329\\
26000	9223\\
30000	9090.9\\
32000	9125.3\\
34000	9043.1\\
36000	8996.1\\
38000	9010.5\\
42000	8888.4\\
44000	8843\\
46000	8825.3\\
48000	8844.2\\
50000	8740.8\\
56000	8631.1\\
58000	8645.1\\
60000	8612.6\\
62000	8635.5\\
64000	8682.6\\
66000	8673.3\\
68000	8557.8\\
70000	8523.3\\
72000	8648.3\\
78000	8707.4\\
80000	8628.4\\
82000	8576.4\\
84000	8581\\
86000	8619\\
88000	8697.5\\
94000	8636.7\\
96000	8687.4\\
1e+05	8591.8\\
};
\addplot [color=mycolor12, forget plot]
  table[row sep=crcr]{%
0	4395.3\\
2000	6383.4\\
4000	9089\\
6000	10711\\
8000	11818\\
10000	12460\\
12000	12798\\
16000	13337\\
20000	13462\\
22000	13606\\
24000	13636\\
28000	13472\\
30000	13610\\
32000	13884\\
34000	14003\\
36000	14079\\
38000	14117\\
40000	14112\\
44000	13900\\
52000	14116\\
54000	14246\\
56000	14317\\
58000	14349\\
60000	14206\\
62000	14252\\
66000	14233\\
70000	14068\\
72000	14058\\
76000	13911\\
80000	14111\\
82000	14056\\
86000	14082\\
88000	14141\\
92000	14085\\
94000	14226\\
96000	14436\\
98000	14428\\
1e+05	14203\\
};
\addplot [color=mycolor13, forget plot]
  table[row sep=crcr]{%
0	4395.3\\
2000	6044.3\\
4000	8129.8\\
6000	9340.9\\
8000	9955\\
10000	10291\\
14000	10436\\
16000	10515\\
20000	10576\\
22000	10721\\
26000	10473\\
28000	10427\\
32000	10539\\
34000	10572\\
38000	10544\\
40000	10640\\
42000	10664\\
44000	10613\\
56000	10841\\
58000	10780\\
62000	10737\\
64000	10747\\
66000	10675\\
68000	10748\\
72000	10765\\
76000	10900\\
78000	10829\\
80000	10794\\
84000	10884\\
86000	10828\\
92000	10822\\
94000	10908\\
98000	10824\\
1e+05	10901\\
};
\addplot [color=mycolor14, forget plot]
  table[row sep=crcr]{%
0	4395.3\\
2000	6886.3\\
4000	10583\\
6000	13230\\
8000	15287\\
10000	16672\\
12000	17762\\
14000	18734\\
16000	19478\\
18000	20028\\
20000	20456\\
22000	20832\\
24000	21303\\
26000	21546\\
28000	21721\\
30000	21787\\
32000	21777\\
34000	21869\\
36000	22173\\
40000	22368\\
46000	22509\\
48000	22743\\
50000	22817\\
52000	22634\\
56000	22828\\
58000	23031\\
60000	23057\\
62000	23005\\
66000	23119\\
68000	22927\\
70000	23132\\
72000	23160\\
74000	23053\\
76000	23264\\
80000	23445\\
86000	23234\\
88000	23300\\
90000	23282\\
92000	23032\\
94000	23123\\
96000	23049\\
1e+05	23026\\
};
\addplot [color=mycolor15, forget plot]
  table[row sep=crcr]{%
0	4395.3\\
2000	6282.8\\
4000	8937\\
6000	10813\\
8000	12102\\
10000	13020\\
12000	13567\\
16000	14443\\
18000	14824\\
20000	15010\\
24000	15274\\
26000	15459\\
28000	15481\\
30000	15705\\
34000	15755\\
36000	15746\\
40000	15585\\
46000	15990\\
48000	16064\\
50000	16207\\
60000	16259\\
62000	16148\\
64000	16180\\
66000	16264\\
70000	16281\\
72000	16468\\
74000	16579\\
76000	16639\\
78000	16578\\
80000	16459\\
82000	16412\\
84000	16281\\
90000	16380\\
94000	16270\\
96000	16275\\
98000	16335\\
1e+05	16201\\
};
\addplot [color=mycolor16, forget plot]
  table[row sep=crcr]{%
0	4395.3\\
2000	5763\\
4000	7424.1\\
6000	8156.1\\
8000	8397.7\\
10000	8436.4\\
12000	8407.3\\
14000	8320.1\\
16000	8208.1\\
20000	7942.4\\
22000	7848.9\\
26000	7577.6\\
28000	7441.2\\
30000	7373.6\\
32000	7342.6\\
34000	7277.1\\
36000	7240.7\\
38000	7156.1\\
42000	7034.1\\
44000	7065.4\\
46000	7018.3\\
48000	7042.4\\
50000	6982.2\\
52000	7023.4\\
54000	7003\\
56000	6950.3\\
62000	7098.8\\
64000	7079.6\\
66000	7092.7\\
68000	7047.8\\
70000	7024.4\\
72000	7033.3\\
74000	6999.5\\
76000	7005.7\\
78000	6965.4\\
80000	6903.4\\
82000	6916\\
84000	6888\\
86000	6882.6\\
88000	6958.9\\
90000	6971.9\\
92000	6872.8\\
94000	6869.2\\
96000	6936.6\\
98000	6929\\
1e+05	6967\\
};
\addplot [color=mycolor17, forget plot]
  table[row sep=crcr]{%
0	4395.3\\
2000	6341.9\\
4000	8969.6\\
6000	10667\\
8000	11722\\
10000	12498\\
14000	13485\\
16000	13864\\
18000	14050\\
20000	14318\\
28000	14537\\
30000	14672\\
32000	14697\\
34000	14820\\
36000	15010\\
40000	15114\\
42000	15254\\
44000	15130\\
46000	15116\\
48000	15041\\
52000	15068\\
54000	15033\\
58000	15071\\
62000	15324\\
64000	15205\\
66000	15146\\
70000	15260\\
72000	15423\\
74000	15465\\
76000	15334\\
78000	15329\\
82000	15451\\
86000	15451\\
88000	15392\\
90000	15456\\
92000	15575\\
94000	15398\\
98000	15426\\
1e+05	15489\\
};
\addplot [color=mycolor18, forget plot]
  table[row sep=crcr]{%
0	4395.3\\
4000	10355\\
6000	12939\\
8000	14893\\
10000	16388\\
12000	17638\\
14000	18602\\
18000	20090\\
20000	20612\\
22000	20986\\
24000	21147\\
26000	21413\\
28000	21787\\
32000	22257\\
34000	22324\\
36000	22266\\
38000	22474\\
40000	22556\\
44000	22887\\
50000	22915\\
52000	22889\\
54000	23133\\
56000	23001\\
58000	22937\\
60000	22964\\
64000	23174\\
66000	23101\\
68000	23142\\
70000	23373\\
76000	23514\\
78000	23459\\
80000	23474\\
82000	23608\\
92000	23520\\
96000	23596\\
1e+05	23570\\
};
\addplot [color=mycolor19, forget plot]
  table[row sep=crcr]{%
0	4395.3\\
2000	5972.1\\
4000	8221.2\\
6000	9773.6\\
8000	10874\\
10000	11696\\
12000	12290\\
14000	12641\\
16000	12901\\
18000	13056\\
20000	13377\\
22000	13549\\
24000	13663\\
26000	13835\\
30000	13775\\
32000	13747\\
36000	13914\\
40000	14142\\
42000	14089\\
44000	14129\\
46000	13913\\
50000	14021\\
52000	14141\\
54000	14218\\
56000	14385\\
58000	14430\\
60000	14319\\
66000	14146\\
68000	14282\\
70000	14367\\
72000	14235\\
74000	14218\\
76000	14341\\
78000	14247\\
82000	14350\\
84000	14276\\
86000	14344\\
88000	14540\\
90000	14595\\
92000	14438\\
94000	14378\\
96000	14455\\
98000	14433\\
1e+05	14452\\
};
\addplot [color=mycolor20, forget plot]
  table[row sep=crcr]{%
0	4395.3\\
2000	6275.1\\
4000	9279.7\\
6000	11370\\
8000	12879\\
10000	13993\\
12000	14809\\
14000	15414\\
16000	15885\\
18000	16294\\
20000	16648\\
24000	17032\\
26000	17341\\
30000	17707\\
32000	17635\\
34000	17639\\
40000	17962\\
42000	17844\\
44000	17912\\
46000	18134\\
50000	18272\\
54000	18262\\
56000	18374\\
58000	18373\\
60000	18175\\
62000	18174\\
66000	18401\\
68000	18334\\
70000	18482\\
72000	18527\\
76000	18432\\
80000	18422\\
82000	18290\\
84000	18395\\
92000	18583\\
94000	18423\\
1e+05	18490\\
};
\addplot [color=mycolor21, forget plot]
  table[row sep=crcr]{%
0	4395.3\\
2000	5238\\
4000	6398.4\\
6000	7109\\
8000	7565\\
10000	7921.6\\
12000	8186.6\\
14000	8401.6\\
16000	8557.4\\
18000	8680.3\\
24000	8964.7\\
26000	8958.8\\
28000	9013.6\\
30000	9041.9\\
32000	9031.7\\
34000	9052.1\\
36000	9167.4\\
38000	9200.3\\
40000	9327.7\\
42000	9314\\
44000	9259.3\\
46000	9298.5\\
48000	9276.6\\
50000	9301.9\\
54000	9286.5\\
56000	9294.6\\
58000	9331.3\\
60000	9401\\
64000	9447.2\\
66000	9403.3\\
68000	9387.1\\
70000	9434.6\\
72000	9442.3\\
74000	9486.1\\
76000	9579.6\\
78000	9518.1\\
80000	9486.3\\
84000	9527\\
86000	9502.2\\
88000	9553.8\\
90000	9547.3\\
92000	9495.9\\
94000	9513.1\\
96000	9565.8\\
98000	9505\\
1e+05	9493.8\\
};
\addplot [color=mycolor22, forget plot]
  table[row sep=crcr]{%
0	4395.3\\
2000	7240.1\\
4000	11772\\
6000	15273\\
8000	17913\\
10000	19870\\
12000	21445\\
14000	22645\\
16000	23531\\
18000	24187\\
22000	25070\\
24000	25647\\
26000	26053\\
28000	26181\\
34000	26121\\
36000	26212\\
38000	26159\\
44000	26216\\
50000	25997\\
56000	25790\\
58000	25687\\
60000	25836\\
62000	25703\\
70000	25496\\
72000	25553\\
74000	25517\\
76000	25322\\
80000	25200\\
82000	25327\\
88000	25292\\
92000	25007\\
94000	25100\\
1e+05	25047\\
};
\addplot [color=mycolor23, forget plot]
  table[row sep=crcr]{%
0	4395.3\\
2000	5744.4\\
4000	7678\\
6000	8973.6\\
8000	9853.1\\
10000	10486\\
12000	10892\\
14000	11144\\
18000	11541\\
22000	11776\\
26000	11972\\
28000	12009\\
30000	12104\\
32000	12123\\
34000	12087\\
36000	12117\\
40000	12054\\
42000	12074\\
44000	12138\\
48000	12161\\
50000	12088\\
54000	12106\\
58000	12058\\
62000	12101\\
68000	12198\\
72000	12095\\
74000	12103\\
76000	12190\\
78000	12202\\
84000	12142\\
86000	12140\\
88000	12070\\
94000	12015\\
96000	12033\\
98000	12111\\
1e+05	12092\\
};
\addplot [color=mycolor24, forget plot]
  table[row sep=crcr]{%
0	4395.3\\
2000	5059.4\\
4000	5872.3\\
6000	6204.4\\
8000	6227.2\\
10000	6103.1\\
12000	6024.1\\
14000	5907.8\\
16000	5760.7\\
18000	5668.9\\
22000	5400.8\\
24000	5336.8\\
26000	5223.6\\
28000	5166.9\\
30000	5190.4\\
32000	5063.2\\
34000	4979.6\\
36000	4978.5\\
40000	4913.6\\
42000	4914.9\\
44000	4900.7\\
46000	4869.1\\
48000	4821.8\\
52000	4833.8\\
54000	4807.8\\
56000	4827.8\\
60000	4814.7\\
62000	4851.9\\
64000	4820.8\\
66000	4828.7\\
68000	4856.1\\
72000	4816.3\\
74000	4764.9\\
76000	4773.9\\
78000	4766.8\\
80000	4834.5\\
82000	4848.2\\
84000	4830.1\\
86000	4783.7\\
88000	4773.7\\
92000	4852.8\\
96000	4775.1\\
98000	4762\\
1e+05	4822.1\\
};
\addplot [color=mycolor25, forget plot]
  table[row sep=crcr]{%
0	4395.3\\
2000	5635.9\\
4000	7491.6\\
6000	8716\\
8000	9501\\
10000	10107\\
12000	10563\\
14000	10846\\
18000	11186\\
20000	11142\\
22000	11252\\
24000	11440\\
26000	11423\\
38000	11807\\
44000	11748\\
46000	11801\\
50000	11705\\
54000	11908\\
56000	11864\\
60000	11945\\
62000	11851\\
64000	11893\\
66000	11887\\
68000	11938\\
70000	11924\\
72000	12007\\
78000	12060\\
80000	12172\\
82000	12110\\
86000	12123\\
88000	12146\\
90000	12110\\
92000	12181\\
94000	12194\\
96000	12141\\
98000	12162\\
1e+05	12116\\
};
\addplot [color=mycolor26, forget plot]
  table[row sep=crcr]{%
0	4395.3\\
2000	5286.7\\
4000	6321.8\\
6000	6948\\
8000	7231\\
12000	7484.9\\
14000	7495.5\\
20000	7708.8\\
22000	7680.4\\
24000	7713.6\\
26000	7693.1\\
28000	7780\\
34000	7768.1\\
36000	7797.9\\
38000	7905.7\\
40000	7861.4\\
42000	7895.4\\
44000	7967.2\\
46000	7889\\
48000	7989.7\\
50000	8041.5\\
56000	8117.7\\
58000	8002.3\\
60000	7959.3\\
62000	7878.2\\
64000	7966.8\\
66000	7971\\
68000	8042.2\\
72000	8053.3\\
74000	8158.1\\
76000	8160\\
78000	8108.6\\
80000	8106\\
82000	8140.6\\
86000	8171.5\\
88000	8106.9\\
90000	8138\\
96000	8082.4\\
98000	8077.9\\
1e+05	8167\\
};
\addplot [color=mycolor27, forget plot]
  table[row sep=crcr]{%
0	4395.3\\
2000	5678.1\\
4000	7463\\
6000	8542.3\\
8000	9240.9\\
10000	9613.6\\
12000	9845.7\\
14000	10002\\
16000	10129\\
20000	10128\\
22000	10067\\
26000	10024\\
30000	9866.2\\
32000	9924.5\\
34000	9779.6\\
36000	9714.6\\
40000	9749.6\\
42000	9744.5\\
44000	9639.3\\
46000	9578\\
48000	9612.2\\
50000	9725.4\\
52000	9712.3\\
54000	9787.6\\
58000	9720\\
60000	9703.6\\
64000	9775.1\\
66000	9757.7\\
68000	9665.7\\
70000	9709.7\\
72000	9662.2\\
74000	9695.6\\
76000	9694.2\\
78000	9806.4\\
80000	9729.5\\
82000	9731.6\\
84000	9673.3\\
86000	9577.5\\
88000	9562.3\\
90000	9575.5\\
92000	9450.9\\
94000	9465.3\\
98000	9550\\
1e+05	9528.4\\
};
\addplot [color=mycolor28, forget plot]
  table[row sep=crcr]{%
0	4395.3\\
2000	5074.3\\
4000	5923.4\\
6000	6339.4\\
8000	6483\\
10000	6567.8\\
12000	6572.9\\
14000	6512.9\\
16000	6368.7\\
18000	6296.7\\
20000	6250.3\\
22000	6130\\
24000	6145.1\\
26000	6110.3\\
28000	5959.1\\
30000	5829.9\\
32000	5854.6\\
34000	5792\\
36000	5774.9\\
38000	5777.5\\
40000	5692.9\\
44000	5579.7\\
48000	5617.7\\
50000	5651.6\\
52000	5585.1\\
60000	5571.1\\
62000	5547.8\\
64000	5566.8\\
66000	5612.5\\
72000	5551.2\\
74000	5480.8\\
76000	5458\\
78000	5414.9\\
80000	5466.1\\
82000	5471.8\\
86000	5542.8\\
88000	5550.1\\
90000	5475.3\\
92000	5444.5\\
96000	5462.5\\
98000	5429.2\\
1e+05	5465.8\\
};
\addplot [color=mycolor29, forget plot]
  table[row sep=crcr]{%
0	4395.3\\
2000	5862.6\\
4000	7995.5\\
6000	9313.4\\
8000	10197\\
10000	10827\\
12000	11196\\
14000	11413\\
18000	11778\\
20000	11895\\
22000	11764\\
24000	11748\\
26000	11865\\
28000	11857\\
30000	11902\\
32000	11767\\
34000	11681\\
36000	11783\\
38000	11806\\
46000	11574\\
50000	11660\\
52000	11596\\
54000	11583\\
56000	11505\\
58000	11484\\
62000	11346\\
64000	11252\\
66000	11347\\
68000	11488\\
70000	11588\\
74000	11544\\
76000	11457\\
78000	11434\\
80000	11539\\
82000	11514\\
86000	11346\\
88000	11286\\
90000	11319\\
92000	11285\\
94000	11437\\
96000	11421\\
98000	11346\\
1e+05	11390\\
};
\addplot [color=mycolor30, forget plot]
  table[row sep=crcr]{%
0	4395.3\\
2000	4916.6\\
4000	5667.4\\
6000	6170.5\\
8000	6517.3\\
10000	6775.3\\
12000	6922.5\\
14000	7045\\
16000	7136.7\\
18000	7262.4\\
22000	7393.2\\
24000	7432\\
28000	7480.2\\
30000	7503.7\\
32000	7480.6\\
34000	7580.8\\
36000	7609.1\\
42000	7569.7\\
44000	7601.9\\
48000	7566.9\\
50000	7591.8\\
52000	7640.8\\
54000	7635.3\\
60000	7769\\
62000	7733.9\\
66000	7739\\
68000	7828.6\\
70000	7850.2\\
72000	7789.6\\
78000	7743.4\\
82000	7745.4\\
84000	7801.2\\
86000	7804.7\\
88000	7862.8\\
90000	7898.6\\
92000	7856.5\\
94000	7728.8\\
96000	7703.2\\
98000	7715.9\\
1e+05	7801.8\\
};
\addplot [color=mycolor31, forget plot]
  table[row sep=crcr]{%
0	4395.3\\
2000	5193.2\\
4000	6227.8\\
6000	6840.8\\
8000	7113.7\\
12000	7303.3\\
14000	7286.7\\
16000	7216.1\\
18000	7190.2\\
20000	7140.4\\
24000	6925\\
26000	6939.4\\
28000	6906.2\\
34000	6666.7\\
36000	6647.2\\
40000	6713.8\\
42000	6703.8\\
48000	6782.1\\
50000	6785.8\\
52000	6769.6\\
54000	6688.1\\
56000	6691.8\\
58000	6629.1\\
66000	6703.4\\
68000	6649.8\\
70000	6621.5\\
72000	6622\\
74000	6641.7\\
76000	6640.8\\
78000	6580.6\\
80000	6573.1\\
82000	6500.9\\
84000	6465.4\\
86000	6467.3\\
90000	6528.1\\
92000	6559\\
96000	6490\\
98000	6442.3\\
1e+05	6457.4\\
};
\addplot [color=mycolor32, forget plot]
  table[row sep=crcr]{%
0	4395.3\\
2000	5700.5\\
4000	7606.1\\
6000	8917.8\\
8000	9797.6\\
10000	10335\\
14000	10924\\
16000	11113\\
18000	11266\\
26000	11154\\
28000	11170\\
32000	11112\\
34000	11117\\
44000	10865\\
48000	10689\\
50000	10719\\
52000	10707\\
54000	10598\\
56000	10568\\
58000	10569\\
60000	10532\\
64000	10571\\
66000	10554\\
68000	10589\\
70000	10653\\
72000	10750\\
74000	10762\\
76000	10703\\
80000	10684\\
84000	10554\\
88000	10609\\
90000	10574\\
94000	10595\\
96000	10581\\
98000	10490\\
1e+05	10465\\
};
\addplot [color=mycolor33, forget plot]
  table[row sep=crcr]{%
0	4395.3\\
2000	4877.4\\
4000	5633.9\\
6000	6163.1\\
8000	6537.2\\
12000	6972.8\\
14000	7147.9\\
16000	7253.5\\
20000	7263.5\\
24000	7521.7\\
26000	7556.1\\
28000	7622.7\\
30000	7728.7\\
32000	7773.2\\
36000	7812.1\\
38000	7758.5\\
42000	7896.7\\
44000	7898.3\\
46000	7968.5\\
48000	8000.3\\
50000	8003.5\\
52000	8083.4\\
54000	8076.4\\
56000	8101.6\\
58000	8080.9\\
62000	7992\\
64000	7972.7\\
66000	8055.5\\
68000	8095.1\\
70000	8084.3\\
72000	8042.1\\
74000	8064.1\\
76000	8167.1\\
78000	8203.5\\
80000	8305.9\\
82000	8229\\
84000	8237.2\\
86000	8207.3\\
88000	8255.6\\
90000	8165.6\\
92000	8197.5\\
94000	8256.7\\
96000	8287.1\\
1e+05	8237.3\\
};
\addplot [color=mycolor34, forget plot]
  table[row sep=crcr]{%
0	4395.3\\
2000	4922\\
4000	5417.2\\
6000	5517.1\\
8000	5404.6\\
10000	5267\\
12000	5106\\
14000	4990.8\\
16000	4909.2\\
18000	4774.2\\
22000	4557.3\\
24000	4415.5\\
26000	4301.9\\
30000	4200.8\\
32000	4188.4\\
34000	4217\\
36000	4221.8\\
38000	4190.9\\
40000	4138.2\\
42000	4120.9\\
44000	4092.2\\
48000	4081.6\\
50000	4086.6\\
52000	4071.8\\
54000	4073.6\\
56000	4060.5\\
60000	3971.6\\
62000	3999.2\\
66000	3996.6\\
68000	4055\\
72000	4042.3\\
74000	4059.3\\
76000	4012.4\\
82000	4028.2\\
84000	4017.5\\
86000	4058.3\\
88000	4076\\
90000	4031.8\\
96000	3974.3\\
1e+05	4052.7\\
};
\addplot [color=mycolor35]
  table[row sep=crcr]{%
0	4395.3\\
2000	5376.2\\
4000	6803.5\\
6000	7785\\
8000	8428.4\\
10000	8889\\
12000	9194.8\\
14000	9394\\
16000	9563.5\\
18000	9631.8\\
24000	9718.7\\
26000	9769.2\\
28000	9766.1\\
30000	9799.3\\
32000	9751.7\\
36000	9754.5\\
38000	9780.7\\
42000	9770.6\\
44000	9798.7\\
46000	9773.2\\
50000	9813.6\\
54000	9778.2\\
56000	9793.6\\
58000	9843\\
60000	9839.9\\
62000	9878.2\\
66000	9821\\
74000	9928\\
76000	9911.2\\
78000	9948.3\\
80000	10016\\
84000	9894.7\\
86000	9872.7\\
88000	9900.4\\
90000	9890.6\\
98000	10033\\
1e+05	9986.3\\
};
\addlegendentry{Classe k=83}

\addplot [color=mycolor36, forget plot]
  table[row sep=crcr]{%
0	4395.3\\
4000	5315.2\\
6000	5484\\
8000	5547.6\\
10000	5456\\
12000	5402\\
18000	5076.5\\
20000	5035.3\\
22000	4912.9\\
24000	4854.5\\
28000	4720.1\\
32000	4691.6\\
34000	4664.2\\
36000	4598.6\\
38000	4603.7\\
42000	4455.2\\
44000	4443.5\\
46000	4447.7\\
48000	4415.9\\
50000	4405\\
52000	4417\\
54000	4416.5\\
56000	4393.8\\
58000	4358.8\\
62000	4396.5\\
68000	4361.1\\
72000	4373.4\\
74000	4333.5\\
76000	4329.7\\
78000	4344.5\\
80000	4389.6\\
84000	4437\\
90000	4324.2\\
94000	4377.2\\
98000	4376.1\\
1e+05	4318.6\\
};
\addplot [color=mycolor37, forget plot]
  table[row sep=crcr]{%
0	4395.3\\
4000	4465\\
6000	4413.9\\
8000	4339.3\\
10000	4294.3\\
12000	4224.1\\
14000	4208.9\\
16000	4137.1\\
20000	4023.5\\
22000	3950.1\\
26000	3869.7\\
28000	3855.5\\
30000	3812.6\\
32000	3788.3\\
34000	3774.6\\
38000	3775\\
40000	3747.9\\
44000	3734.2\\
46000	3749.9\\
48000	3713.9\\
52000	3712.8\\
58000	3677\\
60000	3692\\
64000	3702.1\\
66000	3681.4\\
68000	3671.5\\
72000	3725.3\\
74000	3716.5\\
76000	3719.3\\
80000	3668.8\\
84000	3714.1\\
86000	3697.1\\
88000	3713.4\\
90000	3668\\
92000	3680\\
96000	3686.5\\
98000	3661.1\\
1e+05	3703\\
};
\addplot [color=mycolor38, forget plot]
  table[row sep=crcr]{%
0	4395.3\\
2000	4190.2\\
4000	3875.3\\
6000	3716.2\\
12000	3456.9\\
14000	3396.9\\
16000	3371.5\\
18000	3331.5\\
20000	3324.1\\
22000	3293.2\\
24000	3274.5\\
26000	3270.6\\
28000	3236.9\\
30000	3281.5\\
34000	3246.9\\
36000	3186\\
38000	3162.7\\
40000	3156.9\\
42000	3138.7\\
44000	3158.5\\
46000	3169.5\\
48000	3191.5\\
52000	3134.1\\
54000	3154.6\\
56000	3134\\
58000	3158.6\\
60000	3152.2\\
62000	3178.5\\
64000	3173.6\\
66000	3139.7\\
68000	3135.8\\
72000	3083.7\\
74000	3091.4\\
76000	3120.9\\
78000	3072.9\\
80000	3089.3\\
82000	3132.2\\
84000	3124.4\\
86000	3106.5\\
90000	3044.4\\
94000	3108.9\\
1e+05	3132.5\\
};
\addplot [color=mycolor39, forget plot]
  table[row sep=crcr]{%
0	4395.3\\
2000	5757.9\\
4000	7764.1\\
6000	9187.2\\
8000	10188\\
10000	10901\\
12000	11437\\
14000	11878\\
16000	12156\\
18000	12316\\
22000	12520\\
24000	12518\\
26000	12575\\
28000	12565\\
30000	12443\\
36000	12426\\
38000	12455\\
40000	12576\\
42000	12616\\
44000	12565\\
46000	12454\\
48000	12427\\
50000	12444\\
54000	12330\\
58000	12244\\
60000	12180\\
62000	12230\\
66000	12253\\
68000	12226\\
72000	12079\\
78000	11910\\
84000	11917\\
86000	11802\\
88000	11790\\
92000	11890\\
96000	11856\\
1e+05	11839\\
};
\addplot [color=mycolor40, forget plot]
  table[row sep=crcr]{%
0	4395.3\\
2000	4491.1\\
4000	4637.5\\
8000	4887.5\\
10000	4944.9\\
14000	5113.8\\
22000	5275.2\\
26000	5408.4\\
32000	5477.9\\
40000	5552.6\\
42000	5568.4\\
46000	5626.8\\
48000	5632.6\\
50000	5670.8\\
56000	5669.3\\
58000	5695.9\\
60000	5703.5\\
62000	5691.4\\
64000	5645\\
68000	5672.3\\
76000	5684.1\\
78000	5669.6\\
80000	5681.1\\
82000	5661.8\\
84000	5669.9\\
88000	5711.7\\
94000	5688.2\\
96000	5740.5\\
1e+05	5739.2\\
};
\addplot [color=mycolor41, forget plot]
  table[row sep=crcr]{%
0	4395.3\\
2000	4372.5\\
4000	4424.4\\
6000	4462.8\\
8000	4532\\
10000	4584.8\\
12000	4660.6\\
16000	4843.3\\
20000	4999.2\\
24000	5034.1\\
28000	5069.3\\
30000	5102.9\\
32000	5169.8\\
34000	5196.3\\
36000	5256\\
38000	5263.7\\
42000	5249.1\\
44000	5250.3\\
46000	5230.4\\
48000	5228.6\\
50000	5205.8\\
52000	5243.9\\
54000	5235\\
58000	5278.3\\
60000	5288.9\\
62000	5337.6\\
70000	5306.3\\
72000	5361.1\\
74000	5374\\
76000	5357.5\\
78000	5326.2\\
80000	5352\\
82000	5356\\
84000	5392.6\\
86000	5383.3\\
88000	5341.7\\
90000	5341.9\\
92000	5405.1\\
96000	5416\\
98000	5417.7\\
1e+05	5437.1\\
};
\addplot [color=mycolor42, forget plot]
  table[row sep=crcr]{%
0	4395.3\\
2000	5771.9\\
4000	8002.1\\
6000	9723\\
8000	11225\\
10000	12429\\
12000	13337\\
14000	14192\\
16000	14864\\
18000	15493\\
20000	16026\\
24000	16989\\
26000	17349\\
30000	17964\\
38000	18571\\
40000	18669\\
42000	18887\\
46000	18990\\
52000	19088\\
54000	19212\\
56000	19271\\
58000	19492\\
60000	19439\\
62000	19502\\
66000	19903\\
72000	19595\\
74000	19670\\
76000	19665\\
78000	19787\\
80000	19687\\
86000	19823\\
88000	19748\\
94000	19713\\
96000	19808\\
98000	19838\\
1e+05	19808\\
};
\addplot [color=mycolor43, forget plot]
  table[row sep=crcr]{%
0	4395.3\\
2000	4659.8\\
4000	5043.3\\
6000	5295.4\\
8000	5458.3\\
10000	5544.6\\
12000	5584.4\\
16000	5613.3\\
18000	5587.6\\
20000	5529.8\\
24000	5483.4\\
28000	5389.6\\
32000	5370.8\\
40000	5229.6\\
42000	5209.8\\
46000	5221.2\\
56000	5115.8\\
60000	5129.1\\
64000	5077.1\\
66000	5032.9\\
68000	5047.9\\
70000	5082\\
72000	5080.7\\
76000	5111.7\\
78000	5074.3\\
88000	5077.5\\
90000	5077.9\\
92000	5063.8\\
98000	5085.4\\
1e+05	5069.6\\
};
\addplot [color=mycolor44, forget plot]
  table[row sep=crcr]{%
0	4395.3\\
4000	4648.3\\
6000	4646.9\\
12000	4377.2\\
14000	4261.9\\
20000	4070.7\\
22000	4049.3\\
24000	3974.9\\
26000	3922.1\\
28000	3885.3\\
32000	3740.4\\
34000	3680.5\\
36000	3669\\
40000	3687\\
42000	3656.3\\
46000	3626.6\\
48000	3655.6\\
50000	3674.1\\
52000	3618.3\\
54000	3617.4\\
56000	3651.2\\
58000	3618.4\\
66000	3561.1\\
68000	3562.3\\
72000	3537.6\\
76000	3482.5\\
78000	3471.1\\
80000	3438\\
82000	3471.4\\
84000	3483.2\\
86000	3453.2\\
88000	3387.7\\
90000	3425.8\\
92000	3406\\
96000	3436.3\\
98000	3397.6\\
1e+05	3388.1\\
};
\addplot [color=mycolor45, forget plot]
  table[row sep=crcr]{%
0	4395.3\\
2000	4524.1\\
6000	5085\\
8000	5341.4\\
10000	5545\\
12000	5736.1\\
14000	5894.2\\
18000	6116.5\\
20000	6262.7\\
22000	6342.1\\
26000	6391.5\\
30000	6500.8\\
32000	6512.1\\
36000	6618.5\\
42000	6676.2\\
44000	6723.1\\
46000	6734.3\\
50000	6688.4\\
54000	6706.7\\
56000	6702.6\\
60000	6775.3\\
62000	6789.2\\
64000	6825.2\\
68000	6822\\
70000	6862.4\\
74000	6828.5\\
76000	6782.2\\
78000	6786.7\\
82000	6827\\
90000	6872.2\\
92000	6922.2\\
94000	6911.6\\
96000	6932.1\\
98000	6973.7\\
1e+05	6937.3\\
};
\addplot [color=mycolor46, forget plot]
  table[row sep=crcr]{%
0	4395.3\\
2000	4269.3\\
4000	4108.3\\
8000	3965.8\\
12000	4020.9\\
14000	4037\\
16000	4022.2\\
18000	4022.8\\
20000	4047\\
22000	4053.5\\
26000	4136.4\\
28000	4086.1\\
30000	4056\\
32000	4041.3\\
34000	4048.9\\
36000	4002.9\\
38000	4001.1\\
44000	4089.8\\
50000	4109.1\\
54000	4054.8\\
56000	4048.9\\
58000	4074.4\\
60000	4086.6\\
62000	4037\\
68000	4085.4\\
70000	4055.6\\
76000	4081.3\\
80000	4070.8\\
82000	4095.8\\
92000	4117.8\\
96000	4170.2\\
1e+05	4129.4\\
};
\addplot [color=mycolor47, forget plot]
  table[row sep=crcr]{%
0	4395.3\\
2000	4593.2\\
4000	4911.1\\
6000	5175.7\\
8000	5371.8\\
10000	5482.8\\
12000	5538.3\\
14000	5535.6\\
16000	5509.6\\
20000	5522.3\\
22000	5562.5\\
26000	5511.4\\
28000	5452\\
30000	5439.1\\
32000	5401.5\\
36000	5291.1\\
38000	5324.7\\
42000	5264.9\\
52000	5289.5\\
54000	5268.4\\
56000	5231.9\\
58000	5237.6\\
60000	5262\\
62000	5223.5\\
64000	5263.1\\
66000	5337.1\\
68000	5325.7\\
70000	5287.3\\
72000	5271.1\\
74000	5289.7\\
78000	5249.4\\
80000	5257\\
82000	5297.2\\
86000	5321.1\\
92000	5282.8\\
96000	5219.7\\
1e+05	5308.3\\
};
\addplot [color=mycolor48, forget plot]
  table[row sep=crcr]{%
0	4395.3\\
2000	4203.7\\
4000	3944.6\\
6000	3800.7\\
8000	3690.3\\
10000	3566.7\\
12000	3490.4\\
14000	3439\\
18000	3351.7\\
26000	3183.8\\
30000	3134.8\\
34000	3121.1\\
38000	3077.8\\
40000	3075.9\\
42000	3050.8\\
44000	3044.9\\
46000	3025.4\\
48000	3016.5\\
52000	3019.4\\
54000	3012.3\\
56000	3023.4\\
58000	3023.4\\
60000	3001.5\\
68000	2988.2\\
70000	2993.2\\
72000	2988.6\\
74000	2997.2\\
82000	2990.6\\
84000	3009.6\\
88000	2982.7\\
90000	2994.5\\
96000	2978.1\\
98000	2967.7\\
1e+05	2974.4\\
};
\addplot [color=mycolor49, forget plot]
  table[row sep=crcr]{%
0	4395.3\\
2000	4599.3\\
4000	4833.1\\
8000	5032\\
10000	5059.1\\
12000	4949.5\\
14000	4895.7\\
18000	4755.8\\
20000	4746.7\\
22000	4697.7\\
26000	4632.4\\
28000	4566.5\\
30000	4567.4\\
32000	4627.3\\
34000	4603.6\\
36000	4556.3\\
42000	4488.2\\
46000	4466.3\\
48000	4475.4\\
50000	4425.8\\
52000	4458.2\\
54000	4455.6\\
58000	4513.7\\
62000	4445.3\\
64000	4453.9\\
66000	4417.6\\
68000	4439.9\\
70000	4479.6\\
72000	4499.6\\
76000	4427.9\\
80000	4415.7\\
82000	4397.3\\
84000	4407.8\\
88000	4363.4\\
90000	4436.4\\
92000	4444\\
94000	4466.7\\
96000	4414.6\\
98000	4390.2\\
1e+05	4398\\
};
\addplot [color=mycolor50, forget plot]
  table[row sep=crcr]{%
0	4395.3\\
2000	3901.3\\
4000	3297.3\\
6000	2955\\
8000	2772.4\\
10000	2664.8\\
12000	2572\\
14000	2497.5\\
16000	2445.4\\
18000	2433.2\\
20000	2442.7\\
22000	2468.5\\
24000	2446.5\\
26000	2435.2\\
30000	2462\\
32000	2460.4\\
40000	2513.5\\
42000	2511.1\\
44000	2529.7\\
46000	2503\\
48000	2503.7\\
50000	2495.5\\
58000	2489.2\\
60000	2526.6\\
62000	2535.9\\
68000	2533.6\\
70000	2524.2\\
72000	2506.9\\
76000	2508.6\\
80000	2550.7\\
88000	2521.1\\
90000	2526.2\\
92000	2546.7\\
94000	2553.6\\
96000	2552.3\\
98000	2528.9\\
1e+05	2548.6\\
};
\addplot [color=mycolor51, forget plot]
  table[row sep=crcr]{%
0	4395.3\\
2000	4629.9\\
4000	5063.3\\
6000	5423.9\\
8000	5694.5\\
10000	5799.6\\
12000	5927.2\\
14000	6011.6\\
16000	6020.8\\
18000	6103.9\\
20000	6153.2\\
22000	6111.8\\
24000	6095\\
26000	6117.4\\
28000	6115.3\\
30000	6019.7\\
32000	5979.3\\
36000	5962.5\\
38000	6002.8\\
42000	6016.7\\
46000	5947.1\\
48000	5872.4\\
50000	5871\\
52000	5910\\
54000	5905.5\\
56000	5829.2\\
58000	5833.6\\
60000	5854.8\\
62000	5931.3\\
64000	5978.8\\
66000	5956.3\\
68000	5961.3\\
70000	5898.6\\
74000	5887.2\\
76000	5944.9\\
78000	6022.2\\
80000	6044.7\\
84000	6002\\
88000	6044.4\\
90000	5949.4\\
92000	5887.9\\
96000	6002.3\\
98000	6041\\
1e+05	6009.2\\
};
\addplot [color=mycolor52, forget plot]
  table[row sep=crcr]{%
0	4395.3\\
2000	4422.3\\
4000	4499.3\\
6000	4496.6\\
8000	4463.1\\
14000	4275.7\\
16000	4199.7\\
20000	4138.4\\
22000	4134.9\\
26000	4090.3\\
28000	4015.4\\
30000	3991\\
34000	3892.5\\
38000	3915.5\\
40000	3855.9\\
44000	3832.9\\
46000	3783.8\\
50000	3755.6\\
52000	3764.6\\
54000	3801.3\\
56000	3820.1\\
60000	3832.9\\
62000	3795.8\\
64000	3777.1\\
66000	3702.9\\
68000	3711.8\\
70000	3737.9\\
72000	3716.1\\
74000	3728.7\\
78000	3658.2\\
80000	3636\\
82000	3636.2\\
84000	3624\\
86000	3646\\
88000	3642.1\\
92000	3586.2\\
94000	3597.6\\
96000	3550.2\\
98000	3563.4\\
1e+05	3621.1\\
};
\addplot [color=mycolor53, forget plot]
  table[row sep=crcr]{%
0	4395.3\\
2000	4802\\
4000	5402.4\\
6000	5825.5\\
8000	6123.8\\
10000	6285.3\\
12000	6406.5\\
14000	6455.5\\
16000	6411\\
18000	6431.3\\
20000	6399.2\\
22000	6404.2\\
28000	6319.5\\
32000	6348.9\\
34000	6269.4\\
36000	6214\\
38000	6235.3\\
42000	6180.2\\
44000	6219.4\\
48000	6159.9\\
50000	6155.4\\
52000	6102.8\\
54000	6067.6\\
62000	5991.2\\
66000	5914.2\\
68000	5932.2\\
72000	5883.6\\
80000	5946.2\\
84000	5883\\
86000	5942\\
88000	5963.3\\
90000	5957.8\\
92000	5993.1\\
94000	5971.4\\
96000	5910.7\\
98000	5900.5\\
1e+05	5851.1\\
};
\addplot [color=mycolor54, forget plot]
  table[row sep=crcr]{%
0	4395.3\\
2000	4441.2\\
4000	4552.3\\
6000	4609.9\\
8000	4696.4\\
10000	4729.4\\
12000	4803.4\\
14000	4850.4\\
18000	4879.9\\
20000	4884.8\\
22000	4864.5\\
24000	4871\\
26000	4893.6\\
30000	4841\\
32000	4865.4\\
36000	4838.2\\
38000	4838\\
42000	4788.2\\
44000	4739.2\\
46000	4729.1\\
50000	4660.9\\
52000	4641.9\\
54000	4652.2\\
58000	4641.5\\
62000	4684.5\\
64000	4680.1\\
66000	4636.5\\
68000	4639.9\\
70000	4617.8\\
72000	4621\\
74000	4602\\
78000	4602.1\\
82000	4606.2\\
90000	4633.1\\
92000	4612.6\\
94000	4626\\
98000	4613.5\\
1e+05	4647.5\\
};
\addplot [color=mycolor55, forget plot]
  table[row sep=crcr]{%
0	4395.3\\
2000	4468.6\\
4000	4616.3\\
6000	4784.3\\
8000	4937.6\\
10000	5050.3\\
14000	5209.9\\
18000	5349.1\\
24000	5483.6\\
26000	5494.5\\
30000	5551.5\\
34000	5576.6\\
38000	5544.3\\
42000	5642.6\\
44000	5616.6\\
48000	5661.1\\
50000	5631.5\\
62000	5678.1\\
66000	5733.8\\
68000	5694.5\\
70000	5693.8\\
74000	5728.8\\
76000	5704.4\\
80000	5683.7\\
82000	5711.1\\
84000	5712.6\\
86000	5742.6\\
88000	5745\\
90000	5792.7\\
94000	5782.5\\
96000	5798.4\\
98000	5766.7\\
1e+05	5775.8\\
};
\addplot [color=mycolor56, forget plot]
  table[row sep=crcr]{%
0	4395.3\\
2000	4385.1\\
6000	4465.7\\
8000	4513.8\\
12000	4516.9\\
14000	4494.6\\
16000	4457.4\\
18000	4433\\
20000	4432.3\\
28000	4353.3\\
30000	4300.8\\
32000	4280\\
34000	4292.1\\
36000	4241\\
40000	4176.9\\
42000	4169.9\\
44000	4145.5\\
46000	4142.7\\
48000	4119.3\\
66000	4082.7\\
70000	4058.5\\
72000	4076.4\\
76000	4061.8\\
78000	4045.3\\
84000	4068.7\\
86000	4068.6\\
88000	4080.2\\
92000	4066.1\\
94000	4076.7\\
98000	4050.2\\
1e+05	4058.9\\
};
\addplot [color=mycolor57, forget plot]
  table[row sep=crcr]{%
0	4395.3\\
2000	4592\\
4000	4956.5\\
6000	5302.6\\
8000	5587.5\\
10000	5814.6\\
14000	6178.3\\
16000	6331.8\\
18000	6461.6\\
20000	6536.4\\
22000	6572.8\\
24000	6650.2\\
26000	6669.4\\
28000	6716.7\\
30000	6744.4\\
32000	6752.5\\
36000	6830.4\\
38000	6785.5\\
44000	6835.2\\
48000	6873.5\\
54000	6836.6\\
56000	6866.8\\
58000	6842.2\\
64000	6847.6\\
68000	6797.2\\
70000	6822.2\\
72000	6795.1\\
76000	6776.9\\
78000	6756.9\\
80000	6765.2\\
84000	6736.1\\
86000	6746.4\\
88000	6732.7\\
92000	6740.8\\
96000	6736.5\\
98000	6717.6\\
1e+05	6722.5\\
};
\addplot [color=mycolor58, forget plot]
  table[row sep=crcr]{%
0	4395.3\\
2000	4354.3\\
8000	4261.4\\
10000	4224.9\\
12000	4225.6\\
14000	4148.9\\
16000	4095.2\\
18000	4082.4\\
20000	4024.3\\
24000	3889.2\\
26000	3847\\
28000	3863.8\\
30000	3849.1\\
32000	3783.9\\
36000	3802.1\\
38000	3798.6\\
42000	3780.5\\
44000	3777.4\\
46000	3791.3\\
48000	3767\\
50000	3723\\
52000	3693.5\\
54000	3704.6\\
62000	3700.2\\
68000	3706\\
72000	3700.1\\
74000	3669.3\\
78000	3641.3\\
80000	3662.9\\
82000	3673.1\\
84000	3650.7\\
86000	3653.5\\
92000	3636.6\\
94000	3642.4\\
96000	3675.9\\
98000	3686.6\\
1e+05	3672.9\\
};
\addplot [color=mycolor59, forget plot]
  table[row sep=crcr]{%
0	4395.3\\
2000	4353.1\\
4000	4366.8\\
6000	4439\\
10000	4548.3\\
16000	4703.5\\
18000	4715.1\\
20000	4773.2\\
22000	4815.7\\
24000	4809.3\\
26000	4853.2\\
28000	4879.6\\
32000	4895.2\\
34000	4921.2\\
40000	4958.7\\
46000	4934.4\\
54000	4952.8\\
56000	4979.9\\
64000	4970.1\\
66000	4949\\
68000	4966.8\\
70000	4964.8\\
80000	5052.8\\
86000	5014\\
88000	5019.6\\
90000	5040.5\\
92000	5081.9\\
94000	5057.8\\
1e+05	5043.1\\
};
\addplot [color=mycolor60, forget plot]
  table[row sep=crcr]{%
0	4395.3\\
2000	4222.9\\
4000	4010.7\\
6000	3903\\
8000	3849.4\\
10000	3816\\
12000	3766.7\\
16000	3756.7\\
22000	3788.4\\
24000	3774.9\\
30000	3779.7\\
32000	3759.3\\
36000	3740.6\\
38000	3711.6\\
50000	3685\\
52000	3672.2\\
54000	3684.1\\
56000	3671.2\\
58000	3689.7\\
62000	3680.5\\
64000	3673.2\\
66000	3655.1\\
72000	3642.3\\
74000	3661.4\\
76000	3668.3\\
78000	3664.7\\
80000	3648.4\\
82000	3656\\
90000	3611.7\\
92000	3617.3\\
96000	3592.5\\
1e+05	3608\\
};
\addplot [color=mycolor61, forget plot]
  table[row sep=crcr]{%
0	4395.3\\
2000	4112.4\\
4000	3740.8\\
6000	3525.5\\
8000	3405\\
10000	3338.6\\
12000	3310.3\\
14000	3270.4\\
16000	3244.9\\
24000	3214.9\\
28000	3206.4\\
30000	3212.5\\
40000	3170.2\\
44000	3191.3\\
46000	3178\\
50000	3173.5\\
54000	3193.6\\
56000	3193.4\\
58000	3183\\
60000	3185.2\\
66000	3159.5\\
68000	3168.6\\
70000	3157.9\\
72000	3167.7\\
74000	3195.1\\
76000	3206.1\\
78000	3203.5\\
80000	3214\\
82000	3213.1\\
84000	3226.7\\
86000	3227.5\\
88000	3214.1\\
92000	3212.7\\
98000	3194.7\\
1e+05	3215.3\\
};
\addplot [color=mycolor62, forget plot]
  table[row sep=crcr]{%
0	4395.3\\
2000	4353.8\\
4000	4359.5\\
12000	4669.9\\
16000	4768.9\\
18000	4843.3\\
20000	4902\\
22000	4937.3\\
26000	4960.8\\
28000	4947.5\\
32000	4974.9\\
34000	4994.4\\
42000	5010.9\\
44000	5028.1\\
46000	5025.8\\
48000	4997.6\\
52000	5049.7\\
56000	5050.2\\
58000	5060.4\\
64000	5005.9\\
66000	5004.8\\
70000	5065.1\\
74000	5028.4\\
76000	5032.5\\
78000	5061.9\\
86000	5039.3\\
88000	5046\\
90000	5019.6\\
92000	5030.8\\
98000	5015.2\\
1e+05	5025.4\\
};
\addplot [color=mycolor63, forget plot]
  table[row sep=crcr]{%
0	4395.3\\
2000	4118.3\\
4000	3705.7\\
6000	3444.6\\
8000	3265.3\\
10000	3141.4\\
12000	3049.1\\
14000	2975.7\\
20000	2785.6\\
24000	2714.6\\
26000	2703.6\\
30000	2663.9\\
32000	2644.1\\
34000	2615.4\\
36000	2606.6\\
38000	2588.5\\
42000	2581.3\\
44000	2557.8\\
46000	2542.1\\
48000	2543.2\\
50000	2532.1\\
52000	2529.7\\
54000	2517.6\\
56000	2519.6\\
58000	2514\\
60000	2500.8\\
70000	2510.4\\
74000	2494.6\\
86000	2512.5\\
92000	2484.5\\
94000	2480.5\\
96000	2484.4\\
1e+05	2479\\
};
\addplot [color=mycolor64, forget plot]
  table[row sep=crcr]{%
0	4395.3\\
2000	4064.6\\
4000	3644.1\\
6000	3405.3\\
8000	3271.1\\
10000	3187.4\\
14000	3128.3\\
22000	3113.9\\
24000	3102.1\\
30000	3108.1\\
32000	3120.8\\
34000	3106.4\\
36000	3101.4\\
40000	3113.8\\
44000	3092\\
52000	3091.6\\
54000	3107.4\\
56000	3111\\
58000	3125.6\\
60000	3114.7\\
70000	3124.4\\
72000	3133.5\\
74000	3108.2\\
76000	3111\\
78000	3103.1\\
80000	3104.8\\
82000	3097.6\\
86000	3097.1\\
88000	3093.5\\
92000	3107.2\\
94000	3124\\
96000	3108.7\\
1e+05	3107.6\\
};
\addplot [color=mycolor65, forget plot]
  table[row sep=crcr]{%
0	4395.3\\
2000	4044.4\\
4000	3587.3\\
6000	3293.7\\
8000	3127.6\\
10000	2992\\
12000	2891.1\\
14000	2831.2\\
18000	2779.2\\
20000	2739.2\\
24000	2677.1\\
26000	2666.1\\
28000	2680.2\\
34000	2632.5\\
38000	2611.6\\
40000	2598.4\\
42000	2599.5\\
44000	2615.2\\
48000	2602.3\\
50000	2591.8\\
52000	2589.5\\
56000	2599.7\\
58000	2591.2\\
60000	2590.7\\
62000	2603.1\\
64000	2591.3\\
68000	2599\\
70000	2586.3\\
76000	2601.1\\
84000	2590.5\\
88000	2593.8\\
90000	2590.3\\
92000	2615.5\\
94000	2623.9\\
1e+05	2598.8\\
};
\addplot [color=mycolor66, forget plot]
  table[row sep=crcr]{%
0	4395.3\\
2000	4132.2\\
4000	3803.9\\
6000	3622.4\\
8000	3516.2\\
10000	3456.2\\
14000	3434.1\\
18000	3453.5\\
28000	3434.7\\
30000	3433.3\\
32000	3444.1\\
34000	3433.4\\
40000	3440.4\\
42000	3460.3\\
46000	3463.4\\
48000	3448.2\\
54000	3448.7\\
56000	3449.3\\
58000	3465.1\\
66000	3482.8\\
68000	3483.6\\
70000	3473.6\\
74000	3498.6\\
76000	3492.5\\
78000	3500.4\\
96000	3479.8\\
1e+05	3510.3\\
};
\addplot [color=mycolor67, forget plot]
  table[row sep=crcr]{%
0	4395.3\\
2000	4220.2\\
4000	3993\\
6000	3864.2\\
8000	3767.9\\
10000	3685\\
12000	3638.9\\
14000	3610.5\\
18000	3524\\
20000	3492\\
22000	3475\\
24000	3446.2\\
30000	3420.8\\
34000	3379.1\\
36000	3373.1\\
38000	3383.4\\
40000	3374.6\\
44000	3344.1\\
48000	3349.6\\
52000	3322.4\\
58000	3318.6\\
60000	3309.5\\
66000	3318.3\\
68000	3301.1\\
70000	3297.9\\
74000	3278.2\\
78000	3264.6\\
90000	3272.1\\
94000	3260.7\\
96000	3277.3\\
1e+05	3280.8\\
};
\addplot [color=mycolor67]
  table[row sep=crcr]{%
0	4395.3\\
2000	3998.6\\
4000	3470.6\\
6000	3157.1\\
8000	2971.6\\
10000	2860.8\\
12000	2790.2\\
14000	2738.8\\
16000	2713.8\\
18000	2697.1\\
22000	2686.3\\
30000	2685.2\\
38000	2703.4\\
40000	2686.6\\
42000	2680.5\\
44000	2683.1\\
46000	2694.9\\
60000	2702.8\\
66000	2688.1\\
68000	2694.5\\
76000	2686.1\\
78000	2692.5\\
80000	2685.3\\
88000	2691.9\\
90000	2703.4\\
1e+05	2696.7\\
};
\addlegendentry{Classe k=43}

\addplot [color=mycolor68, forget plot]
  table[row sep=crcr]{%
0	4395.3\\
2000	4106.3\\
4000	3697\\
6000	3425.9\\
8000	3241.6\\
10000	3119.2\\
12000	3033.5\\
14000	2962.6\\
16000	2909.2\\
20000	2834.7\\
28000	2735.4\\
32000	2718.5\\
38000	2664.7\\
44000	2628.1\\
48000	2620.6\\
54000	2608.5\\
58000	2604.4\\
66000	2575.5\\
76000	2574.4\\
78000	2571.8\\
82000	2576.4\\
88000	2558\\
1e+05	2551.7\\
};
\addplot [color=mycolor69, forget plot]
  table[row sep=crcr]{%
0	4395.3\\
2000	4150.6\\
4000	3796\\
6000	3558.5\\
8000	3409.1\\
10000	3299.7\\
12000	3222.8\\
18000	3067\\
20000	3035.1\\
24000	3004.3\\
38000	2930.7\\
44000	2933.8\\
46000	2934.3\\
48000	2950\\
50000	2954.7\\
52000	2944.9\\
60000	2941.1\\
64000	2927.1\\
68000	2942.1\\
70000	2935.1\\
72000	2951.2\\
1e+05	2973.7\\
};
\addplot [color=mycolor70, forget plot]
  table[row sep=crcr]{%
0	4395.3\\
2000	3982.4\\
4000	3387\\
6000	2989.9\\
8000	2724.5\\
10000	2541.7\\
12000	2408.6\\
14000	2324.8\\
16000	2258.8\\
20000	2174.4\\
22000	2157\\
24000	2155.4\\
26000	2139.2\\
28000	2144.4\\
34000	2125.4\\
36000	2122.8\\
38000	2127.3\\
40000	2140\\
42000	2128.8\\
46000	2141.8\\
48000	2152.7\\
50000	2153.1\\
54000	2141.6\\
60000	2149.8\\
72000	2156.6\\
74000	2167.2\\
78000	2165.3\\
80000	2178.1\\
82000	2178\\
84000	2167\\
86000	2167.1\\
88000	2160.6\\
92000	2172.1\\
94000	2171.9\\
96000	2164.1\\
1e+05	2173.6\\
};
\addplot [color=mycolor71, forget plot]
  table[row sep=crcr]{%
0	4395.3\\
2000	4041.7\\
4000	3565.7\\
6000	3272.6\\
8000	3081.4\\
10000	2963.3\\
12000	2882.2\\
14000	2824.1\\
16000	2784.7\\
20000	2725.6\\
22000	2703.5\\
30000	2654.2\\
32000	2662.2\\
34000	2662.6\\
38000	2651.6\\
40000	2651.8\\
44000	2639.4\\
46000	2641\\
52000	2626.8\\
56000	2623\\
62000	2603.4\\
70000	2606.9\\
72000	2596.4\\
78000	2606.9\\
80000	2601.3\\
88000	2610.6\\
92000	2603.6\\
96000	2608.2\\
98000	2609.5\\
1e+05	2602.3\\
};
\addplot [color=mycolor72, forget plot]
  table[row sep=crcr]{%
0	4395.3\\
2000	4080.1\\
4000	3615.9\\
6000	3304.3\\
8000	3090.8\\
10000	2923.7\\
12000	2795.4\\
14000	2693.8\\
16000	2625.5\\
22000	2476.9\\
24000	2434.1\\
26000	2405.6\\
28000	2387.2\\
30000	2358.5\\
34000	2330.6\\
36000	2320.2\\
38000	2297.8\\
40000	2287.3\\
42000	2291.6\\
44000	2289.1\\
46000	2302.1\\
50000	2304.4\\
58000	2291.2\\
62000	2286.6\\
66000	2295.1\\
74000	2285.6\\
76000	2288.2\\
78000	2280.7\\
86000	2296.6\\
88000	2276.4\\
96000	2274.5\\
98000	2262.8\\
1e+05	2266.2\\
};
\addplot [color=mycolor73, forget plot]
  table[row sep=crcr]{%
0	4395.3\\
2000	4072.4\\
4000	3624.3\\
6000	3341.3\\
8000	3159.6\\
10000	3042.4\\
12000	2962.5\\
14000	2905.9\\
16000	2863\\
18000	2830.9\\
22000	2806.2\\
26000	2777.7\\
28000	2764.7\\
32000	2772.3\\
34000	2766.4\\
38000	2778.8\\
42000	2751.4\\
44000	2745.1\\
46000	2749.2\\
50000	2733.8\\
52000	2735.4\\
58000	2714\\
66000	2706.8\\
72000	2713.9\\
82000	2699.6\\
86000	2704.8\\
88000	2696.1\\
1e+05	2703.5\\
};
\addplot [color=mycolor74, forget plot]
  table[row sep=crcr]{%
0	4395.3\\
2000	3960.5\\
4000	3340.5\\
6000	2905.4\\
8000	2602.3\\
10000	2403.1\\
12000	2246.7\\
14000	2132.4\\
16000	2038.9\\
18000	1961.4\\
22000	1854.3\\
24000	1813.8\\
26000	1779.3\\
28000	1760\\
30000	1732.2\\
32000	1714.6\\
52000	1616.9\\
54000	1610.9\\
56000	1616.2\\
62000	1611.9\\
64000	1608.7\\
66000	1614.9\\
68000	1614.8\\
70000	1620.4\\
72000	1618\\
74000	1620.6\\
78000	1616.6\\
82000	1623.3\\
84000	1614.6\\
94000	1616.1\\
96000	1610.2\\
1e+05	1610.6\\
};
\addplot [color=mycolor75, forget plot]
  table[row sep=crcr]{%
0	4395.3\\
2000	4075.6\\
4000	3616.5\\
6000	3297.9\\
8000	3064.4\\
10000	2897.3\\
12000	2771.7\\
14000	2675.2\\
16000	2598.7\\
18000	2541.9\\
20000	2501.8\\
28000	2385.3\\
32000	2349.4\\
36000	2341.9\\
40000	2323\\
60000	2301.4\\
64000	2307.2\\
68000	2296.7\\
72000	2299.9\\
76000	2287.6\\
84000	2294.1\\
88000	2275.4\\
96000	2283.4\\
98000	2279.2\\
1e+05	2282.2\\
};
\addplot [color=mycolor76, forget plot]
  table[row sep=crcr]{%
0	4395.3\\
2000	4012.3\\
4000	3469.9\\
6000	3108.6\\
8000	2859.6\\
10000	2693.7\\
12000	2570.2\\
14000	2482\\
16000	2411.5\\
20000	2321.3\\
22000	2287.7\\
24000	2266.3\\
26000	2253.9\\
30000	2210.3\\
32000	2190.2\\
44000	2142.9\\
48000	2141.1\\
58000	2135.7\\
60000	2123.2\\
66000	2117.7\\
68000	2123.5\\
72000	2114\\
78000	2117.2\\
82000	2099.6\\
88000	2105.4\\
90000	2109.1\\
1e+05	2097.6\\
};
\addplot [color=mycolor77, forget plot]
  table[row sep=crcr]{%
0	4395.3\\
2000	3982.9\\
4000	3380.4\\
6000	2960.1\\
8000	2672.8\\
10000	2454.9\\
12000	2299\\
14000	2191.4\\
16000	2101.9\\
20000	1972.8\\
24000	1893.6\\
26000	1864.4\\
30000	1830.5\\
34000	1793.2\\
36000	1773\\
40000	1754.7\\
44000	1742.4\\
50000	1733.1\\
54000	1709.5\\
58000	1704.3\\
62000	1691.4\\
66000	1673.6\\
72000	1667.1\\
76000	1661.2\\
82000	1658.2\\
88000	1667.7\\
1e+05	1677.4\\
};
\addplot [color=mycolor78, forget plot]
  table[row sep=crcr]{%
0	4395.3\\
2000	4108.1\\
4000	3671.4\\
6000	3363.6\\
8000	3146.3\\
10000	3000\\
12000	2890.9\\
16000	2743.3\\
18000	2689.6\\
26000	2573.7\\
32000	2519.7\\
36000	2503.1\\
40000	2489.7\\
44000	2474.5\\
46000	2462.6\\
50000	2459.2\\
52000	2461.8\\
64000	2428.9\\
66000	2428.7\\
70000	2406\\
80000	2412.8\\
82000	2409.4\\
84000	2397.2\\
88000	2405.7\\
90000	2406\\
94000	2386.5\\
1e+05	2396.3\\
};
\addplot [color=mycolor79, forget plot]
  table[row sep=crcr]{%
0	4395.3\\
2000	4045.6\\
4000	3524.3\\
6000	3174.5\\
8000	2929.9\\
10000	2761.9\\
12000	2647.1\\
14000	2566.1\\
18000	2444\\
20000	2410.8\\
22000	2387\\
28000	2349.7\\
32000	2342.2\\
36000	2338.4\\
38000	2321.6\\
42000	2315.7\\
46000	2313.9\\
48000	2310.9\\
50000	2300.5\\
52000	2304.6\\
68000	2284.6\\
70000	2271.1\\
74000	2259.8\\
76000	2260.3\\
78000	2270.5\\
94000	2281.8\\
1e+05	2248\\
};
\addplot [color=mycolor79, forget plot]
  table[row sep=crcr]{%
0	4395.3\\
2000	4086.3\\
4000	3628.6\\
6000	3301.3\\
8000	3067.9\\
10000	2897.8\\
12000	2767.6\\
14000	2657.1\\
16000	2573\\
20000	2448.5\\
26000	2333.7\\
32000	2244.9\\
36000	2209.3\\
38000	2189.7\\
40000	2178.8\\
42000	2176.2\\
44000	2164.4\\
48000	2158.7\\
52000	2140\\
54000	2128.8\\
56000	2132.6\\
58000	2128.9\\
60000	2132.1\\
70000	2112.1\\
74000	2097.5\\
84000	2108.5\\
86000	2111.8\\
90000	2106\\
92000	2111.5\\
98000	2109.6\\
1e+05	2098.5\\
};
\addplot [color=mycolor80, forget plot]
  table[row sep=crcr]{%
0	4395.3\\
2000	3979.9\\
4000	3362.7\\
6000	2925.7\\
8000	2610.1\\
10000	2390.1\\
12000	2224.2\\
14000	2095.5\\
16000	2000.8\\
18000	1931.8\\
20000	1875.2\\
22000	1830.9\\
30000	1729.8\\
36000	1689\\
46000	1668.2\\
48000	1658.3\\
56000	1644.6\\
60000	1647.3\\
62000	1653.5\\
72000	1650.4\\
76000	1643.9\\
78000	1646.5\\
82000	1643.3\\
84000	1642.3\\
86000	1647.1\\
1e+05	1641.1\\
};
\addplot [color=mycolor81, forget plot]
  table[row sep=crcr]{%
0	4395.3\\
2000	3988.4\\
4000	3376.2\\
6000	2945.2\\
8000	2627.8\\
10000	2400\\
12000	2220.8\\
14000	2085.8\\
16000	1980\\
18000	1894.8\\
20000	1832.3\\
22000	1783.2\\
26000	1707.8\\
28000	1677\\
32000	1635.8\\
36000	1608.4\\
38000	1597.4\\
42000	1586.1\\
58000	1544.1\\
64000	1541.9\\
68000	1553\\
76000	1539.3\\
94000	1517.6\\
98000	1521.2\\
1e+05	1517.6\\
};
\addplot [color=mycolor82, forget plot]
  table[row sep=crcr]{%
0	4395.3\\
2000	4034.5\\
4000	3490.1\\
6000	3103\\
8000	2824.2\\
10000	2616.8\\
12000	2459.2\\
14000	2338\\
16000	2252.3\\
18000	2180\\
22000	2073.8\\
24000	2028.8\\
28000	1965.7\\
30000	1945.5\\
44000	1858\\
46000	1860.2\\
54000	1838.3\\
56000	1822.9\\
60000	1816.6\\
64000	1799.2\\
72000	1790.8\\
80000	1790.9\\
84000	1795.1\\
96000	1783.8\\
1e+05	1790.8\\
};
\addplot [color=mycolor83, forget plot]
  table[row sep=crcr]{%
0	4395.3\\
2000	4052.6\\
4000	3540\\
6000	3170.6\\
8000	2891.2\\
10000	2678.6\\
12000	2517.9\\
14000	2388.5\\
16000	2285.6\\
18000	2212.2\\
20000	2148.7\\
22000	2099.3\\
26000	2026.6\\
30000	1973.7\\
38000	1913\\
48000	1883\\
50000	1884.6\\
62000	1856.4\\
66000	1854\\
68000	1856\\
76000	1838.4\\
94000	1820\\
1e+05	1810.7\\
};
\addplot [color=mycolor83, forget plot]
  table[row sep=crcr]{%
0	4395.3\\
2000	3949.2\\
4000	3258\\
6000	2760.9\\
8000	2410\\
10000	2153.9\\
12000	1967.2\\
14000	1824.2\\
16000	1718.2\\
18000	1642.5\\
20000	1579.5\\
24000	1491.7\\
26000	1460.2\\
30000	1416.8\\
38000	1371.6\\
40000	1363.8\\
48000	1353.7\\
60000	1339.3\\
62000	1334.6\\
66000	1335.4\\
72000	1341.1\\
74000	1336.9\\
86000	1335.3\\
88000	1331.6\\
94000	1331.8\\
98000	1330.7\\
1e+05	1323.2\\
};
\addplot [color=mycolor84, forget plot]
  table[row sep=crcr]{%
0	4395.3\\
2000	3980.1\\
4000	3343.3\\
6000	2872.8\\
8000	2530.8\\
10000	2283.5\\
12000	2095\\
14000	1944.5\\
16000	1824.6\\
18000	1736\\
20000	1669.7\\
22000	1612.2\\
26000	1526.3\\
28000	1499.2\\
32000	1459.7\\
36000	1427.9\\
38000	1420.7\\
40000	1405.2\\
50000	1372\\
52000	1373.1\\
54000	1368.4\\
60000	1367.9\\
62000	1359.8\\
66000	1360.9\\
68000	1352.1\\
74000	1351.7\\
76000	1355.4\\
78000	1351\\
90000	1351.5\\
92000	1346.4\\
96000	1353.9\\
1e+05	1347\\
};
\addplot [color=mycolor85, forget plot]
  table[row sep=crcr]{%
0	4395.3\\
2000	4033.2\\
4000	3468.5\\
6000	3055\\
8000	2736\\
10000	2490.9\\
12000	2310.9\\
14000	2167.2\\
16000	2050.8\\
18000	1954.5\\
20000	1879.1\\
22000	1816.9\\
28000	1685.8\\
30000	1649.5\\
34000	1606.2\\
38000	1575.3\\
42000	1539.6\\
44000	1528.1\\
50000	1521.6\\
52000	1512.5\\
54000	1509\\
58000	1494.7\\
60000	1490.2\\
62000	1495.7\\
66000	1487\\
70000	1489.3\\
72000	1486.1\\
74000	1474.3\\
76000	1467.9\\
80000	1461.8\\
82000	1457.4\\
84000	1462.8\\
88000	1458.8\\
94000	1469.9\\
96000	1470.7\\
98000	1463.7\\
1e+05	1461.6\\
};
\addplot [color=mycolor86, forget plot]
  table[row sep=crcr]{%
0	4395.3\\
2000	3994.2\\
4000	3367.4\\
6000	2911.2\\
8000	2565.8\\
10000	2301.3\\
12000	2103.1\\
14000	1946.1\\
16000	1822.9\\
18000	1727.3\\
20000	1658.7\\
24000	1551.7\\
26000	1507.1\\
28000	1471.5\\
30000	1441.9\\
32000	1421.3\\
38000	1371.4\\
40000	1368.5\\
46000	1344.5\\
50000	1324.9\\
54000	1321.4\\
56000	1315.7\\
66000	1308.8\\
68000	1300.6\\
76000	1289.1\\
78000	1288.9\\
80000	1282.5\\
82000	1282.3\\
88000	1270.7\\
90000	1267\\
94000	1264\\
96000	1261.8\\
98000	1264.1\\
1e+05	1262.8\\
};
\addplot [color=mycolor87, forget plot]
  table[row sep=crcr]{%
0	4395.3\\
2000	3976.2\\
4000	3327\\
6000	2823.2\\
8000	2441\\
10000	2137.7\\
12000	1906.6\\
14000	1724\\
16000	1566.5\\
18000	1446.1\\
20000	1357.8\\
22000	1287\\
24000	1227.4\\
26000	1180.1\\
28000	1143.9\\
30000	1115.7\\
32000	1100\\
34000	1081.5\\
36000	1068\\
38000	1063\\
40000	1050.6\\
42000	1044.7\\
44000	1044.7\\
46000	1041.3\\
48000	1034.9\\
52000	1039.5\\
54000	1034.4\\
56000	1022.7\\
58000	1022.6\\
64000	1046.9\\
66000	1042.7\\
68000	1043.3\\
70000	1051.2\\
72000	1050.7\\
74000	1057.7\\
80000	1043.6\\
82000	1044.8\\
88000	1036.4\\
92000	1037.9\\
96000	1032.5\\
98000	1034.6\\
1e+05	1042\\
};
\addplot [color=mycolor88, forget plot]
  table[row sep=crcr]{%
0	4395.3\\
2000	4029.4\\
4000	3449.9\\
6000	2998.5\\
8000	2653.9\\
10000	2384.3\\
12000	2174.4\\
14000	2007\\
16000	1877.4\\
18000	1771.8\\
20000	1689.9\\
22000	1623.4\\
26000	1518.3\\
28000	1474.1\\
32000	1420.5\\
34000	1396.3\\
36000	1378\\
40000	1353.6\\
50000	1323.8\\
54000	1320.4\\
56000	1314.9\\
60000	1311.5\\
64000	1298.5\\
68000	1290\\
72000	1287.5\\
74000	1285.4\\
76000	1287.7\\
80000	1278.7\\
84000	1279.7\\
94000	1258.7\\
1e+05	1263.5\\
};
\addplot [color=mycolor89, forget plot]
  table[row sep=crcr]{%
0	4395.3\\
2000	4021.5\\
4000	3418.2\\
6000	2946.4\\
8000	2574.7\\
10000	2268.5\\
12000	2025.1\\
14000	1832.6\\
16000	1677.8\\
18000	1555.8\\
20000	1455\\
22000	1368.7\\
24000	1300.3\\
26000	1242.4\\
30000	1143.2\\
32000	1107.6\\
36000	1055.3\\
40000	1016\\
42000	1001.7\\
50000	959.85\\
52000	952.53\\
56000	949.57\\
66000	926.84\\
70000	921.85\\
74000	914.15\\
80000	918.78\\
84000	914.82\\
92000	915.96\\
96000	907.38\\
98000	909.04\\
1e+05	907.8\\
};
\addplot [color=mycolor90, forget plot]
  table[row sep=crcr]{%
0	4395.3\\
2000	3957.4\\
4000	3250.8\\
6000	2724.8\\
8000	2325.5\\
10000	2015\\
12000	1779.1\\
14000	1594.7\\
16000	1450\\
18000	1341.6\\
20000	1254.3\\
22000	1179.8\\
24000	1124.7\\
26000	1075.2\\
28000	1038.6\\
32000	990.73\\
34000	970.97\\
38000	944.15\\
40000	935.44\\
44000	912.84\\
48000	900.58\\
54000	887.35\\
58000	887.22\\
62000	881.58\\
64000	878.89\\
68000	880.01\\
70000	882.57\\
72000	881.61\\
74000	885.82\\
78000	885.62\\
84000	889.82\\
88000	882.39\\
90000	886.85\\
94000	889.51\\
98000	890.22\\
1e+05	891.59\\
};
\addplot [color=mycolor91, forget plot]
  table[row sep=crcr]{%
0	4395.3\\
2000	4017.6\\
4000	3436.8\\
6000	2957.8\\
8000	2574.8\\
10000	2269.4\\
12000	2015.8\\
14000	1807.4\\
16000	1637.1\\
18000	1497.7\\
20000	1375.8\\
22000	1279.6\\
24000	1198.9\\
26000	1138.4\\
28000	1084\\
30000	1035.6\\
32000	994.83\\
34000	958.59\\
44000	851.08\\
46000	834.75\\
52000	803.87\\
54000	797.67\\
56000	796.05\\
58000	790.85\\
62000	788.8\\
64000	790.79\\
70000	781.36\\
74000	782.21\\
76000	775.01\\
80000	773.35\\
82000	774.21\\
84000	770.97\\
90000	769.06\\
92000	772.46\\
96000	772.75\\
98000	767.33\\
1e+05	767.68\\
};
\addplot [color=mycolor92, forget plot]
  table[row sep=crcr]{%
0	4395.3\\
2000	4011.4\\
4000	3400.3\\
6000	2924.7\\
8000	2551.6\\
10000	2255.7\\
12000	2017.5\\
14000	1828.3\\
16000	1683.8\\
18000	1558.6\\
20000	1465.2\\
22000	1388.1\\
24000	1324.6\\
26000	1268.4\\
30000	1180.1\\
34000	1131.7\\
38000	1084.4\\
40000	1073.2\\
42000	1059\\
46000	1036.1\\
48000	1021\\
50000	1011.6\\
56000	999.41\\
58000	999.89\\
62000	990.83\\
64000	986.56\\
66000	990.14\\
68000	989.36\\
70000	985.66\\
74000	983.25\\
76000	980.86\\
78000	984.11\\
86000	974.99\\
92000	958.6\\
94000	954.66\\
96000	955.75\\
98000	959.61\\
1e+05	960.02\\
};
\addplot [color=mycolor93, forget plot]
  table[row sep=crcr]{%
0	4395.3\\
2000	4059.8\\
4000	3511.7\\
6000	3048\\
8000	2668.9\\
10000	2370.5\\
12000	2124.8\\
14000	1923.8\\
16000	1756.7\\
20000	1497.5\\
22000	1382.2\\
24000	1299.3\\
26000	1228.3\\
28000	1164.5\\
30000	1116.7\\
34000	1048.5\\
36000	1011.7\\
38000	979.53\\
44000	910.32\\
46000	897.78\\
48000	877.32\\
52000	860.17\\
56000	840.05\\
68000	827.04\\
72000	813.17\\
74000	814.88\\
78000	805.9\\
80000	805.3\\
82000	801.42\\
84000	800.68\\
86000	803.37\\
88000	798.12\\
94000	791.09\\
1e+05	804.12\\
};
\addplot [color=mycolor94, forget plot]
  table[row sep=crcr]{%
0	4395.3\\
2000	4067\\
4000	3508.1\\
6000	3062.1\\
8000	2697\\
10000	2387.2\\
12000	2133.1\\
14000	1920.6\\
16000	1751\\
18000	1604.1\\
20000	1478.5\\
22000	1376.1\\
24000	1292.3\\
26000	1217\\
28000	1154.4\\
32000	1052.2\\
34000	1011.3\\
36000	980.34\\
42000	903.93\\
44000	882.77\\
48000	851.37\\
54000	824.22\\
56000	815.28\\
60000	803.43\\
70000	781.25\\
72000	782.42\\
78000	772.05\\
88000	762.42\\
92000	759.58\\
96000	762.99\\
98000	761.24\\
1e+05	762.07\\
};
\addplot [color=mycolor95, forget plot]
  table[row sep=crcr]{%
0	4395.3\\
2000	3998.1\\
4000	3351\\
6000	2833.2\\
8000	2424.1\\
10000	2094.7\\
12000	1817.4\\
14000	1595.6\\
16000	1415.6\\
18000	1266.2\\
20000	1137\\
22000	1032.3\\
24000	945.75\\
26000	872.4\\
28000	811.17\\
30000	759.48\\
34000	673.66\\
38000	613.35\\
40000	590.19\\
42000	572.2\\
46000	540.73\\
48000	527.51\\
50000	516.57\\
52000	511.18\\
56000	497.88\\
64000	479.18\\
68000	472.48\\
72000	467.04\\
82000	463.23\\
84000	462.2\\
88000	465.29\\
92000	465.03\\
94000	462.13\\
96000	461.21\\
98000	457.53\\
1e+05	458.53\\
};
\addplot [color=mycolor96, forget plot]
  table[row sep=crcr]{%
0	4395.3\\
2000	4025.7\\
4000	3423.7\\
6000	2960.1\\
8000	2580.3\\
12000	1985.4\\
14000	1754.8\\
16000	1563.2\\
18000	1406.6\\
20000	1274.9\\
22000	1165.6\\
26000	984.43\\
28000	921.12\\
32000	816.33\\
34000	775.84\\
40000	674.86\\
42000	641.76\\
44000	621.5\\
46000	607.06\\
50000	583.08\\
52000	568.22\\
54000	555.3\\
58000	540.89\\
62000	527.49\\
64000	523.34\\
66000	523.65\\
68000	522.37\\
70000	517.88\\
72000	515.83\\
76000	513.74\\
80000	508.03\\
84000	505.58\\
88000	492.98\\
92000	489.08\\
94000	483.91\\
96000	483.35\\
98000	486.84\\
1e+05	484.62\\
};
\addplot [color=mycolor97, forget plot]
  table[row sep=crcr]{%
0	4395.3\\
2000	3944.3\\
4000	3240.5\\
6000	2685.1\\
8000	2263\\
10000	1915.9\\
12000	1641.7\\
14000	1424.7\\
16000	1246.5\\
18000	1099.5\\
20000	977.9\\
22000	878.92\\
24000	798.33\\
26000	731.72\\
28000	673.48\\
30000	627.82\\
32000	592.21\\
34000	554.42\\
36000	527.15\\
38000	503.4\\
40000	484.37\\
42000	464.04\\
44000	448.02\\
52000	406.06\\
54000	401.38\\
56000	399.18\\
60000	385.9\\
64000	382.32\\
66000	383.46\\
68000	382.78\\
70000	380.03\\
74000	371.45\\
76000	368.33\\
78000	363.44\\
80000	361.36\\
82000	363.27\\
84000	362.25\\
88000	364.83\\
92000	364.04\\
94000	363.05\\
98000	352.69\\
1e+05	350.22\\
};
\addplot [color=mycolor98]
  table[row sep=crcr]{%
0	4395.3\\
2000	4076.8\\
4000	3541.8\\
6000	3090\\
10000	2383.4\\
12000	2110.2\\
14000	1877.1\\
16000	1687\\
20000	1382.2\\
24000	1166\\
26000	1073.2\\
28000	991.35\\
30000	920.13\\
32000	863.24\\
36000	773.7\\
38000	732.71\\
40000	697.71\\
44000	643.29\\
46000	622.66\\
48000	605.21\\
52000	569.57\\
58000	539.67\\
60000	533.65\\
62000	530.52\\
64000	521.34\\
66000	516.68\\
70000	505.01\\
82000	483.66\\
88000	481.56\\
90000	477.78\\
94000	478.39\\
96000	474.09\\
1e+05	469.46\\
};
\addlegendentry{Classe k=10}

            % \input{../../../../.tex/20/KS/SizeEvolutionsfig1.tex}

            \nextgroupplot[%
                % xmin=0,xmax=1e+05,
                % ymin=100,ymax=1.4483e+05,
            ] % This file was created by matlab2tikz.
%
\definecolor{mycolor1}{rgb}{0.00146,0.00047,0.01387}%
\definecolor{mycolor2}{rgb}{0.01166,0.00942,0.06346}%
\definecolor{mycolor3}{rgb}{0.02943,0.02150,0.11462}%
\definecolor{mycolor4}{rgb}{0.06134,0.03659,0.17764}%
\definecolor{mycolor5}{rgb}{0.08741,0.04456,0.22481}%
\definecolor{mycolor6}{rgb}{0.12291,0.04754,0.28162}%
\definecolor{mycolor7}{rgb}{0.14907,0.04547,0.31709}%
\definecolor{mycolor8}{rgb}{0.18343,0.04033,0.35497}%
\definecolor{mycolor9}{rgb}{0.21795,0.03662,0.38352}%
\definecolor{mycolor10}{rgb}{0.24497,0.03705,0.40001}%
\definecolor{mycolor11}{rgb}{0.27135,0.04092,0.41198}%
\definecolor{mycolor12}{rgb}{0.29076,0.04564,0.41864}%
\definecolor{mycolor13}{rgb}{0.31628,0.05349,0.42512}%
\definecolor{mycolor14}{rgb}{0.33522,0.06006,0.42852}%
\definecolor{mycolor15}{rgb}{0.36028,0.06925,0.43150}%
\definecolor{mycolor16}{rgb}{0.37900,0.07625,0.43272}%
\definecolor{mycolor17}{rgb}{0.39767,0.08326,0.43318}%
\definecolor{mycolor18}{rgb}{0.41633,0.09020,0.43294}%
\definecolor{mycolor19}{rgb}{0.42877,0.09479,0.43241}%
\definecolor{mycolor20}{rgb}{0.44743,0.10160,0.43108}%
\definecolor{mycolor21}{rgb}{0.46610,0.10832,0.42912}%
\definecolor{mycolor22}{rgb}{0.47856,0.11276,0.42747}%
\definecolor{mycolor23}{rgb}{0.49726,0.11938,0.42449}%
\definecolor{mycolor24}{rgb}{0.50973,0.12377,0.42216}%
\definecolor{mycolor25}{rgb}{0.52221,0.12815,0.41955}%
\definecolor{mycolor26}{rgb}{0.53468,0.13253,0.41667}%
\definecolor{mycolor27}{rgb}{0.55339,0.13913,0.41183}%
\definecolor{mycolor28}{rgb}{0.56585,0.14357,0.40826}%
\definecolor{mycolor29}{rgb}{0.57830,0.14804,0.40441}%
\definecolor{mycolor30}{rgb}{0.59073,0.15256,0.40029}%
\definecolor{mycolor31}{rgb}{0.60314,0.15715,0.39589}%
\definecolor{mycolor32}{rgb}{0.61551,0.16182,0.39122}%
\definecolor{mycolor33}{rgb}{0.62169,0.16418,0.38878}%
\definecolor{mycolor34}{rgb}{0.63400,0.16899,0.38370}%
\definecolor{mycolor35}{rgb}{0.64626,0.17391,0.37836}%
\definecolor{mycolor36}{rgb}{0.65846,0.17896,0.37275}%
\definecolor{mycolor37}{rgb}{0.66454,0.18154,0.36985}%
\definecolor{mycolor38}{rgb}{0.67664,0.18681,0.36385}%
\definecolor{mycolor39}{rgb}{0.68865,0.19224,0.35760}%
\definecolor{mycolor40}{rgb}{0.69463,0.19502,0.35439}%
\definecolor{mycolor41}{rgb}{0.70650,0.20073,0.34778}%
\definecolor{mycolor42}{rgb}{0.71240,0.20366,0.34438}%
\definecolor{mycolor43}{rgb}{0.72410,0.20967,0.33742}%
\definecolor{mycolor44}{rgb}{0.72991,0.21276,0.33386}%
\definecolor{mycolor45}{rgb}{0.74142,0.21911,0.32658}%
\definecolor{mycolor46}{rgb}{0.74713,0.22238,0.32286}%
\definecolor{mycolor47}{rgb}{0.75279,0.22571,0.31909}%
\definecolor{mycolor48}{rgb}{0.76401,0.23255,0.31140}%
\definecolor{mycolor49}{rgb}{0.76956,0.23608,0.30749}%
\definecolor{mycolor50}{rgb}{0.77506,0.23967,0.30353}%
\definecolor{mycolor51}{rgb}{0.79129,0.25086,0.29139}%
\definecolor{mycolor52}{rgb}{0.79661,0.25473,0.28726}%
\definecolor{mycolor53}{rgb}{0.80187,0.25867,0.28310}%
\definecolor{mycolor54}{rgb}{0.81224,0.26679,0.27466}%
\definecolor{mycolor55}{rgb}{0.81734,0.27095,0.27039}%
\definecolor{mycolor56}{rgb}{0.82737,0.27952,0.26175}%
\definecolor{mycolor57}{rgb}{0.83230,0.28391,0.25738}%
\definecolor{mycolor58}{rgb}{0.83717,0.28839,0.25299}%
\definecolor{mycolor59}{rgb}{0.84671,0.29756,0.24411}%
\definecolor{mycolor60}{rgb}{0.85138,0.30226,0.23964}%
\definecolor{mycolor61}{rgb}{0.85599,0.30704,0.23513}%
\definecolor{mycolor62}{rgb}{0.86053,0.31189,0.23061}%
\definecolor{mycolor63}{rgb}{0.87374,0.32691,0.21689}%
\definecolor{mycolor64}{rgb}{0.87800,0.33206,0.21227}%
\definecolor{mycolor65}{rgb}{0.89034,0.34796,0.19829}%
\definecolor{mycolor66}{rgb}{0.89819,0.35891,0.18886}%
\definecolor{mycolor67}{rgb}{0.90200,0.36449,0.18412}%
\definecolor{mycolor68}{rgb}{0.90573,0.37014,0.17935}%
\definecolor{mycolor69}{rgb}{0.90939,0.37586,0.17456}%
\definecolor{mycolor70}{rgb}{0.91297,0.38164,0.16975}%
\definecolor{mycolor71}{rgb}{0.91646,0.38748,0.16492}%
\definecolor{mycolor72}{rgb}{0.91988,0.39339,0.16007}%
\definecolor{mycolor73}{rgb}{0.92322,0.39936,0.15519}%
\definecolor{mycolor74}{rgb}{0.92647,0.40539,0.15029}%
\definecolor{mycolor75}{rgb}{0.92964,0.41148,0.14537}%
\definecolor{mycolor76}{rgb}{0.93274,0.41763,0.14042}%
\definecolor{mycolor77}{rgb}{0.93575,0.42383,0.13544}%
\definecolor{mycolor78}{rgb}{0.93868,0.43009,0.13044}%
\definecolor{mycolor79}{rgb}{0.94152,0.43640,0.12541}%
\definecolor{mycolor80}{rgb}{0.94429,0.44277,0.12035}%
\definecolor{mycolor81}{rgb}{0.94956,0.45566,0.11016}%
\definecolor{mycolor82}{rgb}{0.95208,0.46218,0.10503}%
\definecolor{mycolor83}{rgb}{0.96129,0.48872,0.08429}%
\definecolor{mycolor84}{rgb}{0.96540,0.50225,0.07386}%
\definecolor{mycolor85}{rgb}{0.96732,0.50908,0.06866}%
\definecolor{mycolor86}{rgb}{0.97259,0.52980,0.05332}%
\definecolor{mycolor87}{rgb}{0.97568,0.54380,0.04362}%
\definecolor{mycolor88}{rgb}{0.98082,0.57221,0.02851}%
\definecolor{mycolor89}{rgb}{0.98189,0.57939,0.02625}%
\definecolor{mycolor90}{rgb}{0.98378,0.59385,0.02377}%
\definecolor{mycolor91}{rgb}{0.98532,0.60842,0.02420}%
\definecolor{mycolor92}{rgb}{0.98696,0.63048,0.03091}%
\definecolor{mycolor93}{rgb}{0.98793,0.66025,0.05175}%
\definecolor{mycolor94}{rgb}{0.98771,0.68281,0.07249}%
\definecolor{mycolor95}{rgb}{0.97750,0.78226,0.18592}%
\definecolor{mycolor96}{rgb}{0.96624,0.83619,0.26153}%
\definecolor{mycolor97}{rgb}{0.96252,0.85148,0.28555}%
\definecolor{mycolor98}{rgb}{0.98836,0.99836,0.64492}%
%
\addplot [color=mycolor1]
  table[row sep=crcr]{%
0	4395.3\\
2000	16346\\
4000	36407\\
6000	52907\\
8000	66810\\
10000	78978\\
12000	89344\\
14000	97986\\
16000	1.0542e+05\\
18000	1.1154e+05\\
20000	1.1678e+05\\
22000	1.2118e+05\\
24000	1.2485e+05\\
26000	1.2825e+05\\
28000	1.3097e+05\\
32000	1.3542e+05\\
38000	1.3893e+05\\
40000	1.3999e+05\\
46000	1.4173e+05\\
52000	1.4271e+05\\
56000	1.4304e+05\\
62000	1.4336e+05\\
82000	1.4453e+05\\
84000	1.4472e+05\\
86000	1.4438e+05\\
1e+05	1.4466e+05\\
};
\addlegendentry{Classe k=283}

\addplot [color=mycolor2, forget plot]
  table[row sep=crcr]{%
0	4395.3\\
2000	8966.7\\
4000	14858\\
6000	18167\\
8000	20461\\
10000	21856\\
12000	22844\\
14000	23356\\
16000	23387\\
18000	23734\\
22000	24168\\
24000	24428\\
26000	24431\\
28000	24851\\
32000	24757\\
34000	24910\\
36000	24963\\
38000	24835\\
40000	24577\\
44000	25060\\
46000	25035\\
48000	25221\\
52000	25376\\
54000	25061\\
58000	25360\\
60000	25348\\
66000	25075\\
68000	24843\\
72000	25188\\
74000	25027\\
76000	25016\\
80000	25184\\
84000	24813\\
86000	25085\\
90000	24989\\
92000	25054\\
96000	24833\\
98000	24998\\
1e+05	24980\\
};
\addplot [color=mycolor3, forget plot]
  table[row sep=crcr]{%
0	4395.3\\
2000	10003\\
4000	18079\\
6000	23624\\
8000	27494\\
10000	30138\\
12000	32190\\
14000	33457\\
16000	34658\\
18000	35483\\
20000	36100\\
24000	36774\\
28000	37020\\
30000	37600\\
32000	37786\\
34000	37534\\
40000	38083\\
42000	38260\\
48000	37992\\
50000	37638\\
52000	37434\\
54000	37707\\
56000	37748\\
58000	37983\\
60000	38100\\
64000	38048\\
68000	38579\\
74000	38271\\
76000	38217\\
82000	38558\\
86000	38596\\
90000	38524\\
96000	38139\\
1e+05	38113\\
};
\addplot [color=mycolor4, forget plot]
  table[row sep=crcr]{%
0	4395.3\\
2000	7142\\
4000	10426\\
6000	12010\\
8000	12914\\
10000	13264\\
12000	13574\\
14000	13587\\
16000	13736\\
18000	13841\\
20000	13875\\
22000	14034\\
26000	13887\\
28000	13876\\
30000	13792\\
32000	13888\\
34000	14077\\
36000	13982\\
38000	13991\\
40000	13960\\
42000	14062\\
44000	13962\\
48000	14061\\
50000	14014\\
52000	14092\\
54000	14025\\
56000	13861\\
58000	13978\\
60000	13924\\
62000	13977\\
64000	13816\\
66000	13894\\
68000	14220\\
72000	14123\\
74000	14006\\
76000	14001\\
78000	14119\\
80000	14028\\
82000	13854\\
84000	13804\\
86000	13908\\
88000	13951\\
90000	13844\\
92000	13875\\
94000	13964\\
96000	13907\\
98000	13791\\
1e+05	13863\\
};
\addplot [color=mycolor5, forget plot]
  table[row sep=crcr]{%
0	4395.3\\
2000	9770.4\\
4000	18613\\
6000	25846\\
8000	31687\\
10000	36548\\
12000	40344\\
14000	44006\\
16000	46743\\
18000	49159\\
20000	50978\\
24000	54077\\
26000	55379\\
30000	57168\\
40000	59367\\
44000	60124\\
46000	60283\\
48000	60212\\
56000	60694\\
60000	61305\\
62000	61398\\
66000	60950\\
68000	60703\\
72000	61009\\
78000	60256\\
86000	60876\\
88000	60854\\
92000	60193\\
96000	60587\\
98000	60478\\
1e+05	60620\\
};
\addplot [color=mycolor6, forget plot]
  table[row sep=crcr]{%
0	4395.3\\
2000	9463.1\\
4000	17618\\
6000	24266\\
8000	29647\\
10000	33750\\
12000	37163\\
14000	40086\\
16000	42469\\
18000	44556\\
20000	46325\\
26000	49737\\
34000	52154\\
36000	52676\\
38000	52848\\
42000	52861\\
46000	53021\\
48000	53264\\
54000	53275\\
56000	53051\\
58000	53046\\
66000	53518\\
70000	53418\\
72000	53286\\
74000	53701\\
76000	53786\\
82000	53294\\
86000	53597\\
92000	53774\\
94000	53273\\
96000	53123\\
1e+05	53404\\
};
\addplot [color=mycolor7, forget plot]
  table[row sep=crcr]{%
0	4395.3\\
2000	7895.1\\
4000	13208\\
6000	17014\\
8000	19722\\
10000	21616\\
12000	22980\\
14000	24046\\
16000	24717\\
18000	25106\\
20000	25653\\
22000	26011\\
24000	26027\\
26000	26357\\
28000	26528\\
30000	26513\\
40000	27094\\
42000	27323\\
44000	27293\\
46000	27018\\
48000	27244\\
50000	27345\\
52000	27008\\
54000	26805\\
58000	27077\\
60000	27161\\
64000	27168\\
66000	27339\\
68000	27327\\
72000	27157\\
76000	27440\\
78000	27416\\
80000	27230\\
82000	27220\\
86000	27380\\
88000	27403\\
90000	27578\\
92000	27299\\
96000	27247\\
98000	27506\\
1e+05	27427\\
};
\addplot [color=mycolor8, forget plot]
  table[row sep=crcr]{%
0	4395.3\\
2000	7116.2\\
4000	11044\\
6000	13719\\
8000	15613\\
10000	16801\\
12000	17633\\
14000	18210\\
16000	18677\\
18000	18972\\
20000	19218\\
26000	19665\\
28000	19697\\
30000	19836\\
34000	19772\\
38000	19732\\
40000	19635\\
42000	19381\\
46000	19865\\
52000	20266\\
54000	20003\\
56000	19959\\
60000	20032\\
62000	20136\\
64000	20105\\
66000	19977\\
72000	20056\\
74000	19886\\
76000	19823\\
78000	19837\\
80000	19778\\
82000	19787\\
88000	20003\\
90000	20010\\
94000	19757\\
98000	19655\\
1e+05	19814\\
};
\addplot [color=mycolor9, forget plot]
  table[row sep=crcr]{%
0	4395.3\\
2000	6447.9\\
4000	9190.5\\
6000	10929\\
8000	11985\\
10000	12555\\
12000	12820\\
14000	13133\\
18000	13488\\
22000	13688\\
24000	13671\\
28000	13735\\
36000	13653\\
38000	13560\\
40000	13639\\
44000	13619\\
46000	13569\\
48000	13724\\
50000	13732\\
52000	13829\\
54000	13882\\
56000	13896\\
58000	13758\\
62000	13713\\
64000	13738\\
66000	13660\\
70000	13800\\
72000	13661\\
78000	13959\\
80000	13907\\
82000	13734\\
84000	13693\\
88000	13736\\
90000	13839\\
92000	13828\\
94000	13677\\
98000	13703\\
1e+05	13681\\
};
\addplot [color=mycolor10, forget plot]
  table[row sep=crcr]{%
0	4395.3\\
2000	6764.8\\
4000	10291\\
6000	12871\\
8000	14809\\
10000	16203\\
12000	17227\\
14000	18060\\
16000	18714\\
18000	19258\\
20000	19624\\
26000	20502\\
28000	20743\\
30000	20903\\
38000	21283\\
42000	21041\\
44000	21105\\
46000	21287\\
50000	21378\\
52000	21358\\
54000	21461\\
56000	21466\\
58000	21322\\
64000	21237\\
66000	21296\\
70000	21098\\
72000	21052\\
76000	21157\\
78000	21342\\
80000	21385\\
82000	21353\\
86000	21530\\
88000	21476\\
90000	21481\\
92000	21346\\
94000	21307\\
96000	21335\\
98000	21244\\
1e+05	21353\\
};
\addplot [color=mycolor11, forget plot]
  table[row sep=crcr]{%
0	4395.3\\
2000	5586.2\\
4000	7069.9\\
6000	7853.9\\
8000	8244.3\\
10000	8496.7\\
12000	8614.2\\
14000	8613.4\\
16000	8582.7\\
18000	8623.9\\
20000	8580.3\\
22000	8507.3\\
26000	8675.1\\
28000	8605.9\\
30000	8579.6\\
32000	8613.1\\
34000	8582.2\\
38000	8613\\
42000	8403.5\\
46000	8521\\
48000	8544.8\\
50000	8533.3\\
52000	8576.9\\
56000	8489.6\\
58000	8544.3\\
60000	8476.6\\
66000	8468\\
68000	8556.3\\
70000	8537\\
72000	8440.7\\
76000	8487\\
80000	8442.5\\
82000	8477.5\\
84000	8567.9\\
88000	8677.1\\
90000	8579.9\\
92000	8626\\
94000	8610.2\\
96000	8664.1\\
1e+05	8548.3\\
};
\addplot [color=mycolor12, forget plot]
  table[row sep=crcr]{%
0	4395.3\\
2000	6312.6\\
4000	9096.1\\
6000	10813\\
8000	11930\\
10000	12791\\
12000	13291\\
14000	13644\\
18000	13960\\
20000	14025\\
22000	14031\\
24000	14186\\
28000	14018\\
32000	14350\\
36000	14285\\
40000	14446\\
42000	14483\\
44000	14361\\
46000	14419\\
50000	14436\\
52000	14520\\
54000	14532\\
56000	14435\\
58000	14583\\
60000	14507\\
62000	14534\\
64000	14493\\
68000	14584\\
70000	14524\\
72000	14274\\
74000	14260\\
76000	14418\\
78000	14421\\
80000	14313\\
86000	14278\\
90000	14370\\
92000	14430\\
96000	14607\\
1e+05	14384\\
};
\addplot [color=mycolor13, forget plot]
  table[row sep=crcr]{%
0	4395.3\\
4000	7895.4\\
6000	9147.7\\
8000	9754.3\\
10000	10213\\
12000	10496\\
16000	10754\\
18000	10947\\
22000	11145\\
26000	11106\\
28000	11007\\
30000	11086\\
32000	11013\\
36000	11004\\
38000	10900\\
40000	10867\\
42000	10973\\
44000	10874\\
48000	10792\\
50000	10900\\
52000	10806\\
60000	10969\\
62000	10932\\
64000	11064\\
66000	11156\\
72000	11121\\
74000	11127\\
76000	11102\\
80000	11218\\
82000	11169\\
84000	11082\\
86000	11111\\
88000	11071\\
90000	11212\\
92000	11090\\
96000	11024\\
98000	11111\\
1e+05	11043\\
};
\addplot [color=mycolor14, forget plot]
  table[row sep=crcr]{%
0	4395.3\\
2000	7047.5\\
4000	11111\\
6000	14157\\
8000	16349\\
10000	17944\\
12000	19193\\
14000	20092\\
16000	20848\\
18000	21505\\
20000	21841\\
22000	22297\\
24000	22540\\
28000	22944\\
30000	22954\\
32000	23088\\
36000	23165\\
38000	23264\\
40000	23506\\
44000	23585\\
46000	23502\\
48000	23295\\
50000	23282\\
52000	23561\\
56000	23745\\
60000	23528\\
62000	23305\\
64000	23453\\
66000	23507\\
68000	23411\\
70000	23423\\
72000	23513\\
74000	23468\\
78000	23521\\
86000	23203\\
88000	23167\\
90000	23317\\
92000	23197\\
94000	23202\\
96000	23367\\
98000	23221\\
1e+05	23418\\
};
\addplot [color=mycolor15, forget plot]
  table[row sep=crcr]{%
0	4395.3\\
2000	6510.2\\
4000	9526.1\\
6000	11426\\
8000	12831\\
10000	13820\\
12000	14499\\
16000	15784\\
18000	16039\\
20000	16150\\
24000	16156\\
26000	16194\\
28000	16323\\
30000	16341\\
32000	16264\\
36000	16502\\
40000	16544\\
44000	16482\\
46000	16600\\
48000	16514\\
54000	16537\\
56000	16445\\
58000	16476\\
64000	16447\\
66000	16515\\
70000	16549\\
72000	16672\\
74000	16578\\
76000	16545\\
78000	16435\\
80000	16670\\
82000	16632\\
84000	16542\\
88000	16615\\
90000	16669\\
94000	16714\\
96000	16571\\
98000	16560\\
1e+05	16701\\
};
\addplot [color=mycolor16, forget plot]
  table[row sep=crcr]{%
0	4395.3\\
2000	5142.7\\
4000	6073.2\\
6000	6668.9\\
8000	6927\\
10000	7041.2\\
12000	7097.6\\
14000	7078.7\\
20000	7114.5\\
22000	7189.2\\
28000	7032.9\\
32000	7038.2\\
38000	6981.3\\
40000	7010.9\\
42000	7081.7\\
44000	7114\\
46000	7029.2\\
48000	6966.9\\
52000	7002.1\\
54000	6978.5\\
58000	7033.7\\
64000	6926.1\\
68000	6978.2\\
72000	7029.9\\
74000	7007.1\\
76000	6875.1\\
78000	7018.9\\
80000	7122.8\\
82000	7167.6\\
86000	7013.1\\
90000	6948.2\\
92000	6961.3\\
94000	7012.6\\
96000	6971.1\\
1e+05	7028.2\\
};
\addplot [color=mycolor17, forget plot]
  table[row sep=crcr]{%
0	4395.3\\
4000	9139.6\\
6000	11153\\
8000	12497\\
10000	13340\\
12000	13894\\
14000	14381\\
16000	14640\\
22000	15121\\
24000	15241\\
32000	15490\\
34000	15600\\
36000	15509\\
38000	15628\\
40000	15619\\
42000	15707\\
44000	15680\\
46000	15572\\
48000	15523\\
52000	15737\\
56000	15564\\
60000	15460\\
62000	15385\\
68000	15402\\
72000	15545\\
74000	15717\\
78000	15755\\
82000	15893\\
88000	15481\\
90000	15470\\
92000	15584\\
96000	15662\\
98000	15638\\
1e+05	15535\\
};
\addplot [color=mycolor18, forget plot]
  table[row sep=crcr]{%
0	4395.3\\
2000	6892.9\\
4000	10664\\
6000	13583\\
8000	15785\\
10000	17529\\
12000	18861\\
14000	19960\\
16000	20865\\
20000	21994\\
22000	22540\\
24000	22709\\
26000	22769\\
28000	23025\\
30000	23116\\
32000	23105\\
34000	23320\\
38000	23555\\
40000	23562\\
42000	23639\\
44000	23830\\
46000	23949\\
50000	23904\\
52000	23685\\
56000	23598\\
58000	23633\\
60000	23775\\
62000	24014\\
64000	24075\\
68000	23949\\
70000	23965\\
72000	23693\\
74000	23748\\
76000	23715\\
78000	23596\\
80000	23343\\
82000	23335\\
86000	23557\\
88000	23798\\
90000	23845\\
92000	23729\\
96000	23659\\
98000	23596\\
1e+05	23626\\
};
\addplot [color=mycolor19, forget plot]
  table[row sep=crcr]{%
0	4395.3\\
2000	6149.9\\
4000	8737.9\\
6000	10500\\
8000	11582\\
10000	12398\\
12000	12979\\
14000	13227\\
16000	13578\\
18000	13832\\
22000	14104\\
24000	14325\\
26000	14348\\
28000	14327\\
30000	14383\\
32000	14571\\
34000	14672\\
38000	14735\\
42000	14853\\
44000	14905\\
48000	14701\\
52000	14641\\
58000	14729\\
66000	14618\\
68000	14696\\
70000	14709\\
72000	14569\\
74000	14570\\
76000	14680\\
80000	14516\\
84000	14602\\
86000	14476\\
88000	14418\\
90000	14410\\
96000	14598\\
98000	14671\\
1e+05	14562\\
};
\addplot [color=mycolor20, forget plot]
  table[row sep=crcr]{%
0	4395.3\\
2000	6478.4\\
4000	9683.9\\
6000	12001\\
8000	13653\\
10000	14945\\
14000	16460\\
16000	16991\\
18000	17308\\
20000	17733\\
22000	18061\\
24000	18223\\
28000	18246\\
30000	18297\\
32000	18555\\
38000	18679\\
42000	18432\\
44000	18531\\
46000	18464\\
48000	18345\\
50000	18429\\
52000	18667\\
54000	18712\\
56000	18656\\
58000	18660\\
60000	18806\\
62000	18880\\
64000	18780\\
66000	18864\\
70000	18810\\
72000	18877\\
78000	18760\\
80000	18643\\
82000	18680\\
84000	18769\\
90000	18591\\
92000	18650\\
96000	18464\\
1e+05	18427\\
};
\addplot [color=mycolor21, forget plot]
  table[row sep=crcr]{%
0	4395.3\\
2000	5340.6\\
4000	6774.7\\
6000	7694.6\\
8000	8316.5\\
10000	8661.1\\
12000	8990.6\\
14000	9135.3\\
16000	9310.3\\
22000	9465.6\\
24000	9496.4\\
26000	9443.6\\
28000	9418\\
30000	9430.3\\
32000	9508.1\\
34000	9437.1\\
36000	9423.9\\
38000	9468.3\\
40000	9540.2\\
44000	9568.1\\
48000	9524.9\\
52000	9514.7\\
54000	9472.6\\
56000	9461.5\\
62000	9573.9\\
64000	9567.5\\
66000	9485.4\\
68000	9496.5\\
72000	9596.5\\
74000	9555.5\\
80000	9498.5\\
82000	9546.4\\
84000	9538.9\\
88000	9577.8\\
90000	9554.2\\
92000	9641.1\\
94000	9678.1\\
96000	9630.7\\
1e+05	9683.8\\
};
\addplot [color=mycolor22, forget plot]
  table[row sep=crcr]{%
0	4395.3\\
2000	6766.3\\
4000	10607\\
6000	13538\\
8000	15874\\
10000	17622\\
12000	19035\\
14000	20117\\
16000	20994\\
18000	21756\\
20000	22404\\
22000	22872\\
24000	23257\\
26000	23506\\
28000	23845\\
32000	24209\\
34000	24403\\
38000	24592\\
40000	24421\\
46000	24612\\
50000	24392\\
56000	24415\\
58000	24192\\
62000	24246\\
64000	24423\\
66000	24444\\
68000	24289\\
80000	24176\\
84000	24384\\
86000	24315\\
88000	24173\\
90000	24193\\
94000	24434\\
96000	24328\\
98000	24034\\
1e+05	24076\\
};
\addplot [color=mycolor23, forget plot]
  table[row sep=crcr]{%
0	4395.3\\
2000	5718.2\\
4000	7640.9\\
6000	8971.1\\
8000	9860.8\\
10000	10494\\
12000	10971\\
14000	11305\\
16000	11514\\
18000	11680\\
22000	11877\\
24000	11950\\
26000	11953\\
34000	12133\\
38000	12137\\
42000	11996\\
46000	12034\\
48000	12132\\
54000	12191\\
58000	12045\\
60000	11999\\
64000	12113\\
70000	12056\\
74000	12136\\
76000	12123\\
78000	12065\\
82000	12135\\
90000	11989\\
92000	12067\\
98000	12162\\
1e+05	12142\\
};
\addplot [color=mycolor24, forget plot]
  table[row sep=crcr]{%
0	4395.3\\
2000	4562.7\\
4000	4799.2\\
6000	4942.9\\
10000	4997.8\\
12000	5057.6\\
14000	5014.3\\
16000	5038.8\\
18000	5013.1\\
20000	4969.4\\
26000	4948.2\\
28000	4923.6\\
30000	4877.6\\
32000	4859.3\\
34000	4823.8\\
36000	4817\\
38000	4776.6\\
40000	4769.1\\
44000	4733.4\\
46000	4828\\
50000	4862.5\\
52000	4832.2\\
54000	4853.7\\
56000	4890.8\\
62000	4887.1\\
64000	4851.1\\
66000	4837.1\\
70000	4832.8\\
72000	4768.1\\
74000	4727.8\\
76000	4722.5\\
78000	4742.8\\
80000	4818.1\\
84000	4790.8\\
94000	4783.3\\
98000	4744\\
1e+05	4797.6\\
};
\addplot [color=mycolor25, forget plot]
  table[row sep=crcr]{%
0	4395.3\\
2000	5773.8\\
4000	7879.1\\
6000	9223.8\\
8000	10142\\
10000	10851\\
12000	11260\\
14000	11527\\
16000	11759\\
18000	11811\\
20000	11939\\
22000	11939\\
24000	12076\\
28000	12123\\
30000	12246\\
32000	12213\\
34000	12217\\
38000	11986\\
42000	12301\\
44000	12408\\
46000	12386\\
48000	12322\\
50000	12183\\
52000	12109\\
54000	12173\\
56000	12139\\
60000	12183\\
62000	12205\\
64000	12128\\
66000	12232\\
68000	12277\\
70000	12209\\
72000	12199\\
76000	12100\\
78000	12051\\
80000	12155\\
82000	12192\\
84000	12130\\
86000	11956\\
88000	12014\\
94000	12036\\
96000	12013\\
1e+05	12278\\
};
\addplot [color=mycolor26, forget plot]
  table[row sep=crcr]{%
0	4395.3\\
2000	5292\\
4000	6474.1\\
6000	7186.8\\
8000	7606.6\\
10000	7868.2\\
12000	8057.6\\
14000	8207.7\\
16000	8242.2\\
18000	8207\\
22000	8275.5\\
26000	8198\\
28000	8224.9\\
30000	8206.7\\
32000	8262.5\\
34000	8256.4\\
36000	8138.1\\
38000	8113.8\\
40000	8132\\
42000	8187.7\\
44000	8199.3\\
46000	8247.1\\
48000	8204.8\\
50000	8234.4\\
52000	8241.2\\
54000	8136.9\\
58000	8258.2\\
60000	8256.7\\
62000	8164\\
64000	8097.5\\
66000	8203.7\\
68000	8178.2\\
70000	8178.9\\
74000	8314.6\\
80000	8223.4\\
82000	8237.3\\
84000	8355.9\\
88000	8304.3\\
94000	8322.6\\
96000	8343.5\\
98000	8450.2\\
1e+05	8438.6\\
};
\addplot [color=mycolor27, forget plot]
  table[row sep=crcr]{%
0	4395.3\\
2000	5516.3\\
4000	7062.4\\
6000	7997\\
8000	8599.5\\
10000	8984\\
12000	9246.3\\
14000	9470.2\\
20000	9483.5\\
22000	9523.9\\
24000	9601.3\\
26000	9631.8\\
30000	9766.1\\
32000	9662.7\\
34000	9626.8\\
36000	9662.6\\
40000	9807.2\\
42000	9773.7\\
44000	9798.3\\
50000	9804.1\\
52000	9735.4\\
56000	9774\\
58000	9735.4\\
60000	9724.8\\
62000	9785.1\\
66000	9718.2\\
68000	9689.9\\
70000	9579\\
72000	9674\\
74000	9742.2\\
78000	9748.5\\
82000	9558.9\\
86000	9419.6\\
88000	9473.2\\
90000	9498.4\\
92000	9596.4\\
94000	9657.1\\
1e+05	9723.3\\
};
\addplot [color=mycolor28, forget plot]
  table[row sep=crcr]{%
0	4395.3\\
2000	4744.8\\
4000	5186.3\\
6000	5417.2\\
8000	5558.9\\
10000	5606.4\\
14000	5602.4\\
16000	5617\\
18000	5656.2\\
20000	5584.5\\
26000	5547.6\\
28000	5468.6\\
30000	5436.9\\
32000	5441.1\\
34000	5499.7\\
36000	5516.8\\
38000	5476.4\\
40000	5461.2\\
42000	5529.7\\
44000	5564.9\\
46000	5543.1\\
48000	5485.7\\
50000	5446.2\\
52000	5518.2\\
54000	5514.5\\
56000	5535.2\\
58000	5527.1\\
60000	5562.6\\
62000	5521.7\\
64000	5529.5\\
66000	5440\\
68000	5430.4\\
70000	5506\\
72000	5513.8\\
74000	5546.1\\
76000	5537.1\\
78000	5506\\
80000	5561.4\\
82000	5538.8\\
84000	5466.8\\
86000	5479.7\\
90000	5550.3\\
94000	5558\\
98000	5520.7\\
1e+05	5500.9\\
};
\addplot [color=mycolor29, forget plot]
  table[row sep=crcr]{%
0	4395.3\\
2000	5653.3\\
4000	7480.7\\
6000	8737\\
8000	9606.9\\
10000	10187\\
12000	10638\\
14000	10872\\
16000	11012\\
20000	11208\\
22000	11196\\
24000	11230\\
26000	11346\\
28000	11363\\
30000	11413\\
32000	11399\\
34000	11303\\
36000	11348\\
38000	11433\\
40000	11460\\
42000	11424\\
44000	11494\\
46000	11392\\
48000	11444\\
50000	11448\\
52000	11540\\
54000	11522\\
56000	11577\\
60000	11535\\
64000	11583\\
68000	11452\\
70000	11453\\
72000	11502\\
78000	11344\\
80000	11383\\
82000	11373\\
86000	11474\\
88000	11399\\
90000	11417\\
92000	11532\\
94000	11581\\
96000	11557\\
98000	11584\\
1e+05	11515\\
};
\addplot [color=mycolor30, forget plot]
  table[row sep=crcr]{%
0	4395.3\\
2000	5140.5\\
4000	6166.5\\
6000	6763.5\\
8000	7173.5\\
10000	7479.4\\
12000	7658.9\\
14000	7769.3\\
22000	7866.2\\
24000	7915\\
46000	7832.5\\
48000	7905.8\\
54000	7864.6\\
70000	7912.6\\
72000	7868.8\\
76000	7859.3\\
78000	7939.3\\
80000	7883.1\\
84000	7876.4\\
88000	7937.3\\
90000	7887.2\\
92000	7928.3\\
94000	7938.2\\
96000	7894\\
98000	7919.5\\
1e+05	7882.5\\
};
\addplot [color=mycolor31, forget plot]
  table[row sep=crcr]{%
0	4395.3\\
2000	4936.4\\
4000	5637.9\\
6000	6043.9\\
8000	6305.3\\
10000	6429.7\\
12000	6484.8\\
14000	6472.5\\
16000	6504\\
20000	6523.4\\
22000	6595\\
24000	6550.6\\
26000	6591.6\\
28000	6594.6\\
30000	6530.2\\
32000	6522.5\\
34000	6560.4\\
36000	6517.5\\
42000	6578.9\\
44000	6483.2\\
46000	6517.9\\
48000	6523.8\\
50000	6511.4\\
54000	6415.2\\
56000	6412.4\\
58000	6470.1\\
60000	6434.6\\
62000	6380.4\\
68000	6563.7\\
70000	6565.8\\
72000	6619.4\\
74000	6571\\
76000	6574.7\\
80000	6547\\
84000	6382.7\\
94000	6534.6\\
96000	6521.8\\
98000	6449.4\\
1e+05	6523.4\\
};
\addplot [color=mycolor32, forget plot]
  table[row sep=crcr]{%
0	4395.3\\
2000	5520.4\\
4000	7129.4\\
6000	8176.2\\
8000	8871.7\\
10000	9298.6\\
12000	9612.7\\
14000	9981.9\\
18000	10245\\
22000	10421\\
24000	10436\\
28000	10521\\
32000	10495\\
34000	10419\\
38000	10473\\
44000	10322\\
46000	10424\\
50000	10361\\
52000	10399\\
54000	10404\\
58000	10282\\
64000	10403\\
70000	10253\\
76000	10362\\
78000	10321\\
80000	10190\\
82000	10236\\
88000	10169\\
90000	10287\\
92000	10371\\
96000	10334\\
98000	10283\\
1e+05	10283\\
};
\addplot [color=mycolor33, forget plot]
  table[row sep=crcr]{%
0	4395.3\\
2000	5224.4\\
4000	6368.8\\
6000	7036.7\\
8000	7458.8\\
10000	7713\\
12000	7918.2\\
14000	8078.8\\
18000	8133.9\\
20000	8100.6\\
24000	8111.2\\
26000	8163.6\\
30000	8076.7\\
34000	8166.2\\
36000	8103.7\\
38000	8102\\
42000	8132.5\\
44000	8147.7\\
46000	8211.5\\
48000	8215.8\\
50000	8104\\
52000	8094.6\\
54000	8177.2\\
56000	8215.5\\
60000	8136.1\\
62000	8235.9\\
64000	8237.3\\
68000	8079\\
70000	8152.2\\
76000	8143.7\\
78000	8170\\
80000	8135.9\\
84000	8208.5\\
88000	8349.8\\
90000	8272.6\\
94000	8287.2\\
98000	8310\\
1e+05	8227.5\\
};
\addplot [color=mycolor34, forget plot]
  table[row sep=crcr]{%
0	4395.3\\
2000	4412.2\\
6000	4351.7\\
8000	4303.3\\
10000	4279.9\\
12000	4323.5\\
16000	4180\\
18000	4134.2\\
20000	4118.9\\
22000	4136.1\\
24000	4166.5\\
26000	4173.5\\
28000	4204.9\\
30000	4136.8\\
32000	4082.7\\
34000	4103.3\\
36000	4106.6\\
38000	4066.1\\
40000	4084.3\\
42000	4079.1\\
44000	4091.9\\
46000	4090.1\\
48000	4106.1\\
50000	4045.2\\
52000	4048.7\\
54000	4087.1\\
58000	4091.8\\
60000	4176.4\\
62000	4178.2\\
64000	4198\\
66000	4179.6\\
72000	4038.4\\
74000	4078.1\\
76000	4104.2\\
82000	4037.3\\
84000	4055.7\\
86000	4101.2\\
88000	4100.3\\
90000	4084\\
94000	4118.5\\
96000	4076.5\\
1e+05	4085\\
};
\addplot [color=mycolor35]
  table[row sep=crcr]{%
0	4395.3\\
2000	5383.3\\
4000	6839.9\\
6000	7812.3\\
8000	8454.8\\
10000	9030\\
12000	9350.1\\
14000	9604.3\\
16000	9790.8\\
18000	9854.4\\
26000	9981\\
30000	10044\\
32000	9968.5\\
38000	9943\\
42000	10057\\
44000	10008\\
46000	9989.5\\
48000	10064\\
52000	10067\\
54000	10159\\
58000	10113\\
60000	10070\\
64000	10119\\
68000	10050\\
70000	10055\\
72000	10120\\
80000	10187\\
84000	10144\\
86000	10151\\
90000	10060\\
98000	10175\\
1e+05	10128\\
};
\addlegendentry{Classe k=83}

\addplot [color=mycolor36, forget plot]
  table[row sep=crcr]{%
0	4395.3\\
2000	4483.4\\
4000	4558.8\\
6000	4559.3\\
8000	4531.8\\
10000	4476.5\\
12000	4476.1\\
14000	4493.7\\
16000	4493.4\\
18000	4480.7\\
20000	4423.1\\
22000	4407.7\\
26000	4328.4\\
28000	4329.8\\
30000	4364.7\\
32000	4373.5\\
34000	4335.6\\
38000	4372.8\\
40000	4337.1\\
42000	4352.6\\
44000	4326.2\\
48000	4327.9\\
50000	4353.7\\
52000	4332.8\\
54000	4360.9\\
56000	4308.5\\
58000	4291.2\\
60000	4309.9\\
62000	4345.7\\
64000	4417.5\\
66000	4420.5\\
68000	4392\\
70000	4394.4\\
72000	4374.3\\
74000	4391.2\\
78000	4400.2\\
80000	4419.9\\
82000	4460.2\\
84000	4428\\
86000	4444.5\\
88000	4446.4\\
90000	4335.5\\
92000	4292.3\\
94000	4314.4\\
1e+05	4287.9\\
};
\addplot [color=mycolor37, forget plot]
  table[row sep=crcr]{%
0	4395.3\\
2000	4267.7\\
4000	4127.8\\
6000	4088.2\\
8000	4034.7\\
10000	3963.7\\
14000	3913.1\\
16000	3860.3\\
18000	3842.7\\
20000	3805.7\\
22000	3812.7\\
24000	3779.1\\
26000	3764.2\\
28000	3708.2\\
30000	3728.8\\
40000	3713.8\\
42000	3737.5\\
50000	3727.5\\
52000	3751.9\\
54000	3744\\
56000	3715.1\\
58000	3752.3\\
60000	3773.9\\
62000	3769.6\\
64000	3751.5\\
66000	3755.7\\
68000	3733.3\\
70000	3692.7\\
72000	3691.8\\
74000	3729.6\\
76000	3696.9\\
78000	3689.6\\
82000	3718.9\\
84000	3702.9\\
86000	3727.3\\
88000	3696.4\\
90000	3700.9\\
92000	3725.1\\
98000	3726.2\\
1e+05	3747.1\\
};
\addplot [color=mycolor38, forget plot]
  table[row sep=crcr]{%
0	4395.3\\
2000	4120.6\\
4000	3799.9\\
6000	3570.8\\
8000	3423.2\\
10000	3379.1\\
12000	3315\\
14000	3292\\
16000	3252.5\\
18000	3229.6\\
20000	3156.2\\
24000	3121.5\\
26000	3084.6\\
28000	3127.8\\
30000	3114.7\\
32000	3124.9\\
34000	3122.9\\
36000	3101.8\\
38000	3058.7\\
40000	3049.9\\
42000	3089.2\\
44000	3074.6\\
46000	3078.6\\
48000	3039.4\\
52000	3098.9\\
54000	3127.8\\
56000	3133.3\\
58000	3093.4\\
60000	3122.6\\
62000	3114.9\\
64000	3078.9\\
66000	3081.6\\
68000	3093.3\\
70000	3123.5\\
72000	3095.9\\
76000	3140.5\\
78000	3144.5\\
80000	3161.9\\
82000	3145.1\\
84000	3083.5\\
86000	3102.1\\
90000	3123.5\\
92000	3113.8\\
94000	3086.5\\
96000	3100.6\\
98000	3129.1\\
1e+05	3116\\
};
\addplot [color=mycolor39, forget plot]
  table[row sep=crcr]{%
0	4395.3\\
2000	5487.5\\
4000	7147.7\\
6000	8328.7\\
8000	9274.8\\
10000	9873.3\\
12000	10274\\
14000	10542\\
16000	10851\\
18000	11060\\
20000	11232\\
24000	11411\\
26000	11329\\
28000	11283\\
32000	11324\\
34000	11326\\
36000	11384\\
40000	11616\\
42000	11646\\
46000	11578\\
48000	11639\\
50000	11568\\
54000	11477\\
56000	11544\\
58000	11577\\
60000	11660\\
66000	11675\\
78000	11360\\
80000	11515\\
84000	11618\\
86000	11623\\
88000	11701\\
90000	11645\\
92000	11722\\
96000	11574\\
98000	11424\\
1e+05	11510\\
};
\addplot [color=mycolor40, forget plot]
  table[row sep=crcr]{%
0	4395.3\\
2000	4655.4\\
4000	5041.3\\
6000	5358.1\\
8000	5551.2\\
10000	5705.1\\
14000	5759.1\\
16000	5808.7\\
18000	5836.7\\
20000	5805.8\\
24000	5828.6\\
28000	5761.7\\
34000	5788.6\\
38000	5758.5\\
40000	5819.1\\
42000	5846.5\\
44000	5808.8\\
46000	5732.4\\
48000	5764.9\\
54000	5777.1\\
56000	5795.9\\
62000	5782.8\\
64000	5752.2\\
66000	5737.8\\
70000	5778.7\\
72000	5791\\
78000	5751.5\\
80000	5770.1\\
84000	5768.7\\
86000	5783.4\\
90000	5854.1\\
94000	5824.6\\
96000	5794.8\\
98000	5783.1\\
1e+05	5744.2\\
};
\addplot [color=mycolor41, forget plot]
  table[row sep=crcr]{%
0	4395.3\\
2000	4539.3\\
4000	4756.3\\
6000	4966.1\\
8000	5131\\
10000	5242.6\\
14000	5387.8\\
16000	5388.7\\
18000	5428\\
22000	5421\\
24000	5406.1\\
26000	5433.9\\
28000	5437.3\\
30000	5404.5\\
36000	5425.6\\
40000	5416.2\\
42000	5452\\
50000	5396.7\\
52000	5411.7\\
60000	5406.5\\
64000	5446.4\\
66000	5476.9\\
68000	5448\\
70000	5481.1\\
72000	5433.8\\
74000	5409.6\\
76000	5406.9\\
78000	5437.7\\
88000	5464.7\\
90000	5421.6\\
96000	5431.8\\
98000	5395.5\\
1e+05	5421\\
};
\addplot [color=mycolor42, forget plot]
  table[row sep=crcr]{%
0	4395.3\\
2000	5990.6\\
4000	8627.8\\
6000	10671\\
8000	12468\\
10000	13925\\
12000	15085\\
14000	16049\\
16000	16696\\
18000	17297\\
20000	17801\\
22000	18208\\
26000	18897\\
28000	19063\\
30000	18927\\
38000	19422\\
42000	19708\\
44000	19783\\
46000	19737\\
50000	19878\\
58000	19818\\
68000	19914\\
70000	19797\\
72000	19959\\
80000	20076\\
82000	20075\\
86000	19957\\
90000	20078\\
92000	19954\\
96000	20085\\
98000	20164\\
1e+05	20099\\
};
\addplot [color=mycolor43, forget plot]
  table[row sep=crcr]{%
0	4395.3\\
2000	4557.9\\
4000	4780.2\\
6000	4925.2\\
8000	5029.1\\
12000	5090.4\\
16000	5066.8\\
18000	5032\\
20000	5035.1\\
22000	5055.1\\
26000	5038.1\\
28000	5025.4\\
30000	5041.6\\
34000	4987.3\\
36000	5007.7\\
38000	5069.1\\
42000	4996.3\\
46000	5003.9\\
48000	4965.3\\
50000	4966.4\\
52000	4949\\
54000	4966.7\\
56000	5012.7\\
58000	5019.1\\
60000	4992.8\\
62000	5001.7\\
64000	4970.1\\
66000	4958.7\\
68000	4965.3\\
70000	4994.1\\
72000	4994.2\\
74000	4967.2\\
76000	5009.8\\
78000	4994.2\\
80000	5021.6\\
82000	5017.4\\
86000	4971\\
88000	5022.3\\
90000	5058.8\\
92000	5029\\
94000	5039\\
98000	5001.6\\
1e+05	5032.4\\
};
\addplot [color=mycolor44, forget plot]
  table[row sep=crcr]{%
0	4395.3\\
2000	4279.8\\
4000	4026.1\\
6000	3835.3\\
8000	3733\\
10000	3689.4\\
12000	3636.1\\
14000	3571.7\\
16000	3541\\
18000	3531.1\\
22000	3459.7\\
24000	3468.6\\
28000	3453.3\\
30000	3422.3\\
32000	3448.3\\
34000	3502.8\\
36000	3481.8\\
38000	3480.7\\
40000	3468.5\\
42000	3427.2\\
44000	3401.9\\
46000	3450.3\\
48000	3442.7\\
50000	3420.1\\
52000	3418.8\\
54000	3445.5\\
56000	3489.4\\
58000	3470.3\\
62000	3415.7\\
64000	3414.3\\
66000	3425.6\\
68000	3401.9\\
70000	3447.6\\
72000	3417.1\\
76000	3417.2\\
78000	3443.7\\
80000	3422\\
82000	3383.1\\
84000	3392.8\\
86000	3384.3\\
88000	3401.8\\
92000	3471.6\\
94000	3448.4\\
96000	3452.4\\
1e+05	3492\\
};
\addplot [color=mycolor45, forget plot]
  table[row sep=crcr]{%
0	4395.3\\
2000	4890.1\\
4000	5613.7\\
6000	6105\\
8000	6415.6\\
10000	6622.8\\
12000	6755.4\\
14000	6848.5\\
16000	6896\\
20000	6929.1\\
26000	6931.3\\
28000	6952.7\\
30000	6934.1\\
32000	6946.4\\
40000	6912.7\\
42000	6884.1\\
44000	6885.7\\
46000	6932.8\\
48000	6920.4\\
50000	6973\\
54000	6950.7\\
56000	6962.6\\
58000	7034.8\\
60000	6986.6\\
62000	6970.4\\
64000	7007.3\\
66000	7011.6\\
68000	6969.8\\
70000	6983.9\\
74000	6917.1\\
76000	6951.5\\
78000	6915\\
80000	6932.3\\
82000	6969.3\\
92000	6897.6\\
94000	6845.3\\
98000	6878.3\\
1e+05	6953.7\\
};
\addplot [color=mycolor46, forget plot]
  table[row sep=crcr]{%
0	4395.3\\
2000	4337.4\\
4000	4296\\
6000	4312.5\\
8000	4313.5\\
12000	4352.9\\
14000	4310.1\\
22000	4219.6\\
24000	4243.5\\
28000	4227.1\\
30000	4218.8\\
32000	4144.4\\
34000	4115.6\\
36000	4151.3\\
38000	4130.5\\
40000	4153.3\\
44000	4158.3\\
46000	4152.4\\
48000	4109.7\\
50000	4119.7\\
52000	4144.1\\
54000	4156.2\\
56000	4203.1\\
58000	4190.9\\
64000	4106.4\\
68000	4148\\
70000	4130.6\\
72000	4132.4\\
74000	4109.6\\
80000	4174.6\\
82000	4188.8\\
88000	4120.7\\
90000	4130.5\\
94000	4082\\
96000	4125\\
98000	4116.6\\
1e+05	4126.4\\
};
\addplot [color=mycolor47, forget plot]
  table[row sep=crcr]{%
0	4395.3\\
2000	4563.9\\
4000	4781\\
6000	4950.8\\
8000	5084.1\\
10000	5192.8\\
12000	5270\\
14000	5305.3\\
18000	5326.3\\
28000	5335.8\\
36000	5275.5\\
40000	5353.9\\
42000	5365.6\\
44000	5332.2\\
46000	5326.8\\
48000	5287.9\\
50000	5311.9\\
52000	5312.1\\
56000	5364.7\\
58000	5368.2\\
60000	5342.8\\
62000	5353.4\\
64000	5323.8\\
68000	5298.1\\
70000	5259.2\\
72000	5273.7\\
74000	5258.6\\
80000	5269.9\\
84000	5250.6\\
86000	5291\\
88000	5273.1\\
1e+05	5345.7\\
};
\addplot [color=mycolor48, forget plot]
  table[row sep=crcr]{%
0	4395.3\\
2000	4116\\
4000	3756.6\\
6000	3533.4\\
8000	3389.5\\
10000	3286.3\\
14000	3164.6\\
16000	3121.1\\
18000	3113.4\\
20000	3090.6\\
26000	3069.9\\
28000	3042.7\\
30000	3026.3\\
32000	3036.1\\
34000	3022\\
36000	3023.6\\
38000	3012.5\\
40000	3029.3\\
44000	3006.5\\
46000	2994.7\\
52000	3023.7\\
56000	3033\\
58000	3042.2\\
60000	3030\\
62000	2997.9\\
64000	2989\\
70000	3008.1\\
74000	3012.3\\
76000	3022.9\\
78000	3020.3\\
80000	3001.7\\
82000	2993.8\\
86000	3003.4\\
88000	3003.8\\
90000	2994\\
92000	2995.8\\
94000	2987.2\\
98000	2996.4\\
1e+05	2985.9\\
};
\addplot [color=mycolor49, forget plot]
  table[row sep=crcr]{%
0	4395.3\\
4000	4556.8\\
8000	4623.9\\
10000	4607.2\\
12000	4548.5\\
14000	4549.2\\
16000	4573.8\\
18000	4557.8\\
20000	4489.8\\
26000	4481.3\\
28000	4486.5\\
30000	4450.1\\
32000	4436.5\\
38000	4446.9\\
40000	4479.8\\
42000	4444.4\\
44000	4460.9\\
46000	4453.3\\
48000	4506.5\\
50000	4503\\
52000	4439.1\\
54000	4411.1\\
56000	4368.4\\
58000	4347.2\\
62000	4384.9\\
64000	4381.2\\
66000	4390.4\\
68000	4437.9\\
70000	4424.2\\
74000	4425.5\\
76000	4415.9\\
78000	4454.8\\
80000	4472.2\\
84000	4475.3\\
86000	4439.5\\
90000	4406.9\\
94000	4492.4\\
96000	4513.3\\
98000	4500.3\\
1e+05	4425.3\\
};
\addplot [color=mycolor50, forget plot]
  table[row sep=crcr]{%
0	4395.3\\
2000	3991.4\\
4000	3473.8\\
6000	3181.6\\
8000	3010.6\\
10000	2878.2\\
14000	2731.6\\
18000	2678.5\\
20000	2663.7\\
22000	2622.1\\
26000	2593.5\\
32000	2588.3\\
34000	2584.7\\
36000	2558.1\\
40000	2554.7\\
44000	2527.9\\
46000	2521.2\\
50000	2544.7\\
54000	2577.6\\
56000	2571.8\\
58000	2584.9\\
62000	2563.5\\
64000	2573.2\\
66000	2554.8\\
68000	2545.4\\
70000	2568.7\\
74000	2528.2\\
76000	2520.4\\
78000	2552.1\\
80000	2560.9\\
82000	2559.4\\
84000	2570.1\\
86000	2566.7\\
88000	2535.1\\
90000	2544.4\\
92000	2566.3\\
94000	2545.1\\
98000	2562.4\\
1e+05	2543.9\\
};
\addplot [color=mycolor51, forget plot]
  table[row sep=crcr]{%
0	4395.3\\
2000	4705.9\\
4000	5190\\
6000	5526.1\\
8000	5725.2\\
10000	5882.9\\
12000	6004\\
16000	5996.6\\
18000	6029.1\\
24000	5976.3\\
26000	5991.4\\
28000	6078.3\\
30000	6106.4\\
32000	6172\\
34000	6217.3\\
36000	6106.1\\
40000	6044.3\\
42000	5982.4\\
44000	5990.3\\
46000	5946.9\\
50000	5954.8\\
52000	5937.6\\
54000	5971.1\\
56000	5956.2\\
60000	5969.8\\
62000	5954.8\\
66000	6019.3\\
68000	5947.6\\
70000	6008.7\\
74000	6007.7\\
78000	5978.2\\
80000	5960.9\\
82000	5970.4\\
84000	5905.1\\
86000	5934.5\\
90000	6028.3\\
92000	6022.2\\
94000	6037.2\\
96000	5966\\
1e+05	6033.1\\
};
\addplot [color=mycolor52, forget plot]
  table[row sep=crcr]{%
0	4395.3\\
2000	4297.9\\
4000	4129.1\\
6000	4009.2\\
8000	3924.1\\
12000	3825.2\\
14000	3804.2\\
16000	3754.7\\
20000	3633.8\\
24000	3691.6\\
26000	3669.6\\
28000	3633.2\\
30000	3608.1\\
32000	3621.1\\
34000	3587.6\\
36000	3584.4\\
38000	3610\\
40000	3617.1\\
42000	3636.6\\
44000	3600.3\\
48000	3614.3\\
50000	3580\\
52000	3596.6\\
54000	3602\\
58000	3590.9\\
62000	3634.4\\
64000	3630\\
66000	3576.6\\
68000	3588.7\\
70000	3624.2\\
72000	3612.3\\
74000	3633.1\\
78000	3589.8\\
80000	3587.1\\
84000	3612.9\\
86000	3625.6\\
88000	3665.3\\
92000	3638.2\\
94000	3610.4\\
96000	3610.7\\
98000	3625.3\\
1e+05	3594.3\\
};
\addplot [color=mycolor53, forget plot]
  table[row sep=crcr]{%
0	4395.3\\
2000	4670.4\\
4000	5090\\
6000	5430\\
8000	5599.1\\
10000	5681.1\\
16000	5846.3\\
18000	5794.3\\
20000	5790.1\\
22000	5747.6\\
26000	5730.4\\
28000	5750.9\\
32000	5718.5\\
34000	5760.9\\
36000	5701.1\\
44000	5643.8\\
48000	5681.8\\
50000	5690.8\\
54000	5639.3\\
56000	5662.3\\
58000	5631.4\\
60000	5645.9\\
64000	5741.7\\
66000	5745\\
68000	5775.3\\
70000	5749.6\\
74000	5773.6\\
78000	5695.9\\
80000	5735.5\\
82000	5710.4\\
84000	5706.7\\
86000	5732\\
88000	5708.8\\
90000	5736.3\\
94000	5695.6\\
98000	5702.8\\
1e+05	5659.9\\
};
\addplot [color=mycolor54, forget plot]
  table[row sep=crcr]{%
0	4395.3\\
2000	4418.4\\
4000	4474.8\\
6000	4555.3\\
10000	4629.3\\
16000	4698.4\\
18000	4649.2\\
22000	4651.5\\
24000	4647.3\\
30000	4686.4\\
34000	4687.4\\
36000	4694.1\\
38000	4663.5\\
40000	4663.5\\
42000	4640.8\\
44000	4633.5\\
46000	4643.8\\
48000	4672\\
50000	4645\\
54000	4647.8\\
56000	4628.8\\
62000	4642.5\\
64000	4653.5\\
66000	4632.5\\
72000	4653.2\\
74000	4684.1\\
80000	4633.5\\
82000	4643.8\\
86000	4630.5\\
88000	4624.5\\
90000	4634.2\\
92000	4624.6\\
96000	4629.8\\
1e+05	4581.8\\
};
\addplot [color=mycolor55, forget plot]
  table[row sep=crcr]{%
0	4395.3\\
2000	4633.7\\
4000	5003.8\\
6000	5246.8\\
8000	5409.5\\
10000	5543.6\\
12000	5649.1\\
14000	5727.3\\
18000	5774.3\\
22000	5769.4\\
30000	5807.5\\
40000	5750.4\\
50000	5737\\
52000	5764.4\\
60000	5751.2\\
62000	5753.8\\
66000	5823.9\\
70000	5827.5\\
72000	5819\\
74000	5786\\
76000	5775.9\\
80000	5833.1\\
82000	5829.1\\
84000	5854.9\\
88000	5843.7\\
94000	5865.2\\
98000	5884.1\\
1e+05	5889.3\\
};
\addplot [color=mycolor56, forget plot]
  table[row sep=crcr]{%
0	4395.3\\
4000	4320.8\\
6000	4307.2\\
10000	4248.3\\
14000	4193.9\\
16000	4166.4\\
18000	4153.8\\
20000	4172.9\\
26000	4165\\
28000	4111.2\\
30000	4081.5\\
32000	4094.1\\
38000	4029.9\\
42000	4058.6\\
44000	4042.8\\
46000	4040.2\\
54000	4092\\
56000	4074.1\\
58000	4069.9\\
62000	4034.9\\
64000	4070.2\\
70000	4095.6\\
74000	4086\\
78000	4066.8\\
82000	4078.3\\
84000	4075.6\\
86000	4048.2\\
90000	4048.7\\
92000	4055.3\\
94000	4043.4\\
98000	4054.8\\
1e+05	4029.7\\
};
\addplot [color=mycolor57, forget plot]
  table[row sep=crcr]{%
0	4395.3\\
2000	4648.8\\
4000	5089.1\\
6000	5433.2\\
8000	5673.8\\
10000	5876.4\\
14000	6199.8\\
16000	6305.3\\
20000	6443\\
24000	6534.1\\
34000	6662.5\\
36000	6659.6\\
40000	6684.6\\
44000	6667.9\\
46000	6680.6\\
48000	6671.3\\
52000	6685.6\\
56000	6699.6\\
58000	6715.8\\
60000	6698.4\\
64000	6698.7\\
66000	6662.7\\
70000	6669.2\\
72000	6701.1\\
76000	6674.2\\
82000	6710.4\\
92000	6687.3\\
1e+05	6711.9\\
};
\addplot [color=mycolor58, forget plot]
  table[row sep=crcr]{%
0	4395.3\\
2000	4308.2\\
4000	4185.9\\
6000	4093.2\\
8000	4014.7\\
10000	3962.3\\
14000	3905.8\\
18000	3796.8\\
22000	3725.4\\
24000	3708.6\\
26000	3704\\
28000	3668\\
30000	3643\\
32000	3687\\
34000	3675.1\\
38000	3676.7\\
40000	3661.7\\
44000	3675.8\\
46000	3641.7\\
48000	3656.4\\
50000	3657\\
52000	3689.1\\
58000	3670.5\\
60000	3651.6\\
62000	3659.6\\
64000	3657.5\\
66000	3629.4\\
68000	3638.5\\
70000	3666.9\\
72000	3683\\
74000	3673.5\\
76000	3645.6\\
78000	3597.4\\
82000	3583.2\\
84000	3624.4\\
86000	3642.7\\
90000	3637.6\\
92000	3641.8\\
94000	3667\\
96000	3650.8\\
1e+05	3679.9\\
};
\addplot [color=mycolor59, forget plot]
  table[row sep=crcr]{%
0	4395.3\\
2000	4511\\
4000	4710.3\\
6000	4839.6\\
8000	4933.3\\
12000	5059\\
14000	5066.8\\
16000	5090.4\\
24000	5104.2\\
26000	5095.7\\
28000	5067.5\\
30000	5075.6\\
32000	5065.5\\
34000	5071.5\\
36000	5103.6\\
38000	5109.1\\
46000	5058.3\\
50000	5102.5\\
52000	5107.7\\
54000	5089.6\\
58000	5112.5\\
60000	5107.9\\
62000	5086.4\\
64000	5086\\
66000	5105.1\\
68000	5099.7\\
70000	5113\\
74000	5108.3\\
76000	5122.4\\
78000	5109.1\\
80000	5131.2\\
82000	5127.2\\
84000	5102.4\\
88000	5104\\
92000	5123.2\\
94000	5126.7\\
96000	5087\\
1e+05	5062.2\\
};
\addplot [color=mycolor60, forget plot]
  table[row sep=crcr]{%
0	4395.3\\
2000	4220\\
4000	3983.9\\
6000	3887.1\\
8000	3818.1\\
10000	3768.8\\
16000	3681.2\\
26000	3622.9\\
28000	3617\\
32000	3630.5\\
36000	3612.5\\
38000	3615.8\\
40000	3604.5\\
42000	3618.8\\
46000	3610.2\\
48000	3598.4\\
50000	3599.3\\
54000	3616.2\\
56000	3612.5\\
58000	3633.1\\
62000	3647.8\\
64000	3623.9\\
66000	3615.5\\
68000	3593.3\\
74000	3617.9\\
76000	3605.4\\
78000	3614.5\\
80000	3600\\
88000	3611\\
90000	3613.5\\
96000	3582.2\\
1e+05	3568.3\\
};
\addplot [color=mycolor61, forget plot]
  table[row sep=crcr]{%
0	4395.3\\
2000	4207.3\\
4000	3955\\
6000	3779.9\\
8000	3656.4\\
10000	3557.1\\
12000	3482.1\\
18000	3358.5\\
24000	3295.6\\
26000	3290.1\\
32000	3227.5\\
36000	3234.2\\
40000	3214.3\\
46000	3204.7\\
50000	3222.5\\
54000	3212.5\\
56000	3213.9\\
60000	3195.1\\
62000	3191.4\\
64000	3197.8\\
66000	3218.2\\
68000	3223.8\\
70000	3205\\
72000	3200.6\\
78000	3223.9\\
80000	3204.2\\
82000	3204.4\\
84000	3194.6\\
88000	3196.7\\
90000	3180.3\\
94000	3193.9\\
98000	3199.8\\
1e+05	3204.9\\
};
\addplot [color=mycolor62, forget plot]
  table[row sep=crcr]{%
0	4395.3\\
2000	4421.2\\
4000	4524.3\\
6000	4645.8\\
8000	4737.4\\
14000	4928.9\\
16000	4938.6\\
18000	4969.4\\
22000	4984.2\\
34000	5020.5\\
38000	5025.7\\
44000	4996.4\\
48000	4985.7\\
54000	5025.3\\
58000	5026.6\\
60000	5020.9\\
62000	5035.9\\
66000	5008.7\\
68000	5007.8\\
74000	5034.1\\
76000	5043.6\\
78000	5024.8\\
84000	5034.5\\
86000	5011.8\\
94000	5041.7\\
98000	5028\\
1e+05	5036.9\\
};
\addplot [color=mycolor63, forget plot]
  table[row sep=crcr]{%
0	4395.3\\
2000	4025.3\\
4000	3532.3\\
6000	3212.6\\
8000	2994.4\\
10000	2873.1\\
12000	2790.2\\
14000	2721.3\\
16000	2676.5\\
20000	2615.1\\
22000	2582.1\\
24000	2563.9\\
28000	2547.1\\
32000	2506.4\\
36000	2504.7\\
38000	2497.9\\
40000	2475.7\\
42000	2476\\
44000	2486.8\\
46000	2487.3\\
52000	2468\\
54000	2475.4\\
58000	2470.8\\
60000	2480.1\\
64000	2472.2\\
74000	2493.3\\
78000	2485.2\\
80000	2492\\
82000	2487.8\\
84000	2472.8\\
88000	2496.8\\
90000	2486.6\\
96000	2486.8\\
98000	2492.6\\
1e+05	2491.3\\
};
\addplot [color=mycolor64, forget plot]
  table[row sep=crcr]{%
0	4395.3\\
2000	4134.5\\
4000	3798.9\\
6000	3599.6\\
8000	3456.5\\
10000	3370.3\\
14000	3290.3\\
16000	3262.5\\
18000	3250.9\\
24000	3182.8\\
26000	3158.6\\
28000	3159.8\\
32000	3140.2\\
34000	3143.6\\
36000	3132.5\\
38000	3133.7\\
40000	3124.1\\
44000	3133.3\\
46000	3140.4\\
50000	3124.7\\
58000	3117.5\\
62000	3085.7\\
66000	3108.4\\
68000	3102.6\\
70000	3116.7\\
74000	3114.8\\
76000	3119.2\\
80000	3105.7\\
82000	3104.7\\
84000	3121.6\\
88000	3120.6\\
92000	3103.6\\
94000	3114.4\\
98000	3105.3\\
1e+05	3117.9\\
};
\addplot [color=mycolor65, forget plot]
  table[row sep=crcr]{%
0	4395.3\\
2000	4051.4\\
4000	3579.4\\
6000	3294\\
8000	3111\\
10000	3000.8\\
12000	2910.2\\
14000	2853.6\\
16000	2819.4\\
18000	2795.7\\
20000	2764.4\\
32000	2652.7\\
36000	2659.9\\
42000	2644.7\\
44000	2620.1\\
46000	2622.8\\
52000	2605.4\\
58000	2604.1\\
60000	2620.8\\
64000	2633.7\\
72000	2628\\
76000	2637.8\\
86000	2609.6\\
90000	2624.9\\
94000	2620.4\\
96000	2637.3\\
1e+05	2637.4\\
};
\addplot [color=mycolor66, forget plot]
  table[row sep=crcr]{%
0	4395.3\\
2000	4224.8\\
4000	3984.2\\
6000	3832\\
8000	3734.9\\
12000	3624.3\\
16000	3570.8\\
18000	3549.2\\
24000	3508.8\\
28000	3521.4\\
32000	3502.6\\
46000	3497.8\\
48000	3490.3\\
52000	3493.8\\
54000	3492.8\\
58000	3505\\
64000	3491.4\\
68000	3494.7\\
78000	3488.4\\
80000	3480.1\\
84000	3491.8\\
88000	3479.4\\
90000	3480.1\\
92000	3490.9\\
96000	3466.1\\
1e+05	3463.7\\
};
\addplot [color=mycolor67, forget plot]
  table[row sep=crcr]{%
0	4395.3\\
2000	4209\\
4000	3964.6\\
6000	3804.3\\
8000	3696.1\\
10000	3602.4\\
14000	3478.4\\
16000	3437.2\\
20000	3394.6\\
24000	3347.6\\
26000	3341.4\\
28000	3311\\
30000	3303.1\\
32000	3309.3\\
40000	3279\\
42000	3263.2\\
46000	3256.2\\
48000	3258.7\\
50000	3245.5\\
52000	3261.3\\
56000	3255.9\\
58000	3245.4\\
60000	3221.4\\
62000	3208\\
66000	3229.7\\
68000	3249.1\\
70000	3250\\
72000	3271\\
76000	3245.5\\
80000	3256.5\\
84000	3246.4\\
88000	3228.9\\
90000	3229.9\\
92000	3220.8\\
94000	3230.6\\
96000	3250.3\\
1e+05	3239.9\\
};
\addplot [color=mycolor67]
  table[row sep=crcr]{%
0	4395.3\\
2000	4082.6\\
4000	3650.7\\
6000	3370\\
8000	3195.1\\
10000	3076.4\\
12000	2988.4\\
14000	2915.2\\
18000	2829.4\\
20000	2808\\
24000	2781.3\\
28000	2752.6\\
40000	2709.4\\
60000	2678.1\\
66000	2673\\
68000	2688.9\\
74000	2701.5\\
76000	2698.2\\
78000	2717.1\\
84000	2693\\
88000	2716.3\\
98000	2691.8\\
1e+05	2689\\
};
\addlegendentry{Classe k=43}

\addplot [color=mycolor68, forget plot]
  table[row sep=crcr]{%
0	4395.3\\
2000	4102.2\\
4000	3659.9\\
6000	3358.9\\
8000	3155.4\\
10000	3006.8\\
12000	2894.9\\
14000	2817.6\\
16000	2760.7\\
18000	2713.3\\
24000	2621\\
26000	2601.5\\
30000	2587.6\\
34000	2584.4\\
36000	2566.8\\
42000	2552.8\\
46000	2559.8\\
48000	2556.6\\
50000	2563.5\\
54000	2554\\
62000	2550.8\\
68000	2530.6\\
70000	2528.8\\
72000	2535\\
76000	2531.4\\
82000	2540.4\\
86000	2549.3\\
88000	2545.8\\
96000	2562.3\\
98000	2561.2\\
1e+05	2551.6\\
};
\addplot [color=mycolor69, forget plot]
  table[row sep=crcr]{%
0	4395.3\\
2000	4176.5\\
4000	3853.4\\
6000	3642\\
8000	3491.2\\
10000	3383.7\\
12000	3306.3\\
14000	3243.5\\
18000	3162.1\\
24000	3054.6\\
30000	3050.7\\
32000	3028.5\\
42000	2993.5\\
46000	2993.2\\
50000	2993.3\\
58000	2973\\
62000	2977.3\\
64000	2975.4\\
68000	2997.3\\
72000	2977\\
74000	2981.2\\
80000	2959.3\\
86000	2986\\
90000	2976.7\\
92000	2987.8\\
96000	2986\\
1e+05	2974.1\\
};
\addplot [color=mycolor70, forget plot]
  table[row sep=crcr]{%
0	4395.3\\
2000	4022.7\\
4000	3494.8\\
6000	3136.4\\
8000	2893\\
10000	2714.6\\
12000	2588.5\\
14000	2494.2\\
16000	2428.1\\
18000	2375.6\\
22000	2309.8\\
24000	2274.8\\
30000	2221.3\\
34000	2217.9\\
38000	2210.7\\
40000	2213.3\\
42000	2207.2\\
46000	2185.1\\
52000	2195.4\\
54000	2184.1\\
56000	2184.2\\
58000	2190.6\\
62000	2187\\
64000	2193.4\\
70000	2188.6\\
72000	2210.2\\
74000	2194.5\\
76000	2186.2\\
80000	2187.2\\
82000	2176.8\\
84000	2173.3\\
88000	2179.5\\
94000	2188.7\\
96000	2182.5\\
1e+05	2192.8\\
};
\addplot [color=mycolor71, forget plot]
  table[row sep=crcr]{%
0	4395.3\\
2000	4080.6\\
4000	3629.9\\
6000	3335.1\\
8000	3135.6\\
10000	2992.8\\
12000	2898.6\\
14000	2821.9\\
16000	2763.9\\
18000	2719.2\\
20000	2687.9\\
26000	2627\\
34000	2572.2\\
46000	2567\\
52000	2558\\
56000	2571.5\\
64000	2548.3\\
68000	2544.1\\
74000	2561.7\\
80000	2573.7\\
84000	2577.5\\
90000	2557.5\\
94000	2561.8\\
98000	2562.7\\
1e+05	2561\\
};
\addplot [color=mycolor72, forget plot]
  table[row sep=crcr]{%
0	4395.3\\
2000	4046.3\\
4000	3548.4\\
6000	3217.6\\
8000	2982.3\\
10000	2826.6\\
12000	2702.2\\
14000	2611.5\\
16000	2534.1\\
18000	2482.8\\
20000	2440.2\\
24000	2396.5\\
28000	2367.7\\
32000	2325.6\\
36000	2314.4\\
40000	2313.8\\
48000	2288.8\\
50000	2288.7\\
52000	2273.1\\
54000	2263.9\\
56000	2266.5\\
58000	2260\\
60000	2278.5\\
62000	2275.6\\
64000	2263.3\\
66000	2271.1\\
68000	2287.4\\
72000	2277.6\\
78000	2259.7\\
90000	2263.3\\
94000	2265.2\\
96000	2266.1\\
98000	2274.1\\
1e+05	2272.1\\
};
\addplot [color=mycolor73, forget plot]
  table[row sep=crcr]{%
0	4395.3\\
2000	4113.7\\
4000	3705\\
6000	3419.6\\
8000	3228.4\\
10000	3093.7\\
12000	2999.9\\
14000	2922.2\\
16000	2866.5\\
18000	2830.3\\
24000	2745.8\\
34000	2701.7\\
36000	2704\\
42000	2688.3\\
66000	2688\\
70000	2679.1\\
74000	2677.9\\
78000	2687.8\\
84000	2678.5\\
86000	2669.6\\
92000	2672.1\\
96000	2670.8\\
1e+05	2679.7\\
};
\addplot [color=mycolor74, forget plot]
  table[row sep=crcr]{%
0	4395.3\\
2000	3928.9\\
4000	3246.8\\
6000	2790.8\\
8000	2482.3\\
10000	2247.3\\
12000	2083.1\\
14000	1969.7\\
16000	1894.2\\
18000	1831\\
20000	1776.8\\
22000	1739.7\\
26000	1691.4\\
32000	1644.6\\
34000	1637\\
38000	1641.2\\
42000	1636.7\\
44000	1630.9\\
50000	1626.8\\
52000	1615.4\\
54000	1609.6\\
66000	1619.2\\
70000	1620.6\\
74000	1621.4\\
76000	1613.7\\
80000	1614\\
84000	1611.2\\
92000	1608.8\\
1e+05	1617.7\\
};
\addplot [color=mycolor75, forget plot]
  table[row sep=crcr]{%
0	4395.3\\
2000	4092.4\\
4000	3634.4\\
6000	3317.9\\
8000	3083.1\\
10000	2905.7\\
12000	2757.1\\
14000	2658.4\\
16000	2580.8\\
20000	2464.4\\
24000	2403.6\\
28000	2366.1\\
34000	2313.6\\
46000	2290.4\\
50000	2293.6\\
52000	2287.8\\
56000	2291\\
76000	2275.1\\
84000	2285\\
92000	2292.6\\
96000	2298.1\\
1e+05	2294.8\\
};
\addplot [color=mycolor76, forget plot]
  table[row sep=crcr]{%
0	4395.3\\
2000	4056.8\\
4000	3538\\
6000	3163.7\\
8000	2905.4\\
10000	2711.1\\
12000	2565.7\\
14000	2458.6\\
16000	2378.5\\
18000	2322.6\\
22000	2245.1\\
26000	2189.6\\
30000	2155.3\\
34000	2142.4\\
38000	2117\\
42000	2102.4\\
48000	2100.4\\
50000	2096.4\\
52000	2098.6\\
54000	2094\\
56000	2098.3\\
62000	2091.2\\
66000	2087.2\\
68000	2096.9\\
76000	2091.1\\
78000	2081.6\\
82000	2085.2\\
86000	2081.7\\
88000	2087.8\\
92000	2086.9\\
94000	2081.6\\
98000	2087.2\\
1e+05	2082.2\\
};
\addplot [color=mycolor77, forget plot]
  table[row sep=crcr]{%
0	4395.3\\
2000	3974.6\\
4000	3345.7\\
6000	2900.2\\
8000	2584.1\\
10000	2354.5\\
12000	2181.4\\
14000	2052\\
16000	1965.4\\
18000	1891.8\\
20000	1837.4\\
24000	1758.5\\
26000	1730.4\\
36000	1666.3\\
40000	1655.8\\
42000	1651.5\\
46000	1653.6\\
50000	1644.2\\
52000	1645\\
54000	1640.1\\
56000	1640.5\\
58000	1647.4\\
64000	1648.8\\
66000	1640.7\\
70000	1642.4\\
72000	1640.9\\
76000	1652.5\\
80000	1642.8\\
96000	1643\\
98000	1645.9\\
1e+05	1640.7\\
};
\addplot [color=mycolor78, forget plot]
  table[row sep=crcr]{%
0	4395.3\\
2000	4113.5\\
4000	3698.9\\
6000	3380.1\\
8000	3156.1\\
10000	2981.9\\
12000	2857.3\\
16000	2668.4\\
18000	2607.3\\
20000	2559.9\\
22000	2524.3\\
30000	2422.1\\
34000	2400.3\\
40000	2407\\
42000	2396.8\\
44000	2398.2\\
46000	2382.7\\
50000	2373.7\\
54000	2358.7\\
56000	2365.1\\
62000	2361.3\\
66000	2371.8\\
68000	2365.1\\
72000	2367.6\\
76000	2353.5\\
84000	2362.3\\
88000	2360.1\\
92000	2360.1\\
1e+05	2356.7\\
};
\addplot [color=mycolor79, forget plot]
  table[row sep=crcr]{%
0	4395.3\\
2000	4081\\
4000	3612\\
6000	3276.8\\
8000	3036\\
10000	2849.3\\
12000	2719.5\\
14000	2618.1\\
16000	2538.9\\
20000	2433.5\\
22000	2393.2\\
26000	2342.9\\
28000	2317.9\\
36000	2264.9\\
38000	2255.1\\
40000	2251.7\\
42000	2254.6\\
44000	2246.5\\
46000	2250.2\\
50000	2237.2\\
52000	2243.4\\
56000	2232\\
66000	2241.9\\
80000	2223.7\\
82000	2239\\
84000	2246.8\\
90000	2242.4\\
94000	2249.7\\
1e+05	2246.6\\
};
\addplot [color=mycolor79, forget plot]
  table[row sep=crcr]{%
0	4395.3\\
2000	4105.3\\
4000	3648.8\\
6000	3291.5\\
8000	3027.6\\
10000	2826.2\\
12000	2672.5\\
14000	2561.7\\
16000	2475.9\\
18000	2401.7\\
22000	2292.5\\
24000	2247.5\\
26000	2217.6\\
32000	2156.7\\
36000	2138.7\\
40000	2128.3\\
42000	2115.7\\
44000	2112.7\\
46000	2103.5\\
48000	2104.9\\
52000	2093.7\\
54000	2091.2\\
56000	2100.7\\
58000	2099.6\\
62000	2109.5\\
66000	2100.8\\
70000	2109.4\\
78000	2101.2\\
80000	2105\\
86000	2084.8\\
88000	2071.1\\
90000	2068\\
96000	2093.7\\
1e+05	2081.6\\
};
\addplot [color=mycolor80, forget plot]
  table[row sep=crcr]{%
0	4395.3\\
2000	4013.4\\
4000	3423.1\\
6000	2998.3\\
8000	2675.2\\
10000	2438.1\\
12000	2256.5\\
14000	2114.3\\
16000	2009.3\\
18000	1928.8\\
22000	1818.9\\
24000	1786.5\\
30000	1714.9\\
34000	1688.8\\
38000	1669.9\\
54000	1635.7\\
62000	1635.6\\
68000	1623.6\\
78000	1613.9\\
84000	1623.3\\
88000	1618.8\\
90000	1627.8\\
98000	1623.9\\
1e+05	1622.8\\
};
\addplot [color=mycolor81, forget plot]
  table[row sep=crcr]{%
0	4395.3\\
2000	3986.9\\
4000	3379.7\\
6000	2941\\
8000	2608.8\\
10000	2360.1\\
12000	2173.6\\
14000	2032.6\\
16000	1928.3\\
18000	1841.7\\
20000	1779.9\\
22000	1730.4\\
26000	1656.2\\
30000	1607.1\\
32000	1588.5\\
36000	1566\\
46000	1539\\
52000	1528.6\\
58000	1525.3\\
62000	1517.8\\
66000	1529.1\\
72000	1531.7\\
78000	1529.2\\
82000	1525.9\\
84000	1531\\
88000	1527.4\\
92000	1528.2\\
96000	1521.8\\
1e+05	1527.1\\
};
\addplot [color=mycolor82, forget plot]
  table[row sep=crcr]{%
0	4395.3\\
2000	4034.6\\
4000	3479.4\\
6000	3066.5\\
8000	2763\\
10000	2542.4\\
12000	2374.9\\
14000	2241\\
16000	2141.8\\
18000	2063.1\\
20000	2003.8\\
22000	1958\\
24000	1920.6\\
28000	1865.1\\
30000	1851\\
32000	1829.9\\
38000	1803.6\\
40000	1791.8\\
44000	1783.9\\
48000	1786\\
52000	1779.4\\
90000	1771.9\\
92000	1775.9\\
94000	1774.1\\
98000	1779.8\\
1e+05	1775.1\\
};
\addplot [color=mycolor83, forget plot]
  table[row sep=crcr]{%
0	4395.3\\
2000	4063.2\\
4000	3545.7\\
6000	3162.8\\
8000	2872.7\\
10000	2647.4\\
12000	2471.7\\
14000	2334.6\\
16000	2229.1\\
18000	2148.1\\
20000	2085.6\\
22000	2032.9\\
26000	1955\\
30000	1902.4\\
34000	1865.3\\
38000	1836.8\\
42000	1819.8\\
54000	1808.6\\
56000	1813.1\\
64000	1807.4\\
70000	1799.5\\
72000	1806.4\\
80000	1801.7\\
84000	1803.1\\
86000	1805.7\\
90000	1800.9\\
94000	1795.8\\
98000	1804.1\\
1e+05	1807.6\\
};
\addplot [color=mycolor83, forget plot]
  table[row sep=crcr]{%
0	4395.3\\
2000	3970.7\\
4000	3322.8\\
6000	2842.6\\
8000	2484.1\\
10000	2223.1\\
12000	2029.5\\
14000	1872.6\\
16000	1750.6\\
18000	1663.2\\
20000	1598.4\\
22000	1544.3\\
24000	1498.7\\
28000	1436.1\\
32000	1395.7\\
34000	1380.9\\
36000	1379.3\\
46000	1347.3\\
50000	1339.2\\
52000	1341.9\\
62000	1336\\
64000	1339.5\\
72000	1331.5\\
76000	1335.9\\
78000	1336.5\\
82000	1329.3\\
94000	1318.2\\
96000	1323.6\\
1e+05	1319.7\\
};
\addplot [color=mycolor84, forget plot]
  table[row sep=crcr]{%
0	4395.3\\
2000	4027\\
4000	3427.2\\
6000	2950.5\\
8000	2611.5\\
10000	2344.5\\
12000	2125.4\\
14000	1959.3\\
16000	1830.4\\
18000	1729.3\\
20000	1662.4\\
22000	1603.4\\
26000	1511.7\\
28000	1480.2\\
30000	1460.3\\
32000	1436.1\\
34000	1408\\
38000	1381.8\\
40000	1377.1\\
44000	1372.3\\
46000	1361.7\\
50000	1351.7\\
52000	1357.3\\
54000	1358.7\\
56000	1356.1\\
62000	1360.1\\
64000	1359.6\\
66000	1355\\
70000	1355.2\\
74000	1351.2\\
76000	1350\\
80000	1361.8\\
84000	1362.1\\
86000	1358.4\\
88000	1359.1\\
92000	1349.5\\
96000	1344.7\\
1e+05	1351\\
};
\addplot [color=mycolor85, forget plot]
  table[row sep=crcr]{%
0	4395.3\\
2000	4007.1\\
4000	3408.8\\
6000	2957\\
8000	2618.5\\
10000	2369.3\\
12000	2182.2\\
14000	2032.7\\
16000	1922.3\\
18000	1832.5\\
20000	1761.9\\
24000	1653.2\\
26000	1607.5\\
30000	1553.5\\
36000	1505.3\\
44000	1484.8\\
48000	1482.1\\
50000	1486.4\\
52000	1486.3\\
58000	1465.9\\
64000	1466\\
70000	1469\\
74000	1467\\
84000	1456.3\\
90000	1464.5\\
92000	1461\\
96000	1464\\
98000	1471.9\\
1e+05	1475.7\\
};
\addplot [color=mycolor86, forget plot]
  table[row sep=crcr]{%
0	4395.3\\
2000	4025.8\\
4000	3432.9\\
6000	2981.3\\
8000	2619.7\\
10000	2333.4\\
12000	2109.5\\
14000	1935.3\\
16000	1797.5\\
18000	1691.7\\
20000	1603.5\\
22000	1534.1\\
26000	1436.4\\
28000	1403.2\\
32000	1355.7\\
34000	1338.7\\
36000	1312.8\\
40000	1287.9\\
48000	1276.1\\
50000	1277.4\\
54000	1261.5\\
56000	1261.7\\
58000	1257\\
60000	1247.5\\
62000	1255.2\\
64000	1255.9\\
66000	1248.5\\
68000	1253\\
70000	1261.4\\
78000	1249.7\\
82000	1249.4\\
84000	1240.7\\
90000	1238.1\\
94000	1237.4\\
98000	1231.8\\
1e+05	1234\\
};
\addplot [color=mycolor87, forget plot]
  table[row sep=crcr]{%
0	4395.3\\
2000	3963.4\\
4000	3319.4\\
6000	2823.4\\
8000	2455.7\\
10000	2162.5\\
12000	1936.3\\
14000	1770.1\\
16000	1629.4\\
18000	1506.4\\
20000	1414.1\\
22000	1358.6\\
24000	1296.2\\
26000	1250.7\\
28000	1217\\
30000	1191.6\\
32000	1171.1\\
36000	1122.9\\
38000	1100.8\\
44000	1075.5\\
48000	1047.6\\
50000	1040.8\\
54000	1036.7\\
58000	1037\\
62000	1055\\
64000	1057.4\\
66000	1056\\
68000	1062.5\\
74000	1062.8\\
76000	1056\\
78000	1062.7\\
80000	1064.3\\
90000	1055\\
92000	1058.2\\
96000	1058.1\\
98000	1064.7\\
1e+05	1067.3\\
};
\addplot [color=mycolor88, forget plot]
  table[row sep=crcr]{%
0	4395.3\\
2000	4056.3\\
4000	3501.7\\
6000	3067.3\\
8000	2714.5\\
10000	2427.9\\
12000	2207.2\\
14000	2031\\
16000	1887.3\\
18000	1772.9\\
20000	1672.6\\
22000	1587.6\\
24000	1528.3\\
26000	1476\\
28000	1440.2\\
32000	1385.9\\
34000	1362.2\\
40000	1308.1\\
42000	1295.7\\
44000	1288.8\\
46000	1286\\
58000	1246.9\\
62000	1247.3\\
64000	1249.2\\
68000	1242.5\\
76000	1239.5\\
78000	1242.8\\
86000	1241.6\\
88000	1242.4\\
90000	1238.9\\
92000	1240.7\\
96000	1248.7\\
1e+05	1250.9\\
};
\addplot [color=mycolor89, forget plot]
  table[row sep=crcr]{%
0	4395.3\\
2000	3987.4\\
4000	3332.9\\
6000	2824.6\\
8000	2440\\
10000	2131.2\\
12000	1880.6\\
14000	1691.8\\
16000	1540.6\\
18000	1421.8\\
20000	1323.3\\
22000	1240.5\\
24000	1172.1\\
26000	1122.7\\
30000	1055.1\\
34000	1007.3\\
38000	972.32\\
42000	951.94\\
46000	936.4\\
48000	930.84\\
50000	921.97\\
56000	912.04\\
58000	911.72\\
60000	915.2\\
66000	909.71\\
68000	911.32\\
70000	907.67\\
72000	912.22\\
74000	912.32\\
76000	908.47\\
84000	910.78\\
88000	916.49\\
96000	914.6\\
1e+05	916.51\\
};
\addplot [color=mycolor90, forget plot]
  table[row sep=crcr]{%
0	4395.3\\
2000	3991.6\\
4000	3350\\
6000	2848.4\\
8000	2449.7\\
10000	2144.7\\
12000	1902\\
14000	1707.1\\
16000	1552.7\\
18000	1426.6\\
20000	1323.3\\
22000	1242.1\\
24000	1172.8\\
26000	1114.9\\
28000	1070.3\\
30000	1039.3\\
32000	1013.9\\
34000	992.24\\
36000	974.76\\
38000	960.46\\
40000	952.58\\
44000	929.46\\
50000	905.74\\
54000	899.97\\
58000	898.38\\
62000	885.19\\
64000	882.26\\
68000	879.88\\
72000	879.32\\
76000	880.84\\
84000	873.57\\
88000	872.87\\
92000	868.01\\
94000	868.69\\
1e+05	880.89\\
};
\addplot [color=mycolor91, forget plot]
  table[row sep=crcr]{%
0	4395.3\\
2000	4007\\
4000	3372.3\\
6000	2865.8\\
8000	2458.8\\
10000	2128.1\\
12000	1862\\
14000	1656.6\\
16000	1500.7\\
18000	1363.9\\
20000	1253.8\\
22000	1168.4\\
24000	1098.6\\
26000	1038.3\\
28000	987.98\\
30000	947.1\\
32000	911.92\\
36000	859.81\\
38000	840.69\\
44000	810.29\\
50000	790.84\\
52000	788.65\\
56000	777.49\\
58000	777.27\\
60000	771.59\\
70000	770.57\\
72000	772.77\\
74000	768.76\\
78000	765.92\\
80000	771.1\\
82000	767.71\\
84000	768.12\\
86000	761.2\\
94000	765.23\\
98000	768.13\\
1e+05	765.64\\
};
\addplot [color=mycolor92, forget plot]
  table[row sep=crcr]{%
0	4395.3\\
2000	4051.3\\
4000	3481.5\\
6000	3021\\
8000	2650.8\\
10000	2350.6\\
12000	2106.8\\
14000	1904.7\\
16000	1742.9\\
18000	1609.6\\
20000	1493.2\\
22000	1400.3\\
24000	1328.4\\
26000	1265.4\\
28000	1213\\
30000	1169.8\\
36000	1079.4\\
38000	1055\\
40000	1036.3\\
44000	1011.6\\
46000	999\\
62000	956.92\\
66000	955.66\\
70000	953.1\\
82000	945.69\\
84000	950.28\\
86000	951.36\\
92000	941.91\\
96000	945.9\\
98000	942.48\\
1e+05	946.18\\
};
\addplot [color=mycolor93, forget plot]
  table[row sep=crcr]{%
0	4395.3\\
2000	4043.8\\
4000	3451.7\\
6000	2977\\
8000	2596.9\\
10000	2292.7\\
12000	2037.4\\
14000	1828.8\\
16000	1659.7\\
18000	1521.4\\
20000	1405\\
22000	1305.1\\
24000	1220.5\\
26000	1150.5\\
30000	1041.2\\
32000	1000.4\\
34000	970.24\\
38000	925.76\\
40000	900.49\\
42000	880.65\\
46000	856.08\\
48000	846.92\\
56000	831.67\\
58000	827.33\\
62000	826.51\\
64000	820.69\\
66000	818.88\\
70000	822.39\\
74000	822.74\\
76000	814.64\\
78000	809.62\\
80000	811.62\\
82000	818.41\\
84000	816.37\\
86000	808.6\\
90000	816.95\\
92000	817.58\\
96000	812.04\\
98000	808.33\\
1e+05	809.99\\
};
\addplot [color=mycolor94, forget plot]
  table[row sep=crcr]{%
0	4395.3\\
2000	4061.9\\
4000	3506.9\\
6000	3048.7\\
8000	2660\\
10000	2346.6\\
12000	2087.9\\
14000	1871.3\\
16000	1690.8\\
18000	1539.1\\
20000	1414.2\\
22000	1311\\
24000	1223.2\\
26000	1150.6\\
28000	1087.3\\
30000	1037.2\\
32000	996.14\\
36000	932.78\\
38000	906.74\\
42000	864.4\\
44000	844.44\\
46000	829.91\\
54000	796.15\\
56000	790.44\\
68000	773.4\\
72000	772.05\\
78000	765.38\\
80000	762.61\\
82000	765.34\\
86000	764.15\\
90000	762.89\\
94000	767.93\\
98000	767.95\\
1e+05	764.72\\
};
\addplot [color=mycolor95, forget plot]
  table[row sep=crcr]{%
0	4395.3\\
2000	4009.7\\
4000	3374\\
6000	2853.3\\
8000	2424.4\\
10000	2078.4\\
12000	1791.2\\
14000	1565.2\\
16000	1379.2\\
18000	1228.8\\
20000	1103\\
22000	994.9\\
24000	907.99\\
26000	832.72\\
30000	719.89\\
32000	679.26\\
34000	644.73\\
38000	588.93\\
42000	548.03\\
46000	518.41\\
52000	497.39\\
54000	493.07\\
56000	486.62\\
62000	481.51\\
66000	474.77\\
70000	476.38\\
76000	468.93\\
80000	466.05\\
82000	462.12\\
88000	463.98\\
90000	462.08\\
92000	461.48\\
94000	462.44\\
98000	455.14\\
1e+05	454.86\\
};
\addplot [color=mycolor96, forget plot]
  table[row sep=crcr]{%
0	4395.3\\
2000	4029.4\\
4000	3454.5\\
8000	2591.9\\
10000	2249.3\\
12000	1969.4\\
14000	1732.6\\
16000	1531.5\\
18000	1367.2\\
20000	1232\\
22000	1117.6\\
24000	1023.1\\
28000	868.01\\
30000	817.36\\
32000	772.11\\
36000	694.51\\
40000	641.56\\
46000	587.85\\
48000	577.21\\
50000	562.13\\
52000	550.74\\
54000	541.14\\
60000	520.77\\
62000	515.25\\
64000	508.29\\
66000	503.54\\
68000	496.8\\
70000	493.86\\
72000	497.36\\
74000	493.4\\
78000	494.29\\
80000	489.48\\
82000	486.48\\
84000	486.14\\
86000	481.76\\
92000	483.08\\
96000	489.38\\
98000	485.28\\
1e+05	484.34\\
};
\addplot [color=mycolor97, forget plot]
  table[row sep=crcr]{%
0	4395.3\\
2000	4038.5\\
4000	3423.6\\
6000	2914.7\\
8000	2496.1\\
10000	2148.9\\
12000	1857.7\\
16000	1412.2\\
18000	1241.3\\
20000	1100\\
22000	984.12\\
24000	885.16\\
26000	804.3\\
30000	679.47\\
32000	631.7\\
34000	591.28\\
36000	555.67\\
38000	526.54\\
42000	480.76\\
44000	462.35\\
46000	447.39\\
52000	410.12\\
56000	396.79\\
58000	391.5\\
62000	377.12\\
64000	373.69\\
66000	368.42\\
68000	364.89\\
70000	365.19\\
74000	362.75\\
76000	356.25\\
78000	354.62\\
82000	347.55\\
84000	344.72\\
86000	344.53\\
88000	342.64\\
90000	343.68\\
94000	342.84\\
96000	345.29\\
1e+05	344.95\\
};
\addplot [color=mycolor98]
  table[row sep=crcr]{%
0	4395.3\\
2000	4106.8\\
4000	3608.7\\
6000	3160.1\\
10000	2473.6\\
12000	2193.1\\
16000	1747.6\\
18000	1570.1\\
20000	1421.5\\
22000	1293.8\\
24000	1185.8\\
28000	1013.5\\
30000	943.19\\
32000	883.63\\
36000	781.7\\
40000	711.74\\
42000	681.32\\
44000	656.35\\
46000	634.63\\
48000	617.19\\
52000	577.38\\
54000	566.41\\
56000	552.58\\
58000	537.26\\
64000	511.56\\
68000	498.55\\
72000	488.11\\
74000	483.13\\
78000	477.18\\
82000	472.45\\
88000	470.83\\
90000	473.11\\
98000	463.82\\
1e+05	465.25\\
};
\addlegendentry{Classe k=10}

            % \input{../../../../.tex/20/KS/SizeEvolutionsfig2.tex}
        \end{groupplot}

        \begin{groupplot}[
            group style={
                group name=Row2,
                plotGroup,
                group size=2 by 1,
            },
            plotRow,
            row5Keys,
            xTickScaleLabelEmpty,
            plotLegendSE,
            %
            xmin=0,xmax=1e+05,
            ymin=100,ymax=1.5151e+05,
        ]
            \nextgroupplot[%
                at={($(Row1 c1r1.south west)-(0,\yPlotSep em)-(0,\plotHeight cm)$)},
                % xmin=0,xmax=1e+05,
                % ymin=100,ymax=1.4995e+05,
            ] % This file was created by matlab2tikz.
%
\definecolor{mycolor1}{rgb}{0.00146,0.00047,0.01387}%
\definecolor{mycolor2}{rgb}{0.01166,0.00942,0.06346}%
\definecolor{mycolor3}{rgb}{0.02943,0.02150,0.11462}%
\definecolor{mycolor4}{rgb}{0.06134,0.03659,0.17764}%
\definecolor{mycolor5}{rgb}{0.08741,0.04456,0.22481}%
\definecolor{mycolor6}{rgb}{0.12291,0.04754,0.28162}%
\definecolor{mycolor7}{rgb}{0.14907,0.04547,0.31709}%
\definecolor{mycolor8}{rgb}{0.18343,0.04033,0.35497}%
\definecolor{mycolor9}{rgb}{0.21795,0.03662,0.38352}%
\definecolor{mycolor10}{rgb}{0.24497,0.03705,0.40001}%
\definecolor{mycolor11}{rgb}{0.27135,0.04092,0.41198}%
\definecolor{mycolor12}{rgb}{0.29076,0.04564,0.41864}%
\definecolor{mycolor13}{rgb}{0.31628,0.05349,0.42512}%
\definecolor{mycolor14}{rgb}{0.33522,0.06006,0.42852}%
\definecolor{mycolor15}{rgb}{0.36028,0.06925,0.43150}%
\definecolor{mycolor16}{rgb}{0.37900,0.07625,0.43272}%
\definecolor{mycolor17}{rgb}{0.39767,0.08326,0.43318}%
\definecolor{mycolor18}{rgb}{0.41633,0.09020,0.43294}%
\definecolor{mycolor19}{rgb}{0.42877,0.09479,0.43241}%
\definecolor{mycolor20}{rgb}{0.44743,0.10160,0.43108}%
\definecolor{mycolor21}{rgb}{0.46610,0.10832,0.42912}%
\definecolor{mycolor22}{rgb}{0.47856,0.11276,0.42747}%
\definecolor{mycolor23}{rgb}{0.49726,0.11938,0.42449}%
\definecolor{mycolor24}{rgb}{0.50973,0.12377,0.42216}%
\definecolor{mycolor25}{rgb}{0.52221,0.12815,0.41955}%
\definecolor{mycolor26}{rgb}{0.53468,0.13253,0.41667}%
\definecolor{mycolor27}{rgb}{0.55339,0.13913,0.41183}%
\definecolor{mycolor28}{rgb}{0.56585,0.14357,0.40826}%
\definecolor{mycolor29}{rgb}{0.57830,0.14804,0.40441}%
\definecolor{mycolor30}{rgb}{0.59073,0.15256,0.40029}%
\definecolor{mycolor31}{rgb}{0.60314,0.15715,0.39589}%
\definecolor{mycolor32}{rgb}{0.61551,0.16182,0.39122}%
\definecolor{mycolor33}{rgb}{0.62169,0.16418,0.38878}%
\definecolor{mycolor34}{rgb}{0.63400,0.16899,0.38370}%
\definecolor{mycolor35}{rgb}{0.64626,0.17391,0.37836}%
\definecolor{mycolor36}{rgb}{0.65846,0.17896,0.37275}%
\definecolor{mycolor37}{rgb}{0.66454,0.18154,0.36985}%
\definecolor{mycolor38}{rgb}{0.67664,0.18681,0.36385}%
\definecolor{mycolor39}{rgb}{0.68865,0.19224,0.35760}%
\definecolor{mycolor40}{rgb}{0.69463,0.19502,0.35439}%
\definecolor{mycolor41}{rgb}{0.70650,0.20073,0.34778}%
\definecolor{mycolor42}{rgb}{0.71240,0.20366,0.34438}%
\definecolor{mycolor43}{rgb}{0.72410,0.20967,0.33742}%
\definecolor{mycolor44}{rgb}{0.72991,0.21276,0.33386}%
\definecolor{mycolor45}{rgb}{0.74142,0.21911,0.32658}%
\definecolor{mycolor46}{rgb}{0.74713,0.22238,0.32286}%
\definecolor{mycolor47}{rgb}{0.75279,0.22571,0.31909}%
\definecolor{mycolor48}{rgb}{0.76401,0.23255,0.31140}%
\definecolor{mycolor49}{rgb}{0.76956,0.23608,0.30749}%
\definecolor{mycolor50}{rgb}{0.77506,0.23967,0.30353}%
\definecolor{mycolor51}{rgb}{0.79129,0.25086,0.29139}%
\definecolor{mycolor52}{rgb}{0.79661,0.25473,0.28726}%
\definecolor{mycolor53}{rgb}{0.80187,0.25867,0.28310}%
\definecolor{mycolor54}{rgb}{0.81224,0.26679,0.27466}%
\definecolor{mycolor55}{rgb}{0.81734,0.27095,0.27039}%
\definecolor{mycolor56}{rgb}{0.82737,0.27952,0.26175}%
\definecolor{mycolor57}{rgb}{0.83230,0.28391,0.25738}%
\definecolor{mycolor58}{rgb}{0.83717,0.28839,0.25299}%
\definecolor{mycolor59}{rgb}{0.84671,0.29756,0.24411}%
\definecolor{mycolor60}{rgb}{0.85138,0.30226,0.23964}%
\definecolor{mycolor61}{rgb}{0.85599,0.30704,0.23513}%
\definecolor{mycolor62}{rgb}{0.86053,0.31189,0.23061}%
\definecolor{mycolor63}{rgb}{0.87374,0.32691,0.21689}%
\definecolor{mycolor64}{rgb}{0.87800,0.33206,0.21227}%
\definecolor{mycolor65}{rgb}{0.89034,0.34796,0.19829}%
\definecolor{mycolor66}{rgb}{0.89819,0.35891,0.18886}%
\definecolor{mycolor67}{rgb}{0.90200,0.36449,0.18412}%
\definecolor{mycolor68}{rgb}{0.90573,0.37014,0.17935}%
\definecolor{mycolor69}{rgb}{0.90939,0.37586,0.17456}%
\definecolor{mycolor70}{rgb}{0.91297,0.38164,0.16975}%
\definecolor{mycolor71}{rgb}{0.91646,0.38748,0.16492}%
\definecolor{mycolor72}{rgb}{0.91988,0.39339,0.16007}%
\definecolor{mycolor73}{rgb}{0.92322,0.39936,0.15519}%
\definecolor{mycolor74}{rgb}{0.92647,0.40539,0.15029}%
\definecolor{mycolor75}{rgb}{0.92964,0.41148,0.14537}%
\definecolor{mycolor76}{rgb}{0.93274,0.41763,0.14042}%
\definecolor{mycolor77}{rgb}{0.93575,0.42383,0.13544}%
\definecolor{mycolor78}{rgb}{0.93868,0.43009,0.13044}%
\definecolor{mycolor79}{rgb}{0.94152,0.43640,0.12541}%
\definecolor{mycolor80}{rgb}{0.94429,0.44277,0.12035}%
\definecolor{mycolor81}{rgb}{0.94956,0.45566,0.11016}%
\definecolor{mycolor82}{rgb}{0.95208,0.46218,0.10503}%
\definecolor{mycolor83}{rgb}{0.96129,0.48872,0.08429}%
\definecolor{mycolor84}{rgb}{0.96540,0.50225,0.07386}%
\definecolor{mycolor85}{rgb}{0.96732,0.50908,0.06866}%
\definecolor{mycolor86}{rgb}{0.97259,0.52980,0.05332}%
\definecolor{mycolor87}{rgb}{0.97568,0.54380,0.04362}%
\definecolor{mycolor88}{rgb}{0.98082,0.57221,0.02851}%
\definecolor{mycolor89}{rgb}{0.98189,0.57939,0.02625}%
\definecolor{mycolor90}{rgb}{0.98378,0.59385,0.02377}%
\definecolor{mycolor91}{rgb}{0.98532,0.60842,0.02420}%
\definecolor{mycolor92}{rgb}{0.98696,0.63048,0.03091}%
\definecolor{mycolor93}{rgb}{0.98793,0.66025,0.05175}%
\definecolor{mycolor94}{rgb}{0.98771,0.68281,0.07249}%
\definecolor{mycolor95}{rgb}{0.97750,0.78226,0.18592}%
\definecolor{mycolor96}{rgb}{0.96624,0.83619,0.26153}%
\definecolor{mycolor97}{rgb}{0.96252,0.85148,0.28555}%
\definecolor{mycolor98}{rgb}{0.98836,0.99836,0.64492}%
%
\addplot [color=mycolor1]
  table[row sep=crcr]{%
0	4395.3\\
2000	69720\\
4000	1.2164e+05\\
6000	1.3492e+05\\
8000	1.4333e+05\\
10000	1.4528e+05\\
12000	1.4635e+05\\
14000	1.4586e+05\\
16000	1.4747e+05\\
20000	1.4566e+05\\
22000	1.4664e+05\\
24000	1.4684e+05\\
26000	1.4515e+05\\
28000	1.4419e+05\\
30000	1.4512e+05\\
32000	1.4558e+05\\
34000	1.4521e+05\\
36000	1.4549e+05\\
38000	1.4856e+05\\
40000	1.4589e+05\\
42000	1.46e+05\\
44000	1.4808e+05\\
48000	1.4824e+05\\
50000	1.486e+05\\
54000	1.4532e+05\\
56000	1.4494e+05\\
58000	1.4413e+05\\
62000	1.4813e+05\\
64000	1.4609e+05\\
66000	1.4794e+05\\
68000	1.4817e+05\\
70000	1.4752e+05\\
72000	1.4764e+05\\
74000	1.4843e+05\\
76000	1.4744e+05\\
78000	1.4711e+05\\
80000	1.4831e+05\\
82000	1.4731e+05\\
84000	1.4452e+05\\
86000	1.4634e+05\\
88000	1.4667e+05\\
90000	1.478e+05\\
94000	1.4872e+05\\
96000	1.4995e+05\\
1e+05	1.4855e+05\\
};
\addlegendentry{Classe k=283}

\addplot [color=mycolor2, forget plot]
  table[row sep=crcr]{%
0	4395.3\\
2000	22158\\
4000	26065\\
6000	26205\\
8000	25146\\
10000	25747\\
12000	25399\\
14000	26150\\
16000	25410\\
18000	25656\\
20000	25399\\
22000	25290\\
24000	25265\\
26000	25569\\
28000	25397\\
30000	25748\\
36000	25498\\
38000	24942\\
40000	25327\\
42000	25491\\
44000	25256\\
46000	24703\\
48000	25416\\
50000	25325\\
52000	25753\\
54000	25389\\
56000	25390\\
58000	25622\\
62000	24469\\
64000	25254\\
66000	25358\\
68000	25005\\
70000	25501\\
72000	24700\\
74000	24785\\
76000	24740\\
78000	25684\\
80000	24707\\
82000	24824\\
84000	25495\\
86000	25142\\
88000	25581\\
90000	25069\\
92000	25659\\
94000	25135\\
96000	25554\\
98000	25422\\
1e+05	25416\\
};
\addplot [color=mycolor3, forget plot]
  table[row sep=crcr]{%
0	4395.3\\
2000	29474\\
4000	37811\\
6000	38124\\
8000	37952\\
10000	38371\\
12000	38985\\
14000	37954\\
16000	39462\\
18000	39636\\
20000	39061\\
22000	38655\\
24000	38572\\
26000	39148\\
28000	39143\\
30000	38000\\
32000	38388\\
34000	38590\\
36000	38345\\
38000	37523\\
40000	38746\\
42000	37844\\
44000	38254\\
46000	38805\\
48000	38607\\
50000	39333\\
52000	38350\\
54000	37675\\
58000	38920\\
60000	38214\\
62000	38820\\
64000	38562\\
66000	38074\\
68000	38789\\
70000	38433\\
72000	38738\\
74000	39444\\
76000	38896\\
78000	39018\\
80000	38927\\
82000	38456\\
84000	38377\\
86000	38088\\
88000	38864\\
92000	39252\\
94000	38311\\
96000	39224\\
98000	39001\\
1e+05	38520\\
};
\addplot [color=mycolor4, forget plot]
  table[row sep=crcr]{%
0	4395.3\\
2000	14006\\
4000	14925\\
6000	14401\\
10000	14025\\
12000	14066\\
16000	13832\\
18000	14049\\
20000	14153\\
22000	13763\\
24000	14105\\
26000	14076\\
28000	14122\\
30000	14006\\
32000	14327\\
34000	14301\\
36000	14086\\
38000	14392\\
40000	14086\\
42000	14282\\
46000	14130\\
48000	14330\\
50000	14335\\
52000	13885\\
54000	14467\\
56000	14198\\
58000	14097\\
60000	14128\\
62000	13670\\
64000	14238\\
66000	13959\\
70000	14302\\
72000	14527\\
74000	14191\\
76000	14229\\
78000	13897\\
80000	14273\\
82000	14048\\
84000	14251\\
86000	13940\\
88000	14244\\
90000	14032\\
92000	14039\\
94000	13807\\
96000	13881\\
98000	13791\\
1e+05	14016\\
};
\addplot [color=mycolor5, forget plot]
  table[row sep=crcr]{%
0	4395.3\\
2000	34575\\
4000	56739\\
6000	62999\\
8000	63073\\
10000	63436\\
12000	63418\\
14000	62998\\
16000	61638\\
20000	61329\\
24000	61645\\
28000	61084\\
32000	62426\\
34000	62183\\
36000	62925\\
38000	60106\\
40000	60124\\
42000	60501\\
44000	60275\\
46000	60284\\
50000	61473\\
52000	60809\\
54000	62535\\
56000	63607\\
58000	61685\\
60000	62151\\
62000	61255\\
64000	61624\\
66000	62513\\
68000	62159\\
70000	61206\\
72000	60655\\
74000	61430\\
76000	60929\\
78000	61511\\
80000	61233\\
82000	62102\\
86000	62349\\
88000	62859\\
90000	62962\\
92000	62166\\
94000	61756\\
96000	61545\\
98000	63038\\
1e+05	63497\\
};
\addplot [color=mycolor6, forget plot]
  table[row sep=crcr]{%
0	4395.3\\
2000	29633\\
4000	46969\\
6000	50855\\
8000	52445\\
10000	53173\\
12000	54519\\
14000	52988\\
16000	53577\\
18000	54726\\
20000	54630\\
22000	55101\\
26000	54119\\
28000	53506\\
30000	53545\\
32000	54631\\
34000	55396\\
36000	53731\\
40000	55744\\
42000	54434\\
44000	53859\\
46000	54765\\
48000	54823\\
50000	54331\\
52000	55143\\
54000	55243\\
56000	54537\\
58000	55496\\
60000	54404\\
62000	54133\\
64000	54254\\
70000	55650\\
72000	54915\\
74000	53601\\
76000	54531\\
78000	54228\\
80000	54544\\
82000	53971\\
84000	54774\\
86000	54649\\
88000	53263\\
90000	54189\\
92000	53556\\
94000	53492\\
96000	54917\\
98000	54797\\
1e+05	53888\\
};
\addplot [color=mycolor7, forget plot]
  table[row sep=crcr]{%
0	4395.3\\
2000	18357\\
4000	24643\\
6000	26419\\
8000	27279\\
10000	27187\\
12000	26675\\
14000	27008\\
16000	27138\\
18000	27403\\
20000	26849\\
22000	27019\\
24000	27104\\
26000	27716\\
28000	27844\\
30000	27782\\
32000	27233\\
34000	26806\\
36000	27702\\
38000	27965\\
44000	27400\\
46000	27794\\
48000	27354\\
50000	27569\\
52000	27270\\
54000	27182\\
56000	27254\\
58000	27448\\
60000	28451\\
64000	27604\\
66000	27862\\
68000	27242\\
70000	27109\\
72000	27570\\
74000	27361\\
76000	27385\\
78000	27003\\
80000	27657\\
82000	27546\\
84000	27048\\
86000	26951\\
88000	27463\\
90000	27006\\
92000	27452\\
94000	27625\\
96000	26596\\
98000	27061\\
1e+05	27142\\
};
\addplot [color=mycolor8, forget plot]
  table[row sep=crcr]{%
0	4395.3\\
2000	17234\\
4000	20965\\
6000	21081\\
8000	20680\\
10000	20346\\
12000	20738\\
16000	20315\\
18000	20170\\
20000	20225\\
22000	20583\\
24000	20726\\
26000	20155\\
28000	20629\\
30000	20236\\
32000	19924\\
34000	19850\\
36000	19852\\
38000	20301\\
40000	20292\\
42000	20082\\
44000	20207\\
46000	20219\\
50000	20101\\
52000	20481\\
54000	20122\\
56000	19692\\
58000	20498\\
60000	20901\\
62000	20552\\
64000	19852\\
66000	19868\\
68000	19786\\
70000	20376\\
72000	20459\\
76000	20452\\
78000	19917\\
80000	20063\\
82000	20521\\
84000	20597\\
86000	19891\\
88000	19955\\
90000	20481\\
92000	19773\\
94000	20289\\
96000	19650\\
98000	19794\\
1e+05	20112\\
};
\addplot [color=mycolor9, forget plot]
  table[row sep=crcr]{%
0	4395.3\\
2000	11613\\
4000	13366\\
6000	13310\\
8000	13663\\
12000	13561\\
14000	13318\\
16000	13858\\
18000	14202\\
22000	13885\\
24000	14014\\
26000	13973\\
28000	13747\\
30000	13659\\
32000	13808\\
34000	13790\\
36000	13812\\
38000	14239\\
40000	13989\\
42000	14254\\
44000	14322\\
46000	13983\\
48000	14126\\
50000	13953\\
52000	13513\\
56000	13749\\
58000	14033\\
60000	13953\\
62000	14249\\
64000	13858\\
66000	13934\\
70000	13894\\
74000	13505\\
76000	13666\\
78000	13742\\
80000	14002\\
82000	13817\\
84000	13883\\
86000	14145\\
88000	13831\\
90000	13836\\
92000	14074\\
94000	14423\\
96000	13996\\
98000	13745\\
1e+05	13804\\
};
\addplot [color=mycolor10, forget plot]
  table[row sep=crcr]{%
0	4395.3\\
2000	14970\\
4000	20355\\
6000	20756\\
8000	21411\\
10000	21571\\
12000	21105\\
14000	21374\\
16000	21272\\
18000	21874\\
20000	21964\\
24000	21394\\
26000	21388\\
28000	21733\\
30000	21218\\
32000	21687\\
34000	21596\\
36000	21719\\
38000	21918\\
40000	21562\\
42000	21805\\
44000	21396\\
46000	21357\\
50000	21870\\
52000	21395\\
54000	21795\\
56000	21762\\
58000	21533\\
62000	21363\\
64000	21402\\
66000	20876\\
68000	21650\\
70000	21562\\
72000	21653\\
74000	21415\\
76000	21710\\
78000	21524\\
80000	21608\\
82000	21498\\
84000	21854\\
86000	21727\\
88000	21431\\
90000	21028\\
92000	21223\\
94000	21520\\
96000	21389\\
98000	21338\\
1e+05	21862\\
};
\addplot [color=mycolor11, forget plot]
  table[row sep=crcr]{%
0	4395.3\\
2000	9463\\
4000	9563.5\\
6000	9092.2\\
8000	8873.3\\
10000	8971.8\\
12000	8702\\
14000	8891\\
16000	8487.3\\
18000	8461.9\\
20000	8503.3\\
22000	8504.2\\
24000	8668.9\\
26000	8589.9\\
28000	8687.1\\
30000	8737.9\\
32000	8552.8\\
34000	8751.9\\
36000	8419.3\\
38000	8567.6\\
40000	8448.8\\
42000	8506.8\\
44000	8798.8\\
46000	8486.5\\
48000	8627.9\\
50000	8704.5\\
52000	8689.9\\
54000	8512.1\\
56000	8411.7\\
58000	8643.9\\
60000	8330.4\\
62000	8729.5\\
64000	8531.7\\
68000	8622.6\\
70000	8621.1\\
72000	8505.8\\
74000	8725.1\\
76000	8754.3\\
78000	8329.3\\
80000	8313.6\\
82000	8518.7\\
84000	8525.7\\
86000	8355.7\\
88000	8393.4\\
90000	8578.4\\
92000	8623.7\\
94000	8285.1\\
96000	8415.7\\
98000	8575.4\\
1e+05	8528.1\\
};
\addplot [color=mycolor12, forget plot]
  table[row sep=crcr]{%
0	4395.3\\
2000	11498\\
4000	13305\\
6000	14070\\
8000	14135\\
12000	13942\\
14000	14435\\
16000	14349\\
18000	14604\\
20000	14565\\
22000	14420\\
24000	14571\\
26000	14528\\
28000	14611\\
30000	14529\\
32000	14207\\
34000	14700\\
36000	14189\\
38000	14768\\
40000	14366\\
44000	14624\\
46000	14282\\
48000	14112\\
50000	14529\\
52000	14668\\
54000	14144\\
56000	14553\\
58000	14389\\
60000	14273\\
62000	14635\\
66000	14524\\
68000	14392\\
70000	14685\\
72000	14698\\
78000	14427\\
80000	14225\\
82000	15016\\
84000	14588\\
86000	14527\\
88000	14360\\
90000	14269\\
92000	14031\\
94000	14838\\
96000	14385\\
98000	14463\\
1e+05	14311\\
};
\addplot [color=mycolor13, forget plot]
  table[row sep=crcr]{%
0	4395.3\\
2000	9559.8\\
4000	10655\\
8000	10576\\
10000	10780\\
12000	10950\\
14000	11039\\
16000	11084\\
18000	10778\\
20000	10791\\
22000	10700\\
24000	10817\\
26000	11037\\
28000	11093\\
30000	11113\\
32000	11219\\
34000	11223\\
36000	10980\\
38000	11283\\
40000	11259\\
42000	10734\\
46000	10958\\
48000	11125\\
50000	11170\\
52000	11287\\
54000	10948\\
56000	11155\\
58000	11033\\
62000	10842\\
64000	10833\\
66000	10489\\
68000	10966\\
70000	11040\\
72000	11071\\
74000	11005\\
76000	11069\\
78000	11102\\
80000	10993\\
82000	11119\\
84000	11115\\
86000	10903\\
88000	10840\\
90000	11219\\
92000	11308\\
96000	11021\\
98000	11091\\
1e+05	10838\\
};
\addplot [color=mycolor14, forget plot]
  table[row sep=crcr]{%
0	4395.3\\
2000	15315\\
4000	21464\\
6000	23048\\
10000	23645\\
12000	23276\\
14000	23775\\
16000	23107\\
18000	23721\\
20000	24057\\
22000	24006\\
24000	23775\\
26000	23756\\
28000	23415\\
30000	23350\\
32000	23650\\
34000	23861\\
36000	23751\\
38000	23873\\
40000	23268\\
42000	23931\\
44000	23334\\
46000	23629\\
48000	23617\\
50000	24104\\
52000	23277\\
54000	24193\\
56000	24170\\
60000	24390\\
64000	23978\\
66000	23617\\
68000	23179\\
70000	24120\\
72000	23663\\
74000	23148\\
76000	23387\\
78000	23100\\
80000	23998\\
82000	23292\\
84000	23660\\
88000	24004\\
92000	24195\\
94000	23556\\
96000	23746\\
98000	23560\\
1e+05	23787\\
};
\addplot [color=mycolor15, forget plot]
  table[row sep=crcr]{%
0	4395.3\\
2000	11978\\
4000	15738\\
6000	16079\\
8000	15853\\
10000	16466\\
12000	16434\\
14000	16868\\
16000	17052\\
18000	16334\\
20000	17087\\
22000	16928\\
24000	17013\\
26000	16424\\
28000	16918\\
30000	17127\\
32000	16700\\
34000	16728\\
36000	16623\\
38000	16444\\
40000	16599\\
44000	16642\\
46000	16826\\
48000	16938\\
50000	16890\\
52000	16762\\
54000	16214\\
56000	16450\\
58000	16636\\
60000	16738\\
62000	17001\\
64000	16399\\
66000	16613\\
68000	16653\\
70000	15987\\
72000	16175\\
74000	16841\\
76000	16312\\
78000	16961\\
80000	16792\\
82000	16996\\
84000	16781\\
86000	16616\\
88000	16405\\
92000	16733\\
94000	16491\\
96000	16551\\
98000	16392\\
1e+05	17099\\
};
\addplot [color=mycolor16, forget plot]
  table[row sep=crcr]{%
0	4395.3\\
2000	7892.4\\
4000	7385.6\\
6000	6930.9\\
8000	7098\\
10000	7079.4\\
14000	6991.7\\
16000	7045\\
18000	6999.2\\
20000	7082.2\\
22000	7235.4\\
26000	6964.1\\
28000	7136.9\\
30000	6726.5\\
32000	6822.9\\
34000	7012.8\\
36000	6919.2\\
38000	7238.6\\
40000	7133.7\\
42000	6963.7\\
44000	7252.7\\
48000	6887.8\\
50000	7046.6\\
52000	6763.9\\
54000	6913.1\\
56000	7025.5\\
58000	6919.2\\
62000	7050.4\\
64000	6968.6\\
66000	6908.2\\
68000	6988.4\\
70000	6979.4\\
72000	7111.4\\
74000	7206.9\\
76000	6933.6\\
78000	7135.6\\
80000	7166.7\\
82000	6974.7\\
84000	6835.6\\
88000	6986.7\\
90000	6964\\
92000	6905.7\\
94000	6931.3\\
96000	6916.9\\
98000	6847.3\\
1e+05	6952.9\\
};
\addplot [color=mycolor17, forget plot]
  table[row sep=crcr]{%
0	4395.3\\
2000	11510\\
4000	14454\\
6000	15488\\
8000	15423\\
12000	15872\\
14000	15637\\
16000	15810\\
18000	15344\\
20000	15302\\
22000	15777\\
24000	16424\\
26000	15636\\
28000	15867\\
30000	15923\\
32000	15652\\
34000	16240\\
36000	15559\\
38000	15614\\
40000	15910\\
42000	15586\\
44000	15462\\
46000	15963\\
48000	15839\\
52000	15828\\
54000	15913\\
56000	15654\\
58000	16111\\
60000	15744\\
62000	15545\\
64000	15489\\
66000	15981\\
68000	15691\\
70000	15569\\
74000	16296\\
76000	15951\\
78000	15858\\
80000	15430\\
82000	15839\\
84000	15490\\
86000	16156\\
90000	15608\\
92000	15589\\
96000	16436\\
98000	16100\\
1e+05	15667\\
};
\addplot [color=mycolor18, forget plot]
  table[row sep=crcr]{%
0	4395.3\\
2000	14975\\
4000	21530\\
6000	22861\\
8000	22962\\
10000	24071\\
12000	23834\\
14000	23689\\
16000	24007\\
18000	23454\\
20000	24283\\
22000	23593\\
24000	23518\\
26000	24007\\
28000	23552\\
30000	24286\\
32000	23731\\
34000	23724\\
36000	23881\\
38000	23783\\
40000	23864\\
42000	23466\\
44000	23680\\
46000	24054\\
48000	24082\\
50000	23823\\
52000	23904\\
54000	24406\\
56000	24028\\
58000	23888\\
60000	24002\\
62000	24473\\
64000	24047\\
66000	23254\\
68000	23528\\
70000	23605\\
72000	23773\\
74000	24377\\
76000	24327\\
78000	23890\\
82000	23832\\
84000	24253\\
86000	24193\\
88000	23552\\
90000	24218\\
92000	24327\\
94000	23558\\
96000	23533\\
98000	23590\\
1e+05	23743\\
};
\addplot [color=mycolor19, forget plot]
  table[row sep=crcr]{%
0	4395.3\\
2000	11113\\
4000	13847\\
6000	14105\\
8000	14246\\
10000	14449\\
12000	14488\\
14000	14290\\
16000	14663\\
18000	14926\\
20000	14798\\
22000	14610\\
24000	14338\\
26000	14438\\
28000	14444\\
30000	14734\\
32000	14515\\
34000	14426\\
36000	14864\\
38000	14790\\
40000	15112\\
42000	14785\\
44000	14920\\
46000	14726\\
52000	14537\\
54000	14555\\
56000	14886\\
58000	14881\\
60000	14589\\
62000	14614\\
64000	14500\\
66000	15010\\
68000	14836\\
70000	14804\\
72000	14831\\
74000	14635\\
76000	15001\\
78000	14835\\
80000	14774\\
82000	14771\\
84000	14982\\
86000	14561\\
88000	14443\\
90000	14769\\
92000	14827\\
94000	14465\\
96000	14612\\
1e+05	14382\\
};
\addplot [color=mycolor20, forget plot]
  table[row sep=crcr]{%
0	4395.3\\
2000	12593\\
4000	16870\\
6000	18043\\
10000	18197\\
12000	18869\\
14000	19198\\
16000	18516\\
20000	18850\\
22000	18520\\
24000	18403\\
26000	18714\\
28000	18866\\
30000	18905\\
32000	19081\\
34000	19002\\
36000	19065\\
38000	18712\\
40000	19218\\
42000	19216\\
44000	19606\\
46000	19051\\
48000	18892\\
50000	18577\\
52000	19288\\
54000	18498\\
56000	19094\\
58000	18837\\
60000	18818\\
64000	19111\\
66000	18642\\
68000	18865\\
70000	18617\\
72000	19267\\
76000	18599\\
78000	18805\\
80000	19180\\
82000	19157\\
84000	19590\\
86000	18645\\
88000	18573\\
90000	18438\\
92000	18631\\
94000	18662\\
96000	18794\\
98000	18557\\
1e+05	19184\\
};
\addplot [color=mycolor21, forget plot]
  table[row sep=crcr]{%
0	4395.3\\
2000	7630\\
4000	8911\\
6000	9222.6\\
8000	9384\\
10000	9415.1\\
12000	9486.9\\
14000	9635.2\\
18000	9495.5\\
20000	9807.4\\
22000	9636.3\\
24000	10178\\
26000	9816.3\\
28000	9745.9\\
30000	9859.2\\
32000	9636.6\\
34000	9612\\
36000	9634.4\\
38000	9561.4\\
40000	9647.1\\
42000	9703.2\\
44000	9710.7\\
46000	9464.8\\
48000	9297.4\\
50000	9591.6\\
52000	9696.5\\
54000	9674.7\\
56000	9465.4\\
58000	9642.7\\
60000	9626\\
62000	9764.4\\
64000	9833.8\\
66000	9804.6\\
68000	9634\\
70000	9626.5\\
72000	9725.5\\
74000	9503\\
76000	9773.5\\
78000	9611.3\\
82000	9598\\
84000	9775.1\\
88000	9692.1\\
90000	9594.8\\
92000	9708.9\\
94000	9571.6\\
96000	9605.5\\
98000	9514.2\\
1e+05	9642.9\\
};
\addplot [color=mycolor22, forget plot]
  table[row sep=crcr]{%
0	4395.3\\
2000	17882\\
4000	25882\\
6000	26357\\
8000	25755\\
10000	25385\\
12000	25207\\
14000	24605\\
16000	25168\\
18000	25369\\
20000	25144\\
22000	24542\\
24000	24655\\
26000	24585\\
28000	24855\\
30000	25028\\
32000	24916\\
34000	24626\\
36000	24884\\
38000	24387\\
40000	24821\\
42000	24511\\
44000	24912\\
46000	24713\\
48000	24407\\
50000	24444\\
52000	24687\\
54000	24451\\
56000	24401\\
58000	24703\\
60000	24510\\
62000	24520\\
64000	25012\\
66000	24587\\
68000	25078\\
70000	25205\\
72000	24925\\
76000	24829\\
78000	24409\\
80000	25230\\
82000	25145\\
84000	24930\\
86000	24907\\
88000	24230\\
90000	24301\\
92000	24577\\
94000	24459\\
96000	24629\\
98000	24477\\
1e+05	24417\\
};
\addplot [color=mycolor23, forget plot]
  table[row sep=crcr]{%
0	4395.3\\
2000	9679\\
4000	11940\\
6000	12105\\
10000	12006\\
12000	12129\\
14000	12084\\
16000	12110\\
18000	12262\\
20000	12194\\
22000	12035\\
24000	12213\\
26000	12168\\
28000	12059\\
30000	12174\\
32000	12175\\
34000	11897\\
38000	12158\\
40000	12238\\
42000	12200\\
44000	12417\\
46000	12208\\
48000	12275\\
50000	12106\\
52000	12051\\
54000	12181\\
56000	12395\\
58000	12189\\
62000	11931\\
64000	12072\\
66000	12090\\
68000	12349\\
70000	12283\\
72000	12363\\
74000	12184\\
76000	11972\\
78000	12203\\
80000	12245\\
82000	12378\\
84000	12163\\
86000	12407\\
88000	12387\\
90000	12094\\
92000	12197\\
94000	12153\\
96000	12050\\
98000	11899\\
1e+05	11862\\
};
\addplot [color=mycolor24, forget plot]
  table[row sep=crcr]{%
0	4395.3\\
2000	5868.9\\
4000	5315.6\\
6000	4978.9\\
8000	4725.7\\
10000	4771.3\\
12000	4744\\
14000	4905.9\\
16000	4798.2\\
18000	4917.8\\
20000	4644.3\\
22000	4867.7\\
24000	4909\\
26000	4808.4\\
28000	4815.5\\
30000	4855\\
32000	4811.5\\
34000	4891.9\\
38000	4871.7\\
40000	4837.9\\
42000	4768.4\\
44000	4841.8\\
46000	4875.1\\
48000	4889\\
50000	4735.8\\
52000	4743.7\\
54000	4856.1\\
56000	4895.2\\
58000	4837\\
60000	4920\\
62000	4768.7\\
64000	4934\\
66000	4748.9\\
68000	4886.2\\
70000	4734.7\\
72000	4764.5\\
74000	4663.7\\
76000	4889.4\\
78000	4887.9\\
80000	4699.9\\
82000	4869.8\\
84000	4895.8\\
86000	4859.2\\
88000	4806.8\\
90000	4586.6\\
92000	4579.7\\
94000	4711.2\\
96000	4868.3\\
98000	4779.5\\
1e+05	4744.1\\
};
\addplot [color=mycolor25, forget plot]
  table[row sep=crcr]{%
0	4395.3\\
2000	9448.4\\
4000	11565\\
6000	11494\\
8000	12099\\
10000	12222\\
14000	12115\\
16000	12166\\
18000	12154\\
20000	12350\\
22000	12318\\
24000	12368\\
26000	12328\\
30000	11916\\
32000	12481\\
34000	12350\\
36000	12526\\
38000	12194\\
40000	12043\\
42000	12393\\
44000	12380\\
46000	12503\\
48000	12482\\
50000	12026\\
52000	12456\\
54000	12548\\
56000	11868\\
58000	12169\\
60000	12380\\
62000	12005\\
64000	12422\\
66000	12420\\
68000	12123\\
70000	11789\\
72000	12193\\
74000	12476\\
76000	12284\\
78000	12358\\
80000	12233\\
82000	12358\\
84000	11985\\
86000	11972\\
88000	12161\\
90000	12130\\
92000	11996\\
94000	12322\\
96000	12291\\
98000	12226\\
1e+05	12345\\
};
\addplot [color=mycolor26, forget plot]
  table[row sep=crcr]{%
0	4395.3\\
2000	7087.7\\
4000	7839.9\\
8000	7976.1\\
10000	8148\\
12000	8111.8\\
14000	8037.1\\
16000	8408.8\\
18000	8189\\
20000	8228\\
22000	8141.6\\
24000	7944.8\\
26000	8005.5\\
28000	8209.2\\
30000	8100.6\\
32000	8470.1\\
34000	8094.9\\
36000	7990.4\\
38000	8303.5\\
40000	8285.6\\
42000	8382.5\\
44000	8119.2\\
46000	8190.8\\
54000	8321.6\\
58000	8211.9\\
60000	8298.4\\
64000	8403.3\\
66000	8285.1\\
68000	8105.1\\
70000	8248.2\\
72000	8091.1\\
74000	8256.7\\
76000	8112.1\\
78000	8191.1\\
80000	8052.1\\
82000	8464.8\\
84000	8324.4\\
86000	8370.4\\
88000	8248.5\\
90000	8203\\
92000	8221.5\\
94000	8197.5\\
96000	8337.6\\
98000	8301.1\\
1e+05	8367.4\\
};
\addplot [color=mycolor27, forget plot]
  table[row sep=crcr]{%
0	4395.3\\
2000	9113.1\\
4000	10142\\
6000	9459.8\\
8000	9431.5\\
10000	9600.5\\
12000	9955.9\\
14000	9939.1\\
16000	9693.1\\
18000	9681\\
20000	9824.1\\
22000	9708.2\\
24000	9834.4\\
26000	9450.7\\
28000	10029\\
30000	10011\\
32000	9655.6\\
34000	9943.3\\
36000	9705.5\\
38000	9401.2\\
40000	9709.5\\
42000	9707.3\\
44000	9228.6\\
48000	9480.5\\
50000	9438.7\\
52000	9875.7\\
54000	9682.5\\
56000	9773.3\\
58000	9743.8\\
60000	9604.7\\
62000	9766\\
64000	9656.1\\
66000	9666\\
68000	9419\\
70000	9539.3\\
72000	9465.3\\
74000	9597.1\\
76000	9602.5\\
78000	9918.4\\
80000	9545.7\\
84000	9600.4\\
86000	9798.7\\
88000	9895.9\\
90000	9747.8\\
92000	9885.5\\
94000	9678.5\\
98000	9941.7\\
1e+05	9724.2\\
};
\addplot [color=mycolor28, forget plot]
  table[row sep=crcr]{%
0	4395.3\\
2000	6394.4\\
4000	5905.5\\
6000	5652.7\\
8000	5562.1\\
10000	5498.4\\
12000	5690.1\\
14000	5579.4\\
16000	5521.2\\
18000	5514.9\\
20000	5444.6\\
22000	5321.5\\
24000	5405.7\\
26000	5396.1\\
28000	5504.8\\
30000	5435.5\\
32000	5453.2\\
34000	5658.9\\
36000	5750.2\\
38000	5505.8\\
40000	5513.4\\
42000	5467.2\\
44000	5549.2\\
50000	5437.4\\
54000	5347.2\\
58000	5511\\
60000	5683.2\\
62000	5451.5\\
64000	5513.2\\
66000	5497.6\\
68000	5651.2\\
70000	5528.6\\
72000	5493.3\\
74000	5489.4\\
76000	5426\\
78000	5671.9\\
80000	5463\\
82000	5433.3\\
84000	5574.2\\
86000	5539.6\\
88000	5464.3\\
90000	5409.5\\
92000	5563.1\\
94000	5404.1\\
96000	5404\\
98000	5420.5\\
1e+05	5484.1\\
};
\addplot [color=mycolor29, forget plot]
  table[row sep=crcr]{%
0	4395.3\\
2000	10175\\
4000	11934\\
6000	11681\\
8000	11247\\
10000	11327\\
12000	11667\\
14000	11527\\
16000	11547\\
18000	11361\\
20000	11426\\
22000	11262\\
24000	11614\\
26000	11642\\
28000	11297\\
30000	11251\\
32000	11389\\
34000	11450\\
36000	11411\\
40000	11581\\
42000	11407\\
44000	11424\\
48000	11399\\
50000	11631\\
52000	11915\\
54000	11629\\
56000	11703\\
58000	11418\\
62000	11573\\
64000	11448\\
66000	11703\\
68000	11386\\
70000	11396\\
72000	11581\\
74000	11103\\
76000	11826\\
78000	12012\\
80000	11290\\
82000	11411\\
84000	11655\\
86000	11745\\
88000	11733\\
90000	11377\\
92000	11675\\
94000	11428\\
96000	11616\\
98000	11472\\
1e+05	11261\\
};
\addplot [color=mycolor30, forget plot]
  table[row sep=crcr]{%
0	4395.3\\
2000	6520.8\\
4000	7463.4\\
6000	7663.5\\
8000	7756.8\\
10000	7886\\
12000	7649.7\\
14000	7830.4\\
16000	7753.6\\
18000	7985.2\\
20000	7870.4\\
22000	7793.5\\
26000	7836.9\\
28000	7762.4\\
30000	7738.8\\
32000	7919.4\\
34000	8064.5\\
36000	7874.9\\
38000	7809.3\\
40000	7882.1\\
42000	7795.7\\
46000	7756.9\\
48000	7663.4\\
52000	8000.1\\
54000	7879.1\\
56000	7930\\
58000	7774.3\\
60000	7979.8\\
62000	7996\\
64000	7921.3\\
66000	7967.8\\
68000	7976.6\\
70000	7858.1\\
72000	7888.8\\
74000	7986.1\\
76000	8009\\
78000	7821.8\\
80000	7778.2\\
86000	7974.1\\
88000	7757.9\\
90000	7704.1\\
92000	7842.9\\
94000	8096.9\\
96000	7836.4\\
98000	7970.3\\
1e+05	7933.2\\
};
\addplot [color=mycolor31, forget plot]
  table[row sep=crcr]{%
0	4395.3\\
2000	6824.3\\
4000	6969.4\\
6000	6594.4\\
8000	6589.6\\
10000	6717\\
12000	6559.5\\
14000	6540.7\\
16000	6679.8\\
18000	6457.4\\
20000	6449.7\\
22000	6334.3\\
24000	6335.9\\
26000	6546\\
28000	6474.3\\
30000	6529.4\\
32000	6471.6\\
34000	6521.1\\
36000	6530.3\\
38000	6648.9\\
40000	6503.3\\
42000	6481.9\\
44000	6518.4\\
46000	6494.9\\
48000	6327.1\\
50000	6576.2\\
52000	6645.7\\
54000	6566.8\\
56000	6593.6\\
60000	6289.9\\
62000	6298.3\\
64000	6419.5\\
66000	6419.3\\
68000	6605.3\\
70000	6558.8\\
72000	6604.7\\
74000	6694.9\\
76000	6574.6\\
78000	6422.6\\
80000	6409.1\\
82000	6454.1\\
84000	6434.1\\
86000	6449.5\\
88000	6589.8\\
90000	6268.5\\
92000	6296.9\\
94000	6439\\
96000	6443.8\\
98000	6591.3\\
1e+05	6579.4\\
};
\addplot [color=mycolor32, forget plot]
  table[row sep=crcr]{%
0	4395.3\\
2000	9687.1\\
4000	11192\\
6000	11208\\
8000	10820\\
10000	10547\\
12000	10819\\
16000	10609\\
18000	10207\\
20000	10153\\
22000	10291\\
24000	10299\\
26000	10443\\
28000	10724\\
32000	10500\\
34000	10622\\
36000	10198\\
38000	10164\\
40000	10236\\
42000	10500\\
44000	10170\\
46000	10197\\
48000	10517\\
50000	10288\\
52000	10646\\
56000	10543\\
58000	10441\\
60000	10411\\
62000	10345\\
64000	10568\\
66000	10482\\
68000	10117\\
70000	10075\\
72000	10291\\
74000	10232\\
76000	10456\\
80000	10551\\
82000	10440\\
84000	10595\\
86000	10560\\
88000	10447\\
90000	10297\\
92000	10216\\
94000	10442\\
96000	10523\\
98000	10532\\
1e+05	10376\\
};
\addplot [color=mycolor33, forget plot]
  table[row sep=crcr]{%
0	4395.3\\
2000	6420.7\\
4000	7581\\
6000	7845.2\\
8000	8145.2\\
10000	8151.5\\
12000	8205.1\\
14000	8008.5\\
16000	8069.3\\
20000	8231.7\\
22000	8111.3\\
24000	8031.6\\
26000	8042.5\\
28000	8129.4\\
30000	8080.8\\
32000	8274.4\\
34000	8239.8\\
36000	8363.8\\
38000	8353.2\\
40000	8187.6\\
42000	8326.3\\
44000	8182.3\\
46000	8240.2\\
48000	8055.3\\
50000	8224.1\\
52000	7885.4\\
54000	8349.4\\
56000	8025.7\\
58000	8014.4\\
60000	8358\\
62000	8517.4\\
64000	8391.7\\
66000	8118.7\\
68000	8071.2\\
70000	8174.3\\
72000	8205.2\\
74000	7958.7\\
76000	8245.1\\
78000	8013.8\\
80000	7963.7\\
82000	8308.4\\
84000	8101.7\\
86000	8175.4\\
88000	8448.3\\
90000	8167.8\\
92000	8192\\
94000	8263.5\\
96000	8310.3\\
98000	8122.3\\
1e+05	8259.3\\
};
\addplot [color=mycolor34, forget plot]
  table[row sep=crcr]{%
0	4395.3\\
2000	5141.3\\
4000	4456.7\\
6000	4091.6\\
8000	4050.8\\
10000	4026.3\\
12000	3959.5\\
14000	4022.2\\
16000	3933.3\\
18000	4022.7\\
20000	4062.6\\
22000	4142.6\\
24000	4014\\
26000	3926.3\\
28000	4146.9\\
30000	4051.6\\
32000	3979.9\\
34000	4110\\
36000	4094.3\\
38000	4022.8\\
40000	4244.4\\
42000	4103.2\\
44000	4210.4\\
46000	4163.4\\
48000	3996.1\\
50000	4173.2\\
54000	4088.8\\
56000	4025.7\\
58000	4128.6\\
60000	4058.6\\
62000	4033.9\\
64000	4189.4\\
66000	3988.5\\
68000	4080.6\\
70000	4092.8\\
72000	4039.7\\
74000	4064.1\\
76000	4133.9\\
80000	4186.9\\
82000	4019.1\\
84000	4024.9\\
86000	4150.8\\
88000	4061.4\\
90000	4055.6\\
92000	4201.4\\
94000	4233.8\\
96000	4141.9\\
98000	4106.7\\
1e+05	4167.5\\
};
\addplot [color=mycolor35]
  table[row sep=crcr]{%
0	4395.3\\
2000	8455.7\\
4000	9866.1\\
6000	9874.2\\
8000	9987.1\\
10000	9843.4\\
12000	9889\\
14000	10117\\
16000	10041\\
18000	9910.8\\
20000	10114\\
22000	9883.9\\
24000	10038\\
26000	10237\\
28000	10254\\
30000	10000\\
32000	10181\\
34000	9947.1\\
36000	10266\\
38000	10043\\
40000	10240\\
42000	10216\\
46000	10093\\
48000	10122\\
50000	10029\\
52000	10130\\
54000	10186\\
56000	10075\\
58000	10265\\
60000	10069\\
62000	10047\\
64000	10117\\
66000	9957.9\\
70000	10079\\
72000	10177\\
74000	10145\\
76000	10057\\
80000	10097\\
82000	10152\\
84000	10112\\
86000	10163\\
90000	10016\\
92000	10019\\
94000	10341\\
96000	10228\\
98000	10182\\
1e+05	10202\\
};
\addlegendentry{Classe k=83}

\addplot [color=mycolor36, forget plot]
  table[row sep=crcr]{%
0	4395.3\\
2000	5427.8\\
4000	4749.8\\
6000	4334.5\\
8000	4468.8\\
10000	4393.7\\
12000	4383\\
14000	4229.4\\
16000	4339.9\\
18000	4228.9\\
20000	4267.2\\
22000	4405.4\\
24000	4307.6\\
26000	4304.1\\
28000	4259.8\\
30000	4290.4\\
32000	4287.6\\
34000	4345.5\\
38000	4337.5\\
40000	4268.4\\
42000	4346.4\\
44000	4356.6\\
46000	4300.5\\
48000	4265.7\\
50000	4305.5\\
54000	4315.7\\
56000	4423\\
58000	4308.6\\
60000	4411.3\\
62000	4320.5\\
64000	4357.6\\
66000	4304.7\\
68000	4444.4\\
70000	4298.5\\
72000	4286.9\\
74000	4303\\
76000	4401.8\\
78000	4340.2\\
80000	4440.5\\
82000	4140.1\\
86000	4274\\
88000	4334.1\\
90000	4415.1\\
92000	4442.8\\
94000	4258.6\\
96000	4301\\
98000	4270.8\\
1e+05	4325.9\\
};
\addplot [color=mycolor37, forget plot]
  table[row sep=crcr]{%
0	4395.3\\
2000	4293.8\\
4000	3886.3\\
6000	3761.9\\
8000	3717.3\\
10000	3708.8\\
12000	3623.7\\
14000	3721.3\\
16000	3689.3\\
18000	3642\\
20000	3721.9\\
24000	3612.7\\
26000	3673.1\\
28000	3716.7\\
30000	3693.2\\
32000	3749.1\\
34000	3675\\
36000	3646.3\\
38000	3716.3\\
40000	3683\\
42000	3639.4\\
44000	3684\\
46000	3629.9\\
48000	3645.5\\
50000	3641.3\\
52000	3673.8\\
54000	3723\\
56000	3662.9\\
58000	3719\\
60000	3663.2\\
62000	3579.4\\
64000	3634.9\\
66000	3662.6\\
68000	3671.3\\
70000	3615\\
72000	3665.7\\
74000	3638.5\\
76000	3700.6\\
78000	3630.7\\
80000	3730.7\\
82000	3753.4\\
86000	3785.8\\
88000	3701.5\\
90000	3713.9\\
92000	3750\\
94000	3670.5\\
96000	3626\\
98000	3707.1\\
1e+05	3662.3\\
};
\addplot [color=mycolor38, forget plot]
  table[row sep=crcr]{%
0	4395.3\\
2000	3608.5\\
4000	3208\\
6000	3129.2\\
8000	3171.2\\
10000	3116.4\\
12000	2996\\
14000	3065.3\\
16000	3213.7\\
18000	3129.4\\
20000	3101\\
22000	3050.9\\
24000	3096.7\\
26000	3059\\
28000	3031.6\\
32000	3108\\
34000	2931.8\\
36000	3120.2\\
38000	3076.1\\
40000	2950.8\\
42000	3046\\
44000	3070.4\\
46000	3117.6\\
48000	3100.4\\
50000	3113.8\\
52000	2972.3\\
56000	2999\\
58000	3019.4\\
60000	2975.5\\
62000	3061.8\\
64000	3100.8\\
66000	3103.1\\
68000	3176\\
70000	3064.5\\
72000	3058.7\\
74000	2953.8\\
76000	3116.5\\
78000	3016.2\\
80000	2994.9\\
82000	3043.3\\
84000	3057.1\\
86000	3173.8\\
88000	2995\\
90000	2971.6\\
92000	3034.4\\
94000	3047.6\\
96000	3028.7\\
98000	3136.9\\
1e+05	3104.1\\
};
\addplot [color=mycolor39, forget plot]
  table[row sep=crcr]{%
0	4395.3\\
2000	10218\\
4000	12553\\
6000	12620\\
8000	12405\\
10000	12380\\
12000	11911\\
14000	12285\\
16000	11688\\
18000	11824\\
20000	11749\\
22000	11737\\
24000	11598\\
26000	11774\\
28000	11831\\
30000	11801\\
32000	11459\\
34000	11705\\
36000	11919\\
38000	11818\\
42000	11846\\
44000	11782\\
46000	11526\\
48000	11714\\
50000	11701\\
52000	11396\\
54000	11780\\
56000	11573\\
58000	11665\\
60000	11366\\
62000	11775\\
64000	11975\\
66000	11565\\
68000	11944\\
70000	12033\\
72000	11763\\
74000	11979\\
76000	11693\\
78000	11530\\
80000	11759\\
82000	11641\\
86000	11909\\
88000	11669\\
90000	11499\\
92000	11664\\
98000	11579\\
1e+05	11383\\
};
\addplot [color=mycolor40, forget plot]
  table[row sep=crcr]{%
0	4395.3\\
2000	4930.3\\
4000	5317.7\\
6000	5491.5\\
8000	5694.1\\
10000	5653.2\\
12000	5754.2\\
14000	5718.8\\
16000	5880.8\\
18000	5797.2\\
20000	5750.2\\
22000	5749.7\\
24000	5849.3\\
26000	5811.2\\
28000	5716.7\\
30000	5836.8\\
32000	5844.2\\
34000	5756.8\\
36000	5789.2\\
38000	5729.2\\
40000	5868.1\\
42000	5873.9\\
44000	5767.1\\
46000	5813.5\\
48000	5669.5\\
50000	5814.6\\
52000	5777\\
54000	5820.6\\
58000	5777.2\\
60000	5854.5\\
62000	5824.5\\
64000	5709.4\\
66000	5791\\
68000	5797.3\\
70000	5783.6\\
72000	5659.4\\
74000	5674.4\\
76000	5742.7\\
78000	5786.1\\
80000	5780.5\\
82000	5668.4\\
84000	5764\\
86000	5724.9\\
88000	5731.8\\
90000	5782.2\\
92000	5716.8\\
94000	5779.1\\
96000	5713.1\\
98000	5773.5\\
1e+05	5800.3\\
};
\addplot [color=mycolor41, forget plot]
  table[row sep=crcr]{%
0	4395.3\\
2000	4567.4\\
4000	5122.1\\
6000	5238.8\\
8000	5281.7\\
10000	5383.8\\
14000	5557.3\\
16000	5317.7\\
18000	5393.4\\
20000	5394.2\\
22000	5330.7\\
24000	5229.9\\
26000	5422\\
28000	5415.1\\
30000	5514.7\\
32000	5381\\
34000	5463.4\\
36000	5425\\
38000	5290.7\\
40000	5336.4\\
42000	5534.3\\
44000	5545.4\\
46000	5502.4\\
48000	5410.4\\
50000	5359.9\\
52000	5413.8\\
54000	5336.4\\
56000	5333.7\\
58000	5314.3\\
60000	5179.7\\
62000	5367.7\\
64000	5412.2\\
66000	5556.4\\
68000	5537.4\\
70000	5493.6\\
72000	5368.2\\
74000	5323.9\\
76000	5412\\
78000	5518.6\\
80000	5480.3\\
82000	5395.4\\
84000	5460.4\\
86000	5441.8\\
88000	5558.4\\
90000	5367.2\\
92000	5329.8\\
94000	5490.6\\
96000	5354.2\\
98000	5417.4\\
1e+05	5322.7\\
};
\addplot [color=mycolor42, forget plot]
  table[row sep=crcr]{%
0	4395.3\\
2000	11790\\
4000	17771\\
6000	19057\\
8000	19372\\
12000	20418\\
14000	20086\\
16000	19554\\
18000	20171\\
20000	20706\\
22000	21005\\
24000	20899\\
26000	20384\\
28000	20698\\
30000	20663\\
34000	19707\\
36000	20002\\
38000	19837\\
40000	20257\\
42000	19982\\
44000	20147\\
46000	20707\\
48000	20421\\
50000	19879\\
52000	20277\\
56000	20132\\
58000	19989\\
60000	20563\\
62000	20348\\
64000	19940\\
66000	19833\\
68000	19819\\
70000	20140\\
72000	19931\\
74000	19392\\
76000	19734\\
78000	20325\\
80000	20310\\
82000	20065\\
84000	19728\\
86000	20358\\
88000	19956\\
90000	20459\\
92000	20905\\
94000	20511\\
96000	19972\\
98000	19977\\
1e+05	19901\\
};
\addplot [color=mycolor43, forget plot]
  table[row sep=crcr]{%
0	4395.3\\
2000	5353.2\\
4000	5384.6\\
6000	5309.6\\
10000	5195.8\\
14000	4956.6\\
16000	5054.7\\
18000	5003.6\\
20000	5016.7\\
22000	5070.1\\
24000	4991.2\\
26000	4935.9\\
28000	4895.9\\
30000	5013.3\\
32000	5049\\
34000	5001\\
36000	4974.2\\
38000	5037.3\\
42000	5001.1\\
44000	5017.3\\
46000	5005.8\\
48000	4908.6\\
50000	4939.4\\
52000	4923.2\\
54000	5014.9\\
56000	4985.1\\
58000	5007.4\\
60000	4914.2\\
62000	5028.2\\
64000	5001.1\\
66000	5032.2\\
68000	4952.9\\
70000	4937.2\\
72000	4959.5\\
74000	4931.8\\
76000	4963.2\\
78000	4963.9\\
80000	4936.9\\
82000	4962.2\\
84000	5005.8\\
86000	4921.2\\
92000	4984.2\\
96000	4918.2\\
98000	4983.8\\
1e+05	5023.5\\
};
\addplot [color=mycolor44, forget plot]
  table[row sep=crcr]{%
0	4395.3\\
2000	4417.5\\
4000	3965.6\\
6000	3661.9\\
8000	3685\\
10000	3503.1\\
12000	3600.4\\
14000	3493\\
16000	3565.7\\
18000	3524.4\\
20000	3432.9\\
22000	3375.2\\
24000	3297.9\\
26000	3363\\
28000	3476.2\\
30000	3449\\
34000	3356.8\\
36000	3431.3\\
38000	3549.2\\
40000	3395.4\\
42000	3403.2\\
44000	3370.4\\
46000	3431\\
48000	3549.5\\
50000	3346.1\\
52000	3512.6\\
54000	3495.4\\
56000	3374.9\\
58000	3368.3\\
62000	3424.2\\
64000	3412.8\\
66000	3448.8\\
68000	3392.1\\
70000	3395.1\\
72000	3353\\
74000	3437\\
76000	3383.4\\
78000	3414.3\\
80000	3332.5\\
82000	3462.2\\
84000	3312.3\\
86000	3326.2\\
88000	3294.4\\
90000	3483.7\\
92000	3504.5\\
94000	3488.5\\
96000	3447.8\\
98000	3483.9\\
1e+05	3333.5\\
};
\addplot [color=mycolor45, forget plot]
  table[row sep=crcr]{%
0	4395.3\\
2000	5462.8\\
4000	6393.1\\
6000	6675.8\\
8000	6667.2\\
10000	6762.6\\
12000	6801\\
14000	6810\\
16000	6925.7\\
18000	6965.7\\
20000	6918.7\\
22000	6891.4\\
24000	6918.4\\
26000	6819.2\\
28000	6940.2\\
30000	6991.9\\
32000	6860.5\\
36000	6888.1\\
38000	6944.7\\
42000	7102.5\\
44000	7086.1\\
46000	6865.8\\
52000	6851.7\\
54000	6783.3\\
56000	7011.7\\
58000	6848.1\\
60000	6894.6\\
62000	6853.5\\
66000	6924\\
68000	7005.6\\
70000	6953.7\\
72000	6924.8\\
74000	6764.4\\
76000	7032.3\\
78000	6865.6\\
82000	7065.5\\
84000	6879.7\\
86000	6912.8\\
88000	6850.1\\
90000	6886.6\\
92000	6773.3\\
94000	7010.3\\
96000	6894.3\\
98000	7018.6\\
1e+05	6994.1\\
};
\addplot [color=mycolor46, forget plot]
  table[row sep=crcr]{%
0	4395.3\\
2000	4019.1\\
4000	4122.4\\
6000	4083.7\\
8000	3953.1\\
10000	4113.7\\
14000	4131.9\\
18000	4013\\
20000	4267.4\\
22000	4181.4\\
24000	4077.6\\
26000	4135.2\\
28000	4076.4\\
30000	4158.8\\
32000	4085.4\\
34000	4112.7\\
36000	4186.7\\
38000	4120.9\\
42000	4098.1\\
44000	4061.1\\
46000	4053.4\\
48000	4087.4\\
50000	4044\\
52000	4122.3\\
54000	4154\\
56000	4153.7\\
58000	4233.7\\
60000	4172.3\\
62000	4239.1\\
64000	4141.9\\
66000	4063.4\\
68000	4036.5\\
70000	4105.6\\
74000	4029.5\\
76000	4199.6\\
78000	4057.9\\
80000	4175\\
86000	4099\\
88000	4147\\
92000	4174.1\\
94000	4094.1\\
96000	4112.2\\
98000	4118.7\\
1e+05	4088.6\\
};
\addplot [color=mycolor47, forget plot]
  table[row sep=crcr]{%
0	4395.3\\
2000	5214\\
4000	5478.9\\
6000	5311.8\\
12000	5296.5\\
14000	5238.9\\
16000	5319.6\\
18000	5287.8\\
20000	5292.7\\
24000	5358.5\\
26000	5205.8\\
28000	5243.1\\
30000	5178.8\\
32000	5019.8\\
34000	5291\\
36000	5446.2\\
38000	5294.8\\
40000	5303.9\\
42000	5377.2\\
44000	5269.8\\
46000	5392.9\\
48000	5304.2\\
52000	5422.8\\
54000	5426.8\\
56000	5477\\
58000	5378.2\\
60000	5237.6\\
62000	5192.7\\
64000	5302.4\\
66000	5203.6\\
68000	5220\\
70000	5267.1\\
72000	5379.2\\
74000	5414.9\\
76000	5307\\
78000	5288.7\\
80000	5285.6\\
82000	5426.2\\
84000	5241.2\\
86000	5124.1\\
88000	5432.6\\
90000	5275.9\\
92000	5321.7\\
94000	5330\\
96000	5377.8\\
98000	5246.6\\
1e+05	5280.7\\
};
\addplot [color=mycolor48, forget plot]
  table[row sep=crcr]{%
0	4395.3\\
2000	3655.6\\
4000	3205.1\\
6000	3048.9\\
8000	2987.6\\
10000	2957.2\\
12000	2938.3\\
14000	3026.3\\
16000	3027.3\\
18000	2961.3\\
22000	2968.4\\
24000	2974.9\\
28000	3022.4\\
30000	2998.6\\
32000	2941.5\\
34000	2997.5\\
36000	3022.6\\
38000	2971.9\\
40000	2998\\
42000	2906.7\\
44000	2905.1\\
46000	2997.6\\
48000	3004.6\\
50000	2962.6\\
52000	3014.3\\
54000	2971\\
56000	2972.6\\
58000	2926.9\\
62000	2929.2\\
64000	2955.1\\
66000	2920.4\\
68000	2932.5\\
70000	2964.1\\
74000	2932.1\\
76000	2999.5\\
78000	3008.4\\
80000	2997.8\\
84000	3003.5\\
86000	2946.3\\
88000	2931\\
90000	2928.1\\
94000	2994.7\\
96000	2926.5\\
98000	2982.3\\
1e+05	2949.2\\
};
\addplot [color=mycolor49, forget plot]
  table[row sep=crcr]{%
0	4395.3\\
2000	4941.3\\
4000	4736\\
6000	4592\\
8000	4412.6\\
10000	4473.2\\
12000	4291.6\\
16000	4474.3\\
18000	4477.8\\
20000	4521.2\\
22000	4333.4\\
24000	4364.4\\
26000	4473.1\\
28000	4218.4\\
30000	4594.2\\
32000	4541.2\\
34000	4300.9\\
36000	4459.5\\
38000	4476.1\\
40000	4383.1\\
42000	4565\\
44000	4484.2\\
48000	4522.1\\
50000	4417.5\\
52000	4403.9\\
54000	4369.5\\
56000	4378.9\\
58000	4501.7\\
60000	4401.4\\
62000	4340.7\\
64000	4386.9\\
66000	4414.1\\
68000	4415.6\\
70000	4343.8\\
72000	4321.6\\
74000	4570.6\\
78000	4285.9\\
80000	4441.9\\
82000	4447.9\\
84000	4384.3\\
86000	4337\\
88000	4249.9\\
90000	4290\\
92000	4398.1\\
94000	4488.6\\
96000	4442.9\\
98000	4432\\
1e+05	4329\\
};
\addplot [color=mycolor50, forget plot]
  table[row sep=crcr]{%
0	4395.3\\
2000	2820.2\\
4000	2424.4\\
6000	2374.3\\
8000	2459.4\\
10000	2495.2\\
12000	2522.6\\
14000	2558.9\\
16000	2500.1\\
18000	2527.2\\
20000	2504.1\\
22000	2468.9\\
24000	2517.8\\
26000	2520.5\\
28000	2531.9\\
30000	2590.4\\
32000	2466.9\\
34000	2507.9\\
36000	2531.3\\
38000	2495\\
40000	2527.4\\
42000	2491.8\\
44000	2542.6\\
46000	2568.3\\
48000	2513\\
52000	2520.2\\
54000	2526.1\\
56000	2543.7\\
58000	2499.3\\
62000	2525.4\\
64000	2524.9\\
66000	2589\\
68000	2528.3\\
70000	2574\\
72000	2486.4\\
74000	2530.2\\
76000	2514.2\\
78000	2481.7\\
80000	2525.9\\
82000	2492.1\\
84000	2520.8\\
86000	2517.6\\
88000	2547.8\\
96000	2497.6\\
98000	2540.4\\
1e+05	2556.6\\
};
\addplot [color=mycolor51, forget plot]
  table[row sep=crcr]{%
0	4395.3\\
2000	5563.7\\
4000	5862.2\\
6000	5954.8\\
8000	6199.5\\
10000	5668.2\\
12000	5866.8\\
14000	5912.1\\
16000	5880.4\\
18000	6209.1\\
20000	5939.5\\
22000	5893.5\\
24000	5930\\
26000	6068.3\\
28000	5955\\
30000	5990.3\\
32000	6171.5\\
34000	6006.8\\
36000	5879\\
38000	5887.4\\
42000	6003.2\\
44000	5896.2\\
46000	6042.9\\
48000	6116.6\\
52000	6067.8\\
54000	5982.8\\
56000	6082.3\\
58000	5819.3\\
60000	6021\\
62000	5953.3\\
64000	5773.9\\
66000	5921.5\\
68000	5962.1\\
70000	5946.8\\
72000	6077.3\\
74000	6012.2\\
76000	5978.6\\
78000	6309.9\\
80000	5826.3\\
82000	5907\\
84000	5842\\
86000	5876\\
88000	5993.8\\
90000	6085.5\\
92000	6016.5\\
94000	5843.6\\
98000	6055.4\\
1e+05	5960.2\\
};
\addplot [color=mycolor52, forget plot]
  table[row sep=crcr]{%
0	4395.3\\
2000	4390.2\\
4000	3920\\
6000	3837.2\\
8000	3607.6\\
10000	3511.5\\
12000	3658.4\\
14000	3687.1\\
16000	3623.4\\
18000	3536.4\\
20000	3530.8\\
22000	3623.9\\
24000	3519.7\\
26000	3593.7\\
28000	3583.8\\
30000	3678.8\\
32000	3592.1\\
34000	3662.9\\
36000	3469.7\\
38000	3512\\
40000	3655\\
44000	3569.4\\
46000	3553.2\\
48000	3634.2\\
50000	3517.5\\
52000	3493.2\\
56000	3505.9\\
58000	3487\\
60000	3585.5\\
62000	3502.6\\
64000	3480.9\\
66000	3512.2\\
68000	3640.5\\
70000	3590.3\\
72000	3468.9\\
74000	3531.8\\
76000	3572.9\\
78000	3471.5\\
80000	3633.2\\
82000	3470.2\\
84000	3568.7\\
90000	3641.9\\
92000	3534.5\\
94000	3653.5\\
96000	3565.5\\
98000	3575.8\\
1e+05	3555.8\\
};
\addplot [color=mycolor53, forget plot]
  table[row sep=crcr]{%
0	4395.3\\
2000	6142.3\\
4000	6688.2\\
6000	6254.7\\
8000	5781.4\\
10000	5891.7\\
12000	6140.4\\
14000	5918.9\\
16000	5563.1\\
18000	5542.3\\
20000	5778.4\\
22000	5806.9\\
24000	5678.4\\
26000	5670.6\\
28000	5864.7\\
30000	5812.6\\
32000	5778.2\\
36000	6067.9\\
38000	5774.5\\
40000	5682.4\\
42000	5778.7\\
44000	5757.3\\
46000	5655.3\\
50000	5758.8\\
52000	5719.6\\
54000	5613.7\\
56000	5686.2\\
58000	5686.9\\
60000	5660.8\\
62000	5615.6\\
64000	5695.3\\
66000	5725\\
68000	5531.4\\
70000	5744.9\\
72000	5575.7\\
74000	5640.9\\
76000	5768.7\\
78000	5760.7\\
80000	5771.7\\
82000	5700.3\\
84000	5577.6\\
86000	5708.1\\
88000	5784.7\\
90000	5885\\
92000	5689.7\\
94000	5690.2\\
96000	5801.8\\
98000	5789.4\\
1e+05	5532.4\\
};
\addplot [color=mycolor54, forget plot]
  table[row sep=crcr]{%
0	4395.3\\
2000	4816.8\\
4000	4838.9\\
6000	4738.6\\
8000	4730.5\\
10000	4609.6\\
12000	4606.6\\
14000	4562.1\\
16000	4653.7\\
18000	4609.8\\
20000	4548\\
22000	4621.1\\
24000	4627.8\\
26000	4724.9\\
28000	4628.7\\
30000	4709.7\\
32000	4646.9\\
34000	4718\\
36000	4640.9\\
38000	4464.4\\
40000	4558\\
42000	4682\\
44000	4696.1\\
46000	4747.1\\
48000	4598\\
50000	4608.5\\
52000	4667.5\\
54000	4622.2\\
58000	4633.3\\
60000	4560.4\\
62000	4580.9\\
64000	4525.6\\
66000	4634.3\\
68000	4646.1\\
70000	4523.4\\
72000	4577.5\\
74000	4462.7\\
76000	4516.6\\
78000	4473.7\\
80000	4525.1\\
82000	4682.6\\
84000	4560.5\\
86000	4680.2\\
88000	4651.5\\
90000	4564.9\\
94000	4575.1\\
96000	4711.8\\
98000	4695.7\\
1e+05	4818.6\\
};
\addplot [color=mycolor55, forget plot]
  table[row sep=crcr]{%
0	4395.3\\
2000	4921.3\\
4000	5490.4\\
6000	5558.2\\
8000	5516.6\\
10000	5651.8\\
12000	5757.8\\
14000	5685.3\\
16000	5796.6\\
20000	5587.2\\
22000	5782.5\\
24000	5878.6\\
26000	5896.3\\
28000	5712.7\\
30000	5791.1\\
32000	5741.5\\
34000	5884.3\\
36000	5832.7\\
38000	5808.9\\
40000	5764.1\\
42000	5880.7\\
44000	5854.6\\
46000	5724.1\\
48000	5782\\
50000	5898.6\\
52000	5824.8\\
54000	5823.8\\
56000	5793.3\\
58000	5786.9\\
60000	5835.7\\
62000	5969.3\\
64000	5728\\
66000	5774.2\\
68000	5842.8\\
70000	5697.7\\
72000	5742.6\\
74000	5829.7\\
76000	5832.8\\
78000	5698.8\\
80000	5644.1\\
82000	5741.8\\
84000	5759.2\\
86000	5833.4\\
88000	5759.9\\
90000	5706.5\\
92000	5823.5\\
94000	5769.8\\
96000	5815.7\\
98000	5903.4\\
1e+05	5852.9\\
};
\addplot [color=mycolor56, forget plot]
  table[row sep=crcr]{%
0	4395.3\\
2000	4395.4\\
4000	4408.1\\
6000	4217.4\\
8000	4074.2\\
10000	3991.4\\
12000	3984.9\\
14000	4079.8\\
16000	3996.7\\
18000	3952\\
20000	4009.2\\
22000	4031.6\\
26000	3988.3\\
28000	4082.7\\
32000	4016.5\\
34000	4096.1\\
36000	4140.2\\
38000	4100\\
40000	3989.7\\
44000	4054.4\\
46000	4110.9\\
48000	4063.1\\
50000	4050.7\\
52000	4089.7\\
54000	4150.3\\
56000	4008.8\\
58000	4011.9\\
60000	3937\\
62000	4046.3\\
64000	4034.1\\
66000	3996.7\\
68000	4007.9\\
70000	4060.4\\
72000	4086\\
74000	4028.3\\
76000	4059.4\\
78000	4106.8\\
80000	4060.4\\
82000	3981.4\\
84000	3988.1\\
86000	4058.7\\
88000	4000.7\\
90000	4047\\
92000	3972.6\\
94000	3967.3\\
96000	4013.6\\
98000	3981.8\\
1e+05	4005.1\\
};
\addplot [color=mycolor57, forget plot]
  table[row sep=crcr]{%
0	4395.3\\
2000	5714.9\\
4000	6757.9\\
6000	7025.1\\
8000	6915.5\\
10000	6896.1\\
12000	6839.2\\
16000	6667.6\\
18000	6784.1\\
20000	6753.5\\
22000	6616.2\\
24000	6651.9\\
26000	6788.6\\
28000	6859.5\\
30000	6757.8\\
32000	6690.4\\
34000	6576.1\\
36000	6538.3\\
38000	6587.7\\
40000	6616.9\\
42000	6589\\
44000	6698.4\\
46000	6581.6\\
48000	6617.5\\
50000	6695.7\\
52000	6755.7\\
54000	6724.5\\
56000	6781.2\\
58000	6879.1\\
60000	6748.4\\
62000	6674.9\\
64000	6827.4\\
66000	6770.2\\
68000	6669.3\\
70000	6798.1\\
72000	6799.7\\
76000	6720\\
78000	6768.2\\
80000	6690\\
82000	6667.6\\
84000	6529.4\\
86000	6600.2\\
88000	6744.7\\
90000	6731.6\\
92000	6619.7\\
96000	6686.4\\
98000	6764.1\\
1e+05	6645.1\\
};
\addplot [color=mycolor58, forget plot]
  table[row sep=crcr]{%
0	4395.3\\
2000	4152.2\\
4000	3907.5\\
6000	3665.1\\
8000	3575.4\\
10000	3580.5\\
12000	3564.6\\
14000	3653.9\\
16000	3569.4\\
18000	3629.7\\
20000	3520\\
22000	3555.7\\
24000	3570.8\\
26000	3748.5\\
30000	3594.2\\
32000	3614.4\\
36000	3596.9\\
38000	3558.5\\
40000	3582.6\\
42000	3538.8\\
44000	3548.6\\
46000	3589.6\\
48000	3647.5\\
50000	3675.2\\
52000	3594.2\\
54000	3530.1\\
56000	3585\\
58000	3603.2\\
62000	3538.5\\
64000	3664.4\\
66000	3694.5\\
68000	3740.4\\
72000	3622.4\\
74000	3587\\
76000	3637.8\\
78000	3608.6\\
80000	3609.3\\
84000	3554.7\\
86000	3629.1\\
88000	3679.4\\
90000	3680\\
92000	3703.5\\
94000	3645.9\\
96000	3736.5\\
98000	3627.1\\
1e+05	3600.6\\
};
\addplot [color=mycolor59, forget plot]
  table[row sep=crcr]{%
0	4395.3\\
2000	4627.5\\
4000	4850.8\\
6000	4964.2\\
8000	5021.2\\
10000	4996\\
12000	4952.7\\
14000	4991.8\\
16000	4975.9\\
18000	5043\\
20000	4952.2\\
22000	5062.7\\
24000	5138.9\\
26000	5099.3\\
28000	5029.3\\
30000	5150\\
32000	5143.6\\
34000	5021\\
36000	5190.4\\
38000	5054\\
40000	5075.8\\
42000	5181.5\\
44000	4969\\
48000	4949.4\\
50000	5006.3\\
52000	5086.6\\
54000	5064.6\\
56000	5186.1\\
58000	5111.3\\
60000	5202.8\\
62000	5013.3\\
64000	5103.2\\
66000	4953.3\\
68000	5079.6\\
70000	5094.3\\
72000	5001.7\\
74000	5113.7\\
78000	5135.4\\
80000	5072.8\\
82000	5175.4\\
84000	5164.9\\
86000	5060.9\\
88000	5132\\
90000	5037.6\\
92000	5041\\
94000	5029.9\\
96000	5160.8\\
98000	5114.8\\
1e+05	5090.4\\
};
\addplot [color=mycolor60, forget plot]
  table[row sep=crcr]{%
0	4395.3\\
2000	3856.1\\
4000	3727.4\\
6000	3690.1\\
8000	3696.9\\
10000	3602.5\\
12000	3521.2\\
14000	3585\\
16000	3578.6\\
18000	3623.5\\
20000	3649.8\\
22000	3566.7\\
24000	3543.4\\
26000	3553.8\\
28000	3598.8\\
30000	3603.8\\
32000	3658.2\\
36000	3568.9\\
38000	3617.6\\
40000	3628.7\\
42000	3579.7\\
44000	3651\\
46000	3638.1\\
48000	3572\\
50000	3555.4\\
52000	3563.8\\
54000	3583.4\\
58000	3539.3\\
60000	3665.9\\
62000	3649.1\\
64000	3586.5\\
66000	3580.7\\
68000	3502.2\\
72000	3659.3\\
74000	3618\\
76000	3675.4\\
78000	3604.7\\
80000	3632.9\\
82000	3600.8\\
84000	3649.3\\
86000	3668.2\\
88000	3660.5\\
90000	3619.3\\
92000	3595\\
94000	3494.9\\
96000	3572.5\\
98000	3579.4\\
1e+05	3669\\
};
\addplot [color=mycolor61, forget plot]
  table[row sep=crcr]{%
0	4395.3\\
2000	3461.7\\
4000	3224.1\\
8000	3140.2\\
10000	3115.8\\
12000	3129.7\\
14000	3154.3\\
16000	3211.5\\
18000	3168.1\\
22000	3143.7\\
24000	3180.6\\
26000	3192.7\\
28000	3233.5\\
30000	3126.7\\
32000	3115.8\\
34000	3145.9\\
36000	3195.4\\
38000	3169.2\\
40000	3186.7\\
42000	3233.7\\
44000	3157.8\\
48000	3192.4\\
50000	3138.7\\
52000	3144.2\\
54000	3159.1\\
56000	3192.5\\
58000	3164.2\\
60000	3178.9\\
62000	3108.3\\
64000	3122.9\\
66000	3216.4\\
68000	3181.6\\
70000	3196.1\\
72000	3179.2\\
74000	3175.5\\
76000	3106\\
78000	3184.1\\
80000	3209.4\\
82000	3168.9\\
84000	3187.1\\
86000	3227.4\\
88000	3206.5\\
90000	3195.8\\
92000	3160.5\\
94000	3156.8\\
96000	3162\\
98000	3117\\
1e+05	3135.8\\
};
\addplot [color=mycolor62, forget plot]
  table[row sep=crcr]{%
0	4395.3\\
2000	4515.8\\
4000	4898.5\\
6000	5016\\
8000	5008.9\\
10000	4986.6\\
12000	5045.3\\
16000	4953.3\\
18000	4910.5\\
20000	5052.2\\
22000	5074.4\\
24000	4963.1\\
26000	5040\\
28000	4966.5\\
30000	4977.3\\
32000	5007.7\\
34000	4994.4\\
36000	4956.4\\
38000	4989.7\\
40000	4989.8\\
42000	5028.7\\
44000	4963.7\\
46000	4975.1\\
48000	5074.9\\
50000	5006.1\\
52000	4983.2\\
54000	5055.5\\
56000	4994.4\\
58000	4977.9\\
60000	4930.3\\
62000	4981.1\\
64000	5097\\
66000	5063.5\\
68000	4935.6\\
70000	4982.7\\
72000	4985.1\\
74000	5008.8\\
76000	4933.7\\
80000	4967.2\\
82000	4984.9\\
84000	5023.8\\
86000	4916.2\\
88000	4999.5\\
90000	4956.1\\
92000	4986.4\\
96000	4996.6\\
98000	4984.3\\
1e+05	5019.5\\
};
\addplot [color=mycolor63, forget plot]
  table[row sep=crcr]{%
0	4395.3\\
2000	3237.4\\
4000	2683.2\\
6000	2523.5\\
10000	2538.3\\
12000	2486.5\\
14000	2455.4\\
18000	2414.6\\
20000	2449\\
22000	2465.1\\
24000	2465.1\\
26000	2485.4\\
28000	2465\\
30000	2468.4\\
32000	2512.2\\
34000	2488.6\\
36000	2473.6\\
38000	2485.9\\
40000	2508.6\\
42000	2507.8\\
44000	2462.3\\
46000	2479.3\\
48000	2472.6\\
50000	2419.4\\
52000	2445.9\\
54000	2414.3\\
56000	2429.8\\
58000	2464.7\\
60000	2457.2\\
62000	2438.3\\
64000	2447.3\\
66000	2497.3\\
68000	2502.5\\
70000	2472.8\\
72000	2483\\
74000	2443.3\\
76000	2412.8\\
78000	2477.2\\
80000	2465.6\\
82000	2477.6\\
84000	2465.1\\
86000	2496.2\\
88000	2462.2\\
90000	2443.3\\
92000	2462.9\\
94000	2452.6\\
96000	2468.9\\
98000	2399.3\\
1e+05	2458.9\\
};
\addplot [color=mycolor64, forget plot]
  table[row sep=crcr]{%
0	4395.3\\
2000	3340.9\\
4000	3127.1\\
6000	3055.9\\
8000	3113.7\\
10000	3066.3\\
12000	3045.1\\
14000	3052.2\\
16000	3099.6\\
18000	3046.6\\
20000	3103.9\\
22000	3138.8\\
24000	3057.6\\
28000	3105.9\\
30000	3106.1\\
32000	3065.9\\
34000	3100.5\\
36000	3048.5\\
38000	3083.5\\
40000	3081.4\\
42000	3058\\
44000	3100.6\\
46000	3098.3\\
48000	3064.8\\
50000	3109.4\\
52000	3141.4\\
54000	3102.5\\
58000	3109\\
60000	3097.3\\
62000	3095\\
64000	3078.6\\
66000	3092.7\\
68000	3115.9\\
72000	3103.6\\
74000	3123.6\\
76000	3079.4\\
78000	3106.7\\
82000	3121.9\\
84000	3062.7\\
86000	3074.5\\
90000	3122.1\\
92000	3074.2\\
94000	3087.8\\
96000	3057.2\\
98000	3070.2\\
1e+05	3065.2\\
};
\addplot [color=mycolor65, forget plot]
  table[row sep=crcr]{%
0	4395.3\\
2000	3185.2\\
4000	2643.9\\
6000	2624\\
8000	2530.8\\
10000	2533.6\\
12000	2519.5\\
14000	2559.8\\
16000	2566.7\\
18000	2560.8\\
20000	2610.7\\
22000	2595.6\\
26000	2595.2\\
30000	2632.2\\
32000	2567.1\\
34000	2586.6\\
36000	2649.2\\
38000	2602.3\\
42000	2593.5\\
44000	2624.4\\
46000	2599.9\\
50000	2565.7\\
52000	2534.5\\
54000	2597.2\\
56000	2622.8\\
58000	2628.4\\
60000	2590.6\\
62000	2567.8\\
64000	2598.6\\
66000	2589.1\\
68000	2609.7\\
70000	2586.3\\
72000	2590.8\\
74000	2618.7\\
76000	2588.5\\
78000	2586.3\\
80000	2620.1\\
82000	2567.2\\
84000	2588.6\\
86000	2559.4\\
88000	2556.9\\
90000	2615.6\\
94000	2580.9\\
96000	2592.7\\
98000	2615.7\\
1e+05	2587.8\\
};
\addplot [color=mycolor66, forget plot]
  table[row sep=crcr]{%
0	4395.3\\
2000	3592.2\\
4000	3400.3\\
6000	3461.6\\
8000	3411.7\\
10000	3420.1\\
12000	3439.5\\
14000	3441\\
16000	3460.7\\
18000	3451\\
20000	3460.3\\
22000	3489\\
24000	3478.5\\
26000	3488.1\\
28000	3526.1\\
30000	3467.2\\
32000	3484.9\\
34000	3517.4\\
36000	3438.8\\
38000	3488.8\\
40000	3457.8\\
44000	3475.3\\
46000	3457.1\\
48000	3505.9\\
50000	3435.6\\
52000	3481.6\\
54000	3441.5\\
56000	3446.4\\
58000	3463\\
60000	3465.5\\
62000	3516.8\\
64000	3440.1\\
66000	3505.9\\
68000	3479.1\\
70000	3433.6\\
72000	3425.6\\
76000	3452.2\\
78000	3424.9\\
80000	3414.7\\
82000	3485.6\\
84000	3488.8\\
86000	3448.6\\
88000	3482.3\\
92000	3456.1\\
94000	3503.7\\
96000	3466.6\\
98000	3507.8\\
1e+05	3487.1\\
};
\addplot [color=mycolor67, forget plot]
  table[row sep=crcr]{%
0	4395.3\\
2000	3695.8\\
4000	3444.1\\
8000	3297.3\\
10000	3341.1\\
12000	3255.2\\
14000	3213.7\\
16000	3234.2\\
18000	3200.6\\
20000	3249.9\\
22000	3276.4\\
24000	3212.6\\
26000	3221.3\\
28000	3267.8\\
30000	3228.1\\
32000	3173.5\\
34000	3147.7\\
36000	3177.4\\
38000	3226.4\\
40000	3240.2\\
44000	3207.9\\
46000	3260.2\\
48000	3248.6\\
50000	3250.5\\
52000	3214.7\\
54000	3228.2\\
56000	3272.2\\
58000	3286.1\\
60000	3242.5\\
62000	3221\\
66000	3251.4\\
68000	3238.5\\
70000	3248.2\\
72000	3190.5\\
74000	3268\\
76000	3258.4\\
78000	3217.6\\
82000	3195.7\\
84000	3272\\
88000	3208.5\\
90000	3222.4\\
92000	3248.2\\
94000	3221.8\\
98000	3193.8\\
1e+05	3267.9\\
};
\addplot [color=mycolor67]
  table[row sep=crcr]{%
0	4395.3\\
2000	3024.2\\
4000	2683.4\\
6000	2684.1\\
8000	2692.8\\
10000	2627.6\\
12000	2638.8\\
14000	2680.3\\
16000	2689.1\\
18000	2687.9\\
20000	2637.4\\
22000	2677.9\\
24000	2677.9\\
26000	2704.1\\
28000	2677.4\\
30000	2685.5\\
32000	2708.4\\
34000	2699.6\\
36000	2671.9\\
38000	2678.6\\
40000	2638.7\\
42000	2688.1\\
46000	2685\\
48000	2668.9\\
50000	2666.4\\
52000	2652.3\\
54000	2645.6\\
56000	2675.3\\
58000	2676.7\\
62000	2649.3\\
64000	2687.8\\
66000	2650\\
68000	2639.4\\
70000	2662.1\\
72000	2662.1\\
74000	2674.5\\
76000	2696.4\\
80000	2642.2\\
82000	2664\\
86000	2677.1\\
88000	2712.5\\
90000	2665.3\\
92000	2645.9\\
96000	2696.2\\
98000	2678\\
1e+05	2636.7\\
};
\addlegendentry{Classe k=43}

\addplot [color=mycolor68, forget plot]
  table[row sep=crcr]{%
0	4395.3\\
2000	3263.6\\
4000	2733.9\\
6000	2636.1\\
8000	2579.9\\
12000	2531\\
14000	2517.2\\
16000	2545\\
18000	2524.4\\
20000	2482\\
22000	2503.9\\
24000	2495.6\\
26000	2494.4\\
28000	2504\\
30000	2490.9\\
32000	2503.3\\
34000	2500.3\\
36000	2516.8\\
38000	2505.1\\
42000	2500.7\\
44000	2491.8\\
48000	2537.6\\
50000	2511.9\\
52000	2515.3\\
56000	2502.1\\
58000	2493.9\\
62000	2529.2\\
64000	2488.7\\
66000	2519.8\\
70000	2526.8\\
72000	2486.9\\
74000	2526\\
76000	2529\\
78000	2501.3\\
84000	2487.7\\
86000	2532.9\\
88000	2523.7\\
90000	2526.4\\
92000	2501\\
94000	2505.9\\
98000	2502\\
1e+05	2507.1\\
};
\addplot [color=mycolor69, forget plot]
  table[row sep=crcr]{%
0	4395.3\\
2000	3423.4\\
4000	2964\\
6000	2907.4\\
8000	2887.3\\
10000	2901.8\\
12000	2945.9\\
14000	2952.1\\
16000	2913\\
18000	2937.1\\
22000	2959.9\\
24000	2951.2\\
26000	2909\\
30000	2945.3\\
32000	2937.9\\
34000	2945.1\\
36000	3001.3\\
38000	2956.2\\
40000	2973.1\\
42000	2947\\
44000	2959.4\\
46000	2935.5\\
50000	2916.2\\
52000	2984.5\\
54000	3012.4\\
56000	2971.2\\
58000	2942.1\\
60000	2927.7\\
62000	2970.1\\
64000	2976.8\\
66000	2950.7\\
70000	2963.2\\
72000	2920.7\\
74000	2898.9\\
76000	2961.1\\
80000	2923.2\\
84000	2927.9\\
86000	2915.1\\
88000	2975.2\\
90000	2999.5\\
92000	2960.1\\
94000	2965.4\\
96000	2952.7\\
98000	2973\\
1e+05	2936.6\\
};
\addplot [color=mycolor70, forget plot]
  table[row sep=crcr]{%
0	4395.3\\
2000	2775.7\\
4000	2109.3\\
6000	2093.7\\
8000	2108.2\\
10000	2084.2\\
14000	2163.3\\
18000	2137.1\\
20000	2151\\
22000	2156.9\\
24000	2123.8\\
26000	2164.4\\
30000	2133.2\\
32000	2170.3\\
34000	2133.4\\
36000	2178.4\\
38000	2184\\
40000	2165.6\\
42000	2226.3\\
44000	2168.6\\
46000	2179.7\\
48000	2124.3\\
50000	2148.3\\
52000	2140.4\\
54000	2166.1\\
58000	2166.2\\
60000	2176.6\\
62000	2127.8\\
64000	2176.8\\
66000	2183.5\\
68000	2136.1\\
70000	2115.2\\
72000	2139\\
74000	2154.2\\
76000	2115.9\\
78000	2147.4\\
80000	2153.7\\
82000	2130.2\\
84000	2163.9\\
86000	2169.9\\
88000	2157.9\\
90000	2166\\
92000	2140.3\\
94000	2127.3\\
96000	2143.2\\
98000	2165.9\\
1e+05	2148.4\\
};
\addplot [color=mycolor71, forget plot]
  table[row sep=crcr]{%
0	4395.3\\
2000	3125.6\\
4000	2676.9\\
6000	2621\\
10000	2579.8\\
12000	2547.2\\
14000	2545.8\\
16000	2532.2\\
18000	2555.5\\
20000	2556.2\\
22000	2541.5\\
24000	2542.6\\
26000	2572.6\\
28000	2535.4\\
30000	2533.7\\
32000	2515.7\\
34000	2518.5\\
36000	2510.4\\
38000	2531.6\\
40000	2533.5\\
42000	2508.4\\
44000	2532.4\\
46000	2512.5\\
48000	2522.5\\
50000	2560.9\\
52000	2522.9\\
54000	2528.9\\
58000	2527.3\\
60000	2534.9\\
62000	2515.4\\
64000	2527.8\\
66000	2508.9\\
68000	2504.1\\
70000	2526.8\\
72000	2526.1\\
76000	2553.9\\
80000	2502.6\\
82000	2510.7\\
84000	2551.9\\
86000	2537\\
92000	2546.4\\
94000	2535.5\\
96000	2536\\
98000	2515.4\\
1e+05	2511.4\\
};
\addplot [color=mycolor72, forget plot]
  table[row sep=crcr]{%
0	4395.3\\
2000	3027.2\\
4000	2342.5\\
6000	2250.1\\
8000	2248.9\\
10000	2228.6\\
12000	2261.2\\
14000	2268.5\\
16000	2238.7\\
18000	2253.9\\
20000	2230.4\\
22000	2264.2\\
24000	2261.4\\
26000	2237.5\\
30000	2225.7\\
32000	2198.5\\
34000	2232.6\\
36000	2284\\
38000	2264.9\\
40000	2257.5\\
42000	2267\\
44000	2246.8\\
46000	2238.5\\
48000	2254.2\\
50000	2237.1\\
52000	2257.2\\
54000	2244.2\\
56000	2225\\
58000	2250.7\\
60000	2234.4\\
62000	2241.8\\
64000	2268.1\\
68000	2219.9\\
70000	2226.2\\
72000	2223.1\\
74000	2278.5\\
76000	2269.5\\
78000	2192.5\\
80000	2224.5\\
82000	2245.4\\
84000	2273.1\\
86000	2270.5\\
88000	2223.4\\
90000	2279.5\\
92000	2252.1\\
94000	2275.8\\
96000	2251.6\\
1e+05	2300.5\\
};
\addplot [color=mycolor73, forget plot]
  table[row sep=crcr]{%
0	4395.3\\
2000	3200.3\\
4000	2791.7\\
6000	2705.1\\
8000	2725.1\\
10000	2689\\
12000	2670.6\\
16000	2688.8\\
18000	2679.7\\
20000	2649.4\\
22000	2684\\
24000	2681\\
26000	2702.4\\
28000	2653.5\\
30000	2661.5\\
32000	2650.5\\
34000	2629.9\\
36000	2634.9\\
40000	2689.3\\
42000	2622.6\\
44000	2614.4\\
46000	2627.1\\
48000	2622.6\\
50000	2610.5\\
52000	2620.7\\
54000	2621.2\\
58000	2653.9\\
64000	2649.1\\
68000	2672.7\\
70000	2646.2\\
74000	2693.2\\
76000	2651.4\\
78000	2667.7\\
80000	2658.3\\
82000	2663.8\\
84000	2688.8\\
86000	2677\\
88000	2652.9\\
90000	2668.8\\
92000	2618\\
94000	2669.6\\
96000	2625.4\\
98000	2574.2\\
1e+05	2625.6\\
};
\addplot [color=mycolor74, forget plot]
  table[row sep=crcr]{%
0	4395.3\\
2000	2667.9\\
4000	1786.5\\
6000	1646.7\\
8000	1600.7\\
10000	1583.2\\
12000	1576.2\\
14000	1598\\
16000	1589.1\\
18000	1602.4\\
20000	1580.6\\
22000	1587.7\\
24000	1572.9\\
26000	1565.2\\
28000	1572.7\\
30000	1589.6\\
32000	1588.1\\
34000	1571.5\\
36000	1595.5\\
38000	1576.3\\
40000	1598.7\\
42000	1603\\
44000	1585.1\\
46000	1592\\
48000	1608.3\\
50000	1586.5\\
54000	1593.5\\
56000	1612.1\\
58000	1600.6\\
60000	1604.2\\
62000	1579.8\\
64000	1580.9\\
66000	1589.9\\
68000	1570\\
70000	1602.6\\
72000	1603\\
74000	1587.1\\
78000	1587.5\\
80000	1585.1\\
82000	1568.8\\
84000	1589.4\\
86000	1585.1\\
88000	1598.9\\
92000	1568.3\\
94000	1562.4\\
96000	1578.5\\
98000	1570.2\\
1e+05	1588.1\\
};
\addplot [color=mycolor75, forget plot]
  table[row sep=crcr]{%
0	4395.3\\
2000	3089.4\\
4000	2404\\
6000	2321.5\\
8000	2265.1\\
10000	2269.8\\
12000	2262.2\\
14000	2282.4\\
16000	2237.7\\
20000	2225.3\\
22000	2245.9\\
26000	2245.4\\
28000	2281.5\\
30000	2277.7\\
32000	2283.5\\
36000	2246.3\\
38000	2251.9\\
40000	2248.6\\
42000	2263.8\\
44000	2268.3\\
46000	2285.7\\
48000	2294.1\\
50000	2244.3\\
52000	2269.9\\
58000	2244.8\\
60000	2252.7\\
62000	2274.8\\
64000	2265.5\\
66000	2268\\
68000	2233.7\\
72000	2260.4\\
74000	2254.7\\
76000	2277.9\\
80000	2261.7\\
82000	2243.9\\
84000	2254.1\\
86000	2245.9\\
88000	2222.8\\
90000	2249.1\\
94000	2258.2\\
96000	2213\\
98000	2246.2\\
1e+05	2251.5\\
};
\addplot [color=mycolor76, forget plot]
  table[row sep=crcr]{%
0	4395.3\\
2000	2902.1\\
4000	2243.6\\
6000	2163.6\\
8000	2146.1\\
10000	2092.8\\
12000	2074.4\\
14000	2070.9\\
16000	2074.2\\
18000	2055.2\\
20000	2046.4\\
22000	2078.7\\
26000	2088.5\\
28000	2075.9\\
30000	2086\\
32000	2057.8\\
34000	2070.6\\
36000	2031.8\\
38000	2031.6\\
40000	2054.7\\
42000	2039.8\\
44000	2055.1\\
46000	2034.6\\
48000	2034.5\\
50000	2071.6\\
52000	2082.5\\
54000	2055.2\\
56000	2034.9\\
60000	2067.4\\
62000	2059.8\\
64000	2092.8\\
66000	2103.7\\
68000	2069.2\\
70000	2071.1\\
72000	2086.7\\
74000	2064.1\\
76000	2066.5\\
78000	2086.9\\
82000	2045\\
84000	2063.3\\
86000	2057.4\\
88000	2060.6\\
94000	2044.5\\
96000	2063.6\\
98000	2055.6\\
1e+05	2022.9\\
};
\addplot [color=mycolor77, forget plot]
  table[row sep=crcr]{%
0	4395.3\\
2000	2697.3\\
4000	1862\\
6000	1731.9\\
8000	1684.8\\
10000	1660.5\\
12000	1622.6\\
14000	1604.9\\
16000	1612.9\\
18000	1614.1\\
20000	1622.1\\
22000	1602.4\\
24000	1596.5\\
26000	1612.1\\
28000	1611.3\\
32000	1628.1\\
34000	1630.8\\
38000	1608.9\\
40000	1604.8\\
42000	1629.3\\
44000	1596.8\\
46000	1609.4\\
48000	1627.8\\
50000	1628.3\\
52000	1607.5\\
54000	1602.4\\
58000	1623.1\\
60000	1649.3\\
62000	1586\\
66000	1609.2\\
68000	1612.8\\
70000	1592\\
76000	1632.6\\
78000	1610.4\\
80000	1619.5\\
82000	1600.9\\
84000	1609.4\\
86000	1600.4\\
88000	1601.1\\
90000	1641.3\\
92000	1616.2\\
96000	1606.2\\
98000	1606.2\\
1e+05	1612.2\\
};
\addplot [color=mycolor78, forget plot]
  table[row sep=crcr]{%
0	4395.3\\
2000	3165.3\\
4000	2543.5\\
6000	2469\\
8000	2379.5\\
10000	2360.8\\
12000	2350\\
14000	2374.9\\
16000	2337\\
18000	2336.4\\
20000	2321\\
22000	2375\\
24000	2361.9\\
26000	2338.1\\
28000	2331.7\\
30000	2332.5\\
32000	2374.9\\
34000	2363.4\\
36000	2358.6\\
38000	2378.5\\
40000	2342.5\\
42000	2359\\
44000	2330.6\\
46000	2338.6\\
50000	2328.4\\
54000	2334.9\\
56000	2297.3\\
58000	2319.5\\
60000	2305.6\\
62000	2356\\
64000	2324\\
66000	2337.1\\
68000	2331.6\\
70000	2355.8\\
74000	2356.1\\
76000	2297.9\\
80000	2358\\
84000	2341.1\\
86000	2321\\
88000	2322\\
90000	2341.2\\
94000	2354.7\\
96000	2336.7\\
1e+05	2327.6\\
};
\addplot [color=mycolor79, forget plot]
  table[row sep=crcr]{%
0	4395.3\\
2000	2973.4\\
4000	2364.8\\
6000	2306.8\\
8000	2259.8\\
10000	2206\\
12000	2185.6\\
16000	2216.3\\
22000	2195.9\\
26000	2194.8\\
30000	2231.8\\
32000	2213\\
34000	2229.6\\
36000	2218.2\\
38000	2178.9\\
40000	2167.6\\
42000	2195.5\\
44000	2198.8\\
48000	2217.5\\
50000	2171.9\\
52000	2167.4\\
54000	2191.6\\
56000	2174.2\\
58000	2164.5\\
60000	2221\\
62000	2197.3\\
64000	2215.8\\
66000	2205.9\\
68000	2176.1\\
70000	2128.2\\
72000	2199.3\\
74000	2229.1\\
76000	2206.8\\
78000	2234.2\\
80000	2202.8\\
82000	2219.1\\
84000	2179.3\\
86000	2168.4\\
90000	2208.7\\
92000	2207.9\\
94000	2189.2\\
96000	2193.6\\
98000	2187.7\\
1e+05	2208.6\\
};
\addplot [color=mycolor79, forget plot]
  table[row sep=crcr]{%
0	4395.3\\
2000	3022.6\\
4000	2292.9\\
6000	2167\\
8000	2068.1\\
10000	2052.2\\
12000	2059.7\\
14000	2049.8\\
16000	2087\\
20000	2061.8\\
22000	2100.2\\
24000	2057.8\\
26000	2074.8\\
28000	2069.5\\
30000	2071.3\\
32000	2052.3\\
34000	2053.9\\
36000	2070.6\\
38000	2044.9\\
40000	2075.9\\
42000	2098.3\\
44000	2059.8\\
46000	2088.3\\
48000	2067.5\\
50000	2036.2\\
52000	2053.7\\
54000	2038.5\\
56000	2043.3\\
58000	2061.8\\
60000	2059.6\\
62000	2076.3\\
64000	2085.3\\
66000	2056.2\\
68000	2058.3\\
70000	2088\\
72000	2090.7\\
74000	2044.4\\
76000	2036.5\\
80000	2074.8\\
82000	2057.1\\
86000	2078.4\\
88000	2077.5\\
90000	2052.4\\
96000	2056.5\\
98000	2050.2\\
1e+05	2029.6\\
};
\addplot [color=mycolor80, forget plot]
  table[row sep=crcr]{%
0	4395.3\\
2000	2604.9\\
4000	1753\\
6000	1659.5\\
8000	1646.9\\
10000	1611.7\\
12000	1605.8\\
14000	1609.7\\
16000	1630.4\\
18000	1619.5\\
22000	1585.2\\
24000	1582.2\\
26000	1604\\
30000	1603.8\\
32000	1593.1\\
34000	1619.9\\
36000	1603.4\\
38000	1625.5\\
40000	1590.2\\
42000	1603.2\\
44000	1590.4\\
46000	1594.1\\
48000	1591.8\\
50000	1606.9\\
52000	1605.6\\
54000	1580.6\\
56000	1584.1\\
58000	1609.2\\
60000	1600.5\\
62000	1604.2\\
64000	1615.7\\
66000	1606\\
68000	1611.3\\
70000	1608.4\\
72000	1584.9\\
74000	1593\\
76000	1594.5\\
78000	1627.9\\
80000	1587.8\\
82000	1575\\
84000	1595.1\\
86000	1602.9\\
88000	1595.2\\
90000	1604.9\\
92000	1624.4\\
94000	1599.9\\
96000	1590.1\\
98000	1573.3\\
1e+05	1586.4\\
};
\addplot [color=mycolor81, forget plot]
  table[row sep=crcr]{%
0	4395.3\\
2000	2596.9\\
4000	1691.1\\
6000	1540.7\\
8000	1492.1\\
10000	1519.1\\
12000	1511.7\\
14000	1496.1\\
16000	1509.2\\
18000	1498.9\\
22000	1494.3\\
26000	1492.9\\
28000	1497.2\\
30000	1488.5\\
32000	1502.2\\
34000	1488.6\\
36000	1501.1\\
38000	1501.1\\
40000	1480\\
42000	1486.4\\
44000	1480.9\\
46000	1504\\
50000	1459.9\\
52000	1469\\
54000	1484.8\\
56000	1485.4\\
58000	1497.5\\
60000	1487.6\\
62000	1485.9\\
64000	1510.9\\
66000	1516.6\\
68000	1494.5\\
70000	1488.2\\
72000	1511.2\\
74000	1504.6\\
76000	1471\\
78000	1506.7\\
80000	1508.1\\
86000	1475.3\\
88000	1509.5\\
90000	1488.7\\
92000	1475.1\\
94000	1493\\
98000	1494.3\\
1e+05	1488.6\\
};
\addplot [color=mycolor82, forget plot]
  table[row sep=crcr]{%
0	4395.3\\
2000	2783.4\\
4000	1977.7\\
6000	1843.6\\
8000	1800.1\\
10000	1795.5\\
12000	1766.1\\
14000	1762.3\\
16000	1747\\
18000	1744.2\\
20000	1716.4\\
22000	1719.3\\
24000	1733.4\\
26000	1734.6\\
28000	1746.3\\
30000	1734.1\\
34000	1768.7\\
38000	1741.8\\
40000	1711.5\\
44000	1742.9\\
46000	1745.5\\
48000	1735.9\\
50000	1735.2\\
54000	1749.1\\
56000	1734.3\\
60000	1757.2\\
62000	1735.4\\
64000	1728.2\\
66000	1707.6\\
68000	1742.7\\
70000	1748.1\\
72000	1731.6\\
76000	1718\\
78000	1734\\
80000	1741.1\\
82000	1729.5\\
84000	1739.5\\
86000	1734.3\\
88000	1717.3\\
90000	1729.2\\
92000	1724.6\\
94000	1734.5\\
96000	1739.4\\
98000	1733.6\\
1e+05	1737.3\\
};
\addplot [color=mycolor83, forget plot]
  table[row sep=crcr]{%
0	4395.3\\
2000	2883.7\\
4000	2027.6\\
6000	1878.9\\
8000	1832.7\\
10000	1808.2\\
14000	1806.6\\
16000	1803.2\\
18000	1777.9\\
20000	1761\\
22000	1784.2\\
24000	1757.6\\
26000	1777.6\\
28000	1786.7\\
32000	1758.7\\
34000	1757.7\\
36000	1749.4\\
38000	1766.3\\
42000	1744.4\\
44000	1777.4\\
46000	1803.2\\
50000	1772.7\\
54000	1774.1\\
56000	1776.9\\
58000	1765.6\\
60000	1773.6\\
62000	1787.8\\
64000	1756\\
66000	1778.1\\
68000	1770.7\\
74000	1772.5\\
76000	1782.4\\
78000	1782.7\\
80000	1766.9\\
82000	1781.7\\
84000	1789.1\\
86000	1763.8\\
88000	1753.1\\
90000	1780.9\\
92000	1786\\
94000	1765.1\\
96000	1756.2\\
98000	1771\\
1e+05	1775.5\\
};
\addplot [color=mycolor83, forget plot]
  table[row sep=crcr]{%
0	4395.3\\
2000	2392\\
4000	1439.5\\
6000	1333.4\\
8000	1329.3\\
10000	1316.6\\
12000	1321.1\\
16000	1311.4\\
18000	1317.7\\
20000	1335.8\\
22000	1315.8\\
24000	1308.4\\
26000	1292.6\\
30000	1310.5\\
32000	1298.9\\
34000	1309.9\\
36000	1298.2\\
38000	1298.4\\
42000	1311.6\\
44000	1310.4\\
46000	1317.5\\
48000	1312.7\\
50000	1317.7\\
52000	1310.3\\
54000	1332.5\\
58000	1289.5\\
60000	1300\\
62000	1298.4\\
64000	1313.7\\
66000	1318.8\\
68000	1299.2\\
70000	1308.9\\
72000	1298.2\\
74000	1295.9\\
76000	1320.4\\
78000	1315.9\\
80000	1302.3\\
82000	1294.2\\
84000	1282.3\\
86000	1287.2\\
88000	1299.6\\
94000	1294.3\\
96000	1286.7\\
98000	1326.7\\
1e+05	1313.3\\
};
\addplot [color=mycolor84, forget plot]
  table[row sep=crcr]{%
0	4395.3\\
2000	2502.7\\
4000	1508.1\\
6000	1342.7\\
8000	1346.1\\
10000	1340.1\\
12000	1310.5\\
14000	1324.1\\
20000	1323.1\\
22000	1317.4\\
24000	1331.5\\
26000	1291.2\\
28000	1326\\
30000	1316.6\\
32000	1314.2\\
34000	1331.3\\
36000	1317\\
38000	1329\\
40000	1329.8\\
42000	1316.8\\
44000	1329\\
46000	1279.5\\
48000	1306.2\\
50000	1296.4\\
54000	1292.5\\
56000	1287.5\\
58000	1311.1\\
60000	1301.2\\
62000	1306.2\\
64000	1303.9\\
66000	1325.7\\
68000	1330.6\\
70000	1313.8\\
72000	1334.3\\
76000	1287.8\\
78000	1328.5\\
80000	1340.8\\
82000	1328.9\\
86000	1323.3\\
88000	1306.4\\
90000	1325.6\\
92000	1296\\
94000	1324.3\\
96000	1290.7\\
98000	1317\\
1e+05	1314\\
};
\addplot [color=mycolor85, forget plot]
  table[row sep=crcr]{%
0	4395.3\\
2000	2726.1\\
4000	1712\\
6000	1512.7\\
8000	1509.1\\
10000	1465\\
12000	1453.9\\
14000	1451.6\\
16000	1438.8\\
18000	1450.7\\
20000	1479.5\\
22000	1486.1\\
24000	1462.3\\
26000	1459.2\\
28000	1437\\
30000	1434.1\\
32000	1460.6\\
34000	1450.9\\
36000	1458.5\\
38000	1452.4\\
40000	1414.7\\
44000	1452.1\\
46000	1448.9\\
48000	1463.3\\
50000	1426.7\\
52000	1425.7\\
54000	1447.9\\
56000	1391.7\\
60000	1399.7\\
62000	1416.2\\
64000	1439.6\\
66000	1415\\
68000	1457.4\\
70000	1416.9\\
72000	1434.3\\
74000	1421.3\\
76000	1416.4\\
78000	1434.4\\
80000	1433.6\\
82000	1408.1\\
84000	1409.1\\
86000	1420.7\\
88000	1455.2\\
90000	1412.7\\
92000	1444.5\\
94000	1446.5\\
96000	1427\\
98000	1417.8\\
1e+05	1434.5\\
};
\addplot [color=mycolor86, forget plot]
  table[row sep=crcr]{%
0	4395.3\\
2000	2541.1\\
4000	1526.3\\
6000	1354.3\\
8000	1301.8\\
10000	1266.2\\
12000	1221\\
14000	1248.9\\
16000	1220.9\\
20000	1229.7\\
22000	1208.4\\
24000	1207.7\\
26000	1215.7\\
28000	1205.1\\
32000	1218.8\\
34000	1212.3\\
36000	1237.6\\
38000	1216\\
40000	1215.6\\
42000	1250.6\\
44000	1227\\
46000	1222.9\\
48000	1213.7\\
52000	1225.1\\
54000	1236.5\\
56000	1241.3\\
58000	1238.7\\
60000	1213\\
62000	1226.4\\
64000	1208.3\\
66000	1204.5\\
68000	1191.6\\
70000	1191.2\\
74000	1225.5\\
78000	1192.9\\
82000	1213.2\\
84000	1225.3\\
86000	1212.3\\
88000	1237\\
90000	1245.9\\
92000	1227.4\\
94000	1205.3\\
96000	1212.1\\
98000	1222.8\\
1e+05	1196.3\\
};
\addplot [color=mycolor87, forget plot]
  table[row sep=crcr]{%
0	4395.3\\
2000	2393.8\\
4000	1209.5\\
6000	1053.6\\
8000	1006.3\\
10000	1030.4\\
12000	1020.1\\
14000	1005.5\\
16000	1020.1\\
18000	1027.6\\
20000	1002.2\\
22000	1050.2\\
24000	1053.7\\
26000	1031.1\\
28000	1021.9\\
30000	1035.8\\
32000	1042.6\\
34000	1015.1\\
38000	1030.6\\
40000	1047.2\\
42000	1078.8\\
44000	1024.6\\
48000	992.93\\
50000	1041\\
52000	1036.2\\
54000	1023.5\\
56000	1042.9\\
58000	1074.5\\
60000	1056.3\\
62000	1047.3\\
64000	1069.1\\
66000	1060.2\\
68000	1055.6\\
70000	974.2\\
72000	977.32\\
74000	1021\\
76000	990.17\\
78000	1010.1\\
80000	1077.4\\
82000	1019.9\\
84000	1022.5\\
86000	1032.5\\
88000	1062.4\\
90000	1051.6\\
92000	1035.2\\
94000	1034.7\\
96000	1029.6\\
98000	1034.5\\
1e+05	1013.2\\
};
\addplot [color=mycolor88, forget plot]
  table[row sep=crcr]{%
0	4395.3\\
2000	2616.6\\
4000	1544.1\\
6000	1350.1\\
8000	1312.4\\
10000	1288.4\\
12000	1254.7\\
14000	1239.2\\
16000	1235.2\\
18000	1213.2\\
20000	1213.7\\
22000	1221.6\\
24000	1241.1\\
28000	1208.4\\
30000	1230.3\\
32000	1215.8\\
34000	1214.7\\
36000	1185.7\\
38000	1191.9\\
42000	1217.4\\
44000	1229.3\\
46000	1192.3\\
48000	1206.5\\
52000	1221.3\\
54000	1212.3\\
56000	1217.3\\
58000	1235\\
60000	1217.2\\
62000	1183.4\\
64000	1207.6\\
66000	1196.1\\
68000	1223.1\\
72000	1236.4\\
74000	1231.5\\
76000	1246.3\\
78000	1242.4\\
80000	1251.3\\
82000	1238.9\\
84000	1207.1\\
86000	1198\\
88000	1221\\
92000	1233.9\\
94000	1213.9\\
98000	1214.3\\
1e+05	1236.2\\
};
\addplot [color=mycolor89, forget plot]
  table[row sep=crcr]{%
0	4395.3\\
2000	2507\\
4000	1230\\
6000	968.05\\
8000	911.39\\
10000	898.86\\
12000	892.73\\
14000	890.2\\
16000	895.3\\
18000	889.95\\
20000	911.29\\
22000	897.84\\
24000	874.02\\
26000	869.86\\
28000	889.95\\
32000	886.07\\
34000	875.88\\
36000	861.06\\
38000	880.63\\
40000	888.11\\
42000	890.06\\
44000	883.74\\
46000	889.81\\
48000	872.47\\
50000	893.59\\
52000	883.83\\
54000	889.74\\
56000	901.27\\
58000	892.8\\
60000	893.37\\
62000	883.7\\
64000	895.43\\
68000	872.76\\
72000	865.52\\
74000	866.81\\
76000	874.35\\
78000	873.48\\
84000	882.78\\
86000	878.24\\
88000	899.94\\
92000	876.25\\
94000	876.95\\
96000	891.6\\
98000	879.71\\
1e+05	882.84\\
};
\addplot [color=mycolor90, forget plot]
  table[row sep=crcr]{%
0	4395.3\\
2000	2309.3\\
4000	1075.7\\
6000	892.69\\
8000	861.11\\
10000	845.91\\
12000	854.28\\
20000	839.66\\
22000	862.52\\
24000	838.93\\
26000	862.16\\
28000	861.52\\
30000	852.97\\
32000	856.9\\
34000	853.98\\
36000	868.82\\
38000	858.91\\
42000	850.34\\
44000	859.9\\
46000	854.31\\
48000	857.27\\
50000	825.82\\
52000	856.98\\
54000	841.95\\
56000	861.79\\
58000	855.62\\
60000	864.81\\
62000	826.26\\
64000	846.21\\
66000	847.85\\
68000	853.91\\
70000	878.36\\
74000	843.66\\
76000	837.67\\
78000	849.37\\
80000	837.92\\
82000	849.88\\
84000	869.22\\
86000	850.53\\
88000	847.64\\
90000	867.36\\
92000	876.57\\
94000	838.12\\
96000	853.88\\
98000	878.56\\
1e+05	863.39\\
};
\addplot [color=mycolor91, forget plot]
  table[row sep=crcr]{%
0	4395.3\\
2000	2469.2\\
4000	1149.9\\
6000	860.56\\
8000	786.81\\
10000	768.46\\
14000	742.09\\
16000	719.78\\
18000	739.02\\
20000	739.66\\
22000	744.84\\
24000	766.88\\
26000	748.68\\
28000	728.02\\
30000	730.68\\
32000	739.47\\
34000	738.99\\
36000	729.97\\
38000	756.54\\
40000	744.48\\
42000	742.3\\
44000	749.39\\
46000	728.55\\
48000	755.77\\
50000	758.51\\
52000	741.68\\
54000	757.54\\
56000	738.92\\
58000	733.2\\
60000	767\\
64000	765.88\\
66000	756.95\\
68000	735.13\\
70000	752.71\\
72000	748.93\\
74000	750.99\\
76000	742.83\\
78000	775.28\\
80000	753.43\\
82000	755.24\\
84000	732.7\\
86000	760.83\\
88000	754.45\\
90000	730.96\\
92000	728.09\\
96000	750.79\\
98000	745.47\\
1e+05	749.9\\
};
\addplot [color=mycolor92, forget plot]
  table[row sep=crcr]{%
0	4395.3\\
2000	2521.5\\
4000	1294\\
6000	1028.7\\
8000	965.86\\
10000	944.2\\
12000	935.78\\
14000	947.42\\
16000	923.95\\
18000	921.74\\
20000	908.78\\
22000	921.57\\
24000	928.48\\
26000	918.8\\
28000	912.24\\
30000	924.85\\
32000	948.38\\
34000	929.08\\
36000	928.45\\
38000	921.12\\
40000	910.56\\
42000	904.22\\
44000	922.89\\
46000	936.76\\
48000	923.12\\
50000	921.13\\
52000	923.76\\
54000	929.69\\
56000	916.96\\
58000	919.84\\
60000	909.36\\
62000	938.61\\
64000	943.76\\
68000	920.98\\
70000	932.56\\
72000	922.89\\
74000	932.28\\
76000	928.52\\
78000	939.73\\
80000	926.48\\
82000	931.96\\
86000	928.09\\
88000	915.42\\
90000	921.03\\
94000	954.6\\
96000	924.71\\
98000	918.6\\
1e+05	940.4\\
};
\addplot [color=mycolor93, forget plot]
  table[row sep=crcr]{%
0	4395.3\\
2000	2599.6\\
4000	1259.3\\
6000	911.46\\
8000	820.37\\
10000	783.04\\
12000	772.58\\
14000	782.19\\
16000	812.68\\
18000	798.11\\
20000	764.14\\
22000	763.38\\
24000	772.81\\
26000	766.94\\
28000	758.59\\
30000	793.69\\
32000	774.36\\
34000	784.8\\
36000	757.69\\
38000	769.97\\
40000	784.74\\
42000	769.51\\
44000	809.93\\
46000	804.51\\
48000	775.2\\
50000	788.08\\
52000	784.67\\
54000	830.14\\
56000	797.43\\
58000	753.36\\
60000	769.23\\
62000	759.05\\
64000	768.6\\
66000	792.52\\
68000	771.61\\
70000	753.59\\
72000	733.08\\
74000	764.4\\
76000	754.75\\
78000	786.5\\
80000	797.47\\
82000	812.66\\
84000	799.57\\
86000	796.94\\
88000	806.69\\
90000	764.39\\
92000	780.73\\
96000	757.82\\
98000	770.2\\
1e+05	762.23\\
};
\addplot [color=mycolor94, forget plot]
  table[row sep=crcr]{%
0	4395.3\\
2000	2603.7\\
4000	1228.9\\
6000	901.59\\
8000	785.94\\
10000	740.57\\
12000	730.55\\
14000	731.22\\
16000	749.88\\
18000	758.34\\
20000	757.13\\
22000	747.59\\
24000	740.48\\
26000	746.68\\
30000	739.9\\
32000	740.9\\
34000	744.99\\
36000	752.01\\
38000	763.64\\
40000	755.78\\
42000	743.91\\
44000	747.82\\
46000	735.25\\
50000	747.91\\
52000	726.66\\
54000	732.64\\
56000	753.2\\
58000	732.07\\
60000	738.25\\
62000	733.92\\
64000	746.89\\
66000	745.66\\
68000	724.06\\
70000	743.36\\
72000	735.31\\
74000	738.85\\
76000	727.77\\
78000	722.34\\
80000	733.48\\
82000	726.2\\
84000	727.28\\
86000	750.68\\
88000	753.39\\
90000	735.28\\
92000	734.17\\
94000	722.31\\
96000	717.77\\
98000	733.97\\
1e+05	745.39\\
};
\addplot [color=mycolor95, forget plot]
  table[row sep=crcr]{%
0	4395.3\\
2000	2304.5\\
4000	891.79\\
6000	551.52\\
8000	488.89\\
10000	473.89\\
12000	460.79\\
14000	450.25\\
16000	442.4\\
24000	450.51\\
26000	447.11\\
28000	457.17\\
30000	457.47\\
32000	444.87\\
36000	457.86\\
38000	447.03\\
40000	445.15\\
42000	454.09\\
44000	438.01\\
46000	443.71\\
48000	455.14\\
50000	446.67\\
52000	450.87\\
54000	449.96\\
56000	459.6\\
58000	462.33\\
60000	444.85\\
62000	435.04\\
64000	446.47\\
66000	442\\
68000	443.62\\
70000	455.92\\
72000	451.43\\
74000	442.19\\
76000	445.89\\
78000	453.91\\
80000	437.24\\
82000	433.21\\
84000	455.05\\
86000	447.84\\
88000	449.93\\
90000	439.91\\
92000	437.18\\
94000	445.03\\
96000	435.3\\
98000	446.35\\
1e+05	439.39\\
};
\addplot [color=mycolor96, forget plot]
  table[row sep=crcr]{%
0	4395.3\\
2000	2525.1\\
4000	1066.9\\
6000	647.49\\
8000	524.97\\
10000	494.42\\
12000	463.64\\
14000	455.38\\
16000	458.83\\
18000	465.74\\
20000	453.78\\
22000	461.54\\
24000	442.78\\
26000	459.85\\
30000	463.9\\
32000	471.36\\
34000	460.72\\
36000	468.54\\
38000	486.8\\
40000	483.96\\
42000	451.76\\
44000	451.07\\
46000	462.54\\
48000	468.16\\
50000	461.06\\
52000	488.42\\
54000	494.69\\
56000	452.52\\
58000	438.79\\
60000	456.07\\
62000	466.61\\
64000	462.13\\
66000	478.47\\
68000	486.51\\
70000	479.18\\
72000	480.54\\
74000	493.23\\
76000	479.56\\
78000	459\\
80000	461.98\\
82000	453.89\\
84000	471.47\\
86000	466.14\\
88000	478.19\\
90000	486.22\\
94000	487.62\\
96000	476.47\\
98000	449.3\\
1e+05	470.67\\
};
\addplot [color=mycolor97, forget plot]
  table[row sep=crcr]{%
0	4395.3\\
2000	2124.8\\
4000	745.52\\
6000	473.51\\
8000	386.52\\
10000	359.73\\
12000	355.76\\
14000	339.79\\
16000	330.25\\
18000	344.19\\
20000	344.83\\
22000	330.15\\
24000	346.18\\
26000	353.12\\
28000	327.52\\
30000	335.42\\
32000	316.58\\
34000	326.15\\
36000	329.14\\
38000	328.93\\
40000	351.3\\
42000	341.1\\
44000	339.39\\
46000	352.63\\
48000	345.69\\
50000	332.77\\
52000	346.41\\
54000	347.89\\
56000	342.61\\
58000	338.66\\
60000	337.5\\
62000	352.92\\
64000	344.09\\
66000	344.36\\
68000	323.43\\
70000	338.27\\
72000	346.21\\
74000	333.33\\
76000	338.29\\
80000	327.85\\
82000	316.98\\
84000	323.17\\
86000	333.92\\
88000	314.91\\
90000	319.24\\
92000	329.41\\
94000	330.13\\
96000	347.88\\
98000	342.28\\
1e+05	338.41\\
};
\addplot [color=mycolor98]
  table[row sep=crcr]{%
0	4395.3\\
2000	2622\\
4000	1152\\
6000	697.7\\
8000	551.62\\
10000	490.49\\
12000	464.43\\
14000	457.95\\
16000	474.97\\
18000	458.57\\
20000	452.01\\
22000	440.35\\
24000	435.99\\
26000	427.31\\
28000	423.27\\
30000	422.04\\
32000	436.48\\
34000	456.19\\
36000	473.38\\
38000	461.34\\
40000	458.01\\
42000	447.22\\
44000	440.88\\
46000	445.94\\
48000	442.94\\
50000	453.49\\
52000	465.77\\
56000	445.99\\
58000	444.56\\
60000	453.25\\
62000	446.6\\
64000	453.74\\
66000	447.31\\
68000	443.33\\
70000	446.88\\
72000	456.73\\
74000	447.56\\
76000	433.07\\
78000	440.65\\
80000	433.35\\
84000	441.65\\
86000	456.26\\
88000	464.97\\
90000	461.46\\
92000	468.37\\
94000	464.01\\
96000	464.45\\
98000	468.37\\
1e+05	449.64\\
};
\addlegendentry{Classe k=10}

            % \input{../../../../.tex/20/KS/SizeEvolutionsfig3.tex}
            \legend{}

            \nextgroupplot[%
                % xmin=0,xmax=1e+05,
                % ymin=100,ymax=1.5151e+05,
            ] % This file was created by matlab2tikz.
%
\definecolor{mycolor1}{rgb}{0.00146,0.00047,0.01387}%
\definecolor{mycolor2}{rgb}{0.01166,0.00942,0.06346}%
\definecolor{mycolor3}{rgb}{0.02943,0.02150,0.11462}%
\definecolor{mycolor4}{rgb}{0.06134,0.03659,0.17764}%
\definecolor{mycolor5}{rgb}{0.08741,0.04456,0.22481}%
\definecolor{mycolor6}{rgb}{0.12291,0.04754,0.28162}%
\definecolor{mycolor7}{rgb}{0.14907,0.04547,0.31709}%
\definecolor{mycolor8}{rgb}{0.18343,0.04033,0.35497}%
\definecolor{mycolor9}{rgb}{0.21795,0.03662,0.38352}%
\definecolor{mycolor10}{rgb}{0.24497,0.03705,0.40001}%
\definecolor{mycolor11}{rgb}{0.27135,0.04092,0.41198}%
\definecolor{mycolor12}{rgb}{0.29076,0.04564,0.41864}%
\definecolor{mycolor13}{rgb}{0.31628,0.05349,0.42512}%
\definecolor{mycolor14}{rgb}{0.33522,0.06006,0.42852}%
\definecolor{mycolor15}{rgb}{0.36028,0.06925,0.43150}%
\definecolor{mycolor16}{rgb}{0.37900,0.07625,0.43272}%
\definecolor{mycolor17}{rgb}{0.39767,0.08326,0.43318}%
\definecolor{mycolor18}{rgb}{0.41633,0.09020,0.43294}%
\definecolor{mycolor19}{rgb}{0.42877,0.09479,0.43241}%
\definecolor{mycolor20}{rgb}{0.44743,0.10160,0.43108}%
\definecolor{mycolor21}{rgb}{0.46610,0.10832,0.42912}%
\definecolor{mycolor22}{rgb}{0.47856,0.11276,0.42747}%
\definecolor{mycolor23}{rgb}{0.49726,0.11938,0.42449}%
\definecolor{mycolor24}{rgb}{0.50973,0.12377,0.42216}%
\definecolor{mycolor25}{rgb}{0.52221,0.12815,0.41955}%
\definecolor{mycolor26}{rgb}{0.53468,0.13253,0.41667}%
\definecolor{mycolor27}{rgb}{0.55339,0.13913,0.41183}%
\definecolor{mycolor28}{rgb}{0.56585,0.14357,0.40826}%
\definecolor{mycolor29}{rgb}{0.57830,0.14804,0.40441}%
\definecolor{mycolor30}{rgb}{0.59073,0.15256,0.40029}%
\definecolor{mycolor31}{rgb}{0.60314,0.15715,0.39589}%
\definecolor{mycolor32}{rgb}{0.61551,0.16182,0.39122}%
\definecolor{mycolor33}{rgb}{0.62169,0.16418,0.38878}%
\definecolor{mycolor34}{rgb}{0.63400,0.16899,0.38370}%
\definecolor{mycolor35}{rgb}{0.64626,0.17391,0.37836}%
\definecolor{mycolor36}{rgb}{0.65846,0.17896,0.37275}%
\definecolor{mycolor37}{rgb}{0.66454,0.18154,0.36985}%
\definecolor{mycolor38}{rgb}{0.67664,0.18681,0.36385}%
\definecolor{mycolor39}{rgb}{0.68865,0.19224,0.35760}%
\definecolor{mycolor40}{rgb}{0.69463,0.19502,0.35439}%
\definecolor{mycolor41}{rgb}{0.70650,0.20073,0.34778}%
\definecolor{mycolor42}{rgb}{0.71240,0.20366,0.34438}%
\definecolor{mycolor43}{rgb}{0.72410,0.20967,0.33742}%
\definecolor{mycolor44}{rgb}{0.72991,0.21276,0.33386}%
\definecolor{mycolor45}{rgb}{0.74142,0.21911,0.32658}%
\definecolor{mycolor46}{rgb}{0.74713,0.22238,0.32286}%
\definecolor{mycolor47}{rgb}{0.75279,0.22571,0.31909}%
\definecolor{mycolor48}{rgb}{0.76401,0.23255,0.31140}%
\definecolor{mycolor49}{rgb}{0.76956,0.23608,0.30749}%
\definecolor{mycolor50}{rgb}{0.77506,0.23967,0.30353}%
\definecolor{mycolor51}{rgb}{0.79129,0.25086,0.29139}%
\definecolor{mycolor52}{rgb}{0.79661,0.25473,0.28726}%
\definecolor{mycolor53}{rgb}{0.80187,0.25867,0.28310}%
\definecolor{mycolor54}{rgb}{0.81224,0.26679,0.27466}%
\definecolor{mycolor55}{rgb}{0.81734,0.27095,0.27039}%
\definecolor{mycolor56}{rgb}{0.82737,0.27952,0.26175}%
\definecolor{mycolor57}{rgb}{0.83230,0.28391,0.25738}%
\definecolor{mycolor58}{rgb}{0.83717,0.28839,0.25299}%
\definecolor{mycolor59}{rgb}{0.84671,0.29756,0.24411}%
\definecolor{mycolor60}{rgb}{0.85138,0.30226,0.23964}%
\definecolor{mycolor61}{rgb}{0.85599,0.30704,0.23513}%
\definecolor{mycolor62}{rgb}{0.86053,0.31189,0.23061}%
\definecolor{mycolor63}{rgb}{0.87374,0.32691,0.21689}%
\definecolor{mycolor64}{rgb}{0.87800,0.33206,0.21227}%
\definecolor{mycolor65}{rgb}{0.89034,0.34796,0.19829}%
\definecolor{mycolor66}{rgb}{0.89819,0.35891,0.18886}%
\definecolor{mycolor67}{rgb}{0.90200,0.36449,0.18412}%
\definecolor{mycolor68}{rgb}{0.90573,0.37014,0.17935}%
\definecolor{mycolor69}{rgb}{0.90939,0.37586,0.17456}%
\definecolor{mycolor70}{rgb}{0.91297,0.38164,0.16975}%
\definecolor{mycolor71}{rgb}{0.91646,0.38748,0.16492}%
\definecolor{mycolor72}{rgb}{0.91988,0.39339,0.16007}%
\definecolor{mycolor73}{rgb}{0.92322,0.39936,0.15519}%
\definecolor{mycolor74}{rgb}{0.92647,0.40539,0.15029}%
\definecolor{mycolor75}{rgb}{0.92964,0.41148,0.14537}%
\definecolor{mycolor76}{rgb}{0.93274,0.41763,0.14042}%
\definecolor{mycolor77}{rgb}{0.93575,0.42383,0.13544}%
\definecolor{mycolor78}{rgb}{0.93868,0.43009,0.13044}%
\definecolor{mycolor79}{rgb}{0.94152,0.43640,0.12541}%
\definecolor{mycolor80}{rgb}{0.94429,0.44277,0.12035}%
\definecolor{mycolor81}{rgb}{0.94956,0.45566,0.11016}%
\definecolor{mycolor82}{rgb}{0.95208,0.46218,0.10503}%
\definecolor{mycolor83}{rgb}{0.96129,0.48872,0.08429}%
\definecolor{mycolor84}{rgb}{0.96540,0.50225,0.07386}%
\definecolor{mycolor85}{rgb}{0.96732,0.50908,0.06866}%
\definecolor{mycolor86}{rgb}{0.97259,0.52980,0.05332}%
\definecolor{mycolor87}{rgb}{0.97568,0.54380,0.04362}%
\definecolor{mycolor88}{rgb}{0.98082,0.57221,0.02851}%
\definecolor{mycolor89}{rgb}{0.98189,0.57939,0.02625}%
\definecolor{mycolor90}{rgb}{0.98378,0.59385,0.02377}%
\definecolor{mycolor91}{rgb}{0.98532,0.60842,0.02420}%
\definecolor{mycolor92}{rgb}{0.98696,0.63048,0.03091}%
\definecolor{mycolor93}{rgb}{0.98793,0.66025,0.05175}%
\definecolor{mycolor94}{rgb}{0.98771,0.68281,0.07249}%
\definecolor{mycolor95}{rgb}{0.97750,0.78226,0.18592}%
\definecolor{mycolor96}{rgb}{0.96624,0.83619,0.26153}%
\definecolor{mycolor97}{rgb}{0.96252,0.85148,0.28555}%
\definecolor{mycolor98}{rgb}{0.98836,0.99836,0.64492}%
%
\addplot [color=mycolor1]
  table[row sep=crcr]{%
0	4395.3\\
2000	71780\\
4000	1.2955e+05\\
6000	1.4472e+05\\
8000	1.4692e+05\\
10000	1.4473e+05\\
12000	1.4525e+05\\
14000	1.4911e+05\\
16000	1.4879e+05\\
18000	1.4999e+05\\
20000	1.4911e+05\\
22000	1.4885e+05\\
24000	1.491e+05\\
26000	1.4804e+05\\
28000	1.4743e+05\\
30000	1.4789e+05\\
32000	1.4647e+05\\
34000	1.4876e+05\\
36000	1.4931e+05\\
38000	1.4578e+05\\
40000	1.4397e+05\\
42000	1.4452e+05\\
44000	1.4572e+05\\
46000	1.4979e+05\\
48000	1.5103e+05\\
50000	1.4822e+05\\
52000	1.4641e+05\\
54000	1.4628e+05\\
56000	1.5073e+05\\
58000	1.4885e+05\\
62000	1.4597e+05\\
64000	1.475e+05\\
66000	1.5001e+05\\
70000	1.4716e+05\\
72000	1.4905e+05\\
76000	1.4789e+05\\
78000	1.489e+05\\
80000	1.4669e+05\\
82000	1.4537e+05\\
84000	1.4495e+05\\
86000	1.4791e+05\\
88000	1.4558e+05\\
90000	1.4657e+05\\
92000	1.4885e+05\\
94000	1.5002e+05\\
96000	1.4924e+05\\
98000	1.4635e+05\\
1e+05	1.4782e+05\\
};
\addlegendentry{Classe k=283}

\addplot [color=mycolor2, forget plot]
  table[row sep=crcr]{%
0	4395.3\\
2000	19786\\
4000	25183\\
6000	26062\\
8000	25381\\
10000	25000\\
12000	25409\\
14000	25570\\
18000	24862\\
20000	25570\\
22000	25484\\
24000	25872\\
26000	25298\\
28000	25774\\
30000	25013\\
32000	25346\\
34000	24977\\
36000	25274\\
38000	25262\\
40000	25562\\
42000	25327\\
44000	26044\\
46000	25865\\
48000	24449\\
50000	25118\\
52000	25672\\
54000	25887\\
56000	25517\\
58000	25308\\
60000	25664\\
62000	25055\\
64000	25604\\
66000	25205\\
68000	25603\\
70000	25374\\
72000	25467\\
74000	25298\\
76000	25057\\
78000	25078\\
80000	25568\\
82000	25991\\
84000	25286\\
86000	26186\\
88000	25743\\
90000	25049\\
92000	25583\\
94000	25351\\
98000	25337\\
1e+05	25306\\
};
\addplot [color=mycolor3, forget plot]
  table[row sep=crcr]{%
0	4395.3\\
2000	27626\\
4000	37322\\
6000	39122\\
8000	38762\\
10000	39655\\
12000	38182\\
14000	38178\\
16000	38399\\
18000	38481\\
20000	37464\\
22000	38723\\
24000	38423\\
26000	38536\\
28000	38233\\
32000	39390\\
34000	39346\\
36000	38521\\
38000	38522\\
40000	38749\\
42000	38838\\
44000	39413\\
46000	38350\\
48000	38123\\
50000	38076\\
52000	38181\\
54000	37629\\
56000	37976\\
58000	39078\\
60000	39173\\
62000	38774\\
64000	38789\\
66000	38473\\
70000	38490\\
72000	37677\\
74000	38079\\
76000	38160\\
78000	38867\\
80000	38919\\
82000	39590\\
84000	38848\\
86000	38493\\
88000	38785\\
90000	38485\\
92000	38034\\
98000	39547\\
1e+05	39266\\
};
\addplot [color=mycolor4, forget plot]
  table[row sep=crcr]{%
0	4395.3\\
2000	12431\\
4000	14245\\
8000	14134\\
10000	14298\\
12000	14008\\
14000	14399\\
16000	13990\\
18000	13885\\
20000	13958\\
22000	13917\\
24000	14109\\
26000	14179\\
28000	14030\\
30000	14204\\
32000	14201\\
34000	13909\\
36000	14176\\
38000	14072\\
40000	14139\\
42000	14070\\
44000	13953\\
46000	14057\\
48000	14023\\
50000	14287\\
52000	14419\\
54000	13878\\
56000	14112\\
58000	13799\\
60000	13968\\
62000	13996\\
64000	14139\\
66000	13854\\
68000	14667\\
70000	14002\\
72000	14176\\
74000	13972\\
76000	14284\\
78000	14122\\
80000	14112\\
82000	14034\\
84000	14033\\
86000	13602\\
88000	14262\\
90000	14049\\
92000	14125\\
94000	14090\\
98000	14174\\
1e+05	13994\\
};
\addplot [color=mycolor5, forget plot]
  table[row sep=crcr]{%
0	4395.3\\
2000	33577\\
4000	55951\\
6000	60552\\
8000	60329\\
10000	61012\\
12000	62738\\
14000	61076\\
16000	62688\\
18000	61298\\
20000	63226\\
22000	61914\\
24000	60945\\
28000	62223\\
30000	61409\\
32000	62206\\
34000	63483\\
36000	61848\\
38000	62532\\
40000	61887\\
42000	63625\\
46000	61207\\
48000	61283\\
50000	61778\\
52000	62090\\
54000	60433\\
56000	60497\\
58000	61646\\
60000	60624\\
62000	61310\\
64000	62197\\
66000	62179\\
70000	61657\\
72000	62160\\
74000	61694\\
76000	61543\\
78000	60954\\
80000	60807\\
82000	62631\\
84000	62997\\
86000	61884\\
88000	62907\\
90000	63144\\
92000	63040\\
94000	61699\\
96000	60568\\
98000	61261\\
1e+05	60210\\
};
\addplot [color=mycolor6, forget plot]
  table[row sep=crcr]{%
0	4395.3\\
2000	30230\\
4000	48800\\
6000	53495\\
10000	54467\\
12000	54099\\
14000	52709\\
18000	53554\\
20000	53507\\
22000	53709\\
26000	53551\\
28000	54628\\
30000	54194\\
32000	53140\\
34000	53967\\
36000	54277\\
38000	53754\\
40000	54060\\
42000	54602\\
44000	54604\\
46000	53279\\
48000	54698\\
50000	54379\\
52000	53784\\
54000	55007\\
56000	53525\\
58000	52766\\
60000	54042\\
62000	53743\\
68000	55104\\
70000	55141\\
72000	54578\\
76000	54456\\
78000	55471\\
80000	55407\\
84000	54362\\
86000	53296\\
88000	54331\\
90000	53114\\
92000	54330\\
94000	54512\\
96000	56051\\
98000	54313\\
1e+05	53630\\
};
\addplot [color=mycolor7, forget plot]
  table[row sep=crcr]{%
0	4395.3\\
2000	19381\\
4000	26791\\
6000	27554\\
8000	26869\\
10000	27654\\
12000	27824\\
14000	27729\\
16000	27515\\
20000	26696\\
22000	27504\\
24000	26870\\
28000	27751\\
30000	27558\\
32000	26976\\
34000	27508\\
36000	28217\\
38000	28066\\
40000	27715\\
42000	27123\\
44000	27135\\
46000	27381\\
48000	27189\\
50000	27165\\
52000	27660\\
54000	27560\\
56000	27795\\
58000	26772\\
60000	27398\\
62000	28157\\
64000	27625\\
66000	26526\\
68000	27204\\
70000	27408\\
72000	27770\\
74000	27932\\
76000	27601\\
80000	28371\\
82000	27545\\
84000	27277\\
86000	27614\\
90000	27639\\
92000	27246\\
94000	27577\\
96000	28231\\
98000	28038\\
1e+05	27599\\
};
\addplot [color=mycolor8, forget plot]
  table[row sep=crcr]{%
0	4395.3\\
2000	15373\\
4000	19465\\
6000	19955\\
8000	19901\\
10000	20155\\
12000	20261\\
14000	19922\\
16000	19956\\
18000	20119\\
20000	20345\\
22000	20252\\
24000	19808\\
26000	19806\\
28000	20444\\
30000	20429\\
32000	20254\\
34000	20401\\
36000	20410\\
38000	19914\\
40000	19667\\
42000	19947\\
44000	20675\\
46000	19913\\
48000	19658\\
50000	20255\\
52000	20311\\
54000	19808\\
56000	20097\\
58000	20130\\
60000	20391\\
62000	20729\\
64000	20371\\
66000	19955\\
68000	20268\\
70000	19924\\
72000	19822\\
74000	19963\\
76000	19868\\
78000	19936\\
80000	20094\\
82000	20459\\
84000	20579\\
86000	19793\\
88000	20028\\
90000	20329\\
92000	19561\\
94000	19701\\
96000	19473\\
98000	19794\\
1e+05	20270\\
};
\addplot [color=mycolor9, forget plot]
  table[row sep=crcr]{%
0	4395.3\\
2000	11641\\
4000	13824\\
6000	13792\\
8000	13650\\
10000	13741\\
12000	14288\\
14000	14374\\
16000	13970\\
18000	13824\\
20000	13379\\
22000	13862\\
24000	13925\\
28000	13851\\
30000	13925\\
32000	14039\\
34000	13600\\
36000	13564\\
38000	14108\\
40000	13870\\
42000	14035\\
44000	14047\\
46000	13720\\
48000	13767\\
50000	14153\\
52000	14146\\
54000	13967\\
56000	14156\\
58000	13970\\
60000	13710\\
62000	14183\\
64000	14099\\
66000	13769\\
68000	14114\\
70000	14284\\
72000	13863\\
74000	13663\\
76000	14302\\
78000	13477\\
80000	13770\\
82000	13905\\
84000	13978\\
86000	14234\\
90000	13834\\
92000	14128\\
94000	13830\\
96000	13439\\
98000	14064\\
1e+05	14197\\
};
\addplot [color=mycolor10, forget plot]
  table[row sep=crcr]{%
0	4395.3\\
2000	15125\\
4000	20453\\
6000	21428\\
8000	21753\\
10000	21858\\
12000	21750\\
14000	21417\\
16000	21747\\
18000	21859\\
20000	21729\\
22000	21430\\
24000	21697\\
26000	21581\\
28000	21706\\
30000	21630\\
32000	21304\\
34000	21304\\
36000	21603\\
38000	21511\\
40000	21728\\
44000	21155\\
46000	21496\\
48000	21356\\
50000	21749\\
52000	21530\\
54000	21397\\
56000	21031\\
58000	21198\\
60000	21214\\
62000	21604\\
64000	21436\\
66000	21771\\
68000	21956\\
70000	21368\\
72000	21645\\
74000	21320\\
76000	21180\\
78000	20976\\
80000	21821\\
82000	22050\\
84000	21422\\
86000	21258\\
88000	21305\\
90000	21854\\
92000	21571\\
94000	21423\\
96000	21045\\
98000	21292\\
1e+05	21145\\
};
\addplot [color=mycolor11, forget plot]
  table[row sep=crcr]{%
0	4395.3\\
2000	8023.8\\
4000	8562.4\\
6000	8327.3\\
8000	8601.1\\
10000	8648.4\\
12000	8227.6\\
14000	8591.1\\
16000	8926.7\\
18000	8620.1\\
20000	8489.7\\
22000	8407.3\\
24000	8576.8\\
26000	8589.5\\
28000	8436.4\\
30000	8492.8\\
32000	8585.2\\
34000	8651.8\\
36000	8526.7\\
38000	8631.2\\
40000	8479.3\\
42000	8844.6\\
44000	8597\\
48000	8672.8\\
50000	8547.7\\
54000	8570.5\\
56000	8677.2\\
58000	8698.2\\
60000	8535.7\\
62000	8687.4\\
64000	8634.3\\
66000	8524.4\\
68000	8575.8\\
70000	8494.8\\
72000	8692.4\\
74000	8758.9\\
76000	8413.8\\
78000	8369.6\\
82000	8738.3\\
84000	8438.1\\
86000	8613.2\\
88000	8404.4\\
90000	8712.6\\
92000	8523.3\\
94000	8696.7\\
96000	8662.4\\
98000	8579.7\\
1e+05	8543.1\\
};
\addplot [color=mycolor12, forget plot]
  table[row sep=crcr]{%
0	4395.3\\
2000	11932\\
4000	14199\\
8000	14676\\
10000	14541\\
12000	14957\\
14000	14631\\
16000	14605\\
18000	14391\\
20000	14270\\
22000	14192\\
24000	14423\\
26000	14306\\
28000	14453\\
30000	14811\\
32000	14541\\
34000	14352\\
36000	14437\\
40000	14275\\
42000	14412\\
44000	14451\\
46000	14612\\
48000	14279\\
50000	14207\\
52000	14394\\
54000	14316\\
56000	14540\\
58000	14910\\
60000	14345\\
62000	14366\\
64000	14139\\
66000	14665\\
68000	14011\\
72000	14922\\
74000	14918\\
76000	14320\\
78000	14524\\
80000	14634\\
82000	14993\\
84000	14419\\
86000	14703\\
88000	14069\\
90000	14256\\
92000	14307\\
94000	14240\\
96000	14704\\
1e+05	14360\\
};
\addplot [color=mycolor13, forget plot]
  table[row sep=crcr]{%
0	4395.3\\
2000	9522.2\\
4000	11129\\
6000	11115\\
8000	10984\\
10000	10999\\
12000	11054\\
14000	10957\\
16000	11044\\
18000	11359\\
20000	11198\\
22000	10921\\
24000	11159\\
26000	11163\\
28000	10689\\
30000	10867\\
32000	11146\\
34000	10857\\
36000	10852\\
38000	11109\\
40000	11011\\
42000	11146\\
44000	11018\\
46000	11125\\
48000	10906\\
50000	10833\\
52000	11358\\
54000	11106\\
56000	11243\\
58000	10992\\
60000	11214\\
62000	11177\\
64000	11221\\
66000	10979\\
68000	11332\\
70000	11066\\
72000	11053\\
74000	11303\\
76000	11000\\
78000	10949\\
80000	10983\\
82000	10560\\
84000	10698\\
86000	10685\\
88000	11006\\
90000	10852\\
92000	11063\\
94000	10977\\
96000	10783\\
98000	10929\\
1e+05	11035\\
};
\addplot [color=mycolor14, forget plot]
  table[row sep=crcr]{%
0	4395.3\\
2000	16277\\
4000	22442\\
6000	23212\\
8000	23465\\
12000	23619\\
14000	24052\\
16000	23860\\
18000	23249\\
20000	23786\\
22000	23359\\
24000	23934\\
26000	24308\\
28000	24240\\
30000	23427\\
32000	23361\\
34000	23586\\
36000	23557\\
38000	23604\\
40000	24614\\
42000	23343\\
44000	23840\\
46000	23966\\
48000	23757\\
50000	23627\\
52000	23828\\
54000	23406\\
56000	23334\\
58000	24162\\
60000	24344\\
62000	23868\\
64000	24025\\
68000	23631\\
70000	23549\\
72000	23690\\
74000	24212\\
76000	24172\\
78000	23233\\
80000	23566\\
82000	23985\\
84000	23947\\
86000	24265\\
88000	23885\\
90000	24483\\
92000	23707\\
94000	23823\\
96000	24110\\
98000	23892\\
1e+05	24139\\
};
\addplot [color=mycolor15, forget plot]
  table[row sep=crcr]{%
0	4395.3\\
2000	12645\\
4000	16237\\
6000	16467\\
8000	16899\\
10000	16909\\
12000	16676\\
14000	16303\\
16000	16239\\
18000	17163\\
20000	16481\\
22000	16647\\
24000	17329\\
26000	16645\\
30000	17104\\
34000	16417\\
36000	16454\\
38000	17136\\
40000	17093\\
42000	16712\\
44000	16606\\
46000	16574\\
48000	16455\\
50000	16746\\
52000	17198\\
54000	16568\\
56000	16833\\
58000	16916\\
60000	17489\\
62000	17241\\
64000	16381\\
66000	17126\\
68000	16464\\
70000	16977\\
72000	16688\\
74000	16768\\
76000	16523\\
78000	16453\\
80000	16647\\
82000	16023\\
84000	16551\\
86000	16251\\
88000	16741\\
90000	16949\\
92000	16996\\
94000	16660\\
96000	16503\\
98000	17024\\
1e+05	16599\\
};
\addplot [color=mycolor16, forget plot]
  table[row sep=crcr]{%
0	4395.3\\
2000	6559.3\\
4000	7043.2\\
6000	7130.2\\
8000	7002.4\\
10000	6951.2\\
12000	7008.9\\
14000	6777.3\\
16000	7150.8\\
18000	7078.3\\
20000	7043.2\\
22000	7038\\
24000	6871.6\\
28000	6777.9\\
30000	7059.4\\
32000	7020.1\\
34000	6842\\
36000	7061.7\\
38000	7066\\
40000	6769.3\\
42000	7104\\
44000	6957.8\\
46000	6953.6\\
48000	7004.5\\
50000	6785\\
52000	7062.7\\
54000	6959.4\\
56000	6979.5\\
60000	7223.5\\
62000	7256.3\\
64000	6987.7\\
66000	6755.8\\
68000	7117.5\\
70000	6954.6\\
72000	6831.9\\
74000	7123.4\\
76000	7005.7\\
78000	7222.5\\
80000	7112.6\\
82000	7175.4\\
84000	7007.1\\
86000	7037.3\\
88000	6811.9\\
90000	7015.1\\
92000	6938.8\\
96000	6959.4\\
98000	7127.7\\
1e+05	6958.1\\
};
\addplot [color=mycolor17, forget plot]
  table[row sep=crcr]{%
0	4395.3\\
2000	12275\\
4000	15655\\
6000	15789\\
8000	15664\\
10000	16138\\
12000	15817\\
14000	15942\\
16000	15725\\
18000	16097\\
20000	15846\\
22000	16064\\
24000	15908\\
26000	15664\\
28000	15899\\
30000	15833\\
32000	16060\\
34000	15599\\
36000	15472\\
38000	16280\\
40000	15887\\
42000	15597\\
44000	15409\\
46000	15896\\
48000	15991\\
50000	15867\\
52000	15849\\
54000	16009\\
56000	15980\\
58000	16314\\
60000	16204\\
62000	15779\\
64000	15272\\
66000	15903\\
68000	15891\\
70000	15688\\
72000	15253\\
74000	15958\\
76000	15882\\
78000	15431\\
80000	15442\\
84000	15674\\
86000	15358\\
88000	15456\\
90000	15770\\
94000	16062\\
98000	15689\\
1e+05	15564\\
};
\addplot [color=mycolor18, forget plot]
  table[row sep=crcr]{%
0	4395.3\\
2000	16360\\
4000	22677\\
6000	23294\\
8000	24102\\
10000	23860\\
12000	23399\\
14000	23918\\
16000	23474\\
18000	24253\\
20000	23829\\
22000	23775\\
24000	23966\\
26000	24608\\
28000	23798\\
30000	24179\\
32000	24189\\
34000	23893\\
36000	23777\\
40000	24283\\
42000	24386\\
44000	24130\\
46000	25107\\
48000	24160\\
50000	24012\\
54000	24339\\
56000	23828\\
58000	24560\\
60000	23858\\
62000	24285\\
64000	24234\\
66000	24489\\
68000	24890\\
70000	23876\\
72000	23589\\
74000	24707\\
76000	24501\\
80000	23702\\
82000	23707\\
84000	24323\\
86000	24691\\
88000	24119\\
90000	23212\\
92000	23829\\
94000	24306\\
96000	24418\\
98000	24772\\
1e+05	23842\\
};
\addplot [color=mycolor19, forget plot]
  table[row sep=crcr]{%
0	4395.3\\
2000	11433\\
4000	14522\\
6000	14727\\
8000	14783\\
10000	15023\\
12000	14576\\
14000	14878\\
16000	15373\\
18000	14911\\
20000	14315\\
22000	14901\\
24000	14902\\
28000	14089\\
30000	14641\\
32000	14799\\
34000	14527\\
36000	14573\\
38000	14706\\
40000	14655\\
42000	14552\\
44000	14930\\
46000	14388\\
48000	14454\\
50000	14466\\
52000	14286\\
54000	14859\\
56000	14505\\
58000	14462\\
60000	14809\\
62000	14847\\
64000	14970\\
66000	13978\\
70000	14433\\
72000	14403\\
74000	14244\\
76000	14306\\
78000	14801\\
80000	14690\\
82000	14687\\
84000	14817\\
86000	14627\\
88000	15106\\
90000	14586\\
92000	14540\\
94000	14630\\
98000	14094\\
1e+05	14531\\
};
\addplot [color=mycolor20, forget plot]
  table[row sep=crcr]{%
0	4395.3\\
2000	13511\\
4000	18193\\
6000	18913\\
8000	19151\\
10000	18469\\
12000	18724\\
14000	18743\\
16000	18640\\
18000	18961\\
20000	19420\\
22000	19408\\
24000	18571\\
26000	18805\\
28000	18915\\
30000	18473\\
34000	19127\\
36000	19215\\
38000	19020\\
40000	18758\\
42000	19008\\
44000	19479\\
48000	18767\\
50000	19034\\
52000	19031\\
54000	18750\\
56000	19394\\
58000	19259\\
60000	18555\\
62000	18765\\
64000	18639\\
66000	18089\\
68000	18704\\
70000	19666\\
72000	18736\\
74000	18394\\
76000	19083\\
78000	19185\\
80000	18813\\
82000	18673\\
84000	19036\\
86000	18908\\
88000	18569\\
90000	18780\\
92000	18259\\
96000	18294\\
98000	18747\\
1e+05	19051\\
};
\addplot [color=mycolor21, forget plot]
  table[row sep=crcr]{%
0	4395.3\\
2000	7984.1\\
4000	9215.9\\
6000	9496\\
8000	9310.4\\
10000	9621.1\\
12000	9652.1\\
14000	9742.9\\
16000	9625.8\\
18000	9776\\
20000	9309.6\\
22000	9334\\
24000	9488\\
26000	9829.6\\
28000	9655.6\\
30000	9397\\
36000	9695.3\\
38000	9598\\
40000	9652.3\\
42000	9835.4\\
46000	9572.6\\
48000	9495.1\\
52000	9793.5\\
54000	9462.5\\
56000	9558.1\\
58000	9573\\
60000	9556\\
62000	9718.4\\
64000	9727.4\\
66000	9504.1\\
68000	9467.6\\
70000	9660.7\\
72000	9613.8\\
74000	9723.5\\
76000	9578.5\\
78000	9604.9\\
80000	9602.9\\
82000	9459.5\\
84000	9585.6\\
86000	9459\\
88000	9718.8\\
92000	9586.1\\
94000	9930.8\\
96000	9677\\
98000	9617.5\\
1e+05	9729.8\\
};
\addplot [color=mycolor22, forget plot]
  table[row sep=crcr]{%
0	4395.3\\
2000	16457\\
4000	24255\\
6000	24538\\
8000	24670\\
10000	24470\\
12000	25009\\
14000	24138\\
16000	24376\\
18000	24450\\
20000	24995\\
22000	25098\\
24000	24012\\
26000	24258\\
28000	25009\\
30000	24903\\
32000	25248\\
34000	24376\\
36000	25002\\
38000	25168\\
40000	24372\\
42000	25142\\
44000	25058\\
46000	25359\\
50000	25436\\
52000	24554\\
54000	24898\\
56000	24723\\
58000	25094\\
60000	24355\\
62000	24145\\
64000	24439\\
66000	23471\\
68000	24239\\
70000	24899\\
72000	24908\\
76000	24554\\
78000	24923\\
80000	25425\\
82000	24854\\
84000	24677\\
86000	24783\\
88000	24776\\
90000	24868\\
92000	24343\\
94000	24015\\
96000	24712\\
98000	24955\\
1e+05	24799\\
};
\addplot [color=mycolor23, forget plot]
  table[row sep=crcr]{%
0	4395.3\\
2000	9917.3\\
4000	12083\\
6000	12074\\
8000	12275\\
10000	12119\\
12000	12082\\
16000	12346\\
20000	12065\\
22000	12004\\
24000	12280\\
26000	12152\\
28000	11983\\
30000	12081\\
32000	12302\\
34000	12304\\
36000	12199\\
38000	12165\\
40000	11961\\
42000	12010\\
44000	12321\\
46000	12318\\
48000	11988\\
50000	12296\\
52000	12147\\
54000	12212\\
58000	12007\\
60000	12145\\
62000	12027\\
64000	12292\\
66000	12089\\
68000	11927\\
70000	12072\\
72000	12153\\
74000	11986\\
76000	12284\\
78000	12194\\
80000	12053\\
82000	12217\\
84000	12264\\
86000	12070\\
90000	12113\\
92000	12214\\
94000	12400\\
96000	12168\\
98000	12121\\
1e+05	12142\\
};
\addplot [color=mycolor24, forget plot]
  table[row sep=crcr]{%
0	4395.3\\
2000	4928.9\\
4000	5014.3\\
6000	4892.8\\
8000	4797.2\\
10000	4850\\
12000	4809.1\\
14000	4865.5\\
16000	4870\\
18000	4807.6\\
20000	4781.9\\
22000	4854.3\\
24000	4789.5\\
26000	4749.3\\
28000	4789\\
30000	4806.7\\
32000	4855.7\\
34000	4860\\
36000	4890.1\\
38000	4778.7\\
40000	4845.2\\
42000	4844.1\\
46000	4790.4\\
50000	4850.3\\
52000	4790.1\\
54000	4912.6\\
56000	4765.6\\
58000	4893.5\\
60000	4810.3\\
62000	4656.9\\
64000	4845.9\\
66000	4946.3\\
68000	4967\\
70000	4894\\
72000	4766.6\\
74000	4616.3\\
76000	4757.9\\
78000	4725\\
80000	4774.7\\
82000	4703.1\\
84000	4841.7\\
86000	4794.8\\
88000	4768.2\\
90000	4658.6\\
92000	4973\\
94000	4785.7\\
96000	4778.8\\
98000	4798.7\\
1e+05	4857\\
};
\addplot [color=mycolor25, forget plot]
  table[row sep=crcr]{%
0	4395.3\\
2000	9971.8\\
4000	12197\\
8000	12338\\
10000	12097\\
12000	12196\\
14000	12159\\
16000	11992\\
18000	11867\\
20000	12177\\
22000	12053\\
24000	12484\\
26000	12023\\
28000	12029\\
30000	12193\\
32000	11992\\
34000	12429\\
36000	12404\\
38000	12620\\
40000	12328\\
42000	12283\\
44000	12299\\
46000	11963\\
48000	12671\\
50000	12563\\
52000	12349\\
56000	12324\\
58000	12173\\
60000	12172\\
62000	12344\\
64000	12200\\
66000	12598\\
68000	11950\\
70000	11894\\
72000	12174\\
74000	12410\\
76000	12579\\
78000	12070\\
80000	12243\\
82000	12187\\
84000	12459\\
86000	12380\\
88000	12167\\
90000	12170\\
92000	12113\\
94000	12166\\
96000	12446\\
98000	12344\\
1e+05	12549\\
};
\addplot [color=mycolor26, forget plot]
  table[row sep=crcr]{%
0	4395.3\\
2000	7490.5\\
4000	8265.7\\
6000	8207.4\\
8000	8400.6\\
10000	8407.7\\
12000	8356.3\\
14000	8204.8\\
16000	8224.1\\
18000	8084.4\\
20000	8130.2\\
22000	8245.6\\
24000	8482.1\\
26000	8254.9\\
28000	8093.3\\
30000	8393.8\\
32000	8329.2\\
34000	8154\\
36000	8089.3\\
38000	8668.4\\
40000	8407.2\\
42000	8312.7\\
44000	8268\\
46000	8305.9\\
48000	7900.2\\
50000	7732.8\\
52000	8106\\
54000	8274\\
56000	8196.7\\
58000	8298.3\\
60000	8192.8\\
62000	8353.2\\
64000	8174.6\\
66000	8150.1\\
68000	8304.6\\
70000	8269.6\\
72000	8138.3\\
76000	8282.8\\
78000	8245.7\\
80000	8157\\
82000	8488.7\\
84000	8220.2\\
86000	8069.1\\
88000	8188.6\\
90000	8265.1\\
92000	8200.1\\
94000	8402\\
96000	8076.2\\
98000	8319.4\\
1e+05	8434.1\\
};
\addplot [color=mycolor27, forget plot]
  table[row sep=crcr]{%
0	4395.3\\
2000	8342.2\\
4000	9667.3\\
6000	9545.2\\
8000	9662\\
10000	9589\\
12000	9650.3\\
14000	9555.1\\
16000	10034\\
18000	9860.1\\
20000	9608.6\\
22000	9707.4\\
24000	9494.1\\
26000	9877.1\\
28000	9467.4\\
30000	9582.9\\
32000	9538.1\\
34000	9668.5\\
36000	9714.9\\
38000	9576.1\\
40000	9733.3\\
42000	9855.3\\
44000	9727.4\\
46000	9919.4\\
48000	9715.7\\
50000	9861.1\\
52000	9867\\
54000	9663.4\\
56000	9670\\
58000	9563\\
60000	9570.5\\
62000	9614.6\\
64000	9784.1\\
66000	9767.8\\
68000	9892.3\\
70000	10142\\
72000	9349\\
74000	9384\\
78000	9763.4\\
80000	9583.9\\
82000	9736\\
84000	9941.4\\
88000	9482.9\\
90000	9851.1\\
92000	9642.3\\
94000	9711.4\\
96000	9983\\
98000	9931.7\\
1e+05	9804.7\\
};
\addplot [color=mycolor28, forget plot]
  table[row sep=crcr]{%
0	4395.3\\
2000	5353.7\\
4000	5404.5\\
6000	5433.2\\
8000	5506.2\\
10000	5468.3\\
12000	5467.3\\
14000	5574.5\\
16000	5495.6\\
18000	5618\\
20000	5603.6\\
22000	5420.5\\
24000	5471\\
26000	5539\\
28000	5532.6\\
30000	5621.5\\
32000	5358.7\\
36000	5587.6\\
38000	5507.4\\
40000	5539.5\\
42000	5607.6\\
44000	5291.5\\
46000	5312.7\\
48000	5431.7\\
50000	5751.5\\
52000	5503.4\\
54000	5483.2\\
56000	5402\\
58000	5533.1\\
60000	5551\\
62000	5447.2\\
64000	5368.9\\
66000	5406.3\\
68000	5364.6\\
70000	5470.9\\
72000	5558.5\\
74000	5345.7\\
76000	5498.7\\
78000	5571.1\\
80000	5585.1\\
82000	5496.6\\
84000	5586.5\\
88000	5560.9\\
90000	5445.4\\
94000	5431.2\\
96000	5551.6\\
98000	5494.2\\
1e+05	5263.5\\
};
\addplot [color=mycolor29, forget plot]
  table[row sep=crcr]{%
0	4395.3\\
2000	9439.2\\
4000	11204\\
6000	11235\\
8000	11672\\
10000	11724\\
14000	11200\\
16000	11398\\
20000	11979\\
22000	11811\\
24000	11588\\
26000	11322\\
28000	11822\\
30000	11535\\
32000	11350\\
34000	11258\\
36000	11528\\
38000	11496\\
46000	11536\\
48000	11123\\
50000	11436\\
52000	11567\\
54000	11639\\
56000	11769\\
58000	11441\\
60000	11620\\
64000	11576\\
66000	11336\\
68000	11265\\
70000	11532\\
72000	11285\\
74000	11104\\
76000	11807\\
78000	11811\\
80000	11640\\
82000	11287\\
84000	11903\\
86000	11641\\
88000	11667\\
90000	12027\\
92000	11560\\
94000	11553\\
96000	11440\\
98000	11386\\
1e+05	11665\\
};
\addplot [color=mycolor30, forget plot]
  table[row sep=crcr]{%
0	4395.3\\
2000	7130.9\\
4000	7902.1\\
6000	7975.2\\
8000	8107.5\\
12000	7911.7\\
14000	7855\\
16000	7841.6\\
18000	7922.7\\
20000	7789.5\\
22000	7803.1\\
24000	7882.1\\
26000	7860.6\\
28000	7805\\
30000	7897\\
32000	8078.2\\
34000	7817.8\\
36000	7914.2\\
38000	7957.3\\
40000	8101.2\\
42000	7839.4\\
46000	7677.7\\
48000	7751.4\\
52000	7772.4\\
54000	7836.2\\
56000	7808.6\\
58000	8040.6\\
62000	7973.6\\
64000	7844\\
68000	7799.1\\
70000	8050.2\\
72000	7876.5\\
74000	7822.2\\
76000	7818.9\\
78000	7881.5\\
80000	7989.4\\
82000	7765.4\\
84000	7977.5\\
86000	7792.2\\
88000	7880.4\\
90000	7681.9\\
92000	7925.8\\
94000	7880.3\\
96000	7954.1\\
98000	7815.2\\
1e+05	7881.3\\
};
\addplot [color=mycolor31, forget plot]
  table[row sep=crcr]{%
0	4395.3\\
2000	6195.7\\
4000	6602.6\\
6000	6576.2\\
8000	6581.3\\
10000	6774.5\\
12000	6466.3\\
14000	6545.2\\
16000	6573\\
18000	6529.4\\
20000	6446.8\\
22000	6417.5\\
24000	6689.3\\
28000	6289.1\\
30000	6405.3\\
32000	6563.1\\
34000	6517.7\\
36000	6711.1\\
38000	6565.4\\
40000	6573.5\\
42000	6492.8\\
44000	6624.5\\
46000	6445.6\\
48000	6466.6\\
50000	6327.7\\
52000	6303.2\\
54000	6516.7\\
56000	6405.6\\
58000	6477.5\\
60000	6483\\
62000	6641.1\\
64000	6353\\
66000	6192.6\\
68000	6378.4\\
70000	6241.1\\
72000	6500.4\\
74000	6488.1\\
76000	6594.8\\
78000	6537.8\\
80000	6372\\
82000	6592.5\\
84000	6613.3\\
86000	6696.8\\
88000	6453.1\\
90000	6337.6\\
92000	6458.7\\
94000	6367.5\\
96000	6662.8\\
98000	6366.4\\
1e+05	6058.5\\
};
\addplot [color=mycolor32, forget plot]
  table[row sep=crcr]{%
0	4395.3\\
2000	8816.1\\
4000	10434\\
6000	10437\\
8000	10348\\
10000	10455\\
12000	10400\\
14000	10850\\
16000	10471\\
18000	10618\\
20000	10265\\
22000	10402\\
24000	10134\\
28000	10519\\
30000	10094\\
32000	10273\\
34000	10262\\
36000	10291\\
38000	10156\\
40000	10270\\
42000	10484\\
44000	10575\\
46000	10266\\
48000	10364\\
50000	10161\\
52000	10431\\
56000	10571\\
58000	10575\\
60000	10449\\
62000	10561\\
64000	9889\\
66000	10242\\
68000	10673\\
70000	10545\\
72000	10728\\
74000	9989.9\\
76000	10715\\
78000	10473\\
80000	10555\\
82000	10677\\
84000	10412\\
86000	10501\\
88000	10299\\
90000	10273\\
92000	10438\\
94000	10567\\
96000	10571\\
98000	10826\\
1e+05	10591\\
};
\addplot [color=mycolor33, forget plot]
  table[row sep=crcr]{%
0	4395.3\\
2000	7327.1\\
4000	8262.3\\
6000	8530.2\\
8000	8162.2\\
10000	8171.8\\
12000	8224.3\\
14000	8027.7\\
16000	8114.6\\
18000	8092.3\\
20000	8204.7\\
22000	8232.9\\
24000	8144.6\\
26000	8152.5\\
28000	8106.3\\
30000	8263.8\\
32000	8083.8\\
34000	8231.8\\
36000	8261.3\\
38000	8238.9\\
40000	8314.2\\
44000	8106.6\\
46000	8140.3\\
48000	8342.2\\
50000	8120.8\\
52000	7940.9\\
54000	8097.3\\
56000	8150.3\\
58000	8294.7\\
60000	8175.2\\
62000	8398.7\\
64000	8226.9\\
66000	8604.7\\
68000	8221.5\\
72000	8014.6\\
74000	8226.1\\
76000	8171.6\\
78000	8046.8\\
80000	7978.6\\
82000	8045.7\\
84000	8146.9\\
86000	8368\\
88000	8276.1\\
90000	8277.2\\
92000	8093.2\\
94000	8335.9\\
96000	8267.9\\
98000	8059\\
1e+05	8222.1\\
};
\addplot [color=mycolor34, forget plot]
  table[row sep=crcr]{%
0	4395.3\\
2000	4350.6\\
6000	3977.7\\
8000	3968.3\\
10000	4006.1\\
12000	4238.5\\
14000	4253.7\\
16000	4101.6\\
18000	4142\\
20000	4119.1\\
22000	4065.2\\
24000	4122.4\\
26000	3992.5\\
28000	4038.2\\
30000	4186.4\\
32000	4117.3\\
34000	3983\\
36000	3975.6\\
38000	3918.1\\
40000	4168.2\\
42000	4079.4\\
44000	4034.8\\
46000	4149.2\\
48000	4051.4\\
50000	4081.7\\
52000	4060.4\\
54000	4102.9\\
58000	4148.5\\
60000	4070.3\\
62000	3950.1\\
64000	4068.7\\
66000	3971\\
70000	4033.7\\
72000	4031.8\\
74000	4095.3\\
76000	4083.3\\
78000	4046.8\\
80000	4080.6\\
82000	3963.9\\
84000	4049.5\\
86000	4024\\
88000	4085.3\\
90000	4061.8\\
92000	4169.4\\
94000	4077\\
96000	4069.9\\
98000	4184.6\\
1e+05	3935.5\\
};
\addplot [color=mycolor35]
  table[row sep=crcr]{%
0	4395.3\\
2000	8447.4\\
4000	9950.6\\
6000	10100\\
10000	10107\\
12000	10244\\
14000	10065\\
16000	10029\\
18000	9931.8\\
20000	10072\\
22000	10023\\
24000	10204\\
26000	10179\\
28000	10074\\
30000	10240\\
32000	10194\\
34000	9999\\
36000	10043\\
38000	9985.3\\
40000	10018\\
42000	10268\\
44000	10076\\
46000	10274\\
48000	10155\\
50000	10228\\
52000	9933.8\\
54000	10291\\
58000	10125\\
60000	10070\\
62000	10142\\
66000	10086\\
68000	10158\\
70000	10177\\
72000	10029\\
78000	10038\\
80000	9880.5\\
82000	9952.8\\
84000	10181\\
86000	10248\\
90000	10118\\
92000	10185\\
94000	10126\\
96000	10184\\
98000	10284\\
1e+05	10274\\
};
\addlegendentry{Classe k=83}

\addplot [color=mycolor36, forget plot]
  table[row sep=crcr]{%
0	4395.3\\
2000	4549\\
4000	4371.8\\
6000	4381.9\\
10000	4266.9\\
12000	4258.2\\
14000	4414.9\\
16000	4236.8\\
18000	4371.2\\
20000	4306.2\\
22000	4256.2\\
24000	4330.8\\
26000	4374\\
28000	4232.3\\
30000	4334.4\\
32000	4346.6\\
34000	4183.5\\
36000	4339\\
38000	4202.6\\
40000	4390.3\\
42000	4250.1\\
44000	4304.4\\
46000	4301.6\\
48000	4329.7\\
50000	4204.7\\
52000	4289.6\\
54000	4340.3\\
58000	4285.5\\
60000	4377.7\\
64000	4324.3\\
66000	4255.9\\
68000	4267.8\\
70000	4260.9\\
72000	4310.5\\
74000	4242\\
76000	4268.3\\
78000	4227.1\\
80000	4373.8\\
82000	4234.4\\
84000	4263.1\\
88000	4275.8\\
90000	4306\\
92000	4321.8\\
96000	4296.6\\
98000	4334.4\\
1e+05	4299.5\\
};
\addplot [color=mycolor37, forget plot]
  table[row sep=crcr]{%
0	4395.3\\
2000	4031.6\\
4000	3776.7\\
6000	3671.6\\
8000	3690.4\\
10000	3655.6\\
12000	3814.1\\
14000	3727.2\\
16000	3674.4\\
22000	3743.9\\
24000	3574.5\\
26000	3614\\
28000	3731.7\\
30000	3656.2\\
32000	3717.2\\
34000	3695.4\\
36000	3724.6\\
40000	3688.3\\
42000	3736.4\\
44000	3652.1\\
46000	3701.2\\
48000	3662.5\\
50000	3676.1\\
52000	3622.5\\
54000	3703.3\\
56000	3691.5\\
58000	3701.7\\
60000	3693.9\\
64000	3755.3\\
66000	3660\\
68000	3680.8\\
70000	3669.2\\
72000	3699.6\\
76000	3718.3\\
78000	3654\\
80000	3708.6\\
82000	3678.4\\
84000	3628\\
86000	3645.1\\
88000	3703.8\\
90000	3677.5\\
92000	3692\\
94000	3686.5\\
96000	3727.7\\
98000	3684.2\\
1e+05	3746.7\\
};
\addplot [color=mycolor38, forget plot]
  table[row sep=crcr]{%
0	4395.3\\
2000	3558.1\\
4000	3104\\
6000	3155.4\\
8000	3001.5\\
10000	3026.5\\
12000	2994.8\\
14000	3047.9\\
16000	2892.9\\
18000	2981.3\\
20000	3180\\
22000	3130.5\\
24000	3033.6\\
26000	3087.6\\
28000	3068.9\\
30000	3076.8\\
32000	3133.1\\
34000	3148.2\\
36000	3123.8\\
38000	3131\\
40000	3022.3\\
42000	3122.6\\
44000	3105.6\\
46000	3001.9\\
48000	3087.1\\
50000	3135.6\\
52000	3101.4\\
54000	3034.8\\
56000	3113.3\\
58000	3039.7\\
62000	3142.4\\
64000	3116.1\\
66000	3101.3\\
68000	3127.4\\
70000	3085\\
72000	3215.9\\
74000	3161.6\\
76000	3017.1\\
78000	3040.2\\
80000	2959.8\\
82000	3035.8\\
84000	3081.3\\
86000	3155.8\\
88000	3050.4\\
90000	3057\\
92000	3023.9\\
94000	3065.7\\
96000	3050.5\\
98000	3089.1\\
1e+05	3109.3\\
};
\addplot [color=mycolor39, forget plot]
  table[row sep=crcr]{%
0	4395.3\\
2000	9407.6\\
4000	11621\\
6000	11535\\
8000	11964\\
10000	11415\\
12000	11372\\
14000	11811\\
16000	11397\\
18000	11746\\
20000	11589\\
22000	11657\\
26000	11691\\
28000	11655\\
30000	11975\\
32000	11622\\
34000	11413\\
36000	11374\\
38000	11711\\
40000	11910\\
42000	11398\\
44000	11170\\
46000	11338\\
48000	11203\\
50000	11826\\
52000	11877\\
56000	11747\\
58000	11658\\
60000	11865\\
62000	11509\\
64000	11893\\
66000	11755\\
68000	11765\\
70000	11948\\
72000	11500\\
74000	11758\\
76000	11945\\
78000	11484\\
80000	11856\\
82000	11946\\
84000	11804\\
86000	11923\\
88000	11466\\
90000	11338\\
92000	11518\\
94000	11616\\
96000	11922\\
98000	11896\\
1e+05	12177\\
};
\addplot [color=mycolor40, forget plot]
  table[row sep=crcr]{%
0	4395.3\\
2000	5423.7\\
4000	5735.3\\
6000	5840.8\\
8000	5756.7\\
10000	5867\\
12000	5841.7\\
14000	5889.1\\
16000	5740.1\\
18000	5740.9\\
20000	5808.2\\
22000	5749.7\\
24000	5809.8\\
26000	5708.7\\
28000	5803.2\\
32000	5650.9\\
34000	5704.6\\
36000	5863.9\\
38000	5767.5\\
40000	5740.2\\
42000	5782\\
44000	5800.2\\
46000	5875.8\\
48000	5864.7\\
50000	5832.6\\
52000	5769.3\\
54000	5769.2\\
56000	5877.9\\
58000	5853.1\\
60000	5753\\
62000	5857.9\\
64000	5784.6\\
66000	5792.8\\
68000	5677.6\\
70000	5794.6\\
72000	5837.8\\
74000	5722.1\\
78000	5794.4\\
80000	5777.5\\
82000	5721.6\\
84000	5816.4\\
86000	5735.1\\
88000	5789\\
90000	5755.4\\
92000	5844.1\\
96000	5784.3\\
98000	5733.8\\
1e+05	5842.3\\
};
\addplot [color=mycolor41, forget plot]
  table[row sep=crcr]{%
0	4395.3\\
2000	5324.8\\
4000	5493.8\\
6000	5467.4\\
8000	5484.7\\
10000	5402.7\\
14000	5446.3\\
16000	5255.9\\
18000	5290.4\\
20000	5255.6\\
22000	5330.8\\
24000	5277.9\\
26000	5365.7\\
28000	5336.8\\
30000	5454\\
32000	5347.1\\
34000	5308.2\\
36000	5365.3\\
38000	5391.8\\
40000	5481.6\\
44000	5282.9\\
46000	5338.1\\
48000	5483\\
52000	5388.8\\
54000	5527\\
56000	5424.1\\
58000	5447.4\\
60000	5492.1\\
62000	5475.1\\
66000	5319\\
70000	5347\\
72000	5464.7\\
74000	5395.8\\
76000	5456.2\\
78000	5380\\
80000	5375.8\\
82000	5263.4\\
84000	5485.9\\
86000	5392.6\\
88000	5363.8\\
90000	5536.2\\
92000	5538.2\\
94000	5467.9\\
96000	5478.6\\
1e+05	5341.4\\
};
\addplot [color=mycolor42, forget plot]
  table[row sep=crcr]{%
0	4395.3\\
2000	12790\\
4000	19010\\
6000	19718\\
8000	20872\\
10000	21236\\
12000	19869\\
14000	19299\\
16000	19646\\
18000	19308\\
20000	19836\\
22000	19905\\
24000	20237\\
26000	20263\\
28000	19795\\
30000	20466\\
32000	20818\\
34000	20011\\
36000	20004\\
38000	20710\\
40000	20741\\
42000	19925\\
44000	20814\\
46000	19921\\
48000	20264\\
50000	20041\\
52000	20162\\
54000	20604\\
56000	20410\\
62000	20552\\
64000	20093\\
66000	20235\\
68000	20264\\
70000	19929\\
74000	20144\\
76000	20387\\
78000	20155\\
80000	20183\\
82000	20029\\
84000	20382\\
86000	19776\\
88000	20323\\
90000	20016\\
92000	19374\\
94000	19662\\
96000	20132\\
98000	20128\\
1e+05	19846\\
};
\addplot [color=mycolor43, forget plot]
  table[row sep=crcr]{%
0	4395.3\\
2000	4880.4\\
4000	4965.6\\
6000	4963.8\\
8000	4924.8\\
10000	4809.3\\
12000	4898.7\\
14000	4900.3\\
16000	4944.2\\
18000	4907\\
20000	5040.1\\
22000	5024.7\\
24000	4920.2\\
26000	4911.5\\
28000	4944.8\\
30000	4936.7\\
32000	4959.1\\
34000	5002.5\\
36000	4958.6\\
42000	4963.1\\
44000	4896.7\\
46000	4999.4\\
50000	4987.7\\
52000	4936.6\\
54000	4979.3\\
56000	4998.2\\
58000	4990.4\\
60000	5039.1\\
62000	4921.3\\
64000	4851.2\\
66000	4998.9\\
68000	5011.8\\
70000	4964.4\\
74000	5051\\
78000	4940.4\\
80000	5054.4\\
82000	4934.4\\
84000	4939.9\\
88000	5126\\
90000	4999.5\\
92000	4964\\
94000	4853.4\\
96000	4912.8\\
98000	5018.3\\
1e+05	4972.4\\
};
\addplot [color=mycolor44, forget plot]
  table[row sep=crcr]{%
0	4395.3\\
2000	3710.9\\
4000	3525.3\\
6000	3452.2\\
10000	3470.9\\
14000	3280\\
16000	3368.4\\
18000	3273.1\\
20000	3408.6\\
22000	3411\\
24000	3390.4\\
26000	3455.8\\
30000	3404\\
32000	3373.6\\
34000	3447.1\\
36000	3456.9\\
38000	3413.2\\
40000	3349.2\\
42000	3368\\
44000	3240.7\\
46000	3385.8\\
50000	3395.6\\
52000	3375.2\\
54000	3503.4\\
56000	3381.5\\
58000	3329.6\\
60000	3443.2\\
62000	3404.8\\
64000	3337.6\\
66000	3358.5\\
68000	3307.9\\
70000	3436.4\\
72000	3318.4\\
74000	3398.9\\
76000	3491.1\\
78000	3503.9\\
80000	3382.5\\
84000	3399.1\\
86000	3378.6\\
88000	3301.2\\
92000	3426.5\\
94000	3386\\
96000	3316\\
98000	3339.6\\
1e+05	3411.8\\
};
\addplot [color=mycolor45, forget plot]
  table[row sep=crcr]{%
0	4395.3\\
2000	6280.3\\
4000	6869.9\\
6000	6903.8\\
8000	6876\\
10000	6926.8\\
12000	6900.3\\
16000	7043.5\\
18000	6924.4\\
20000	6878.5\\
22000	6897.3\\
24000	6794.9\\
26000	6847.5\\
32000	7150.5\\
34000	7104.3\\
36000	6982.6\\
40000	6858.2\\
42000	6858.2\\
44000	7005.1\\
46000	7031.5\\
50000	6878.6\\
52000	6967.7\\
56000	6815.7\\
58000	6981.6\\
60000	7057.8\\
62000	6996.5\\
64000	6870.4\\
66000	6939.8\\
68000	6850.4\\
70000	6865.5\\
72000	6913.4\\
74000	7119.6\\
76000	7030.2\\
78000	7129\\
80000	6894.9\\
82000	7076.9\\
84000	7000.9\\
86000	6828.5\\
88000	6903.8\\
90000	6853.3\\
92000	6853.6\\
94000	6897\\
96000	6810.1\\
98000	6875.5\\
1e+05	6897.3\\
};
\addplot [color=mycolor46, forget plot]
  table[row sep=crcr]{%
0	4395.3\\
2000	4261.8\\
4000	4145.2\\
6000	4083.1\\
8000	4195\\
10000	4088.3\\
14000	4031.9\\
16000	4111.2\\
18000	4133.8\\
20000	4100.2\\
22000	4128.1\\
24000	4016.9\\
26000	4116.8\\
28000	4167.6\\
30000	4152.1\\
32000	4164.2\\
34000	4152.1\\
36000	4106.7\\
38000	4226.2\\
40000	4094.5\\
42000	4193.9\\
44000	4052.9\\
46000	4146.2\\
48000	4152.3\\
50000	4136.1\\
52000	4225.7\\
54000	4116.2\\
56000	4139.3\\
58000	4090.1\\
60000	4199.6\\
62000	4065.7\\
64000	4029.1\\
66000	4065.8\\
68000	4119.6\\
70000	4154.4\\
72000	3976.7\\
74000	4061.6\\
76000	4123.8\\
78000	4224.4\\
80000	4235.9\\
82000	4110.2\\
84000	4154.3\\
86000	4108.8\\
88000	4136.2\\
90000	4151.7\\
92000	4229.4\\
94000	4074.6\\
96000	4094.5\\
98000	4148.5\\
1e+05	4136.9\\
};
\addplot [color=mycolor47, forget plot]
  table[row sep=crcr]{%
0	4395.3\\
2000	5169.5\\
4000	5367.5\\
6000	5298.1\\
8000	5304.8\\
10000	5347.6\\
12000	5284.9\\
14000	5278.7\\
16000	5152.4\\
18000	5174.5\\
20000	5332.2\\
26000	5344.7\\
28000	5243.4\\
30000	5317.9\\
32000	5209.8\\
34000	5219.4\\
36000	5142.6\\
38000	5177.1\\
40000	5469.8\\
42000	5339.2\\
44000	5433\\
46000	5233.8\\
48000	5328.9\\
50000	5312.8\\
52000	5378.1\\
54000	5383.1\\
56000	5249.1\\
58000	5288.5\\
60000	5278.8\\
62000	5230.9\\
64000	5484.2\\
66000	5481.5\\
68000	5342.4\\
70000	5339.1\\
72000	5284.8\\
74000	5301.6\\
78000	5377.2\\
80000	5274.3\\
82000	5326.2\\
84000	5324.6\\
86000	5211\\
92000	5251.5\\
94000	5348.4\\
96000	5258.6\\
98000	5343.1\\
1e+05	5334.4\\
};
\addplot [color=mycolor48, forget plot]
  table[row sep=crcr]{%
0	4395.3\\
2000	3403.3\\
4000	3042.5\\
6000	2961.5\\
8000	2932.9\\
10000	2940.6\\
12000	3020\\
14000	3026.4\\
16000	2965.8\\
18000	2936.6\\
22000	2995.3\\
24000	2939.9\\
26000	2992.3\\
30000	2986.9\\
32000	2922.9\\
34000	2987.3\\
36000	2940.1\\
38000	2986.4\\
40000	2998.2\\
42000	2937.9\\
44000	2960.6\\
46000	2993\\
48000	2948.6\\
50000	2947.9\\
52000	2958.3\\
56000	2992\\
58000	2931.7\\
60000	2924.7\\
62000	3021.4\\
64000	2909.3\\
66000	2952.6\\
70000	2999.7\\
72000	2987.9\\
74000	2947.9\\
76000	2955.4\\
78000	3009.2\\
80000	2931.5\\
84000	2948.5\\
86000	2983.8\\
88000	2993.2\\
90000	2965\\
92000	2975.4\\
96000	2953.5\\
98000	2908.6\\
1e+05	2990.9\\
};
\addplot [color=mycolor49, forget plot]
  table[row sep=crcr]{%
0	4395.3\\
2000	4643.9\\
4000	4438.3\\
6000	4367.1\\
8000	4414.3\\
10000	4276\\
12000	4428.4\\
14000	4344.4\\
16000	4352.1\\
18000	4387.6\\
20000	4515.6\\
22000	4264.9\\
24000	4371\\
26000	4458.5\\
28000	4188.5\\
30000	4398.3\\
32000	4390.7\\
34000	4359.8\\
36000	4359.2\\
38000	4312.9\\
40000	4472.3\\
44000	4365.9\\
48000	4269.1\\
50000	4337.6\\
52000	4257.6\\
54000	4451.3\\
56000	4342.5\\
58000	4310.1\\
60000	4406.1\\
62000	4380.7\\
64000	4397.4\\
66000	4344.4\\
68000	4261.1\\
70000	4410.8\\
72000	4481.7\\
74000	4322.4\\
76000	4359.8\\
78000	4353.8\\
80000	4498\\
82000	4475.1\\
84000	4401.5\\
86000	4285.5\\
88000	4550.1\\
90000	4365.3\\
94000	4333.4\\
96000	4265.1\\
98000	4258.8\\
1e+05	4330.3\\
};
\addplot [color=mycolor50, forget plot]
  table[row sep=crcr]{%
0	4395.3\\
2000	3006.4\\
4000	2578.9\\
6000	2532.9\\
8000	2551.3\\
10000	2485.7\\
12000	2505.1\\
16000	2492.3\\
18000	2527.3\\
20000	2554.2\\
22000	2485.5\\
24000	2525.9\\
26000	2485.7\\
28000	2504.5\\
36000	2511.3\\
38000	2481.6\\
40000	2502.4\\
42000	2541.5\\
44000	2523.5\\
46000	2549.4\\
48000	2532.7\\
50000	2537.6\\
52000	2490.8\\
54000	2565.8\\
56000	2495.6\\
58000	2513.7\\
60000	2486.2\\
62000	2527.4\\
64000	2506.5\\
66000	2537.5\\
68000	2581.2\\
70000	2509.8\\
72000	2493\\
74000	2567.4\\
76000	2505.7\\
78000	2475.6\\
80000	2564.8\\
82000	2532.5\\
86000	2612.4\\
88000	2563.9\\
90000	2619.1\\
92000	2513.9\\
94000	2505.8\\
96000	2580.9\\
98000	2490\\
1e+05	2510.1\\
};
\addplot [color=mycolor51, forget plot]
  table[row sep=crcr]{%
0	4395.3\\
2000	5663.1\\
4000	6125.1\\
6000	5921.1\\
8000	5903.3\\
12000	5842.8\\
14000	5992.3\\
16000	5988.2\\
18000	5959.8\\
20000	6063.1\\
22000	6032\\
24000	6031\\
26000	5916.9\\
28000	5876.8\\
30000	6007.8\\
32000	5845.3\\
34000	5873.7\\
36000	5961.8\\
38000	5886.2\\
40000	6018.8\\
42000	6036.6\\
44000	5770.4\\
46000	6010.3\\
48000	5768.6\\
50000	5947.4\\
52000	5902.9\\
54000	5924.6\\
56000	6082.5\\
58000	6226\\
60000	6143\\
62000	5993.2\\
64000	6154.6\\
66000	6132.1\\
68000	5847.1\\
70000	6030\\
72000	6000.3\\
74000	5749.8\\
76000	5852.1\\
78000	6023.5\\
80000	5833.8\\
82000	5865.2\\
84000	6045.1\\
86000	6063.1\\
88000	6215.1\\
90000	5998.6\\
92000	5980.6\\
94000	5819\\
96000	5846.8\\
98000	6257\\
1e+05	5999.5\\
};
\addplot [color=mycolor52, forget plot]
  table[row sep=crcr]{%
0	4395.3\\
2000	4007.9\\
4000	3620.3\\
6000	3476.4\\
8000	3424.6\\
10000	3487.8\\
12000	3647.7\\
14000	3584.8\\
16000	3423.9\\
18000	3464.9\\
20000	3430.5\\
22000	3508.9\\
24000	3476.3\\
26000	3573.6\\
30000	3450.8\\
32000	3575.3\\
34000	3558.9\\
36000	3483.8\\
38000	3540.1\\
40000	3554.4\\
42000	3546.2\\
44000	3624.6\\
48000	3426.1\\
50000	3485.9\\
52000	3534.8\\
54000	3428.4\\
56000	3403.5\\
58000	3627.5\\
60000	3723.1\\
64000	3549.2\\
66000	3665.8\\
70000	3489.6\\
72000	3534.5\\
74000	3522.3\\
76000	3474.4\\
78000	3480.9\\
80000	3624.2\\
82000	3611.7\\
84000	3497.2\\
86000	3579.3\\
88000	3499.9\\
90000	3613.4\\
92000	3586.9\\
94000	3504.8\\
96000	3505.5\\
98000	3583.6\\
1e+05	3467.5\\
};
\addplot [color=mycolor53, forget plot]
  table[row sep=crcr]{%
0	4395.3\\
2000	5300.8\\
4000	5706\\
6000	5608.1\\
8000	5703.5\\
12000	5741.5\\
14000	5729.7\\
16000	5527.8\\
18000	5702.2\\
20000	5704.8\\
22000	5838.1\\
24000	5769.7\\
26000	5588.7\\
28000	5780\\
30000	5758.8\\
32000	5808.7\\
34000	5639.4\\
36000	5714.7\\
38000	5921.5\\
40000	5926.1\\
42000	5725\\
44000	5887.5\\
46000	5858.4\\
48000	5812.8\\
50000	5785.4\\
52000	5849.8\\
54000	5684.2\\
56000	5599.6\\
58000	5685.6\\
60000	5632.4\\
62000	5560.8\\
64000	5766.7\\
66000	5823.6\\
68000	5967.3\\
70000	5859.2\\
72000	5936.7\\
76000	5870.8\\
78000	5804\\
80000	5771.8\\
84000	5645.5\\
86000	5748.3\\
88000	5785.8\\
90000	5706\\
92000	5786.8\\
94000	5523.1\\
96000	5590.6\\
98000	5562.7\\
1e+05	5804\\
};
\addplot [color=mycolor54, forget plot]
  table[row sep=crcr]{%
0	4395.3\\
2000	4610.9\\
4000	4674.2\\
6000	4528.2\\
8000	4598.6\\
10000	4495.3\\
12000	4681.9\\
14000	4669.7\\
16000	4687.5\\
20000	4457.8\\
22000	4611.8\\
24000	4551\\
26000	4668.3\\
28000	4677.4\\
30000	4554\\
32000	4675.8\\
34000	4631.5\\
40000	4589.5\\
44000	4613.8\\
50000	4598.8\\
52000	4634.9\\
54000	4696.5\\
56000	4628.8\\
58000	4529.7\\
60000	4587.9\\
62000	4632.7\\
64000	4717.1\\
66000	4734\\
68000	4512.9\\
70000	4643.9\\
72000	4647.7\\
74000	4682.1\\
76000	4672\\
78000	4712.2\\
80000	4677.2\\
82000	4706.1\\
86000	4571.6\\
88000	4564.7\\
90000	4626.2\\
92000	4629.2\\
94000	4604.5\\
96000	4652.9\\
98000	4579.8\\
1e+05	4488.8\\
};
\addplot [color=mycolor55, forget plot]
  table[row sep=crcr]{%
0	4395.3\\
2000	5376.5\\
4000	5744.1\\
6000	5800.2\\
10000	5650.8\\
12000	5826.3\\
16000	5713.6\\
18000	5761.6\\
20000	5852.5\\
22000	5794.8\\
24000	5777.7\\
26000	5843.3\\
28000	5797.2\\
32000	5853.4\\
34000	5814.7\\
36000	5747.8\\
40000	5731.8\\
42000	5847.5\\
44000	5792.2\\
46000	5703.4\\
48000	5885.1\\
50000	5763.2\\
52000	5856.2\\
54000	5789.9\\
56000	5842\\
58000	5727.4\\
60000	5748.9\\
62000	5822.2\\
64000	5735.3\\
66000	5694.7\\
68000	5628.9\\
70000	5679.8\\
72000	5840.6\\
74000	5842.2\\
76000	5710.3\\
78000	5771.9\\
80000	5636.3\\
82000	5771\\
84000	5784.1\\
86000	5677.3\\
88000	5903.9\\
90000	5838.1\\
92000	5828\\
94000	5920.4\\
96000	5825.5\\
98000	5802.7\\
1e+05	5667.1\\
};
\addplot [color=mycolor56, forget plot]
  table[row sep=crcr]{%
0	4395.3\\
2000	4218.8\\
6000	4032.4\\
8000	4029.8\\
10000	4066\\
12000	4049.4\\
14000	4110.4\\
16000	3950.5\\
18000	3991.7\\
20000	3952.8\\
22000	4103.7\\
24000	4112.7\\
26000	4017.9\\
28000	4008.2\\
30000	3954.1\\
32000	4006.3\\
34000	3976.8\\
38000	4073.3\\
40000	4104.8\\
42000	4023.6\\
44000	4115.7\\
46000	3997\\
48000	4001.8\\
50000	4028.5\\
52000	3992.4\\
54000	4015.5\\
56000	4001.8\\
58000	4020\\
62000	4020.9\\
64000	4071.4\\
66000	4040\\
68000	4036.3\\
70000	4015.4\\
72000	4050\\
74000	4018.3\\
76000	4083.1\\
78000	4066.7\\
80000	3963.3\\
82000	3984\\
84000	4032.2\\
88000	4026.2\\
92000	4091.3\\
94000	4065.3\\
96000	4072\\
98000	4053.6\\
1e+05	4048.9\\
};
\addplot [color=mycolor57, forget plot]
  table[row sep=crcr]{%
0	4395.3\\
2000	5675.9\\
4000	6549.7\\
6000	6651.6\\
8000	6651.3\\
10000	6672.8\\
12000	6621.4\\
14000	6667.6\\
16000	6762.4\\
18000	6735.5\\
20000	6734.4\\
22000	6640.4\\
24000	6607.2\\
26000	6793.8\\
28000	6850\\
30000	6734.2\\
32000	6728.1\\
34000	6748.2\\
36000	6592.6\\
38000	6642.3\\
42000	6798.6\\
44000	6701\\
46000	6567.3\\
48000	6611.5\\
50000	6712.5\\
52000	6733.2\\
54000	6658.8\\
56000	6618.5\\
58000	6541.5\\
60000	6659.9\\
64000	6763.5\\
66000	6641.4\\
70000	6678.3\\
76000	6695.7\\
78000	6721.3\\
80000	6847.2\\
82000	6780.8\\
84000	6685.5\\
86000	6662.2\\
88000	6719.8\\
92000	6661\\
94000	6727.7\\
96000	6628.3\\
98000	6698.3\\
1e+05	6693.9\\
};
\addplot [color=mycolor58, forget plot]
  table[row sep=crcr]{%
0	4395.3\\
2000	4013.4\\
4000	3681.3\\
6000	3606.8\\
8000	3669\\
10000	3572.6\\
12000	3631.1\\
14000	3562.6\\
16000	3620.4\\
18000	3559.2\\
20000	3639\\
22000	3570.1\\
24000	3622.6\\
26000	3543.4\\
28000	3609.3\\
30000	3741.1\\
32000	3609.3\\
36000	3614.4\\
38000	3630.3\\
44000	3534.2\\
46000	3521.7\\
48000	3634.7\\
50000	3651.6\\
52000	3550.3\\
54000	3624.8\\
56000	3625.6\\
58000	3714.4\\
60000	3644.8\\
62000	3487.4\\
64000	3661.8\\
68000	3629.3\\
70000	3688.4\\
72000	3665.9\\
74000	3672.7\\
76000	3645.5\\
78000	3602.3\\
80000	3616.9\\
82000	3657\\
84000	3647.8\\
86000	3582\\
88000	3549.8\\
90000	3706.2\\
92000	3587.2\\
96000	3555.9\\
98000	3630.7\\
1e+05	3632.1\\
};
\addplot [color=mycolor59, forget plot]
  table[row sep=crcr]{%
0	4395.3\\
2000	4855.3\\
4000	5040.8\\
6000	5039.8\\
8000	5097.9\\
10000	5117.8\\
12000	5115.6\\
14000	5074.5\\
16000	5015.8\\
18000	5046.3\\
20000	5212.6\\
22000	5055.3\\
24000	5047.8\\
26000	5070.4\\
28000	5008.6\\
30000	5015.1\\
32000	5125.8\\
34000	5090.2\\
36000	5138.2\\
38000	5106.8\\
40000	4941.8\\
42000	4944.1\\
44000	5132.6\\
46000	5087.2\\
48000	5006.8\\
52000	5055.1\\
54000	5120.4\\
56000	5041.5\\
58000	5081.8\\
60000	5021.4\\
64000	5079.1\\
66000	5111.8\\
68000	5065\\
70000	5074.8\\
72000	5158.1\\
74000	5039.3\\
76000	4992.6\\
78000	5158.5\\
80000	5220.1\\
82000	5128.5\\
84000	4981.4\\
86000	5062.7\\
88000	4982.5\\
90000	4926.1\\
92000	5018.2\\
94000	5041.1\\
98000	5170.4\\
1e+05	5099.8\\
};
\addplot [color=mycolor60, forget plot]
  table[row sep=crcr]{%
0	4395.3\\
2000	3855.1\\
4000	3611.2\\
6000	3570.1\\
10000	3665.5\\
12000	3607.8\\
14000	3649.3\\
16000	3600.7\\
18000	3522.6\\
20000	3652.7\\
22000	3615\\
24000	3706.1\\
26000	3609\\
28000	3573.9\\
30000	3584.6\\
32000	3579.7\\
34000	3596.9\\
36000	3574\\
38000	3630.1\\
42000	3523.1\\
44000	3565\\
46000	3496.5\\
48000	3583.4\\
50000	3572.4\\
52000	3530.8\\
54000	3615.5\\
56000	3556.9\\
58000	3541.2\\
60000	3569\\
62000	3628.4\\
64000	3580.5\\
66000	3653.7\\
68000	3604.4\\
70000	3652.1\\
72000	3590\\
74000	3574.3\\
76000	3519.3\\
78000	3582.2\\
80000	3604.9\\
82000	3614.1\\
84000	3611.5\\
86000	3570.4\\
88000	3616.6\\
90000	3601.4\\
92000	3523.9\\
94000	3582.5\\
96000	3504.2\\
98000	3560.7\\
1e+05	3554.7\\
};
\addplot [color=mycolor61, forget plot]
  table[row sep=crcr]{%
0	4395.3\\
2000	3578.7\\
4000	3204.1\\
6000	3175.3\\
8000	3161\\
10000	3181.8\\
12000	3189.4\\
14000	3153\\
16000	3193.3\\
18000	3139.5\\
20000	3156.4\\
22000	3188.1\\
24000	3146.7\\
26000	3144.2\\
28000	3191.9\\
30000	3192.9\\
32000	3153.6\\
34000	3167.5\\
36000	3154.3\\
38000	3187.1\\
42000	3183.5\\
44000	3156.8\\
46000	3184.3\\
48000	3170.9\\
50000	3130.1\\
52000	3187.2\\
54000	3184.1\\
56000	3205.2\\
58000	3194.5\\
60000	3155.1\\
62000	3146.5\\
64000	3181.3\\
66000	3195.1\\
68000	3169\\
70000	3185.5\\
72000	3151.5\\
74000	3144.8\\
76000	3172.1\\
78000	3155.5\\
80000	3178.7\\
82000	3124.5\\
84000	3134.6\\
86000	3191.7\\
88000	3171\\
90000	3184.4\\
92000	3152.7\\
98000	3140.7\\
1e+05	3178\\
};
\addplot [color=mycolor62, forget plot]
  table[row sep=crcr]{%
0	4395.3\\
2000	4797.7\\
4000	5010.9\\
6000	5005.3\\
8000	5094.4\\
12000	5008.5\\
14000	5099.1\\
16000	4979.8\\
20000	5036.3\\
22000	5025.3\\
24000	5030.7\\
26000	4954.7\\
28000	5105.5\\
30000	5048.8\\
32000	5079.7\\
34000	4985.8\\
36000	5013.3\\
38000	4976\\
40000	5012.4\\
42000	5084.2\\
44000	4998\\
46000	5084.9\\
48000	5013.9\\
50000	5055.3\\
52000	5035.8\\
54000	4957.4\\
56000	5025.7\\
58000	4996\\
60000	4991.9\\
62000	5099.8\\
64000	5019.7\\
66000	5000.8\\
68000	5007.9\\
70000	5029.3\\
72000	5070.7\\
74000	4961.4\\
76000	4986.6\\
78000	5106.4\\
80000	5056.6\\
82000	4971.7\\
84000	4983.7\\
86000	5047.2\\
90000	5014.9\\
92000	5015.6\\
94000	4981.4\\
96000	5000.6\\
98000	4938.2\\
1e+05	4959\\
};
\addplot [color=mycolor63, forget plot]
  table[row sep=crcr]{%
0	4395.3\\
2000	3073.2\\
4000	2548\\
6000	2482.3\\
8000	2462.9\\
10000	2469.1\\
12000	2466.7\\
14000	2489.5\\
16000	2442.5\\
18000	2463.9\\
20000	2508.5\\
22000	2408.7\\
24000	2463.3\\
26000	2465.4\\
28000	2493.7\\
30000	2440.8\\
32000	2422\\
34000	2434.3\\
36000	2470.6\\
38000	2441.1\\
40000	2425.3\\
42000	2436.8\\
44000	2471.6\\
46000	2451\\
50000	2467.9\\
52000	2449.1\\
54000	2445.4\\
56000	2413.5\\
58000	2459.8\\
60000	2476.5\\
62000	2459.7\\
64000	2472.7\\
66000	2515.9\\
68000	2460.4\\
70000	2473.7\\
72000	2449.3\\
74000	2449.6\\
76000	2418.8\\
78000	2438\\
82000	2398.5\\
84000	2458\\
86000	2483.8\\
88000	2419.5\\
90000	2424.7\\
92000	2398\\
94000	2415\\
96000	2404.2\\
1e+05	2456.8\\
};
\addplot [color=mycolor64, forget plot]
  table[row sep=crcr]{%
0	4395.3\\
2000	3491.7\\
4000	3163.4\\
6000	3078.8\\
8000	3054.2\\
10000	3089.1\\
16000	3044.9\\
18000	3070.7\\
20000	3065\\
22000	3050\\
24000	3119.3\\
26000	3113.8\\
28000	3071.1\\
30000	3050.7\\
32000	3058.3\\
34000	3043.6\\
36000	3050\\
38000	3088.1\\
40000	3074.7\\
42000	3073.4\\
44000	3057.4\\
46000	3092.7\\
48000	3094.7\\
50000	3113.7\\
52000	3048.1\\
54000	3101.1\\
56000	3076\\
60000	3105.9\\
62000	3057.4\\
64000	3084.4\\
70000	3087.8\\
72000	3126\\
74000	3116.7\\
78000	3073.9\\
80000	3118.7\\
84000	3041\\
86000	3074\\
88000	3077\\
90000	3097.5\\
92000	3047.3\\
94000	3062.4\\
96000	3102.3\\
98000	3096.4\\
1e+05	3135.9\\
};
\addplot [color=mycolor65, forget plot]
  table[row sep=crcr]{%
0	4395.3\\
2000	3134.5\\
4000	2673.4\\
6000	2614.2\\
8000	2604.8\\
10000	2582.5\\
12000	2595.5\\
14000	2623.2\\
16000	2628.4\\
18000	2572.2\\
20000	2606.2\\
24000	2583.1\\
26000	2608.1\\
28000	2564.6\\
30000	2534.6\\
32000	2586.9\\
36000	2561.3\\
40000	2611.2\\
42000	2575.2\\
44000	2613\\
46000	2609.7\\
48000	2578.6\\
50000	2627.4\\
52000	2627.4\\
54000	2564.9\\
56000	2580.8\\
58000	2607.7\\
60000	2659.8\\
62000	2597.1\\
64000	2580.2\\
68000	2575.8\\
70000	2595.5\\
72000	2603.6\\
76000	2574.8\\
78000	2500.3\\
80000	2584.8\\
82000	2574.5\\
84000	2601.6\\
86000	2611.4\\
88000	2579.8\\
90000	2605.7\\
92000	2559.7\\
94000	2598.2\\
96000	2556.1\\
98000	2565.2\\
1e+05	2561.9\\
};
\addplot [color=mycolor66, forget plot]
  table[row sep=crcr]{%
0	4395.3\\
2000	3745\\
4000	3497.8\\
6000	3481.6\\
8000	3502\\
10000	3509.6\\
12000	3502.1\\
14000	3455.3\\
16000	3491.8\\
20000	3515\\
22000	3496.3\\
24000	3461.3\\
26000	3482.9\\
28000	3462.3\\
30000	3488.7\\
32000	3451.8\\
34000	3456.2\\
36000	3445.4\\
38000	3446.5\\
40000	3501.8\\
42000	3526.1\\
44000	3452.6\\
46000	3448.1\\
48000	3463\\
50000	3437.4\\
52000	3471.9\\
54000	3477.8\\
60000	3414.8\\
62000	3466.3\\
64000	3450.3\\
66000	3376.3\\
68000	3386.2\\
70000	3422.5\\
72000	3477.4\\
74000	3473.8\\
76000	3480.8\\
78000	3462\\
80000	3466.9\\
82000	3461.9\\
84000	3390.1\\
86000	3475.6\\
88000	3416.7\\
90000	3396.9\\
92000	3433.4\\
94000	3431.3\\
96000	3470.5\\
98000	3464.5\\
1e+05	3527.3\\
};
\addplot [color=mycolor67, forget plot]
  table[row sep=crcr]{%
0	4395.3\\
2000	3705.5\\
4000	3346.4\\
6000	3187.5\\
8000	3258.3\\
12000	3249.9\\
14000	3261\\
16000	3211.2\\
18000	3226.7\\
20000	3201.6\\
22000	3251.6\\
24000	3262.2\\
28000	3138.6\\
30000	3165.1\\
32000	3230.5\\
34000	3220.8\\
36000	3181\\
38000	3252.7\\
40000	3270.5\\
42000	3235.7\\
44000	3181.9\\
46000	3235.9\\
48000	3273.1\\
50000	3298.2\\
52000	3237.9\\
54000	3216.1\\
56000	3207.7\\
58000	3242.3\\
60000	3248.2\\
62000	3217.8\\
64000	3201.1\\
66000	3227\\
68000	3180.6\\
70000	3245.7\\
74000	3193.3\\
76000	3213.9\\
78000	3250.5\\
80000	3205.4\\
86000	3266.6\\
88000	3204.3\\
92000	3227.4\\
94000	3209\\
96000	3173.4\\
98000	3235.6\\
1e+05	3247.2\\
};
\addplot [color=mycolor67]
  table[row sep=crcr]{%
0	4395.3\\
2000	3215.4\\
4000	2700.4\\
6000	2652.6\\
8000	2648.3\\
12000	2670\\
16000	2702\\
18000	2687.9\\
20000	2663.3\\
22000	2672.8\\
24000	2707.1\\
26000	2668.1\\
28000	2693.3\\
30000	2674.4\\
32000	2681.7\\
36000	2713.1\\
38000	2673.2\\
40000	2677.4\\
42000	2668.3\\
44000	2694.1\\
46000	2649.4\\
48000	2669.1\\
50000	2606.5\\
54000	2712.6\\
56000	2672.5\\
58000	2656.1\\
60000	2669.5\\
62000	2668.2\\
64000	2650.6\\
66000	2675.8\\
68000	2668.3\\
70000	2674.7\\
72000	2668.8\\
74000	2699\\
76000	2713.2\\
78000	2685.4\\
80000	2647.4\\
82000	2633.8\\
84000	2631.7\\
86000	2649.4\\
88000	2658.3\\
90000	2689.6\\
92000	2687.1\\
94000	2628.9\\
96000	2645.6\\
1e+05	2711.3\\
};
\addlegendentry{Classe k=43}

\addplot [color=mycolor68, forget plot]
  table[row sep=crcr]{%
0	4395.3\\
2000	3159.8\\
4000	2599.7\\
6000	2513.1\\
8000	2501.6\\
12000	2520.2\\
14000	2503.4\\
16000	2520.8\\
20000	2497.2\\
22000	2526.4\\
28000	2532.4\\
30000	2515.8\\
32000	2512.7\\
34000	2494.2\\
36000	2499.7\\
38000	2486.1\\
40000	2513\\
42000	2516.7\\
46000	2479.8\\
48000	2517.4\\
50000	2496.2\\
52000	2503.2\\
54000	2523.7\\
56000	2475.9\\
58000	2503.5\\
60000	2480.7\\
62000	2471.8\\
64000	2476\\
66000	2494.9\\
68000	2494.2\\
70000	2516.6\\
72000	2488.4\\
74000	2507.7\\
76000	2479.8\\
78000	2511.2\\
80000	2477.7\\
84000	2520.3\\
88000	2515.4\\
90000	2552.1\\
94000	2486.1\\
98000	2505.9\\
1e+05	2487.9\\
};
\addplot [color=mycolor69, forget plot]
  table[row sep=crcr]{%
0	4395.3\\
2000	3492.3\\
4000	3024.5\\
6000	2932.4\\
8000	2928.4\\
10000	2949.4\\
12000	2925.2\\
14000	2932\\
16000	2960.2\\
18000	2926.1\\
20000	2927.2\\
22000	2953.7\\
24000	2938.2\\
28000	2933\\
30000	2970.1\\
32000	2959.9\\
34000	2973.8\\
36000	2936.9\\
38000	2966.8\\
40000	2963.8\\
42000	2937.2\\
44000	2945.4\\
46000	2944.1\\
48000	2961\\
50000	2936.9\\
52000	2928.1\\
54000	2942\\
56000	2977.7\\
58000	2944.7\\
60000	2945.3\\
62000	2928.1\\
66000	2948\\
70000	2934\\
72000	2964\\
74000	2980.9\\
76000	2978\\
78000	2927.5\\
80000	2929.9\\
82000	2969.6\\
84000	2997.5\\
86000	3003.1\\
88000	2957.1\\
90000	2981.8\\
92000	2981\\
94000	2945.6\\
96000	2976.8\\
98000	2920.5\\
1e+05	2959.6\\
};
\addplot [color=mycolor70, forget plot]
  table[row sep=crcr]{%
0	4395.3\\
2000	2906.4\\
4000	2211.9\\
6000	2149\\
8000	2166.1\\
12000	2183.7\\
14000	2125.3\\
16000	2152.3\\
18000	2146.5\\
20000	2165.1\\
22000	2152.7\\
24000	2184.4\\
26000	2144.2\\
28000	2147.3\\
30000	2156.8\\
32000	2140.9\\
34000	2157.4\\
36000	2133.8\\
38000	2175.1\\
40000	2126.4\\
42000	2171.9\\
44000	2151.6\\
46000	2143.8\\
48000	2164.4\\
50000	2143.5\\
52000	2137.8\\
58000	2180.6\\
60000	2157.1\\
62000	2141.6\\
68000	2136.7\\
70000	2167.1\\
72000	2121.7\\
74000	2118\\
76000	2101.6\\
78000	2176.2\\
80000	2169\\
82000	2192.3\\
84000	2170.7\\
86000	2172.5\\
88000	2192.5\\
90000	2168.3\\
92000	2185.2\\
94000	2193.8\\
96000	2160.6\\
98000	2209\\
1e+05	2211\\
};
\addplot [color=mycolor71, forget plot]
  table[row sep=crcr]{%
0	4395.3\\
2000	3151.1\\
4000	2629.2\\
6000	2539\\
8000	2518.5\\
10000	2514.8\\
12000	2574.7\\
14000	2553.7\\
18000	2534.8\\
20000	2540.2\\
22000	2516.5\\
26000	2540.4\\
28000	2505.7\\
30000	2505.1\\
32000	2524.1\\
36000	2523.7\\
38000	2507\\
40000	2534\\
42000	2521.4\\
44000	2522.2\\
48000	2556.5\\
50000	2537.3\\
52000	2544.2\\
54000	2520.4\\
56000	2538.6\\
58000	2501\\
60000	2537.3\\
62000	2540.6\\
64000	2533.2\\
66000	2553\\
68000	2528.3\\
70000	2496.1\\
72000	2536.4\\
74000	2541.3\\
78000	2520.7\\
80000	2530.8\\
82000	2526.7\\
88000	2560.1\\
90000	2550.6\\
92000	2530.7\\
94000	2538.2\\
96000	2514.9\\
1e+05	2520.5\\
};
\addplot [color=mycolor72, forget plot]
  table[row sep=crcr]{%
0	4395.3\\
2000	2998.4\\
4000	2315.4\\
6000	2272.3\\
8000	2269.4\\
10000	2237.7\\
12000	2233.9\\
14000	2297.4\\
18000	2301.7\\
20000	2215.6\\
22000	2246.8\\
24000	2246.2\\
26000	2238.4\\
28000	2275.4\\
30000	2249.5\\
32000	2271.1\\
36000	2242.1\\
38000	2280.2\\
40000	2311.9\\
42000	2233.1\\
44000	2226.6\\
46000	2208.8\\
48000	2181.4\\
52000	2232.2\\
54000	2226.5\\
56000	2207.8\\
58000	2229.7\\
60000	2232.2\\
62000	2243\\
64000	2204.2\\
66000	2236.4\\
68000	2230.5\\
70000	2245.9\\
72000	2195\\
74000	2254.6\\
76000	2219.9\\
78000	2254\\
80000	2208.4\\
82000	2250.2\\
86000	2217.4\\
88000	2261.6\\
90000	2273.7\\
92000	2230.4\\
94000	2257.4\\
98000	2233.2\\
1e+05	2258.3\\
};
\addplot [color=mycolor73, forget plot]
  table[row sep=crcr]{%
0	4395.3\\
2000	3211.4\\
4000	2726.8\\
6000	2646.7\\
8000	2635.2\\
10000	2658.4\\
14000	2654.5\\
16000	2672\\
18000	2672.3\\
20000	2650.2\\
22000	2644.3\\
24000	2658.5\\
26000	2640.2\\
30000	2661\\
34000	2655.4\\
36000	2677.8\\
38000	2647.1\\
40000	2667.9\\
42000	2640.9\\
44000	2624.9\\
46000	2629.3\\
48000	2649.6\\
50000	2678.8\\
52000	2685.8\\
54000	2682\\
56000	2629.6\\
58000	2650.5\\
60000	2679.2\\
62000	2683.1\\
66000	2622\\
68000	2655.8\\
74000	2652.7\\
76000	2640.9\\
78000	2621.4\\
80000	2629.7\\
82000	2654.9\\
84000	2664.2\\
88000	2671.2\\
90000	2634.7\\
92000	2618.8\\
94000	2637\\
96000	2633.8\\
98000	2670.6\\
1e+05	2672.8\\
};
\addplot [color=mycolor74, forget plot]
  table[row sep=crcr]{%
0	4395.3\\
2000	2479.4\\
4000	1681.6\\
6000	1612.7\\
8000	1595.3\\
12000	1596.7\\
14000	1585.4\\
18000	1597.3\\
20000	1580.9\\
22000	1555.1\\
24000	1589.5\\
26000	1597.7\\
28000	1599.7\\
30000	1572.9\\
32000	1567.4\\
34000	1597.5\\
38000	1567.9\\
40000	1589.7\\
42000	1581.3\\
44000	1589.3\\
46000	1573.6\\
48000	1582.1\\
50000	1556.7\\
54000	1599.2\\
56000	1586.6\\
58000	1590.5\\
62000	1589.1\\
64000	1578\\
66000	1586.5\\
68000	1584.6\\
70000	1587.8\\
72000	1572.4\\
74000	1590.8\\
76000	1574.6\\
78000	1563.1\\
80000	1556.8\\
84000	1587.2\\
86000	1582.3\\
90000	1593\\
92000	1573.2\\
94000	1593.5\\
98000	1574.5\\
1e+05	1580.9\\
};
\addplot [color=mycolor75, forget plot]
  table[row sep=crcr]{%
0	4395.3\\
2000	3039.6\\
4000	2346.4\\
6000	2286.9\\
8000	2281.8\\
10000	2284.7\\
12000	2230.6\\
14000	2251.9\\
18000	2231.8\\
20000	2245.7\\
22000	2271\\
24000	2250.4\\
26000	2251.1\\
28000	2234.1\\
30000	2236.9\\
32000	2277.1\\
34000	2279.4\\
36000	2261.6\\
38000	2250.6\\
42000	2277.7\\
44000	2248.3\\
46000	2231.4\\
48000	2258.9\\
52000	2236.1\\
54000	2256.6\\
56000	2253.1\\
58000	2264\\
60000	2246.4\\
62000	2246.2\\
64000	2264.9\\
66000	2264.5\\
68000	2233.2\\
72000	2257.5\\
74000	2289.6\\
78000	2245.1\\
84000	2263.9\\
86000	2230.6\\
88000	2244.7\\
90000	2241.7\\
92000	2221.9\\
94000	2250.7\\
96000	2273.5\\
1e+05	2261.8\\
};
\addplot [color=mycolor76, forget plot]
  table[row sep=crcr]{%
0	4395.3\\
2000	2934.7\\
4000	2187.2\\
6000	2087.8\\
8000	2058.3\\
10000	2072.5\\
12000	2074\\
14000	2061\\
16000	2087.5\\
18000	2084.9\\
20000	2053.1\\
22000	2081.3\\
24000	2042.4\\
26000	2072.1\\
28000	2055.1\\
30000	2069.1\\
32000	2068.2\\
34000	2048.9\\
36000	2050.5\\
38000	2061.1\\
40000	2055.2\\
42000	2070.3\\
46000	2038.7\\
48000	2041.5\\
50000	2063.2\\
52000	2066.3\\
54000	2087\\
56000	2073.8\\
58000	2042.4\\
62000	2055.5\\
64000	2088.5\\
66000	2091.7\\
68000	2075.8\\
70000	2047\\
74000	2064.9\\
78000	2054.3\\
82000	2084.7\\
84000	2057.4\\
90000	2049.3\\
92000	2055.5\\
94000	2048\\
96000	2062\\
1e+05	2042.1\\
};
\addplot [color=mycolor77, forget plot]
  table[row sep=crcr]{%
0	4395.3\\
2000	2572.7\\
4000	1731.4\\
6000	1636.4\\
8000	1609\\
10000	1600.4\\
12000	1616.7\\
14000	1625.5\\
16000	1596.3\\
18000	1599.6\\
20000	1625.6\\
22000	1623.8\\
24000	1606.6\\
26000	1615.1\\
28000	1600\\
30000	1626.7\\
32000	1625.7\\
34000	1603.9\\
36000	1606.7\\
38000	1633.3\\
42000	1636.3\\
44000	1618.5\\
46000	1612.5\\
48000	1584.7\\
50000	1623.7\\
52000	1637.7\\
54000	1596.4\\
56000	1599.3\\
58000	1616.6\\
60000	1591.3\\
62000	1579.2\\
64000	1609\\
66000	1624\\
68000	1618.1\\
70000	1598.9\\
72000	1602.2\\
74000	1621.4\\
76000	1609.9\\
78000	1618.8\\
80000	1621\\
82000	1606.8\\
84000	1630.4\\
86000	1632\\
88000	1603.3\\
90000	1603.4\\
92000	1598.1\\
94000	1586.2\\
96000	1608.8\\
98000	1614.2\\
1e+05	1624.2\\
};
\addplot [color=mycolor78, forget plot]
  table[row sep=crcr]{%
0	4395.3\\
2000	3133.7\\
4000	2464\\
6000	2376.2\\
8000	2319.7\\
12000	2309.2\\
14000	2316.1\\
18000	2368.3\\
20000	2338.2\\
24000	2322\\
28000	2361.7\\
30000	2330.1\\
32000	2361.5\\
34000	2295.7\\
36000	2324.6\\
38000	2306.1\\
40000	2374\\
42000	2356.9\\
44000	2350.6\\
46000	2364.7\\
48000	2327.9\\
50000	2312\\
56000	2361.8\\
58000	2320.6\\
60000	2331\\
62000	2278.1\\
64000	2311.9\\
66000	2331.1\\
68000	2374.3\\
70000	2357.5\\
74000	2312.5\\
76000	2320.8\\
78000	2337.7\\
80000	2287.1\\
82000	2317.7\\
84000	2312.7\\
86000	2321.7\\
88000	2370.9\\
90000	2368.8\\
92000	2385.9\\
94000	2349.8\\
96000	2324.5\\
98000	2359.6\\
1e+05	2368.2\\
};
\addplot [color=mycolor79, forget plot]
  table[row sep=crcr]{%
0	4395.3\\
2000	3030.3\\
4000	2327.5\\
6000	2253.9\\
8000	2192.9\\
16000	2214.9\\
18000	2173.2\\
20000	2205\\
22000	2174\\
24000	2197.9\\
26000	2204.7\\
28000	2170.5\\
30000	2174.8\\
32000	2123.8\\
34000	2189.4\\
40000	2203.6\\
42000	2172.7\\
44000	2172.3\\
46000	2221.3\\
50000	2188.7\\
52000	2219.7\\
54000	2187.3\\
56000	2208.8\\
58000	2197.2\\
60000	2174\\
62000	2201.2\\
64000	2195.7\\
66000	2175.7\\
68000	2188.2\\
70000	2177.6\\
78000	2194.5\\
80000	2205.9\\
82000	2193.1\\
86000	2204.1\\
88000	2177.7\\
90000	2171.8\\
92000	2245.8\\
94000	2191.4\\
96000	2181.3\\
98000	2188.9\\
1e+05	2232.8\\
};
\addplot [color=mycolor79, forget plot]
  table[row sep=crcr]{%
0	4395.3\\
2000	2967.2\\
4000	2191.8\\
6000	2084.1\\
8000	2092.5\\
10000	2054.1\\
12000	2068\\
14000	2095.2\\
16000	2113.6\\
20000	2043.5\\
22000	2090.1\\
24000	2061.8\\
28000	2060.6\\
30000	2031.8\\
32000	2021.6\\
34000	2056.4\\
38000	2063.9\\
40000	2081.5\\
42000	2087.8\\
44000	2072.7\\
46000	2048.5\\
48000	2077.6\\
50000	2044.5\\
52000	2032.4\\
54000	2069.1\\
56000	2063.9\\
60000	2039.2\\
62000	2072.1\\
64000	2076.5\\
66000	2063\\
68000	2029.2\\
70000	2076.4\\
72000	2054.1\\
74000	2081.3\\
76000	2072.6\\
78000	2082.9\\
80000	2061.8\\
84000	2042.8\\
86000	2065.8\\
88000	2064.8\\
90000	2123.1\\
92000	2073.3\\
94000	2091.9\\
96000	2089\\
98000	2057.7\\
1e+05	2039.6\\
};
\addplot [color=mycolor80, forget plot]
  table[row sep=crcr]{%
0	4395.3\\
2000	2656.9\\
4000	1741.8\\
6000	1606.4\\
8000	1584.2\\
10000	1603.4\\
12000	1605.4\\
16000	1585.5\\
18000	1589.7\\
20000	1610.4\\
22000	1598.3\\
24000	1591.4\\
26000	1609.4\\
28000	1588.9\\
30000	1579.5\\
32000	1593.4\\
36000	1582.6\\
38000	1594.2\\
42000	1588.7\\
44000	1607.7\\
46000	1605.9\\
48000	1614.1\\
50000	1592.7\\
52000	1594.4\\
54000	1607.7\\
56000	1578.8\\
58000	1574.6\\
60000	1627.4\\
62000	1603.7\\
64000	1598.4\\
66000	1607.5\\
68000	1595.7\\
70000	1577.9\\
72000	1590.5\\
74000	1578\\
76000	1596.7\\
78000	1587\\
80000	1604.2\\
82000	1604.7\\
84000	1593.5\\
86000	1597.5\\
88000	1616.6\\
90000	1596.8\\
92000	1603.8\\
94000	1590.7\\
96000	1586.2\\
98000	1595.7\\
1e+05	1595.2\\
};
\addplot [color=mycolor81, forget plot]
  table[row sep=crcr]{%
0	4395.3\\
2000	2600.6\\
4000	1649.4\\
6000	1525.4\\
8000	1511.1\\
14000	1500.9\\
16000	1476.5\\
18000	1522.2\\
20000	1512.7\\
22000	1488.9\\
24000	1485.1\\
26000	1487.8\\
28000	1497.1\\
30000	1483.9\\
32000	1497.5\\
34000	1493.8\\
36000	1484.4\\
40000	1504.3\\
42000	1499.4\\
44000	1481.3\\
46000	1471.2\\
48000	1478.5\\
50000	1470.6\\
54000	1503.9\\
56000	1501.3\\
58000	1482\\
60000	1499.4\\
62000	1493.4\\
64000	1512.8\\
66000	1511.8\\
68000	1499.9\\
70000	1493.3\\
72000	1481.4\\
74000	1481.9\\
76000	1498.5\\
78000	1491.6\\
80000	1489.3\\
82000	1481.6\\
86000	1489.4\\
88000	1484.2\\
90000	1493.9\\
92000	1492.2\\
94000	1498.5\\
1e+05	1500.5\\
};
\addplot [color=mycolor82, forget plot]
  table[row sep=crcr]{%
0	4395.3\\
2000	2755.9\\
4000	1896.3\\
6000	1732.7\\
8000	1737.4\\
10000	1763.7\\
12000	1741.5\\
14000	1740.5\\
18000	1724.4\\
20000	1731\\
22000	1748.8\\
24000	1745.2\\
26000	1751.4\\
28000	1765.7\\
30000	1730.1\\
32000	1740.6\\
34000	1760.8\\
36000	1769.7\\
38000	1753.7\\
40000	1705.3\\
42000	1708.1\\
44000	1731.9\\
48000	1750\\
50000	1752.9\\
54000	1701.6\\
56000	1728.6\\
62000	1723.6\\
64000	1734.7\\
66000	1739.8\\
68000	1732.8\\
70000	1715.5\\
72000	1742.4\\
74000	1732.4\\
76000	1743\\
78000	1727.2\\
80000	1723.7\\
86000	1736.2\\
88000	1758.3\\
90000	1757.4\\
92000	1743.7\\
94000	1745.5\\
96000	1726.3\\
1e+05	1721.7\\
};
\addplot [color=mycolor83, forget plot]
  table[row sep=crcr]{%
0	4395.3\\
2000	2866.3\\
4000	1954.5\\
6000	1816.8\\
8000	1770.6\\
10000	1780.9\\
12000	1755.7\\
18000	1790.6\\
20000	1750.9\\
22000	1768.6\\
24000	1780.8\\
26000	1765.7\\
28000	1786\\
30000	1776.9\\
32000	1776.8\\
34000	1791.2\\
36000	1742.1\\
38000	1745.3\\
40000	1766\\
42000	1761.3\\
46000	1762.1\\
48000	1788.1\\
50000	1758.8\\
52000	1782.5\\
54000	1772\\
56000	1752.7\\
58000	1781.3\\
60000	1769.8\\
64000	1778.5\\
66000	1759.7\\
68000	1758.8\\
70000	1781.9\\
72000	1766.9\\
74000	1765\\
76000	1750.3\\
80000	1798.2\\
82000	1784.2\\
88000	1761.9\\
90000	1764.3\\
92000	1777.4\\
94000	1769.7\\
96000	1782.3\\
98000	1762.2\\
1e+05	1752.9\\
};
\addplot [color=mycolor83, forget plot]
  table[row sep=crcr]{%
0	4395.3\\
2000	2500.7\\
4000	1461.5\\
6000	1344.4\\
8000	1294.9\\
10000	1301.4\\
12000	1296.8\\
14000	1281.1\\
16000	1291.5\\
18000	1289.1\\
20000	1299.5\\
22000	1305.7\\
24000	1298.3\\
26000	1311.1\\
28000	1296.2\\
30000	1299.9\\
32000	1298.6\\
34000	1305.4\\
36000	1303.6\\
38000	1296.9\\
40000	1300.8\\
42000	1286.1\\
44000	1299.8\\
46000	1290\\
48000	1294.2\\
50000	1347.5\\
52000	1319.3\\
54000	1317.5\\
56000	1300\\
58000	1301\\
62000	1291.8\\
64000	1278.4\\
68000	1318.7\\
72000	1321\\
74000	1305.1\\
76000	1307.8\\
78000	1297.2\\
80000	1302.8\\
82000	1292.6\\
86000	1293.5\\
88000	1282.4\\
90000	1297.2\\
92000	1316.8\\
94000	1310.4\\
96000	1298.1\\
1e+05	1286.9\\
};
\addplot [color=mycolor84, forget plot]
  table[row sep=crcr]{%
0	4395.3\\
2000	2563\\
4000	1516.3\\
6000	1384.9\\
8000	1322.5\\
10000	1329\\
12000	1313.6\\
14000	1303.7\\
16000	1316.3\\
20000	1334.6\\
22000	1326.6\\
24000	1303.3\\
26000	1324.7\\
28000	1323.1\\
30000	1311.7\\
32000	1332.1\\
34000	1320.4\\
36000	1293.3\\
38000	1295.3\\
42000	1324.7\\
44000	1316.9\\
46000	1314.7\\
48000	1341.7\\
50000	1311.2\\
52000	1333.3\\
54000	1335.3\\
56000	1330.2\\
58000	1347.7\\
60000	1347.2\\
62000	1315.5\\
64000	1319.3\\
66000	1315.9\\
68000	1325.8\\
70000	1294.8\\
72000	1299.2\\
76000	1301.5\\
78000	1313.1\\
80000	1317.5\\
82000	1317.3\\
86000	1305.9\\
90000	1328.2\\
92000	1317.3\\
94000	1292.5\\
96000	1304.1\\
98000	1311.4\\
1e+05	1308.4\\
};
\addplot [color=mycolor85, forget plot]
  table[row sep=crcr]{%
0	4395.3\\
2000	2616\\
4000	1645.6\\
6000	1466.6\\
8000	1485.5\\
10000	1442.5\\
12000	1432.1\\
14000	1443.2\\
16000	1419.4\\
18000	1434.2\\
20000	1417.7\\
24000	1417.6\\
26000	1463.9\\
28000	1433.4\\
30000	1440.4\\
32000	1430.7\\
34000	1429.1\\
36000	1422.2\\
38000	1445.1\\
40000	1437.6\\
44000	1431.7\\
46000	1451.5\\
48000	1465.6\\
50000	1395.6\\
52000	1425.4\\
54000	1469.3\\
58000	1456.1\\
60000	1420.3\\
62000	1463.7\\
64000	1436.7\\
66000	1429.1\\
70000	1430.7\\
72000	1433.2\\
74000	1419.8\\
76000	1427.9\\
78000	1415.3\\
80000	1424.1\\
82000	1420.2\\
84000	1455.3\\
86000	1460.5\\
88000	1457\\
92000	1414.7\\
94000	1401.5\\
96000	1424.2\\
98000	1426\\
1e+05	1417.8\\
};
\addplot [color=mycolor86, forget plot]
  table[row sep=crcr]{%
0	4395.3\\
2000	2601\\
4000	1458.2\\
6000	1263.2\\
10000	1188\\
12000	1210.5\\
14000	1192.7\\
16000	1197.8\\
18000	1215.3\\
20000	1248.6\\
22000	1240.9\\
24000	1213.5\\
26000	1213\\
28000	1219.2\\
32000	1205.7\\
34000	1221.5\\
38000	1227.4\\
40000	1223.7\\
42000	1214.6\\
44000	1217.6\\
46000	1216.5\\
48000	1204.2\\
52000	1231\\
54000	1218.9\\
56000	1216.6\\
58000	1224.5\\
60000	1244.9\\
62000	1216.5\\
64000	1211.4\\
66000	1210.6\\
68000	1219.8\\
70000	1223.8\\
72000	1201.8\\
74000	1205.7\\
76000	1204.3\\
80000	1232.7\\
82000	1203.4\\
86000	1203.1\\
88000	1223.1\\
90000	1211.5\\
92000	1234.8\\
96000	1219\\
98000	1225.7\\
1e+05	1206.8\\
};
\addplot [color=mycolor87, forget plot]
  table[row sep=crcr]{%
0	4395.3\\
2000	2415.7\\
4000	1234.2\\
6000	1053.9\\
8000	1051.3\\
10000	1037.9\\
14000	1029.7\\
16000	1036.1\\
18000	1070.8\\
20000	1042\\
22000	1024.7\\
24000	1016.4\\
26000	1043.7\\
28000	1065.5\\
30000	1042.7\\
32000	1047.6\\
34000	1022.8\\
36000	1046\\
38000	1053.3\\
40000	1040.8\\
42000	1044.4\\
44000	1020.9\\
46000	1029.5\\
48000	1058.4\\
52000	1039.3\\
54000	991.93\\
56000	1012.3\\
58000	1007.1\\
60000	1024.7\\
62000	1021.5\\
64000	1034.6\\
66000	1040.4\\
68000	1057.8\\
70000	1041.7\\
72000	1021.6\\
74000	1042.6\\
76000	1020.9\\
78000	1022.2\\
80000	1049.4\\
82000	1036.7\\
86000	1041.5\\
88000	1053.7\\
90000	1038.3\\
92000	1036\\
94000	1052.1\\
96000	1040.6\\
98000	1055.7\\
1e+05	1033.2\\
};
\addplot [color=mycolor88, forget plot]
  table[row sep=crcr]{%
0	4395.3\\
2000	2642.7\\
4000	1474.1\\
6000	1246.5\\
8000	1202.4\\
10000	1222.4\\
12000	1224.6\\
14000	1211.2\\
16000	1216.1\\
18000	1210.1\\
20000	1188.3\\
22000	1203.1\\
24000	1201.6\\
26000	1229.8\\
28000	1227.3\\
30000	1213.5\\
32000	1205.5\\
34000	1203.6\\
36000	1209.4\\
38000	1207.2\\
40000	1221.6\\
44000	1228.6\\
46000	1227.9\\
48000	1210.1\\
50000	1212.3\\
52000	1195.9\\
54000	1202.6\\
56000	1199\\
62000	1248.6\\
64000	1240\\
66000	1209.6\\
68000	1203.5\\
70000	1222.7\\
72000	1223.9\\
74000	1209\\
76000	1215.9\\
80000	1202.4\\
82000	1228.9\\
84000	1199.5\\
90000	1177.9\\
94000	1241.2\\
96000	1219.4\\
1e+05	1193.2\\
};
\addplot [color=mycolor89, forget plot]
  table[row sep=crcr]{%
0	4395.3\\
2000	2399.9\\
4000	1123.8\\
6000	919.33\\
8000	895.92\\
10000	890.86\\
12000	890.33\\
14000	902.53\\
16000	888.89\\
20000	890.31\\
22000	902.16\\
24000	882.92\\
26000	871.29\\
28000	865.87\\
30000	891.02\\
36000	856.95\\
40000	884\\
42000	881.81\\
44000	890.7\\
46000	866.48\\
48000	867.42\\
50000	864.61\\
52000	887.92\\
54000	886.87\\
56000	903.71\\
58000	896.57\\
60000	867.89\\
62000	864.07\\
66000	884.99\\
68000	884.44\\
70000	879.12\\
72000	866.68\\
74000	888.37\\
76000	873.62\\
82000	876.7\\
88000	881.44\\
90000	878.29\\
92000	895.42\\
94000	888.11\\
98000	897.14\\
1e+05	892.86\\
};
\addplot [color=mycolor90, forget plot]
  table[row sep=crcr]{%
0	4395.3\\
2000	2377.1\\
4000	1113.3\\
6000	897.67\\
8000	858.07\\
10000	842.42\\
14000	845.65\\
16000	838.07\\
18000	847.21\\
20000	853.52\\
22000	839.89\\
24000	853.42\\
26000	834.31\\
28000	843.83\\
30000	833.69\\
32000	845.05\\
34000	850.35\\
36000	846.8\\
38000	848.63\\
40000	853.9\\
44000	844.92\\
46000	839.48\\
48000	850.2\\
50000	856.35\\
52000	858.41\\
54000	844.98\\
56000	848.23\\
60000	860.47\\
62000	847.01\\
64000	845.82\\
66000	832.28\\
68000	855.79\\
70000	859.01\\
72000	842.79\\
74000	862.88\\
76000	855.95\\
78000	860.39\\
80000	845.26\\
84000	861.96\\
86000	865.54\\
88000	883.05\\
90000	875.47\\
92000	860.21\\
94000	851.76\\
96000	870.01\\
98000	839.86\\
1e+05	837.39\\
};
\addplot [color=mycolor91, forget plot]
  table[row sep=crcr]{%
0	4395.3\\
2000	2365.3\\
4000	1070.3\\
6000	807.88\\
8000	780.76\\
10000	744.11\\
12000	732.29\\
14000	755.96\\
16000	756.53\\
18000	751.13\\
20000	741.77\\
22000	740.6\\
24000	719.6\\
26000	725.21\\
28000	727.15\\
30000	743.54\\
32000	729.95\\
34000	742.58\\
36000	737.51\\
38000	728.47\\
40000	752.86\\
42000	728.61\\
44000	750.35\\
46000	753.27\\
48000	738.73\\
50000	770.77\\
52000	756.37\\
54000	763.02\\
58000	745.64\\
60000	720.39\\
62000	732.18\\
64000	725.87\\
66000	742.13\\
68000	741.82\\
70000	754.72\\
72000	748\\
74000	729.8\\
76000	747.29\\
78000	752.15\\
80000	750.74\\
82000	764.93\\
84000	739.89\\
88000	759.09\\
90000	734.61\\
92000	719.28\\
94000	744.92\\
96000	738.61\\
98000	742.6\\
1e+05	754.48\\
};
\addplot [color=mycolor92, forget plot]
  table[row sep=crcr]{%
0	4395.3\\
2000	2550.5\\
4000	1277.9\\
6000	992.29\\
8000	924.61\\
10000	897.96\\
12000	922.37\\
14000	913.57\\
16000	938.51\\
18000	924.95\\
20000	941.16\\
22000	925.1\\
30000	925.77\\
32000	935.49\\
36000	913.6\\
40000	933.69\\
44000	903.04\\
46000	906.3\\
48000	939.47\\
50000	934.59\\
52000	935.27\\
54000	922.1\\
56000	916.81\\
58000	901.28\\
60000	888.77\\
62000	925.24\\
66000	905.78\\
68000	918.32\\
70000	928.36\\
72000	931.79\\
74000	909.27\\
76000	914.24\\
78000	915.68\\
80000	924.11\\
82000	921.28\\
84000	924.68\\
86000	916.25\\
88000	901.1\\
92000	915.2\\
94000	907.49\\
98000	924.43\\
1e+05	928.18\\
};
\addplot [color=mycolor93, forget plot]
  table[row sep=crcr]{%
0	4395.3\\
2000	2551.8\\
4000	1143.9\\
6000	854.17\\
8000	792.62\\
10000	784.27\\
12000	772.29\\
14000	764.84\\
16000	761.17\\
18000	764.16\\
20000	747.37\\
22000	755.37\\
24000	777.74\\
26000	814.21\\
28000	814.95\\
30000	777.79\\
32000	770.66\\
34000	792.66\\
36000	803.83\\
38000	772.69\\
40000	770.72\\
42000	792.7\\
44000	778.88\\
46000	779.14\\
48000	745.34\\
54000	769.54\\
56000	769.26\\
60000	746.78\\
62000	745.96\\
66000	799.62\\
68000	780.01\\
72000	765.39\\
74000	774.23\\
76000	791.96\\
78000	776.27\\
80000	801.35\\
82000	824.65\\
84000	786.08\\
86000	782.37\\
88000	788.16\\
90000	784.76\\
92000	750.45\\
94000	746.71\\
96000	771.54\\
98000	780.19\\
1e+05	807.5\\
};
\addplot [color=mycolor94, forget plot]
  table[row sep=crcr]{%
0	4395.3\\
2000	2579.7\\
4000	1188.1\\
6000	871.85\\
8000	781.95\\
12000	726.32\\
20000	741.87\\
22000	754.58\\
26000	750.58\\
28000	738.77\\
30000	743.85\\
32000	739.48\\
34000	725.15\\
36000	724.65\\
38000	719.2\\
40000	720.28\\
42000	717.74\\
44000	718.96\\
46000	738.2\\
48000	741.47\\
50000	736.76\\
52000	742.15\\
54000	740.81\\
56000	715.66\\
58000	729.64\\
60000	728.35\\
62000	755.08\\
64000	745.37\\
66000	733.28\\
68000	732.84\\
70000	727.09\\
72000	736.7\\
74000	734.74\\
76000	742.42\\
78000	734.04\\
80000	730.84\\
82000	738.42\\
84000	729.89\\
86000	750.13\\
88000	738.53\\
90000	730.66\\
92000	756.19\\
94000	754.97\\
96000	736.08\\
98000	740.53\\
1e+05	748.58\\
};
\addplot [color=mycolor95, forget plot]
  table[row sep=crcr]{%
0	4395.3\\
2000	2343.2\\
4000	835.61\\
6000	549.55\\
8000	442.83\\
10000	439\\
12000	437.36\\
14000	441.95\\
16000	440.6\\
18000	440.53\\
20000	438.06\\
24000	449.24\\
28000	445.2\\
30000	448.45\\
32000	443.4\\
34000	440.07\\
36000	446.97\\
38000	440.21\\
40000	443.6\\
42000	451.89\\
44000	446.76\\
46000	451.18\\
48000	448.46\\
50000	438.13\\
52000	444.28\\
54000	454.78\\
56000	441.49\\
58000	445.54\\
60000	440.8\\
62000	432.12\\
64000	435.67\\
66000	447.06\\
68000	445.99\\
72000	439.56\\
74000	432.79\\
76000	442.63\\
78000	437.23\\
80000	433.87\\
82000	440.63\\
84000	431.02\\
88000	442.92\\
90000	447.19\\
92000	449.45\\
94000	429.27\\
96000	422.86\\
1e+05	448.72\\
};
\addplot [color=mycolor96, forget plot]
  table[row sep=crcr]{%
0	4395.3\\
2000	2486.9\\
4000	955.93\\
6000	622.55\\
8000	504.73\\
10000	468.37\\
12000	467.99\\
14000	478.39\\
16000	490.41\\
18000	489.45\\
20000	479.27\\
22000	475.21\\
24000	465.67\\
26000	471.3\\
28000	479.98\\
30000	461.43\\
32000	485.82\\
34000	497.88\\
36000	502.48\\
38000	482.89\\
40000	472.39\\
42000	478.81\\
44000	457.8\\
48000	491.8\\
50000	491.23\\
52000	484.31\\
54000	446.61\\
56000	462.36\\
58000	489.21\\
60000	480.93\\
62000	456.26\\
64000	456.9\\
66000	476.96\\
68000	488.95\\
70000	455.53\\
72000	469.45\\
74000	476.28\\
76000	473.37\\
78000	482.45\\
80000	472.52\\
82000	446.71\\
84000	444.41\\
86000	481.15\\
88000	476.91\\
90000	463.83\\
92000	461.47\\
94000	450.11\\
96000	445.84\\
98000	454.5\\
1e+05	446.42\\
};
\addplot [color=mycolor97, forget plot]
  table[row sep=crcr]{%
0	4395.3\\
2000	2393.8\\
4000	853.12\\
6000	478.62\\
8000	361.14\\
10000	352.44\\
12000	357.53\\
14000	332.28\\
16000	324.32\\
18000	304.12\\
20000	326.82\\
22000	332.69\\
24000	333.92\\
26000	324.31\\
28000	330.81\\
30000	322.16\\
32000	326.38\\
34000	326.42\\
36000	322.79\\
38000	324.55\\
42000	321.81\\
44000	335.07\\
46000	340.41\\
48000	330.23\\
50000	326.28\\
52000	311.96\\
54000	319.42\\
56000	332.32\\
58000	318.92\\
60000	314.57\\
62000	323.3\\
64000	339.78\\
66000	335.14\\
68000	328.87\\
70000	323.71\\
72000	314.11\\
74000	326.5\\
76000	321.22\\
78000	327.18\\
80000	330.56\\
82000	327.9\\
84000	334.86\\
86000	350.88\\
88000	327.43\\
90000	336.55\\
92000	319\\
94000	321.67\\
96000	343.23\\
98000	335.88\\
1e+05	331.4\\
};
\addplot [color=mycolor98]
  table[row sep=crcr]{%
0	4395.3\\
2000	2672.6\\
4000	1116.9\\
6000	661.14\\
8000	525.46\\
10000	473.85\\
12000	452.69\\
14000	459.12\\
16000	444.79\\
18000	463.29\\
20000	451.08\\
22000	464.81\\
24000	449.22\\
26000	453.05\\
28000	466.23\\
30000	460.47\\
32000	448.05\\
34000	445.59\\
36000	456.54\\
38000	453.11\\
40000	456.32\\
42000	443.71\\
44000	438.37\\
46000	448.83\\
48000	451.65\\
50000	440.47\\
52000	446.38\\
54000	439.39\\
56000	440.76\\
58000	437.83\\
60000	443.22\\
62000	460.31\\
66000	456.41\\
68000	458.12\\
70000	448.12\\
72000	451.37\\
74000	460.74\\
76000	452.19\\
78000	457.68\\
80000	442.02\\
82000	446.32\\
84000	445.01\\
86000	458.34\\
88000	455.19\\
90000	458.73\\
92000	459.42\\
94000	466.39\\
96000	455.2\\
98000	432.46\\
1e+05	448.61\\
};
\addlegendentry{Classe k=10}

            % \input{../../../../.tex/20/KS/SizeEvolutionsfig4.tex}
            \legend{}
        \end{groupplot}

        \begin{groupplot}[
            group style={
                group name=Row3,
                plotGroup,
                group size=2 by 1,
            },
            plotRow,
            row5Keys,
            xTickScaleLabelEmpty,
            plotLegendSE,
            %
            xmin=0,xmax=1e+05,
            ymin=100,ymax=1.5501e+05,
        ]
            \nextgroupplot[%
                at={($(Row2 c1r1.south west)-(0,\yPlotSep em)-(0,\plotHeight cm)$)},
                % xmin=0,xmax=1e+05,
                % ymin=100,ymax=1.5126e+05,
            ] % This file was created by matlab2tikz.
%
\definecolor{mycolor1}{rgb}{0.00146,0.00047,0.01387}%
\definecolor{mycolor2}{rgb}{0.01166,0.00942,0.06346}%
\definecolor{mycolor3}{rgb}{0.02943,0.02150,0.11462}%
\definecolor{mycolor4}{rgb}{0.06134,0.03659,0.17764}%
\definecolor{mycolor5}{rgb}{0.08741,0.04456,0.22481}%
\definecolor{mycolor6}{rgb}{0.12291,0.04754,0.28162}%
\definecolor{mycolor7}{rgb}{0.14907,0.04547,0.31709}%
\definecolor{mycolor8}{rgb}{0.18343,0.04033,0.35497}%
\definecolor{mycolor9}{rgb}{0.21795,0.03662,0.38352}%
\definecolor{mycolor10}{rgb}{0.24497,0.03705,0.40001}%
\definecolor{mycolor11}{rgb}{0.27135,0.04092,0.41198}%
\definecolor{mycolor12}{rgb}{0.29076,0.04564,0.41864}%
\definecolor{mycolor13}{rgb}{0.31628,0.05349,0.42512}%
\definecolor{mycolor14}{rgb}{0.33522,0.06006,0.42852}%
\definecolor{mycolor15}{rgb}{0.36028,0.06925,0.43150}%
\definecolor{mycolor16}{rgb}{0.37900,0.07625,0.43272}%
\definecolor{mycolor17}{rgb}{0.39767,0.08326,0.43318}%
\definecolor{mycolor18}{rgb}{0.41633,0.09020,0.43294}%
\definecolor{mycolor19}{rgb}{0.42877,0.09479,0.43241}%
\definecolor{mycolor20}{rgb}{0.44743,0.10160,0.43108}%
\definecolor{mycolor21}{rgb}{0.46610,0.10832,0.42912}%
\definecolor{mycolor22}{rgb}{0.47856,0.11276,0.42747}%
\definecolor{mycolor23}{rgb}{0.49726,0.11938,0.42449}%
\definecolor{mycolor24}{rgb}{0.50973,0.12377,0.42216}%
\definecolor{mycolor25}{rgb}{0.52221,0.12815,0.41955}%
\definecolor{mycolor26}{rgb}{0.53468,0.13253,0.41667}%
\definecolor{mycolor27}{rgb}{0.55339,0.13913,0.41183}%
\definecolor{mycolor28}{rgb}{0.56585,0.14357,0.40826}%
\definecolor{mycolor29}{rgb}{0.57830,0.14804,0.40441}%
\definecolor{mycolor30}{rgb}{0.59073,0.15256,0.40029}%
\definecolor{mycolor31}{rgb}{0.60314,0.15715,0.39589}%
\definecolor{mycolor32}{rgb}{0.61551,0.16182,0.39122}%
\definecolor{mycolor33}{rgb}{0.62169,0.16418,0.38878}%
\definecolor{mycolor34}{rgb}{0.63400,0.16899,0.38370}%
\definecolor{mycolor35}{rgb}{0.64626,0.17391,0.37836}%
\definecolor{mycolor36}{rgb}{0.65846,0.17896,0.37275}%
\definecolor{mycolor37}{rgb}{0.66454,0.18154,0.36985}%
\definecolor{mycolor38}{rgb}{0.67664,0.18681,0.36385}%
\definecolor{mycolor39}{rgb}{0.68865,0.19224,0.35760}%
\definecolor{mycolor40}{rgb}{0.69463,0.19502,0.35439}%
\definecolor{mycolor41}{rgb}{0.70650,0.20073,0.34778}%
\definecolor{mycolor42}{rgb}{0.71240,0.20366,0.34438}%
\definecolor{mycolor43}{rgb}{0.72410,0.20967,0.33742}%
\definecolor{mycolor44}{rgb}{0.72991,0.21276,0.33386}%
\definecolor{mycolor45}{rgb}{0.74142,0.21911,0.32658}%
\definecolor{mycolor46}{rgb}{0.74713,0.22238,0.32286}%
\definecolor{mycolor47}{rgb}{0.75279,0.22571,0.31909}%
\definecolor{mycolor48}{rgb}{0.76401,0.23255,0.31140}%
\definecolor{mycolor49}{rgb}{0.76956,0.23608,0.30749}%
\definecolor{mycolor50}{rgb}{0.77506,0.23967,0.30353}%
\definecolor{mycolor51}{rgb}{0.79129,0.25086,0.29139}%
\definecolor{mycolor52}{rgb}{0.79661,0.25473,0.28726}%
\definecolor{mycolor53}{rgb}{0.80187,0.25867,0.28310}%
\definecolor{mycolor54}{rgb}{0.81224,0.26679,0.27466}%
\definecolor{mycolor55}{rgb}{0.81734,0.27095,0.27039}%
\definecolor{mycolor56}{rgb}{0.82737,0.27952,0.26175}%
\definecolor{mycolor57}{rgb}{0.83230,0.28391,0.25738}%
\definecolor{mycolor58}{rgb}{0.83717,0.28839,0.25299}%
\definecolor{mycolor59}{rgb}{0.84671,0.29756,0.24411}%
\definecolor{mycolor60}{rgb}{0.85138,0.30226,0.23964}%
\definecolor{mycolor61}{rgb}{0.85599,0.30704,0.23513}%
\definecolor{mycolor62}{rgb}{0.86053,0.31189,0.23061}%
\definecolor{mycolor63}{rgb}{0.87374,0.32691,0.21689}%
\definecolor{mycolor64}{rgb}{0.87800,0.33206,0.21227}%
\definecolor{mycolor65}{rgb}{0.89034,0.34796,0.19829}%
\definecolor{mycolor66}{rgb}{0.89819,0.35891,0.18886}%
\definecolor{mycolor67}{rgb}{0.90200,0.36449,0.18412}%
\definecolor{mycolor68}{rgb}{0.90573,0.37014,0.17935}%
\definecolor{mycolor69}{rgb}{0.90939,0.37586,0.17456}%
\definecolor{mycolor70}{rgb}{0.91297,0.38164,0.16975}%
\definecolor{mycolor71}{rgb}{0.91646,0.38748,0.16492}%
\definecolor{mycolor72}{rgb}{0.91988,0.39339,0.16007}%
\definecolor{mycolor73}{rgb}{0.92322,0.39936,0.15519}%
\definecolor{mycolor74}{rgb}{0.92647,0.40539,0.15029}%
\definecolor{mycolor75}{rgb}{0.92964,0.41148,0.14537}%
\definecolor{mycolor76}{rgb}{0.93274,0.41763,0.14042}%
\definecolor{mycolor77}{rgb}{0.93575,0.42383,0.13544}%
\definecolor{mycolor78}{rgb}{0.93868,0.43009,0.13044}%
\definecolor{mycolor79}{rgb}{0.94152,0.43640,0.12541}%
\definecolor{mycolor80}{rgb}{0.94429,0.44277,0.12035}%
\definecolor{mycolor81}{rgb}{0.94956,0.45566,0.11016}%
\definecolor{mycolor82}{rgb}{0.95208,0.46218,0.10503}%
\definecolor{mycolor83}{rgb}{0.96129,0.48872,0.08429}%
\definecolor{mycolor84}{rgb}{0.96540,0.50225,0.07386}%
\definecolor{mycolor85}{rgb}{0.96732,0.50908,0.06866}%
\definecolor{mycolor86}{rgb}{0.97259,0.52980,0.05332}%
\definecolor{mycolor87}{rgb}{0.97568,0.54380,0.04362}%
\definecolor{mycolor88}{rgb}{0.98082,0.57221,0.02851}%
\definecolor{mycolor89}{rgb}{0.98189,0.57939,0.02625}%
\definecolor{mycolor90}{rgb}{0.98378,0.59385,0.02377}%
\definecolor{mycolor91}{rgb}{0.98532,0.60842,0.02420}%
\definecolor{mycolor92}{rgb}{0.98696,0.63048,0.03091}%
\definecolor{mycolor93}{rgb}{0.98793,0.66025,0.05175}%
\definecolor{mycolor94}{rgb}{0.98771,0.68281,0.07249}%
\definecolor{mycolor95}{rgb}{0.97750,0.78226,0.18592}%
\definecolor{mycolor96}{rgb}{0.96624,0.83619,0.26153}%
\definecolor{mycolor97}{rgb}{0.96252,0.85148,0.28555}%
\definecolor{mycolor98}{rgb}{0.98836,0.99836,0.64492}%
%
\addplot [color=mycolor1]
  table[row sep=crcr]{%
0	4395.3\\
2000	94531\\
4000	1.3812e+05\\
6000	1.4618e+05\\
8000	1.4644e+05\\
10000	1.518e+05\\
12000	1.5096e+05\\
14000	1.4711e+05\\
16000	1.486e+05\\
22000	1.4849e+05\\
24000	1.4743e+05\\
26000	1.5247e+05\\
28000	1.4977e+05\\
30000	1.5046e+05\\
32000	1.5041e+05\\
34000	1.5086e+05\\
36000	1.4748e+05\\
38000	1.4779e+05\\
42000	1.5019e+05\\
44000	1.5024e+05\\
48000	1.4928e+05\\
52000	1.5211e+05\\
54000	1.5087e+05\\
56000	1.4886e+05\\
58000	1.4852e+05\\
60000	1.5199e+05\\
62000	1.5248e+05\\
64000	1.5038e+05\\
66000	1.5023e+05\\
68000	1.4946e+05\\
70000	1.5143e+05\\
74000	1.4911e+05\\
76000	1.4704e+05\\
78000	1.4751e+05\\
80000	1.4988e+05\\
86000	1.4752e+05\\
88000	1.4856e+05\\
90000	1.5016e+05\\
92000	1.4887e+05\\
96000	1.4997e+05\\
98000	1.5196e+05\\
1e+05	1.5023e+05\\
};
\addlegendentry{Classe k=283}

\addplot [color=mycolor2, forget plot]
  table[row sep=crcr]{%
0	4395.3\\
2000	24749\\
4000	25274\\
6000	25731\\
8000	25914\\
10000	25398\\
12000	25490\\
14000	25957\\
16000	25833\\
18000	25441\\
20000	26182\\
26000	25182\\
30000	25732\\
32000	25641\\
34000	25018\\
36000	25548\\
38000	25684\\
40000	25621\\
42000	25749\\
44000	25556\\
46000	25646\\
50000	25503\\
52000	26027\\
54000	24895\\
56000	25976\\
58000	25879\\
60000	25349\\
62000	25505\\
64000	25920\\
66000	25613\\
68000	25893\\
72000	25818\\
74000	25624\\
76000	25708\\
80000	26152\\
82000	25776\\
86000	25598\\
88000	25387\\
90000	25987\\
92000	26370\\
94000	25515\\
96000	25741\\
98000	25505\\
1e+05	25974\\
};
\addplot [color=mycolor3, forget plot]
  table[row sep=crcr]{%
0	4395.3\\
2000	34244\\
4000	38146\\
6000	38613\\
8000	38358\\
10000	38762\\
12000	39421\\
14000	39933\\
16000	39944\\
18000	39549\\
20000	39778\\
24000	39446\\
26000	39017\\
28000	38744\\
30000	37950\\
32000	38942\\
34000	39014\\
36000	39448\\
38000	39232\\
42000	38574\\
44000	39458\\
46000	39557\\
48000	39397\\
50000	38527\\
52000	39192\\
54000	39004\\
56000	40154\\
58000	39139\\
60000	38609\\
62000	39503\\
64000	38568\\
66000	38964\\
68000	38821\\
70000	38281\\
72000	39510\\
74000	39566\\
76000	38725\\
78000	38850\\
80000	39375\\
82000	39557\\
84000	37957\\
86000	38615\\
88000	39857\\
90000	40296\\
92000	39002\\
96000	38537\\
1e+05	38462\\
};
\addplot [color=mycolor4, forget plot]
  table[row sep=crcr]{%
0	4395.3\\
2000	14833\\
4000	14334\\
6000	14174\\
8000	14434\\
10000	14447\\
12000	14184\\
14000	14242\\
16000	14402\\
18000	14159\\
20000	14309\\
22000	14404\\
24000	14125\\
26000	14179\\
28000	14126\\
30000	14209\\
32000	14374\\
34000	13965\\
36000	14177\\
38000	13880\\
40000	14237\\
42000	13943\\
44000	13807\\
46000	14442\\
48000	14138\\
52000	14188\\
54000	14320\\
56000	14030\\
58000	14149\\
60000	14039\\
62000	14020\\
64000	14257\\
66000	13957\\
68000	14225\\
70000	14233\\
72000	14091\\
74000	14046\\
76000	14088\\
78000	14592\\
80000	14570\\
82000	13916\\
84000	14200\\
86000	14137\\
88000	14025\\
90000	14212\\
92000	14449\\
94000	14529\\
96000	14063\\
98000	14065\\
1e+05	14114\\
};
\addplot [color=mycolor5, forget plot]
  table[row sep=crcr]{%
0	4395.3\\
2000	46056\\
4000	62952\\
6000	64504\\
8000	63313\\
10000	61933\\
12000	62389\\
14000	62191\\
16000	63089\\
18000	63796\\
20000	63335\\
22000	61418\\
24000	62529\\
26000	60443\\
28000	61917\\
30000	62368\\
32000	62477\\
34000	62150\\
36000	63688\\
38000	63944\\
40000	62237\\
42000	64235\\
44000	62510\\
46000	62207\\
48000	60784\\
50000	61103\\
52000	59948\\
54000	61339\\
56000	63325\\
60000	62524\\
62000	62565\\
64000	63379\\
66000	62065\\
68000	63041\\
70000	62395\\
72000	63483\\
74000	64011\\
76000	64111\\
78000	63881\\
80000	62437\\
82000	63916\\
84000	64036\\
86000	61751\\
88000	62922\\
92000	61379\\
94000	62029\\
96000	62926\\
1e+05	62250\\
};
\addplot [color=mycolor6, forget plot]
  table[row sep=crcr]{%
0	4395.3\\
2000	38102\\
4000	53152\\
6000	52998\\
8000	53900\\
10000	53481\\
12000	55058\\
14000	55439\\
16000	55079\\
18000	56693\\
20000	53741\\
22000	55019\\
26000	55196\\
28000	54363\\
30000	54780\\
32000	55945\\
34000	55767\\
36000	55979\\
38000	56491\\
40000	55815\\
42000	54410\\
44000	53273\\
46000	54675\\
48000	55195\\
50000	55300\\
52000	56354\\
54000	54835\\
56000	54826\\
58000	55415\\
62000	56142\\
64000	56112\\
66000	54009\\
68000	55184\\
70000	54963\\
72000	54970\\
74000	56949\\
76000	57163\\
78000	58004\\
80000	57167\\
82000	56541\\
84000	56442\\
86000	55383\\
88000	54007\\
90000	56854\\
92000	55447\\
94000	55491\\
96000	55939\\
98000	54465\\
1e+05	54757\\
};
\addplot [color=mycolor7, forget plot]
  table[row sep=crcr]{%
0	4395.3\\
2000	22170\\
4000	26878\\
6000	28427\\
8000	27404\\
10000	27052\\
12000	27603\\
14000	26733\\
16000	27795\\
18000	27727\\
20000	28318\\
22000	27572\\
24000	28339\\
26000	27254\\
28000	27671\\
30000	27667\\
32000	27037\\
34000	27656\\
36000	28154\\
38000	28028\\
40000	27554\\
42000	27898\\
44000	27577\\
46000	28416\\
48000	28234\\
50000	28932\\
52000	28026\\
54000	27749\\
56000	28528\\
58000	27959\\
60000	28444\\
62000	27496\\
64000	28132\\
66000	27929\\
68000	27576\\
70000	28360\\
72000	27890\\
74000	27564\\
76000	27849\\
78000	27437\\
80000	27739\\
82000	27354\\
84000	28473\\
86000	28106\\
88000	27887\\
92000	27236\\
94000	28275\\
96000	28887\\
98000	28871\\
1e+05	27891\\
};
\addplot [color=mycolor8, forget plot]
  table[row sep=crcr]{%
0	4395.3\\
2000	20042\\
4000	21604\\
6000	20441\\
8000	20388\\
10000	20716\\
12000	20451\\
14000	20291\\
16000	20474\\
18000	19994\\
20000	20395\\
22000	20537\\
24000	20309\\
26000	20645\\
28000	19691\\
30000	20784\\
32000	20712\\
36000	20236\\
38000	20209\\
40000	20588\\
42000	19722\\
44000	20389\\
46000	20682\\
48000	20266\\
50000	20238\\
52000	19863\\
54000	19430\\
56000	20399\\
58000	19993\\
60000	20292\\
62000	20381\\
64000	20828\\
66000	20476\\
68000	19900\\
70000	19827\\
74000	20040\\
76000	19974\\
78000	19423\\
82000	20115\\
84000	20411\\
86000	20510\\
88000	20422\\
90000	20107\\
92000	20226\\
94000	20452\\
96000	20442\\
98000	20193\\
1e+05	20397\\
};
\addplot [color=mycolor9, forget plot]
  table[row sep=crcr]{%
0	4395.3\\
2000	13068\\
4000	13855\\
6000	13923\\
8000	14098\\
10000	13796\\
12000	13609\\
14000	13936\\
16000	14312\\
18000	14090\\
20000	13950\\
22000	13737\\
24000	13884\\
26000	13871\\
28000	14070\\
30000	13967\\
32000	14121\\
34000	13975\\
38000	13611\\
40000	14432\\
42000	14116\\
44000	14509\\
46000	14117\\
48000	14164\\
52000	13656\\
54000	14097\\
56000	14250\\
58000	13593\\
60000	14350\\
62000	13969\\
64000	14112\\
66000	14087\\
68000	14133\\
70000	13761\\
72000	14353\\
74000	14394\\
76000	13725\\
78000	14140\\
80000	13651\\
82000	13927\\
86000	14090\\
88000	13778\\
90000	13524\\
92000	13936\\
94000	14421\\
96000	13878\\
98000	13771\\
1e+05	13785\\
};
\addplot [color=mycolor10, forget plot]
  table[row sep=crcr]{%
0	4395.3\\
2000	17670\\
4000	21156\\
6000	21204\\
8000	21666\\
10000	21459\\
12000	21921\\
14000	22535\\
16000	22223\\
18000	21430\\
20000	21392\\
22000	21948\\
24000	21465\\
26000	21892\\
28000	22087\\
30000	21824\\
32000	21673\\
34000	21843\\
36000	21564\\
38000	21698\\
40000	21649\\
42000	21710\\
44000	22102\\
46000	21487\\
48000	21670\\
50000	21674\\
52000	21973\\
54000	21920\\
56000	21624\\
58000	22101\\
60000	22079\\
62000	21973\\
64000	21640\\
66000	22081\\
68000	22129\\
70000	21808\\
72000	21412\\
74000	21779\\
76000	21883\\
78000	21656\\
80000	22002\\
82000	21893\\
84000	21648\\
86000	21856\\
88000	21557\\
90000	21541\\
92000	22258\\
94000	22128\\
96000	21569\\
98000	22169\\
1e+05	21420\\
};
\addplot [color=mycolor11, forget plot]
  table[row sep=crcr]{%
0	4395.3\\
2000	9728\\
4000	9050\\
6000	8800.4\\
8000	8716.4\\
10000	8537.9\\
14000	8682.7\\
16000	8909.5\\
18000	8779.6\\
22000	8603.3\\
24000	8672.6\\
26000	8829.4\\
28000	8527.8\\
30000	8409.7\\
32000	8469.5\\
34000	8571.1\\
36000	8738\\
38000	8882.9\\
40000	8838\\
42000	8573.7\\
44000	8561.6\\
46000	8681.5\\
50000	8652.4\\
52000	8418.4\\
54000	8469.5\\
56000	8806.4\\
58000	8825.5\\
60000	8674.6\\
62000	8604.5\\
64000	8403.9\\
66000	8686\\
68000	8767.2\\
70000	8335.8\\
72000	8605.4\\
76000	8323.8\\
80000	8397.3\\
82000	8571.9\\
84000	8650.3\\
88000	8361.2\\
90000	8591.6\\
92000	8741.9\\
96000	8661.1\\
98000	8431.7\\
1e+05	8602.1\\
};
\addplot [color=mycolor12, forget plot]
  table[row sep=crcr]{%
0	4395.3\\
2000	12907\\
4000	14668\\
8000	14521\\
10000	14758\\
12000	14551\\
16000	14581\\
18000	14841\\
20000	14645\\
22000	14730\\
24000	14650\\
26000	14334\\
28000	14741\\
30000	14816\\
32000	14190\\
34000	14340\\
36000	14748\\
38000	14933\\
40000	14670\\
42000	14943\\
44000	14534\\
46000	14864\\
48000	14301\\
50000	14950\\
52000	14163\\
54000	14855\\
56000	14858\\
58000	14252\\
60000	14277\\
62000	14505\\
64000	14601\\
66000	14440\\
68000	14552\\
70000	15055\\
72000	14842\\
74000	14852\\
78000	14308\\
80000	14523\\
82000	14507\\
84000	14861\\
86000	14929\\
88000	14792\\
90000	14467\\
92000	14815\\
94000	14610\\
96000	14505\\
98000	13846\\
1e+05	14588\\
};
\addplot [color=mycolor13, forget plot]
  table[row sep=crcr]{%
0	4395.3\\
2000	10105\\
4000	11064\\
6000	10937\\
8000	11401\\
10000	11049\\
12000	11163\\
14000	11143\\
16000	10893\\
18000	10689\\
20000	11102\\
22000	11227\\
24000	11070\\
26000	11265\\
28000	10898\\
30000	11212\\
32000	11253\\
34000	11357\\
36000	11329\\
38000	11396\\
40000	11316\\
44000	11326\\
46000	11065\\
48000	10896\\
50000	11145\\
52000	10972\\
54000	11234\\
56000	11101\\
58000	11036\\
60000	10806\\
62000	11099\\
64000	11107\\
68000	11019\\
70000	11174\\
72000	11013\\
74000	11207\\
76000	11074\\
78000	10974\\
80000	10712\\
82000	10922\\
84000	11052\\
86000	11133\\
88000	11001\\
90000	10966\\
92000	11129\\
94000	10988\\
96000	10790\\
98000	11053\\
1e+05	10742\\
};
\addplot [color=mycolor14, forget plot]
  table[row sep=crcr]{%
0	4395.3\\
2000	18238\\
4000	23005\\
6000	23837\\
8000	24215\\
10000	24825\\
14000	23907\\
16000	23689\\
18000	24239\\
20000	23634\\
22000	24199\\
24000	24395\\
26000	24520\\
28000	24052\\
30000	24129\\
32000	24649\\
34000	23657\\
36000	24233\\
38000	24252\\
40000	24454\\
42000	23671\\
44000	24541\\
46000	23315\\
48000	24801\\
50000	24272\\
52000	24106\\
54000	24384\\
56000	23627\\
58000	24199\\
60000	24081\\
62000	24037\\
64000	24461\\
66000	24512\\
68000	23765\\
70000	23838\\
72000	24302\\
74000	24466\\
78000	23803\\
80000	25237\\
82000	24188\\
84000	24204\\
86000	25136\\
88000	24383\\
90000	23245\\
92000	24068\\
94000	24685\\
96000	23853\\
1e+05	24592\\
};
\addplot [color=mycolor15, forget plot]
  table[row sep=crcr]{%
0	4395.3\\
2000	13963\\
4000	16386\\
6000	16230\\
8000	16500\\
10000	16897\\
12000	16732\\
14000	16421\\
16000	17019\\
18000	16673\\
20000	16712\\
22000	16672\\
24000	17082\\
26000	16936\\
28000	17159\\
30000	16611\\
32000	17211\\
34000	16436\\
36000	16392\\
38000	17017\\
40000	16884\\
42000	16256\\
44000	16510\\
46000	17240\\
48000	17158\\
50000	16794\\
52000	16686\\
54000	16959\\
56000	16957\\
58000	16889\\
60000	17012\\
62000	16798\\
64000	16760\\
68000	17204\\
70000	16737\\
74000	16991\\
76000	16588\\
78000	17406\\
80000	16568\\
84000	17391\\
86000	16558\\
88000	17148\\
90000	17278\\
92000	17036\\
94000	16964\\
96000	17161\\
98000	17059\\
1e+05	16848\\
};
\addplot [color=mycolor16, forget plot]
  table[row sep=crcr]{%
0	4395.3\\
2000	8107.8\\
4000	7151.7\\
6000	7067.3\\
8000	6899.7\\
10000	6811.1\\
12000	6912.1\\
14000	6803\\
16000	6856.2\\
18000	7029.2\\
22000	7039.6\\
24000	6993.3\\
26000	7116.8\\
28000	7090.3\\
32000	6971.1\\
34000	7125.8\\
36000	7015.6\\
38000	7068.8\\
40000	6860.2\\
42000	6921.7\\
44000	7050.7\\
46000	6900.1\\
48000	7061.5\\
50000	7082.8\\
52000	6983.7\\
54000	7107.5\\
56000	7019.7\\
58000	6852.8\\
60000	7003.8\\
62000	6851\\
64000	7009.3\\
66000	7039.7\\
68000	7154.8\\
70000	7099\\
72000	6984.3\\
74000	7171.8\\
76000	6982.3\\
78000	7081.7\\
80000	6894.5\\
82000	7001.4\\
84000	6821\\
86000	7187.6\\
88000	7113.9\\
90000	7007.1\\
92000	6863.4\\
94000	7106.6\\
96000	6852.5\\
98000	6928.7\\
1e+05	6858\\
};
\addplot [color=mycolor17, forget plot]
  table[row sep=crcr]{%
0	4395.3\\
2000	13239\\
4000	14941\\
6000	15982\\
8000	15755\\
10000	16092\\
12000	15821\\
14000	16128\\
16000	16006\\
18000	15717\\
20000	16440\\
22000	15796\\
24000	15717\\
26000	15841\\
28000	16090\\
30000	15671\\
32000	15426\\
34000	16396\\
38000	15899\\
40000	15970\\
42000	15868\\
44000	15989\\
46000	16228\\
48000	16353\\
50000	15804\\
52000	15866\\
54000	16218\\
56000	15392\\
58000	15863\\
60000	16140\\
62000	16303\\
64000	15808\\
68000	16143\\
70000	15727\\
72000	16271\\
74000	16052\\
76000	16242\\
80000	16256\\
82000	16105\\
84000	16495\\
86000	15813\\
88000	15654\\
90000	15282\\
92000	16097\\
94000	15708\\
96000	15939\\
98000	15957\\
1e+05	16100\\
};
\addplot [color=mycolor18, forget plot]
  table[row sep=crcr]{%
0	4395.3\\
2000	18488\\
4000	24692\\
8000	24435\\
10000	24538\\
12000	24080\\
14000	24433\\
16000	23665\\
18000	24330\\
20000	24550\\
22000	24982\\
24000	24581\\
26000	24884\\
28000	23898\\
30000	23985\\
32000	24339\\
34000	24114\\
36000	23216\\
38000	23750\\
40000	24442\\
42000	24125\\
44000	24757\\
46000	24121\\
48000	24196\\
50000	24025\\
52000	23573\\
56000	24195\\
58000	24129\\
60000	24303\\
62000	24158\\
64000	24564\\
68000	23743\\
70000	24197\\
72000	24281\\
74000	23437\\
76000	23739\\
78000	24619\\
80000	23970\\
82000	25131\\
84000	24872\\
86000	23963\\
88000	23993\\
90000	23191\\
92000	23359\\
94000	24108\\
96000	24467\\
98000	23203\\
1e+05	24037\\
};
\addplot [color=mycolor19, forget plot]
  table[row sep=crcr]{%
0	4395.3\\
2000	12378\\
4000	14795\\
6000	14856\\
8000	14499\\
10000	15106\\
12000	14845\\
14000	14915\\
18000	14431\\
20000	15033\\
22000	14952\\
24000	14735\\
26000	14848\\
28000	14649\\
30000	14831\\
32000	14603\\
34000	14858\\
36000	14436\\
38000	14679\\
40000	14818\\
42000	14444\\
44000	14851\\
46000	14606\\
48000	14781\\
54000	14837\\
56000	14684\\
58000	15152\\
60000	15061\\
62000	14881\\
64000	14634\\
68000	14443\\
70000	14675\\
72000	14637\\
74000	14685\\
76000	14645\\
78000	14665\\
80000	14811\\
82000	14784\\
84000	14609\\
86000	15230\\
88000	15178\\
90000	14764\\
92000	15062\\
94000	14654\\
96000	14173\\
98000	14312\\
1e+05	15362\\
};
\addplot [color=mycolor20, forget plot]
  table[row sep=crcr]{%
0	4395.3\\
2000	15163\\
4000	18098\\
6000	18463\\
8000	18900\\
10000	18899\\
12000	18551\\
14000	18617\\
16000	18870\\
18000	19044\\
20000	18846\\
22000	19176\\
24000	19257\\
26000	19265\\
28000	18935\\
30000	19311\\
32000	19126\\
34000	19020\\
36000	19804\\
38000	18946\\
40000	18487\\
42000	19712\\
44000	18969\\
46000	18616\\
48000	19091\\
50000	19096\\
52000	19415\\
56000	18933\\
58000	18753\\
60000	19545\\
62000	18800\\
64000	19546\\
68000	19022\\
70000	19153\\
72000	19069\\
74000	18743\\
76000	19309\\
78000	19531\\
80000	18347\\
82000	19380\\
84000	18986\\
86000	19299\\
88000	19114\\
90000	19237\\
92000	19120\\
94000	19309\\
96000	19305\\
98000	19715\\
1e+05	19002\\
};
\addplot [color=mycolor21, forget plot]
  table[row sep=crcr]{%
0	4395.3\\
2000	8330.8\\
4000	9755.8\\
6000	9683.8\\
8000	9909.4\\
12000	9569.5\\
14000	9778.8\\
16000	9751.3\\
18000	9602\\
20000	9904.1\\
22000	9763.5\\
24000	9472.6\\
26000	9484.2\\
28000	9753.1\\
30000	9623.8\\
32000	9777\\
36000	9616.3\\
38000	9645.8\\
40000	9844.6\\
42000	9549.1\\
44000	9565.6\\
46000	9692.6\\
48000	9689.5\\
50000	9616.1\\
52000	9673\\
54000	9625.9\\
56000	9744.2\\
58000	9394.5\\
60000	9513.7\\
62000	9490.3\\
64000	9672.4\\
66000	9637.7\\
68000	9695.5\\
70000	9574.3\\
72000	9739.6\\
74000	9468.3\\
76000	9626.5\\
78000	9510.2\\
80000	9569\\
84000	9591.4\\
86000	9794.8\\
88000	9746.9\\
90000	9770.7\\
92000	9734.6\\
94000	9611.7\\
98000	9716.5\\
1e+05	9583\\
};
\addplot [color=mycolor22, forget plot]
  table[row sep=crcr]{%
0	4395.3\\
2000	21781\\
4000	26239\\
6000	26291\\
10000	25670\\
12000	26220\\
14000	25319\\
16000	25135\\
18000	25504\\
20000	24166\\
22000	24585\\
24000	25095\\
26000	25225\\
28000	25172\\
30000	25719\\
32000	25536\\
34000	25600\\
36000	25999\\
38000	25029\\
40000	24555\\
42000	25067\\
44000	24977\\
46000	24068\\
48000	24803\\
50000	25434\\
52000	25319\\
54000	25014\\
56000	25097\\
58000	24928\\
60000	24392\\
62000	23549\\
64000	23978\\
66000	23892\\
68000	25280\\
70000	25297\\
72000	24587\\
74000	24432\\
76000	25612\\
78000	25332\\
80000	25239\\
82000	24427\\
84000	24793\\
86000	25093\\
88000	24514\\
90000	25336\\
92000	24904\\
94000	24919\\
98000	24616\\
1e+05	25033\\
};
\addplot [color=mycolor23, forget plot]
  table[row sep=crcr]{%
0	4395.3\\
2000	10887\\
4000	12471\\
6000	12649\\
10000	12204\\
12000	12217\\
14000	12287\\
16000	12225\\
18000	12355\\
22000	12012\\
24000	12313\\
26000	12245\\
28000	12301\\
30000	12136\\
32000	12226\\
34000	12366\\
36000	11973\\
38000	12041\\
40000	12279\\
42000	12265\\
44000	12418\\
46000	12381\\
48000	12508\\
50000	12468\\
52000	12066\\
54000	12078\\
60000	12322\\
62000	12046\\
64000	12287\\
66000	12467\\
68000	12217\\
70000	12157\\
72000	12054\\
74000	12195\\
76000	12415\\
78000	12151\\
80000	12047\\
82000	12149\\
84000	11955\\
86000	12036\\
90000	12576\\
92000	12348\\
94000	12217\\
96000	12265\\
98000	12375\\
1e+05	12248\\
};
\addplot [color=mycolor24, forget plot]
  table[row sep=crcr]{%
0	4395.3\\
2000	5526.5\\
4000	4904.5\\
6000	4690.3\\
8000	4738.4\\
10000	4860.1\\
12000	4828.4\\
14000	4860.9\\
16000	4933.9\\
18000	5056\\
20000	4801.1\\
22000	4778.1\\
24000	4707.9\\
26000	4814.1\\
28000	4890.7\\
30000	4807.7\\
32000	4812.4\\
36000	4851.9\\
38000	4923.4\\
40000	4825.2\\
42000	4929.6\\
44000	4812.1\\
46000	4931.6\\
48000	4746.1\\
50000	4996.5\\
52000	5015.9\\
54000	4886.4\\
56000	4854.8\\
58000	4809.2\\
60000	4898.3\\
62000	4738.4\\
64000	4682.9\\
66000	4708.3\\
68000	4802.4\\
70000	4973.2\\
72000	4767.9\\
74000	4877.5\\
76000	4661.1\\
78000	4813.2\\
80000	4860.5\\
82000	4646.5\\
84000	4745.4\\
86000	4754.4\\
88000	4608.6\\
90000	4868.6\\
92000	4863.6\\
94000	4807.4\\
96000	4884.4\\
98000	4717.4\\
1e+05	4766.7\\
};
\addplot [color=mycolor25, forget plot]
  table[row sep=crcr]{%
0	4395.3\\
2000	10710\\
4000	12175\\
6000	12209\\
8000	12097\\
10000	12395\\
12000	12180\\
14000	12052\\
16000	12193\\
22000	12065\\
24000	12404\\
26000	12140\\
28000	12324\\
30000	12324\\
32000	12650\\
34000	12038\\
36000	11940\\
38000	11982\\
40000	12382\\
42000	12103\\
44000	12470\\
46000	12494\\
48000	12339\\
50000	12426\\
52000	12806\\
54000	12273\\
56000	12641\\
58000	12786\\
60000	12718\\
62000	12465\\
64000	12177\\
66000	12635\\
68000	12186\\
70000	12523\\
72000	12520\\
74000	12827\\
76000	12445\\
78000	12746\\
80000	12434\\
84000	11982\\
86000	12877\\
88000	12532\\
90000	12157\\
92000	12360\\
94000	12471\\
96000	12408\\
98000	12539\\
1e+05	12468\\
};
\addplot [color=mycolor26, forget plot]
  table[row sep=crcr]{%
0	4395.3\\
2000	7429.1\\
4000	8022.3\\
6000	8113.5\\
8000	8307.5\\
12000	8134\\
14000	8161.8\\
16000	8516.5\\
18000	8167.4\\
20000	8389.2\\
22000	8316.4\\
24000	8182.8\\
26000	8354.2\\
28000	8493\\
30000	8231.1\\
32000	8136.6\\
34000	8173.7\\
36000	8339.3\\
38000	8348.6\\
40000	8395.6\\
42000	8124.6\\
44000	8275.1\\
46000	8258.1\\
48000	8602.6\\
50000	8429.7\\
52000	8568.8\\
54000	8325.9\\
56000	8124.4\\
58000	8607.3\\
60000	8205.3\\
62000	8192\\
64000	8029\\
66000	8388.1\\
68000	8510.6\\
72000	8699.1\\
74000	8313.9\\
76000	8167.3\\
78000	8205.1\\
80000	8080.6\\
82000	8206.6\\
84000	8282\\
86000	8275.8\\
88000	8358.2\\
90000	8025.6\\
94000	8254.4\\
96000	8353.6\\
98000	8228.3\\
1e+05	8352.6\\
};
\addplot [color=mycolor27, forget plot]
  table[row sep=crcr]{%
0	4395.3\\
2000	9509.9\\
4000	9515.8\\
6000	9924.2\\
8000	9984.1\\
10000	9603.8\\
12000	9934.6\\
14000	9553.2\\
16000	9843.1\\
18000	9857.7\\
20000	9598.9\\
22000	9723.1\\
24000	9516.1\\
26000	9649.6\\
28000	9409.7\\
30000	9888\\
32000	9666.9\\
34000	9960.9\\
36000	10071\\
38000	9574.5\\
40000	9743.3\\
42000	9751.6\\
44000	9511.7\\
46000	9427.9\\
48000	9601\\
50000	9549.9\\
52000	9959.8\\
54000	9978.4\\
56000	9760.2\\
58000	9452.3\\
60000	9488.4\\
62000	9667.5\\
64000	9779\\
66000	9728.1\\
68000	9752.7\\
70000	9697.4\\
72000	10103\\
74000	9659\\
76000	9986.1\\
78000	9457.7\\
80000	9998.3\\
82000	9967.2\\
84000	9615\\
86000	9784.6\\
88000	10009\\
90000	9850.2\\
92000	9983\\
94000	9750.9\\
96000	9861.2\\
98000	9748.9\\
1e+05	9767.4\\
};
\addplot [color=mycolor28, forget plot]
  table[row sep=crcr]{%
0	4395.3\\
2000	6256.2\\
4000	5678\\
6000	5617.4\\
8000	5635.8\\
10000	5556.4\\
12000	5359.8\\
14000	5503.4\\
16000	5355\\
18000	5689.1\\
20000	5321.9\\
22000	5440.2\\
24000	5492\\
26000	5633\\
30000	5523.9\\
32000	5661\\
34000	5436.1\\
36000	5535.9\\
38000	5422.5\\
40000	5590.9\\
42000	5500.2\\
44000	5538.9\\
46000	5456.9\\
48000	5327.8\\
50000	5221.5\\
52000	5514.2\\
54000	5420.4\\
56000	5432.2\\
58000	5641.5\\
60000	5469.5\\
62000	5513\\
64000	5452.6\\
66000	5533\\
68000	5571.3\\
70000	5352.8\\
72000	5475.4\\
74000	5351.8\\
76000	5481.9\\
78000	5647.6\\
80000	5433.3\\
84000	5309.8\\
86000	5520.9\\
88000	5594.2\\
90000	5304.1\\
92000	5476.3\\
94000	5484.8\\
96000	5388.4\\
98000	5342.6\\
1e+05	5419.2\\
};
\addplot [color=mycolor29, forget plot]
  table[row sep=crcr]{%
0	4395.3\\
2000	11020\\
4000	11509\\
6000	11662\\
8000	11417\\
10000	11559\\
12000	11315\\
14000	11437\\
16000	11364\\
18000	11417\\
20000	11006\\
24000	11379\\
26000	11498\\
28000	11895\\
30000	11511\\
32000	11752\\
34000	11356\\
36000	11453\\
38000	11760\\
42000	11624\\
44000	11633\\
46000	11831\\
48000	11927\\
50000	11584\\
52000	11571\\
54000	11772\\
56000	11495\\
58000	11146\\
60000	11704\\
62000	12065\\
64000	11290\\
66000	11702\\
68000	11263\\
70000	11488\\
72000	11509\\
74000	11671\\
76000	11505\\
78000	11176\\
80000	10995\\
82000	11531\\
84000	12004\\
86000	11479\\
88000	11632\\
90000	11724\\
92000	11390\\
94000	11259\\
96000	11727\\
98000	11460\\
1e+05	11509\\
};
\addplot [color=mycolor30, forget plot]
  table[row sep=crcr]{%
0	4395.3\\
2000	6968.9\\
4000	7736.2\\
6000	7923\\
8000	7900.5\\
10000	7773.7\\
12000	8052.6\\
14000	7806.4\\
16000	7668.6\\
18000	7697.5\\
20000	7949.2\\
22000	8042.3\\
24000	7803.3\\
26000	7856.6\\
28000	7828.8\\
30000	7755.3\\
34000	7970.9\\
36000	7875.2\\
38000	7941.4\\
40000	7834.3\\
46000	8003.7\\
48000	7835.7\\
50000	7870.6\\
52000	7788.1\\
54000	8034\\
56000	7804\\
58000	8033.4\\
62000	7826.3\\
64000	7962.7\\
66000	7780\\
70000	7960.6\\
74000	7935.2\\
76000	8015\\
78000	7986.6\\
80000	7804.2\\
82000	7741.7\\
84000	8016.6\\
86000	7678\\
88000	7744.3\\
90000	7756.3\\
92000	7793.4\\
94000	7752.8\\
96000	8002.9\\
98000	7906.2\\
1e+05	7889.8\\
};
\addplot [color=mycolor31, forget plot]
  table[row sep=crcr]{%
0	4395.3\\
2000	7008.5\\
4000	6774.8\\
6000	6262.1\\
8000	6346\\
10000	6375.4\\
12000	6271.4\\
14000	6470.6\\
16000	6321\\
18000	6460.9\\
20000	6285.8\\
22000	6243.7\\
24000	6526.4\\
26000	6435.1\\
28000	6476.2\\
30000	6402.9\\
32000	6353.8\\
34000	6495.7\\
36000	6341.2\\
38000	6263.3\\
40000	6534.4\\
42000	6446.4\\
44000	6527.2\\
46000	6436.8\\
50000	6458.5\\
52000	6332.4\\
54000	6360.2\\
56000	6542.9\\
58000	6682.2\\
60000	6546.3\\
62000	6377.9\\
64000	6582.3\\
66000	6570.9\\
68000	6378.4\\
70000	6818.4\\
72000	6632.1\\
76000	6369.7\\
78000	6379.3\\
80000	6504.1\\
82000	6400.6\\
84000	6507.9\\
86000	6559.4\\
88000	6312.1\\
90000	6713.4\\
92000	6663.2\\
94000	6524.4\\
96000	6586.2\\
98000	6268.7\\
1e+05	6578.1\\
};
\addplot [color=mycolor32, forget plot]
  table[row sep=crcr]{%
0	4395.3\\
2000	10345\\
4000	11184\\
6000	10630\\
8000	10759\\
10000	10711\\
12000	10980\\
14000	10376\\
16000	10466\\
18000	10620\\
20000	10161\\
22000	10852\\
24000	10743\\
26000	10864\\
28000	10480\\
30000	10281\\
32000	10485\\
34000	10468\\
36000	10693\\
38000	10608\\
42000	10240\\
44000	10724\\
46000	10636\\
48000	10266\\
50000	10086\\
52000	10567\\
54000	10672\\
56000	10368\\
58000	10325\\
60000	10564\\
62000	10273\\
64000	10592\\
66000	10206\\
70000	10403\\
72000	10338\\
74000	10402\\
76000	10574\\
78000	10528\\
80000	10362\\
82000	10457\\
84000	10355\\
86000	10593\\
88000	10541\\
90000	10398\\
92000	10446\\
94000	10375\\
96000	10598\\
98000	10670\\
1e+05	10149\\
};
\addplot [color=mycolor33, forget plot]
  table[row sep=crcr]{%
0	4395.3\\
2000	6887.6\\
4000	7934.1\\
6000	7942.6\\
8000	7881.6\\
10000	8294.5\\
12000	8255\\
14000	8626.3\\
16000	7920.1\\
18000	8151.1\\
20000	8085\\
22000	8390.9\\
24000	8317.3\\
26000	8012.8\\
28000	8075.8\\
30000	8263.1\\
32000	8376.2\\
34000	7823.4\\
36000	8015.7\\
38000	8256.7\\
40000	7910\\
42000	8247.7\\
44000	8302.8\\
46000	8121\\
50000	8367\\
52000	8242\\
54000	8276.6\\
56000	8102.9\\
58000	8483.5\\
60000	8123.5\\
62000	8226.4\\
64000	8114.2\\
66000	8032.6\\
68000	8175.3\\
70000	8101.2\\
72000	8295.8\\
74000	8215.9\\
76000	8057.2\\
78000	8197.8\\
80000	8393.1\\
82000	7865.5\\
84000	8166.3\\
86000	8428.8\\
88000	8189.9\\
90000	8104.2\\
92000	8052.7\\
94000	8254.6\\
96000	8241.5\\
98000	8129.4\\
1e+05	8159.8\\
};
\addplot [color=mycolor34, forget plot]
  table[row sep=crcr]{%
0	4395.3\\
2000	4866.4\\
4000	4038.3\\
6000	3930.5\\
8000	4057.1\\
10000	4120.5\\
12000	4046.4\\
16000	4044.3\\
18000	3951.6\\
20000	4131.6\\
22000	4208.2\\
26000	4028.3\\
30000	4153.1\\
32000	4077.4\\
34000	4104.7\\
36000	3993.3\\
38000	4002.9\\
40000	4144.2\\
42000	4142.2\\
44000	4099.4\\
46000	3994.1\\
48000	4145.1\\
50000	4002.9\\
52000	4200.7\\
54000	4003.1\\
56000	4174.2\\
58000	3941.7\\
60000	4078.7\\
62000	4056.6\\
64000	4093.9\\
66000	3991.3\\
68000	4032.9\\
70000	4265.2\\
72000	4092.9\\
74000	4090.3\\
76000	3948.2\\
78000	3911.2\\
80000	4167.3\\
82000	3980.7\\
84000	4004.7\\
86000	3972.9\\
88000	4032.5\\
90000	4073.9\\
92000	4101.7\\
94000	4071.7\\
96000	4081.9\\
98000	4057.6\\
1e+05	3922.5\\
};
\addplot [color=mycolor35]
  table[row sep=crcr]{%
0	4395.3\\
2000	9096.1\\
4000	9866.2\\
8000	10261\\
12000	10087\\
14000	10206\\
16000	10116\\
20000	10137\\
22000	10222\\
24000	10400\\
26000	10073\\
28000	10193\\
32000	10087\\
34000	10240\\
36000	10296\\
38000	10065\\
40000	10258\\
44000	9898\\
46000	10105\\
48000	10112\\
50000	10289\\
52000	10185\\
54000	10255\\
58000	10327\\
60000	9995.6\\
62000	10221\\
64000	9962.3\\
66000	10034\\
68000	10151\\
70000	10104\\
72000	10208\\
74000	9800.3\\
76000	9996.1\\
78000	9962.8\\
82000	10319\\
84000	10347\\
88000	10248\\
90000	10163\\
92000	10373\\
94000	10222\\
98000	10224\\
1e+05	10336\\
};
\addlegendentry{Classe k=83}

\addplot [color=mycolor36, forget plot]
  table[row sep=crcr]{%
0	4395.3\\
2000	5047.5\\
4000	4455.1\\
6000	4365\\
8000	4181.5\\
10000	4341.6\\
12000	4098.7\\
14000	4127.7\\
16000	4203\\
18000	4333.1\\
20000	4220.3\\
22000	4291.1\\
24000	4254.2\\
26000	4435.8\\
28000	4312.7\\
30000	4422.8\\
32000	4473.9\\
34000	4143.4\\
36000	4298.1\\
38000	4315.8\\
40000	4594.8\\
42000	4451\\
44000	4453.6\\
46000	4227.7\\
48000	4388\\
50000	4341.7\\
52000	4229.4\\
54000	4330\\
56000	4456.5\\
58000	4310.6\\
60000	4214.2\\
64000	4276.3\\
66000	4329.5\\
68000	4370.9\\
70000	4325.5\\
72000	4187.2\\
74000	4346.9\\
76000	4262.5\\
78000	4194.9\\
80000	4424.2\\
82000	4362.4\\
84000	4239.6\\
86000	4199.3\\
88000	4346\\
90000	4285.6\\
92000	4334.7\\
94000	4216.9\\
96000	4315.8\\
98000	4326.4\\
1e+05	4288.8\\
};
\addplot [color=mycolor37, forget plot]
  table[row sep=crcr]{%
0	4395.3\\
2000	4089.2\\
4000	3689.2\\
6000	3604.9\\
8000	3667.5\\
10000	3557.1\\
12000	3574.7\\
14000	3648.5\\
16000	3817.3\\
18000	3685.6\\
20000	3628.1\\
22000	3714.3\\
24000	3666.4\\
26000	3690.5\\
28000	3579.3\\
30000	3674.5\\
32000	3677.2\\
34000	3620.8\\
38000	3618.1\\
40000	3528.4\\
42000	3677.9\\
44000	3610\\
46000	3578.5\\
48000	3598.3\\
50000	3735.3\\
52000	3615.4\\
54000	3607.5\\
56000	3643.5\\
58000	3693.4\\
60000	3584.2\\
62000	3697.1\\
64000	3595.3\\
66000	3648.1\\
68000	3614.3\\
70000	3658.6\\
72000	3765.2\\
74000	3674.6\\
76000	3759.3\\
78000	3632.2\\
80000	3681.4\\
82000	3657.4\\
86000	3639.5\\
88000	3680.4\\
90000	3670.9\\
92000	3631.5\\
94000	3671.2\\
96000	3614.7\\
1e+05	3676.9\\
};
\addplot [color=mycolor38, forget plot]
  table[row sep=crcr]{%
0	4395.3\\
2000	3427.6\\
4000	3172.3\\
6000	3008.2\\
8000	3125.1\\
10000	3131.6\\
12000	2968.4\\
14000	3037.5\\
16000	3000.5\\
18000	3106.6\\
20000	3048.6\\
22000	3009.8\\
24000	2995.4\\
26000	3109\\
28000	3018.7\\
30000	3138\\
32000	3054.6\\
36000	3013.1\\
38000	3064.3\\
40000	2976.4\\
42000	3019.7\\
44000	3108.1\\
46000	2992.1\\
50000	2908\\
52000	3005.8\\
54000	3001.8\\
56000	2984.2\\
60000	3079.2\\
62000	2976.5\\
64000	2894.6\\
66000	3048.4\\
68000	3047.2\\
70000	3082.8\\
72000	3037.7\\
74000	2941.5\\
76000	3023.9\\
78000	3002.1\\
80000	3041.4\\
82000	3046\\
84000	3030.8\\
86000	3073.4\\
88000	3129.1\\
90000	3129.4\\
92000	2951.3\\
94000	2941.6\\
96000	3025.3\\
98000	3016.4\\
1e+05	3097.7\\
};
\addplot [color=mycolor39, forget plot]
  table[row sep=crcr]{%
0	4395.3\\
2000	11633\\
4000	12676\\
6000	12278\\
10000	11656\\
16000	11892\\
18000	11830\\
20000	11831\\
22000	11899\\
24000	12233\\
26000	11938\\
28000	11709\\
30000	12010\\
34000	11787\\
36000	11878\\
40000	11688\\
42000	11678\\
44000	12035\\
46000	11752\\
48000	11593\\
50000	11298\\
52000	11551\\
54000	11534\\
56000	12013\\
58000	11648\\
60000	11480\\
62000	11548\\
64000	11769\\
66000	11713\\
68000	11400\\
70000	11498\\
72000	11768\\
74000	12083\\
76000	11876\\
78000	11807\\
80000	12157\\
82000	11727\\
84000	11610\\
86000	12047\\
88000	11948\\
90000	12090\\
92000	12089\\
94000	11875\\
96000	11774\\
98000	12013\\
1e+05	11612\\
};
\addplot [color=mycolor40, forget plot]
  table[row sep=crcr]{%
0	4395.3\\
2000	5184.7\\
4000	5655.9\\
6000	5589.1\\
8000	5677.6\\
10000	5609\\
14000	5918.2\\
16000	5842.2\\
18000	5648.8\\
20000	5851.4\\
22000	5763.1\\
24000	5844.8\\
26000	5734.8\\
28000	5739.3\\
30000	5776.7\\
32000	5861.3\\
34000	5770.9\\
36000	5771\\
38000	5752.9\\
40000	5681.2\\
42000	5850.8\\
44000	5708.3\\
46000	5737.3\\
48000	5664.5\\
50000	5933.1\\
52000	5812.6\\
54000	5619.3\\
56000	5704.7\\
58000	5663.4\\
62000	5702.9\\
64000	5828.2\\
66000	5654.6\\
70000	5761.1\\
74000	5686.3\\
76000	5680.3\\
80000	5747.7\\
82000	5729.9\\
84000	5796.1\\
86000	5767.7\\
88000	5768.4\\
90000	5734\\
92000	5816.2\\
94000	5794.1\\
96000	5741.4\\
98000	5810.7\\
1e+05	5810.9\\
};
\addplot [color=mycolor41, forget plot]
  table[row sep=crcr]{%
0	4395.3\\
2000	4898.5\\
4000	5188.9\\
6000	5224.4\\
8000	5193.7\\
10000	5230.7\\
12000	5344.5\\
14000	5437.2\\
16000	5435.4\\
20000	5299\\
22000	5438.4\\
24000	5350.3\\
26000	5569.1\\
28000	5415.3\\
30000	5526.6\\
32000	5248\\
34000	5367.7\\
36000	5318\\
38000	5340.5\\
40000	5384.2\\
42000	5288.1\\
44000	5494.1\\
46000	5316.9\\
48000	5391.5\\
50000	5310.5\\
52000	5267.2\\
54000	5477.2\\
56000	5308.4\\
58000	5606.6\\
60000	5387.8\\
62000	5334.1\\
64000	5425.3\\
68000	5239.2\\
70000	5579.9\\
72000	5481.9\\
74000	5303.1\\
76000	5317\\
78000	5550.8\\
80000	5374.6\\
82000	5412.9\\
84000	5416.8\\
86000	5354.4\\
88000	5356.2\\
90000	5438.9\\
92000	5366.4\\
94000	5240\\
96000	5377\\
98000	5294.6\\
1e+05	5455.3\\
};
\addplot [color=mycolor42, forget plot]
  table[row sep=crcr]{%
0	4395.3\\
2000	14373\\
4000	18951\\
6000	19658\\
8000	20767\\
10000	19913\\
12000	21007\\
14000	20871\\
18000	20385\\
20000	20963\\
22000	20882\\
24000	20363\\
26000	20041\\
28000	20494\\
30000	20318\\
32000	20624\\
34000	20523\\
38000	20635\\
40000	20944\\
44000	20917\\
46000	19985\\
48000	20359\\
50000	20654\\
52000	20039\\
54000	21024\\
56000	20141\\
58000	20610\\
60000	20427\\
62000	20340\\
64000	20424\\
66000	19866\\
68000	20694\\
70000	20397\\
72000	20649\\
74000	20699\\
76000	20514\\
78000	20203\\
80000	20643\\
82000	20328\\
84000	20236\\
86000	21154\\
88000	20302\\
90000	20219\\
92000	20597\\
94000	20464\\
96000	20208\\
98000	21031\\
1e+05	20312\\
};
\addplot [color=mycolor43, forget plot]
  table[row sep=crcr]{%
0	4395.3\\
2000	5446.8\\
4000	5192\\
6000	5029.6\\
8000	5075.7\\
10000	5063.6\\
12000	5036.8\\
14000	4910.6\\
16000	4867\\
18000	5079.2\\
20000	5084.3\\
22000	4889.1\\
24000	4907.1\\
26000	4963\\
28000	4802.1\\
30000	4882.3\\
32000	4847.1\\
34000	5066.9\\
36000	4914.3\\
38000	4949.8\\
40000	4929.2\\
42000	4947.7\\
44000	4857.7\\
46000	4886.3\\
48000	4890.3\\
50000	4994.2\\
52000	4916.6\\
54000	5005.1\\
56000	5030.4\\
58000	5017.7\\
60000	4814.4\\
64000	4899.5\\
66000	4907.6\\
68000	5032.8\\
70000	4894.8\\
72000	4908.7\\
74000	5030.2\\
76000	4949.3\\
78000	4825.3\\
80000	4833\\
82000	4942.8\\
84000	4959.8\\
86000	5031.9\\
88000	4950.9\\
90000	4913.6\\
92000	5067.9\\
94000	4972\\
96000	5052\\
98000	4987.4\\
1e+05	4954.2\\
};
\addplot [color=mycolor44, forget plot]
  table[row sep=crcr]{%
0	4395.3\\
2000	4133.6\\
4000	3589.6\\
6000	3509.6\\
8000	3457.4\\
10000	3427.3\\
12000	3411.4\\
14000	3322\\
16000	3436.4\\
18000	3407.3\\
20000	3524.1\\
22000	3387.8\\
24000	3329.5\\
26000	3429.4\\
28000	3320.6\\
30000	3379.9\\
32000	3310.3\\
34000	3292.4\\
36000	3388.5\\
38000	3409.7\\
40000	3407.2\\
42000	3225.7\\
44000	3357.4\\
46000	3478.7\\
48000	3382\\
50000	3191.1\\
52000	3179.4\\
54000	3382.2\\
56000	3199.5\\
58000	3369\\
60000	3200.4\\
62000	3401.2\\
64000	3341.4\\
66000	3353\\
68000	3396.9\\
70000	3321.6\\
72000	3362.9\\
74000	3289.9\\
76000	3293.1\\
78000	3426\\
80000	3278.9\\
82000	3327\\
84000	3395.6\\
88000	3512.4\\
90000	3417\\
92000	3349\\
94000	3434.6\\
96000	3353.4\\
1e+05	3391.1\\
};
\addplot [color=mycolor45, forget plot]
  table[row sep=crcr]{%
0	4395.3\\
2000	5672.7\\
4000	6696.3\\
6000	6967.2\\
8000	7011.5\\
10000	6941.8\\
12000	7011.5\\
14000	6919.3\\
16000	7025.2\\
18000	6768.7\\
20000	6929.9\\
22000	6720.4\\
24000	6916.5\\
26000	6818.4\\
28000	7136.9\\
30000	6788.2\\
32000	6853.4\\
34000	6860.4\\
36000	6977.7\\
38000	7062.2\\
40000	6974.6\\
42000	6807.4\\
44000	6848\\
46000	6953.3\\
48000	6998.5\\
50000	6934.2\\
52000	7110.9\\
54000	6873.9\\
56000	6867.6\\
58000	6890.1\\
60000	6814.7\\
64000	7051.7\\
66000	6932.5\\
68000	6967\\
70000	6771.1\\
72000	7112.4\\
74000	7059\\
76000	6893.4\\
78000	6876.7\\
80000	6791.1\\
82000	7098.5\\
84000	6977.7\\
86000	7091.1\\
88000	7002.8\\
90000	6874.4\\
92000	6826.2\\
94000	6955\\
96000	6994\\
98000	6872\\
1e+05	6725.6\\
};
\addplot [color=mycolor46, forget plot]
  table[row sep=crcr]{%
0	4395.3\\
2000	4056.6\\
4000	3985.4\\
6000	3998.5\\
8000	4153.8\\
10000	4006.6\\
12000	4146\\
14000	4169.5\\
16000	4148\\
18000	4244.4\\
20000	4021\\
22000	4100.3\\
24000	4041.1\\
26000	4108\\
28000	4262.7\\
30000	4245.2\\
32000	4242.1\\
34000	4141.3\\
38000	4035.6\\
42000	4123.2\\
44000	3904.3\\
46000	4178.5\\
48000	4101\\
50000	4084.3\\
52000	4026.9\\
54000	4216.4\\
56000	4182.7\\
58000	4239.1\\
60000	4089.7\\
62000	4089.6\\
64000	4150.5\\
66000	4039.9\\
68000	4114.6\\
70000	4132\\
72000	4008.1\\
74000	4069.7\\
78000	4034.2\\
80000	4156.1\\
82000	4186.9\\
84000	4048.7\\
86000	4050.6\\
90000	4027\\
92000	4054.5\\
94000	4160.8\\
96000	4080.7\\
98000	4136.6\\
1e+05	4087.4\\
};
\addplot [color=mycolor47, forget plot]
  table[row sep=crcr]{%
0	4395.3\\
2000	5509.7\\
4000	5366.4\\
6000	5330.8\\
8000	5368.8\\
10000	5241.4\\
12000	5247.3\\
14000	5373.4\\
16000	5286\\
18000	5394.1\\
20000	5228.5\\
22000	5273.4\\
24000	5376.6\\
26000	5347.7\\
28000	5211.6\\
30000	5310.2\\
32000	5351.5\\
34000	5287.8\\
38000	5296.5\\
40000	5082.2\\
42000	5182.2\\
44000	5395.7\\
46000	5427.3\\
48000	5432.7\\
50000	5399.5\\
52000	5141.4\\
54000	5281.4\\
56000	5356.3\\
58000	5392.9\\
60000	5203.3\\
62000	5289.1\\
64000	5230.1\\
66000	5360.7\\
68000	5377.2\\
70000	5320.9\\
72000	5227.9\\
74000	5224.1\\
76000	5281.6\\
78000	5370.3\\
80000	5329.6\\
82000	5371.7\\
84000	5336.5\\
88000	5387.8\\
90000	5211.3\\
92000	5370.1\\
94000	5271.6\\
96000	5208.6\\
98000	5248.9\\
1e+05	5251\\
};
\addplot [color=mycolor48, forget plot]
  table[row sep=crcr]{%
0	4395.3\\
2000	3417.5\\
4000	3000.6\\
8000	2944.1\\
12000	2987.1\\
14000	2890.3\\
16000	2916.4\\
18000	2955.3\\
20000	2945\\
22000	2917.1\\
24000	2955.9\\
26000	2937.7\\
28000	2896.4\\
30000	2950.1\\
32000	2947.6\\
34000	2904.2\\
36000	2942.4\\
38000	2954.5\\
40000	2945.7\\
42000	2982.9\\
44000	2972.5\\
46000	2917.6\\
48000	2953.1\\
50000	2954.4\\
52000	2931.1\\
60000	2930.5\\
64000	2925.4\\
66000	2964.4\\
68000	2978.3\\
70000	2922\\
72000	2937.7\\
76000	2914.3\\
78000	2946.6\\
80000	2953.4\\
82000	2972.5\\
84000	2922.5\\
86000	2959.5\\
88000	2939.4\\
90000	2984.9\\
92000	2913.4\\
94000	2900.4\\
1e+05	2968.4\\
};
\addplot [color=mycolor49, forget plot]
  table[row sep=crcr]{%
0	4395.3\\
2000	4900.7\\
4000	4475.4\\
6000	4306.4\\
8000	4319.8\\
10000	4498.3\\
12000	4357.5\\
14000	4369.1\\
16000	4363.1\\
18000	4504.8\\
20000	4369.6\\
22000	4401.1\\
24000	4344.8\\
26000	4318.4\\
28000	4312\\
30000	4187.7\\
32000	4405.8\\
34000	4451.5\\
36000	4205.1\\
40000	4414.8\\
42000	4225\\
44000	4229.6\\
46000	4409.1\\
48000	4436.2\\
50000	4241.7\\
52000	4441.8\\
54000	4530.3\\
56000	4298.8\\
58000	4446.7\\
60000	4380.4\\
62000	4521.2\\
64000	4370.8\\
66000	4454.7\\
68000	4279.6\\
70000	4340.1\\
72000	4228.1\\
74000	4407.9\\
76000	4326.5\\
78000	4373.4\\
80000	4494.2\\
82000	4419.6\\
84000	4470.4\\
86000	4302.7\\
88000	4345.3\\
90000	4406.8\\
92000	4287.9\\
94000	4319.7\\
96000	4392.5\\
98000	4419.9\\
1e+05	4326.4\\
};
\addplot [color=mycolor50, forget plot]
  table[row sep=crcr]{%
0	4395.3\\
2000	2603\\
4000	2481.6\\
6000	2411.1\\
12000	2494.1\\
14000	2481.8\\
16000	2483.2\\
18000	2399.9\\
20000	2482.5\\
22000	2549.9\\
24000	2561\\
26000	2440\\
28000	2431.1\\
30000	2472.7\\
32000	2468.4\\
36000	2443.7\\
38000	2522\\
40000	2403.8\\
42000	2467.6\\
44000	2506.3\\
46000	2444.8\\
50000	2528.4\\
52000	2521.5\\
54000	2482.2\\
56000	2542.7\\
58000	2478.1\\
60000	2463.7\\
62000	2486.3\\
64000	2450.6\\
66000	2525.2\\
68000	2544.5\\
70000	2476.5\\
72000	2439.2\\
74000	2475.6\\
76000	2441.5\\
78000	2475.2\\
80000	2431.1\\
82000	2498.2\\
86000	2520.7\\
88000	2479.2\\
90000	2459.8\\
92000	2490.1\\
94000	2534.7\\
96000	2501.6\\
98000	2482.9\\
1e+05	2404.2\\
};
\addplot [color=mycolor51, forget plot]
  table[row sep=crcr]{%
0	4395.3\\
2000	5451.9\\
4000	5811.3\\
6000	5775.6\\
8000	5771.7\\
10000	5853.1\\
12000	5877.8\\
14000	6024.9\\
16000	5735.5\\
18000	5975\\
20000	6002.3\\
22000	5979.9\\
24000	6041.1\\
28000	5991.2\\
30000	5945.9\\
32000	5871\\
34000	5868.6\\
36000	5992.3\\
38000	5787.3\\
40000	6058.2\\
42000	5875\\
44000	6124.2\\
46000	6054.3\\
48000	5823.8\\
50000	6047.4\\
52000	6038.8\\
54000	6056.1\\
56000	5670.2\\
58000	6053.5\\
60000	5937.1\\
62000	6351.3\\
64000	5939.9\\
66000	6002.8\\
68000	5953.8\\
70000	5827.8\\
72000	5889.3\\
74000	5777.1\\
76000	5947.1\\
78000	6032.5\\
80000	5854.1\\
82000	5966.9\\
84000	5854\\
86000	6019.2\\
88000	5891.2\\
90000	6013.6\\
92000	6107.7\\
94000	6121.9\\
96000	6107.6\\
98000	6131.7\\
1e+05	5970.6\\
};
\addplot [color=mycolor52, forget plot]
  table[row sep=crcr]{%
0	4395.3\\
2000	4061.7\\
4000	3718.7\\
6000	3636.7\\
8000	3588\\
10000	3502.7\\
12000	3540.4\\
14000	3477.5\\
16000	3468.7\\
18000	3431.9\\
20000	3502.6\\
22000	3475.4\\
24000	3507\\
26000	3589.5\\
28000	3572.8\\
30000	3527.7\\
32000	3573.9\\
34000	3699.9\\
36000	3501.3\\
38000	3561.6\\
42000	3671\\
44000	3435.3\\
46000	3584.3\\
48000	3404.8\\
50000	3578.6\\
52000	3525.6\\
54000	3409.2\\
56000	3437.9\\
58000	3481.4\\
60000	3556.9\\
62000	3540.1\\
64000	3502.9\\
66000	3594.3\\
68000	3504.9\\
70000	3436.3\\
72000	3556.4\\
74000	3547.3\\
76000	3667.8\\
78000	3435.9\\
80000	3476.7\\
82000	3553.2\\
84000	3542.4\\
86000	3609.1\\
88000	3584.5\\
90000	3574.1\\
92000	3354.2\\
94000	3468.4\\
96000	3572.9\\
98000	3655.2\\
1e+05	3565.5\\
};
\addplot [color=mycolor53, forget plot]
  table[row sep=crcr]{%
0	4395.3\\
2000	6101.8\\
4000	6333.5\\
6000	6063.2\\
8000	5730.2\\
10000	5669.4\\
14000	5776.1\\
16000	5777.5\\
18000	5547\\
20000	5673.5\\
22000	5763.4\\
24000	5963.6\\
26000	5765\\
28000	5719.6\\
32000	5737.5\\
34000	5619.3\\
36000	5805.3\\
38000	5709.2\\
40000	5672.6\\
42000	5874\\
44000	5494.9\\
46000	5829.9\\
48000	5773.5\\
50000	5795\\
52000	5661.2\\
54000	5586.6\\
56000	5553\\
58000	5805.8\\
60000	5790.2\\
62000	5554.4\\
64000	5501.8\\
66000	5556.4\\
68000	5795.1\\
70000	5494.3\\
72000	5324\\
74000	5527.7\\
76000	5799.3\\
78000	5776.6\\
80000	5728.3\\
82000	5504.2\\
84000	5775.2\\
86000	5738.8\\
88000	5868.8\\
90000	5628.5\\
92000	5722\\
94000	5544.6\\
96000	5759.1\\
98000	5869.7\\
1e+05	5896.1\\
};
\addplot [color=mycolor54, forget plot]
  table[row sep=crcr]{%
0	4395.3\\
2000	4703.4\\
4000	4629.4\\
6000	4588.4\\
8000	4488.3\\
10000	4549.9\\
12000	4515.9\\
14000	4504.2\\
18000	4522.7\\
20000	4638.4\\
22000	4613.4\\
24000	4539.6\\
26000	4580.5\\
28000	4521.6\\
30000	4648.2\\
32000	4625\\
34000	4555.8\\
36000	4660.5\\
38000	4492.8\\
40000	4620.4\\
42000	4685.7\\
44000	4543.7\\
46000	4575.2\\
48000	4474.4\\
50000	4569.2\\
52000	4542.2\\
54000	4546.3\\
56000	4591.7\\
58000	4516.8\\
60000	4570.9\\
62000	4557.3\\
64000	4646.8\\
66000	4704.6\\
68000	4778.7\\
70000	4684.5\\
72000	4571.9\\
74000	4616.4\\
76000	4736.2\\
78000	4570\\
80000	4539.2\\
82000	4693\\
84000	4602.6\\
88000	4522.4\\
90000	4658\\
92000	4428.3\\
94000	4434.5\\
96000	4505.3\\
98000	4550.4\\
1e+05	4516.5\\
};
\addplot [color=mycolor55, forget plot]
  table[row sep=crcr]{%
0	4395.3\\
2000	5193.5\\
4000	5634.1\\
6000	5748.6\\
8000	5637.3\\
10000	5921\\
12000	5892.2\\
16000	5746.4\\
18000	5649\\
20000	5866.9\\
22000	5732.1\\
24000	5824.6\\
26000	5668.8\\
28000	5787.5\\
32000	5696.1\\
34000	5690.8\\
36000	5720.5\\
38000	5665.5\\
40000	5665.5\\
42000	5875\\
44000	5787.7\\
48000	5729.8\\
52000	5799.9\\
56000	5693\\
58000	5849.5\\
60000	5828.8\\
62000	5884.9\\
64000	5830.3\\
68000	5693.9\\
70000	5776.2\\
72000	5796.3\\
74000	5781.9\\
76000	5817.1\\
78000	5895.6\\
80000	5802.5\\
82000	5738.1\\
84000	5753.5\\
86000	5837.4\\
88000	5719.2\\
90000	5801.2\\
92000	5790.1\\
94000	5642.6\\
96000	5655.7\\
98000	5788.3\\
1e+05	5750.4\\
};
\addplot [color=mycolor56, forget plot]
  table[row sep=crcr]{%
0	4395.3\\
2000	4248.5\\
4000	4118.9\\
6000	4077.1\\
8000	4065\\
10000	3998.3\\
12000	3911.7\\
14000	4108\\
24000	4015.1\\
26000	3998.6\\
30000	4052.6\\
32000	3979.3\\
34000	4046.9\\
38000	4023.3\\
40000	3920.8\\
42000	3970.7\\
44000	4064.6\\
46000	4048.8\\
48000	4129.3\\
50000	4032.8\\
52000	3950.6\\
54000	4069.3\\
56000	4022.8\\
58000	4068.4\\
60000	3971.9\\
62000	4007.4\\
64000	3948.9\\
66000	4029.2\\
68000	3968.8\\
70000	4018.8\\
72000	4113.1\\
74000	3986.5\\
76000	4028.8\\
78000	4014.5\\
80000	4092.6\\
82000	3933.4\\
84000	3934.6\\
86000	4024.3\\
88000	4016.3\\
90000	4032.7\\
92000	3993.4\\
94000	4101.3\\
96000	4061.1\\
98000	4040.8\\
1e+05	4079\\
};
\addplot [color=mycolor57, forget plot]
  table[row sep=crcr]{%
0	4395.3\\
2000	6092\\
4000	6766.6\\
6000	6594.5\\
8000	6528.9\\
10000	6700.4\\
12000	6453.8\\
14000	6639.1\\
16000	6716.8\\
18000	6644.6\\
20000	6625.8\\
22000	6672.6\\
24000	6819.8\\
26000	6768.2\\
28000	6872.6\\
30000	6704.4\\
32000	6677.9\\
34000	6701.8\\
36000	6680.9\\
38000	6583.7\\
40000	6653.2\\
42000	6801.9\\
44000	6663.5\\
48000	6771.9\\
50000	6612.8\\
52000	6743.3\\
54000	6569.7\\
56000	6772.7\\
58000	6653.8\\
60000	6766.5\\
62000	6719.3\\
64000	6809.8\\
68000	6575.5\\
70000	6826.9\\
72000	6659.2\\
74000	6639.3\\
76000	6655.4\\
78000	6572.5\\
80000	6656.1\\
82000	6567.1\\
84000	6709.5\\
86000	6780\\
88000	6759.6\\
90000	6619.6\\
92000	6723\\
94000	6638.5\\
96000	6581\\
98000	6738.4\\
1e+05	6639.2\\
};
\addplot [color=mycolor58, forget plot]
  table[row sep=crcr]{%
0	4395.3\\
2000	3950.3\\
4000	3741.6\\
6000	3624\\
8000	3580.3\\
10000	3648.7\\
12000	3572.1\\
16000	3595.3\\
18000	3586.4\\
20000	3648.6\\
22000	3591.1\\
24000	3516.8\\
26000	3557.9\\
28000	3516.4\\
30000	3648.8\\
34000	3529.9\\
36000	3658.6\\
38000	3553\\
40000	3651.3\\
42000	3643.7\\
44000	3578.1\\
46000	3618.5\\
48000	3506.6\\
50000	3537\\
52000	3588.1\\
54000	3535.3\\
56000	3512.3\\
58000	3562.5\\
60000	3440.5\\
62000	3598.4\\
64000	3568.4\\
66000	3727.6\\
68000	3709.8\\
70000	3574.4\\
72000	3574.5\\
76000	3600.8\\
78000	3563.9\\
80000	3556.2\\
82000	3645.8\\
84000	3528.4\\
88000	3588.8\\
90000	3581.1\\
92000	3547.4\\
94000	3622.2\\
98000	3563.7\\
1e+05	3549.7\\
};
\addplot [color=mycolor59, forget plot]
  table[row sep=crcr]{%
0	4395.3\\
2000	4669.1\\
4000	4898.9\\
6000	5012.6\\
8000	4974.4\\
10000	5149.3\\
12000	5012.2\\
14000	4985.1\\
18000	5128.3\\
20000	5110.2\\
22000	5028.9\\
24000	5159.8\\
26000	4958\\
28000	5125.3\\
30000	4904.2\\
32000	5049\\
34000	5095.5\\
36000	4973.2\\
38000	5170.9\\
40000	5065.8\\
44000	5064.4\\
46000	4928.8\\
48000	4955.9\\
50000	5128.7\\
52000	4961.2\\
54000	5062.5\\
56000	5061.3\\
58000	5025.3\\
60000	5123.2\\
62000	5139\\
64000	5087.3\\
66000	5052.8\\
68000	5103.9\\
70000	5115.3\\
72000	4996.1\\
74000	5067.8\\
76000	5192\\
78000	5108.3\\
80000	4987.2\\
82000	5112.9\\
84000	5014.4\\
86000	4948.9\\
88000	4999.2\\
90000	5073.2\\
92000	5033.4\\
94000	5098.2\\
96000	5111.4\\
98000	4962.5\\
1e+05	5104.2\\
};
\addplot [color=mycolor60, forget plot]
  table[row sep=crcr]{%
0	4395.3\\
2000	3868.7\\
6000	3627.8\\
10000	3563.9\\
12000	3738.3\\
14000	3642.5\\
16000	3588.1\\
18000	3487.4\\
20000	3528.2\\
22000	3615.6\\
24000	3479.5\\
26000	3574.3\\
28000	3569.3\\
30000	3668.5\\
32000	3575\\
34000	3473.9\\
36000	3543.8\\
38000	3538.4\\
40000	3652.4\\
48000	3568.9\\
50000	3513.1\\
52000	3607.8\\
54000	3513.2\\
56000	3494.5\\
58000	3607.7\\
60000	3428\\
62000	3524.5\\
64000	3495.2\\
66000	3498.1\\
68000	3564.9\\
70000	3535.5\\
72000	3524.3\\
74000	3489.2\\
76000	3516.4\\
78000	3592.3\\
80000	3546.8\\
86000	3562.3\\
88000	3446.4\\
90000	3632.2\\
92000	3568.9\\
94000	3482.5\\
96000	3571.2\\
98000	3601.5\\
1e+05	3513\\
};
\addplot [color=mycolor61, forget plot]
  table[row sep=crcr]{%
0	4395.3\\
2000	3327.2\\
4000	3152\\
6000	3111.6\\
8000	3135.1\\
10000	3127.5\\
12000	3006.4\\
14000	3126\\
16000	3137.5\\
22000	3128.8\\
26000	3113.5\\
28000	3155.6\\
30000	3131.9\\
32000	3122.4\\
36000	3181.1\\
40000	3129.3\\
42000	3203.2\\
44000	3155\\
46000	3130.3\\
48000	3141.3\\
50000	3114.5\\
54000	3149.7\\
56000	3232.5\\
58000	3180.1\\
60000	3161.8\\
62000	3106.8\\
64000	3106\\
66000	3093.7\\
68000	3145.8\\
70000	3167.2\\
72000	3132.1\\
74000	3174.7\\
76000	3135.3\\
78000	3198.4\\
80000	3093.4\\
82000	3168.7\\
84000	3123.7\\
86000	3121.8\\
88000	3130.5\\
90000	3104.2\\
92000	3140\\
94000	3104.4\\
96000	3135.5\\
98000	3138.4\\
1e+05	3156.2\\
};
\addplot [color=mycolor62, forget plot]
  table[row sep=crcr]{%
0	4395.3\\
2000	4743.2\\
4000	4965.1\\
6000	5070.6\\
8000	5081.9\\
12000	4956.1\\
14000	5095.4\\
16000	5035.6\\
18000	4886.9\\
20000	4972.7\\
22000	5028.5\\
24000	4986.5\\
26000	5096.3\\
30000	4977.1\\
32000	4900.5\\
34000	5026\\
36000	4922.1\\
38000	4977.7\\
40000	5001.5\\
42000	5048.6\\
44000	4918.9\\
46000	5022.7\\
48000	5044\\
50000	4962.4\\
52000	4996.6\\
54000	5085.5\\
56000	4912.8\\
58000	4982.4\\
60000	4928.1\\
62000	5054.4\\
64000	5026.1\\
66000	5084.7\\
68000	4942.1\\
70000	4980.6\\
72000	5069.7\\
74000	4983.3\\
76000	5032.4\\
82000	5016.1\\
84000	4932.5\\
86000	4964.6\\
88000	5019.7\\
90000	5006.5\\
92000	5018.5\\
94000	4933.4\\
96000	4973.6\\
98000	4885.4\\
1e+05	5049.7\\
};
\addplot [color=mycolor63, forget plot]
  table[row sep=crcr]{%
0	4395.3\\
2000	2983.7\\
4000	2495.2\\
6000	2419.1\\
8000	2431.3\\
10000	2404.3\\
12000	2339.1\\
14000	2436.2\\
16000	2422.5\\
18000	2440.6\\
20000	2426.1\\
22000	2447.2\\
24000	2404.5\\
26000	2369.3\\
28000	2424.8\\
30000	2465.1\\
32000	2426.4\\
34000	2395.4\\
38000	2407.3\\
40000	2445.4\\
42000	2412.6\\
44000	2425.3\\
46000	2418.3\\
50000	2442.4\\
52000	2439.4\\
54000	2471.4\\
56000	2492.1\\
58000	2417\\
62000	2393.1\\
64000	2364.7\\
66000	2415.7\\
68000	2390.2\\
70000	2430.7\\
72000	2449.5\\
74000	2451.7\\
76000	2437.4\\
78000	2439.4\\
80000	2462.6\\
82000	2373.4\\
84000	2402.4\\
86000	2474.1\\
88000	2388.8\\
90000	2451.6\\
92000	2459.3\\
96000	2434.8\\
98000	2379.2\\
1e+05	2399.4\\
};
\addplot [color=mycolor64, forget plot]
  table[row sep=crcr]{%
0	4395.3\\
2000	3204.9\\
4000	3061.8\\
6000	3056.6\\
8000	3019.1\\
10000	3072.4\\
12000	3036.2\\
14000	3083.3\\
16000	3041.7\\
18000	3040.6\\
20000	3061.5\\
22000	3117.6\\
24000	3038.9\\
26000	3009.8\\
28000	3034.8\\
30000	3007.6\\
32000	3113.8\\
34000	3044.8\\
36000	3123.7\\
38000	3067.5\\
40000	3138.1\\
42000	3066.3\\
44000	3052.8\\
46000	3068.1\\
48000	3127\\
50000	3082.9\\
52000	3112\\
54000	3068.3\\
56000	3049.6\\
58000	3053.7\\
60000	3039.2\\
62000	3046.2\\
64000	3070.1\\
66000	3016\\
68000	3033.7\\
70000	3039\\
72000	3075.3\\
74000	3049.7\\
76000	3096.8\\
80000	3016\\
82000	3058.7\\
84000	3114.6\\
88000	3057.6\\
90000	3057.5\\
92000	3089.3\\
94000	3046.1\\
96000	3073.2\\
98000	2995\\
1e+05	3062.5\\
};
\addplot [color=mycolor65, forget plot]
  table[row sep=crcr]{%
0	4395.3\\
2000	2869.6\\
4000	2560.2\\
6000	2492.9\\
8000	2508.8\\
10000	2517.8\\
12000	2556.6\\
14000	2522.1\\
16000	2530.4\\
18000	2581.3\\
20000	2608.7\\
22000	2560.6\\
24000	2578.1\\
26000	2580.6\\
28000	2555\\
30000	2538.4\\
32000	2586.2\\
34000	2555\\
36000	2543\\
40000	2556.5\\
44000	2502.5\\
46000	2601.4\\
48000	2547.7\\
50000	2594.1\\
52000	2547.4\\
54000	2555.7\\
56000	2576.9\\
58000	2542.9\\
60000	2497.5\\
62000	2565.8\\
66000	2563.6\\
68000	2552.4\\
70000	2529.4\\
72000	2569.8\\
74000	2579.3\\
76000	2518.3\\
80000	2552.7\\
82000	2589.7\\
84000	2560.8\\
86000	2509.7\\
88000	2579.6\\
90000	2544.9\\
92000	2574\\
94000	2594.8\\
96000	2534\\
98000	2586.7\\
1e+05	2565.4\\
};
\addplot [color=mycolor66, forget plot]
  table[row sep=crcr]{%
0	4395.3\\
2000	3537.5\\
4000	3407.1\\
6000	3461.5\\
8000	3478.8\\
10000	3444.4\\
12000	3445\\
14000	3398.3\\
16000	3453.6\\
18000	3435\\
20000	3401\\
22000	3456.3\\
24000	3441.7\\
26000	3404.6\\
28000	3459\\
30000	3465.5\\
32000	3413.7\\
34000	3461.7\\
36000	3438.7\\
38000	3510.6\\
42000	3464.2\\
44000	3352.6\\
46000	3392.6\\
48000	3395.9\\
50000	3378.3\\
52000	3410.1\\
54000	3377.2\\
56000	3443.6\\
58000	3467.4\\
60000	3433.5\\
62000	3378.2\\
64000	3395.1\\
66000	3457.2\\
68000	3403.3\\
70000	3412.3\\
76000	3492.5\\
78000	3457.6\\
80000	3451.6\\
82000	3384\\
86000	3395.5\\
88000	3450.4\\
90000	3411.4\\
92000	3418.8\\
94000	3439.2\\
96000	3382.5\\
98000	3456.6\\
1e+05	3472.1\\
};
\addplot [color=mycolor67, forget plot]
  table[row sep=crcr]{%
0	4395.3\\
2000	3625.2\\
4000	3227.7\\
6000	3206.4\\
8000	3203.1\\
10000	3126.9\\
12000	3205.5\\
14000	3139\\
16000	3147.1\\
18000	3232.8\\
20000	3218.4\\
22000	3189.9\\
24000	3184.7\\
26000	3213.3\\
28000	3184\\
30000	3211.6\\
32000	3187.9\\
34000	3193.3\\
36000	3208\\
38000	3235.6\\
40000	3197.8\\
42000	3174.6\\
44000	3235.4\\
46000	3232.4\\
48000	3181.5\\
50000	3196.2\\
52000	3161.8\\
54000	3290.9\\
56000	3249.9\\
58000	3237.9\\
60000	3166\\
62000	3197.5\\
64000	3190.1\\
66000	3153.3\\
68000	3225.5\\
70000	3253\\
72000	3183.8\\
74000	3160.6\\
76000	3123.3\\
78000	3182.7\\
80000	3186.2\\
86000	3155\\
88000	3232.8\\
90000	3247.5\\
92000	3209.1\\
94000	3235.1\\
96000	3164\\
98000	3187.8\\
1e+05	3190.5\\
};
\addplot [color=mycolor67]
  table[row sep=crcr]{%
0	4395.3\\
2000	2828.3\\
4000	2619.7\\
6000	2660.2\\
12000	2632.7\\
16000	2651.6\\
18000	2634.3\\
20000	2655\\
22000	2695.7\\
24000	2679.7\\
26000	2674.4\\
30000	2612.2\\
32000	2629.5\\
34000	2613.7\\
36000	2659.5\\
38000	2672.8\\
40000	2646.3\\
42000	2658.4\\
44000	2657.8\\
46000	2686.9\\
50000	2609.8\\
52000	2610\\
54000	2680.3\\
56000	2666.3\\
58000	2640\\
60000	2599.8\\
62000	2651.3\\
64000	2629\\
66000	2658.6\\
68000	2650.6\\
72000	2619.5\\
74000	2631.2\\
76000	2624.1\\
78000	2663.3\\
80000	2680.8\\
82000	2635\\
84000	2625.4\\
86000	2679.3\\
88000	2644.2\\
90000	2576.2\\
92000	2627.6\\
96000	2634.5\\
98000	2636.5\\
1e+05	2626.6\\
};
\addlegendentry{Classe k=43}

\addplot [color=mycolor68, forget plot]
  table[row sep=crcr]{%
0	4395.3\\
2000	2924\\
4000	2605.9\\
6000	2498.9\\
8000	2506.2\\
10000	2476.6\\
12000	2482.2\\
14000	2508\\
16000	2492.1\\
20000	2481.9\\
22000	2451.6\\
24000	2476\\
26000	2509.2\\
28000	2466.9\\
30000	2530\\
32000	2480.6\\
36000	2486.7\\
38000	2468.4\\
40000	2494.2\\
44000	2487.4\\
46000	2473.3\\
48000	2477\\
50000	2514.7\\
54000	2460.7\\
56000	2464.7\\
58000	2483.2\\
60000	2489.5\\
64000	2521.4\\
66000	2499\\
68000	2499.5\\
70000	2464.7\\
72000	2498.9\\
74000	2467.7\\
76000	2481.6\\
78000	2447.4\\
80000	2496.5\\
82000	2493.4\\
84000	2481.6\\
86000	2493.7\\
88000	2491.9\\
90000	2467.2\\
92000	2457\\
94000	2483.1\\
96000	2496.4\\
98000	2476.5\\
1e+05	2496.9\\
};
\addplot [color=mycolor69, forget plot]
  table[row sep=crcr]{%
0	4395.3\\
2000	3134.6\\
4000	2859.6\\
6000	2865.1\\
8000	2935.4\\
12000	2898.4\\
14000	2936.2\\
16000	2920.6\\
18000	2883.6\\
20000	2927.9\\
22000	2928.6\\
24000	2919.4\\
26000	2852.6\\
28000	2908.2\\
30000	2951.8\\
32000	2914.8\\
36000	2886.4\\
38000	2939.9\\
40000	2901.6\\
42000	2929.5\\
44000	2930.6\\
46000	2916.7\\
48000	2932.2\\
50000	2906.2\\
52000	2919.9\\
54000	2886.3\\
56000	2877.7\\
58000	2909.1\\
60000	2895.7\\
62000	2922.7\\
64000	2936.2\\
66000	2968.4\\
68000	2934.5\\
70000	2981\\
76000	2918.3\\
78000	2931.8\\
80000	2898.4\\
82000	2930.3\\
84000	2925.1\\
86000	2931.3\\
88000	2916.8\\
92000	2925.9\\
94000	2946.7\\
96000	2938.3\\
98000	2897.7\\
1e+05	2926\\
};
\addplot [color=mycolor70, forget plot]
  table[row sep=crcr]{%
0	4395.3\\
2000	2394.4\\
4000	2082.5\\
6000	2113.4\\
8000	2118.7\\
10000	2133.7\\
12000	2114.1\\
16000	2122.6\\
18000	2140.8\\
20000	2106.2\\
22000	2125.2\\
24000	2136.2\\
26000	2072.4\\
28000	2141.6\\
30000	2128.9\\
32000	2100.7\\
34000	2097.7\\
36000	2116.8\\
38000	2096.8\\
40000	2124.5\\
42000	2118\\
44000	2151.9\\
46000	2110.1\\
48000	2117.1\\
50000	2136.8\\
52000	2109.1\\
54000	2130\\
56000	2095.6\\
58000	2129.3\\
60000	2126.8\\
62000	2149.4\\
64000	2157.3\\
66000	2130.4\\
68000	2136.6\\
70000	2166.3\\
72000	2119\\
74000	2116\\
76000	2101.1\\
82000	2155\\
84000	2093.5\\
90000	2107.7\\
92000	2071.2\\
94000	2115.9\\
96000	2140.9\\
98000	2152.1\\
1e+05	2184.1\\
};
\addplot [color=mycolor71, forget plot]
  table[row sep=crcr]{%
0	4395.3\\
2000	2896.1\\
4000	2609\\
6000	2544.9\\
8000	2510.6\\
10000	2520\\
12000	2486.7\\
14000	2504.7\\
16000	2498.2\\
18000	2509.8\\
20000	2478.7\\
24000	2515.7\\
26000	2487.2\\
28000	2483.1\\
30000	2512.3\\
34000	2549.4\\
36000	2498.2\\
38000	2496.3\\
40000	2477.9\\
42000	2498.6\\
46000	2518.1\\
48000	2534.9\\
50000	2490.1\\
54000	2524.3\\
56000	2504.1\\
58000	2519.9\\
60000	2504.8\\
62000	2500.2\\
64000	2465.7\\
66000	2473.1\\
68000	2495.9\\
70000	2503.4\\
72000	2469.5\\
74000	2487.4\\
78000	2495.9\\
80000	2493.4\\
82000	2468.7\\
84000	2516\\
86000	2469.9\\
88000	2514.5\\
94000	2472\\
96000	2482.2\\
1e+05	2513.5\\
};
\addplot [color=mycolor72, forget plot]
  table[row sep=crcr]{%
0	4395.3\\
2000	2704.6\\
4000	2182.9\\
6000	2239.6\\
8000	2197.4\\
10000	2220.1\\
12000	2223.6\\
14000	2217.8\\
18000	2249.1\\
20000	2201.8\\
22000	2202.3\\
26000	2234.7\\
28000	2219.7\\
30000	2196\\
32000	2209.6\\
34000	2177.1\\
36000	2226.7\\
38000	2192.8\\
40000	2232.6\\
42000	2184.6\\
44000	2177.4\\
46000	2230.1\\
48000	2186.4\\
50000	2211.4\\
52000	2209.7\\
54000	2158.2\\
56000	2244.9\\
58000	2212.2\\
60000	2230.4\\
62000	2242.1\\
64000	2196.5\\
66000	2192.7\\
68000	2229.8\\
72000	2183.8\\
76000	2240.8\\
78000	2201.4\\
80000	2209\\
82000	2157.1\\
84000	2212.7\\
86000	2211.1\\
88000	2258\\
90000	2240.3\\
92000	2215.6\\
94000	2163.3\\
96000	2276.2\\
98000	2231.4\\
1e+05	2203.4\\
};
\addplot [color=mycolor73, forget plot]
  table[row sep=crcr]{%
0	4395.3\\
2000	2962\\
4000	2679.2\\
6000	2642.2\\
8000	2617.8\\
10000	2631.4\\
12000	2657.9\\
14000	2601.5\\
16000	2604.2\\
18000	2617.9\\
20000	2662.2\\
24000	2626.5\\
26000	2629.7\\
28000	2692.6\\
30000	2657.6\\
32000	2590.3\\
34000	2626.2\\
38000	2638.6\\
40000	2600.2\\
42000	2619.3\\
44000	2606.2\\
46000	2620.2\\
50000	2662.8\\
52000	2682.2\\
54000	2660.3\\
56000	2583.8\\
58000	2573.7\\
60000	2633.8\\
62000	2603.9\\
64000	2595.8\\
66000	2661.9\\
68000	2601.3\\
70000	2593.8\\
74000	2609.8\\
76000	2611.3\\
78000	2632\\
80000	2630.5\\
82000	2616.5\\
84000	2667.2\\
86000	2628.3\\
88000	2656.3\\
90000	2628.8\\
92000	2624\\
94000	2607.3\\
96000	2605.4\\
1e+05	2615.9\\
};
\addplot [color=mycolor74, forget plot]
  table[row sep=crcr]{%
0	4395.3\\
2000	2198.3\\
4000	1590.4\\
6000	1573.1\\
8000	1586.4\\
10000	1562.2\\
12000	1550.8\\
14000	1526.3\\
16000	1544.1\\
18000	1550.3\\
20000	1567.6\\
22000	1571.9\\
24000	1530.2\\
26000	1559.9\\
28000	1550.6\\
30000	1564.9\\
32000	1543.7\\
36000	1548.3\\
38000	1574\\
40000	1572.9\\
42000	1566\\
44000	1567.4\\
46000	1549.2\\
48000	1554.7\\
50000	1567.8\\
52000	1548.1\\
54000	1549.5\\
56000	1557.6\\
58000	1525.6\\
60000	1536.6\\
62000	1541.5\\
64000	1562.4\\
66000	1561.2\\
68000	1534.4\\
70000	1537.7\\
72000	1564.5\\
74000	1556.7\\
78000	1556.1\\
80000	1578.6\\
82000	1548.5\\
84000	1545.2\\
86000	1529.1\\
88000	1565.5\\
90000	1549.6\\
92000	1544.3\\
94000	1555.7\\
96000	1557.7\\
98000	1526.3\\
1e+05	1552.8\\
};
\addplot [color=mycolor75, forget plot]
  table[row sep=crcr]{%
0	4395.3\\
2000	2726.7\\
4000	2276.9\\
6000	2244.6\\
8000	2280.4\\
10000	2210.4\\
12000	2256.1\\
14000	2226.6\\
16000	2228.9\\
18000	2197.2\\
20000	2216.4\\
22000	2223.8\\
24000	2238.9\\
26000	2219.1\\
28000	2260.3\\
30000	2209.8\\
32000	2227.3\\
34000	2210.4\\
38000	2208.3\\
40000	2241.4\\
42000	2220.6\\
44000	2212.3\\
46000	2228.4\\
48000	2205.2\\
50000	2200\\
54000	2242.3\\
56000	2212\\
58000	2238.7\\
62000	2210.8\\
64000	2166.7\\
66000	2223.3\\
68000	2246.4\\
70000	2210.2\\
72000	2198.7\\
74000	2227.2\\
76000	2196.3\\
78000	2238.8\\
80000	2226.3\\
82000	2227.5\\
84000	2220.3\\
86000	2204\\
88000	2209.2\\
90000	2198.8\\
92000	2223.7\\
94000	2274\\
96000	2199.6\\
98000	2235.3\\
1e+05	2226.9\\
};
\addplot [color=mycolor76, forget plot]
  table[row sep=crcr]{%
0	4395.3\\
2000	2567.6\\
4000	2096.5\\
8000	2032.5\\
10000	2039.8\\
12000	2018.6\\
20000	2029.4\\
22000	2042.2\\
24000	2009.7\\
28000	2031.7\\
32000	2005.7\\
34000	2036.6\\
36000	2006.9\\
40000	2022.9\\
42000	1999.2\\
46000	2043.7\\
50000	2031.2\\
54000	2058.6\\
56000	2024.8\\
58000	2009.5\\
60000	2031.7\\
62000	2012\\
64000	1998.9\\
66000	2027.2\\
70000	2025.7\\
72000	1999.7\\
74000	2020.5\\
76000	2032.1\\
78000	2026.4\\
80000	2002.9\\
82000	2038.9\\
84000	2020.7\\
86000	2051.6\\
88000	2049.4\\
90000	2014\\
92000	2009.4\\
94000	2040.8\\
96000	2015.5\\
98000	2022.8\\
1e+05	2064.2\\
};
\addplot [color=mycolor77, forget plot]
  table[row sep=crcr]{%
0	4395.3\\
2000	2250.8\\
4000	1675.1\\
6000	1624.4\\
10000	1574.5\\
12000	1579\\
14000	1572.1\\
16000	1575.3\\
18000	1591.8\\
20000	1560.7\\
24000	1580.8\\
26000	1566.6\\
28000	1596.4\\
30000	1563.3\\
32000	1557.1\\
34000	1580.7\\
36000	1571.4\\
38000	1570.9\\
40000	1579.8\\
42000	1569.8\\
44000	1593.2\\
46000	1553.4\\
48000	1592.5\\
50000	1565.3\\
54000	1573.9\\
56000	1587.2\\
58000	1556.4\\
60000	1570.5\\
62000	1561.6\\
64000	1607.4\\
66000	1577.9\\
68000	1590.1\\
70000	1583.3\\
72000	1567.3\\
74000	1580.3\\
76000	1568.4\\
78000	1568.3\\
80000	1573.3\\
82000	1566.6\\
84000	1570.2\\
86000	1579.8\\
88000	1567.3\\
90000	1563.5\\
92000	1569.2\\
94000	1591.2\\
96000	1596.5\\
98000	1571.6\\
1e+05	1569.7\\
};
\addplot [color=mycolor78, forget plot]
  table[row sep=crcr]{%
0	4395.3\\
2000	2840.2\\
4000	2385.9\\
6000	2343.7\\
8000	2319.4\\
10000	2339.9\\
12000	2318\\
14000	2311.5\\
18000	2274.2\\
20000	2319.8\\
22000	2271.9\\
24000	2316.3\\
26000	2337.6\\
28000	2348.1\\
30000	2293.4\\
32000	2326.1\\
34000	2335.6\\
36000	2308.4\\
38000	2340.9\\
40000	2325.2\\
42000	2270.3\\
44000	2277.5\\
46000	2341.9\\
48000	2340.1\\
50000	2287.8\\
52000	2300.5\\
56000	2261.1\\
58000	2319.2\\
60000	2310.5\\
62000	2320.5\\
64000	2257.7\\
66000	2278.1\\
68000	2318.7\\
70000	2307\\
72000	2319.9\\
74000	2317.2\\
76000	2343.8\\
78000	2296.1\\
80000	2259.6\\
82000	2273\\
84000	2279.3\\
86000	2277.5\\
88000	2322.4\\
90000	2304.8\\
92000	2277.4\\
94000	2289.3\\
96000	2319.6\\
98000	2315.2\\
1e+05	2317.9\\
};
\addplot [color=mycolor79, forget plot]
  table[row sep=crcr]{%
0	4395.3\\
2000	2689.7\\
4000	2316.9\\
6000	2225.1\\
10000	2218.4\\
12000	2195.6\\
14000	2233.2\\
16000	2204.4\\
18000	2212.2\\
22000	2149.7\\
24000	2190.8\\
26000	2194.1\\
28000	2172.6\\
30000	2158\\
32000	2116.9\\
34000	2172\\
36000	2197.9\\
42000	2150.4\\
44000	2192.3\\
46000	2156.7\\
48000	2162.1\\
50000	2149.5\\
52000	2168.6\\
54000	2159.7\\
56000	2166.2\\
58000	2124.8\\
60000	2170.4\\
62000	2146\\
66000	2217.6\\
68000	2154.6\\
70000	2181.3\\
72000	2161.4\\
74000	2160.8\\
76000	2138.2\\
78000	2166.2\\
80000	2168.9\\
82000	2129.1\\
84000	2173.2\\
86000	2156.9\\
88000	2171.9\\
90000	2208.8\\
92000	2221.3\\
94000	2178.3\\
96000	2198\\
98000	2207.1\\
1e+05	2207.2\\
};
\addplot [color=mycolor79, forget plot]
  table[row sep=crcr]{%
0	4395.3\\
2000	2672.3\\
4000	2067.4\\
6000	1994.7\\
8000	2014.2\\
10000	1990.4\\
12000	2023\\
14000	2040.8\\
16000	1997.8\\
18000	2028.8\\
20000	2070\\
22000	2026.1\\
24000	2034.5\\
26000	2013.1\\
28000	2060.6\\
30000	2007.6\\
32000	2010.7\\
34000	2040.3\\
36000	2008.6\\
38000	2051.8\\
40000	2062.5\\
42000	2048.6\\
44000	2078\\
46000	2010.6\\
48000	2013.1\\
50000	2005.7\\
52000	2026.2\\
54000	2022.6\\
56000	2031\\
58000	2007.6\\
60000	2029.2\\
62000	2008.4\\
64000	2050.1\\
66000	2050.3\\
68000	2027.5\\
70000	1994.3\\
72000	1993.3\\
74000	2007.5\\
76000	1983.8\\
78000	1992.2\\
82000	2045.7\\
84000	2012.4\\
86000	2057.9\\
88000	2044.9\\
90000	2064.5\\
92000	2017.9\\
94000	2025.9\\
96000	2052.2\\
98000	1982\\
1e+05	1989.9\\
};
\addplot [color=mycolor80, forget plot]
  table[row sep=crcr]{%
0	4395.3\\
2000	2172.5\\
4000	1637.8\\
6000	1585.9\\
8000	1577.5\\
12000	1587.6\\
14000	1602.6\\
16000	1568.8\\
18000	1553.3\\
20000	1567.7\\
22000	1576.3\\
24000	1553.6\\
26000	1585.6\\
28000	1561.7\\
30000	1565\\
32000	1576.3\\
34000	1578\\
36000	1591\\
38000	1555.7\\
40000	1560.9\\
44000	1549.1\\
46000	1585.7\\
50000	1554.7\\
52000	1564.6\\
56000	1597.5\\
58000	1561.2\\
62000	1578.8\\
64000	1571.3\\
66000	1571.9\\
68000	1588.9\\
72000	1557.1\\
74000	1557.4\\
76000	1570.3\\
78000	1551.7\\
80000	1579.1\\
82000	1566.7\\
84000	1571.4\\
86000	1563.3\\
88000	1572.5\\
90000	1574.5\\
92000	1557.8\\
96000	1568.8\\
98000	1547.4\\
1e+05	1557\\
};
\addplot [color=mycolor81, forget plot]
  table[row sep=crcr]{%
0	4395.3\\
2000	2130.1\\
4000	1490.3\\
6000	1456.5\\
8000	1461.4\\
10000	1475.9\\
12000	1453.9\\
14000	1450.9\\
16000	1452.9\\
18000	1444.6\\
20000	1475.8\\
22000	1468.2\\
24000	1484.7\\
28000	1453.5\\
30000	1453.4\\
32000	1441.5\\
34000	1452.3\\
36000	1484.2\\
38000	1469.1\\
40000	1467.7\\
44000	1448.4\\
48000	1474.9\\
50000	1492.6\\
52000	1490.7\\
54000	1484.5\\
58000	1480.5\\
60000	1488.7\\
62000	1455.7\\
66000	1483.9\\
68000	1511.5\\
70000	1457.8\\
72000	1487.4\\
76000	1458.4\\
80000	1462.4\\
84000	1469.4\\
86000	1454.3\\
88000	1467.9\\
90000	1453.3\\
92000	1481\\
94000	1459.6\\
96000	1429.5\\
98000	1464.7\\
1e+05	1477.2\\
};
\addplot [color=mycolor82, forget plot]
  table[row sep=crcr]{%
0	4395.3\\
2000	2370.7\\
4000	1818.8\\
6000	1734.9\\
10000	1711.7\\
12000	1717.8\\
14000	1712.4\\
16000	1697\\
18000	1716.1\\
20000	1716.5\\
22000	1688\\
24000	1715.8\\
26000	1701.4\\
28000	1707.6\\
30000	1725.3\\
32000	1703.9\\
34000	1706.6\\
38000	1723.4\\
40000	1670.2\\
42000	1714.5\\
46000	1747.4\\
48000	1705.8\\
50000	1699.8\\
52000	1713.6\\
54000	1687.9\\
56000	1695.3\\
58000	1717.8\\
64000	1717.3\\
66000	1727.7\\
68000	1705.6\\
72000	1677.9\\
74000	1695.9\\
76000	1724\\
78000	1724.2\\
80000	1706.7\\
84000	1681.4\\
86000	1689.6\\
88000	1709.1\\
90000	1701.4\\
92000	1706.6\\
94000	1717.2\\
96000	1710\\
98000	1710.3\\
1e+05	1743.1\\
};
\addplot [color=mycolor83, forget plot]
  table[row sep=crcr]{%
0	4395.3\\
2000	2453.2\\
4000	1858\\
6000	1809.4\\
8000	1773.5\\
10000	1746.8\\
16000	1739.3\\
18000	1746.5\\
20000	1736.1\\
26000	1757.2\\
28000	1753.2\\
30000	1769.4\\
32000	1731.2\\
36000	1757.3\\
38000	1709.9\\
40000	1717.9\\
42000	1718.4\\
44000	1750.4\\
48000	1731.7\\
52000	1732\\
54000	1734\\
56000	1718.3\\
58000	1737.5\\
60000	1712.2\\
62000	1734.2\\
64000	1727.1\\
66000	1748.2\\
68000	1749.8\\
72000	1727.1\\
78000	1730.8\\
80000	1730\\
82000	1741.3\\
86000	1736\\
88000	1745.7\\
90000	1770\\
92000	1764\\
96000	1727.4\\
98000	1742.4\\
1e+05	1724.1\\
};
\addplot [color=mycolor83, forget plot]
  table[row sep=crcr]{%
0	4395.3\\
2000	1920.4\\
4000	1289.2\\
6000	1292.1\\
8000	1275.5\\
10000	1272.2\\
12000	1262.7\\
14000	1271.8\\
16000	1273\\
18000	1283\\
20000	1264.3\\
22000	1280.3\\
24000	1253.7\\
26000	1270.9\\
28000	1272.2\\
30000	1264.4\\
32000	1260.7\\
34000	1266.7\\
36000	1303.4\\
42000	1268.7\\
44000	1284.7\\
46000	1276.3\\
48000	1281.7\\
50000	1259.9\\
52000	1252\\
54000	1272.8\\
56000	1284.6\\
58000	1279.9\\
60000	1287.4\\
62000	1256.8\\
64000	1281.2\\
66000	1331.8\\
68000	1281.7\\
70000	1270.2\\
72000	1264.7\\
76000	1286.2\\
78000	1275.8\\
80000	1284.2\\
82000	1285.2\\
84000	1250.8\\
86000	1273.9\\
88000	1306.7\\
90000	1294.4\\
92000	1304\\
94000	1290.7\\
96000	1269.9\\
98000	1269.9\\
1e+05	1276.4\\
};
\addplot [color=mycolor84, forget plot]
  table[row sep=crcr]{%
0	4395.3\\
2000	1990\\
4000	1371.1\\
6000	1270.2\\
10000	1292\\
12000	1284.3\\
14000	1291.5\\
16000	1291.6\\
18000	1300.1\\
20000	1267.5\\
22000	1242.9\\
24000	1298.4\\
26000	1264.5\\
28000	1266.9\\
30000	1282\\
34000	1295.6\\
36000	1322\\
38000	1283.7\\
40000	1302.8\\
42000	1293\\
44000	1258.9\\
46000	1282.6\\
48000	1261.4\\
50000	1297.9\\
52000	1314.9\\
54000	1267.7\\
56000	1271\\
58000	1252.6\\
60000	1285.9\\
62000	1284.4\\
64000	1278.4\\
66000	1279.3\\
68000	1269.7\\
70000	1283.6\\
72000	1263.1\\
74000	1302.1\\
76000	1322.6\\
78000	1283.5\\
80000	1273.3\\
82000	1279.3\\
84000	1308.4\\
86000	1290.6\\
88000	1246.4\\
90000	1295.1\\
92000	1275.7\\
94000	1312.9\\
96000	1302.4\\
98000	1306.5\\
1e+05	1286.3\\
};
\addplot [color=mycolor85, forget plot]
  table[row sep=crcr]{%
0	4395.3\\
2000	2212.2\\
4000	1490.5\\
6000	1444.3\\
8000	1455.3\\
10000	1405.5\\
12000	1435.2\\
14000	1422.4\\
16000	1394.6\\
18000	1402.2\\
20000	1457.4\\
22000	1407.4\\
26000	1362.3\\
28000	1414.7\\
30000	1384.8\\
32000	1430.2\\
34000	1400.6\\
38000	1422.6\\
40000	1388.8\\
42000	1403.6\\
44000	1438\\
48000	1385.1\\
50000	1439.2\\
52000	1404.9\\
54000	1412.3\\
56000	1428.6\\
60000	1405.9\\
62000	1408\\
64000	1439.5\\
66000	1414.8\\
68000	1417.1\\
70000	1398.6\\
72000	1445.8\\
74000	1436.2\\
76000	1433.8\\
78000	1424\\
80000	1407.4\\
82000	1438.5\\
84000	1410.7\\
88000	1388.6\\
92000	1413\\
94000	1404.9\\
96000	1430.8\\
98000	1379.5\\
1e+05	1414\\
};
\addplot [color=mycolor86, forget plot]
  table[row sep=crcr]{%
0	4395.3\\
2000	1980.1\\
4000	1262.3\\
6000	1200.4\\
8000	1228.3\\
10000	1180.1\\
12000	1177.2\\
16000	1188.2\\
18000	1202.6\\
20000	1198.1\\
22000	1202.3\\
24000	1215.2\\
26000	1191.7\\
28000	1175.3\\
30000	1175\\
32000	1204.7\\
34000	1179.9\\
36000	1185\\
38000	1174.9\\
40000	1203.7\\
42000	1243.5\\
44000	1202.2\\
46000	1198.4\\
48000	1187\\
50000	1186.1\\
52000	1189.8\\
54000	1207.5\\
56000	1205.3\\
58000	1182\\
60000	1178.9\\
62000	1191.3\\
64000	1196.1\\
66000	1181.3\\
68000	1186.7\\
70000	1181.6\\
72000	1184.2\\
74000	1223.5\\
76000	1225.2\\
78000	1199.1\\
80000	1154\\
82000	1184\\
84000	1190.1\\
86000	1186.6\\
88000	1206.5\\
90000	1185.8\\
92000	1171.4\\
94000	1177.2\\
96000	1176.3\\
98000	1214\\
1e+05	1181.2\\
};
\addplot [color=mycolor87, forget plot]
  table[row sep=crcr]{%
0	4395.3\\
2000	1808.9\\
4000	996.57\\
6000	1026.8\\
8000	1027\\
10000	960.79\\
12000	986.53\\
14000	1029.1\\
16000	1043.8\\
18000	1005.2\\
20000	991.6\\
22000	1032.6\\
26000	983.22\\
30000	992.74\\
32000	972.98\\
34000	978.71\\
36000	994.95\\
38000	986.05\\
40000	990.8\\
42000	1001\\
44000	1005.2\\
46000	1023.2\\
48000	1028\\
50000	1019.8\\
54000	1011.8\\
56000	1002.8\\
58000	1049\\
60000	1022.6\\
64000	1003.7\\
66000	1058.3\\
68000	1036.2\\
70000	1030.9\\
72000	1042.1\\
74000	1006.2\\
78000	980.49\\
80000	1035.1\\
82000	1005.9\\
84000	1012.3\\
86000	1026.4\\
88000	997.64\\
90000	1020.4\\
92000	1002.9\\
94000	1006.4\\
96000	981.8\\
98000	1011.1\\
1e+05	960.27\\
};
\addplot [color=mycolor88, forget plot]
  table[row sep=crcr]{%
0	4395.3\\
2000	2033.5\\
4000	1269\\
6000	1200.5\\
8000	1212\\
10000	1188.2\\
12000	1186.9\\
14000	1201.8\\
16000	1190.3\\
18000	1189.4\\
20000	1160.5\\
24000	1190.3\\
26000	1176.7\\
28000	1176.5\\
30000	1204.5\\
32000	1203.9\\
34000	1182.8\\
36000	1205\\
38000	1208.1\\
40000	1188.2\\
42000	1174\\
44000	1198.9\\
46000	1172.2\\
48000	1153.7\\
50000	1169.4\\
52000	1196\\
54000	1187.7\\
56000	1183.9\\
58000	1172.1\\
60000	1201.5\\
62000	1193.7\\
64000	1180.2\\
66000	1195.1\\
68000	1174.7\\
70000	1173.7\\
72000	1156.6\\
74000	1149.7\\
78000	1211.5\\
80000	1167.1\\
82000	1171.9\\
84000	1194.2\\
86000	1183.3\\
88000	1176.2\\
90000	1201.9\\
92000	1238.2\\
94000	1168.2\\
96000	1197.7\\
98000	1219.4\\
1e+05	1193.3\\
};
\addplot [color=mycolor89, forget plot]
  table[row sep=crcr]{%
0	4395.3\\
2000	1940.9\\
4000	948.01\\
6000	855.39\\
8000	871.42\\
10000	846.68\\
12000	855.41\\
14000	847.31\\
16000	863.03\\
18000	894.33\\
20000	877.19\\
22000	856.51\\
24000	868.24\\
28000	862.99\\
30000	848.05\\
32000	855.28\\
34000	856.2\\
36000	851.43\\
38000	834.69\\
40000	835.11\\
42000	856.88\\
44000	844.13\\
48000	877.64\\
50000	862.84\\
52000	839.15\\
54000	856.15\\
56000	858.39\\
58000	835.46\\
60000	875.68\\
62000	849.53\\
64000	867.38\\
66000	851.13\\
68000	848.83\\
72000	855.62\\
74000	872.27\\
76000	867.41\\
78000	850.02\\
80000	857.22\\
82000	856.36\\
84000	840.75\\
86000	857.07\\
88000	862.3\\
90000	871.11\\
92000	862.31\\
94000	862.18\\
96000	857\\
1e+05	856.63\\
};
\addplot [color=mycolor90, forget plot]
  table[row sep=crcr]{%
0	4395.3\\
2000	1686.3\\
4000	858.74\\
8000	830.3\\
10000	836.2\\
12000	823.74\\
14000	831.61\\
16000	844.62\\
18000	836.81\\
20000	818.04\\
22000	829.01\\
24000	822.85\\
26000	828.09\\
28000	820.06\\
30000	835.9\\
32000	842.94\\
34000	835.12\\
36000	840.15\\
38000	836.32\\
40000	822.66\\
42000	823.76\\
44000	835.9\\
46000	830.45\\
48000	832.14\\
50000	839.5\\
52000	839.82\\
56000	823.64\\
58000	852.2\\
60000	827.95\\
62000	825.48\\
64000	839.7\\
66000	831.32\\
68000	832.75\\
70000	808.45\\
72000	842.31\\
74000	825.2\\
76000	826.49\\
78000	806.33\\
80000	819.04\\
82000	837.64\\
84000	837.37\\
86000	847.08\\
90000	824.02\\
92000	821.84\\
94000	839.58\\
96000	808.89\\
98000	813.87\\
1e+05	829.33\\
};
\addplot [color=mycolor91, forget plot]
  table[row sep=crcr]{%
0	4395.3\\
2000	1863.3\\
4000	811.12\\
6000	720.23\\
8000	712.4\\
10000	717.99\\
12000	705.68\\
14000	735.49\\
16000	726.02\\
18000	751.62\\
20000	753.25\\
22000	745.99\\
24000	724.42\\
26000	732.48\\
28000	721.21\\
30000	716.77\\
32000	717.5\\
34000	737.5\\
36000	702.65\\
38000	717.53\\
40000	723.53\\
42000	688.36\\
44000	699.97\\
46000	753.63\\
48000	735.15\\
50000	712.58\\
52000	725.21\\
56000	696.95\\
60000	724.66\\
62000	702.93\\
64000	732.72\\
66000	744.26\\
68000	736.6\\
70000	715.89\\
72000	729.59\\
74000	710.88\\
76000	709.22\\
78000	727.73\\
80000	706.52\\
82000	734.53\\
84000	715.09\\
86000	740.19\\
88000	740.99\\
90000	723.14\\
92000	721.3\\
94000	691.27\\
96000	699.27\\
98000	711.13\\
1e+05	720.57\\
};
\addplot [color=mycolor92, forget plot]
  table[row sep=crcr]{%
0	4395.3\\
2000	1880.9\\
4000	1004.3\\
6000	919.7\\
8000	910.02\\
10000	911.14\\
12000	876.57\\
14000	897\\
16000	888.46\\
18000	931.41\\
20000	926.47\\
22000	905.48\\
24000	882.09\\
26000	889.99\\
28000	892.21\\
30000	916.01\\
32000	888.99\\
34000	926.91\\
36000	898.07\\
38000	898.62\\
40000	880.72\\
42000	907.77\\
44000	900.75\\
46000	888.87\\
48000	895.21\\
50000	905.63\\
52000	902.29\\
54000	888.54\\
56000	919.63\\
58000	885.24\\
60000	904.17\\
64000	912.88\\
66000	913.68\\
68000	883.97\\
70000	908.53\\
72000	905.59\\
74000	895.21\\
76000	923.77\\
78000	905.08\\
80000	902.31\\
82000	892.61\\
84000	900.37\\
86000	905.31\\
88000	907.26\\
90000	902.62\\
92000	911.56\\
94000	917.25\\
98000	871.29\\
1e+05	904.89\\
};
\addplot [color=mycolor93, forget plot]
  table[row sep=crcr]{%
0	4395.3\\
2000	1995.2\\
4000	887.41\\
6000	783.05\\
8000	765.48\\
10000	756.07\\
12000	770.36\\
14000	798.91\\
18000	757.01\\
20000	747.05\\
22000	748.53\\
24000	742.21\\
26000	751.73\\
28000	751.33\\
30000	753.77\\
32000	749.62\\
34000	750.13\\
36000	771.32\\
38000	762.69\\
40000	726.4\\
42000	761.8\\
44000	774.36\\
46000	756.42\\
48000	767.45\\
50000	744.65\\
52000	746.45\\
54000	744.84\\
56000	770.42\\
58000	728.58\\
60000	774.08\\
62000	764.83\\
64000	768\\
66000	780.78\\
68000	735.89\\
70000	783.81\\
72000	773.71\\
74000	770.47\\
76000	733.06\\
78000	737.75\\
80000	764.64\\
82000	711.1\\
84000	735.85\\
86000	776.17\\
88000	748.7\\
90000	743.54\\
92000	761.04\\
94000	745.85\\
96000	749.31\\
98000	748.22\\
1e+05	767.2\\
};
\addplot [color=mycolor94, forget plot]
  table[row sep=crcr]{%
0	4395.3\\
2000	1955.8\\
4000	842.05\\
6000	751.87\\
8000	745.71\\
10000	733.81\\
14000	718.91\\
16000	708.29\\
18000	729.12\\
20000	733.88\\
22000	720.74\\
24000	701.82\\
26000	728.79\\
28000	727.51\\
30000	706.63\\
32000	697.34\\
34000	697.58\\
36000	725.32\\
38000	726.24\\
40000	717\\
44000	723.05\\
46000	719.4\\
48000	731.77\\
50000	721.95\\
52000	724.23\\
54000	716.79\\
56000	725.58\\
58000	713.6\\
62000	708.19\\
64000	710\\
66000	695.14\\
68000	721.39\\
70000	715.65\\
72000	738.39\\
74000	734.6\\
76000	708.64\\
78000	710.34\\
80000	702.73\\
82000	703.4\\
84000	722.26\\
86000	705.01\\
88000	729.23\\
90000	679.9\\
92000	688.26\\
94000	692.06\\
96000	726.28\\
98000	713.4\\
1e+05	709.72\\
};
\addplot [color=mycolor95, forget plot]
  table[row sep=crcr]{%
0	4395.3\\
2000	1552.9\\
4000	516.54\\
6000	439.22\\
8000	431.49\\
10000	426.97\\
12000	431.52\\
14000	429.19\\
16000	424.92\\
18000	423.06\\
20000	428.05\\
22000	436.24\\
24000	425.04\\
26000	430.12\\
28000	452.52\\
32000	422.81\\
34000	414.58\\
36000	409.76\\
38000	422.94\\
40000	442.73\\
42000	410.3\\
44000	437\\
46000	414.51\\
48000	407.56\\
50000	441.58\\
52000	436.73\\
54000	422.2\\
56000	423.29\\
58000	425.8\\
60000	432.57\\
62000	423.22\\
64000	412.21\\
66000	436.18\\
68000	437.66\\
70000	426.87\\
72000	423.32\\
74000	422.84\\
76000	434.06\\
78000	419.74\\
80000	446.39\\
82000	422.24\\
84000	426.05\\
86000	437.46\\
88000	438.06\\
90000	431.87\\
92000	441.88\\
94000	424.55\\
96000	398.06\\
98000	418.25\\
1e+05	424.99\\
};
\addplot [color=mycolor96, forget plot]
  table[row sep=crcr]{%
0	4395.3\\
2000	1757.6\\
4000	588.79\\
6000	478.22\\
8000	457.06\\
10000	463.44\\
12000	477.42\\
14000	446.87\\
16000	453.12\\
18000	465.24\\
20000	444.77\\
22000	444.77\\
24000	455.92\\
26000	442.6\\
28000	431.88\\
30000	426.76\\
32000	452.59\\
34000	460.03\\
36000	441.05\\
38000	444.51\\
40000	469.85\\
42000	429.53\\
44000	464.8\\
46000	447.11\\
48000	442.68\\
50000	435.41\\
52000	467.61\\
54000	462.07\\
56000	454.04\\
58000	420.67\\
60000	445.82\\
62000	450.53\\
64000	485.22\\
66000	456.07\\
68000	449.39\\
70000	458.78\\
72000	434.52\\
74000	470.72\\
76000	466.84\\
78000	437.88\\
80000	451.33\\
82000	470.49\\
84000	455.45\\
86000	456.75\\
88000	441.92\\
90000	460.36\\
92000	467.64\\
94000	448.1\\
96000	421.58\\
98000	456.28\\
1e+05	463.68\\
};
\addplot [color=mycolor97, forget plot]
  table[row sep=crcr]{%
0	4395.3\\
2000	1521.5\\
4000	434.76\\
6000	341.12\\
8000	320.9\\
10000	335.5\\
12000	334.1\\
14000	313.01\\
16000	318.78\\
18000	327.23\\
20000	317.4\\
22000	313.46\\
24000	328.24\\
26000	322.21\\
28000	323.91\\
30000	321.42\\
32000	329.1\\
34000	315.6\\
36000	298.29\\
38000	319.29\\
40000	308.14\\
42000	317.96\\
44000	324.33\\
46000	322.87\\
48000	348.7\\
50000	323.06\\
52000	320.29\\
54000	313.02\\
56000	316.47\\
58000	322.89\\
60000	306.47\\
62000	304.57\\
64000	325.17\\
66000	327.03\\
68000	317.74\\
70000	323.57\\
72000	319.94\\
74000	328.5\\
76000	327.71\\
78000	331.91\\
80000	323.94\\
82000	346.28\\
84000	352.96\\
86000	329.96\\
88000	339.47\\
90000	322.57\\
92000	318.45\\
94000	326.83\\
96000	323.94\\
98000	317.2\\
1e+05	317.22\\
};
\addplot [color=mycolor98]
  table[row sep=crcr]{%
0	4395.3\\
2000	1874.7\\
4000	596.96\\
6000	467.86\\
8000	416.13\\
10000	431.56\\
12000	444.21\\
14000	428.81\\
16000	424.92\\
18000	431.13\\
20000	452.75\\
22000	434.23\\
24000	422.36\\
26000	426.79\\
28000	415.67\\
30000	444.75\\
32000	444.08\\
34000	436.01\\
36000	445.58\\
38000	429.29\\
40000	439.02\\
42000	435.31\\
44000	439.29\\
46000	434.63\\
50000	429.9\\
52000	437.56\\
54000	413.5\\
56000	420.11\\
58000	416.25\\
60000	428.21\\
62000	412.51\\
64000	401.85\\
66000	429.89\\
68000	457.54\\
70000	434.12\\
72000	400.6\\
74000	409.32\\
76000	429.35\\
78000	445.65\\
82000	412.58\\
84000	432.44\\
88000	436.31\\
90000	418\\
92000	414.12\\
94000	435.71\\
96000	432.64\\
98000	450.87\\
1e+05	454.61\\
};
\addlegendentry{Classe k=10}

            % \input{../../../../.tex/20/KS/SizeEvolutionsfig5.tex}
            \legend{}

            \nextgroupplot[%
                % xmin=0,xmax=1e+05,
                % ymin=100,ymax=1.5501e+05,
            ] % This file was created by matlab2tikz.
%
\definecolor{mycolor1}{rgb}{0.00146,0.00047,0.01387}%
\definecolor{mycolor2}{rgb}{0.01166,0.00942,0.06346}%
\definecolor{mycolor3}{rgb}{0.02943,0.02150,0.11462}%
\definecolor{mycolor4}{rgb}{0.06134,0.03659,0.17764}%
\definecolor{mycolor5}{rgb}{0.08741,0.04456,0.22481}%
\definecolor{mycolor6}{rgb}{0.12291,0.04754,0.28162}%
\definecolor{mycolor7}{rgb}{0.14907,0.04547,0.31709}%
\definecolor{mycolor8}{rgb}{0.18343,0.04033,0.35497}%
\definecolor{mycolor9}{rgb}{0.21795,0.03662,0.38352}%
\definecolor{mycolor10}{rgb}{0.24497,0.03705,0.40001}%
\definecolor{mycolor11}{rgb}{0.27135,0.04092,0.41198}%
\definecolor{mycolor12}{rgb}{0.29076,0.04564,0.41864}%
\definecolor{mycolor13}{rgb}{0.31628,0.05349,0.42512}%
\definecolor{mycolor14}{rgb}{0.33522,0.06006,0.42852}%
\definecolor{mycolor15}{rgb}{0.36028,0.06925,0.43150}%
\definecolor{mycolor16}{rgb}{0.37900,0.07625,0.43272}%
\definecolor{mycolor17}{rgb}{0.39767,0.08326,0.43318}%
\definecolor{mycolor18}{rgb}{0.41633,0.09020,0.43294}%
\definecolor{mycolor19}{rgb}{0.42877,0.09479,0.43241}%
\definecolor{mycolor20}{rgb}{0.44743,0.10160,0.43108}%
\definecolor{mycolor21}{rgb}{0.46610,0.10832,0.42912}%
\definecolor{mycolor22}{rgb}{0.47856,0.11276,0.42747}%
\definecolor{mycolor23}{rgb}{0.49726,0.11938,0.42449}%
\definecolor{mycolor24}{rgb}{0.50973,0.12377,0.42216}%
\definecolor{mycolor25}{rgb}{0.52221,0.12815,0.41955}%
\definecolor{mycolor26}{rgb}{0.53468,0.13253,0.41667}%
\definecolor{mycolor27}{rgb}{0.55339,0.13913,0.41183}%
\definecolor{mycolor28}{rgb}{0.56585,0.14357,0.40826}%
\definecolor{mycolor29}{rgb}{0.57830,0.14804,0.40441}%
\definecolor{mycolor30}{rgb}{0.59073,0.15256,0.40029}%
\definecolor{mycolor31}{rgb}{0.60314,0.15715,0.39589}%
\definecolor{mycolor32}{rgb}{0.61551,0.16182,0.39122}%
\definecolor{mycolor33}{rgb}{0.62169,0.16418,0.38878}%
\definecolor{mycolor34}{rgb}{0.63400,0.16899,0.38370}%
\definecolor{mycolor35}{rgb}{0.64626,0.17391,0.37836}%
\definecolor{mycolor36}{rgb}{0.65846,0.17896,0.37275}%
\definecolor{mycolor37}{rgb}{0.66454,0.18154,0.36985}%
\definecolor{mycolor38}{rgb}{0.67664,0.18681,0.36385}%
\definecolor{mycolor39}{rgb}{0.68865,0.19224,0.35760}%
\definecolor{mycolor40}{rgb}{0.69463,0.19502,0.35439}%
\definecolor{mycolor41}{rgb}{0.70650,0.20073,0.34778}%
\definecolor{mycolor42}{rgb}{0.71240,0.20366,0.34438}%
\definecolor{mycolor43}{rgb}{0.72410,0.20967,0.33742}%
\definecolor{mycolor44}{rgb}{0.72991,0.21276,0.33386}%
\definecolor{mycolor45}{rgb}{0.74142,0.21911,0.32658}%
\definecolor{mycolor46}{rgb}{0.74713,0.22238,0.32286}%
\definecolor{mycolor47}{rgb}{0.75279,0.22571,0.31909}%
\definecolor{mycolor48}{rgb}{0.76401,0.23255,0.31140}%
\definecolor{mycolor49}{rgb}{0.76956,0.23608,0.30749}%
\definecolor{mycolor50}{rgb}{0.77506,0.23967,0.30353}%
\definecolor{mycolor51}{rgb}{0.79129,0.25086,0.29139}%
\definecolor{mycolor52}{rgb}{0.79661,0.25473,0.28726}%
\definecolor{mycolor53}{rgb}{0.80187,0.25867,0.28310}%
\definecolor{mycolor54}{rgb}{0.81224,0.26679,0.27466}%
\definecolor{mycolor55}{rgb}{0.81734,0.27095,0.27039}%
\definecolor{mycolor56}{rgb}{0.82737,0.27952,0.26175}%
\definecolor{mycolor57}{rgb}{0.83230,0.28391,0.25738}%
\definecolor{mycolor58}{rgb}{0.83717,0.28839,0.25299}%
\definecolor{mycolor59}{rgb}{0.84671,0.29756,0.24411}%
\definecolor{mycolor60}{rgb}{0.85138,0.30226,0.23964}%
\definecolor{mycolor61}{rgb}{0.85599,0.30704,0.23513}%
\definecolor{mycolor62}{rgb}{0.86053,0.31189,0.23061}%
\definecolor{mycolor63}{rgb}{0.87374,0.32691,0.21689}%
\definecolor{mycolor64}{rgb}{0.87800,0.33206,0.21227}%
\definecolor{mycolor65}{rgb}{0.89034,0.34796,0.19829}%
\definecolor{mycolor66}{rgb}{0.89819,0.35891,0.18886}%
\definecolor{mycolor67}{rgb}{0.90200,0.36449,0.18412}%
\definecolor{mycolor68}{rgb}{0.90573,0.37014,0.17935}%
\definecolor{mycolor69}{rgb}{0.90939,0.37586,0.17456}%
\definecolor{mycolor70}{rgb}{0.91297,0.38164,0.16975}%
\definecolor{mycolor71}{rgb}{0.91646,0.38748,0.16492}%
\definecolor{mycolor72}{rgb}{0.91988,0.39339,0.16007}%
\definecolor{mycolor73}{rgb}{0.92322,0.39936,0.15519}%
\definecolor{mycolor74}{rgb}{0.92647,0.40539,0.15029}%
\definecolor{mycolor75}{rgb}{0.92964,0.41148,0.14537}%
\definecolor{mycolor76}{rgb}{0.93274,0.41763,0.14042}%
\definecolor{mycolor77}{rgb}{0.93575,0.42383,0.13544}%
\definecolor{mycolor78}{rgb}{0.93868,0.43009,0.13044}%
\definecolor{mycolor79}{rgb}{0.94152,0.43640,0.12541}%
\definecolor{mycolor80}{rgb}{0.94429,0.44277,0.12035}%
\definecolor{mycolor81}{rgb}{0.94956,0.45566,0.11016}%
\definecolor{mycolor82}{rgb}{0.95208,0.46218,0.10503}%
\definecolor{mycolor83}{rgb}{0.96129,0.48872,0.08429}%
\definecolor{mycolor84}{rgb}{0.96540,0.50225,0.07386}%
\definecolor{mycolor85}{rgb}{0.96732,0.50908,0.06866}%
\definecolor{mycolor86}{rgb}{0.97259,0.52980,0.05332}%
\definecolor{mycolor87}{rgb}{0.97568,0.54380,0.04362}%
\definecolor{mycolor88}{rgb}{0.98082,0.57221,0.02851}%
\definecolor{mycolor89}{rgb}{0.98189,0.57939,0.02625}%
\definecolor{mycolor90}{rgb}{0.98378,0.59385,0.02377}%
\definecolor{mycolor91}{rgb}{0.98532,0.60842,0.02420}%
\definecolor{mycolor92}{rgb}{0.98696,0.63048,0.03091}%
\definecolor{mycolor93}{rgb}{0.98793,0.66025,0.05175}%
\definecolor{mycolor94}{rgb}{0.98771,0.68281,0.07249}%
\definecolor{mycolor95}{rgb}{0.97750,0.78226,0.18592}%
\definecolor{mycolor96}{rgb}{0.96624,0.83619,0.26153}%
\definecolor{mycolor97}{rgb}{0.96252,0.85148,0.28555}%
\definecolor{mycolor98}{rgb}{0.98836,0.99836,0.64492}%
%
\addplot [color=mycolor1]
  table[row sep=crcr]{%
0	4395.3\\
2000	1.004e+05\\
4000	1.4618e+05\\
6000	1.5049e+05\\
8000	1.5066e+05\\
10000	1.5033e+05\\
14000	1.546e+05\\
16000	1.5476e+05\\
18000	1.52e+05\\
20000	1.5282e+05\\
22000	1.4743e+05\\
24000	1.5103e+05\\
26000	1.5065e+05\\
28000	1.5069e+05\\
30000	1.4879e+05\\
32000	1.5009e+05\\
36000	1.4985e+05\\
38000	1.5208e+05\\
40000	1.516e+05\\
42000	1.4863e+05\\
44000	1.5039e+05\\
46000	1.5023e+05\\
48000	1.4569e+05\\
50000	1.4957e+05\\
52000	1.5007e+05\\
54000	1.4855e+05\\
56000	1.5087e+05\\
58000	1.5071e+05\\
60000	1.4855e+05\\
62000	1.4713e+05\\
64000	1.487e+05\\
66000	1.5354e+05\\
68000	1.5163e+05\\
70000	1.4696e+05\\
72000	1.5122e+05\\
74000	1.5111e+05\\
76000	1.5001e+05\\
78000	1.5214e+05\\
80000	1.495e+05\\
82000	1.4993e+05\\
84000	1.4962e+05\\
86000	1.5501e+05\\
88000	1.5095e+05\\
90000	1.4952e+05\\
92000	1.4882e+05\\
94000	1.4944e+05\\
96000	1.5102e+05\\
98000	1.4991e+05\\
1e+05	1.5291e+05\\
};
\addlegendentry{Classe k=283}

\addplot [color=mycolor2, forget plot]
  table[row sep=crcr]{%
0	4395.3\\
2000	22623\\
4000	25537\\
6000	24940\\
8000	25704\\
10000	25211\\
12000	25652\\
14000	24993\\
16000	25841\\
18000	25595\\
20000	26118\\
22000	25471\\
24000	25787\\
26000	26466\\
28000	25326\\
30000	25440\\
32000	25883\\
34000	25334\\
36000	25895\\
38000	25935\\
40000	25609\\
42000	25830\\
44000	25661\\
46000	25392\\
48000	26174\\
50000	25417\\
52000	25768\\
54000	25462\\
56000	25751\\
58000	24967\\
60000	26564\\
62000	25169\\
64000	25682\\
66000	24824\\
68000	25696\\
70000	25903\\
72000	25588\\
74000	25485\\
80000	26999\\
82000	25393\\
84000	25658\\
86000	25425\\
88000	25616\\
90000	25463\\
92000	26640\\
94000	25336\\
96000	25998\\
98000	25497\\
1e+05	24669\\
};
\addplot [color=mycolor3, forget plot]
  table[row sep=crcr]{%
0	4395.3\\
2000	32234\\
4000	38575\\
6000	38984\\
8000	39209\\
10000	39858\\
12000	39565\\
14000	39109\\
16000	38294\\
18000	38737\\
20000	40201\\
22000	38066\\
24000	38687\\
26000	39774\\
28000	39892\\
30000	38828\\
32000	39582\\
34000	39989\\
36000	39331\\
38000	38869\\
40000	39389\\
42000	39246\\
44000	38107\\
46000	38780\\
48000	39975\\
50000	38334\\
52000	38585\\
56000	39527\\
58000	38890\\
60000	39594\\
62000	38634\\
64000	39387\\
66000	38914\\
68000	39054\\
70000	39312\\
72000	39360\\
74000	39094\\
76000	37452\\
78000	39024\\
82000	38810\\
84000	39087\\
86000	39662\\
88000	39070\\
90000	39346\\
92000	38406\\
94000	40357\\
96000	39041\\
98000	39351\\
1e+05	38913\\
};
\addplot [color=mycolor4, forget plot]
  table[row sep=crcr]{%
0	4395.3\\
2000	13631\\
4000	13990\\
6000	14460\\
8000	14547\\
10000	14420\\
12000	14382\\
14000	13863\\
16000	14105\\
18000	14282\\
20000	14205\\
22000	14022\\
24000	14038\\
26000	14423\\
28000	14038\\
30000	14080\\
32000	14006\\
34000	14242\\
36000	14678\\
38000	13880\\
40000	14277\\
42000	13934\\
44000	13998\\
46000	14313\\
48000	14085\\
50000	14381\\
52000	14088\\
54000	13874\\
56000	14291\\
58000	14464\\
60000	14299\\
62000	13898\\
64000	14529\\
66000	14573\\
68000	14128\\
70000	14461\\
72000	14054\\
74000	14634\\
76000	14356\\
78000	14286\\
80000	14532\\
84000	14364\\
86000	13806\\
88000	13714\\
90000	13987\\
92000	14143\\
94000	14089\\
96000	14206\\
98000	14141\\
1e+05	14121\\
};
\addplot [color=mycolor5, forget plot]
  table[row sep=crcr]{%
0	4395.3\\
2000	45104\\
4000	62380\\
6000	61594\\
8000	61084\\
10000	63064\\
12000	61019\\
14000	63788\\
16000	61577\\
18000	62924\\
20000	63635\\
24000	62646\\
26000	62152\\
28000	63753\\
32000	63164\\
34000	61506\\
36000	63371\\
38000	63608\\
40000	62984\\
42000	61772\\
44000	63217\\
46000	63185\\
48000	63682\\
50000	62083\\
52000	61699\\
54000	62553\\
56000	61891\\
58000	63089\\
60000	63274\\
64000	63375\\
66000	64528\\
68000	62589\\
70000	63854\\
72000	61105\\
76000	65465\\
78000	62226\\
82000	63106\\
84000	64483\\
86000	61939\\
88000	62992\\
90000	63762\\
92000	62974\\
94000	62803\\
96000	62298\\
98000	64712\\
1e+05	62399\\
};
\addplot [color=mycolor6, forget plot]
  table[row sep=crcr]{%
0	4395.3\\
2000	40307\\
4000	54148\\
6000	57285\\
10000	54721\\
12000	54836\\
14000	55333\\
16000	56410\\
18000	53954\\
20000	52778\\
22000	56404\\
24000	55779\\
26000	53180\\
28000	54072\\
30000	55810\\
32000	54257\\
34000	55162\\
36000	55807\\
38000	56179\\
40000	56849\\
42000	55223\\
46000	56200\\
48000	54804\\
50000	55360\\
52000	53544\\
54000	55253\\
56000	55856\\
60000	55692\\
62000	55279\\
64000	53989\\
66000	52955\\
68000	53827\\
70000	56301\\
72000	53956\\
74000	55803\\
76000	55777\\
78000	53187\\
80000	53974\\
82000	54350\\
84000	58031\\
86000	55109\\
88000	54650\\
90000	55249\\
92000	55076\\
94000	54072\\
96000	54685\\
98000	54515\\
1e+05	55528\\
};
\addplot [color=mycolor7, forget plot]
  table[row sep=crcr]{%
0	4395.3\\
2000	22931\\
4000	28286\\
6000	27837\\
8000	28259\\
10000	27839\\
12000	27227\\
14000	27976\\
16000	27350\\
18000	27598\\
20000	27178\\
22000	27859\\
24000	29217\\
26000	27443\\
28000	27781\\
30000	27540\\
32000	27570\\
34000	27184\\
36000	27678\\
38000	27445\\
40000	28055\\
42000	28355\\
44000	27421\\
46000	27925\\
48000	28861\\
50000	28368\\
52000	28455\\
54000	28891\\
56000	27234\\
60000	27981\\
62000	28572\\
64000	27665\\
66000	26993\\
68000	27932\\
70000	27403\\
72000	29005\\
74000	26727\\
76000	28270\\
78000	27514\\
80000	27149\\
82000	28046\\
84000	27461\\
86000	27866\\
88000	29019\\
90000	27414\\
92000	28062\\
94000	27511\\
96000	27697\\
98000	27762\\
1e+05	28304\\
};
\addplot [color=mycolor8, forget plot]
  table[row sep=crcr]{%
0	4395.3\\
2000	17646\\
4000	20329\\
6000	19877\\
8000	20375\\
10000	20627\\
12000	19865\\
14000	19970\\
16000	19976\\
18000	20889\\
20000	20181\\
22000	20404\\
24000	19698\\
26000	19972\\
28000	20660\\
30000	19715\\
32000	20433\\
34000	19756\\
36000	20282\\
38000	20448\\
40000	20105\\
42000	20471\\
44000	20317\\
46000	20787\\
48000	20584\\
50000	20266\\
52000	20806\\
54000	20216\\
56000	21388\\
58000	20344\\
60000	20278\\
62000	20847\\
64000	20904\\
68000	20356\\
70000	19924\\
72000	20050\\
74000	20769\\
76000	20630\\
78000	19858\\
80000	20411\\
82000	20181\\
84000	20334\\
86000	19998\\
90000	20730\\
92000	20838\\
94000	20603\\
96000	20071\\
98000	20950\\
1e+05	20828\\
};
\addplot [color=mycolor9, forget plot]
  table[row sep=crcr]{%
0	4395.3\\
2000	12707\\
4000	14130\\
6000	13716\\
8000	14094\\
10000	14208\\
12000	14707\\
14000	14067\\
16000	14302\\
18000	13940\\
20000	13741\\
22000	13838\\
24000	13871\\
26000	13766\\
30000	13997\\
32000	13698\\
34000	13728\\
36000	13676\\
38000	14004\\
40000	14163\\
42000	13646\\
44000	14070\\
46000	13976\\
48000	14629\\
50000	13942\\
52000	13937\\
54000	13889\\
56000	13644\\
58000	14283\\
60000	13341\\
62000	13912\\
64000	13830\\
66000	13829\\
68000	14262\\
70000	14361\\
72000	14197\\
74000	13952\\
76000	13821\\
78000	14347\\
80000	13753\\
82000	14211\\
84000	13652\\
86000	13794\\
88000	13456\\
90000	13983\\
92000	14360\\
94000	13816\\
96000	13631\\
98000	13371\\
1e+05	14343\\
};
\addplot [color=mycolor10, forget plot]
  table[row sep=crcr]{%
0	4395.3\\
2000	18135\\
4000	21796\\
6000	21779\\
8000	21032\\
10000	21817\\
12000	21848\\
14000	21099\\
16000	21993\\
18000	22079\\
22000	21450\\
24000	21759\\
26000	21962\\
28000	21693\\
30000	22424\\
32000	22278\\
34000	21787\\
36000	21598\\
38000	21338\\
40000	21818\\
42000	21804\\
46000	22480\\
48000	22239\\
50000	21622\\
54000	21819\\
56000	21606\\
60000	21684\\
62000	22141\\
64000	21650\\
66000	21569\\
68000	22196\\
70000	22351\\
72000	21670\\
76000	21425\\
78000	21849\\
80000	22146\\
82000	22093\\
84000	21591\\
86000	22465\\
88000	21940\\
90000	21770\\
92000	21674\\
94000	21455\\
96000	21504\\
98000	21994\\
1e+05	22035\\
};
\addplot [color=mycolor11, forget plot]
  table[row sep=crcr]{%
0	4395.3\\
2000	8119.8\\
4000	8651.7\\
6000	8603.5\\
8000	8780.3\\
10000	8521.6\\
12000	8457.3\\
14000	8674.4\\
16000	8448.8\\
18000	8762.9\\
20000	8363.4\\
22000	8658.1\\
24000	8463.9\\
26000	8416.4\\
28000	8703.9\\
32000	8409.6\\
34000	8812.7\\
36000	8406.8\\
38000	8712.4\\
40000	8614.5\\
42000	8628.9\\
44000	8500.6\\
46000	8689.5\\
48000	8468.2\\
50000	8566.6\\
52000	8369.1\\
54000	8559\\
56000	8439.2\\
58000	8529.4\\
60000	8298.4\\
62000	8420.6\\
64000	8585.8\\
66000	8566.1\\
68000	8239.9\\
70000	8652.8\\
74000	8664.8\\
76000	8645.5\\
78000	8563.4\\
80000	8618\\
82000	8543.5\\
84000	8734.2\\
86000	8603.5\\
88000	8656.2\\
90000	8805.9\\
92000	8881.1\\
94000	8485\\
98000	8333.8\\
1e+05	8715\\
};
\addplot [color=mycolor12, forget plot]
  table[row sep=crcr]{%
0	4395.3\\
2000	13467\\
4000	14510\\
6000	14811\\
8000	15062\\
10000	14696\\
12000	14566\\
14000	14945\\
16000	14364\\
18000	14815\\
20000	14914\\
22000	15058\\
24000	14373\\
26000	14372\\
28000	14447\\
30000	14929\\
32000	14758\\
34000	14710\\
36000	14295\\
38000	14832\\
40000	14423\\
42000	14913\\
44000	15117\\
46000	14517\\
48000	14301\\
50000	14891\\
52000	14441\\
54000	14748\\
56000	14471\\
60000	15104\\
62000	14752\\
64000	14572\\
66000	14300\\
68000	14467\\
70000	14884\\
72000	15176\\
74000	14205\\
76000	14682\\
78000	14897\\
80000	14385\\
82000	14758\\
84000	14780\\
86000	14363\\
88000	15028\\
90000	14382\\
92000	14208\\
94000	14622\\
96000	14749\\
98000	14993\\
1e+05	14257\\
};
\addplot [color=mycolor13, forget plot]
  table[row sep=crcr]{%
0	4395.3\\
2000	10651\\
4000	11137\\
6000	11096\\
8000	11139\\
10000	11394\\
12000	11146\\
14000	11150\\
16000	10947\\
18000	10836\\
22000	10991\\
24000	10820\\
26000	11294\\
28000	11086\\
30000	10926\\
32000	11235\\
34000	11077\\
36000	10838\\
38000	11244\\
40000	11303\\
42000	11260\\
44000	11360\\
46000	11425\\
48000	11247\\
50000	10857\\
52000	10833\\
56000	10986\\
58000	11234\\
60000	11050\\
62000	11310\\
64000	11518\\
66000	10887\\
68000	11330\\
70000	10936\\
72000	11042\\
74000	10785\\
76000	11137\\
78000	11204\\
80000	11378\\
82000	11018\\
84000	11114\\
86000	11015\\
88000	11211\\
90000	10961\\
92000	11161\\
94000	10770\\
96000	11269\\
98000	10950\\
1e+05	11439\\
};
\addplot [color=mycolor14, forget plot]
  table[row sep=crcr]{%
0	4395.3\\
2000	20122\\
4000	23422\\
10000	24080\\
12000	24437\\
14000	24515\\
16000	24209\\
18000	23653\\
20000	24248\\
22000	24492\\
24000	24081\\
26000	23472\\
28000	24101\\
30000	23883\\
32000	24269\\
34000	24282\\
36000	24012\\
38000	24405\\
42000	24119\\
44000	23728\\
46000	23990\\
48000	24541\\
50000	24147\\
52000	24539\\
54000	23611\\
56000	23508\\
58000	23898\\
60000	24471\\
62000	24463\\
64000	23636\\
66000	23671\\
68000	23932\\
70000	23228\\
72000	24102\\
74000	24349\\
76000	23881\\
78000	24252\\
82000	24240\\
84000	24605\\
86000	24123\\
90000	24388\\
92000	24314\\
94000	24309\\
96000	23852\\
98000	23847\\
1e+05	24522\\
};
\addplot [color=mycolor15, forget plot]
  table[row sep=crcr]{%
0	4395.3\\
2000	14719\\
4000	16870\\
6000	17035\\
8000	16896\\
10000	16388\\
12000	17016\\
14000	17468\\
16000	16928\\
18000	16922\\
20000	17636\\
22000	16809\\
24000	16698\\
26000	16945\\
28000	17049\\
30000	17002\\
32000	16769\\
34000	16799\\
36000	16703\\
38000	17238\\
40000	16269\\
42000	16694\\
44000	16719\\
46000	16521\\
48000	17027\\
50000	16885\\
52000	17261\\
54000	16941\\
56000	17308\\
58000	16337\\
60000	16939\\
62000	16969\\
64000	16830\\
66000	16974\\
68000	16536\\
70000	16890\\
72000	17187\\
74000	17163\\
76000	17194\\
78000	16502\\
80000	17265\\
82000	16851\\
84000	16139\\
88000	16661\\
90000	16861\\
92000	16680\\
96000	16890\\
98000	16340\\
1e+05	16863\\
};
\addplot [color=mycolor16, forget plot]
  table[row sep=crcr]{%
0	4395.3\\
2000	6820.9\\
4000	6760.6\\
6000	7019.7\\
8000	6813.5\\
10000	7050.6\\
12000	6936.7\\
14000	6926.9\\
16000	6966.7\\
18000	6823.9\\
20000	6788.9\\
22000	7001.5\\
24000	6681.8\\
26000	7273\\
28000	6989\\
30000	7137.9\\
34000	6913.7\\
36000	7048.4\\
40000	7264.9\\
42000	7089.8\\
44000	7084.4\\
46000	6789.7\\
48000	7106.6\\
50000	6943.3\\
52000	6872.2\\
54000	7104.2\\
56000	6828.1\\
58000	7095.4\\
60000	6930.3\\
62000	6998.5\\
64000	6961.3\\
66000	6863.3\\
68000	6857.9\\
70000	7376.3\\
72000	6985.3\\
74000	6920.8\\
76000	6781.4\\
78000	6986.2\\
80000	6947.9\\
82000	7062.6\\
84000	7042.6\\
86000	6841\\
88000	7072.8\\
90000	7101.4\\
92000	7215.4\\
94000	7061.5\\
96000	6827.5\\
98000	6882.6\\
1e+05	6961\\
};
\addplot [color=mycolor17, forget plot]
  table[row sep=crcr]{%
0	4395.3\\
2000	13425\\
4000	15798\\
6000	16103\\
8000	15914\\
10000	16340\\
14000	15801\\
16000	15896\\
18000	15629\\
20000	16272\\
22000	15573\\
24000	15601\\
26000	16174\\
28000	15796\\
30000	16061\\
32000	16246\\
34000	15880\\
36000	16285\\
38000	16041\\
40000	16452\\
42000	16593\\
44000	16431\\
46000	16032\\
48000	16016\\
50000	16362\\
52000	15955\\
54000	15432\\
56000	16291\\
58000	16100\\
60000	16096\\
62000	15502\\
64000	16028\\
68000	15678\\
70000	16207\\
72000	15445\\
74000	15929\\
76000	16149\\
78000	15747\\
80000	16042\\
82000	15818\\
84000	15899\\
86000	16078\\
88000	16202\\
90000	16439\\
92000	16066\\
94000	15896\\
96000	15532\\
98000	15852\\
1e+05	15899\\
};
\addplot [color=mycolor18, forget plot]
  table[row sep=crcr]{%
0	4395.3\\
2000	19400\\
4000	24378\\
6000	24065\\
8000	25051\\
10000	24187\\
12000	23644\\
14000	23835\\
16000	24644\\
18000	23870\\
20000	24454\\
24000	25340\\
28000	24199\\
30000	24068\\
32000	24645\\
34000	24259\\
36000	23654\\
38000	24008\\
40000	23812\\
42000	24761\\
44000	23393\\
46000	24802\\
48000	24957\\
50000	23990\\
52000	24421\\
54000	24196\\
56000	24514\\
58000	25085\\
60000	24705\\
62000	24964\\
64000	25142\\
66000	24760\\
68000	24493\\
70000	24370\\
72000	23662\\
74000	23135\\
76000	23737\\
78000	24521\\
80000	23920\\
82000	24595\\
84000	24870\\
86000	24414\\
88000	24530\\
90000	24196\\
92000	23645\\
94000	24302\\
96000	24226\\
98000	24730\\
1e+05	23753\\
};
\addplot [color=mycolor19, forget plot]
  table[row sep=crcr]{%
0	4395.3\\
2000	13276\\
4000	14380\\
6000	14743\\
8000	14340\\
10000	14787\\
12000	15595\\
14000	14634\\
16000	15386\\
18000	15285\\
20000	14903\\
22000	14881\\
24000	14754\\
26000	14550\\
28000	14639\\
30000	14855\\
32000	14783\\
34000	15024\\
36000	14604\\
38000	14662\\
40000	14303\\
42000	14599\\
44000	14839\\
46000	14735\\
48000	14854\\
50000	14533\\
52000	14924\\
54000	14313\\
56000	15041\\
58000	14917\\
60000	14998\\
62000	14827\\
64000	14901\\
66000	14903\\
68000	14957\\
70000	15348\\
72000	14293\\
74000	15099\\
76000	14722\\
78000	14769\\
80000	14519\\
82000	14486\\
84000	14139\\
86000	15154\\
88000	14390\\
92000	14422\\
94000	14885\\
96000	14719\\
98000	14408\\
1e+05	14772\\
};
\addplot [color=mycolor20, forget plot]
  table[row sep=crcr]{%
0	4395.3\\
2000	15822\\
4000	18973\\
6000	19482\\
8000	18771\\
10000	18449\\
12000	19204\\
14000	18908\\
16000	19055\\
18000	19676\\
20000	19549\\
22000	19072\\
24000	19232\\
26000	18882\\
28000	19071\\
30000	18623\\
32000	19160\\
36000	19131\\
38000	18755\\
40000	19427\\
42000	18636\\
44000	19187\\
46000	19597\\
48000	19305\\
50000	19087\\
52000	19155\\
54000	19566\\
56000	19471\\
58000	19163\\
60000	19449\\
62000	19424\\
64000	19512\\
66000	19252\\
68000	19200\\
70000	19422\\
72000	19456\\
74000	19712\\
76000	19242\\
78000	18987\\
80000	18879\\
82000	19212\\
84000	19690\\
86000	18288\\
88000	18951\\
90000	19267\\
92000	18592\\
94000	19826\\
96000	19844\\
98000	18977\\
1e+05	18326\\
};
\addplot [color=mycolor21, forget plot]
  table[row sep=crcr]{%
0	4395.3\\
2000	8975.2\\
4000	9646.7\\
6000	9642.4\\
8000	9666.8\\
10000	9936.7\\
12000	9818.7\\
14000	9449.5\\
16000	9512.1\\
18000	9763.1\\
20000	9476.4\\
22000	9882.2\\
26000	9691.4\\
28000	9463.2\\
32000	9548.4\\
34000	9648.8\\
36000	9496.9\\
38000	9585.2\\
40000	9472.6\\
42000	9662.4\\
48000	9595.3\\
50000	9421.4\\
52000	9696.2\\
54000	9451.3\\
56000	9457.2\\
58000	9753.3\\
60000	9716.9\\
62000	9582.8\\
64000	9696.1\\
66000	9354.7\\
68000	9860.5\\
70000	9794.4\\
72000	9816\\
74000	9758.6\\
76000	9606\\
78000	9691.9\\
80000	9718.9\\
82000	9518.4\\
84000	9672.3\\
86000	9601.9\\
90000	9739.4\\
92000	9458.2\\
94000	9693.2\\
96000	9718.5\\
98000	9355\\
1e+05	9595.4\\
};
\addplot [color=mycolor22, forget plot]
  table[row sep=crcr]{%
0	4395.3\\
2000	19819\\
4000	24601\\
6000	24853\\
8000	25435\\
10000	25312\\
12000	25368\\
14000	25011\\
16000	25066\\
18000	25481\\
20000	24736\\
22000	25808\\
24000	25990\\
26000	24949\\
28000	25302\\
30000	25490\\
32000	25190\\
34000	24767\\
36000	24587\\
40000	24990\\
42000	26101\\
44000	24884\\
50000	25731\\
52000	25032\\
54000	25457\\
56000	25642\\
58000	24945\\
60000	24810\\
62000	26114\\
64000	24755\\
66000	26077\\
68000	24880\\
72000	25444\\
74000	25556\\
76000	25054\\
78000	25236\\
80000	24671\\
82000	25306\\
84000	24279\\
86000	24690\\
88000	25574\\
90000	25331\\
92000	26189\\
94000	25516\\
96000	25432\\
98000	24326\\
1e+05	24569\\
};
\addplot [color=mycolor23, forget plot]
  table[row sep=crcr]{%
0	4395.3\\
2000	10978\\
4000	12408\\
8000	12110\\
10000	11989\\
12000	12163\\
14000	12238\\
16000	12149\\
18000	12256\\
20000	12214\\
22000	12323\\
24000	12264\\
26000	12509\\
28000	12481\\
30000	11997\\
32000	12001\\
34000	12165\\
36000	12046\\
38000	12055\\
40000	12187\\
46000	12390\\
48000	12209\\
50000	12438\\
52000	12159\\
54000	12430\\
56000	12309\\
58000	12283\\
60000	12197\\
64000	12355\\
68000	12226\\
70000	12229\\
72000	12102\\
74000	12105\\
76000	12149\\
78000	12034\\
80000	12263\\
84000	12112\\
86000	12178\\
88000	12138\\
90000	12228\\
96000	12285\\
98000	12376\\
1e+05	12134\\
};
\addplot [color=mycolor24, forget plot]
  table[row sep=crcr]{%
0	4395.3\\
2000	5071.7\\
4000	4800.9\\
6000	4882.5\\
8000	4880.9\\
10000	4828.1\\
12000	4857.8\\
14000	4757.7\\
16000	4883.3\\
18000	4861.8\\
20000	4736.2\\
22000	4869.9\\
24000	4859.1\\
26000	4797.5\\
28000	4701\\
30000	4884.8\\
32000	4892.1\\
34000	4754.4\\
36000	4816.6\\
38000	4859.6\\
40000	4716.2\\
42000	4844.6\\
44000	4724.8\\
46000	4641.4\\
48000	4604.3\\
50000	4715.4\\
52000	4738.9\\
54000	4846.6\\
56000	4789.2\\
58000	4770.9\\
60000	4836.2\\
62000	4780.6\\
64000	4887.1\\
66000	4866.8\\
68000	4748.4\\
70000	4711.1\\
72000	4809.7\\
74000	4738.8\\
76000	4790.7\\
78000	4659.7\\
80000	4852.8\\
82000	4826.6\\
84000	4784.7\\
86000	4705.9\\
88000	4763.6\\
90000	4914.1\\
92000	4959.4\\
94000	4849.5\\
96000	4771.2\\
98000	4760.4\\
1e+05	4840.6\\
};
\addplot [color=mycolor25, forget plot]
  table[row sep=crcr]{%
0	4395.3\\
2000	11056\\
4000	12122\\
6000	12243\\
8000	12655\\
10000	12640\\
12000	12376\\
14000	12348\\
16000	11995\\
18000	12479\\
20000	12373\\
22000	11995\\
24000	12201\\
26000	12256\\
30000	12506\\
32000	12406\\
34000	11968\\
36000	12686\\
38000	12318\\
40000	12364\\
42000	12731\\
46000	12313\\
48000	12604\\
50000	12221\\
52000	12656\\
54000	12582\\
56000	12745\\
58000	12561\\
60000	11829\\
62000	12446\\
64000	12603\\
66000	12585\\
70000	12256\\
72000	12156\\
74000	12113\\
76000	12454\\
78000	12389\\
82000	12485\\
84000	12302\\
86000	12505\\
88000	11839\\
90000	11773\\
92000	11997\\
94000	12034\\
96000	12372\\
98000	12239\\
1e+05	12476\\
};
\addplot [color=mycolor26, forget plot]
  table[row sep=crcr]{%
0	4395.3\\
2000	7705.2\\
4000	8196.6\\
6000	8298.5\\
8000	8032.8\\
10000	8262.7\\
12000	8235.2\\
14000	8064.2\\
16000	8462.2\\
18000	8270.2\\
22000	8452.5\\
24000	8471.1\\
26000	8387.5\\
28000	8045.2\\
30000	8172.3\\
32000	8325.9\\
34000	8223.5\\
36000	7900.9\\
38000	8046.2\\
40000	8126.6\\
42000	8078.4\\
44000	8057.1\\
46000	8333.4\\
50000	8216.6\\
52000	8320.8\\
54000	8345.3\\
56000	8200.2\\
58000	8301.5\\
60000	8171.5\\
62000	8159.3\\
64000	8367.8\\
66000	8451.6\\
68000	8604.7\\
70000	8336.1\\
72000	8141.3\\
74000	8079.4\\
76000	7993.2\\
78000	8067.4\\
80000	8350.8\\
82000	8159.6\\
84000	8228.5\\
86000	8090.3\\
88000	8072.9\\
90000	8217.3\\
92000	7948\\
94000	7838.5\\
96000	8340.4\\
98000	8215.2\\
1e+05	8045.7\\
};
\addplot [color=mycolor27, forget plot]
  table[row sep=crcr]{%
0	4395.3\\
2000	9132.8\\
4000	9619.8\\
6000	10022\\
8000	10075\\
10000	9642.1\\
12000	9752.7\\
14000	9928.7\\
16000	9820.8\\
18000	9889.5\\
24000	9627.4\\
26000	9811.4\\
28000	9616.5\\
30000	9897.1\\
32000	9788.4\\
34000	9787.3\\
38000	9916.4\\
40000	9682.6\\
44000	9724.9\\
46000	9673.5\\
48000	9760.8\\
50000	10011\\
52000	9644.4\\
54000	9788.7\\
56000	9848.1\\
58000	9672.8\\
60000	9727.7\\
62000	9411.4\\
64000	9874.3\\
66000	9668.4\\
68000	9826.4\\
70000	9508.3\\
72000	9718.5\\
74000	10175\\
76000	10030\\
78000	9632.4\\
80000	9448.1\\
82000	9626.5\\
84000	9774.2\\
86000	9762.4\\
88000	9886.1\\
90000	9743.1\\
92000	9574.6\\
94000	9738\\
98000	9543.7\\
1e+05	9698.1\\
};
\addplot [color=mycolor28, forget plot]
  table[row sep=crcr]{%
0	4395.3\\
2000	5412.3\\
6000	5565.8\\
8000	5522.2\\
10000	5641.6\\
12000	5692.6\\
14000	5276.9\\
16000	5482.6\\
18000	5410.5\\
20000	5446.2\\
22000	5649.3\\
24000	5641.2\\
26000	5604.8\\
30000	5476.5\\
32000	5562.3\\
34000	5555.8\\
36000	5504.5\\
38000	5477.8\\
40000	5374.6\\
42000	5476.4\\
44000	5542.2\\
46000	5490.1\\
48000	5490.1\\
50000	5405.6\\
52000	5742.2\\
54000	5570.1\\
58000	5422.5\\
60000	5532.4\\
62000	5619.6\\
64000	5464.8\\
66000	5521\\
68000	5427.7\\
70000	5424.2\\
72000	5616.5\\
74000	5471.7\\
76000	5434.4\\
78000	5599.7\\
80000	5638.6\\
82000	5399.5\\
84000	5563.3\\
86000	5465.6\\
88000	5547.5\\
90000	5430.2\\
92000	5458.1\\
94000	5333.4\\
96000	5652.9\\
98000	5460.3\\
1e+05	5554.2\\
};
\addplot [color=mycolor29, forget plot]
  table[row sep=crcr]{%
0	4395.3\\
2000	10423\\
4000	11410\\
6000	11190\\
8000	11785\\
12000	11479\\
14000	11281\\
18000	11552\\
20000	11650\\
22000	11882\\
24000	11792\\
26000	11432\\
28000	11230\\
30000	11684\\
32000	11959\\
34000	11818\\
36000	11470\\
38000	11536\\
40000	11200\\
42000	11278\\
44000	11317\\
46000	11682\\
48000	11415\\
50000	11287\\
52000	11869\\
54000	11268\\
56000	11688\\
58000	11510\\
60000	11302\\
62000	11448\\
64000	11718\\
66000	11705\\
68000	11126\\
70000	11422\\
72000	11545\\
74000	11938\\
76000	11662\\
78000	11585\\
80000	11726\\
84000	11590\\
86000	11438\\
88000	11708\\
90000	11542\\
94000	11731\\
96000	11118\\
98000	11429\\
1e+05	11675\\
};
\addplot [color=mycolor30, forget plot]
  table[row sep=crcr]{%
0	4395.3\\
2000	7439.4\\
4000	7752.3\\
6000	7863.9\\
8000	7908.3\\
10000	7998\\
12000	7726.7\\
14000	7926.6\\
16000	7787.6\\
18000	8027.4\\
20000	8148\\
22000	8068.9\\
24000	7807.2\\
26000	7826\\
28000	7725.8\\
30000	7885.8\\
34000	8054.3\\
38000	7798.1\\
40000	8118.7\\
42000	7967.7\\
44000	8053.3\\
46000	8103.2\\
48000	7978.6\\
50000	7952.1\\
52000	8163.6\\
54000	7868.9\\
56000	7758.1\\
58000	7782.3\\
62000	7909.5\\
64000	7972.4\\
66000	7784.5\\
68000	7984.1\\
72000	7859.7\\
74000	7877.1\\
76000	7796.2\\
78000	7854.6\\
80000	7958\\
82000	7959.6\\
84000	7839.6\\
86000	7878.3\\
88000	7968.3\\
90000	8235.2\\
92000	7817.7\\
94000	8016.7\\
96000	7789.2\\
98000	7961.4\\
1e+05	7784.5\\
};
\addplot [color=mycolor31, forget plot]
  table[row sep=crcr]{%
0	4395.3\\
2000	6316.8\\
4000	6625.6\\
6000	6402.8\\
8000	6311.8\\
10000	6482.4\\
12000	6446.5\\
14000	6518.1\\
16000	6471.1\\
18000	6563.5\\
20000	6365\\
22000	6470.3\\
26000	6394.8\\
28000	6291.6\\
30000	6731.7\\
32000	6405.1\\
34000	6515\\
36000	6419.9\\
40000	6444.1\\
42000	6442.6\\
44000	6475.3\\
48000	6241.1\\
50000	6639.6\\
52000	6589.3\\
54000	6608.7\\
56000	6341.3\\
58000	6556.8\\
60000	6479\\
62000	6534.9\\
64000	6395.4\\
66000	6213\\
68000	6438.6\\
70000	6457.5\\
72000	6398.7\\
74000	6794\\
76000	6352.9\\
78000	6301.4\\
82000	6718\\
84000	6364.3\\
86000	6580.2\\
88000	6527\\
90000	6539.9\\
92000	6459.3\\
94000	6316.2\\
96000	6595.8\\
98000	6434.2\\
1e+05	6534\\
};
\addplot [color=mycolor32, forget plot]
  table[row sep=crcr]{%
0	4395.3\\
2000	9622.7\\
4000	10333\\
6000	10513\\
8000	10397\\
12000	10578\\
14000	10266\\
16000	10448\\
18000	10228\\
20000	10129\\
22000	10516\\
24000	10259\\
26000	10212\\
28000	10346\\
30000	10371\\
32000	10619\\
34000	10540\\
36000	10517\\
38000	10655\\
40000	10435\\
42000	10476\\
44000	10468\\
46000	10697\\
48000	10254\\
50000	11086\\
52000	10226\\
54000	10488\\
56000	10216\\
58000	10134\\
60000	10331\\
62000	10727\\
64000	10250\\
66000	10148\\
68000	10304\\
70000	10222\\
72000	10406\\
74000	10348\\
76000	10535\\
78000	10359\\
80000	10307\\
82000	10489\\
84000	10581\\
86000	10393\\
88000	10579\\
90000	10671\\
92000	10519\\
94000	10834\\
96000	10509\\
98000	10699\\
1e+05	10698\\
};
\addplot [color=mycolor33, forget plot]
  table[row sep=crcr]{%
0	4395.3\\
2000	7653.5\\
4000	8536.7\\
6000	8449.7\\
8000	8311.2\\
10000	8320.4\\
12000	8103.5\\
14000	8126.3\\
16000	8222.8\\
18000	8369\\
20000	8036.4\\
22000	8056.8\\
24000	8356.8\\
26000	8305.2\\
28000	8160.3\\
30000	8272.9\\
32000	8272.7\\
34000	8405.2\\
36000	8446.5\\
38000	8235\\
40000	7898.9\\
42000	8054.8\\
44000	7857.8\\
46000	7991.5\\
48000	7909\\
50000	8388.4\\
52000	8331\\
54000	8071.5\\
56000	8080.2\\
58000	8217.9\\
60000	8003.5\\
62000	8341.8\\
64000	8382.4\\
70000	7968\\
72000	8266.8\\
74000	8355.6\\
78000	8163.5\\
80000	8219.1\\
82000	8047.7\\
84000	8223.1\\
86000	8264.9\\
88000	8170.6\\
90000	8034.7\\
92000	8057.5\\
96000	8489.5\\
98000	8094.7\\
1e+05	8300.9\\
};
\addplot [color=mycolor34, forget plot]
  table[row sep=crcr]{%
0	4395.3\\
2000	4206.5\\
4000	3953.3\\
6000	3964.9\\
8000	4057.9\\
10000	3991.4\\
12000	4230.4\\
16000	4063.4\\
18000	4116.7\\
20000	4120.4\\
24000	3979.7\\
26000	3966\\
28000	4116.4\\
30000	3873.3\\
32000	3964.3\\
34000	4178.6\\
36000	3938.5\\
38000	3975.1\\
40000	3986\\
42000	4180.8\\
44000	4066.5\\
48000	3968.4\\
50000	4066.3\\
52000	3871.7\\
54000	4097.6\\
56000	4061\\
58000	3926.8\\
60000	3956.4\\
62000	4195.6\\
64000	4008\\
66000	4040\\
68000	4042.2\\
70000	4061.2\\
72000	3948\\
74000	4094.1\\
76000	4061\\
78000	4004.7\\
80000	4095.2\\
82000	3970.2\\
84000	4015.6\\
88000	4014.3\\
90000	4054.6\\
92000	4005.1\\
94000	3898.2\\
96000	4019.7\\
98000	4178.7\\
1e+05	3989.1\\
};
\addplot [color=mycolor35]
  table[row sep=crcr]{%
0	4395.3\\
2000	9172.4\\
4000	10016\\
6000	10206\\
8000	10316\\
10000	10139\\
12000	10238\\
14000	9887.4\\
16000	10182\\
18000	10007\\
20000	10072\\
22000	10054\\
24000	10312\\
26000	10030\\
28000	10026\\
30000	10207\\
32000	10040\\
34000	10134\\
36000	10174\\
38000	10004\\
40000	10152\\
42000	9950.9\\
44000	10198\\
46000	10049\\
48000	10031\\
50000	10235\\
52000	10177\\
54000	10377\\
56000	9861.7\\
58000	10033\\
60000	10264\\
62000	10284\\
64000	10219\\
66000	10114\\
68000	10170\\
70000	10081\\
72000	10156\\
76000	9987.4\\
78000	10109\\
80000	10079\\
84000	9778\\
88000	10184\\
90000	10439\\
92000	10178\\
94000	10110\\
96000	10342\\
1e+05	9863.5\\
};
\addlegendentry{Classe k=83}

\addplot [color=mycolor36, forget plot]
  table[row sep=crcr]{%
0	4395.3\\
2000	4576\\
4000	4530\\
6000	4265.6\\
8000	4305.1\\
10000	4130.9\\
12000	4454.7\\
14000	4492.3\\
16000	4322.1\\
18000	4202.3\\
20000	4336.1\\
22000	4392.3\\
24000	4098.3\\
26000	4199.8\\
30000	4277.8\\
32000	4245.3\\
34000	4463\\
36000	4284.9\\
38000	4159.4\\
40000	4289.5\\
42000	4382.1\\
44000	4134.7\\
46000	4303.3\\
48000	4402.1\\
50000	4242.2\\
52000	4365.9\\
54000	4167\\
56000	4207.8\\
58000	4234.6\\
60000	4205.1\\
62000	4346.7\\
64000	4157.6\\
66000	4232.8\\
68000	4343.3\\
70000	4434.6\\
72000	4316.8\\
74000	4442.3\\
76000	4272.4\\
78000	4276.1\\
80000	4180.5\\
82000	4382.7\\
84000	4259.7\\
86000	4433.4\\
88000	4291.4\\
90000	4322.2\\
92000	4462.9\\
94000	4407.7\\
96000	4383.6\\
1e+05	4232.7\\
};
\addplot [color=mycolor37, forget plot]
  table[row sep=crcr]{%
0	4395.3\\
2000	3878.7\\
4000	3728.1\\
6000	3767\\
8000	3718.3\\
10000	3650.9\\
12000	3662.1\\
14000	3569.9\\
16000	3677.2\\
18000	3609.9\\
20000	3650.5\\
22000	3714.7\\
24000	3656.1\\
26000	3666.1\\
28000	3613.6\\
30000	3678\\
32000	3655\\
34000	3664.3\\
36000	3715.5\\
38000	3729.7\\
40000	3654\\
42000	3675.7\\
44000	3541.2\\
46000	3570.7\\
48000	3686.6\\
50000	3659.7\\
52000	3663.5\\
54000	3562.4\\
56000	3602.3\\
58000	3734.7\\
60000	3664.5\\
62000	3720.8\\
64000	3658.5\\
66000	3631.1\\
68000	3634\\
70000	3578.8\\
72000	3625.8\\
74000	3578\\
76000	3668\\
78000	3718.5\\
80000	3641.9\\
82000	3740.4\\
84000	3645.8\\
86000	3602.5\\
88000	3539.4\\
90000	3637.6\\
92000	3727.3\\
94000	3611\\
96000	3709.5\\
98000	3663.4\\
1e+05	3583.4\\
};
\addplot [color=mycolor38, forget plot]
  table[row sep=crcr]{%
0	4395.3\\
2000	3283.8\\
4000	2896.5\\
6000	3034.8\\
8000	3023.1\\
10000	3115.6\\
14000	2915.7\\
16000	3106.6\\
18000	2989\\
20000	3004.3\\
22000	3008.6\\
24000	3055.2\\
26000	3079.3\\
28000	2993.6\\
30000	2980.5\\
32000	3038.9\\
34000	3079.3\\
36000	3075.6\\
38000	2982.4\\
40000	3058.3\\
42000	3099.7\\
44000	2959.8\\
46000	2914.5\\
48000	3013\\
50000	3002\\
52000	2950.6\\
54000	2963.3\\
56000	2940.9\\
58000	3080.6\\
60000	3042.2\\
62000	2984.2\\
64000	2942.1\\
66000	3081.7\\
68000	3033.4\\
70000	2972.8\\
72000	2904.1\\
74000	3043.3\\
76000	3028.3\\
78000	3004.4\\
80000	3031.8\\
82000	2985.2\\
84000	2905.3\\
86000	3038\\
88000	3012.6\\
90000	3019\\
92000	3005.6\\
94000	3064\\
96000	3136.9\\
98000	3158.7\\
1e+05	3011.7\\
};
\addplot [color=mycolor39, forget plot]
  table[row sep=crcr]{%
0	4395.3\\
2000	10468\\
4000	11648\\
8000	11454\\
10000	11596\\
12000	11604\\
14000	11405\\
16000	11431\\
18000	11872\\
20000	11754\\
22000	12068\\
24000	11865\\
26000	11827\\
28000	11902\\
30000	12019\\
32000	11991\\
34000	11783\\
36000	12011\\
38000	11701\\
40000	11332\\
42000	11864\\
44000	12010\\
46000	11949\\
48000	12009\\
50000	11663\\
52000	12097\\
54000	11834\\
56000	12203\\
58000	11862\\
60000	11440\\
64000	11591\\
66000	11846\\
68000	11888\\
70000	12064\\
72000	11947\\
74000	11581\\
76000	11673\\
78000	11625\\
80000	11660\\
82000	11353\\
84000	11560\\
86000	11717\\
88000	11948\\
90000	11759\\
92000	12121\\
94000	11436\\
96000	11694\\
98000	11855\\
1e+05	12272\\
};
\addplot [color=mycolor40, forget plot]
  table[row sep=crcr]{%
0	4395.3\\
2000	5599.8\\
4000	5941.7\\
8000	5750.8\\
10000	5803.5\\
12000	5809.1\\
14000	5773.4\\
16000	5827.5\\
18000	5817.3\\
20000	5775.9\\
22000	5710.3\\
24000	5782.5\\
30000	5771.1\\
34000	5663.3\\
36000	5676.8\\
40000	5879.8\\
42000	5626.5\\
44000	5706.1\\
48000	5771.6\\
50000	5851.8\\
52000	5686.2\\
54000	5751.3\\
56000	5565.7\\
58000	5825.2\\
60000	5781.7\\
64000	5825.6\\
68000	5759.1\\
70000	5608.3\\
72000	5813.5\\
74000	5809.7\\
76000	5744.5\\
78000	5845.6\\
80000	5805\\
82000	5706.1\\
84000	5790.6\\
88000	5700.3\\
90000	5792.7\\
92000	5630.8\\
96000	5691.7\\
98000	5817.1\\
1e+05	5751.1\\
};
\addplot [color=mycolor41, forget plot]
  table[row sep=crcr]{%
0	4395.3\\
2000	5210.5\\
4000	5557.1\\
6000	5471.5\\
8000	5311.3\\
10000	5396.3\\
12000	5429.4\\
14000	5527.6\\
16000	5309.7\\
18000	5248.6\\
20000	5526.8\\
22000	5367.3\\
24000	5188.9\\
26000	5381\\
28000	5331.3\\
30000	5575.1\\
32000	5352.9\\
34000	5279.9\\
36000	5330.8\\
38000	5448\\
40000	5369.4\\
42000	5489\\
44000	5459.1\\
46000	5282.4\\
48000	5342.2\\
50000	5511.1\\
52000	5467.4\\
54000	5392.6\\
56000	5286.6\\
58000	5305.9\\
60000	5492.4\\
62000	5434.7\\
64000	5257.5\\
66000	5372.2\\
68000	5398.6\\
70000	5355.2\\
72000	5394.9\\
74000	5389\\
76000	5416.7\\
78000	5577.5\\
80000	5391.2\\
82000	5305.1\\
84000	5359.6\\
86000	5387.3\\
88000	5263.2\\
90000	5500.7\\
92000	5435.6\\
94000	5423.7\\
96000	5501.3\\
98000	5413.7\\
1e+05	5403.1\\
};
\addplot [color=mycolor42, forget plot]
  table[row sep=crcr]{%
0	4395.3\\
2000	15538\\
4000	20570\\
6000	20461\\
8000	20588\\
10000	20774\\
12000	20709\\
14000	19950\\
16000	20327\\
18000	21373\\
20000	19797\\
22000	19890\\
24000	20057\\
26000	20723\\
28000	21139\\
30000	20959\\
32000	20113\\
36000	20767\\
38000	19760\\
40000	20260\\
42000	20702\\
44000	20445\\
46000	19786\\
48000	20379\\
50000	20093\\
52000	20154\\
54000	20695\\
56000	20097\\
58000	20724\\
60000	20316\\
62000	20758\\
64000	19668\\
66000	20282\\
68000	19994\\
70000	20403\\
72000	20721\\
74000	20471\\
76000	20569\\
78000	20416\\
80000	21136\\
82000	20595\\
84000	20168\\
86000	20692\\
88000	20351\\
90000	20897\\
92000	20480\\
94000	20771\\
96000	20462\\
98000	20537\\
1e+05	20704\\
};
\addplot [color=mycolor43, forget plot]
  table[row sep=crcr]{%
0	4395.3\\
2000	4992.3\\
4000	4927.2\\
6000	4921.2\\
8000	4839.2\\
12000	4862.3\\
14000	4921.2\\
16000	4824.4\\
18000	4890.5\\
20000	4913.1\\
22000	5134.3\\
24000	4890.9\\
26000	4927.7\\
28000	5105\\
30000	4882.6\\
32000	4882\\
34000	4966.7\\
36000	4986.8\\
38000	4967.8\\
40000	5002\\
42000	4957.8\\
44000	4881\\
46000	4851.1\\
48000	5001\\
50000	5037.9\\
52000	5011\\
54000	4925.4\\
56000	5069.6\\
58000	4910.6\\
62000	4826.6\\
64000	4949.9\\
66000	4863.3\\
68000	4925.6\\
70000	4943.8\\
72000	5009.7\\
74000	4963.7\\
76000	4891.5\\
78000	5024.9\\
80000	5026.3\\
82000	4863.5\\
84000	4980.3\\
86000	4932.7\\
90000	4939\\
92000	4894.9\\
94000	4967.7\\
96000	4973.7\\
98000	5013.7\\
1e+05	4952.6\\
};
\addplot [color=mycolor44, forget plot]
  table[row sep=crcr]{%
0	4395.3\\
2000	3532.4\\
4000	3350.9\\
6000	3321.6\\
8000	3422\\
10000	3326.7\\
12000	3453.1\\
14000	3360.3\\
16000	3387\\
18000	3365.7\\
20000	3321.6\\
22000	3490.1\\
24000	3320.2\\
26000	3359.9\\
30000	3402\\
32000	3259.4\\
34000	3451.6\\
36000	3401.4\\
38000	3398\\
40000	3363.5\\
42000	3384.1\\
44000	3385\\
46000	3290.4\\
48000	3323.4\\
50000	3316.3\\
52000	3278.7\\
54000	3436.5\\
56000	3386.4\\
58000	3511.6\\
60000	3427.3\\
62000	3296.6\\
64000	3287\\
66000	3265.4\\
68000	3374.4\\
70000	3336\\
74000	3419.1\\
76000	3510.7\\
78000	3457.1\\
82000	3426.4\\
84000	3374.4\\
86000	3439.2\\
88000	3471.2\\
90000	3271.8\\
92000	3385\\
96000	3373\\
98000	3315.5\\
1e+05	3389.9\\
};
\addplot [color=mycolor45, forget plot]
  table[row sep=crcr]{%
0	4395.3\\
2000	6549.4\\
4000	6819.4\\
6000	6968.7\\
8000	6857.3\\
10000	6913.2\\
12000	6766.2\\
14000	7023.7\\
16000	7003.9\\
18000	6718.7\\
20000	6804.3\\
24000	6929.8\\
26000	7046.1\\
28000	6974.8\\
30000	6966.4\\
32000	6884.8\\
34000	7102\\
36000	6937.7\\
38000	6841\\
40000	6948.8\\
42000	6814.4\\
44000	6955.6\\
46000	6938\\
48000	6886.5\\
50000	6873.8\\
52000	7050.6\\
54000	6858.9\\
56000	6903.7\\
58000	6981.1\\
64000	6885.2\\
66000	6975.7\\
68000	7012.9\\
70000	6887.7\\
72000	6819.3\\
74000	6672.9\\
76000	6921.6\\
78000	6865\\
80000	6989.3\\
82000	6885.7\\
84000	7093.2\\
86000	6810\\
88000	6963.8\\
90000	6695.2\\
92000	7068.4\\
94000	7046.5\\
96000	6943.4\\
98000	7061.7\\
1e+05	6897.4\\
};
\addplot [color=mycolor46, forget plot]
  table[row sep=crcr]{%
0	4395.3\\
2000	4225.6\\
4000	4138.7\\
6000	4158.6\\
8000	4086.7\\
10000	4098.2\\
12000	4092.3\\
14000	4061.1\\
16000	4048.8\\
18000	4016.8\\
20000	4157.5\\
22000	4110.8\\
24000	4117.4\\
26000	4054.2\\
28000	4122.6\\
30000	4134.1\\
34000	4043.2\\
36000	4166.4\\
38000	4186.6\\
40000	4189.8\\
42000	4121.3\\
46000	3903.9\\
48000	3974.1\\
50000	4005.1\\
52000	4152.2\\
54000	4038.8\\
56000	4227\\
58000	4061.9\\
60000	4030.2\\
62000	4066.3\\
64000	4120.2\\
66000	4016.4\\
68000	3953\\
70000	4087.8\\
72000	4091.7\\
74000	4161.9\\
76000	4108.7\\
78000	4248.1\\
80000	4086.9\\
82000	3967.7\\
84000	4065.6\\
90000	4024.4\\
92000	4206.1\\
94000	4136.6\\
96000	4004.6\\
98000	4118.8\\
1e+05	4124.5\\
};
\addplot [color=mycolor47, forget plot]
  table[row sep=crcr]{%
0	4395.3\\
2000	5223.5\\
4000	5283.1\\
8000	5330.6\\
10000	5323.8\\
12000	5335.9\\
16000	5232\\
18000	5285.4\\
20000	5218\\
22000	5355\\
24000	5412.1\\
26000	5345.8\\
28000	5146.3\\
30000	5164.2\\
32000	5485.5\\
36000	5200.3\\
38000	5226.9\\
40000	5087.6\\
42000	5215.6\\
44000	5322.8\\
46000	5285.4\\
48000	5431\\
50000	5435\\
52000	5244.1\\
54000	5297.8\\
56000	5263.6\\
58000	5130.7\\
62000	5303.4\\
64000	5542.8\\
68000	5249.2\\
70000	5338.5\\
72000	5298.1\\
74000	5420.7\\
76000	5168.1\\
78000	5388.1\\
80000	5367.2\\
82000	5226.6\\
84000	5337.8\\
86000	5431.2\\
90000	5276.5\\
92000	5349.4\\
94000	5340\\
96000	5220.3\\
98000	5257.2\\
1e+05	5070.3\\
};
\addplot [color=mycolor48, forget plot]
  table[row sep=crcr]{%
0	4395.3\\
2000	3165.2\\
4000	2901.9\\
6000	2891.4\\
8000	2939.6\\
12000	2901.9\\
14000	2947.4\\
16000	2956.9\\
18000	2921.7\\
20000	2901.3\\
22000	2943\\
24000	2956.1\\
26000	2953.9\\
28000	2903.2\\
30000	2899.4\\
32000	2932.2\\
34000	2949.9\\
36000	2956.8\\
38000	2939\\
40000	2912.3\\
44000	2925.4\\
46000	2896.1\\
48000	2905.6\\
50000	2973.9\\
52000	2944.1\\
54000	2957.4\\
56000	2910.6\\
58000	2896.3\\
60000	2957\\
62000	2928.8\\
64000	2950\\
66000	2923.1\\
74000	2923.3\\
76000	2947.1\\
78000	2921\\
80000	2945.5\\
84000	2937.2\\
86000	2912.5\\
88000	2925.5\\
90000	2872.7\\
92000	2977.3\\
94000	2982.9\\
96000	2976.1\\
1e+05	2892.2\\
};
\addplot [color=mycolor49, forget plot]
  table[row sep=crcr]{%
0	4395.3\\
2000	4433.8\\
4000	4627.4\\
6000	4349.1\\
8000	4543.5\\
10000	4343\\
12000	4348.6\\
14000	4452.7\\
16000	4465.1\\
18000	4371.3\\
20000	4424.8\\
22000	4227.3\\
24000	4399\\
26000	4298.3\\
28000	4377.3\\
30000	4439.8\\
32000	4386.5\\
34000	4381.1\\
36000	4354.6\\
38000	4548.4\\
40000	4281\\
42000	4398\\
44000	4283.5\\
46000	4436.8\\
48000	4144\\
50000	4276.6\\
52000	4386.5\\
54000	4382.3\\
56000	4321.5\\
58000	4388.1\\
60000	4354.5\\
64000	4358.2\\
66000	4455.1\\
68000	4327.3\\
72000	4357.7\\
74000	4335.5\\
76000	4144.9\\
78000	4358.3\\
80000	4399.6\\
82000	4342.1\\
84000	4168.4\\
86000	4282.7\\
88000	4194.5\\
90000	4223.7\\
92000	4344.3\\
94000	4210.5\\
96000	4401.4\\
98000	4288\\
1e+05	4539.6\\
};
\addplot [color=mycolor50, forget plot]
  table[row sep=crcr]{%
0	4395.3\\
2000	2799.7\\
4000	2436.6\\
6000	2449.2\\
8000	2504.9\\
10000	2533\\
12000	2543.5\\
14000	2486.4\\
16000	2526.6\\
18000	2461.8\\
20000	2440\\
22000	2491.9\\
24000	2498\\
26000	2478.4\\
28000	2497.8\\
30000	2440\\
32000	2506.7\\
34000	2486.1\\
36000	2569\\
38000	2465.5\\
44000	2485.7\\
48000	2532.3\\
50000	2477.2\\
52000	2402.6\\
54000	2476.1\\
56000	2434.7\\
58000	2422.3\\
60000	2476.9\\
62000	2501.2\\
64000	2475.6\\
66000	2521.3\\
68000	2438.2\\
70000	2465.8\\
72000	2561.6\\
74000	2470.1\\
76000	2521.7\\
78000	2503.1\\
80000	2441.8\\
82000	2473.8\\
84000	2491.1\\
86000	2441.6\\
88000	2484.3\\
92000	2510.6\\
94000	2446.1\\
96000	2518.8\\
98000	2482.7\\
1e+05	2462.4\\
};
\addplot [color=mycolor51, forget plot]
  table[row sep=crcr]{%
0	4395.3\\
2000	5752.9\\
4000	6060.5\\
8000	5970.6\\
12000	5703.4\\
14000	5958.2\\
16000	5444.1\\
18000	5895.6\\
20000	5811.8\\
22000	6140.8\\
24000	6371.4\\
26000	6036.9\\
28000	6176.6\\
30000	5849.1\\
32000	5913.7\\
34000	6047.5\\
36000	5971.5\\
38000	5833.4\\
40000	5860.1\\
42000	5635.3\\
44000	5807.2\\
46000	5938.8\\
48000	5976.7\\
50000	5861.5\\
52000	5840.5\\
54000	6090.3\\
56000	6024.5\\
58000	6071.1\\
60000	5889.9\\
62000	5805.9\\
64000	5913\\
66000	5951.6\\
68000	5860.4\\
70000	5927.2\\
74000	5927.4\\
76000	5753\\
78000	5900.8\\
80000	5966\\
82000	6050.1\\
84000	5964.8\\
86000	5942.3\\
88000	5945.3\\
90000	5692.4\\
92000	5996.2\\
94000	5825.2\\
96000	5847.9\\
98000	5908.9\\
1e+05	5846.3\\
};
\addplot [color=mycolor52, forget plot]
  table[row sep=crcr]{%
0	4395.3\\
2000	3889\\
4000	3548.2\\
6000	3463.8\\
8000	3539.9\\
10000	3508.5\\
12000	3549.8\\
14000	3641.4\\
16000	3482.9\\
18000	3483.3\\
20000	3521.2\\
22000	3583.8\\
24000	3406.5\\
26000	3520.9\\
28000	3415.8\\
30000	3561.9\\
32000	3429.9\\
34000	3536.8\\
36000	3601.7\\
38000	3700.1\\
40000	3675.9\\
42000	3451.8\\
44000	3398.1\\
46000	3429.7\\
48000	3489\\
50000	3485.4\\
52000	3435.2\\
54000	3607.4\\
56000	3565.9\\
58000	3415.8\\
60000	3539\\
62000	3596.1\\
64000	3420.7\\
66000	3457\\
68000	3520.9\\
70000	3655.3\\
72000	3650.8\\
74000	3534.6\\
76000	3445.1\\
78000	3472.2\\
80000	3638.9\\
82000	3546.5\\
84000	3573.6\\
86000	3571.8\\
88000	3379.6\\
90000	3631.5\\
92000	3547.7\\
94000	3622.4\\
96000	3489.2\\
98000	3583.1\\
1e+05	3623\\
};
\addplot [color=mycolor53, forget plot]
  table[row sep=crcr]{%
0	4395.3\\
2000	5629.1\\
8000	5759.4\\
10000	5595.3\\
12000	5851.8\\
14000	5794\\
16000	5795.8\\
18000	5679.6\\
20000	5651.6\\
24000	5715\\
26000	5844.5\\
28000	5661.4\\
30000	5703.8\\
32000	5565.1\\
34000	5860.1\\
36000	5724.6\\
38000	5629.7\\
40000	5820.3\\
44000	5654.7\\
46000	5590\\
48000	5769.4\\
50000	5851.1\\
52000	5324.7\\
54000	5730.8\\
56000	5569.8\\
58000	5715.7\\
60000	5681.2\\
62000	5736.1\\
64000	5740.1\\
66000	5814.2\\
68000	5833.9\\
70000	5714.6\\
72000	5696.9\\
74000	5828.8\\
76000	5787.8\\
78000	5644.3\\
80000	5569\\
82000	5652.2\\
84000	5960\\
86000	5680.7\\
88000	5885.8\\
90000	5635\\
92000	5547.7\\
94000	5699.8\\
96000	6044.8\\
98000	5728\\
1e+05	5547.6\\
};
\addplot [color=mycolor54, forget plot]
  table[row sep=crcr]{%
0	4395.3\\
2000	4571.6\\
4000	4555.2\\
8000	4587.3\\
10000	4549.4\\
12000	4565\\
14000	4738.9\\
18000	4525.2\\
20000	4673.1\\
22000	4696.4\\
24000	4471.9\\
26000	4634.9\\
28000	4531.3\\
30000	4613\\
32000	4645.6\\
34000	4570.4\\
36000	4570.8\\
38000	4742.1\\
40000	4680.8\\
42000	4646.8\\
44000	4526.9\\
46000	4572.8\\
48000	4531.8\\
50000	4650.6\\
52000	4574.6\\
54000	4679.6\\
56000	4471.4\\
58000	4601.6\\
60000	4507.8\\
62000	4483.2\\
66000	4540.2\\
68000	4668\\
70000	4529.8\\
72000	4501.7\\
74000	4654.9\\
76000	4676.1\\
78000	4618.3\\
80000	4592.8\\
82000	4637\\
84000	4516.1\\
86000	4589\\
88000	4524.6\\
92000	4528\\
94000	4639.3\\
96000	4524.5\\
98000	4665\\
1e+05	4682.7\\
};
\addplot [color=mycolor55, forget plot]
  table[row sep=crcr]{%
0	4395.3\\
2000	5706.1\\
4000	5831.7\\
6000	5730.3\\
8000	5866.4\\
10000	5675\\
12000	5793.7\\
14000	5757\\
16000	5843.5\\
18000	5702.2\\
20000	5742.3\\
22000	5855.8\\
24000	5719.2\\
26000	5812.6\\
28000	5807.3\\
32000	5669.5\\
34000	5676.7\\
36000	5864.2\\
38000	5703.9\\
40000	5763\\
42000	5724.5\\
44000	5722.5\\
46000	6053.3\\
48000	5805.9\\
50000	5736.4\\
52000	5730\\
54000	5777.1\\
56000	5874.2\\
58000	5698.3\\
60000	5672.6\\
62000	5613.9\\
66000	5856.4\\
70000	5682.9\\
72000	5960.4\\
74000	6001\\
76000	5733.4\\
78000	5832.6\\
80000	5855.7\\
82000	5927.8\\
84000	5809.3\\
88000	5802.3\\
90000	5752.4\\
92000	5719.2\\
94000	5728.1\\
96000	5820.7\\
1e+05	5710.1\\
};
\addplot [color=mycolor56, forget plot]
  table[row sep=crcr]{%
0	4395.3\\
2000	4087.2\\
4000	3982.1\\
6000	3997.7\\
8000	3994\\
10000	4083.7\\
12000	3961.4\\
14000	4055\\
16000	3959.4\\
18000	4021.8\\
20000	4013.2\\
24000	4097\\
26000	4034.3\\
28000	3959.9\\
30000	4027.3\\
32000	4114.5\\
36000	3976.5\\
38000	4018\\
40000	3978.6\\
42000	3925.7\\
44000	3998.4\\
46000	3957.8\\
48000	3944.5\\
50000	3961.3\\
54000	4120.2\\
56000	4000.4\\
58000	3987.7\\
60000	4100.1\\
66000	3932.4\\
68000	3906.9\\
70000	4090.8\\
72000	4062\\
74000	3967.6\\
76000	3914.2\\
78000	4126.3\\
80000	3970.2\\
82000	3983\\
84000	3952.8\\
86000	4018.3\\
88000	3992.8\\
90000	4034.4\\
92000	4062.1\\
94000	3927.4\\
96000	3976.8\\
98000	3953.4\\
1e+05	4060.4\\
};
\addplot [color=mycolor57, forget plot]
  table[row sep=crcr]{%
0	4395.3\\
2000	6054.7\\
4000	6800.9\\
6000	6788.9\\
8000	6715.1\\
10000	6622.6\\
12000	6579.3\\
14000	6638.4\\
16000	6730.5\\
18000	6725.1\\
20000	6798.3\\
22000	6583.5\\
24000	6598.8\\
28000	6713.8\\
30000	6683\\
32000	6827.7\\
34000	6863.1\\
36000	6769.7\\
38000	6610.7\\
40000	6742.9\\
42000	6848\\
44000	6776.3\\
46000	6724.2\\
48000	6763.9\\
50000	6698.5\\
52000	6771.1\\
54000	6626.3\\
56000	6720.1\\
58000	6660\\
60000	6764.5\\
64000	6613.3\\
66000	6599.9\\
68000	6846.4\\
70000	6673.2\\
72000	6753.8\\
74000	6692.3\\
76000	6774.3\\
78000	6598.7\\
80000	6788.2\\
82000	6732.8\\
86000	6454.3\\
88000	6586\\
90000	6617.8\\
92000	6734.9\\
94000	6702.2\\
96000	6446.7\\
98000	6697.1\\
1e+05	6732\\
};
\addplot [color=mycolor58, forget plot]
  table[row sep=crcr]{%
0	4395.3\\
2000	3774.5\\
4000	3667.2\\
6000	3533.8\\
8000	3600\\
10000	3552.1\\
12000	3516.2\\
14000	3654.7\\
16000	3490.3\\
18000	3601.6\\
20000	3503.7\\
22000	3601.1\\
24000	3558.3\\
26000	3567.2\\
28000	3690.7\\
30000	3576.8\\
32000	3590.7\\
34000	3527.1\\
36000	3637.4\\
38000	3697.4\\
40000	3672.7\\
42000	3612.3\\
44000	3729.6\\
46000	3559.5\\
48000	3648.7\\
50000	3651.8\\
52000	3527.1\\
54000	3569.5\\
56000	3567.5\\
60000	3455.7\\
64000	3615.4\\
66000	3621.3\\
68000	3521.9\\
70000	3629\\
72000	3630.7\\
74000	3482.1\\
76000	3487.8\\
78000	3666\\
80000	3647\\
82000	3614.9\\
84000	3629\\
86000	3500.2\\
88000	3538.5\\
90000	3538.4\\
92000	3711.7\\
94000	3598.4\\
96000	3562.8\\
98000	3592.2\\
1e+05	3525.8\\
};
\addplot [color=mycolor59, forget plot]
  table[row sep=crcr]{%
0	4395.3\\
2000	5065.5\\
4000	4924.9\\
6000	5008.1\\
8000	5033.5\\
12000	5194.1\\
14000	5198.8\\
16000	5057.5\\
18000	4958.3\\
20000	5126\\
22000	5062.1\\
24000	5101.6\\
26000	5037.2\\
28000	5161.1\\
30000	5098\\
32000	4885.2\\
34000	5092.4\\
36000	5103.8\\
38000	5083.1\\
40000	5115.7\\
42000	5128\\
44000	5177.7\\
46000	4948.9\\
48000	5084\\
50000	5060.3\\
52000	5139.2\\
54000	5151.1\\
56000	5003.6\\
58000	5112.4\\
60000	5019.1\\
62000	5037.5\\
64000	5331.8\\
66000	5114.6\\
68000	5163.4\\
70000	5034.8\\
72000	4963.7\\
76000	5140.7\\
78000	4983.7\\
80000	5042.7\\
82000	5020.4\\
84000	4964.3\\
86000	5111\\
88000	5058.1\\
90000	5107.8\\
92000	5124.9\\
94000	5044.6\\
96000	5071.9\\
98000	4990.8\\
1e+05	4981.7\\
};
\addplot [color=mycolor60, forget plot]
  table[row sep=crcr]{%
0	4395.3\\
2000	3719.3\\
4000	3628\\
6000	3695.8\\
8000	3517.7\\
10000	3498.1\\
12000	3500.8\\
14000	3534.8\\
16000	3532.7\\
18000	3586.2\\
22000	3576.7\\
24000	3585.4\\
26000	3542.9\\
28000	3475.5\\
30000	3504.5\\
32000	3559.6\\
34000	3578.8\\
36000	3528.8\\
38000	3567.1\\
40000	3569.1\\
42000	3589.8\\
44000	3529.1\\
46000	3596.9\\
48000	3609.7\\
50000	3531.2\\
52000	3573.8\\
54000	3519.7\\
56000	3559.3\\
58000	3558.9\\
60000	3632.3\\
62000	3550.7\\
64000	3559.7\\
66000	3598.2\\
68000	3585.1\\
70000	3593.3\\
72000	3578.5\\
76000	3530.4\\
78000	3527.1\\
80000	3583\\
82000	3597.8\\
84000	3630.4\\
86000	3549.6\\
88000	3533.2\\
90000	3561.5\\
92000	3501.7\\
94000	3542.3\\
96000	3552.8\\
98000	3490.9\\
1e+05	3508.4\\
};
\addplot [color=mycolor61, forget plot]
  table[row sep=crcr]{%
0	4395.3\\
2000	3427.1\\
4000	3124.9\\
6000	3120.1\\
8000	3126.4\\
10000	3146\\
12000	3118.1\\
14000	3100.6\\
18000	3180.3\\
20000	3146.3\\
22000	3150.2\\
24000	3122.6\\
26000	3168.3\\
28000	3089.4\\
30000	3100\\
32000	3122.4\\
36000	3145.9\\
38000	3139.5\\
40000	3120\\
44000	3154.6\\
46000	3151.9\\
48000	3126.3\\
50000	3128.9\\
52000	3102.2\\
54000	3183.2\\
56000	3223.7\\
58000	3177.3\\
60000	3152.7\\
62000	3148.3\\
64000	3128.9\\
66000	3084.2\\
68000	3154.1\\
70000	3097.8\\
72000	3153.9\\
74000	3164.3\\
76000	3125.3\\
78000	3123.6\\
80000	3156.2\\
82000	3113\\
84000	3140.4\\
86000	3150.8\\
88000	3191.6\\
90000	3113.7\\
92000	3170.2\\
94000	3153\\
96000	3167.4\\
98000	3134.4\\
1e+05	3135.2\\
};
\addplot [color=mycolor62, forget plot]
  table[row sep=crcr]{%
0	4395.3\\
2000	4870.6\\
4000	4970\\
6000	5049.2\\
8000	4995.9\\
10000	4988.2\\
12000	5077.3\\
14000	4970.2\\
16000	5011.3\\
18000	4915.5\\
22000	5044.2\\
24000	5012.8\\
26000	5011.2\\
28000	5075\\
32000	5107.8\\
34000	5019.2\\
36000	5001.4\\
38000	5059.2\\
40000	4981.2\\
42000	4920.1\\
44000	4982.1\\
48000	5011.6\\
50000	4979\\
52000	4921.5\\
54000	5081.7\\
56000	4992.7\\
58000	4952.6\\
60000	5037.8\\
64000	5026.5\\
68000	4967.4\\
70000	4974.7\\
72000	5054.1\\
74000	4998.8\\
76000	4969.5\\
78000	5029.3\\
80000	4904.6\\
84000	5025.7\\
86000	5026.6\\
88000	5101.8\\
92000	5013.7\\
94000	5076.9\\
96000	4904.7\\
1e+05	5008.8\\
};
\addplot [color=mycolor63, forget plot]
  table[row sep=crcr]{%
0	4395.3\\
2000	2752.1\\
4000	2412.1\\
6000	2393.2\\
8000	2394.6\\
10000	2426.3\\
12000	2422.2\\
14000	2452.3\\
16000	2376.3\\
18000	2444.9\\
20000	2439\\
22000	2478.7\\
24000	2420.6\\
28000	2443.4\\
30000	2375.2\\
32000	2392.1\\
36000	2405.7\\
38000	2349.6\\
40000	2440.7\\
42000	2429.3\\
44000	2467.8\\
46000	2417.7\\
48000	2393.5\\
50000	2440.6\\
52000	2423\\
54000	2450.8\\
56000	2445.8\\
58000	2469.4\\
60000	2467.1\\
64000	2441.4\\
66000	2423.6\\
68000	2424.4\\
70000	2449.7\\
72000	2402.2\\
74000	2398.7\\
76000	2442.6\\
78000	2422\\
82000	2399.1\\
84000	2421.3\\
86000	2374.9\\
90000	2422.5\\
92000	2402.8\\
94000	2433\\
96000	2441.1\\
98000	2419.4\\
1e+05	2463.8\\
};
\addplot [color=mycolor64, forget plot]
  table[row sep=crcr]{%
0	4395.3\\
2000	3299.2\\
4000	3093.4\\
6000	3027.7\\
10000	3088.8\\
12000	3084.2\\
14000	3028.3\\
16000	3001.3\\
18000	3042.7\\
20000	3028.8\\
22000	3092.6\\
24000	2952.5\\
26000	3073.9\\
30000	3104.8\\
32000	3035.9\\
36000	3055.3\\
38000	3070.5\\
40000	3053.9\\
44000	3059.9\\
46000	3104.8\\
48000	3057.4\\
50000	3127.2\\
52000	3037.7\\
54000	3097\\
56000	3059\\
58000	3010\\
60000	3066.7\\
62000	3054.3\\
64000	3081.9\\
66000	3090\\
68000	3056.8\\
70000	3034\\
72000	3089.9\\
74000	3041.3\\
76000	3050.1\\
78000	3088.2\\
80000	3117.3\\
82000	3065.1\\
84000	3047.4\\
86000	3057.9\\
88000	3008.4\\
90000	3027.9\\
92000	3070.5\\
94000	3091\\
96000	3014.2\\
98000	3065.6\\
1e+05	3021.2\\
};
\addplot [color=mycolor65, forget plot]
  table[row sep=crcr]{%
0	4395.3\\
2000	2892.8\\
4000	2584.3\\
6000	2559.2\\
8000	2654.7\\
10000	2559.6\\
12000	2570.1\\
16000	2534.3\\
18000	2527.7\\
20000	2557.5\\
24000	2539.5\\
26000	2585.8\\
28000	2541.6\\
30000	2605.6\\
32000	2512.4\\
34000	2574.4\\
36000	2629.9\\
38000	2586.5\\
40000	2556.1\\
42000	2550.5\\
44000	2506.5\\
46000	2559.8\\
48000	2579\\
50000	2535.1\\
52000	2597.9\\
54000	2518\\
56000	2510.5\\
58000	2563.2\\
60000	2580.1\\
62000	2498.2\\
64000	2550.2\\
66000	2574.4\\
70000	2549.2\\
72000	2568.8\\
74000	2544.6\\
76000	2562.3\\
78000	2592.4\\
82000	2535\\
84000	2523.9\\
86000	2530.1\\
88000	2599.2\\
90000	2567.8\\
92000	2505.1\\
94000	2490.4\\
96000	2563.3\\
98000	2565.2\\
1e+05	2504.3\\
};
\addplot [color=mycolor66, forget plot]
  table[row sep=crcr]{%
0	4395.3\\
2000	3639.6\\
4000	3474.3\\
6000	3413.5\\
8000	3421.9\\
10000	3404\\
14000	3435.6\\
16000	3473\\
18000	3481.3\\
20000	3423.6\\
22000	3454\\
24000	3402.4\\
26000	3488.5\\
28000	3428.6\\
30000	3431.7\\
32000	3380.2\\
34000	3447.9\\
36000	3441.8\\
38000	3399.2\\
40000	3412.5\\
42000	3436.9\\
44000	3383.2\\
46000	3387.8\\
50000	3511.8\\
52000	3497.3\\
54000	3417.1\\
58000	3411.7\\
60000	3460.5\\
62000	3470.4\\
64000	3415\\
66000	3411.2\\
68000	3435.2\\
70000	3496.8\\
72000	3395.9\\
74000	3409.9\\
76000	3486.8\\
78000	3481.7\\
82000	3455.2\\
84000	3411.3\\
86000	3407\\
88000	3478.5\\
90000	3395.8\\
92000	3422.4\\
96000	3425\\
98000	3490.2\\
1e+05	3412\\
};
\addplot [color=mycolor67, forget plot]
  table[row sep=crcr]{%
0	4395.3\\
2000	3456.5\\
4000	3245.9\\
6000	3247.7\\
8000	3176\\
10000	3162.5\\
12000	3168.3\\
14000	3211.2\\
16000	3166.4\\
18000	3159.7\\
20000	3176.5\\
22000	3182.2\\
24000	3201.2\\
26000	3184.5\\
28000	3181.5\\
30000	3222.9\\
32000	3183.3\\
34000	3189.5\\
36000	3208.6\\
38000	3140.4\\
40000	3228.6\\
42000	3187.9\\
44000	3229.7\\
46000	3226.2\\
48000	3053.4\\
50000	3192.9\\
52000	3165.7\\
54000	3229.2\\
56000	3257.5\\
58000	3229.6\\
60000	3182.1\\
62000	3183.5\\
64000	3153.2\\
66000	3175.6\\
68000	3141.3\\
70000	3162.8\\
72000	3243.4\\
74000	3128.4\\
76000	3192.7\\
80000	3137.7\\
82000	3213.7\\
84000	3222.7\\
86000	3185.8\\
88000	3188\\
90000	3227.2\\
92000	3185.6\\
94000	3193.4\\
98000	3170.3\\
1e+05	3156.2\\
};
\addplot [color=mycolor67]
  table[row sep=crcr]{%
0	4395.3\\
2000	2989.5\\
4000	2670.3\\
6000	2639.2\\
10000	2643.6\\
12000	2671\\
14000	2642.5\\
18000	2629.3\\
20000	2617.1\\
22000	2665.7\\
24000	2638.6\\
26000	2647.7\\
36000	2645.5\\
38000	2601.9\\
44000	2654.7\\
46000	2635.6\\
48000	2604.3\\
50000	2651\\
56000	2646.7\\
58000	2666\\
60000	2659.2\\
62000	2618.7\\
64000	2659.5\\
66000	2581.3\\
68000	2590.8\\
70000	2641.8\\
72000	2607.8\\
76000	2634.4\\
78000	2648.8\\
80000	2619.1\\
82000	2611.8\\
84000	2636.9\\
86000	2613.9\\
90000	2645\\
92000	2682.2\\
94000	2673.2\\
96000	2619.7\\
98000	2651.6\\
1e+05	2635.5\\
};
\addlegendentry{Classe k=43}

\addplot [color=mycolor68, forget plot]
  table[row sep=crcr]{%
0	4395.3\\
2000	2878.9\\
4000	2450\\
6000	2450.1\\
8000	2462.4\\
10000	2458.9\\
12000	2468.3\\
14000	2487.1\\
16000	2457.5\\
18000	2513.1\\
20000	2487.1\\
22000	2477.2\\
24000	2434.2\\
26000	2461.4\\
28000	2475.2\\
30000	2464.4\\
32000	2468.7\\
34000	2440.7\\
36000	2434.3\\
38000	2472.1\\
40000	2459.5\\
42000	2482.6\\
44000	2492.5\\
46000	2453.3\\
48000	2469.1\\
50000	2469.3\\
52000	2454.4\\
54000	2458.8\\
58000	2482.9\\
60000	2452.8\\
62000	2495.9\\
64000	2497.5\\
70000	2483\\
72000	2463.6\\
74000	2427.3\\
76000	2460.5\\
78000	2452.1\\
80000	2485.2\\
82000	2491.4\\
84000	2443.6\\
86000	2431\\
90000	2460.2\\
92000	2494.7\\
94000	2462.6\\
96000	2490.7\\
98000	2460.4\\
1e+05	2472.2\\
};
\addplot [color=mycolor69, forget plot]
  table[row sep=crcr]{%
0	4395.3\\
2000	3282.4\\
4000	2934.2\\
6000	2915.6\\
10000	2898.6\\
12000	2903.2\\
14000	2872.2\\
20000	2881.7\\
22000	2876.8\\
24000	2935.1\\
26000	2964.3\\
28000	2902.4\\
30000	2891.1\\
32000	2951.6\\
34000	2958.5\\
36000	2952.4\\
40000	2849.9\\
42000	2960.1\\
44000	2956.8\\
46000	2924.2\\
48000	2906.5\\
50000	2916.4\\
52000	2942.4\\
54000	2904.2\\
56000	2918.2\\
58000	2921.6\\
60000	2940\\
64000	2887.6\\
66000	2943.7\\
68000	2918.4\\
70000	2914.1\\
72000	2876.1\\
74000	2932.5\\
76000	2916.5\\
78000	2932.2\\
80000	2908.5\\
82000	2961.2\\
84000	2906.4\\
90000	2923.5\\
92000	2900.5\\
94000	2909.6\\
96000	2936.6\\
98000	2954.5\\
1e+05	2934.4\\
};
\addplot [color=mycolor70, forget plot]
  table[row sep=crcr]{%
0	4395.3\\
2000	2524.6\\
4000	2157\\
6000	2100.9\\
8000	2084.2\\
10000	2119.6\\
12000	2140.6\\
14000	2106.6\\
16000	2106.3\\
18000	2128.2\\
20000	2109.9\\
22000	2137.5\\
24000	2081.5\\
26000	2081.5\\
28000	2129.7\\
30000	2147.1\\
32000	2114.1\\
34000	2133.6\\
36000	2108.5\\
38000	2122.5\\
40000	2109.5\\
42000	2134.4\\
44000	2136.6\\
46000	2086.6\\
48000	2115.9\\
50000	2095.7\\
54000	2112.6\\
58000	2134.5\\
60000	2072.1\\
62000	2147.8\\
64000	2120.3\\
66000	2159\\
68000	2126.1\\
70000	2117.6\\
72000	2138.9\\
74000	2127.8\\
76000	2144.1\\
78000	2107.5\\
80000	2141.3\\
82000	2135.5\\
84000	2151.5\\
86000	2117\\
88000	2124.7\\
90000	2147\\
92000	2128.3\\
94000	2127.4\\
96000	2174.3\\
98000	2126.9\\
1e+05	2133.9\\
};
\addplot [color=mycolor71, forget plot]
  table[row sep=crcr]{%
0	4395.3\\
2000	2865\\
4000	2497\\
6000	2511.2\\
8000	2498.1\\
14000	2494.5\\
16000	2478.1\\
18000	2485.4\\
20000	2480.8\\
22000	2465.2\\
24000	2518.1\\
28000	2504.7\\
30000	2520.1\\
32000	2501\\
34000	2530.8\\
36000	2486.9\\
38000	2487.7\\
40000	2503.9\\
42000	2483.3\\
44000	2483.9\\
46000	2509\\
48000	2505.3\\
52000	2522.5\\
54000	2496.2\\
56000	2461.2\\
58000	2460.5\\
60000	2467.5\\
62000	2502.8\\
64000	2505.5\\
66000	2477\\
68000	2524.6\\
70000	2517.3\\
72000	2492.3\\
76000	2475.7\\
78000	2501.1\\
80000	2488\\
82000	2530.2\\
84000	2509.2\\
88000	2487.3\\
90000	2495.2\\
92000	2513.2\\
94000	2483.3\\
96000	2470.5\\
98000	2484\\
1e+05	2505\\
};
\addplot [color=mycolor72, forget plot]
  table[row sep=crcr]{%
0	4395.3\\
2000	2653.5\\
4000	2242.1\\
6000	2180.4\\
8000	2244\\
10000	2275.6\\
12000	2239.3\\
14000	2148.1\\
16000	2232.2\\
18000	2188.2\\
22000	2216.7\\
24000	2208.4\\
26000	2220.5\\
28000	2215.5\\
30000	2192.2\\
32000	2225.5\\
34000	2241.7\\
36000	2232.8\\
38000	2239.3\\
40000	2226.6\\
42000	2196.7\\
44000	2261.5\\
46000	2191.6\\
48000	2241.5\\
50000	2180.6\\
52000	2207.1\\
54000	2258.3\\
56000	2210.5\\
58000	2198\\
62000	2208.5\\
68000	2262.4\\
70000	2218.1\\
74000	2176.7\\
76000	2156.3\\
78000	2213\\
80000	2203.9\\
82000	2217.3\\
84000	2205.9\\
86000	2207\\
88000	2167.3\\
90000	2171.1\\
92000	2248.3\\
94000	2198.2\\
96000	2177.6\\
98000	2191.5\\
1e+05	2154.1\\
};
\addplot [color=mycolor73, forget plot]
  table[row sep=crcr]{%
0	4395.3\\
2000	2961.7\\
4000	2635\\
6000	2638.3\\
8000	2587.2\\
10000	2632.6\\
12000	2579.9\\
14000	2641.2\\
16000	2593.1\\
18000	2617.6\\
20000	2622.7\\
22000	2602.7\\
24000	2615.3\\
26000	2607.7\\
30000	2644.4\\
32000	2590.3\\
34000	2665\\
36000	2631.7\\
38000	2608.6\\
40000	2612.4\\
42000	2628.2\\
44000	2600\\
46000	2634.2\\
48000	2582.6\\
50000	2598.5\\
52000	2636.2\\
54000	2618.1\\
56000	2630.4\\
58000	2587.5\\
60000	2656.8\\
62000	2595.1\\
64000	2633.8\\
68000	2615.8\\
70000	2624.7\\
72000	2652.5\\
74000	2644.4\\
76000	2613.4\\
78000	2618\\
80000	2612.2\\
82000	2591.7\\
84000	2606.7\\
86000	2658.8\\
88000	2581.9\\
96000	2668.5\\
98000	2641.2\\
1e+05	2605.1\\
};
\addplot [color=mycolor74, forget plot]
  table[row sep=crcr]{%
0	4395.3\\
2000	2041.5\\
4000	1569.5\\
6000	1538.8\\
8000	1537.3\\
10000	1548\\
12000	1539.9\\
14000	1522.2\\
16000	1537.6\\
18000	1539.7\\
20000	1531.8\\
22000	1532\\
24000	1537.7\\
26000	1538.2\\
28000	1563.9\\
30000	1551.8\\
32000	1564.6\\
34000	1548.5\\
36000	1566.7\\
38000	1566.7\\
40000	1547.4\\
52000	1567.8\\
54000	1542.6\\
56000	1553.1\\
58000	1535.5\\
60000	1552.5\\
62000	1529.4\\
64000	1544.6\\
66000	1568\\
68000	1561.7\\
70000	1530.3\\
72000	1560.6\\
74000	1545\\
76000	1540.2\\
78000	1543.8\\
80000	1559.5\\
82000	1554.1\\
84000	1567.4\\
88000	1567.1\\
90000	1527.6\\
92000	1548.7\\
94000	1556.4\\
96000	1555\\
98000	1559.4\\
1e+05	1556.3\\
};
\addplot [color=mycolor75, forget plot]
  table[row sep=crcr]{%
0	4395.3\\
2000	2702.6\\
4000	2231.7\\
6000	2190.5\\
8000	2215.6\\
10000	2218.6\\
12000	2203.2\\
18000	2241.5\\
20000	2248.4\\
24000	2213.7\\
26000	2222.1\\
28000	2213.8\\
32000	2218.5\\
34000	2236.8\\
36000	2169.4\\
38000	2206.1\\
40000	2202.1\\
44000	2245.1\\
46000	2194.2\\
48000	2223.3\\
52000	2245.7\\
56000	2211\\
58000	2253.7\\
60000	2225.4\\
62000	2227.8\\
64000	2200.7\\
70000	2242.3\\
72000	2215.8\\
74000	2218.5\\
76000	2274.9\\
78000	2244.5\\
80000	2242\\
82000	2231.6\\
84000	2241.5\\
88000	2209.6\\
90000	2215.7\\
92000	2193.3\\
94000	2232.7\\
96000	2256.3\\
98000	2269.2\\
1e+05	2222.1\\
};
\addplot [color=mycolor76, forget plot]
  table[row sep=crcr]{%
0	4395.3\\
2000	2550.7\\
4000	2047.9\\
6000	2033.6\\
8000	2035.6\\
10000	2056.1\\
12000	2032.6\\
14000	1996.1\\
16000	2025.5\\
18000	2021.4\\
20000	2041.9\\
22000	2043.3\\
26000	2020.2\\
28000	1989.6\\
30000	1993.5\\
32000	2030.6\\
34000	2036.6\\
36000	1998.2\\
38000	2027.1\\
42000	2023.9\\
46000	1996.3\\
48000	2040.7\\
50000	2014.4\\
52000	2032.3\\
54000	2023.2\\
56000	2040.2\\
58000	2036\\
60000	2023.6\\
62000	2053.2\\
64000	2026.7\\
66000	2050.1\\
68000	2033.9\\
70000	2011.8\\
72000	2013.2\\
74000	2031.6\\
76000	2020.4\\
78000	2029.2\\
80000	2013.4\\
82000	2055.8\\
84000	2019.9\\
86000	2046.9\\
88000	2036.4\\
90000	2043.5\\
92000	2026.8\\
94000	2026\\
96000	2012.4\\
98000	2037\\
1e+05	2045.3\\
};
\addplot [color=mycolor77, forget plot]
  table[row sep=crcr]{%
0	4395.3\\
2000	2149.2\\
4000	1587.1\\
6000	1540.9\\
8000	1555.4\\
10000	1553.2\\
12000	1576.9\\
14000	1558\\
16000	1583.7\\
18000	1581\\
20000	1596.6\\
22000	1571.6\\
24000	1570\\
26000	1575.3\\
28000	1573.9\\
30000	1588.2\\
32000	1579.1\\
34000	1553.6\\
36000	1569.1\\
40000	1564.1\\
42000	1579.9\\
44000	1579.9\\
46000	1575.2\\
48000	1584.3\\
50000	1562.8\\
52000	1605.2\\
54000	1580.2\\
56000	1597.5\\
58000	1582.3\\
60000	1572.6\\
62000	1581.4\\
64000	1577.8\\
66000	1565.2\\
68000	1571.9\\
70000	1565.3\\
72000	1595.7\\
74000	1565.7\\
78000	1575.9\\
80000	1551.2\\
82000	1559.9\\
84000	1560.4\\
86000	1572.7\\
88000	1595.5\\
90000	1534.6\\
94000	1588\\
96000	1567.9\\
98000	1578.6\\
1e+05	1584.4\\
};
\addplot [color=mycolor78, forget plot]
  table[row sep=crcr]{%
0	4395.3\\
2000	2789.2\\
4000	2323\\
6000	2306.5\\
8000	2328.7\\
12000	2260.5\\
14000	2309.7\\
16000	2314.6\\
20000	2273.1\\
22000	2282\\
24000	2313.5\\
26000	2275.7\\
28000	2320.6\\
30000	2282.7\\
32000	2338.2\\
34000	2325.9\\
36000	2337.3\\
38000	2326.7\\
40000	2268.6\\
42000	2333.3\\
44000	2301.1\\
46000	2292.6\\
48000	2306.6\\
50000	2263.6\\
52000	2339.6\\
54000	2283.9\\
58000	2272.9\\
60000	2294.4\\
62000	2278.9\\
64000	2311.8\\
66000	2258.7\\
72000	2295.8\\
76000	2305.2\\
78000	2298\\
82000	2337.4\\
84000	2256.2\\
86000	2246.2\\
88000	2304.9\\
90000	2277.6\\
92000	2301.9\\
94000	2344.1\\
96000	2315.4\\
98000	2325.8\\
1e+05	2312\\
};
\addplot [color=mycolor79, forget plot]
  table[row sep=crcr]{%
0	4395.3\\
2000	2667.7\\
4000	2182.6\\
6000	2186.8\\
8000	2143.1\\
10000	2147.3\\
12000	2157.7\\
14000	2184.8\\
18000	2163.9\\
20000	2194.1\\
22000	2206.4\\
26000	2163.2\\
28000	2190.5\\
30000	2191.8\\
32000	2130.4\\
34000	2163.1\\
36000	2183.5\\
38000	2180.7\\
40000	2155.1\\
42000	2156.7\\
44000	2186.9\\
46000	2161.6\\
48000	2155.1\\
50000	2108.4\\
52000	2195.2\\
54000	2181.3\\
56000	2202.2\\
58000	2160.9\\
60000	2130.3\\
62000	2178.9\\
64000	2200\\
66000	2214.1\\
68000	2177.7\\
70000	2191.8\\
72000	2175.6\\
74000	2180.8\\
76000	2129.2\\
78000	2181.1\\
80000	2163.7\\
82000	2162.7\\
84000	2184.2\\
86000	2124.5\\
88000	2192.2\\
92000	2175.3\\
94000	2208.2\\
96000	2214.1\\
98000	2140.1\\
1e+05	2163\\
};
\addplot [color=mycolor79, forget plot]
  table[row sep=crcr]{%
0	4395.3\\
2000	2579.3\\
4000	2069\\
6000	2058.6\\
8000	2022.4\\
10000	2032.1\\
12000	2063.3\\
14000	2044.9\\
16000	2020.3\\
20000	2004.3\\
22000	2088.7\\
24000	2024.1\\
26000	2026.2\\
28000	1975.8\\
30000	2007.7\\
32000	2015.1\\
34000	1960.5\\
36000	2023\\
38000	2021\\
42000	2031.2\\
44000	2020.4\\
46000	2044.6\\
48000	2049.2\\
50000	1973.4\\
52000	2017.8\\
54000	1996.6\\
58000	2016.9\\
60000	2022.3\\
62000	2011.4\\
64000	2041.9\\
66000	2038.7\\
68000	2042.9\\
70000	2014.4\\
72000	2033.2\\
74000	2030.9\\
76000	2021\\
78000	2030.7\\
80000	2028.9\\
84000	2041.6\\
86000	2005.4\\
88000	2018.9\\
90000	2045.6\\
92000	2060\\
94000	2009.2\\
98000	2046.1\\
1e+05	2028.2\\
};
\addplot [color=mycolor80, forget plot]
  table[row sep=crcr]{%
0	4395.3\\
2000	2219.8\\
4000	1595.4\\
6000	1555.7\\
8000	1567.7\\
10000	1597\\
12000	1567.8\\
14000	1583\\
16000	1541.6\\
18000	1569.3\\
20000	1580.3\\
22000	1558.3\\
24000	1543.5\\
26000	1542.3\\
28000	1549.8\\
30000	1566\\
32000	1545.1\\
36000	1583.6\\
38000	1562.7\\
40000	1565.7\\
42000	1562\\
44000	1575.5\\
46000	1535.3\\
48000	1577.3\\
50000	1568.1\\
52000	1548.1\\
54000	1563.2\\
56000	1558.7\\
58000	1582.8\\
60000	1561.5\\
62000	1548.9\\
64000	1552.5\\
66000	1549.1\\
68000	1529.6\\
70000	1554\\
72000	1551\\
74000	1552.3\\
76000	1562.8\\
80000	1546.3\\
82000	1578.3\\
84000	1565.6\\
86000	1561.3\\
88000	1547.9\\
90000	1562.9\\
92000	1563.9\\
94000	1559\\
96000	1586.7\\
98000	1555.8\\
1e+05	1554.4\\
};
\addplot [color=mycolor81, forget plot]
  table[row sep=crcr]{%
0	4395.3\\
2000	2079.4\\
4000	1477.5\\
6000	1476\\
8000	1479.6\\
10000	1438.1\\
12000	1441.7\\
14000	1436.1\\
16000	1448.3\\
18000	1443.7\\
24000	1470.1\\
26000	1468.4\\
30000	1449.2\\
32000	1472.4\\
34000	1480.8\\
36000	1457.9\\
40000	1447.1\\
44000	1471.5\\
46000	1447.2\\
48000	1477.3\\
50000	1494\\
52000	1470\\
60000	1477.7\\
62000	1449.8\\
66000	1471.3\\
68000	1469.5\\
70000	1458.7\\
72000	1483.8\\
74000	1470.5\\
76000	1473.7\\
78000	1463.1\\
80000	1438.9\\
82000	1450.6\\
84000	1472.9\\
86000	1470.4\\
88000	1447.1\\
90000	1460.2\\
92000	1477.8\\
94000	1460.7\\
96000	1476.4\\
98000	1451.9\\
1e+05	1503.5\\
};
\addplot [color=mycolor82, forget plot]
  table[row sep=crcr]{%
0	4395.3\\
2000	2312.1\\
4000	1722.6\\
6000	1672\\
8000	1725.8\\
10000	1702.2\\
12000	1692.9\\
14000	1705.4\\
16000	1668.9\\
18000	1668.2\\
20000	1685.6\\
22000	1694.9\\
24000	1716.2\\
28000	1718.9\\
30000	1698.4\\
32000	1706.8\\
34000	1696.3\\
38000	1714.3\\
40000	1683.6\\
42000	1707.2\\
44000	1682.9\\
48000	1707\\
50000	1706.3\\
52000	1678\\
54000	1697.2\\
56000	1692.1\\
58000	1698.9\\
60000	1719.3\\
62000	1688.6\\
64000	1702.8\\
66000	1701.6\\
68000	1714.7\\
70000	1723\\
72000	1716.2\\
74000	1690.2\\
76000	1737.5\\
78000	1700.5\\
80000	1718.4\\
82000	1700.6\\
86000	1694\\
88000	1701.9\\
90000	1702.3\\
92000	1690.4\\
94000	1694.4\\
96000	1724\\
98000	1707.5\\
1e+05	1705\\
};
\addplot [color=mycolor83, forget plot]
  table[row sep=crcr]{%
0	4395.3\\
2000	2361.9\\
4000	1723.3\\
6000	1736\\
8000	1741.3\\
10000	1741.2\\
12000	1727\\
14000	1735.3\\
16000	1755.9\\
18000	1748.5\\
20000	1728.7\\
22000	1699\\
26000	1735.9\\
28000	1741.8\\
30000	1734.4\\
32000	1746.7\\
34000	1712.5\\
36000	1734.2\\
38000	1738.1\\
40000	1736.1\\
42000	1726.2\\
44000	1753\\
46000	1731.7\\
48000	1732.3\\
50000	1749.1\\
54000	1719.9\\
58000	1746.6\\
60000	1737.3\\
62000	1736.4\\
64000	1717.2\\
66000	1743.3\\
68000	1733.4\\
70000	1754.1\\
72000	1731\\
74000	1692.1\\
76000	1714\\
78000	1718.5\\
80000	1739.5\\
82000	1744\\
84000	1736.5\\
86000	1747.8\\
88000	1747.8\\
90000	1755.4\\
92000	1747.6\\
94000	1754.2\\
96000	1730.2\\
98000	1727.1\\
1e+05	1742.4\\
};
\addplot [color=mycolor83, forget plot]
  table[row sep=crcr]{%
0	4395.3\\
2000	1963.7\\
4000	1297\\
6000	1279.7\\
8000	1275.8\\
10000	1299.2\\
12000	1260.1\\
16000	1270.2\\
18000	1270.4\\
20000	1260.1\\
22000	1268.3\\
24000	1261\\
26000	1278.1\\
28000	1270.9\\
30000	1269.4\\
32000	1274.1\\
34000	1254.3\\
36000	1287.7\\
38000	1286.4\\
40000	1275.1\\
42000	1279\\
44000	1287.3\\
46000	1256.7\\
50000	1284.7\\
52000	1265.8\\
54000	1285.1\\
56000	1281.8\\
58000	1293.8\\
60000	1276.9\\
62000	1282.3\\
64000	1299.8\\
66000	1282.3\\
68000	1258.1\\
70000	1254.1\\
72000	1289.8\\
76000	1270.2\\
78000	1296\\
80000	1263.2\\
82000	1282\\
84000	1277.1\\
86000	1280.8\\
88000	1248.2\\
90000	1249.7\\
92000	1275.4\\
94000	1272.8\\
96000	1263.9\\
98000	1269.2\\
1e+05	1302.9\\
};
\addplot [color=mycolor84, forget plot]
  table[row sep=crcr]{%
0	4395.3\\
2000	2038.1\\
4000	1313.3\\
6000	1278\\
8000	1314.5\\
10000	1279.9\\
12000	1288.8\\
14000	1261.2\\
16000	1282.9\\
18000	1294.3\\
20000	1300.5\\
22000	1292.8\\
24000	1259.6\\
26000	1269.6\\
28000	1269.2\\
30000	1309\\
32000	1284.4\\
34000	1301.4\\
36000	1301.6\\
38000	1284.5\\
40000	1245\\
42000	1307.9\\
44000	1294.1\\
46000	1259.3\\
48000	1267.5\\
50000	1308.6\\
52000	1293.4\\
54000	1303.9\\
56000	1300.3\\
58000	1269.1\\
60000	1285.6\\
62000	1309.6\\
64000	1289.2\\
66000	1278.1\\
68000	1290.9\\
70000	1254.1\\
72000	1292.5\\
74000	1285.6\\
76000	1247.7\\
78000	1240.9\\
80000	1276.9\\
82000	1260.6\\
84000	1269.1\\
86000	1293.4\\
88000	1308.3\\
90000	1277.7\\
92000	1278.6\\
94000	1266.6\\
96000	1285.2\\
98000	1269.6\\
1e+05	1264.4\\
};
\addplot [color=mycolor85, forget plot]
  table[row sep=crcr]{%
0	4395.3\\
2000	2100.5\\
4000	1435.9\\
6000	1389.7\\
8000	1406.8\\
10000	1409.5\\
12000	1403.2\\
14000	1416.8\\
16000	1405\\
18000	1400.5\\
20000	1388.9\\
22000	1422.3\\
24000	1413.3\\
26000	1397.9\\
28000	1436.5\\
32000	1381.5\\
34000	1394\\
36000	1395\\
38000	1362.5\\
40000	1388.1\\
42000	1427.6\\
44000	1404.6\\
46000	1391.3\\
48000	1391.3\\
50000	1424.4\\
52000	1414.6\\
54000	1410\\
56000	1427\\
58000	1416.4\\
60000	1394.4\\
62000	1404.6\\
64000	1401.8\\
66000	1345.5\\
70000	1426.8\\
72000	1405\\
76000	1443.8\\
78000	1418.5\\
80000	1433.8\\
82000	1410.8\\
84000	1405.8\\
86000	1411.1\\
88000	1386\\
90000	1421.6\\
92000	1380.8\\
94000	1396.5\\
96000	1389.8\\
98000	1390.6\\
1e+05	1359.9\\
};
\addplot [color=mycolor86, forget plot]
  table[row sep=crcr]{%
0	4395.3\\
2000	1982.9\\
4000	1192.6\\
6000	1189.5\\
8000	1160.6\\
10000	1160.9\\
12000	1179.7\\
14000	1171\\
16000	1181.6\\
18000	1187.7\\
20000	1200.1\\
22000	1163.6\\
24000	1174.5\\
26000	1196.5\\
28000	1197.7\\
30000	1169.7\\
32000	1177.2\\
34000	1189.4\\
36000	1214.9\\
38000	1215.7\\
40000	1198.7\\
42000	1197.4\\
44000	1175.5\\
46000	1163.7\\
48000	1189.3\\
50000	1181.1\\
52000	1206.8\\
54000	1207.7\\
56000	1173\\
58000	1166.5\\
60000	1175.3\\
62000	1175.1\\
64000	1212.2\\
66000	1177.5\\
68000	1181.8\\
72000	1178.1\\
74000	1204.2\\
76000	1177.3\\
78000	1182.3\\
80000	1167.1\\
82000	1191.3\\
84000	1172.1\\
86000	1185.2\\
88000	1193\\
90000	1207.6\\
92000	1174.7\\
94000	1186.5\\
96000	1171.2\\
98000	1168.3\\
1e+05	1193.6\\
};
\addplot [color=mycolor87, forget plot]
  table[row sep=crcr]{%
0	4395.3\\
2000	1883.4\\
4000	1088.3\\
6000	1029.7\\
8000	960.65\\
10000	991.34\\
12000	1041\\
14000	1005.7\\
16000	1034\\
18000	990.28\\
20000	972.97\\
22000	977.04\\
24000	977.48\\
26000	989.01\\
28000	993.17\\
30000	1010.9\\
32000	1004.1\\
34000	969.68\\
36000	970.44\\
38000	1029.2\\
40000	952.46\\
42000	975.41\\
44000	959.34\\
46000	1017.1\\
48000	1029.4\\
50000	1035.6\\
52000	1058.2\\
54000	1031.5\\
56000	983.31\\
58000	992\\
60000	1026.1\\
62000	1012.4\\
66000	970.55\\
68000	976.19\\
70000	1005.7\\
72000	1012.4\\
76000	1033.8\\
78000	1026.7\\
80000	1014.7\\
82000	1016\\
84000	989.11\\
86000	997.12\\
88000	972.61\\
90000	980.21\\
92000	996.77\\
94000	1050.6\\
96000	1001.4\\
98000	988.8\\
1e+05	1002.7\\
};
\addplot [color=mycolor88, forget plot]
  table[row sep=crcr]{%
0	4395.3\\
2000	2061.7\\
4000	1248.1\\
6000	1178.2\\
8000	1173.7\\
10000	1158.6\\
12000	1183.4\\
14000	1160.9\\
16000	1168.4\\
18000	1155.7\\
22000	1179.7\\
24000	1203.9\\
26000	1186.6\\
28000	1196.3\\
30000	1177\\
32000	1208.4\\
34000	1180.1\\
36000	1185.6\\
38000	1207.1\\
40000	1185.3\\
42000	1222.6\\
44000	1142.3\\
46000	1143.7\\
48000	1193.1\\
50000	1184.8\\
52000	1186.9\\
54000	1206.7\\
56000	1182.3\\
58000	1182.7\\
60000	1173.4\\
64000	1181.4\\
66000	1179\\
68000	1195.9\\
72000	1185.4\\
74000	1189.1\\
76000	1208.4\\
80000	1173.2\\
82000	1196.7\\
86000	1180.8\\
88000	1192.7\\
94000	1178.7\\
96000	1186.6\\
1e+05	1167.5\\
};
\addplot [color=mycolor89, forget plot]
  table[row sep=crcr]{%
0	4395.3\\
2000	1760.5\\
4000	904.65\\
6000	854.3\\
8000	852.99\\
10000	866.48\\
12000	847.8\\
14000	839.83\\
16000	839.06\\
18000	855.32\\
20000	839.74\\
22000	862.88\\
24000	833.47\\
26000	851.15\\
28000	842.67\\
32000	860.19\\
34000	845.42\\
36000	833.94\\
38000	835.88\\
40000	868.31\\
42000	844.98\\
44000	852.19\\
46000	830.86\\
48000	856.49\\
50000	842.75\\
52000	837.88\\
54000	851.28\\
56000	834.33\\
58000	863.94\\
60000	839.85\\
62000	860.32\\
64000	837.79\\
66000	856.03\\
68000	861.38\\
70000	854.61\\
72000	871.87\\
74000	862.56\\
76000	860.6\\
78000	862.16\\
80000	884.43\\
82000	859.41\\
84000	857.57\\
86000	872.65\\
88000	844.66\\
90000	831.93\\
92000	852.21\\
94000	880.51\\
96000	868.59\\
98000	874.81\\
1e+05	843.83\\
};
\addplot [color=mycolor90, forget plot]
  table[row sep=crcr]{%
0	4395.3\\
2000	1776.5\\
4000	889.27\\
6000	857.25\\
8000	829.09\\
10000	811.7\\
12000	838.7\\
14000	834.01\\
16000	833.96\\
18000	815.07\\
22000	830.09\\
24000	824.13\\
26000	843.03\\
28000	821.18\\
30000	826.04\\
32000	834.56\\
34000	856.39\\
36000	844.81\\
38000	841.01\\
40000	830.79\\
42000	847.15\\
44000	825.93\\
46000	830.07\\
48000	810.4\\
50000	809.54\\
52000	821.7\\
54000	837.22\\
56000	843.26\\
58000	841.82\\
60000	833.73\\
62000	828.7\\
64000	806.31\\
66000	827.41\\
68000	857.11\\
70000	843.65\\
72000	826.9\\
74000	834.52\\
78000	820.15\\
80000	829.13\\
82000	852.25\\
84000	816.74\\
86000	805.64\\
88000	862.58\\
90000	841.59\\
92000	849.93\\
94000	854.65\\
96000	846.8\\
98000	842.83\\
1e+05	812.07\\
};
\addplot [color=mycolor91, forget plot]
  table[row sep=crcr]{%
0	4395.3\\
2000	1714.7\\
4000	773.51\\
6000	714.33\\
8000	739.09\\
10000	720.36\\
12000	722.48\\
14000	706.81\\
16000	703.11\\
18000	717.08\\
20000	706.63\\
22000	704.15\\
24000	719.52\\
26000	727.89\\
28000	733.68\\
30000	725.44\\
32000	729.7\\
34000	718.28\\
36000	724.76\\
38000	738.26\\
40000	738.41\\
42000	723.97\\
44000	717.07\\
46000	728.96\\
48000	722.63\\
50000	722.5\\
52000	732.11\\
54000	723.22\\
56000	728.52\\
58000	728.85\\
60000	741.11\\
62000	751.25\\
64000	701.54\\
66000	715.97\\
68000	747.28\\
70000	739.44\\
72000	722.73\\
74000	726.52\\
76000	707.77\\
78000	705.67\\
80000	721.74\\
82000	729.17\\
84000	759.25\\
86000	722.7\\
88000	720.18\\
90000	703.99\\
92000	739.08\\
94000	726.26\\
96000	740.83\\
98000	720.8\\
1e+05	724.45\\
};
\addplot [color=mycolor92, forget plot]
  table[row sep=crcr]{%
0	4395.3\\
2000	1982.3\\
4000	969.42\\
6000	910.78\\
8000	908.19\\
10000	888.44\\
12000	866.21\\
14000	908.72\\
16000	912.7\\
18000	887.31\\
20000	865.25\\
22000	872.71\\
24000	912.95\\
26000	891.5\\
30000	884.64\\
32000	886.28\\
34000	894.52\\
36000	895.99\\
38000	906.95\\
40000	923.35\\
42000	895.96\\
44000	897.98\\
46000	886.27\\
48000	861.56\\
50000	888.91\\
52000	894.8\\
54000	884.45\\
56000	888.29\\
58000	862.1\\
62000	897.06\\
64000	891.44\\
66000	910.78\\
68000	891.63\\
70000	894.21\\
72000	909.84\\
76000	893.29\\
78000	881.04\\
80000	901.79\\
82000	898.87\\
84000	875.41\\
86000	905.08\\
90000	912.04\\
92000	874\\
94000	882.94\\
96000	884.57\\
98000	907.33\\
1e+05	942.27\\
};
\addplot [color=mycolor93, forget plot]
  table[row sep=crcr]{%
0	4395.3\\
2000	1901.2\\
4000	808.15\\
6000	741.84\\
8000	758.22\\
12000	736.87\\
14000	768.98\\
16000	736.15\\
18000	773.01\\
20000	773.5\\
22000	760.9\\
24000	727.23\\
26000	737.96\\
28000	728.34\\
30000	772.53\\
32000	756.7\\
34000	737.96\\
36000	741.57\\
38000	740.69\\
40000	763.87\\
42000	725.56\\
44000	749.06\\
46000	724.32\\
48000	756.71\\
50000	781.52\\
52000	733.1\\
54000	773.8\\
56000	756.71\\
58000	764.15\\
60000	754.98\\
62000	727.61\\
64000	718.34\\
66000	777.2\\
68000	756\\
70000	761.73\\
72000	738.26\\
74000	735.76\\
76000	771.7\\
78000	720.59\\
80000	747.02\\
82000	764.17\\
84000	750.38\\
88000	719.74\\
90000	765.35\\
92000	772.78\\
94000	776.56\\
96000	762.57\\
98000	733.33\\
1e+05	750.78\\
};
\addplot [color=mycolor94, forget plot]
  table[row sep=crcr]{%
0	4395.3\\
2000	1875.9\\
4000	812.82\\
6000	726.3\\
8000	720.29\\
10000	701.92\\
12000	699.93\\
14000	708.17\\
16000	705.45\\
18000	718.45\\
20000	721.6\\
22000	697.16\\
24000	710.47\\
26000	735.87\\
28000	719.16\\
30000	695.21\\
32000	726.36\\
34000	720.61\\
36000	719.45\\
38000	711.19\\
40000	715.2\\
42000	732.08\\
44000	718.58\\
46000	710.51\\
48000	728.82\\
50000	699.42\\
52000	706.75\\
54000	706.25\\
56000	717.54\\
58000	706.57\\
60000	712.1\\
62000	714.36\\
64000	724.14\\
66000	712.38\\
68000	717.24\\
70000	698.29\\
72000	724.98\\
74000	713.29\\
76000	717.8\\
78000	684.6\\
82000	718.86\\
84000	708.8\\
86000	726.83\\
88000	730.71\\
90000	726.45\\
92000	705.65\\
96000	736.46\\
98000	702.13\\
1e+05	702.83\\
};
\addplot [color=mycolor95, forget plot]
  table[row sep=crcr]{%
0	4395.3\\
2000	1624.8\\
4000	517.73\\
6000	433.37\\
8000	422.53\\
10000	418.03\\
12000	430.99\\
14000	437.51\\
16000	428.87\\
18000	429.85\\
20000	433.94\\
22000	414.1\\
24000	406.26\\
26000	405.76\\
28000	410.8\\
30000	431.02\\
32000	416.97\\
34000	421.31\\
36000	429.26\\
38000	431.94\\
40000	422.65\\
42000	400.09\\
44000	433.51\\
46000	442.06\\
48000	420.59\\
50000	418.52\\
52000	427.28\\
54000	421.96\\
56000	434.86\\
58000	414.47\\
60000	399.42\\
62000	425.04\\
64000	415.7\\
66000	409.87\\
68000	425.18\\
70000	427.63\\
72000	422.84\\
74000	428.61\\
76000	439.57\\
78000	435.39\\
80000	425.74\\
82000	424.66\\
84000	437.17\\
86000	426.97\\
88000	422.18\\
90000	413.49\\
92000	435.44\\
94000	434.85\\
96000	425.79\\
98000	432.15\\
1e+05	436.25\\
};
\addplot [color=mycolor96, forget plot]
  table[row sep=crcr]{%
0	4395.3\\
2000	1807\\
4000	594.27\\
6000	476.69\\
8000	474.79\\
12000	459.73\\
14000	456.49\\
16000	485.08\\
18000	461.01\\
20000	461.16\\
22000	462.89\\
24000	423.33\\
26000	439.21\\
28000	453.29\\
30000	455.82\\
32000	454.35\\
34000	441.22\\
36000	447.25\\
38000	459.34\\
40000	463.56\\
42000	484.74\\
44000	428.36\\
46000	423.67\\
48000	449.28\\
50000	447.62\\
52000	479.74\\
54000	452.79\\
56000	424.05\\
58000	439.25\\
60000	462.42\\
62000	415.4\\
64000	427.23\\
66000	474.98\\
68000	443.45\\
70000	431.74\\
72000	409.96\\
74000	442.43\\
76000	474.5\\
78000	458.51\\
80000	436.73\\
82000	423.47\\
84000	454.14\\
86000	449.48\\
88000	451.33\\
90000	438.23\\
94000	434.56\\
96000	440.41\\
98000	433.69\\
1e+05	442.52\\
};
\addplot [color=mycolor97, forget plot]
  table[row sep=crcr]{%
0	4395.3\\
2000	1618.3\\
4000	418.91\\
6000	342.07\\
8000	333.17\\
10000	346.69\\
12000	333.22\\
14000	333.56\\
16000	303.98\\
18000	336.7\\
20000	305.6\\
22000	326.19\\
24000	319.59\\
26000	329.92\\
30000	305.55\\
34000	315.65\\
36000	325.27\\
38000	312.93\\
40000	318.51\\
42000	345.42\\
44000	330\\
46000	326.5\\
48000	301.45\\
50000	311.42\\
52000	303.35\\
54000	306.3\\
56000	313.16\\
58000	306.3\\
60000	297.76\\
62000	331.64\\
64000	313.73\\
66000	313.44\\
68000	347.98\\
70000	317.03\\
72000	322.17\\
74000	322.34\\
76000	298.32\\
78000	316.2\\
80000	293.36\\
82000	302.65\\
84000	321.54\\
86000	335.03\\
88000	321.47\\
90000	326.68\\
92000	303.13\\
94000	295.67\\
96000	308.77\\
98000	293.8\\
1e+05	296.32\\
};
\addplot [color=mycolor98]
  table[row sep=crcr]{%
0	4395.3\\
2000	1950.4\\
4000	610.8\\
6000	471.16\\
8000	438.07\\
10000	431.78\\
12000	421.35\\
14000	419.4\\
16000	439.38\\
18000	431.74\\
20000	422.84\\
22000	420.96\\
24000	439.02\\
26000	432.72\\
28000	441.74\\
30000	449.63\\
32000	429.38\\
34000	448.87\\
36000	434.27\\
38000	428.7\\
40000	436.17\\
42000	440.92\\
44000	423.28\\
46000	415.61\\
48000	428.24\\
50000	418.29\\
52000	437.4\\
54000	427.25\\
56000	419.83\\
58000	439.7\\
60000	438.39\\
62000	431.66\\
64000	406.03\\
66000	403.94\\
68000	410.61\\
70000	424.92\\
72000	432.76\\
74000	419.64\\
76000	443.76\\
78000	440.58\\
80000	439.91\\
82000	430.52\\
84000	454.46\\
86000	466.72\\
88000	423.67\\
90000	425.23\\
92000	430.88\\
96000	421.68\\
98000	443.58\\
1e+05	445.68\\
};
\addlegendentry{Classe k=10}

            % \input{../../../../.tex/20/KS/SizeEvolutionsfig6.tex}
            \legend{}
        \end{groupplot}

        \begin{groupplot}[
            group style={
                group name=Row4,
                plotGroup,
                group size=2 by 1,
            },
            plotRow,
            row5Keys,
            xTickScaleLabelEmpty,
            plotLegendSE,
            %
            xmin=0,xmax=1e+05,
            ymin=100,ymax=1.5689e+05,
        ]
            \nextgroupplot[%
                at={($(Row3 c1r1.south west)-(0,\yPlotSep em)-(0,\plotHeight cm)$)},
                % xmin=0,xmax=1e+05,
                % ymin=100,ymax=1.5698e+05,
            ] % This file was created by matlab2tikz.
%
\definecolor{mycolor1}{rgb}{0.00146,0.00047,0.01387}%
\definecolor{mycolor2}{rgb}{0.01166,0.00942,0.06346}%
\definecolor{mycolor3}{rgb}{0.02943,0.02150,0.11462}%
\definecolor{mycolor4}{rgb}{0.06134,0.03659,0.17764}%
\definecolor{mycolor5}{rgb}{0.08741,0.04456,0.22481}%
\definecolor{mycolor6}{rgb}{0.12291,0.04754,0.28162}%
\definecolor{mycolor7}{rgb}{0.14907,0.04547,0.31709}%
\definecolor{mycolor8}{rgb}{0.18343,0.04033,0.35497}%
\definecolor{mycolor9}{rgb}{0.21795,0.03662,0.38352}%
\definecolor{mycolor10}{rgb}{0.24497,0.03705,0.40001}%
\definecolor{mycolor11}{rgb}{0.27135,0.04092,0.41198}%
\definecolor{mycolor12}{rgb}{0.29076,0.04564,0.41864}%
\definecolor{mycolor13}{rgb}{0.31628,0.05349,0.42512}%
\definecolor{mycolor14}{rgb}{0.33522,0.06006,0.42852}%
\definecolor{mycolor15}{rgb}{0.36028,0.06925,0.43150}%
\definecolor{mycolor16}{rgb}{0.37900,0.07625,0.43272}%
\definecolor{mycolor17}{rgb}{0.39767,0.08326,0.43318}%
\definecolor{mycolor18}{rgb}{0.41633,0.09020,0.43294}%
\definecolor{mycolor19}{rgb}{0.42877,0.09479,0.43241}%
\definecolor{mycolor20}{rgb}{0.44743,0.10160,0.43108}%
\definecolor{mycolor21}{rgb}{0.46610,0.10832,0.42912}%
\definecolor{mycolor22}{rgb}{0.47856,0.11276,0.42747}%
\definecolor{mycolor23}{rgb}{0.49726,0.11938,0.42449}%
\definecolor{mycolor24}{rgb}{0.50973,0.12377,0.42216}%
\definecolor{mycolor25}{rgb}{0.52221,0.12815,0.41955}%
\definecolor{mycolor26}{rgb}{0.53468,0.13253,0.41667}%
\definecolor{mycolor27}{rgb}{0.55339,0.13913,0.41183}%
\definecolor{mycolor28}{rgb}{0.56585,0.14357,0.40826}%
\definecolor{mycolor29}{rgb}{0.57830,0.14804,0.40441}%
\definecolor{mycolor30}{rgb}{0.59073,0.15256,0.40029}%
\definecolor{mycolor31}{rgb}{0.60314,0.15715,0.39589}%
\definecolor{mycolor32}{rgb}{0.61551,0.16182,0.39122}%
\definecolor{mycolor33}{rgb}{0.62169,0.16418,0.38878}%
\definecolor{mycolor34}{rgb}{0.63400,0.16899,0.38370}%
\definecolor{mycolor35}{rgb}{0.64626,0.17391,0.37836}%
\definecolor{mycolor36}{rgb}{0.65846,0.17896,0.37275}%
\definecolor{mycolor37}{rgb}{0.66454,0.18154,0.36985}%
\definecolor{mycolor38}{rgb}{0.67664,0.18681,0.36385}%
\definecolor{mycolor39}{rgb}{0.68865,0.19224,0.35760}%
\definecolor{mycolor40}{rgb}{0.69463,0.19502,0.35439}%
\definecolor{mycolor41}{rgb}{0.70650,0.20073,0.34778}%
\definecolor{mycolor42}{rgb}{0.71240,0.20366,0.34438}%
\definecolor{mycolor43}{rgb}{0.72410,0.20967,0.33742}%
\definecolor{mycolor44}{rgb}{0.72991,0.21276,0.33386}%
\definecolor{mycolor45}{rgb}{0.74142,0.21911,0.32658}%
\definecolor{mycolor46}{rgb}{0.74713,0.22238,0.32286}%
\definecolor{mycolor47}{rgb}{0.75279,0.22571,0.31909}%
\definecolor{mycolor48}{rgb}{0.76401,0.23255,0.31140}%
\definecolor{mycolor49}{rgb}{0.76956,0.23608,0.30749}%
\definecolor{mycolor50}{rgb}{0.77506,0.23967,0.30353}%
\definecolor{mycolor51}{rgb}{0.79129,0.25086,0.29139}%
\definecolor{mycolor52}{rgb}{0.79661,0.25473,0.28726}%
\definecolor{mycolor53}{rgb}{0.80187,0.25867,0.28310}%
\definecolor{mycolor54}{rgb}{0.81224,0.26679,0.27466}%
\definecolor{mycolor55}{rgb}{0.81734,0.27095,0.27039}%
\definecolor{mycolor56}{rgb}{0.82737,0.27952,0.26175}%
\definecolor{mycolor57}{rgb}{0.83230,0.28391,0.25738}%
\definecolor{mycolor58}{rgb}{0.83717,0.28839,0.25299}%
\definecolor{mycolor59}{rgb}{0.84671,0.29756,0.24411}%
\definecolor{mycolor60}{rgb}{0.85138,0.30226,0.23964}%
\definecolor{mycolor61}{rgb}{0.85599,0.30704,0.23513}%
\definecolor{mycolor62}{rgb}{0.86053,0.31189,0.23061}%
\definecolor{mycolor63}{rgb}{0.87374,0.32691,0.21689}%
\definecolor{mycolor64}{rgb}{0.87800,0.33206,0.21227}%
\definecolor{mycolor65}{rgb}{0.89034,0.34796,0.19829}%
\definecolor{mycolor66}{rgb}{0.89819,0.35891,0.18886}%
\definecolor{mycolor67}{rgb}{0.90200,0.36449,0.18412}%
\definecolor{mycolor68}{rgb}{0.90573,0.37014,0.17935}%
\definecolor{mycolor69}{rgb}{0.90939,0.37586,0.17456}%
\definecolor{mycolor70}{rgb}{0.91297,0.38164,0.16975}%
\definecolor{mycolor71}{rgb}{0.91646,0.38748,0.16492}%
\definecolor{mycolor72}{rgb}{0.91988,0.39339,0.16007}%
\definecolor{mycolor73}{rgb}{0.92322,0.39936,0.15519}%
\definecolor{mycolor74}{rgb}{0.92647,0.40539,0.15029}%
\definecolor{mycolor75}{rgb}{0.92964,0.41148,0.14537}%
\definecolor{mycolor76}{rgb}{0.93274,0.41763,0.14042}%
\definecolor{mycolor77}{rgb}{0.93575,0.42383,0.13544}%
\definecolor{mycolor78}{rgb}{0.93868,0.43009,0.13044}%
\definecolor{mycolor79}{rgb}{0.94152,0.43640,0.12541}%
\definecolor{mycolor80}{rgb}{0.94429,0.44277,0.12035}%
\definecolor{mycolor81}{rgb}{0.94956,0.45566,0.11016}%
\definecolor{mycolor82}{rgb}{0.95208,0.46218,0.10503}%
\definecolor{mycolor83}{rgb}{0.96129,0.48872,0.08429}%
\definecolor{mycolor84}{rgb}{0.96540,0.50225,0.07386}%
\definecolor{mycolor85}{rgb}{0.96732,0.50908,0.06866}%
\definecolor{mycolor86}{rgb}{0.97259,0.52980,0.05332}%
\definecolor{mycolor87}{rgb}{0.97568,0.54380,0.04362}%
\definecolor{mycolor88}{rgb}{0.98082,0.57221,0.02851}%
\definecolor{mycolor89}{rgb}{0.98189,0.57939,0.02625}%
\definecolor{mycolor90}{rgb}{0.98378,0.59385,0.02377}%
\definecolor{mycolor91}{rgb}{0.98532,0.60842,0.02420}%
\definecolor{mycolor92}{rgb}{0.98696,0.63048,0.03091}%
\definecolor{mycolor93}{rgb}{0.98793,0.66025,0.05175}%
\definecolor{mycolor94}{rgb}{0.98771,0.68281,0.07249}%
\definecolor{mycolor95}{rgb}{0.97750,0.78226,0.18592}%
\definecolor{mycolor96}{rgb}{0.96624,0.83619,0.26153}%
\definecolor{mycolor97}{rgb}{0.96252,0.85148,0.28555}%
\definecolor{mycolor98}{rgb}{0.98836,0.99836,0.64492}%
%
\addplot [color=mycolor1]
  table[row sep=crcr]{%
0	4395.3\\
2000	1.0922e+05\\
4000	1.4937e+05\\
6000	1.4897e+05\\
10000	1.5192e+05\\
16000	1.5095e+05\\
18000	1.4784e+05\\
22000	1.5219e+05\\
24000	1.5168e+05\\
26000	1.4951e+05\\
28000	1.5306e+05\\
30000	1.517e+05\\
32000	1.4985e+05\\
34000	1.5123e+05\\
36000	1.4854e+05\\
40000	1.5167e+05\\
42000	1.5442e+05\\
44000	1.5563e+05\\
48000	1.524e+05\\
50000	1.4893e+05\\
52000	1.5078e+05\\
54000	1.548e+05\\
56000	1.5525e+05\\
58000	1.5318e+05\\
60000	1.5036e+05\\
62000	1.499e+05\\
64000	1.5018e+05\\
66000	1.5184e+05\\
68000	1.5263e+05\\
70000	1.5079e+05\\
74000	1.5182e+05\\
76000	1.5437e+05\\
78000	1.5592e+05\\
80000	1.5145e+05\\
82000	1.5071e+05\\
84000	1.5138e+05\\
86000	1.5099e+05\\
88000	1.489e+05\\
90000	1.5143e+05\\
1e+05	1.5125e+05\\
};
\addlegendentry{Classe k=283}

\addplot [color=mycolor2, forget plot]
  table[row sep=crcr]{%
0	4395.3\\
2000	26004\\
4000	26062\\
6000	26198\\
8000	26238\\
10000	26018\\
12000	25367\\
14000	25707\\
16000	25879\\
18000	26202\\
22000	25215\\
24000	26450\\
26000	25984\\
28000	26379\\
30000	25791\\
32000	25861\\
34000	25663\\
36000	25818\\
38000	26159\\
40000	26039\\
42000	26224\\
44000	26529\\
46000	25081\\
48000	26003\\
50000	26843\\
52000	25807\\
56000	26240\\
58000	26145\\
60000	26534\\
62000	26752\\
64000	25947\\
66000	26477\\
68000	25707\\
70000	26660\\
72000	25724\\
74000	25842\\
76000	26938\\
78000	26323\\
80000	25810\\
82000	25537\\
84000	26490\\
86000	26058\\
88000	25995\\
90000	25699\\
92000	26689\\
94000	25652\\
96000	26479\\
98000	26627\\
1e+05	25629\\
};
\addplot [color=mycolor3, forget plot]
  table[row sep=crcr]{%
0	4395.3\\
2000	36378\\
4000	38746\\
6000	39325\\
8000	39317\\
14000	39973\\
16000	39186\\
20000	39715\\
22000	38938\\
24000	39494\\
26000	39331\\
28000	39547\\
30000	39166\\
32000	39510\\
34000	39242\\
36000	40227\\
38000	39507\\
40000	40344\\
42000	40256\\
44000	38651\\
46000	40396\\
48000	40168\\
50000	39361\\
52000	39419\\
54000	40037\\
56000	39102\\
58000	38777\\
60000	39661\\
62000	40221\\
64000	38656\\
66000	39313\\
68000	40457\\
70000	40217\\
72000	40242\\
74000	39736\\
76000	39484\\
78000	40155\\
80000	40281\\
82000	39915\\
84000	39401\\
86000	39642\\
90000	39369\\
92000	38178\\
94000	39950\\
96000	38788\\
98000	40212\\
1e+05	40092\\
};
\addplot [color=mycolor4, forget plot]
  table[row sep=crcr]{%
0	4395.3\\
2000	14990\\
4000	14178\\
6000	14494\\
8000	14236\\
10000	14185\\
12000	14612\\
14000	14381\\
16000	14597\\
18000	14583\\
20000	14307\\
22000	14580\\
24000	14249\\
26000	14462\\
28000	14388\\
30000	14259\\
32000	14612\\
34000	14481\\
36000	13718\\
38000	14415\\
42000	14605\\
44000	14490\\
46000	14536\\
48000	14875\\
50000	14419\\
54000	14128\\
56000	14602\\
58000	14437\\
60000	14187\\
62000	14687\\
64000	14320\\
66000	14706\\
68000	14782\\
70000	14204\\
72000	14702\\
74000	14740\\
76000	14396\\
78000	14417\\
80000	14332\\
82000	14518\\
84000	14310\\
86000	14425\\
88000	14475\\
90000	14283\\
92000	14485\\
94000	14185\\
96000	14582\\
98000	14229\\
1e+05	13954\\
};
\addplot [color=mycolor5, forget plot]
  table[row sep=crcr]{%
0	4395.3\\
2000	52470\\
4000	65974\\
6000	64872\\
8000	64454\\
10000	65927\\
12000	61548\\
14000	63524\\
16000	63585\\
18000	65300\\
20000	64230\\
24000	63991\\
26000	64793\\
28000	64726\\
30000	64066\\
32000	62297\\
34000	62786\\
36000	61771\\
38000	64882\\
40000	63569\\
42000	62740\\
44000	62437\\
46000	63412\\
48000	64010\\
50000	65372\\
52000	63250\\
54000	62764\\
56000	65057\\
58000	63588\\
60000	64221\\
62000	63107\\
66000	62829\\
68000	63292\\
70000	62946\\
72000	65760\\
74000	65390\\
76000	64023\\
78000	62184\\
80000	60721\\
82000	63591\\
84000	63060\\
86000	62103\\
88000	63489\\
90000	62487\\
92000	63567\\
94000	64974\\
96000	63175\\
98000	63833\\
1e+05	62730\\
};
\addplot [color=mycolor6, forget plot]
  table[row sep=crcr]{%
0	4395.3\\
2000	41984\\
4000	53801\\
6000	55176\\
8000	57020\\
10000	56768\\
12000	57189\\
14000	57014\\
16000	57293\\
18000	54324\\
22000	56785\\
24000	54749\\
28000	54767\\
30000	56756\\
32000	55245\\
34000	55037\\
36000	55970\\
38000	55783\\
40000	56256\\
42000	54767\\
44000	55272\\
46000	55483\\
48000	55001\\
50000	56925\\
52000	55287\\
54000	54803\\
56000	56202\\
58000	55120\\
60000	56524\\
62000	55165\\
64000	56245\\
66000	56414\\
68000	56036\\
70000	56621\\
72000	54052\\
74000	53475\\
76000	54630\\
78000	55042\\
80000	57973\\
82000	56285\\
84000	56782\\
86000	57838\\
88000	55124\\
90000	54701\\
92000	55978\\
94000	56232\\
96000	56719\\
1e+05	55387\\
};
\addplot [color=mycolor7, forget plot]
  table[row sep=crcr]{%
0	4395.3\\
2000	24178\\
4000	27776\\
6000	27988\\
8000	27869\\
10000	26898\\
12000	28821\\
14000	28999\\
16000	28546\\
18000	27042\\
20000	28622\\
22000	28074\\
24000	27058\\
26000	28109\\
28000	28258\\
30000	27721\\
32000	27860\\
34000	28484\\
36000	28688\\
38000	29218\\
40000	27953\\
42000	28387\\
44000	27901\\
46000	27653\\
48000	28745\\
50000	28189\\
52000	28670\\
54000	28883\\
56000	28896\\
58000	27478\\
60000	28085\\
62000	28536\\
64000	28633\\
66000	27832\\
68000	28355\\
70000	28222\\
72000	28200\\
74000	28728\\
76000	26910\\
78000	27995\\
80000	28309\\
82000	29010\\
84000	27587\\
86000	28506\\
88000	28573\\
90000	27868\\
92000	28452\\
94000	28660\\
96000	29045\\
98000	28762\\
1e+05	27798\\
};
\addplot [color=mycolor8, forget plot]
  table[row sep=crcr]{%
0	4395.3\\
2000	20527\\
4000	21197\\
6000	22136\\
8000	20139\\
12000	20716\\
14000	19703\\
16000	20632\\
18000	20151\\
20000	20774\\
22000	21608\\
24000	20726\\
26000	21079\\
28000	20500\\
30000	20509\\
32000	20318\\
36000	20757\\
38000	20405\\
40000	20478\\
42000	20639\\
44000	20679\\
46000	20943\\
48000	20192\\
50000	20807\\
52000	20518\\
54000	20138\\
56000	20173\\
58000	20845\\
60000	20571\\
62000	20785\\
64000	20336\\
66000	20598\\
68000	20579\\
70000	20382\\
72000	20512\\
74000	20749\\
76000	21654\\
78000	21057\\
80000	20604\\
82000	19796\\
84000	19945\\
86000	20186\\
88000	20538\\
90000	21484\\
92000	20802\\
94000	20507\\
96000	20766\\
98000	20492\\
1e+05	21313\\
};
\addplot [color=mycolor9, forget plot]
  table[row sep=crcr]{%
0	4395.3\\
2000	13069\\
4000	14159\\
6000	13959\\
8000	14182\\
10000	14078\\
12000	14149\\
14000	14096\\
16000	14582\\
18000	14404\\
22000	13727\\
24000	14030\\
26000	13914\\
28000	13936\\
30000	13707\\
32000	14073\\
34000	13700\\
36000	14311\\
38000	13770\\
40000	14123\\
42000	14368\\
44000	14348\\
46000	14621\\
48000	13921\\
50000	14258\\
52000	13886\\
54000	14158\\
56000	14095\\
58000	14375\\
60000	14057\\
62000	13957\\
64000	14226\\
66000	13877\\
68000	14210\\
70000	13882\\
72000	13760\\
74000	14185\\
78000	14102\\
80000	13843\\
82000	14048\\
84000	13984\\
86000	14034\\
88000	14181\\
90000	14063\\
92000	13889\\
94000	14256\\
98000	13971\\
1e+05	14236\\
};
\addplot [color=mycolor10, forget plot]
  table[row sep=crcr]{%
0	4395.3\\
2000	18934\\
4000	21114\\
6000	22337\\
8000	22355\\
10000	22178\\
12000	22120\\
14000	22230\\
16000	21954\\
18000	22668\\
20000	22051\\
24000	22231\\
28000	22107\\
30000	22248\\
32000	21925\\
34000	22249\\
36000	22355\\
38000	22223\\
40000	21725\\
42000	22101\\
44000	21748\\
46000	21553\\
48000	21263\\
50000	22579\\
52000	22037\\
54000	21868\\
56000	21457\\
58000	22156\\
60000	21917\\
62000	21996\\
64000	22429\\
66000	22101\\
68000	22583\\
72000	22201\\
74000	21623\\
76000	21506\\
80000	22470\\
82000	22432\\
84000	22314\\
86000	22423\\
88000	22170\\
90000	22364\\
92000	21559\\
96000	22105\\
98000	22019\\
1e+05	22130\\
};
\addplot [color=mycolor11, forget plot]
  table[row sep=crcr]{%
0	4395.3\\
2000	9660.1\\
4000	8903.8\\
6000	8767.3\\
8000	8710\\
10000	8552.7\\
14000	8586.8\\
16000	8334.4\\
18000	8993.9\\
20000	8823.3\\
24000	8381.4\\
26000	8872.9\\
28000	8629.9\\
30000	8480.2\\
32000	8655.6\\
34000	8725.6\\
36000	8975.7\\
38000	8810.3\\
40000	8953.5\\
42000	8706.1\\
44000	8873.5\\
46000	8879.2\\
48000	8457.9\\
50000	8642.6\\
52000	8693\\
54000	8798.7\\
56000	8452.2\\
58000	8687.9\\
60000	8662.9\\
62000	8731.1\\
64000	8568.6\\
66000	8707.8\\
68000	8762.4\\
70000	8727.4\\
72000	8773.6\\
74000	8787.5\\
76000	8589.8\\
78000	8525\\
80000	8705.6\\
82000	8713.8\\
84000	8484.3\\
86000	8684.5\\
88000	8785.7\\
90000	8681.4\\
92000	8753.5\\
94000	8494.7\\
96000	8654.4\\
98000	8955.1\\
1e+05	8827.6\\
};
\addplot [color=mycolor12, forget plot]
  table[row sep=crcr]{%
0	4395.3\\
2000	13202\\
4000	14046\\
6000	14331\\
8000	14919\\
10000	14943\\
12000	15508\\
16000	14780\\
18000	14906\\
20000	14834\\
22000	15046\\
24000	14301\\
26000	14806\\
28000	14901\\
30000	15173\\
32000	14346\\
34000	14876\\
36000	14960\\
38000	14553\\
40000	14784\\
42000	14903\\
44000	14689\\
46000	14955\\
48000	14827\\
50000	14843\\
52000	14807\\
54000	14709\\
56000	13724\\
58000	14461\\
60000	15311\\
62000	14757\\
64000	14929\\
66000	14481\\
68000	14551\\
70000	15227\\
72000	14121\\
74000	14961\\
76000	14315\\
78000	15109\\
80000	14514\\
82000	15081\\
84000	14876\\
86000	14165\\
88000	14896\\
90000	15062\\
92000	14676\\
94000	14990\\
96000	14589\\
98000	14773\\
1e+05	14909\\
};
\addplot [color=mycolor13, forget plot]
  table[row sep=crcr]{%
0	4395.3\\
2000	10761\\
4000	11030\\
6000	11060\\
10000	11711\\
12000	11151\\
14000	11165\\
16000	10978\\
18000	11332\\
20000	11444\\
22000	11192\\
24000	11216\\
26000	11171\\
28000	11459\\
30000	11177\\
32000	10942\\
34000	10938\\
36000	11064\\
38000	11422\\
40000	10785\\
42000	11351\\
44000	11165\\
46000	10642\\
48000	11263\\
50000	11126\\
54000	11263\\
56000	11176\\
58000	11040\\
60000	11294\\
62000	11004\\
64000	11291\\
66000	10889\\
68000	11163\\
70000	11342\\
72000	11174\\
74000	11100\\
76000	10741\\
78000	10789\\
80000	11395\\
82000	11351\\
84000	11108\\
86000	11181\\
88000	11053\\
90000	11451\\
92000	11233\\
94000	11085\\
96000	11168\\
98000	11000\\
1e+05	10919\\
};
\addplot [color=mycolor14, forget plot]
  table[row sep=crcr]{%
0	4395.3\\
2000	20144\\
4000	23893\\
6000	24466\\
8000	24205\\
10000	23879\\
12000	24547\\
14000	24782\\
16000	24513\\
18000	23608\\
20000	23504\\
22000	24381\\
24000	24698\\
26000	24427\\
28000	24779\\
32000	26047\\
34000	24009\\
36000	25021\\
38000	24966\\
40000	24530\\
42000	23834\\
44000	23774\\
46000	24589\\
48000	24167\\
50000	23652\\
52000	24663\\
54000	23563\\
56000	24073\\
58000	24488\\
60000	24145\\
62000	24640\\
64000	24707\\
66000	23463\\
68000	23894\\
70000	24029\\
72000	24635\\
74000	23986\\
76000	24382\\
78000	23531\\
80000	23854\\
82000	23931\\
84000	24821\\
86000	24165\\
88000	25072\\
92000	23800\\
94000	24164\\
96000	23335\\
98000	24114\\
1e+05	24279\\
};
\addplot [color=mycolor15, forget plot]
  table[row sep=crcr]{%
0	4395.3\\
2000	14867\\
4000	16337\\
6000	16640\\
8000	16452\\
10000	16911\\
12000	17250\\
14000	17191\\
16000	17265\\
18000	17461\\
20000	17344\\
22000	17064\\
24000	16848\\
26000	17131\\
28000	17001\\
30000	17612\\
32000	17173\\
34000	17444\\
36000	17390\\
38000	17061\\
40000	17175\\
42000	16823\\
44000	16884\\
48000	16861\\
50000	16550\\
52000	17074\\
54000	17187\\
56000	16622\\
58000	16852\\
60000	17318\\
62000	17565\\
64000	17250\\
66000	16678\\
68000	16914\\
70000	17068\\
72000	16251\\
74000	17083\\
76000	17194\\
78000	17243\\
80000	16425\\
82000	16404\\
84000	17302\\
86000	17085\\
88000	17114\\
90000	17461\\
92000	17173\\
94000	17228\\
96000	16562\\
98000	17183\\
1e+05	16990\\
};
\addplot [color=mycolor16, forget plot]
  table[row sep=crcr]{%
0	4395.3\\
2000	7696.4\\
4000	7064.4\\
6000	6817.4\\
8000	7124.8\\
10000	6976.6\\
12000	7200.9\\
14000	6811.9\\
16000	6766.6\\
18000	7155.1\\
20000	7062.2\\
22000	6866.5\\
24000	7010.3\\
26000	6892\\
28000	7018.7\\
30000	7010.9\\
32000	6962.3\\
34000	6956.1\\
36000	6871.8\\
38000	6936.7\\
40000	6901.5\\
42000	7193.1\\
44000	6970.1\\
46000	6847.4\\
48000	6964.2\\
50000	7015.5\\
52000	7198.8\\
54000	6815.8\\
56000	7097.7\\
58000	7264.1\\
60000	7154\\
62000	6871\\
64000	7193.6\\
66000	7045.1\\
68000	6718.1\\
72000	7125.1\\
74000	6933.6\\
76000	7174.1\\
78000	6797.9\\
80000	6948.3\\
82000	6896.2\\
84000	7011.7\\
86000	6880.3\\
88000	7093.9\\
90000	7132.3\\
94000	7059.2\\
96000	7057\\
98000	6931.8\\
1e+05	7091.2\\
};
\addplot [color=mycolor17, forget plot]
  table[row sep=crcr]{%
0	4395.3\\
2000	13848\\
4000	15391\\
6000	15818\\
8000	16193\\
10000	15744\\
12000	16315\\
14000	16088\\
16000	16374\\
20000	15895\\
22000	15720\\
24000	16801\\
26000	16168\\
28000	16603\\
30000	16121\\
32000	15941\\
34000	15810\\
36000	15749\\
38000	15751\\
40000	15928\\
42000	15747\\
44000	16095\\
46000	15572\\
48000	16378\\
50000	16029\\
52000	15780\\
54000	16080\\
56000	15736\\
58000	16213\\
60000	16028\\
62000	16324\\
64000	15981\\
66000	15520\\
68000	15691\\
70000	15988\\
72000	15941\\
74000	16181\\
76000	16363\\
78000	15760\\
80000	16498\\
84000	16228\\
86000	16475\\
88000	16623\\
90000	16411\\
92000	16666\\
94000	16392\\
96000	16294\\
98000	16072\\
1e+05	16054\\
};
\addplot [color=mycolor18, forget plot]
  table[row sep=crcr]{%
0	4395.3\\
2000	20309\\
4000	24038\\
6000	23924\\
8000	24278\\
10000	24277\\
12000	24472\\
14000	24249\\
16000	24209\\
18000	25295\\
22000	24094\\
24000	24844\\
26000	24666\\
28000	24684\\
30000	24055\\
32000	23850\\
34000	24283\\
36000	24503\\
38000	25127\\
40000	24931\\
42000	25790\\
44000	23864\\
46000	25258\\
48000	25091\\
50000	24819\\
52000	24893\\
54000	24779\\
56000	24257\\
58000	24684\\
60000	24552\\
62000	25072\\
64000	24475\\
66000	24970\\
68000	24353\\
70000	23987\\
72000	24531\\
74000	25494\\
76000	24515\\
78000	25035\\
82000	25089\\
84000	24165\\
86000	24912\\
88000	24730\\
90000	24418\\
92000	24980\\
94000	24415\\
96000	24586\\
98000	23883\\
1e+05	24856\\
};
\addplot [color=mycolor19, forget plot]
  table[row sep=crcr]{%
0	4395.3\\
2000	13376\\
4000	14291\\
6000	15370\\
8000	14929\\
10000	14814\\
14000	15380\\
16000	14996\\
18000	15198\\
20000	15354\\
22000	15156\\
24000	15026\\
26000	15318\\
28000	14896\\
30000	14716\\
32000	15350\\
34000	15177\\
36000	15133\\
38000	14768\\
40000	15189\\
42000	14420\\
44000	15250\\
46000	14777\\
48000	14554\\
50000	14993\\
52000	14735\\
54000	14534\\
56000	14568\\
58000	15095\\
60000	14842\\
62000	15146\\
64000	14991\\
66000	14960\\
68000	15436\\
70000	14504\\
72000	14726\\
74000	14848\\
76000	15089\\
78000	14546\\
80000	14522\\
82000	14947\\
86000	15071\\
88000	15049\\
90000	15465\\
92000	14723\\
94000	14933\\
96000	14822\\
98000	14467\\
1e+05	14792\\
};
\addplot [color=mycolor20, forget plot]
  table[row sep=crcr]{%
0	4395.3\\
2000	16583\\
4000	19214\\
6000	18510\\
8000	18914\\
10000	18849\\
12000	19507\\
14000	18809\\
16000	19188\\
18000	19652\\
20000	19138\\
22000	19278\\
24000	19342\\
26000	19767\\
28000	19247\\
30000	19607\\
32000	19728\\
34000	19331\\
36000	19428\\
38000	19695\\
40000	19188\\
42000	19146\\
44000	19505\\
46000	18913\\
48000	19316\\
50000	18933\\
52000	19298\\
54000	19745\\
56000	19660\\
58000	19708\\
60000	19960\\
64000	19429\\
66000	19403\\
68000	19310\\
70000	19125\\
72000	18997\\
74000	19807\\
76000	19308\\
78000	19405\\
80000	19089\\
82000	18587\\
84000	19145\\
88000	19981\\
90000	19162\\
92000	19708\\
94000	19143\\
96000	19917\\
98000	18851\\
1e+05	19213\\
};
\addplot [color=mycolor21, forget plot]
  table[row sep=crcr]{%
0	4395.3\\
2000	8698.8\\
4000	9470.2\\
6000	9559\\
8000	9546.4\\
10000	9799\\
12000	9861.2\\
14000	9633.2\\
16000	10096\\
18000	9657.8\\
20000	9759.1\\
22000	9899.1\\
24000	9673.4\\
26000	9606\\
28000	9862.7\\
32000	9635.1\\
34000	9608.2\\
36000	9806.9\\
38000	9716.6\\
40000	9662.4\\
42000	9563.6\\
44000	9702.1\\
46000	9679.6\\
48000	9512.2\\
50000	9780.1\\
52000	9958.3\\
54000	9480.4\\
56000	9910.5\\
58000	9476.4\\
60000	9650.9\\
62000	9573.8\\
64000	9610\\
66000	9700.3\\
68000	9749.1\\
72000	9550.6\\
74000	9563.8\\
76000	9722.3\\
78000	9614.9\\
80000	9604\\
82000	9716.8\\
84000	9774.7\\
86000	9733.1\\
88000	9550\\
90000	9770.2\\
92000	9603.5\\
94000	9802.4\\
96000	9909.2\\
98000	9568.5\\
1e+05	9588.6\\
};
\addplot [color=mycolor22, forget plot]
  table[row sep=crcr]{%
0	4395.3\\
2000	23968\\
4000	26223\\
6000	25431\\
8000	25447\\
10000	25870\\
12000	25265\\
14000	24433\\
16000	24606\\
18000	25284\\
20000	25649\\
22000	24934\\
24000	25240\\
26000	25063\\
28000	24796\\
30000	25598\\
32000	25944\\
34000	25908\\
36000	24957\\
38000	25321\\
40000	25334\\
42000	25078\\
44000	25401\\
46000	25182\\
48000	26574\\
50000	23958\\
52000	25481\\
54000	25597\\
56000	25333\\
58000	24809\\
60000	26064\\
62000	25136\\
64000	25407\\
66000	24802\\
68000	25101\\
70000	26063\\
72000	25779\\
74000	25420\\
76000	24503\\
78000	25989\\
80000	25914\\
82000	25603\\
84000	25417\\
86000	25005\\
88000	25602\\
90000	24806\\
92000	25501\\
94000	25359\\
96000	24517\\
98000	25677\\
1e+05	25444\\
};
\addplot [color=mycolor23, forget plot]
  table[row sep=crcr]{%
0	4395.3\\
2000	11830\\
4000	12431\\
6000	12439\\
8000	12147\\
10000	12123\\
12000	12180\\
16000	12168\\
18000	12378\\
20000	12323\\
22000	12547\\
24000	12593\\
26000	12385\\
28000	12480\\
30000	12341\\
32000	12309\\
34000	12356\\
36000	12468\\
40000	12261\\
42000	12341\\
44000	12306\\
46000	12385\\
48000	11986\\
50000	12503\\
52000	12401\\
54000	12095\\
56000	12418\\
58000	12238\\
60000	12107\\
62000	12331\\
64000	12204\\
68000	12518\\
70000	12064\\
72000	12441\\
74000	12163\\
76000	12117\\
78000	12254\\
80000	12290\\
82000	12588\\
84000	12540\\
86000	12370\\
88000	12433\\
92000	12212\\
98000	12402\\
1e+05	12372\\
};
\addplot [color=mycolor24, forget plot]
  table[row sep=crcr]{%
0	4395.3\\
2000	5455.6\\
4000	4848.6\\
6000	4722.4\\
10000	4807\\
12000	4649.9\\
14000	4745.1\\
18000	4797.6\\
20000	4879.7\\
22000	4772.1\\
24000	4785.2\\
26000	4940.2\\
28000	4752\\
30000	4694.4\\
32000	4865.6\\
34000	4800.1\\
36000	4824.2\\
38000	4662.8\\
40000	4749.1\\
42000	5000.6\\
44000	4859.2\\
46000	4657.4\\
48000	4753.8\\
50000	4788.6\\
52000	4936.9\\
54000	4800.1\\
56000	4711\\
58000	4599\\
60000	4854.6\\
64000	4935.9\\
66000	4717.9\\
68000	4894\\
70000	4752.3\\
72000	4896\\
74000	4955.2\\
76000	4897.9\\
78000	4894.1\\
80000	4678\\
82000	4693\\
84000	4683.8\\
86000	4969\\
88000	4842.6\\
90000	4910.8\\
92000	4607.1\\
94000	4785.4\\
96000	4713.1\\
98000	4711\\
1e+05	4790.8\\
};
\addplot [color=mycolor25, forget plot]
  table[row sep=crcr]{%
0	4395.3\\
2000	11037\\
4000	12261\\
6000	12600\\
8000	12750\\
10000	12810\\
12000	12921\\
14000	12176\\
18000	12452\\
20000	12309\\
22000	12251\\
24000	12542\\
26000	13015\\
28000	12277\\
30000	12582\\
32000	12825\\
34000	12535\\
36000	12844\\
38000	12573\\
40000	12490\\
42000	12046\\
44000	12500\\
46000	12626\\
48000	12791\\
50000	12574\\
52000	12488\\
54000	12036\\
56000	12211\\
58000	12515\\
60000	12310\\
62000	12420\\
64000	12489\\
66000	12003\\
68000	12831\\
70000	12503\\
72000	12345\\
74000	11971\\
76000	12054\\
78000	12195\\
80000	12664\\
82000	12276\\
84000	12649\\
86000	12396\\
90000	12439\\
92000	12542\\
94000	12375\\
96000	12365\\
98000	12402\\
1e+05	12303\\
};
\addplot [color=mycolor26, forget plot]
  table[row sep=crcr]{%
0	4395.3\\
2000	7323.2\\
4000	8219.1\\
6000	8077.4\\
8000	8271.1\\
10000	8367.4\\
12000	8110.4\\
14000	8393.8\\
16000	8436.4\\
18000	8201.9\\
20000	8362.9\\
22000	8243.9\\
24000	8199.5\\
26000	8414.4\\
28000	7813.5\\
30000	8199\\
32000	8212.1\\
34000	8164.5\\
36000	8147.7\\
38000	8406.3\\
40000	8263.4\\
42000	8319.3\\
44000	8338.7\\
46000	8123.6\\
48000	8657.7\\
50000	8300.9\\
52000	8152.3\\
56000	8106.1\\
58000	8167.7\\
60000	8418.4\\
62000	8231\\
64000	8073\\
66000	7958.7\\
68000	8179.6\\
70000	8150.6\\
72000	7992.5\\
74000	8084.7\\
76000	8297.5\\
78000	8087.3\\
80000	8195.6\\
82000	8172.3\\
84000	8500.5\\
86000	8520.9\\
88000	8296.8\\
90000	7981.8\\
94000	8450.3\\
96000	8169.7\\
98000	8716.7\\
1e+05	8280.1\\
};
\addplot [color=mycolor27, forget plot]
  table[row sep=crcr]{%
0	4395.3\\
2000	9980.2\\
4000	9699.8\\
6000	9357.8\\
8000	9780.4\\
10000	10090\\
12000	9686.4\\
14000	9785.7\\
16000	9449.1\\
18000	9518.5\\
20000	9718\\
22000	9764.8\\
24000	9921.1\\
26000	9891.3\\
28000	9388.8\\
30000	9511.2\\
32000	9589.4\\
34000	10065\\
36000	9740.8\\
38000	9854.7\\
40000	9627.9\\
42000	9732.7\\
44000	9929\\
46000	9835.3\\
48000	9645.2\\
50000	9616.8\\
52000	9657.1\\
54000	10131\\
56000	9669.2\\
58000	9614.8\\
60000	9814\\
62000	9775.1\\
64000	9631.7\\
66000	10206\\
70000	9867.7\\
72000	9782.1\\
74000	9664.2\\
76000	9918.5\\
80000	9834.3\\
82000	9862.4\\
84000	9608.9\\
86000	10138\\
88000	9894.8\\
90000	9900.1\\
92000	10113\\
94000	9847.8\\
96000	9679\\
98000	9732.3\\
1e+05	9726.3\\
};
\addplot [color=mycolor28, forget plot]
  table[row sep=crcr]{%
0	4395.3\\
2000	5997\\
4000	5599.2\\
6000	5672.7\\
8000	5512.1\\
10000	5490.7\\
12000	5340.3\\
14000	5389.1\\
16000	5415.5\\
18000	5578.7\\
20000	5435.6\\
22000	5549\\
24000	5361.1\\
26000	5682.6\\
28000	5425.6\\
30000	5523.6\\
32000	5430.1\\
34000	5509.2\\
36000	5450.4\\
38000	5431.9\\
40000	5436.4\\
42000	5662.7\\
44000	5550.4\\
46000	5341.1\\
48000	5490.6\\
50000	5390.6\\
52000	5458.9\\
56000	5376.4\\
58000	5492.4\\
60000	5551.7\\
64000	5403\\
66000	5475.7\\
68000	5608.3\\
70000	5593.2\\
72000	5298.8\\
76000	5763.9\\
78000	5288.4\\
80000	5706.9\\
82000	5478\\
84000	5232\\
86000	5473\\
88000	5556.1\\
90000	5554.2\\
92000	5511.8\\
94000	5750\\
96000	5430.8\\
98000	5335.1\\
1e+05	5598.5\\
};
\addplot [color=mycolor29, forget plot]
  table[row sep=crcr]{%
0	4395.3\\
2000	11052\\
4000	11811\\
6000	11698\\
8000	11844\\
10000	11260\\
12000	11556\\
14000	11681\\
16000	11410\\
18000	12097\\
20000	11799\\
22000	11636\\
24000	11975\\
26000	11505\\
28000	11614\\
30000	12095\\
32000	11688\\
34000	11509\\
36000	11609\\
38000	11464\\
40000	11872\\
42000	11630\\
44000	11484\\
46000	11759\\
48000	11794\\
50000	11724\\
52000	11990\\
54000	11795\\
56000	11333\\
58000	12220\\
60000	11108\\
64000	11896\\
66000	11495\\
68000	11403\\
70000	11675\\
72000	11599\\
74000	11741\\
76000	11678\\
78000	11768\\
80000	11434\\
82000	11956\\
84000	11588\\
86000	11888\\
88000	11971\\
90000	11896\\
92000	11901\\
94000	11481\\
98000	11535\\
1e+05	11824\\
};
\addplot [color=mycolor30, forget plot]
  table[row sep=crcr]{%
0	4395.3\\
2000	7037.5\\
4000	7740.5\\
6000	7663.5\\
8000	7902.7\\
10000	7841.3\\
12000	7973.5\\
14000	7898.6\\
16000	8029.9\\
18000	7765.3\\
20000	7877.5\\
22000	7616.9\\
24000	7963.8\\
26000	7791.6\\
28000	7952.8\\
30000	7841\\
32000	7842.9\\
34000	7894.9\\
36000	7825.6\\
38000	7903.9\\
40000	7877.2\\
42000	7684.6\\
44000	7828.4\\
46000	7875.3\\
48000	7789.7\\
50000	7778\\
52000	7862.8\\
54000	7979.6\\
56000	7606.1\\
58000	7864.5\\
60000	7776.1\\
62000	7993.6\\
64000	7979.6\\
66000	8023.9\\
68000	7916.8\\
70000	7863.7\\
72000	7737.2\\
74000	7931.3\\
76000	7784.2\\
78000	7873.7\\
80000	7872.1\\
82000	8128.2\\
84000	7675.1\\
86000	7744.7\\
88000	8004.4\\
90000	7874.9\\
92000	7794.7\\
94000	8039.9\\
96000	7750.8\\
98000	7966.1\\
1e+05	7942.7\\
};
\addplot [color=mycolor31, forget plot]
  table[row sep=crcr]{%
0	4395.3\\
2000	6806.2\\
4000	6467.1\\
6000	6690.2\\
8000	6576.7\\
10000	6646.1\\
12000	6399.8\\
14000	6681.3\\
16000	6745.6\\
18000	6503.7\\
20000	6532\\
22000	6333.5\\
26000	6796.3\\
28000	6515.7\\
30000	6665.7\\
32000	6489.2\\
34000	6540.3\\
36000	6488.4\\
38000	6653.7\\
40000	6669.5\\
42000	6447.8\\
44000	6301.6\\
46000	6714.1\\
48000	6279.5\\
50000	6613.1\\
52000	6711.3\\
54000	6625.8\\
56000	6560.3\\
58000	6661.7\\
60000	6568.6\\
62000	6363\\
64000	6545.3\\
66000	6544.3\\
68000	6280\\
70000	6483.7\\
72000	6510.5\\
74000	6630.6\\
76000	6605.3\\
78000	6473.9\\
80000	6550.1\\
82000	6303.5\\
84000	6580.3\\
86000	6468.5\\
88000	6637.3\\
90000	6532.2\\
94000	6787.1\\
96000	6456.2\\
98000	6564.3\\
1e+05	6346.3\\
};
\addplot [color=mycolor32, forget plot]
  table[row sep=crcr]{%
0	4395.3\\
2000	10978\\
4000	11213\\
6000	11073\\
8000	10759\\
10000	10545\\
12000	10529\\
14000	10180\\
16000	10306\\
18000	10699\\
20000	10753\\
22000	10412\\
24000	10572\\
26000	10533\\
28000	10295\\
30000	10644\\
32000	10693\\
34000	10421\\
36000	10451\\
38000	10595\\
40000	10799\\
42000	10589\\
44000	10564\\
46000	10783\\
48000	10576\\
50000	10411\\
52000	10730\\
54000	10639\\
56000	10320\\
58000	10704\\
60000	10533\\
62000	10198\\
64000	10831\\
66000	10722\\
68000	10356\\
70000	10641\\
72000	10733\\
74000	10433\\
76000	10570\\
78000	10800\\
80000	10859\\
82000	10667\\
84000	10084\\
86000	10316\\
88000	10901\\
90000	11044\\
92000	10552\\
94000	10218\\
96000	10551\\
98000	10711\\
1e+05	10914\\
};
\addplot [color=mycolor33, forget plot]
  table[row sep=crcr]{%
0	4395.3\\
2000	7403.1\\
4000	8148.2\\
6000	7903.5\\
8000	8203.3\\
10000	7887.1\\
12000	8327.8\\
14000	8071.6\\
16000	8191.9\\
18000	8396\\
20000	8394.6\\
22000	8062.3\\
24000	8290.2\\
26000	8195.2\\
28000	8009.2\\
30000	7868.2\\
32000	8546.9\\
34000	8466.6\\
36000	8189.3\\
38000	8453.5\\
40000	8167.5\\
42000	7997.9\\
44000	8458.1\\
46000	8112.6\\
48000	8254.8\\
50000	8218.5\\
52000	8257.6\\
54000	8116.2\\
56000	8226.4\\
58000	8296.9\\
60000	8145.3\\
62000	8059.3\\
64000	8503.7\\
66000	8133.5\\
68000	8130.6\\
70000	8195.7\\
72000	8297.7\\
74000	8160.6\\
76000	7679.7\\
78000	8180.4\\
80000	8151.6\\
82000	8379.5\\
84000	8280.3\\
86000	7966.2\\
90000	8093.1\\
92000	8000.6\\
94000	8311.3\\
96000	8586.8\\
98000	8292.1\\
1e+05	8474\\
};
\addplot [color=mycolor34, forget plot]
  table[row sep=crcr]{%
0	4395.3\\
2000	4684.5\\
4000	3847.9\\
6000	3910.2\\
10000	4096.8\\
12000	4057.2\\
14000	4099.5\\
16000	4000.4\\
18000	3979.4\\
20000	3924.6\\
22000	3998.9\\
24000	3927.7\\
26000	4085.8\\
28000	3902.3\\
30000	3882.7\\
32000	4152.1\\
34000	3998.5\\
36000	4011.4\\
38000	4171.5\\
40000	3986.6\\
42000	4025.4\\
44000	4104.4\\
48000	4002.4\\
50000	3991.2\\
52000	4102.1\\
54000	3985.5\\
56000	4020.3\\
58000	3942.8\\
60000	4040\\
62000	3950.5\\
64000	4037\\
66000	3904.6\\
68000	3827.5\\
70000	4022.5\\
72000	3963.7\\
74000	4260.2\\
76000	4070.2\\
78000	3943.2\\
80000	4048.7\\
82000	4098.1\\
84000	4177.2\\
86000	3971.7\\
88000	4001.5\\
90000	4199.9\\
92000	4028.7\\
94000	4150.8\\
96000	3962.3\\
98000	3974.3\\
1e+05	4160\\
};
\addplot [color=mycolor35]
  table[row sep=crcr]{%
0	4395.3\\
2000	9482.2\\
4000	9999.5\\
6000	10302\\
8000	10110\\
10000	10269\\
12000	10164\\
14000	10324\\
16000	10072\\
18000	10279\\
20000	10445\\
22000	10261\\
24000	10429\\
26000	10079\\
28000	9973\\
32000	10136\\
34000	10240\\
36000	10120\\
38000	10231\\
42000	10235\\
44000	10440\\
46000	10328\\
48000	10108\\
50000	10445\\
52000	10310\\
54000	10337\\
56000	10024\\
58000	9889.6\\
60000	10458\\
64000	10206\\
66000	10251\\
68000	10054\\
72000	10011\\
74000	10162\\
76000	10070\\
78000	10192\\
80000	10160\\
84000	10217\\
86000	10178\\
88000	9988.6\\
90000	10065\\
92000	10378\\
94000	9982.4\\
96000	10281\\
98000	10250\\
1e+05	10148\\
};
\addlegendentry{Classe k=83}

\addplot [color=mycolor36, forget plot]
  table[row sep=crcr]{%
0	4395.3\\
2000	4794.2\\
4000	4321.3\\
6000	4388.4\\
8000	4288.4\\
10000	4473\\
12000	4354.7\\
14000	4376.6\\
16000	4295.7\\
20000	4444.5\\
22000	4418.3\\
24000	4281.1\\
28000	4174.7\\
30000	4342.4\\
32000	4222.8\\
34000	4445.2\\
36000	4179.8\\
38000	4149.5\\
40000	4172.5\\
42000	4388.7\\
44000	4243.5\\
46000	4332.7\\
48000	4346.7\\
50000	4432.2\\
52000	4431.8\\
54000	4203.3\\
56000	4324.9\\
58000	4368.7\\
60000	4439.9\\
62000	4124.3\\
64000	4172.9\\
66000	4388.6\\
68000	4366.8\\
70000	4372.7\\
72000	4130.6\\
74000	4362\\
76000	4414.2\\
78000	4280.9\\
80000	4313.3\\
82000	4289.8\\
84000	4356.2\\
86000	4288.5\\
88000	4177.9\\
94000	4305.8\\
96000	4193.6\\
1e+05	4207.5\\
};
\addplot [color=mycolor37, forget plot]
  table[row sep=crcr]{%
0	4395.3\\
2000	3874.2\\
4000	3553.8\\
6000	3494.7\\
8000	3599.5\\
10000	3638.9\\
12000	3601.1\\
14000	3678.2\\
16000	3601.1\\
20000	3610.9\\
22000	3592.3\\
24000	3608.8\\
26000	3606.2\\
28000	3637.3\\
32000	3632.4\\
36000	3682.4\\
38000	3543.4\\
40000	3531.4\\
42000	3580.6\\
44000	3676.9\\
46000	3631.1\\
48000	3560.8\\
50000	3565.1\\
52000	3697.9\\
54000	3640.1\\
56000	3663.5\\
58000	3611.4\\
60000	3615.1\\
62000	3659.2\\
64000	3627.7\\
66000	3655.4\\
68000	3695\\
70000	3710\\
72000	3569.3\\
74000	3666.3\\
76000	3623.1\\
78000	3613.9\\
80000	3577.5\\
82000	3602.6\\
84000	3590.7\\
86000	3642.4\\
88000	3552\\
92000	3662.8\\
94000	3688.8\\
98000	3617.9\\
1e+05	3610.9\\
};
\addplot [color=mycolor38, forget plot]
  table[row sep=crcr]{%
0	4395.3\\
2000	3246.8\\
4000	3062.6\\
6000	2867.6\\
8000	3022\\
10000	3002.9\\
12000	2968\\
14000	2871.8\\
16000	3003.5\\
20000	2933.5\\
22000	3047.9\\
24000	2959.8\\
26000	2967.6\\
28000	3153.5\\
30000	2883.5\\
32000	3020.7\\
34000	2935.9\\
36000	2958\\
38000	3051.2\\
40000	3025.9\\
42000	2942.8\\
44000	3095.5\\
46000	3005.3\\
50000	3010.5\\
52000	2926.3\\
54000	2933\\
56000	2857.7\\
58000	3071.2\\
60000	2948.9\\
62000	2962.5\\
64000	2943.4\\
66000	2936.4\\
68000	2965.8\\
70000	3070.8\\
72000	3062.1\\
74000	2903.8\\
76000	3044.6\\
78000	3022.3\\
80000	2978\\
82000	3024.9\\
84000	2977.3\\
86000	3044.6\\
88000	2966.3\\
90000	2995.8\\
92000	2941.6\\
94000	2876.1\\
96000	3115.7\\
98000	3076.7\\
1e+05	2903\\
};
\addplot [color=mycolor39, forget plot]
  table[row sep=crcr]{%
0	4395.3\\
2000	12034\\
4000	12593\\
6000	11733\\
8000	11563\\
10000	11844\\
12000	11853\\
14000	11976\\
16000	11970\\
18000	11882\\
20000	11211\\
22000	11732\\
24000	11443\\
26000	11571\\
28000	11904\\
30000	12078\\
32000	12137\\
34000	12329\\
38000	11804\\
40000	11787\\
42000	11456\\
44000	11983\\
46000	11694\\
48000	11683\\
50000	12108\\
52000	11919\\
54000	11321\\
56000	12125\\
58000	12423\\
60000	11838\\
62000	11846\\
64000	11915\\
66000	12042\\
68000	11539\\
70000	11482\\
72000	12005\\
74000	11797\\
76000	12041\\
78000	11686\\
80000	11739\\
82000	11695\\
84000	11822\\
86000	12157\\
88000	11831\\
92000	11858\\
94000	11781\\
96000	11426\\
98000	11202\\
1e+05	11508\\
};
\addplot [color=mycolor40, forget plot]
  table[row sep=crcr]{%
0	4395.3\\
2000	5149\\
4000	5539.8\\
6000	5717.7\\
8000	5773.1\\
10000	5701.1\\
12000	5768.4\\
14000	5781.5\\
16000	5728.8\\
18000	5749.6\\
20000	5734.8\\
22000	5787.2\\
24000	5811.7\\
26000	5671.6\\
28000	5822.9\\
30000	5622.9\\
34000	5837.8\\
36000	5648.9\\
38000	5562.2\\
40000	5730.9\\
42000	5657.8\\
44000	5653.3\\
46000	5717.5\\
48000	5660.8\\
50000	5734.9\\
52000	5755.5\\
54000	5680.4\\
56000	5717.9\\
58000	5738.8\\
60000	5467\\
62000	5567.1\\
64000	5733.2\\
66000	5741.1\\
68000	5890.1\\
70000	5793.5\\
72000	5817.9\\
76000	5718.1\\
78000	5609.2\\
80000	5690.6\\
82000	5872.9\\
84000	5705.2\\
86000	5811.9\\
88000	5792.7\\
90000	5738.7\\
92000	5829.2\\
94000	5690\\
96000	5744.7\\
98000	5759.9\\
1e+05	5739.9\\
};
\addplot [color=mycolor41, forget plot]
  table[row sep=crcr]{%
0	4395.3\\
2000	4728.5\\
4000	5312.9\\
6000	5265.3\\
8000	5380\\
10000	5453.9\\
12000	5358.5\\
14000	5287.5\\
16000	5445.7\\
18000	5383.6\\
20000	5247.4\\
22000	5399\\
24000	5262.6\\
26000	5348.9\\
28000	5415.3\\
30000	5346.5\\
32000	5428.4\\
34000	5426.2\\
36000	5382.9\\
38000	5146.5\\
40000	5482\\
42000	5238.9\\
44000	5461.9\\
46000	5435.4\\
48000	5537.8\\
50000	5251.4\\
52000	5530.3\\
54000	5535.8\\
56000	5408.1\\
58000	5463.3\\
60000	5256.9\\
62000	5505.9\\
64000	5389\\
68000	5421.5\\
70000	5354\\
72000	5258.4\\
74000	5303.8\\
76000	5304.9\\
78000	5365.7\\
80000	5324.3\\
82000	5318\\
84000	5584.8\\
86000	5274.5\\
88000	5299.1\\
90000	5438.2\\
92000	5540.6\\
94000	5167.1\\
96000	5397.5\\
98000	5603\\
1e+05	5416.7\\
};
\addplot [color=mycolor42, forget plot]
  table[row sep=crcr]{%
0	4395.3\\
2000	15774\\
4000	19877\\
6000	21919\\
8000	20773\\
10000	20623\\
12000	20278\\
14000	21019\\
16000	21150\\
18000	20983\\
20000	20711\\
22000	20755\\
24000	20708\\
26000	21351\\
28000	20451\\
30000	20453\\
32000	20215\\
36000	20704\\
38000	20245\\
40000	20950\\
42000	20850\\
44000	20851\\
46000	21239\\
48000	20185\\
50000	20500\\
52000	21194\\
56000	20611\\
58000	20436\\
60000	20606\\
62000	21082\\
64000	21073\\
68000	20499\\
70000	20636\\
72000	20894\\
74000	21091\\
76000	21593\\
78000	20882\\
80000	20752\\
82000	21330\\
84000	20370\\
86000	21019\\
88000	20494\\
90000	20233\\
92000	20848\\
94000	20554\\
96000	21917\\
98000	20337\\
1e+05	20712\\
};
\addplot [color=mycolor43, forget plot]
  table[row sep=crcr]{%
0	4395.3\\
2000	5433.4\\
4000	5084.7\\
6000	4998.3\\
8000	5059.8\\
10000	4855.1\\
12000	5013.3\\
14000	5040.3\\
16000	5025.6\\
18000	5066.6\\
20000	5043.9\\
22000	4798.7\\
24000	4995.2\\
26000	4930.3\\
28000	4888.4\\
30000	4971.3\\
32000	4974.5\\
34000	5021.2\\
36000	5024.5\\
38000	4896.8\\
40000	4907.4\\
42000	4968.1\\
44000	4974.4\\
46000	4947.4\\
48000	4875.2\\
50000	4963.1\\
52000	4892.4\\
54000	5018.5\\
56000	4981.3\\
58000	5012.2\\
60000	4814.8\\
62000	4952.6\\
64000	4939.1\\
66000	4991.1\\
68000	4896.3\\
70000	4876.4\\
72000	5050.8\\
74000	4870.6\\
76000	4956\\
78000	4940.9\\
80000	5035.8\\
82000	4882.5\\
84000	4869.1\\
86000	5003.1\\
88000	4892.9\\
90000	4893.7\\
92000	5075.9\\
96000	4981.2\\
98000	4917\\
1e+05	5016.3\\
};
\addplot [color=mycolor44, forget plot]
  table[row sep=crcr]{%
0	4395.3\\
2000	4032.9\\
4000	3512.3\\
6000	3447.8\\
8000	3411\\
10000	3412.5\\
14000	3373.5\\
16000	3272.2\\
18000	3338.7\\
20000	3345.4\\
22000	3294.5\\
24000	3320.5\\
26000	3388.1\\
28000	3342.1\\
30000	3456.6\\
32000	3341\\
34000	3400.1\\
36000	3315.9\\
38000	3349.7\\
40000	3371.2\\
42000	3437.6\\
44000	3349.1\\
46000	3375.3\\
48000	3554.5\\
50000	3466.4\\
52000	3317.4\\
54000	3416.1\\
56000	3404.5\\
58000	3316.3\\
60000	3390.4\\
62000	3327.2\\
64000	3386.6\\
66000	3296\\
68000	3417\\
70000	3339.4\\
72000	3406.1\\
74000	3458.2\\
76000	3357.8\\
78000	3339.4\\
80000	3369.6\\
82000	3303.3\\
84000	3318.2\\
88000	3315.1\\
90000	3431.4\\
92000	3382.1\\
94000	3468.8\\
96000	3450.3\\
98000	3441.6\\
1e+05	3490.4\\
};
\addplot [color=mycolor45, forget plot]
  table[row sep=crcr]{%
0	4395.3\\
2000	6011.2\\
4000	6604.4\\
6000	6823.2\\
8000	6923.5\\
10000	6848.2\\
12000	7039.1\\
14000	6876.7\\
16000	6932.9\\
18000	6824.4\\
20000	7064.8\\
22000	6782.2\\
24000	6962.9\\
26000	6843.2\\
28000	6904.9\\
30000	6924.5\\
32000	7060.1\\
34000	7049.4\\
38000	6970.3\\
40000	7020.7\\
42000	6823.3\\
44000	7059.7\\
46000	6840.6\\
48000	7023.3\\
50000	6806.1\\
52000	6854.8\\
54000	6734.8\\
56000	6946.6\\
58000	6892.4\\
60000	6737.3\\
62000	6969.2\\
64000	6759.4\\
66000	7044\\
68000	6811.8\\
70000	6882\\
72000	6904.9\\
74000	6725.7\\
76000	6851.2\\
78000	6796.3\\
80000	6864.8\\
82000	6990.2\\
84000	6851.5\\
88000	6780.5\\
92000	7074.6\\
94000	6961.7\\
96000	6940.8\\
98000	6947.1\\
1e+05	6904.3\\
};
\addplot [color=mycolor46, forget plot]
  table[row sep=crcr]{%
0	4395.3\\
2000	3859.1\\
4000	3933.6\\
6000	4124.9\\
8000	4067.2\\
10000	3992.1\\
12000	4013\\
14000	3993.5\\
16000	4079.1\\
20000	4067\\
22000	4076.9\\
24000	4230.4\\
26000	4040.5\\
28000	4032.1\\
30000	4112.8\\
32000	4162.7\\
34000	3972\\
36000	4070.4\\
38000	4135.7\\
40000	4015.4\\
42000	4012.4\\
44000	4056.4\\
46000	4047.3\\
48000	4096.4\\
50000	4071.5\\
52000	4219.3\\
54000	4205.6\\
56000	4017\\
58000	3937.9\\
60000	4017.9\\
62000	3964.9\\
64000	4067.7\\
66000	3980.8\\
68000	4093.2\\
70000	4041\\
72000	4021\\
74000	4070.6\\
76000	3997.8\\
78000	4180.1\\
80000	4117.9\\
82000	4131.1\\
84000	4052.5\\
86000	4115.5\\
88000	4097.4\\
90000	4115.3\\
92000	4188.3\\
94000	4043.4\\
96000	3995.3\\
98000	3982.4\\
1e+05	4105.7\\
};
\addplot [color=mycolor47, forget plot]
  table[row sep=crcr]{%
0	4395.3\\
2000	5456.5\\
4000	5342.6\\
6000	5339.2\\
8000	5404.3\\
10000	5263.7\\
12000	5297.5\\
14000	5372.5\\
16000	5178.4\\
18000	5211\\
20000	5325.8\\
22000	5489.2\\
24000	5230.2\\
26000	5311.2\\
28000	5135.1\\
30000	5207.9\\
32000	5322.7\\
34000	5204.4\\
36000	5356.2\\
38000	5341.6\\
40000	5214.8\\
42000	5429.4\\
44000	5273.6\\
46000	5302.7\\
48000	5372.6\\
50000	5340.1\\
52000	5400.7\\
54000	5343\\
56000	5434.9\\
58000	5302.3\\
60000	5406.7\\
62000	5327.3\\
64000	5289.3\\
66000	5376.6\\
68000	5320.6\\
70000	5336.2\\
72000	5287.1\\
74000	5475.7\\
76000	5253.8\\
78000	5466.1\\
80000	5420.1\\
82000	5490.2\\
84000	5061.4\\
86000	5289\\
88000	5343.6\\
90000	5472\\
92000	5226.2\\
94000	5249.9\\
96000	5175.6\\
98000	5313.7\\
1e+05	5234.5\\
};
\addplot [color=mycolor48, forget plot]
  table[row sep=crcr]{%
0	4395.3\\
2000	3281.2\\
4000	2922\\
8000	2956.3\\
10000	2897.7\\
12000	2849.8\\
14000	2901.8\\
16000	2900\\
18000	2927.8\\
20000	2967.3\\
22000	2900.7\\
24000	2871.6\\
26000	2959.8\\
28000	2867.7\\
30000	2858.5\\
32000	2903\\
34000	2886.5\\
36000	2958.5\\
38000	2954.2\\
40000	2902.4\\
42000	2918.8\\
46000	2909.5\\
48000	2915.2\\
50000	2956.4\\
52000	2874.8\\
54000	2887.4\\
58000	2959.5\\
60000	2861.8\\
62000	2880.5\\
64000	2910.6\\
66000	2949.9\\
70000	2930.8\\
72000	2871\\
74000	2958.2\\
76000	2915.1\\
78000	2853.5\\
80000	2922.9\\
84000	2886.8\\
86000	2935.6\\
88000	2953.5\\
90000	2856\\
94000	2946.4\\
96000	2994.9\\
98000	2905.4\\
1e+05	2874.5\\
};
\addplot [color=mycolor49, forget plot]
  table[row sep=crcr]{%
0	4395.3\\
2000	4708.4\\
4000	4266.9\\
6000	4340.8\\
8000	4661.8\\
10000	4324\\
12000	4368.5\\
14000	4455.6\\
16000	4277.2\\
18000	4268\\
20000	4372\\
22000	4242\\
24000	4483\\
26000	4333.8\\
30000	4350.7\\
32000	4455.2\\
34000	4339.1\\
36000	4485\\
38000	4281.2\\
40000	4186.6\\
42000	4272.1\\
44000	4390.8\\
46000	4455.2\\
48000	4493.6\\
50000	4372.4\\
52000	4209.9\\
54000	4359.6\\
56000	4280\\
58000	4427.3\\
60000	4283\\
62000	4506\\
66000	4405.1\\
70000	4435\\
72000	4406.4\\
74000	4441\\
76000	4279.3\\
78000	4377.9\\
80000	4357.2\\
82000	4160.6\\
84000	4304.2\\
86000	4519.3\\
88000	4363.8\\
90000	4367.1\\
92000	4496.8\\
96000	4206.3\\
98000	4386.2\\
1e+05	4349.1\\
};
\addplot [color=mycolor50, forget plot]
  table[row sep=crcr]{%
0	4395.3\\
2000	2533.8\\
4000	2393.4\\
6000	2386\\
8000	2488.9\\
10000	2434.5\\
12000	2525.4\\
14000	2413.7\\
16000	2420.5\\
18000	2407.3\\
20000	2481.6\\
22000	2433.7\\
24000	2468\\
26000	2491.1\\
28000	2389.1\\
30000	2439.2\\
32000	2478.8\\
34000	2501.5\\
36000	2471.1\\
38000	2373.7\\
40000	2438.9\\
42000	2546.6\\
44000	2457.9\\
46000	2439.2\\
48000	2365.4\\
50000	2468.6\\
52000	2501.2\\
54000	2478.5\\
56000	2480.6\\
58000	2491.7\\
60000	2390.7\\
62000	2483.2\\
64000	2483.9\\
66000	2441.3\\
68000	2471.6\\
70000	2470.9\\
72000	2501.8\\
74000	2430.2\\
76000	2461.3\\
78000	2423\\
80000	2441.9\\
82000	2545.2\\
84000	2522.9\\
86000	2469.5\\
88000	2407.6\\
90000	2397.8\\
92000	2484.1\\
94000	2451.6\\
96000	2427.5\\
98000	2484.2\\
1e+05	2474.7\\
};
\addplot [color=mycolor51, forget plot]
  table[row sep=crcr]{%
0	4395.3\\
2000	5865.6\\
4000	5890.1\\
6000	5814.7\\
8000	5958.7\\
10000	5973.9\\
12000	6138.9\\
14000	5798.1\\
16000	5709.1\\
18000	5873.1\\
20000	5871\\
22000	5817.6\\
24000	6033.6\\
28000	6168.9\\
30000	6102.5\\
34000	5912.1\\
36000	5959.8\\
38000	6264.3\\
40000	5862.5\\
42000	5856.5\\
46000	6041.4\\
48000	6018.1\\
52000	5890.4\\
54000	6055.4\\
58000	5876.6\\
60000	5954.2\\
62000	5966.4\\
64000	6149.1\\
66000	6017.6\\
68000	5830.9\\
70000	6139.7\\
72000	5691.6\\
74000	6090.1\\
76000	6088.7\\
78000	5991.8\\
80000	5745.9\\
84000	5975.8\\
86000	5780.1\\
88000	5895.7\\
90000	5967.6\\
92000	6112.4\\
94000	5935.6\\
96000	6077.5\\
98000	6162.7\\
1e+05	5574.4\\
};
\addplot [color=mycolor52, forget plot]
  table[row sep=crcr]{%
0	4395.3\\
4000	3663.3\\
6000	3580.3\\
8000	3401.7\\
10000	3386.6\\
12000	3475\\
14000	3404.4\\
16000	3447.1\\
18000	3424.7\\
20000	3565.5\\
22000	3429.8\\
24000	3403.9\\
26000	3564\\
30000	3432.7\\
34000	3653.6\\
36000	3537\\
38000	3587.1\\
40000	3409.1\\
42000	3540.6\\
44000	3531\\
46000	3536.2\\
48000	3621.4\\
50000	3456.3\\
52000	3518.5\\
54000	3537.2\\
56000	3427.2\\
58000	3496.3\\
60000	3586.9\\
62000	3416.7\\
64000	3549.9\\
66000	3559.2\\
70000	3482.7\\
72000	3462.8\\
74000	3589.7\\
76000	3593.7\\
78000	3656.2\\
80000	3514.8\\
82000	3399.2\\
84000	3568.5\\
86000	3403\\
88000	3604.6\\
90000	3536\\
92000	3517.8\\
94000	3445.4\\
96000	3626.9\\
98000	3701.7\\
1e+05	3492.7\\
};
\addplot [color=mycolor53, forget plot]
  table[row sep=crcr]{%
0	4395.3\\
2000	6394.4\\
4000	6354.5\\
6000	5950.3\\
8000	5651\\
10000	5583.8\\
12000	5466\\
14000	5571.7\\
16000	5737.4\\
18000	5523.4\\
20000	6115.8\\
22000	5813.4\\
24000	5780.4\\
26000	5790.1\\
28000	5724.9\\
30000	5590.1\\
32000	5921.6\\
34000	5741.4\\
36000	5951.7\\
38000	5655.7\\
40000	5645.6\\
42000	5785.6\\
44000	5685.2\\
46000	5739.3\\
48000	5880.9\\
50000	5602.2\\
52000	5843.3\\
54000	5640.5\\
56000	5919.7\\
58000	5922.3\\
60000	5728.8\\
62000	5706.7\\
64000	5712.7\\
66000	5555\\
68000	5712.3\\
70000	5512.2\\
72000	5824.9\\
74000	5712.6\\
76000	5799.1\\
78000	5638.1\\
80000	5823.4\\
82000	5781.4\\
84000	5804.2\\
86000	5693.9\\
88000	5663.5\\
90000	5764.8\\
92000	5733.2\\
94000	5666.2\\
96000	5576.6\\
98000	5699.5\\
1e+05	5801.5\\
};
\addplot [color=mycolor54, forget plot]
  table[row sep=crcr]{%
0	4395.3\\
2000	4520.9\\
4000	4455.7\\
6000	4732.1\\
8000	4585.7\\
10000	4473.7\\
12000	4529.2\\
14000	4444.6\\
16000	4642.4\\
18000	4677\\
20000	4686.8\\
22000	4617\\
24000	4512.8\\
26000	4538.8\\
28000	4620.9\\
30000	4627.2\\
32000	4800.1\\
34000	4510.2\\
36000	4458.7\\
38000	4685.7\\
40000	4374\\
42000	4621.8\\
44000	4627.2\\
46000	4727\\
48000	4651.2\\
52000	4769.3\\
54000	4404.4\\
56000	4648.9\\
58000	4600\\
60000	4655.3\\
62000	4657.1\\
64000	4694.6\\
66000	4529.3\\
70000	4628.5\\
72000	4634.9\\
74000	4605.2\\
76000	4718.3\\
78000	4522.4\\
80000	4634.5\\
82000	4641.2\\
84000	4693.6\\
86000	4478.6\\
88000	4554\\
90000	4708.8\\
92000	4622\\
94000	4589.8\\
96000	4727\\
98000	4743\\
1e+05	4464\\
};
\addplot [color=mycolor55, forget plot]
  table[row sep=crcr]{%
0	4395.3\\
2000	5383.3\\
4000	5736.2\\
8000	5846.3\\
10000	5653.1\\
12000	5910.3\\
14000	5823.2\\
16000	5711\\
18000	5761.8\\
20000	5794.9\\
22000	5637.6\\
24000	5975.9\\
26000	5593.1\\
28000	5789.7\\
32000	5847.5\\
36000	5638\\
40000	5729.6\\
42000	5858.3\\
46000	5525.1\\
48000	5751.6\\
50000	5715.6\\
52000	5781.8\\
54000	5732.5\\
56000	5771.4\\
58000	5742\\
60000	5875\\
62000	5750.2\\
64000	5683.5\\
66000	5924.7\\
68000	5839.9\\
70000	5724.1\\
72000	5770\\
74000	5693.4\\
76000	5737.6\\
78000	5732.3\\
80000	5938.1\\
82000	5699.9\\
86000	5661.6\\
88000	5760.3\\
90000	5917\\
92000	5570.5\\
94000	5687.4\\
96000	5658.7\\
98000	5740.5\\
1e+05	5868.7\\
};
\addplot [color=mycolor56, forget plot]
  table[row sep=crcr]{%
0	4395.3\\
2000	4273.7\\
4000	4041.1\\
6000	3995.9\\
8000	3985.3\\
10000	3945\\
12000	4006\\
14000	3946.9\\
16000	3980.9\\
18000	3988.1\\
20000	3907.6\\
22000	3962.5\\
24000	4031.3\\
26000	3966.8\\
28000	3867.5\\
30000	3918.2\\
32000	4041.7\\
34000	4152.6\\
36000	3966.6\\
38000	3992.5\\
40000	3988.3\\
42000	4046.7\\
44000	3967.6\\
46000	4006.7\\
48000	3969.5\\
50000	4066.5\\
52000	3999.6\\
54000	3962.1\\
56000	3943.6\\
58000	3824.5\\
60000	3970.8\\
62000	4051.4\\
64000	4055.5\\
66000	4003.4\\
68000	3918.4\\
70000	3957.7\\
72000	4047.9\\
74000	3929.7\\
76000	3955\\
78000	3994.4\\
80000	3922.2\\
82000	4033.2\\
84000	3941.1\\
86000	3926.7\\
88000	4001.1\\
90000	3869.8\\
92000	4020.7\\
94000	3929.7\\
96000	4005.5\\
98000	4037.5\\
1e+05	4026.5\\
};
\addplot [color=mycolor57, forget plot]
  table[row sep=crcr]{%
0	4395.3\\
2000	6459.4\\
4000	6957.6\\
6000	6753.7\\
8000	6802.9\\
10000	6502.3\\
12000	6598.8\\
14000	6919.1\\
16000	6892.8\\
18000	6818.4\\
20000	6694.3\\
22000	6767.8\\
24000	6710.3\\
26000	6751.7\\
28000	7002.6\\
30000	6830.2\\
32000	6877.9\\
34000	6884\\
38000	6676.5\\
40000	6752.7\\
42000	6724.9\\
44000	6536.3\\
46000	6663\\
48000	6613.8\\
50000	6626.2\\
52000	6499.5\\
54000	6656.6\\
56000	6595.5\\
58000	6753\\
60000	6745.3\\
62000	6637.6\\
64000	6683.1\\
66000	6681.5\\
68000	6606.7\\
70000	6790.6\\
72000	6720.2\\
74000	6744\\
76000	6531.6\\
78000	6529.6\\
80000	6715.3\\
82000	6711.5\\
84000	6865.3\\
86000	6636.3\\
88000	6745.6\\
90000	6835.5\\
92000	6699.5\\
94000	6723.8\\
96000	6845.6\\
98000	6830.6\\
1e+05	6775.2\\
};
\addplot [color=mycolor58, forget plot]
  table[row sep=crcr]{%
0	4395.3\\
2000	3856.5\\
4000	3541.6\\
6000	3541.7\\
8000	3569.1\\
10000	3800.7\\
12000	3586.1\\
14000	3541.7\\
16000	3612.2\\
20000	3554.8\\
22000	3648.9\\
24000	3647\\
26000	3580.4\\
28000	3542\\
30000	3692.3\\
32000	3703.6\\
34000	3622.6\\
36000	3615.5\\
38000	3647.9\\
40000	3646.7\\
42000	3497.3\\
44000	3565.8\\
46000	3496.3\\
48000	3592.2\\
50000	3537.7\\
52000	3553.7\\
54000	3582.3\\
58000	3677.7\\
60000	3636.2\\
62000	3496.6\\
64000	3554.3\\
66000	3572.4\\
68000	3458.4\\
70000	3574.7\\
72000	3535.7\\
74000	3572.9\\
76000	3598.7\\
78000	3581.3\\
80000	3513.8\\
82000	3704.4\\
84000	3451.8\\
86000	3500.5\\
88000	3618.1\\
90000	3513.5\\
92000	3571.1\\
94000	3577.2\\
96000	3676.1\\
98000	3526.7\\
1e+05	3527.5\\
};
\addplot [color=mycolor59, forget plot]
  table[row sep=crcr]{%
0	4395.3\\
2000	4638\\
4000	4958\\
6000	4836.5\\
8000	5106.7\\
10000	5072.9\\
12000	5023.3\\
14000	5057.3\\
16000	5116.1\\
18000	5158.6\\
20000	4987.4\\
22000	5076.9\\
26000	5018.3\\
28000	5046\\
30000	5000.3\\
32000	5057.9\\
34000	5053.9\\
36000	5111.7\\
38000	5083.8\\
40000	4996.9\\
42000	5076.9\\
44000	5120.8\\
46000	5061.5\\
48000	5173.9\\
50000	4945.1\\
52000	5041.4\\
54000	5197\\
56000	5055.2\\
58000	5284.4\\
60000	4956.8\\
62000	5189\\
64000	4923.4\\
66000	5054.7\\
68000	5075.3\\
70000	5036.6\\
72000	5050.1\\
74000	5024.1\\
76000	4969.1\\
78000	4953.3\\
80000	5170.8\\
82000	5007.9\\
84000	5182.8\\
86000	5012.3\\
88000	5064.3\\
90000	4867.3\\
92000	5023.1\\
94000	5071.5\\
96000	5187.2\\
98000	5048.1\\
1e+05	5046.3\\
};
\addplot [color=mycolor60, forget plot]
  table[row sep=crcr]{%
0	4395.3\\
2000	3790.8\\
4000	3637.3\\
6000	3627\\
8000	3469.8\\
10000	3573.7\\
14000	3535.9\\
16000	3550.5\\
18000	3508.4\\
20000	3519.7\\
22000	3548.7\\
24000	3515.3\\
26000	3523.5\\
28000	3463.6\\
30000	3522.1\\
32000	3533.3\\
34000	3623.8\\
36000	3520.9\\
38000	3607.3\\
40000	3596.6\\
42000	3556.8\\
44000	3489.5\\
46000	3512.8\\
48000	3509.8\\
50000	3569.2\\
54000	3458.5\\
56000	3526.5\\
58000	3523\\
60000	3611.3\\
64000	3559.4\\
66000	3466.7\\
68000	3441\\
70000	3534.3\\
72000	3505.6\\
74000	3600.9\\
76000	3539.2\\
78000	3573.7\\
80000	3589.6\\
82000	3427.2\\
84000	3537.7\\
88000	3564.8\\
90000	3591\\
92000	3602.9\\
94000	3632\\
96000	3574.7\\
98000	3561.5\\
1e+05	3563.1\\
};
\addplot [color=mycolor61, forget plot]
  table[row sep=crcr]{%
0	4395.3\\
2000	3261.9\\
4000	3076.1\\
6000	3119.5\\
8000	3085.4\\
10000	3131.6\\
12000	3058\\
14000	3108.9\\
16000	3150.7\\
18000	3110.6\\
20000	3088.2\\
22000	3025.2\\
24000	3116.8\\
26000	3123.7\\
28000	3115.1\\
30000	3153.6\\
32000	3122\\
34000	3128.5\\
36000	3190\\
38000	3070.1\\
40000	3085.1\\
42000	3114.3\\
44000	3094.1\\
46000	3128.8\\
48000	3096.2\\
50000	3095.9\\
52000	3114.5\\
54000	3143.9\\
56000	3037.6\\
58000	3034.1\\
60000	3129.6\\
62000	3101.2\\
64000	3163.5\\
66000	3079.1\\
70000	3082.2\\
72000	3120.1\\
74000	3063.5\\
76000	3089.2\\
80000	3111.7\\
84000	3115.4\\
86000	3075.7\\
88000	3165.3\\
90000	3070.2\\
92000	3149.4\\
94000	3097.6\\
96000	3139.6\\
98000	3118.7\\
1e+05	3156.1\\
};
\addplot [color=mycolor62, forget plot]
  table[row sep=crcr]{%
0	4395.3\\
2000	4743.2\\
4000	5054.4\\
6000	5003.4\\
8000	4915.7\\
10000	4878.2\\
12000	4925\\
14000	5111.8\\
16000	4996.3\\
18000	5028.5\\
20000	5045.6\\
22000	4963.4\\
24000	4908.4\\
26000	5042.3\\
28000	4936.1\\
30000	5024.8\\
32000	4975.9\\
34000	5014.1\\
38000	4973.1\\
40000	4953.6\\
42000	4829\\
44000	4907.9\\
46000	4921.9\\
48000	4974.6\\
50000	4978.9\\
52000	5013.7\\
54000	4910.5\\
56000	5029.7\\
58000	4928.1\\
60000	4945.2\\
62000	4990.7\\
64000	4990.3\\
66000	4962.2\\
68000	4890.7\\
70000	4998.8\\
72000	4924.6\\
76000	4990.6\\
78000	4963.5\\
80000	4841.5\\
82000	4900\\
86000	5107.6\\
88000	5010.6\\
90000	4945\\
92000	4978.2\\
94000	4969\\
96000	4813.3\\
98000	5006.8\\
1e+05	4995.1\\
};
\addplot [color=mycolor63, forget plot]
  table[row sep=crcr]{%
0	4395.3\\
2000	2729.5\\
4000	2405\\
6000	2394.4\\
10000	2433.8\\
12000	2447.4\\
18000	2370.4\\
20000	2416.1\\
22000	2416.3\\
24000	2439.4\\
26000	2395.4\\
28000	2360.3\\
30000	2393.3\\
32000	2374\\
34000	2373\\
36000	2440.6\\
38000	2434.2\\
40000	2362.4\\
42000	2361\\
44000	2393.9\\
46000	2414.5\\
48000	2388\\
50000	2424.6\\
52000	2389.8\\
54000	2391.6\\
56000	2352.7\\
58000	2369.9\\
60000	2432.9\\
62000	2421.5\\
64000	2361.8\\
66000	2358.7\\
68000	2419.8\\
70000	2391.3\\
72000	2417.4\\
74000	2393.4\\
76000	2437.2\\
78000	2394\\
80000	2388.4\\
82000	2395.9\\
84000	2433.7\\
86000	2385.7\\
88000	2388.6\\
90000	2377.7\\
92000	2410.3\\
94000	2418.2\\
98000	2368.6\\
1e+05	2407\\
};
\addplot [color=mycolor64, forget plot]
  table[row sep=crcr]{%
0	4395.3\\
2000	3066\\
4000	2997.4\\
8000	3044.5\\
12000	3044.3\\
14000	3011.7\\
16000	3064.2\\
18000	3037.4\\
22000	3023.5\\
24000	3061.3\\
28000	3037.2\\
32000	3046\\
34000	3003.9\\
36000	2993.7\\
38000	3051.7\\
40000	3036.3\\
42000	3003.4\\
44000	2934.9\\
48000	3002.5\\
50000	3074\\
52000	3074.2\\
56000	3006.3\\
58000	3032.4\\
64000	3014.8\\
66000	3059\\
68000	3025.1\\
70000	3046.1\\
72000	2999.5\\
74000	3015.4\\
76000	2983.5\\
78000	2995.7\\
80000	3032.3\\
82000	3107.9\\
84000	3108.4\\
88000	3025.1\\
90000	3027.6\\
92000	2978.9\\
94000	3020\\
96000	3071.4\\
98000	3105.4\\
1e+05	2928\\
};
\addplot [color=mycolor65, forget plot]
  table[row sep=crcr]{%
0	4395.3\\
2000	2756.5\\
4000	2502.8\\
6000	2488.9\\
8000	2540\\
10000	2567.8\\
12000	2575.2\\
14000	2496.2\\
16000	2471.6\\
18000	2490\\
20000	2569.8\\
22000	2536.8\\
30000	2537.1\\
32000	2521.6\\
34000	2533.1\\
36000	2498.5\\
38000	2487.4\\
40000	2524.2\\
44000	2516.9\\
46000	2558.2\\
48000	2511.6\\
50000	2546.1\\
52000	2528.1\\
54000	2563.3\\
56000	2458.9\\
60000	2557.2\\
64000	2560.2\\
66000	2563.5\\
68000	2543.7\\
72000	2550.5\\
74000	2483.9\\
76000	2498.3\\
78000	2544.6\\
80000	2488.3\\
82000	2539.9\\
84000	2561.1\\
86000	2530.7\\
88000	2546\\
90000	2499.4\\
92000	2514.2\\
94000	2519.9\\
98000	2497.1\\
1e+05	2507\\
};
\addplot [color=mycolor66, forget plot]
  table[row sep=crcr]{%
0	4395.3\\
2000	3430.6\\
4000	3417.8\\
6000	3427.3\\
8000	3449.7\\
10000	3451\\
12000	3401.9\\
14000	3366.1\\
16000	3411.9\\
20000	3428.2\\
22000	3407\\
24000	3413.7\\
26000	3469.7\\
28000	3372\\
30000	3426.7\\
32000	3403.6\\
34000	3396.7\\
36000	3453.3\\
38000	3433.9\\
40000	3443.4\\
42000	3387.7\\
44000	3373.2\\
46000	3414\\
48000	3439.5\\
50000	3425.8\\
52000	3397\\
56000	3437.1\\
58000	3414.2\\
60000	3360.5\\
62000	3424.2\\
64000	3351.7\\
66000	3477.5\\
68000	3387.7\\
72000	3411.2\\
74000	3378.5\\
76000	3381.7\\
78000	3359.9\\
80000	3496.1\\
82000	3413.7\\
84000	3383\\
86000	3402.8\\
88000	3358.9\\
90000	3423\\
92000	3438.9\\
94000	3475.6\\
96000	3371.3\\
98000	3413.4\\
1e+05	3391.1\\
};
\addplot [color=mycolor67, forget plot]
  table[row sep=crcr]{%
0	4395.3\\
2000	3475.3\\
4000	3227.8\\
6000	3131.1\\
8000	3197.3\\
10000	3150.9\\
12000	3139.5\\
14000	3139.8\\
16000	3174.4\\
18000	3125.7\\
22000	3171.9\\
24000	3168.6\\
26000	3216.3\\
30000	3147.4\\
32000	3182.1\\
34000	3140.8\\
36000	3219.8\\
38000	3106.8\\
40000	3248.5\\
42000	3218.5\\
44000	3203.7\\
46000	3204.4\\
48000	3168.6\\
52000	3175.7\\
54000	3216.8\\
56000	3156.7\\
58000	3172.9\\
60000	3107.6\\
62000	3159.6\\
64000	3163.3\\
66000	3234\\
68000	3148.2\\
70000	3217.8\\
72000	3253.7\\
74000	3179.4\\
78000	3182.7\\
80000	3150.6\\
82000	3165.4\\
84000	3208.6\\
86000	3151\\
88000	3216.4\\
90000	3180.9\\
94000	3138.6\\
96000	3137.1\\
98000	3144.9\\
1e+05	3207.9\\
};
\addplot [color=mycolor67]
  table[row sep=crcr]{%
0	4395.3\\
2000	2747.2\\
4000	2625.6\\
6000	2573.2\\
8000	2582.1\\
10000	2557.7\\
14000	2620.8\\
16000	2607.1\\
18000	2640.8\\
20000	2626.1\\
22000	2655.3\\
24000	2617\\
26000	2587.2\\
28000	2635.8\\
30000	2553.4\\
32000	2630.4\\
34000	2598\\
38000	2609.6\\
40000	2642.8\\
42000	2606.3\\
44000	2587.1\\
48000	2586.3\\
50000	2599.3\\
52000	2653.4\\
54000	2611.2\\
56000	2608.6\\
58000	2586.9\\
60000	2593.1\\
62000	2577.7\\
64000	2609.1\\
66000	2627.9\\
72000	2561.9\\
74000	2623.3\\
76000	2618.9\\
78000	2638.8\\
80000	2599.9\\
82000	2572.9\\
84000	2632.5\\
86000	2649.2\\
88000	2608.7\\
90000	2612.4\\
92000	2575.6\\
94000	2563\\
96000	2564.6\\
98000	2606.3\\
1e+05	2614.5\\
};
\addlegendentry{Classe k=43}

\addplot [color=mycolor68, forget plot]
  table[row sep=crcr]{%
0	4395.3\\
2000	2845\\
4000	2550.1\\
6000	2483.8\\
8000	2454.2\\
10000	2409.8\\
12000	2440.1\\
14000	2416\\
16000	2446.2\\
18000	2426.4\\
22000	2448.3\\
24000	2429.4\\
26000	2463.2\\
28000	2472.2\\
30000	2465.8\\
34000	2438.3\\
36000	2442.5\\
38000	2478.2\\
40000	2455.7\\
42000	2464.6\\
44000	2457.1\\
48000	2471.2\\
50000	2442.8\\
52000	2431.8\\
54000	2467.5\\
56000	2424.3\\
58000	2444.1\\
60000	2438.3\\
62000	2415.6\\
64000	2479.7\\
66000	2438.3\\
68000	2445.6\\
70000	2480\\
74000	2466\\
76000	2427.4\\
80000	2503.7\\
82000	2413.6\\
84000	2457.6\\
86000	2419.3\\
90000	2454.2\\
92000	2448.1\\
94000	2465.6\\
96000	2445.6\\
98000	2445.9\\
1e+05	2454.5\\
};
\addplot [color=mycolor69, forget plot]
  table[row sep=crcr]{%
0	4395.3\\
2000	3065.8\\
4000	2789.9\\
6000	2842.7\\
10000	2923.8\\
12000	2895\\
14000	2905.7\\
16000	2811.1\\
18000	2872.3\\
20000	2872.1\\
24000	2894.9\\
26000	2857.3\\
28000	2893.4\\
30000	2879.2\\
32000	2941.5\\
34000	2901.7\\
36000	2906.6\\
38000	2889.8\\
40000	2890.8\\
42000	2922.1\\
44000	2890.2\\
50000	2855.1\\
52000	2881.8\\
54000	2863.1\\
56000	2910.2\\
58000	2908.6\\
60000	2883.9\\
62000	2879.3\\
64000	2910\\
66000	2864.1\\
68000	2933.1\\
70000	2909.9\\
72000	2920.3\\
74000	2848.1\\
76000	2919.8\\
78000	2887.7\\
80000	2901.7\\
82000	2945.9\\
84000	2852.1\\
86000	2901.4\\
88000	2867.6\\
90000	2847.2\\
92000	2850.1\\
94000	2879.5\\
96000	2858.2\\
98000	2880.4\\
1e+05	2894.2\\
};
\addplot [color=mycolor70, forget plot]
  table[row sep=crcr]{%
0	4395.3\\
2000	2252.8\\
4000	2041.9\\
6000	2078\\
8000	2075.4\\
10000	2053.9\\
16000	2101.3\\
18000	2087.8\\
20000	2039.9\\
22000	2053.1\\
24000	2046.7\\
26000	2123.2\\
28000	2034.2\\
30000	2065.6\\
32000	2065.6\\
34000	2056.3\\
36000	2104\\
38000	2086.8\\
40000	2042.5\\
42000	2085\\
44000	2055.7\\
46000	2121.8\\
48000	2136.5\\
50000	2096.8\\
52000	2035.2\\
54000	2054.9\\
56000	2066.2\\
58000	2035\\
60000	2106.7\\
62000	2117.8\\
64000	2090.8\\
66000	2078.4\\
68000	2133.4\\
70000	2087.8\\
72000	2095\\
74000	2043.8\\
76000	2042.4\\
78000	2095.5\\
80000	2086.6\\
82000	2147.6\\
84000	2053.2\\
86000	2097.4\\
88000	2105.6\\
90000	2071.9\\
92000	2110.2\\
94000	2065.2\\
96000	2129.6\\
98000	2135.4\\
1e+05	2120.4\\
};
\addplot [color=mycolor71, forget plot]
  table[row sep=crcr]{%
0	4395.3\\
2000	2763.6\\
4000	2512.6\\
8000	2480.4\\
10000	2491.5\\
12000	2482.7\\
14000	2454.5\\
16000	2462.1\\
18000	2456.3\\
20000	2459.6\\
26000	2439.4\\
28000	2458.8\\
30000	2492.5\\
34000	2462.3\\
36000	2455.7\\
40000	2469.9\\
42000	2497.2\\
46000	2469\\
48000	2483.2\\
50000	2490.2\\
52000	2433.4\\
54000	2443.5\\
56000	2480.3\\
60000	2482.4\\
62000	2497.8\\
64000	2445.9\\
66000	2467.9\\
68000	2480.7\\
70000	2457.4\\
72000	2485.9\\
74000	2485\\
76000	2505.4\\
80000	2458.3\\
82000	2441.3\\
84000	2484.3\\
86000	2468.2\\
88000	2463.7\\
90000	2500.8\\
92000	2468.6\\
1e+05	2483.8\\
};
\addplot [color=mycolor72, forget plot]
  table[row sep=crcr]{%
0	4395.3\\
2000	2574.5\\
4000	2160.4\\
6000	2217\\
8000	2184.6\\
10000	2215.2\\
12000	2171.6\\
16000	2165.4\\
18000	2180.1\\
20000	2167.6\\
22000	2226.4\\
24000	2092.7\\
26000	2140.8\\
28000	2196.8\\
30000	2138.5\\
32000	2137.5\\
34000	2190.1\\
36000	2209.9\\
38000	2183\\
40000	2185.5\\
42000	2154.9\\
44000	2164.3\\
46000	2194.3\\
48000	2176.3\\
50000	2148.7\\
52000	2178.7\\
54000	2175.3\\
58000	2222.9\\
60000	2187.8\\
62000	2161\\
66000	2183.4\\
68000	2163.7\\
70000	2181.3\\
72000	2207.4\\
74000	2159.7\\
76000	2215.3\\
78000	2183.6\\
80000	2251.8\\
82000	2132\\
84000	2184.9\\
86000	2193.8\\
88000	2187.2\\
90000	2207\\
92000	2148.8\\
94000	2169.6\\
96000	2144.7\\
98000	2134.3\\
1e+05	2192.4\\
};
\addplot [color=mycolor73, forget plot]
  table[row sep=crcr]{%
0	4395.3\\
2000	2865\\
4000	2685.7\\
6000	2661.7\\
8000	2589.6\\
10000	2617.9\\
12000	2587.4\\
14000	2610.7\\
16000	2605.2\\
18000	2580.9\\
20000	2573.5\\
22000	2611.5\\
24000	2605.4\\
26000	2580.3\\
28000	2592.2\\
32000	2560.1\\
34000	2608.4\\
36000	2614.1\\
38000	2590.5\\
40000	2618.9\\
42000	2592.9\\
44000	2593.3\\
46000	2625.5\\
48000	2635.6\\
50000	2623.2\\
52000	2564.3\\
54000	2594.8\\
56000	2592.2\\
58000	2618.9\\
60000	2585\\
62000	2578.2\\
64000	2593.1\\
66000	2627.9\\
68000	2585.5\\
70000	2637\\
72000	2596.7\\
74000	2605.2\\
76000	2595.7\\
78000	2551.4\\
80000	2591.6\\
82000	2576.2\\
84000	2610\\
86000	2586.9\\
88000	2572.1\\
92000	2587.4\\
94000	2586\\
96000	2635.9\\
98000	2601.3\\
1e+05	2599.8\\
};
\addplot [color=mycolor74, forget plot]
  table[row sep=crcr]{%
0	4395.3\\
2000	1951.7\\
4000	1505.7\\
6000	1516.9\\
8000	1514.6\\
10000	1537.6\\
12000	1502.7\\
14000	1537.3\\
16000	1541.7\\
18000	1537.7\\
20000	1504.8\\
22000	1496.7\\
24000	1517.8\\
26000	1558.9\\
28000	1524\\
32000	1514.5\\
34000	1524.6\\
36000	1508.5\\
38000	1531.7\\
40000	1502.2\\
42000	1533.4\\
44000	1506.8\\
46000	1533.5\\
48000	1515.7\\
50000	1518.8\\
52000	1534.6\\
54000	1538.9\\
56000	1531.9\\
58000	1490.8\\
60000	1493\\
62000	1508.7\\
64000	1504.4\\
66000	1537.9\\
68000	1536.4\\
70000	1518\\
72000	1542.4\\
74000	1500.2\\
76000	1531.8\\
78000	1539.5\\
82000	1498.7\\
84000	1521.9\\
88000	1536.5\\
90000	1500.5\\
92000	1497.2\\
94000	1512.3\\
96000	1512.9\\
98000	1504.3\\
1e+05	1558.6\\
};
\addplot [color=mycolor75, forget plot]
  table[row sep=crcr]{%
0	4395.3\\
2000	2553.4\\
4000	2229.4\\
6000	2177.2\\
8000	2193.2\\
10000	2201.9\\
12000	2217.2\\
14000	2140.6\\
16000	2178.5\\
18000	2244.6\\
20000	2181.9\\
22000	2218.1\\
24000	2208.9\\
26000	2167\\
28000	2203.6\\
32000	2184.8\\
34000	2183.7\\
36000	2213.8\\
38000	2155.8\\
42000	2159.1\\
44000	2191.3\\
48000	2185.6\\
50000	2181.3\\
52000	2220.9\\
56000	2190.4\\
58000	2195.9\\
60000	2218.3\\
62000	2167.6\\
64000	2189\\
66000	2193.3\\
68000	2185\\
70000	2191.3\\
74000	2228.5\\
76000	2228.6\\
78000	2213.5\\
80000	2177.8\\
82000	2191.7\\
84000	2184\\
86000	2155.2\\
88000	2207.4\\
90000	2176.2\\
92000	2166.7\\
94000	2196.5\\
96000	2181\\
98000	2186.9\\
1e+05	2250.7\\
};
\addplot [color=mycolor76, forget plot]
  table[row sep=crcr]{%
0	4395.3\\
2000	2371.3\\
4000	2043.2\\
6000	1992.7\\
8000	1965.4\\
10000	2003\\
12000	1997.3\\
14000	1998.1\\
16000	1951.6\\
18000	2019.2\\
20000	1996.3\\
24000	1995.2\\
26000	1970.6\\
28000	1991.8\\
30000	1993.6\\
32000	2005.4\\
34000	2009.1\\
36000	1992.7\\
38000	1969.1\\
40000	1998.8\\
44000	2000\\
46000	1963\\
50000	2002.9\\
52000	1980.9\\
54000	2003.1\\
56000	2017.4\\
58000	1978.7\\
62000	2031.6\\
64000	2025.6\\
66000	2054.4\\
68000	1971.8\\
70000	2018.2\\
74000	1981.9\\
76000	2007.5\\
78000	2002.4\\
80000	1968.4\\
82000	2007.1\\
84000	1997.3\\
86000	1935\\
88000	1986.7\\
90000	1986.8\\
92000	1995.8\\
94000	1990.2\\
96000	2008.7\\
1e+05	1986.5\\
};
\addplot [color=mycolor77, forget plot]
  table[row sep=crcr]{%
0	4395.3\\
2000	2076.3\\
4000	1592.1\\
8000	1554.2\\
12000	1525.5\\
14000	1550\\
16000	1548.6\\
20000	1569.3\\
22000	1535\\
24000	1561.8\\
26000	1570.4\\
30000	1540.4\\
32000	1553\\
34000	1571.2\\
36000	1555.4\\
40000	1559.9\\
48000	1531\\
50000	1547.3\\
52000	1555.4\\
54000	1568.1\\
56000	1547.4\\
62000	1556.7\\
64000	1543.9\\
66000	1559.9\\
68000	1538.8\\
70000	1553.3\\
72000	1526.4\\
74000	1533.1\\
76000	1534\\
80000	1566\\
82000	1534.2\\
84000	1541.7\\
86000	1541.2\\
88000	1524.7\\
90000	1550.5\\
92000	1554\\
94000	1543.4\\
96000	1570.1\\
98000	1554.6\\
1e+05	1548.3\\
};
\addplot [color=mycolor78, forget plot]
  table[row sep=crcr]{%
0	4395.3\\
2000	2681.1\\
4000	2340.9\\
6000	2339.1\\
8000	2295.7\\
10000	2261.4\\
12000	2288.5\\
14000	2293.9\\
16000	2277.1\\
18000	2302.7\\
22000	2311\\
24000	2294.6\\
26000	2258.9\\
28000	2271.5\\
30000	2229.3\\
32000	2286.1\\
34000	2234.9\\
36000	2275.4\\
42000	2257\\
44000	2276.5\\
48000	2288.5\\
50000	2249.7\\
52000	2261.3\\
54000	2295.6\\
56000	2268.4\\
58000	2270.2\\
60000	2246.1\\
62000	2251.8\\
64000	2285.5\\
68000	2228.2\\
72000	2306.4\\
74000	2308.9\\
76000	2330.2\\
78000	2324.1\\
80000	2247.3\\
82000	2275.1\\
86000	2287.7\\
88000	2267.1\\
90000	2323.3\\
92000	2259.8\\
94000	2240.6\\
96000	2249.9\\
98000	2242.4\\
1e+05	2284.5\\
};
\addplot [color=mycolor79, forget plot]
  table[row sep=crcr]{%
0	4395.3\\
2000	2482.2\\
4000	2206.4\\
6000	2123.8\\
8000	2111.5\\
10000	2150.3\\
12000	2133.6\\
14000	2172.3\\
16000	2185.9\\
18000	2167.4\\
20000	2105.4\\
22000	2125.2\\
24000	2099.6\\
26000	2135.6\\
28000	2161.5\\
30000	2134.8\\
32000	2133.2\\
34000	2155.5\\
36000	2144.5\\
38000	2160.9\\
42000	2105.3\\
44000	2152.6\\
46000	2124.4\\
48000	2180\\
52000	2110.4\\
54000	2091.8\\
56000	2154.8\\
58000	2188.1\\
60000	2155.4\\
62000	2138.1\\
64000	2135.2\\
66000	2121.2\\
68000	2097.5\\
70000	2103.9\\
72000	2195.7\\
76000	2158.5\\
78000	2123.1\\
80000	2140.8\\
82000	2109.6\\
84000	2124.7\\
86000	2157.2\\
88000	2147\\
90000	2185.8\\
92000	2149.3\\
94000	2159.3\\
96000	2113.7\\
98000	2124.7\\
1e+05	2150.9\\
};
\addplot [color=mycolor79, forget plot]
  table[row sep=crcr]{%
0	4395.3\\
2000	2456.4\\
4000	2056.6\\
6000	1971.8\\
8000	1980.5\\
12000	2026.1\\
14000	1967\\
16000	2032.9\\
18000	1986.1\\
20000	1967\\
22000	1992.3\\
24000	1989.4\\
26000	2010.5\\
28000	2003.8\\
30000	1956.7\\
32000	1962.3\\
34000	2014.8\\
36000	2001.1\\
38000	1993.7\\
40000	1967.2\\
42000	1983.7\\
44000	2034.1\\
46000	1980.8\\
48000	2006.9\\
50000	2005.9\\
52000	1989\\
54000	1993\\
56000	2011.3\\
58000	1965.3\\
60000	2004.3\\
62000	1942.3\\
64000	2001.9\\
66000	1993.6\\
68000	2029.3\\
70000	2079.3\\
72000	2003\\
74000	1987.9\\
78000	1985\\
80000	1997.3\\
82000	1960.8\\
84000	1982.9\\
86000	2031.3\\
88000	2053.6\\
90000	1953.5\\
92000	2007.6\\
94000	2016.8\\
96000	1999.6\\
98000	1998\\
1e+05	2026.8\\
};
\addplot [color=mycolor80, forget plot]
  table[row sep=crcr]{%
0	4395.3\\
2000	1978.3\\
4000	1546.2\\
6000	1538.3\\
8000	1515.2\\
10000	1534.9\\
12000	1526.5\\
14000	1525.4\\
16000	1550.2\\
18000	1546.5\\
20000	1555\\
22000	1527\\
24000	1521.8\\
26000	1527.6\\
28000	1525.5\\
30000	1510\\
32000	1520.6\\
34000	1525.1\\
36000	1544\\
38000	1517.5\\
40000	1537.9\\
42000	1518.2\\
46000	1539.1\\
48000	1511.4\\
50000	1502.5\\
54000	1547.3\\
56000	1536.5\\
58000	1536.5\\
60000	1529.6\\
62000	1552.2\\
64000	1542.8\\
66000	1520.6\\
68000	1532.9\\
72000	1511.3\\
74000	1546.1\\
78000	1524.9\\
80000	1550.2\\
82000	1520.2\\
84000	1529.2\\
86000	1520.8\\
88000	1543.1\\
90000	1536.1\\
92000	1518.7\\
94000	1523.2\\
96000	1539.9\\
98000	1539.2\\
1e+05	1543.6\\
};
\addplot [color=mycolor81, forget plot]
  table[row sep=crcr]{%
0	4395.3\\
2000	1930.9\\
4000	1430.2\\
6000	1450.1\\
8000	1447.9\\
12000	1420.2\\
14000	1449.4\\
18000	1420.4\\
20000	1398.3\\
22000	1451.8\\
24000	1416.3\\
26000	1462.5\\
28000	1441.6\\
30000	1438.4\\
32000	1416.5\\
34000	1413.7\\
36000	1435.2\\
38000	1418.5\\
40000	1427\\
42000	1418.8\\
44000	1466.3\\
46000	1408.8\\
48000	1444\\
50000	1460.8\\
52000	1432.2\\
54000	1422.5\\
56000	1446.1\\
58000	1433.5\\
60000	1456.8\\
62000	1409.9\\
64000	1434.7\\
68000	1421.7\\
70000	1427.7\\
72000	1422.9\\
74000	1432.3\\
76000	1415.3\\
78000	1424.5\\
80000	1408.2\\
82000	1436.9\\
84000	1426.9\\
86000	1407.6\\
88000	1467.6\\
90000	1467.2\\
92000	1455.7\\
94000	1425.8\\
98000	1424.6\\
1e+05	1441.8\\
};
\addplot [color=mycolor82, forget plot]
  table[row sep=crcr]{%
0	4395.3\\
2000	2207.9\\
4000	1726.9\\
8000	1667.4\\
10000	1660.1\\
12000	1693.5\\
14000	1666.8\\
16000	1684.5\\
18000	1649.2\\
20000	1671.2\\
22000	1715.3\\
24000	1714.9\\
26000	1695.3\\
30000	1695.2\\
32000	1711.1\\
34000	1668.9\\
36000	1679.6\\
38000	1660.3\\
40000	1666.7\\
42000	1694.7\\
44000	1689.5\\
46000	1674.3\\
48000	1686.2\\
50000	1688.1\\
52000	1666\\
54000	1694.8\\
56000	1682.9\\
58000	1702.1\\
60000	1711.3\\
62000	1685.5\\
66000	1695.3\\
68000	1722.5\\
70000	1686\\
72000	1728.7\\
74000	1691\\
76000	1693.8\\
78000	1671.5\\
80000	1683.7\\
82000	1700.8\\
84000	1704.7\\
86000	1688.8\\
88000	1684.8\\
90000	1690.5\\
92000	1683.2\\
94000	1653.2\\
96000	1651.1\\
98000	1655.2\\
1e+05	1648.3\\
};
\addplot [color=mycolor83, forget plot]
  table[row sep=crcr]{%
0	4395.3\\
2000	2200.4\\
4000	1773.8\\
6000	1739.9\\
8000	1722.3\\
10000	1716\\
12000	1690.1\\
14000	1734.2\\
16000	1705.1\\
18000	1719.8\\
20000	1725\\
22000	1691.8\\
24000	1695.7\\
26000	1713.2\\
32000	1722.3\\
34000	1705.7\\
36000	1719.2\\
38000	1690.6\\
40000	1705.6\\
42000	1696\\
44000	1714\\
46000	1725.3\\
48000	1713.5\\
50000	1695.8\\
52000	1753.4\\
54000	1718.6\\
56000	1706.3\\
58000	1712\\
60000	1690.3\\
62000	1718.5\\
64000	1728.3\\
66000	1729.7\\
68000	1685\\
70000	1718.8\\
72000	1718.1\\
74000	1705.4\\
76000	1712.3\\
78000	1687.9\\
80000	1686.1\\
82000	1662.7\\
84000	1677\\
86000	1729.2\\
92000	1699.7\\
94000	1714.7\\
96000	1715.2\\
98000	1729.7\\
1e+05	1707.8\\
};
\addplot [color=mycolor83, forget plot]
  table[row sep=crcr]{%
0	4395.3\\
2000	1700.8\\
4000	1281\\
6000	1239.3\\
8000	1248.3\\
10000	1269.2\\
12000	1211\\
14000	1238.7\\
16000	1254.3\\
18000	1252.2\\
20000	1241.8\\
22000	1244\\
24000	1229.8\\
26000	1251\\
28000	1229.5\\
30000	1224.4\\
32000	1250.3\\
34000	1254.2\\
36000	1208.2\\
38000	1232.8\\
40000	1226\\
42000	1247.6\\
44000	1242.4\\
46000	1225.6\\
48000	1280.4\\
50000	1239.5\\
52000	1227.5\\
54000	1252.6\\
56000	1233.6\\
58000	1220.5\\
60000	1247.4\\
62000	1236.6\\
64000	1245.6\\
66000	1235.8\\
68000	1238.6\\
70000	1231.7\\
72000	1246.4\\
74000	1230.4\\
76000	1250.8\\
78000	1238.6\\
80000	1248.9\\
82000	1245.1\\
84000	1255.3\\
86000	1195.9\\
88000	1240.2\\
92000	1258.3\\
94000	1234.5\\
96000	1239.7\\
98000	1230\\
1e+05	1249\\
};
\addplot [color=mycolor84, forget plot]
  table[row sep=crcr]{%
0	4395.3\\
2000	1763.8\\
4000	1254.2\\
6000	1253.6\\
8000	1242.1\\
12000	1255.3\\
14000	1244.6\\
16000	1269.4\\
18000	1279.2\\
20000	1263.8\\
24000	1245.4\\
26000	1238.9\\
28000	1236.8\\
30000	1250.5\\
32000	1240.6\\
34000	1258.9\\
36000	1234.4\\
42000	1278\\
44000	1248.1\\
46000	1273\\
48000	1246\\
50000	1240.9\\
52000	1256.9\\
54000	1248.3\\
56000	1247.3\\
58000	1255\\
60000	1226.7\\
62000	1247.6\\
64000	1253.4\\
66000	1213\\
68000	1233.5\\
70000	1258.5\\
72000	1258.3\\
74000	1253.4\\
76000	1269.5\\
78000	1257.2\\
80000	1258.6\\
82000	1273.4\\
84000	1243.6\\
86000	1271\\
88000	1244\\
90000	1241.8\\
92000	1236.1\\
94000	1264.7\\
96000	1263.8\\
98000	1243.5\\
1e+05	1262.3\\
};
\addplot [color=mycolor85, forget plot]
  table[row sep=crcr]{%
0	4395.3\\
2000	2004.9\\
4000	1462.3\\
6000	1346.3\\
8000	1391.5\\
10000	1400.4\\
12000	1389.6\\
14000	1344.3\\
16000	1388.1\\
18000	1379.1\\
20000	1383.2\\
22000	1357.3\\
24000	1382.6\\
26000	1392.8\\
28000	1368.4\\
30000	1382.9\\
32000	1359.1\\
34000	1357.1\\
36000	1411.9\\
38000	1392.1\\
42000	1378.1\\
44000	1407.8\\
46000	1383.3\\
48000	1327.5\\
50000	1389.1\\
52000	1334.4\\
54000	1375\\
56000	1370.7\\
58000	1429.3\\
60000	1393.7\\
62000	1382.8\\
64000	1398.4\\
66000	1377.8\\
68000	1391.8\\
70000	1389.8\\
72000	1370.2\\
76000	1393.4\\
78000	1385\\
80000	1400.5\\
82000	1352.3\\
84000	1384.4\\
86000	1409.5\\
88000	1378.7\\
90000	1390.9\\
92000	1383.3\\
94000	1379.7\\
96000	1383.3\\
98000	1378.5\\
1e+05	1363.5\\
};
\addplot [color=mycolor86, forget plot]
  table[row sep=crcr]{%
0	4395.3\\
2000	1769.9\\
4000	1231.5\\
6000	1156\\
8000	1144.6\\
10000	1115.6\\
12000	1135.2\\
14000	1196.4\\
18000	1123.7\\
20000	1155.9\\
22000	1141.1\\
24000	1148.2\\
26000	1144.7\\
28000	1130.6\\
30000	1142.6\\
32000	1165.2\\
34000	1170.3\\
36000	1190.5\\
38000	1158.2\\
40000	1165.5\\
42000	1159.8\\
44000	1147.9\\
48000	1160.1\\
50000	1134.4\\
52000	1166.7\\
54000	1167.3\\
56000	1143.6\\
58000	1159.9\\
60000	1147.4\\
62000	1163.2\\
64000	1189.8\\
68000	1142.8\\
70000	1129.3\\
72000	1149.8\\
74000	1164.8\\
76000	1110.7\\
78000	1156.9\\
80000	1144.7\\
82000	1163.3\\
84000	1124.2\\
86000	1146.9\\
88000	1139.3\\
90000	1166.3\\
92000	1142.4\\
94000	1155.9\\
96000	1130.4\\
98000	1141.3\\
1e+05	1160.3\\
};
\addplot [color=mycolor87, forget plot]
  table[row sep=crcr]{%
0	4395.3\\
2000	1481.1\\
4000	957.75\\
6000	1007.3\\
8000	953.5\\
10000	938.25\\
12000	932.45\\
16000	1012.4\\
18000	995.45\\
20000	986.97\\
22000	1005.9\\
24000	932.48\\
26000	928.68\\
28000	950.54\\
30000	968.15\\
32000	994.5\\
34000	959.87\\
36000	985.04\\
38000	967.52\\
40000	1024.8\\
42000	1012.1\\
44000	949.11\\
46000	961.93\\
48000	986.83\\
50000	963.99\\
52000	994.03\\
56000	968.23\\
58000	966.97\\
60000	999.37\\
62000	968.47\\
64000	961.15\\
66000	1008.1\\
68000	981.57\\
70000	951.97\\
72000	951.37\\
74000	978.1\\
76000	990.02\\
78000	950.98\\
80000	979.68\\
82000	967.34\\
84000	899.47\\
86000	1001.9\\
88000	953.74\\
90000	972.14\\
92000	971.72\\
96000	984.6\\
98000	974.72\\
1e+05	973.29\\
};
\addplot [color=mycolor88, forget plot]
  table[row sep=crcr]{%
0	4395.3\\
2000	1809.2\\
4000	1218.2\\
6000	1203.8\\
8000	1174.8\\
10000	1169.3\\
12000	1171\\
14000	1133.9\\
16000	1119.5\\
18000	1167.9\\
20000	1184.6\\
22000	1171.2\\
24000	1147.2\\
26000	1167.7\\
28000	1142.7\\
30000	1158.2\\
32000	1154.5\\
34000	1184.2\\
36000	1174.4\\
38000	1154.6\\
40000	1172.9\\
42000	1140.4\\
44000	1177.8\\
46000	1165.7\\
48000	1180.2\\
50000	1158.1\\
52000	1164.6\\
54000	1159.5\\
56000	1143.8\\
58000	1186.4\\
60000	1206.2\\
62000	1179.8\\
64000	1160.6\\
68000	1201.3\\
70000	1214.4\\
72000	1165.3\\
74000	1171.9\\
76000	1172.9\\
78000	1214.9\\
80000	1176.5\\
82000	1158.4\\
84000	1181.8\\
88000	1154\\
90000	1158.7\\
92000	1149.2\\
94000	1165.1\\
96000	1172.7\\
98000	1141.3\\
1e+05	1157.2\\
};
\addplot [color=mycolor89, forget plot]
  table[row sep=crcr]{%
0	4395.3\\
2000	1548.4\\
4000	851.28\\
6000	823.94\\
8000	811.73\\
10000	842.52\\
12000	841.85\\
14000	834.8\\
16000	817.15\\
18000	830.7\\
22000	809.12\\
24000	806.17\\
26000	820.63\\
28000	811.16\\
30000	831.22\\
32000	842.72\\
34000	850.61\\
36000	828.17\\
40000	840.51\\
42000	831.83\\
44000	816.8\\
46000	812.3\\
48000	844.54\\
50000	823.1\\
52000	828.8\\
54000	848\\
56000	835.6\\
58000	828.76\\
60000	832.75\\
62000	852.85\\
64000	857.69\\
66000	839.83\\
68000	811.34\\
70000	835.88\\
72000	845.6\\
74000	837.37\\
76000	835.89\\
78000	825.16\\
80000	822.89\\
82000	842.26\\
84000	821.16\\
86000	851.77\\
88000	855.47\\
90000	827.69\\
92000	840.97\\
94000	837.37\\
96000	824.76\\
98000	805.27\\
1e+05	847.1\\
};
\addplot [color=mycolor90, forget plot]
  table[row sep=crcr]{%
0	4395.3\\
2000	1356.8\\
4000	810.3\\
6000	797.45\\
8000	796.55\\
10000	807.54\\
12000	815.7\\
14000	817.8\\
16000	833.19\\
18000	803.29\\
20000	814.89\\
22000	811.36\\
24000	787.25\\
26000	808.02\\
28000	782.83\\
30000	818.94\\
32000	809.94\\
34000	811.07\\
36000	796.41\\
38000	807.2\\
40000	791.28\\
42000	791.06\\
44000	812.2\\
46000	813.2\\
48000	830.96\\
50000	804.48\\
52000	811.43\\
54000	771.91\\
56000	797.56\\
58000	819.78\\
60000	823.12\\
62000	804.33\\
64000	808.87\\
66000	810.8\\
68000	839.22\\
70000	791.02\\
72000	807.98\\
74000	794.65\\
76000	806.46\\
78000	803.06\\
80000	790.49\\
82000	794.04\\
84000	807.02\\
86000	816.12\\
88000	817.38\\
90000	840.71\\
92000	828.62\\
94000	824.24\\
96000	794.46\\
98000	792.38\\
1e+05	786.62\\
};
\addplot [color=mycolor91, forget plot]
  table[row sep=crcr]{%
0	4395.3\\
2000	1543.7\\
4000	730.11\\
6000	689.45\\
8000	709.25\\
10000	688.66\\
12000	706.88\\
14000	682.55\\
16000	693.44\\
18000	689.8\\
20000	678.17\\
22000	681.19\\
24000	680.8\\
26000	691.46\\
28000	715.49\\
30000	706.81\\
32000	706.11\\
34000	696.5\\
36000	743.17\\
38000	706.91\\
40000	704.39\\
42000	682.41\\
44000	666.22\\
46000	729.05\\
48000	703.96\\
50000	704.82\\
52000	710.25\\
54000	684.9\\
56000	683.2\\
58000	708.05\\
60000	683.61\\
64000	695.27\\
66000	695.38\\
68000	684.62\\
70000	679.41\\
74000	714.74\\
76000	717.97\\
78000	688.37\\
80000	684.87\\
82000	688.03\\
84000	729.83\\
86000	712.44\\
88000	689.4\\
90000	692.9\\
92000	701.62\\
94000	687.08\\
96000	701.61\\
98000	694.67\\
1e+05	707.49\\
};
\addplot [color=mycolor92, forget plot]
  table[row sep=crcr]{%
0	4395.3\\
2000	1601.8\\
4000	893.99\\
6000	889.1\\
8000	893.13\\
10000	860.65\\
12000	847.45\\
14000	853.44\\
16000	862.27\\
20000	873.44\\
22000	873.32\\
24000	842.41\\
26000	851.78\\
28000	854.1\\
30000	876.91\\
32000	892.33\\
34000	895.85\\
36000	878.1\\
38000	844.06\\
40000	852.88\\
42000	888.94\\
46000	868.3\\
48000	850.36\\
50000	874.55\\
52000	873.45\\
54000	865.33\\
56000	853.49\\
58000	899.48\\
62000	849.64\\
64000	857.04\\
68000	848.79\\
70000	875.96\\
72000	885.99\\
74000	873.84\\
76000	853.14\\
78000	841.51\\
80000	859.94\\
82000	858.76\\
84000	881.43\\
86000	862.95\\
88000	849.32\\
90000	874.19\\
92000	870.4\\
94000	907.33\\
96000	931.53\\
98000	878.74\\
1e+05	874.24\\
};
\addplot [color=mycolor93, forget plot]
  table[row sep=crcr]{%
0	4395.3\\
2000	1572.7\\
4000	762.87\\
6000	740.17\\
8000	794.9\\
10000	765.42\\
12000	752.14\\
14000	746.92\\
16000	745.01\\
18000	735.46\\
20000	747.86\\
22000	725.06\\
24000	705.47\\
26000	749.01\\
28000	715.53\\
30000	743.84\\
32000	699.09\\
34000	730.73\\
36000	751.27\\
38000	746.49\\
40000	771.02\\
42000	741.85\\
44000	728.39\\
46000	701.09\\
48000	751.96\\
50000	715.05\\
52000	756.46\\
54000	737.44\\
56000	772.86\\
58000	710.64\\
60000	717.04\\
62000	738.38\\
64000	706.64\\
66000	749.37\\
70000	715.22\\
72000	722.35\\
74000	717.15\\
76000	719.82\\
78000	740.92\\
80000	704.37\\
82000	702.8\\
84000	723.85\\
86000	734.41\\
88000	719.2\\
90000	768.87\\
92000	713.29\\
94000	740.22\\
96000	696.05\\
98000	732.26\\
1e+05	737.61\\
};
\addplot [color=mycolor94, forget plot]
  table[row sep=crcr]{%
0	4395.3\\
2000	1590.8\\
4000	724.76\\
6000	690.3\\
8000	714.17\\
10000	703.38\\
12000	706.33\\
14000	690.1\\
16000	676.68\\
18000	697.91\\
20000	701.14\\
22000	698.24\\
24000	680.43\\
26000	683.12\\
28000	701.86\\
30000	703.44\\
32000	687.32\\
34000	695.89\\
36000	695.52\\
38000	698.52\\
40000	711.24\\
42000	695.64\\
44000	676.18\\
46000	701.48\\
48000	686.32\\
50000	684.62\\
54000	702.01\\
56000	706.79\\
60000	683.61\\
62000	694.75\\
64000	681.92\\
66000	676.21\\
68000	693.33\\
70000	681.97\\
72000	703.07\\
74000	713.91\\
76000	702.14\\
78000	707.34\\
80000	682.43\\
82000	675.6\\
84000	690.82\\
86000	689.64\\
88000	713.08\\
92000	698.35\\
94000	675.05\\
96000	659.15\\
98000	695.19\\
1e+05	692.18\\
};
\addplot [color=mycolor95, forget plot]
  table[row sep=crcr]{%
0	4395.3\\
2000	1256\\
4000	450.52\\
8000	415.8\\
10000	429.45\\
12000	407.2\\
14000	405.66\\
16000	396.78\\
18000	408.97\\
22000	401.98\\
24000	406.24\\
26000	405.97\\
28000	392.42\\
30000	400.72\\
32000	403.06\\
34000	408.52\\
36000	407.69\\
38000	430.91\\
40000	427.25\\
42000	411.65\\
44000	412.44\\
46000	418.6\\
48000	396.53\\
50000	396.12\\
52000	382.9\\
54000	391.06\\
56000	423.3\\
58000	409.54\\
60000	427.04\\
62000	402.72\\
64000	397.35\\
66000	400.83\\
68000	410.14\\
70000	405.77\\
72000	410.5\\
74000	412.45\\
76000	405.69\\
78000	418.33\\
80000	393.49\\
82000	409.58\\
84000	400.24\\
86000	425.31\\
88000	413.52\\
90000	416.12\\
94000	409.34\\
96000	394.06\\
98000	420.71\\
1e+05	418.83\\
};
\addplot [color=mycolor96, forget plot]
  table[row sep=crcr]{%
0	4395.3\\
2000	1442.1\\
4000	457.41\\
6000	443.51\\
8000	411.21\\
10000	397.04\\
12000	407.35\\
14000	414.21\\
16000	400.72\\
18000	421.67\\
20000	419.04\\
22000	406.94\\
24000	447.47\\
26000	441.07\\
28000	398.14\\
30000	416.3\\
32000	439.17\\
34000	449.57\\
36000	453.52\\
38000	432\\
40000	452.93\\
42000	427.6\\
44000	414.57\\
46000	421.62\\
48000	416.35\\
50000	469.45\\
52000	452.1\\
54000	455.34\\
56000	506.39\\
58000	447.54\\
60000	448.21\\
62000	460.82\\
64000	412.81\\
66000	431.65\\
68000	428.43\\
70000	456.67\\
72000	450.86\\
76000	410.46\\
78000	477.65\\
80000	430.25\\
82000	434.88\\
84000	440.79\\
86000	422.26\\
88000	418.54\\
90000	460.29\\
92000	454.61\\
94000	433.51\\
96000	446.56\\
98000	423.73\\
1e+05	411.91\\
};
\addplot [color=mycolor97, forget plot]
  table[row sep=crcr]{%
0	4395.3\\
2000	1080.2\\
4000	334.46\\
6000	339.5\\
8000	323.95\\
10000	328.51\\
12000	335.07\\
14000	307.21\\
16000	308.08\\
18000	333.21\\
20000	309.79\\
22000	316.19\\
24000	301.49\\
26000	284.06\\
28000	293.5\\
30000	310.94\\
32000	350.01\\
34000	339.8\\
36000	293.45\\
38000	307.92\\
40000	291.77\\
42000	286.22\\
44000	300.56\\
46000	325.7\\
48000	321.43\\
50000	323.18\\
52000	285.24\\
54000	279.95\\
58000	302.3\\
60000	299.12\\
62000	292.29\\
64000	278.89\\
66000	317.33\\
68000	322.6\\
70000	314.04\\
72000	323.11\\
74000	318.89\\
76000	308.42\\
78000	323.88\\
80000	321.24\\
82000	299.9\\
84000	300.1\\
86000	303.75\\
88000	325.19\\
90000	337.43\\
92000	337.94\\
94000	327.91\\
96000	313.43\\
98000	318.23\\
1e+05	316.45\\
};
\addplot [color=mycolor98]
  table[row sep=crcr]{%
0	4395.3\\
2000	1517\\
4000	494.38\\
6000	417.61\\
8000	408.83\\
10000	398.74\\
12000	403.76\\
14000	415.34\\
16000	421.94\\
18000	409.26\\
20000	420.12\\
22000	414.74\\
24000	432.97\\
26000	448.65\\
28000	426.05\\
30000	414.77\\
32000	397.59\\
34000	389.49\\
36000	420.27\\
38000	409.68\\
40000	428.89\\
42000	415.45\\
44000	432.72\\
46000	404.4\\
48000	403.08\\
50000	423.53\\
52000	424.72\\
54000	421.95\\
56000	412.52\\
58000	411.02\\
60000	407.78\\
62000	430.14\\
64000	414.8\\
66000	405.83\\
68000	406.59\\
70000	402.67\\
72000	423.37\\
74000	421.35\\
76000	411.46\\
78000	413\\
80000	406.17\\
82000	428.11\\
84000	419.59\\
86000	424.23\\
88000	422.48\\
90000	411.89\\
92000	430.87\\
96000	441.5\\
98000	422.33\\
1e+05	436.98\\
};
\addlegendentry{Classe k=10}

            % \input{../../../../.tex/20/KS/SizeEvolutionsfig7.tex}
            \legend{}

            \nextgroupplot[%
                % xmin=0,xmax=1e+05,
                % ymin=100,ymax=1.5689e+05,
            ] % This file was created by matlab2tikz.
%
\definecolor{mycolor1}{rgb}{0.00146,0.00047,0.01387}%
\definecolor{mycolor2}{rgb}{0.01166,0.00942,0.06346}%
\definecolor{mycolor3}{rgb}{0.02943,0.02150,0.11462}%
\definecolor{mycolor4}{rgb}{0.06134,0.03659,0.17764}%
\definecolor{mycolor5}{rgb}{0.08741,0.04456,0.22481}%
\definecolor{mycolor6}{rgb}{0.12291,0.04754,0.28162}%
\definecolor{mycolor7}{rgb}{0.14907,0.04547,0.31709}%
\definecolor{mycolor8}{rgb}{0.18343,0.04033,0.35497}%
\definecolor{mycolor9}{rgb}{0.21795,0.03662,0.38352}%
\definecolor{mycolor10}{rgb}{0.24497,0.03705,0.40001}%
\definecolor{mycolor11}{rgb}{0.27135,0.04092,0.41198}%
\definecolor{mycolor12}{rgb}{0.29076,0.04564,0.41864}%
\definecolor{mycolor13}{rgb}{0.31628,0.05349,0.42512}%
\definecolor{mycolor14}{rgb}{0.33522,0.06006,0.42852}%
\definecolor{mycolor15}{rgb}{0.36028,0.06925,0.43150}%
\definecolor{mycolor16}{rgb}{0.37900,0.07625,0.43272}%
\definecolor{mycolor17}{rgb}{0.39767,0.08326,0.43318}%
\definecolor{mycolor18}{rgb}{0.41633,0.09020,0.43294}%
\definecolor{mycolor19}{rgb}{0.42877,0.09479,0.43241}%
\definecolor{mycolor20}{rgb}{0.44743,0.10160,0.43108}%
\definecolor{mycolor21}{rgb}{0.46610,0.10832,0.42912}%
\definecolor{mycolor22}{rgb}{0.47856,0.11276,0.42747}%
\definecolor{mycolor23}{rgb}{0.49726,0.11938,0.42449}%
\definecolor{mycolor24}{rgb}{0.50973,0.12377,0.42216}%
\definecolor{mycolor25}{rgb}{0.52221,0.12815,0.41955}%
\definecolor{mycolor26}{rgb}{0.53468,0.13253,0.41667}%
\definecolor{mycolor27}{rgb}{0.55339,0.13913,0.41183}%
\definecolor{mycolor28}{rgb}{0.56585,0.14357,0.40826}%
\definecolor{mycolor29}{rgb}{0.57830,0.14804,0.40441}%
\definecolor{mycolor30}{rgb}{0.59073,0.15256,0.40029}%
\definecolor{mycolor31}{rgb}{0.60314,0.15715,0.39589}%
\definecolor{mycolor32}{rgb}{0.61551,0.16182,0.39122}%
\definecolor{mycolor33}{rgb}{0.62169,0.16418,0.38878}%
\definecolor{mycolor34}{rgb}{0.63400,0.16899,0.38370}%
\definecolor{mycolor35}{rgb}{0.64626,0.17391,0.37836}%
\definecolor{mycolor36}{rgb}{0.65846,0.17896,0.37275}%
\definecolor{mycolor37}{rgb}{0.66454,0.18154,0.36985}%
\definecolor{mycolor38}{rgb}{0.67664,0.18681,0.36385}%
\definecolor{mycolor39}{rgb}{0.68865,0.19224,0.35760}%
\definecolor{mycolor40}{rgb}{0.69463,0.19502,0.35439}%
\definecolor{mycolor41}{rgb}{0.70650,0.20073,0.34778}%
\definecolor{mycolor42}{rgb}{0.71240,0.20366,0.34438}%
\definecolor{mycolor43}{rgb}{0.72410,0.20967,0.33742}%
\definecolor{mycolor44}{rgb}{0.72991,0.21276,0.33386}%
\definecolor{mycolor45}{rgb}{0.74142,0.21911,0.32658}%
\definecolor{mycolor46}{rgb}{0.74713,0.22238,0.32286}%
\definecolor{mycolor47}{rgb}{0.75279,0.22571,0.31909}%
\definecolor{mycolor48}{rgb}{0.76401,0.23255,0.31140}%
\definecolor{mycolor49}{rgb}{0.76956,0.23608,0.30749}%
\definecolor{mycolor50}{rgb}{0.77506,0.23967,0.30353}%
\definecolor{mycolor51}{rgb}{0.79129,0.25086,0.29139}%
\definecolor{mycolor52}{rgb}{0.79661,0.25473,0.28726}%
\definecolor{mycolor53}{rgb}{0.80187,0.25867,0.28310}%
\definecolor{mycolor54}{rgb}{0.81224,0.26679,0.27466}%
\definecolor{mycolor55}{rgb}{0.81734,0.27095,0.27039}%
\definecolor{mycolor56}{rgb}{0.82737,0.27952,0.26175}%
\definecolor{mycolor57}{rgb}{0.83230,0.28391,0.25738}%
\definecolor{mycolor58}{rgb}{0.83717,0.28839,0.25299}%
\definecolor{mycolor59}{rgb}{0.84671,0.29756,0.24411}%
\definecolor{mycolor60}{rgb}{0.85138,0.30226,0.23964}%
\definecolor{mycolor61}{rgb}{0.85599,0.30704,0.23513}%
\definecolor{mycolor62}{rgb}{0.86053,0.31189,0.23061}%
\definecolor{mycolor63}{rgb}{0.87374,0.32691,0.21689}%
\definecolor{mycolor64}{rgb}{0.87800,0.33206,0.21227}%
\definecolor{mycolor65}{rgb}{0.89034,0.34796,0.19829}%
\definecolor{mycolor66}{rgb}{0.89819,0.35891,0.18886}%
\definecolor{mycolor67}{rgb}{0.90200,0.36449,0.18412}%
\definecolor{mycolor68}{rgb}{0.90573,0.37014,0.17935}%
\definecolor{mycolor69}{rgb}{0.90939,0.37586,0.17456}%
\definecolor{mycolor70}{rgb}{0.91297,0.38164,0.16975}%
\definecolor{mycolor71}{rgb}{0.91646,0.38748,0.16492}%
\definecolor{mycolor72}{rgb}{0.91988,0.39339,0.16007}%
\definecolor{mycolor73}{rgb}{0.92322,0.39936,0.15519}%
\definecolor{mycolor74}{rgb}{0.92647,0.40539,0.15029}%
\definecolor{mycolor75}{rgb}{0.92964,0.41148,0.14537}%
\definecolor{mycolor76}{rgb}{0.93274,0.41763,0.14042}%
\definecolor{mycolor77}{rgb}{0.93575,0.42383,0.13544}%
\definecolor{mycolor78}{rgb}{0.93868,0.43009,0.13044}%
\definecolor{mycolor79}{rgb}{0.94152,0.43640,0.12541}%
\definecolor{mycolor80}{rgb}{0.94429,0.44277,0.12035}%
\definecolor{mycolor81}{rgb}{0.94956,0.45566,0.11016}%
\definecolor{mycolor82}{rgb}{0.95208,0.46218,0.10503}%
\definecolor{mycolor83}{rgb}{0.96129,0.48872,0.08429}%
\definecolor{mycolor84}{rgb}{0.96540,0.50225,0.07386}%
\definecolor{mycolor85}{rgb}{0.96732,0.50908,0.06866}%
\definecolor{mycolor86}{rgb}{0.97259,0.52980,0.05332}%
\definecolor{mycolor87}{rgb}{0.97568,0.54380,0.04362}%
\definecolor{mycolor88}{rgb}{0.98082,0.57221,0.02851}%
\definecolor{mycolor89}{rgb}{0.98189,0.57939,0.02625}%
\definecolor{mycolor90}{rgb}{0.98378,0.59385,0.02377}%
\definecolor{mycolor91}{rgb}{0.98532,0.60842,0.02420}%
\definecolor{mycolor92}{rgb}{0.98696,0.63048,0.03091}%
\definecolor{mycolor93}{rgb}{0.98793,0.66025,0.05175}%
\definecolor{mycolor94}{rgb}{0.98771,0.68281,0.07249}%
\definecolor{mycolor95}{rgb}{0.97750,0.78226,0.18592}%
\definecolor{mycolor96}{rgb}{0.96624,0.83619,0.26153}%
\definecolor{mycolor97}{rgb}{0.96252,0.85148,0.28555}%
\definecolor{mycolor98}{rgb}{0.98836,0.99836,0.64492}%
%
\addplot [color=mycolor1]
  table[row sep=crcr]{%
0	4395.3\\
2000	1.1533e+05\\
4000	1.515e+05\\
6000	1.5239e+05\\
8000	1.5187e+05\\
12000	1.5303e+05\\
14000	1.533e+05\\
16000	1.5689e+05\\
18000	1.5425e+05\\
20000	1.5622e+05\\
22000	1.5645e+05\\
24000	1.5509e+05\\
26000	1.5548e+05\\
28000	1.5415e+05\\
30000	1.5231e+05\\
32000	1.5001e+05\\
34000	1.5196e+05\\
36000	1.5542e+05\\
38000	1.5631e+05\\
40000	1.554e+05\\
42000	1.5211e+05\\
44000	1.5353e+05\\
46000	1.5157e+05\\
48000	1.5256e+05\\
50000	1.5628e+05\\
52000	1.5079e+05\\
54000	1.5088e+05\\
56000	1.5506e+05\\
58000	1.5552e+05\\
60000	1.5378e+05\\
66000	1.5183e+05\\
68000	1.5189e+05\\
72000	1.5425e+05\\
74000	1.4933e+05\\
76000	1.5638e+05\\
78000	1.4781e+05\\
80000	1.4818e+05\\
82000	1.5213e+05\\
84000	1.5424e+05\\
86000	1.539e+05\\
88000	1.5124e+05\\
94000	1.5565e+05\\
98000	1.5025e+05\\
1e+05	1.5429e+05\\
};
\addlegendentry{Classe k=283}

\addplot [color=mycolor2, forget plot]
  table[row sep=crcr]{%
0	4395.3\\
2000	24354\\
4000	26108\\
6000	26206\\
8000	25994\\
12000	25929\\
14000	25734\\
16000	26139\\
18000	25884\\
20000	26761\\
22000	26896\\
24000	26076\\
26000	26352\\
32000	26400\\
34000	25846\\
36000	25569\\
38000	26413\\
40000	26377\\
42000	26869\\
44000	26788\\
46000	26105\\
48000	26091\\
50000	25717\\
52000	25793\\
54000	26103\\
58000	25371\\
60000	26115\\
62000	26045\\
64000	26582\\
66000	26778\\
68000	26414\\
70000	25799\\
72000	26201\\
74000	25649\\
76000	26023\\
78000	25831\\
82000	26483\\
84000	26184\\
86000	26057\\
88000	25557\\
90000	25624\\
92000	26167\\
94000	26380\\
96000	26116\\
98000	25648\\
1e+05	25762\\
};
\addplot [color=mycolor3, forget plot]
  table[row sep=crcr]{%
0	4395.3\\
2000	34821\\
4000	38531\\
6000	39739\\
8000	39932\\
10000	39505\\
12000	40152\\
14000	39604\\
16000	40088\\
18000	39446\\
20000	39280\\
22000	39954\\
24000	40934\\
26000	40402\\
28000	38691\\
30000	38612\\
32000	41830\\
34000	40601\\
36000	41056\\
38000	39518\\
40000	38936\\
42000	40697\\
44000	39975\\
46000	39883\\
48000	40045\\
50000	40442\\
52000	39104\\
54000	41143\\
58000	39434\\
60000	39563\\
62000	39519\\
64000	40868\\
66000	39959\\
68000	40496\\
70000	39790\\
72000	39883\\
74000	40598\\
76000	40716\\
78000	38529\\
80000	39099\\
82000	39280\\
84000	39931\\
86000	39838\\
88000	40297\\
90000	39390\\
92000	40019\\
94000	39054\\
96000	40072\\
98000	39894\\
1e+05	39393\\
};
\addplot [color=mycolor4, forget plot]
  table[row sep=crcr]{%
0	4395.3\\
2000	13810\\
4000	14563\\
6000	14187\\
8000	14503\\
10000	14279\\
12000	14490\\
14000	14456\\
16000	14666\\
18000	14495\\
20000	14058\\
24000	13910\\
26000	14553\\
28000	14316\\
30000	14576\\
32000	14661\\
34000	14402\\
36000	14041\\
38000	14193\\
40000	14616\\
42000	14146\\
44000	14104\\
46000	13827\\
48000	14473\\
50000	14578\\
52000	14184\\
54000	14061\\
56000	14343\\
58000	14806\\
60000	14708\\
62000	14263\\
64000	14530\\
66000	14606\\
68000	14391\\
70000	14726\\
72000	14112\\
74000	14111\\
76000	14191\\
78000	14534\\
80000	14276\\
82000	14431\\
84000	14074\\
86000	14611\\
88000	14281\\
90000	14479\\
92000	14019\\
94000	14025\\
96000	14505\\
98000	14732\\
1e+05	14088\\
};
\addplot [color=mycolor5, forget plot]
  table[row sep=crcr]{%
0	4395.3\\
2000	49689\\
4000	65114\\
6000	63310\\
8000	65597\\
10000	66625\\
12000	64030\\
14000	64132\\
16000	63649\\
18000	61824\\
20000	63566\\
22000	63787\\
24000	63508\\
30000	65848\\
32000	64205\\
34000	62369\\
36000	63736\\
38000	62835\\
40000	64878\\
42000	62291\\
44000	64696\\
46000	64995\\
48000	65785\\
50000	64600\\
52000	64787\\
54000	64576\\
56000	62860\\
58000	62280\\
60000	63923\\
62000	65005\\
64000	63129\\
66000	64433\\
68000	64533\\
70000	65090\\
72000	62835\\
74000	63512\\
76000	63320\\
78000	62482\\
80000	61973\\
84000	62360\\
86000	65127\\
88000	67474\\
90000	64351\\
92000	63271\\
94000	65763\\
96000	63969\\
98000	64244\\
1e+05	62498\\
};
\addplot [color=mycolor6, forget plot]
  table[row sep=crcr]{%
0	4395.3\\
2000	45283\\
4000	55511\\
6000	56229\\
8000	55458\\
10000	57487\\
12000	57445\\
14000	54946\\
16000	53567\\
18000	56417\\
20000	56414\\
22000	56189\\
24000	56893\\
26000	57057\\
28000	55434\\
30000	56922\\
32000	57469\\
34000	57516\\
36000	55941\\
38000	56060\\
40000	55665\\
42000	57222\\
44000	58083\\
46000	55760\\
48000	55566\\
52000	57802\\
54000	56357\\
56000	54587\\
58000	56306\\
60000	55682\\
62000	56910\\
64000	56245\\
66000	55945\\
68000	55290\\
70000	55481\\
72000	56162\\
74000	58968\\
76000	55541\\
78000	58417\\
80000	56983\\
82000	56397\\
84000	56716\\
86000	56512\\
88000	55872\\
90000	56781\\
92000	56041\\
94000	55490\\
96000	55841\\
98000	58902\\
1e+05	57708\\
};
\addplot [color=mycolor7, forget plot]
  table[row sep=crcr]{%
0	4395.3\\
2000	25228\\
4000	27486\\
6000	28763\\
8000	27383\\
10000	28268\\
12000	28711\\
14000	29323\\
16000	29524\\
18000	28804\\
20000	27786\\
22000	27754\\
24000	28662\\
26000	27690\\
28000	28137\\
30000	29026\\
32000	28653\\
34000	28011\\
36000	28184\\
38000	28073\\
40000	28762\\
42000	27562\\
44000	27546\\
46000	29066\\
48000	28567\\
50000	27002\\
52000	28579\\
54000	28786\\
56000	28350\\
58000	27828\\
60000	28746\\
62000	28122\\
64000	29169\\
66000	29019\\
68000	28246\\
70000	28129\\
72000	28237\\
74000	28600\\
76000	28612\\
78000	28434\\
80000	28505\\
82000	29138\\
84000	28843\\
86000	28762\\
88000	27955\\
90000	28337\\
92000	28134\\
94000	28843\\
96000	28042\\
98000	28316\\
1e+05	28963\\
};
\addplot [color=mycolor8, forget plot]
  table[row sep=crcr]{%
0	4395.3\\
2000	18364\\
4000	20347\\
6000	20995\\
8000	20595\\
10000	20636\\
12000	20385\\
14000	20967\\
16000	20782\\
18000	20774\\
20000	20163\\
22000	20658\\
24000	20667\\
26000	20220\\
28000	20287\\
30000	20689\\
32000	20495\\
34000	21018\\
36000	20568\\
38000	20721\\
40000	20596\\
42000	21227\\
44000	20112\\
46000	20809\\
48000	20619\\
50000	20080\\
52000	21021\\
54000	20729\\
56000	20715\\
58000	20846\\
60000	20713\\
62000	20858\\
64000	20731\\
66000	20890\\
68000	20207\\
70000	20498\\
72000	20518\\
74000	20653\\
76000	20983\\
78000	21106\\
80000	19964\\
84000	20430\\
86000	19927\\
90000	21092\\
92000	20534\\
94000	20243\\
96000	20258\\
98000	21358\\
1e+05	20585\\
};
\addplot [color=mycolor9, forget plot]
  table[row sep=crcr]{%
0	4395.3\\
2000	13242\\
4000	14364\\
8000	14398\\
10000	14251\\
12000	14363\\
14000	13790\\
16000	13917\\
18000	13834\\
20000	14008\\
22000	13689\\
24000	13904\\
26000	14082\\
28000	14304\\
30000	13954\\
32000	13712\\
34000	14253\\
36000	14330\\
38000	13851\\
40000	14132\\
42000	14367\\
44000	14013\\
46000	14174\\
48000	14520\\
50000	13806\\
52000	13876\\
54000	14196\\
56000	14468\\
58000	14212\\
62000	13999\\
64000	14505\\
66000	13804\\
68000	13742\\
70000	13774\\
72000	14327\\
74000	14347\\
76000	14013\\
78000	13582\\
80000	13997\\
82000	13568\\
84000	14121\\
86000	13927\\
88000	14163\\
90000	13848\\
92000	14296\\
94000	14181\\
96000	13838\\
98000	14083\\
1e+05	13758\\
};
\addplot [color=mycolor10, forget plot]
  table[row sep=crcr]{%
0	4395.3\\
2000	19429\\
4000	22677\\
6000	22542\\
8000	21529\\
10000	22261\\
12000	22567\\
14000	21938\\
16000	22076\\
18000	22666\\
20000	22104\\
22000	22623\\
24000	22010\\
26000	22470\\
28000	22261\\
30000	22393\\
32000	21574\\
34000	22289\\
36000	22274\\
38000	22058\\
40000	22080\\
42000	22403\\
44000	22338\\
46000	22554\\
48000	22027\\
52000	22439\\
54000	21840\\
56000	22486\\
58000	22407\\
60000	22456\\
62000	22418\\
64000	22520\\
66000	21975\\
68000	22080\\
70000	21978\\
72000	21251\\
74000	21555\\
76000	21343\\
78000	22359\\
80000	22425\\
82000	22757\\
84000	21888\\
86000	22230\\
88000	22246\\
90000	21439\\
92000	22372\\
94000	22060\\
96000	22494\\
1e+05	22288\\
};
\addplot [color=mycolor11, forget plot]
  table[row sep=crcr]{%
0	4395.3\\
2000	8218.2\\
4000	8623.2\\
6000	8813.1\\
8000	8696\\
10000	8719.3\\
12000	8492.4\\
14000	8812.2\\
16000	8417.4\\
18000	8846.9\\
20000	8658.5\\
22000	8889.8\\
24000	8780.5\\
26000	8791\\
28000	8910\\
30000	8653.7\\
32000	8453.1\\
36000	8678.1\\
38000	8518.4\\
40000	8329.5\\
42000	8556.7\\
44000	8886.1\\
46000	8643.4\\
48000	8348.6\\
50000	8418.4\\
52000	8348\\
54000	8685.1\\
56000	8798.6\\
58000	8833.5\\
60000	8764.4\\
62000	8235.8\\
64000	8515.1\\
66000	8706.6\\
68000	8826.2\\
70000	8712.7\\
72000	8532.3\\
74000	8804.2\\
78000	8559.8\\
80000	8746.7\\
82000	8443.9\\
84000	8698.4\\
86000	8434.5\\
88000	8607\\
90000	8574.5\\
92000	8509.8\\
94000	8735.7\\
96000	8751.6\\
98000	8550.9\\
1e+05	8643.9\\
};
\addplot [color=mycolor12, forget plot]
  table[row sep=crcr]{%
0	4395.3\\
2000	13621\\
4000	15050\\
6000	14460\\
8000	14481\\
10000	14560\\
12000	15049\\
14000	14741\\
16000	14757\\
18000	14910\\
20000	14641\\
22000	14692\\
24000	14196\\
26000	15031\\
28000	15082\\
30000	14396\\
32000	14751\\
34000	14564\\
36000	14794\\
38000	14947\\
40000	14918\\
42000	15054\\
44000	14665\\
46000	14208\\
48000	14291\\
50000	14926\\
52000	14982\\
54000	14612\\
56000	14702\\
58000	14702\\
60000	14411\\
62000	14472\\
64000	14692\\
66000	14267\\
68000	15107\\
70000	14865\\
74000	15081\\
76000	14502\\
78000	14604\\
80000	15172\\
82000	15009\\
84000	14655\\
86000	15046\\
88000	14959\\
90000	14467\\
92000	14892\\
94000	14542\\
96000	14713\\
98000	15062\\
1e+05	14669\\
};
\addplot [color=mycolor13, forget plot]
  table[row sep=crcr]{%
0	4395.3\\
2000	10523\\
4000	11508\\
6000	10768\\
8000	11243\\
10000	10747\\
14000	11146\\
16000	11164\\
18000	11253\\
20000	10899\\
22000	11119\\
24000	10887\\
26000	10613\\
28000	11133\\
30000	11182\\
32000	11312\\
34000	11056\\
36000	10841\\
38000	11045\\
40000	10821\\
42000	10834\\
44000	11615\\
46000	11586\\
48000	10935\\
50000	11190\\
52000	11295\\
54000	11151\\
56000	10839\\
58000	11097\\
60000	11141\\
62000	11444\\
64000	11026\\
66000	11076\\
68000	11216\\
70000	10940\\
72000	10992\\
74000	11169\\
76000	10707\\
78000	11229\\
80000	11622\\
82000	11501\\
84000	10937\\
86000	11182\\
88000	11006\\
90000	11453\\
92000	11116\\
1e+05	11518\\
};
\addplot [color=mycolor14, forget plot]
  table[row sep=crcr]{%
0	4395.3\\
2000	20851\\
4000	24673\\
6000	24737\\
8000	24386\\
10000	24813\\
12000	24909\\
14000	24921\\
16000	24301\\
18000	24905\\
20000	24124\\
22000	23610\\
24000	24068\\
26000	24790\\
28000	24220\\
32000	25209\\
34000	25271\\
36000	24166\\
38000	24615\\
40000	24634\\
42000	23606\\
48000	25207\\
50000	23383\\
52000	24542\\
54000	23898\\
56000	24202\\
60000	23515\\
62000	24496\\
64000	23958\\
66000	23614\\
68000	25007\\
70000	24380\\
72000	25212\\
74000	23790\\
76000	24717\\
78000	25020\\
80000	24296\\
82000	24902\\
84000	24582\\
86000	24116\\
88000	23875\\
90000	24179\\
94000	24074\\
96000	23662\\
98000	24493\\
1e+05	24625\\
};
\addplot [color=mycolor15, forget plot]
  table[row sep=crcr]{%
0	4395.3\\
2000	15594\\
4000	17468\\
6000	16600\\
8000	16745\\
12000	16306\\
14000	17062\\
16000	16522\\
18000	16898\\
20000	17583\\
22000	17497\\
26000	16977\\
28000	16978\\
30000	17156\\
32000	17209\\
34000	16678\\
36000	17029\\
38000	16663\\
40000	16784\\
42000	16496\\
44000	16933\\
46000	16726\\
48000	16318\\
50000	16742\\
52000	17298\\
54000	17298\\
56000	16854\\
60000	16948\\
62000	16660\\
64000	17033\\
66000	16829\\
70000	16798\\
72000	16462\\
74000	17358\\
76000	16349\\
78000	17033\\
80000	17476\\
82000	16906\\
84000	16650\\
86000	17219\\
88000	16914\\
90000	16834\\
92000	16894\\
94000	16308\\
96000	17356\\
98000	17130\\
1e+05	17318\\
};
\addplot [color=mycolor16, forget plot]
  table[row sep=crcr]{%
0	4395.3\\
2000	7002.3\\
4000	6890.5\\
6000	6989.9\\
8000	6948.9\\
10000	7214.3\\
12000	7033.1\\
14000	6878.9\\
16000	7124.6\\
18000	7147.8\\
22000	7104.4\\
24000	7099.3\\
26000	6719.6\\
28000	6921.9\\
30000	7067.6\\
32000	7157.2\\
34000	7045.3\\
36000	6903.1\\
38000	7078.8\\
40000	6979.6\\
42000	7170.6\\
44000	6911.6\\
46000	7042.1\\
48000	6981.4\\
50000	6985.2\\
52000	7140.9\\
54000	7055.2\\
58000	7143\\
60000	6945.4\\
62000	6985.9\\
64000	7139.9\\
66000	7012.4\\
68000	7136.8\\
70000	7127.2\\
72000	7026.3\\
74000	6733.3\\
76000	7010.7\\
78000	7099.5\\
80000	7263.4\\
82000	7350.4\\
84000	6904.1\\
86000	7180.8\\
88000	6888.7\\
90000	6990.8\\
92000	7008\\
94000	7081.4\\
96000	6996.2\\
98000	6784.3\\
1e+05	6666.7\\
};
\addplot [color=mycolor17, forget plot]
  table[row sep=crcr]{%
0	4395.3\\
2000	14760\\
4000	16195\\
6000	16286\\
8000	16167\\
10000	15818\\
12000	16110\\
14000	16010\\
16000	15798\\
20000	16282\\
22000	15738\\
24000	16233\\
26000	15582\\
30000	16132\\
32000	16771\\
34000	16389\\
36000	15597\\
38000	16261\\
40000	16173\\
42000	15959\\
44000	16305\\
46000	16264\\
48000	16510\\
50000	15978\\
52000	15717\\
54000	16230\\
56000	16490\\
58000	15726\\
60000	16393\\
62000	16271\\
64000	15418\\
66000	15982\\
68000	16215\\
70000	16627\\
72000	16573\\
74000	15900\\
76000	16315\\
80000	15961\\
82000	16140\\
84000	15567\\
86000	15984\\
88000	16253\\
90000	16119\\
92000	16458\\
94000	15754\\
96000	15719\\
98000	15819\\
1e+05	16282\\
};
\addplot [color=mycolor18, forget plot]
  table[row sep=crcr]{%
0	4395.3\\
2000	22169\\
4000	24896\\
6000	24458\\
8000	24824\\
10000	24609\\
12000	24613\\
14000	24791\\
16000	24269\\
18000	24754\\
20000	24294\\
22000	24766\\
24000	24433\\
28000	24136\\
30000	24163\\
32000	23826\\
36000	24645\\
38000	24384\\
40000	24521\\
42000	25351\\
44000	24360\\
46000	24145\\
48000	24679\\
50000	24728\\
52000	25030\\
54000	24668\\
56000	24733\\
58000	25217\\
60000	24062\\
62000	24199\\
64000	23971\\
66000	24461\\
68000	24799\\
70000	24778\\
72000	23962\\
74000	24412\\
76000	24463\\
78000	24289\\
80000	24624\\
82000	24778\\
84000	24150\\
86000	24252\\
90000	24747\\
92000	24276\\
94000	25219\\
96000	24728\\
98000	24471\\
1e+05	24827\\
};
\addplot [color=mycolor19, forget plot]
  table[row sep=crcr]{%
0	4395.3\\
2000	14347\\
4000	15018\\
6000	14987\\
8000	15098\\
10000	14686\\
12000	15071\\
14000	14863\\
16000	14501\\
18000	15439\\
22000	14821\\
24000	14709\\
26000	14772\\
28000	14656\\
30000	14664\\
32000	14937\\
34000	14721\\
36000	14774\\
38000	15017\\
40000	14598\\
42000	14847\\
44000	14511\\
46000	14699\\
48000	15208\\
50000	15480\\
52000	15108\\
54000	15045\\
56000	14361\\
58000	15131\\
60000	15259\\
62000	15132\\
64000	15299\\
68000	14913\\
70000	14323\\
72000	15156\\
74000	15118\\
76000	14620\\
78000	15592\\
80000	14928\\
82000	14930\\
84000	14873\\
86000	14860\\
88000	15132\\
90000	15050\\
92000	15417\\
94000	14782\\
96000	15010\\
98000	14558\\
1e+05	15349\\
};
\addplot [color=mycolor20, forget plot]
  table[row sep=crcr]{%
0	4395.3\\
2000	17345\\
4000	18422\\
6000	19407\\
10000	19718\\
12000	18927\\
14000	18812\\
16000	18902\\
18000	18776\\
20000	19729\\
22000	19967\\
24000	19704\\
26000	18623\\
28000	20136\\
30000	19008\\
32000	19814\\
34000	19296\\
36000	19298\\
38000	18591\\
40000	19638\\
42000	18948\\
44000	18964\\
46000	18811\\
48000	18386\\
50000	19372\\
52000	18990\\
54000	19327\\
56000	19227\\
58000	19483\\
62000	19435\\
64000	19244\\
66000	18691\\
68000	19869\\
70000	20024\\
72000	19619\\
74000	19638\\
76000	19264\\
78000	19351\\
80000	18825\\
82000	19239\\
84000	19163\\
86000	19628\\
88000	19164\\
90000	19756\\
92000	19386\\
94000	19222\\
96000	19546\\
98000	19816\\
1e+05	19640\\
};
\addplot [color=mycolor21, forget plot]
  table[row sep=crcr]{%
0	4395.3\\
2000	9217.9\\
4000	9782.6\\
6000	9851\\
8000	9524.9\\
10000	9462.5\\
12000	10057\\
14000	9648.4\\
16000	9695.5\\
18000	9556.3\\
20000	9578\\
22000	9790.8\\
24000	9775.4\\
26000	9853.6\\
28000	9690.6\\
30000	9587.4\\
32000	9737.9\\
34000	9801.9\\
36000	9918.2\\
38000	9848.9\\
40000	9739.1\\
42000	9711.3\\
44000	9812.4\\
46000	9715.9\\
48000	9944.5\\
50000	9668\\
52000	9722.3\\
54000	9578.4\\
56000	9824.3\\
58000	9563.9\\
60000	9860.9\\
62000	9667.5\\
64000	9594.6\\
66000	9692.1\\
68000	9705.1\\
72000	9889.4\\
74000	9475.1\\
76000	9425.8\\
78000	9825.1\\
80000	9918\\
82000	9455.6\\
84000	9694.4\\
86000	9704.9\\
88000	9657.1\\
90000	9880.6\\
92000	9539.2\\
94000	9833.3\\
96000	9839.8\\
98000	9479.8\\
1e+05	9964\\
};
\addplot [color=mycolor22, forget plot]
  table[row sep=crcr]{%
0	4395.3\\
2000	21644\\
4000	25161\\
8000	25542\\
10000	25446\\
12000	24798\\
14000	24834\\
16000	24446\\
18000	24962\\
20000	26052\\
22000	25082\\
24000	25632\\
26000	24786\\
28000	26142\\
30000	24501\\
32000	24785\\
34000	25564\\
36000	25424\\
38000	26121\\
40000	24955\\
42000	24976\\
44000	26443\\
46000	25478\\
48000	25116\\
50000	24979\\
52000	25361\\
54000	24811\\
56000	25381\\
58000	25539\\
60000	25425\\
62000	25729\\
64000	25099\\
66000	25588\\
68000	24369\\
70000	25591\\
72000	24487\\
74000	25814\\
76000	25197\\
78000	25049\\
80000	25195\\
82000	25821\\
84000	25994\\
86000	26024\\
88000	26414\\
90000	25575\\
92000	25236\\
94000	26002\\
96000	25675\\
98000	25187\\
1e+05	24901\\
};
\addplot [color=mycolor23, forget plot]
  table[row sep=crcr]{%
0	4395.3\\
2000	11460\\
4000	12160\\
6000	12454\\
8000	12199\\
10000	12216\\
12000	12366\\
14000	12738\\
16000	12284\\
18000	12118\\
20000	12234\\
22000	12116\\
24000	12205\\
26000	11948\\
28000	12291\\
30000	12606\\
32000	12394\\
34000	12473\\
36000	12209\\
38000	12358\\
40000	12239\\
42000	12365\\
44000	12294\\
46000	12608\\
48000	12200\\
50000	12025\\
52000	12036\\
54000	12274\\
58000	12326\\
60000	12462\\
62000	12159\\
64000	12397\\
66000	12206\\
68000	12265\\
70000	12192\\
72000	12277\\
74000	12460\\
76000	12383\\
80000	12334\\
82000	12423\\
86000	12285\\
88000	12343\\
90000	12762\\
92000	12354\\
94000	12213\\
96000	12239\\
98000	12218\\
1e+05	12290\\
};
\addplot [color=mycolor24, forget plot]
  table[row sep=crcr]{%
0	4395.3\\
2000	4974.4\\
4000	4723.7\\
6000	4654\\
8000	4836.3\\
10000	4717.5\\
12000	4842.5\\
14000	4942.4\\
16000	4627.7\\
18000	4748.8\\
20000	4655.7\\
22000	4649.5\\
24000	4774.1\\
26000	4963.5\\
28000	4741.6\\
30000	4557.1\\
34000	4765\\
38000	4760.5\\
40000	4695.5\\
42000	4696\\
44000	4730.7\\
46000	4910.8\\
48000	4774.1\\
50000	4719.9\\
52000	4824.5\\
54000	4701.4\\
56000	4913.3\\
58000	4768.5\\
60000	4893\\
62000	4728.6\\
64000	4903.8\\
68000	4733\\
70000	4845.3\\
72000	4800.5\\
74000	4901.6\\
76000	4885.1\\
78000	4903.6\\
80000	4772.5\\
82000	4928.5\\
84000	5020.8\\
86000	4626.1\\
88000	4861.9\\
90000	4818.2\\
94000	4891.6\\
96000	4781\\
98000	4765.5\\
1e+05	4886.2\\
};
\addplot [color=mycolor25, forget plot]
  table[row sep=crcr]{%
0	4395.3\\
2000	11576\\
4000	12258\\
6000	12640\\
8000	12235\\
10000	12468\\
12000	12106\\
14000	12668\\
16000	12209\\
18000	12280\\
20000	12502\\
22000	12624\\
24000	12262\\
26000	12621\\
28000	12455\\
30000	12837\\
32000	12257\\
34000	12251\\
36000	12587\\
38000	12354\\
40000	12777\\
42000	12024\\
44000	12573\\
46000	12817\\
48000	12561\\
50000	12393\\
52000	12823\\
54000	11919\\
56000	12887\\
58000	12379\\
60000	12004\\
62000	12323\\
64000	12712\\
66000	12605\\
68000	12254\\
70000	12538\\
72000	12157\\
74000	12661\\
76000	12396\\
78000	12332\\
80000	12517\\
82000	12233\\
84000	12457\\
88000	12196\\
90000	12381\\
92000	12832\\
94000	12229\\
96000	13085\\
98000	12299\\
1e+05	12272\\
};
\addplot [color=mycolor26, forget plot]
  table[row sep=crcr]{%
0	4395.3\\
2000	8071.9\\
4000	8257.7\\
6000	8119.7\\
10000	8442.4\\
12000	8267.8\\
14000	8523.6\\
16000	8022.6\\
18000	8244.6\\
20000	8623.3\\
22000	8466.4\\
24000	7914.8\\
26000	8143.3\\
28000	8261\\
30000	8475.2\\
32000	8203.6\\
34000	8124.6\\
36000	8153.4\\
38000	8460.2\\
40000	8098\\
42000	8194.6\\
44000	8020.9\\
46000	8428.8\\
48000	8115.5\\
50000	8515.8\\
52000	8055.7\\
54000	8410.5\\
56000	8373.9\\
58000	8188.8\\
60000	8187.7\\
62000	8463.7\\
64000	7786.3\\
66000	8212.9\\
68000	8103.4\\
70000	8454.9\\
72000	8192.1\\
74000	8201.6\\
76000	8125.3\\
78000	8004.2\\
80000	8020.3\\
82000	8111.8\\
84000	8233.3\\
86000	8508\\
88000	8438.7\\
90000	7950.5\\
92000	8220.4\\
94000	8415.2\\
96000	8446.3\\
98000	8137.4\\
1e+05	8302.7\\
};
\addplot [color=mycolor27, forget plot]
  table[row sep=crcr]{%
0	4395.3\\
2000	9491.7\\
4000	9415.1\\
6000	9494.5\\
8000	9671.5\\
10000	9552.9\\
12000	10145\\
14000	9601.7\\
16000	9614.2\\
18000	9595.7\\
22000	9906.1\\
24000	10014\\
26000	9996.3\\
28000	9889.3\\
30000	9703.2\\
32000	9710.6\\
34000	9877.3\\
36000	9651.8\\
38000	9733.4\\
40000	9577.4\\
42000	9327.1\\
44000	9671.9\\
46000	9715.2\\
48000	9435.4\\
50000	9646\\
52000	9817.6\\
54000	10213\\
56000	9750.6\\
58000	9757.2\\
60000	9487.8\\
62000	9973.8\\
64000	9701.1\\
66000	9393.5\\
68000	9606\\
70000	9734.1\\
72000	9478.9\\
74000	9436.2\\
76000	9847.5\\
80000	9941.2\\
82000	9917.3\\
84000	9774.8\\
86000	9678.8\\
88000	9947.8\\
90000	9790.3\\
92000	9798.7\\
94000	9641.2\\
96000	9835.9\\
98000	9882.3\\
1e+05	9672.7\\
};
\addplot [color=mycolor28, forget plot]
  table[row sep=crcr]{%
0	4395.3\\
2000	5585.8\\
4000	5524.8\\
6000	5445.1\\
10000	5395.9\\
12000	5482.3\\
14000	5592.8\\
16000	5395.4\\
18000	5415.2\\
20000	5452\\
22000	5471.1\\
24000	5391.5\\
26000	5102.7\\
28000	5403.3\\
30000	5684.1\\
32000	5388.1\\
34000	5432.2\\
36000	5519.5\\
38000	5588.3\\
40000	5426.4\\
42000	5440\\
44000	5347.8\\
46000	5544.5\\
48000	5464.5\\
50000	5368.6\\
52000	5329\\
54000	5351.8\\
56000	5568.3\\
58000	5294.6\\
60000	5507.7\\
62000	5256.8\\
64000	5467.1\\
66000	5437.5\\
68000	5307.5\\
70000	5474.7\\
72000	5392.1\\
74000	5567.3\\
76000	5472.7\\
78000	5458.1\\
80000	5334.3\\
82000	5436.2\\
84000	5515.9\\
86000	5282.6\\
88000	5401.3\\
90000	5419\\
92000	5566.3\\
94000	5278.9\\
96000	5540.4\\
98000	5352.3\\
1e+05	5542.3\\
};
\addplot [color=mycolor29, forget plot]
  table[row sep=crcr]{%
0	4395.3\\
2000	10735\\
4000	11789\\
6000	11491\\
8000	11960\\
10000	11539\\
12000	11388\\
14000	11365\\
16000	11841\\
18000	11605\\
20000	11182\\
22000	11240\\
24000	11966\\
26000	11474\\
28000	11590\\
30000	11413\\
32000	11651\\
34000	11675\\
36000	11851\\
38000	11398\\
40000	11834\\
42000	11778\\
44000	11685\\
46000	11641\\
48000	10930\\
50000	11372\\
52000	11330\\
54000	11907\\
56000	11638\\
58000	11456\\
60000	11885\\
62000	11715\\
66000	11862\\
68000	11480\\
70000	11527\\
72000	11517\\
74000	11814\\
76000	11850\\
78000	11468\\
82000	11372\\
84000	11783\\
86000	11264\\
88000	11327\\
90000	11629\\
92000	11403\\
94000	11446\\
96000	11697\\
98000	11662\\
1e+05	11522\\
};
\addplot [color=mycolor30, forget plot]
  table[row sep=crcr]{%
0	4395.3\\
2000	7679\\
4000	7836.9\\
6000	7718.3\\
8000	7882.4\\
10000	7901.2\\
12000	7821.3\\
14000	7873.6\\
16000	7889.9\\
18000	7853.8\\
20000	7790.2\\
22000	7947.8\\
24000	7987\\
26000	7854.6\\
28000	7956.2\\
30000	7902.2\\
32000	7979.4\\
38000	7806.9\\
40000	7711.9\\
42000	8285.3\\
44000	7890.8\\
46000	7809.9\\
48000	7803.7\\
50000	7981.2\\
52000	8004.8\\
54000	8060.6\\
56000	7884.7\\
58000	7988.6\\
60000	7792.4\\
62000	7833.8\\
64000	8069.9\\
66000	8035.2\\
68000	7962.4\\
70000	7720.2\\
72000	7874.9\\
74000	7733.7\\
76000	7999.7\\
78000	7930.9\\
80000	7995.6\\
82000	7881.6\\
84000	7932.9\\
86000	7907.5\\
88000	7935.8\\
90000	7838.7\\
92000	7595.4\\
94000	7996.6\\
96000	7590.3\\
98000	7781.3\\
1e+05	7930.3\\
};
\addplot [color=mycolor31, forget plot]
  table[row sep=crcr]{%
0	4395.3\\
2000	6294.1\\
4000	6820.8\\
6000	6349.1\\
8000	6360.5\\
10000	6503.9\\
12000	6368.6\\
14000	6430.1\\
18000	6389.1\\
20000	6497.7\\
22000	6548.4\\
26000	6339\\
30000	6545.8\\
32000	6586\\
34000	6259\\
36000	6609.3\\
38000	6376.7\\
40000	6426.1\\
44000	6452.4\\
46000	6206\\
48000	6439.7\\
50000	6394\\
52000	6321.4\\
54000	6340.6\\
56000	6462.6\\
58000	6918.4\\
60000	6440.8\\
62000	6588.3\\
64000	6320.3\\
66000	6318.6\\
68000	6443.8\\
70000	6470.2\\
72000	6427.6\\
74000	6595.8\\
76000	6286.9\\
78000	6646.6\\
80000	6436.9\\
82000	6627\\
84000	6512.7\\
86000	6645.7\\
88000	6341.5\\
90000	6345.5\\
92000	6509.7\\
94000	6353.8\\
96000	6448.7\\
98000	6492.7\\
1e+05	6405.5\\
};
\addplot [color=mycolor32, forget plot]
  table[row sep=crcr]{%
0	4395.3\\
2000	9922.4\\
4000	10441\\
6000	10826\\
8000	10207\\
10000	10102\\
12000	10293\\
14000	10647\\
16000	10880\\
18000	10423\\
20000	10429\\
22000	10729\\
24000	10646\\
26000	10451\\
28000	10221\\
30000	10764\\
32000	10115\\
34000	10761\\
36000	10412\\
38000	10513\\
40000	10244\\
42000	10592\\
44000	10551\\
46000	10479\\
48000	10873\\
50000	10334\\
52000	10723\\
54000	10545\\
56000	10853\\
58000	10223\\
60000	10207\\
64000	10622\\
66000	10709\\
68000	10520\\
70000	10420\\
72000	10445\\
74000	10543\\
76000	10088\\
78000	10617\\
80000	10819\\
82000	10328\\
84000	10609\\
86000	10371\\
88000	10364\\
90000	10324\\
92000	10341\\
94000	10572\\
96000	10424\\
98000	10470\\
1e+05	10671\\
};
\addplot [color=mycolor33, forget plot]
  table[row sep=crcr]{%
0	4395.3\\
2000	7691.4\\
4000	8423.7\\
6000	8191.8\\
8000	8157\\
10000	8156.4\\
12000	8185.4\\
14000	8338.1\\
16000	8204.9\\
18000	8211.4\\
20000	8248\\
22000	7923.6\\
24000	8427.9\\
26000	8544.5\\
28000	8436.8\\
30000	8248.9\\
32000	7877.2\\
34000	8573.4\\
38000	7862.3\\
42000	8410\\
44000	7968.6\\
46000	8110.9\\
48000	8188.9\\
50000	8354.9\\
52000	8372.9\\
54000	8144.4\\
56000	8086.8\\
58000	8152.4\\
60000	8269.1\\
62000	8270.8\\
64000	8110.4\\
66000	8197.5\\
68000	8384.9\\
70000	8308.9\\
72000	8158.2\\
74000	8141.9\\
76000	8313\\
78000	8281.9\\
80000	8182\\
82000	8155.1\\
84000	8391.1\\
86000	8268\\
88000	8304.4\\
90000	8292.8\\
92000	8071\\
94000	8061.3\\
96000	8259.2\\
98000	8386.3\\
1e+05	8359.9\\
};
\addplot [color=mycolor34, forget plot]
  table[row sep=crcr]{%
0	4395.3\\
2000	4184\\
4000	4059.5\\
6000	3994\\
8000	4112.8\\
10000	4069\\
12000	3934.8\\
14000	3953.3\\
16000	3853.3\\
18000	3958.5\\
20000	4081.4\\
22000	3986.1\\
24000	4038\\
26000	3989.5\\
28000	4002.3\\
30000	4068.5\\
32000	3816.2\\
34000	4052.9\\
36000	4012.6\\
38000	4035\\
40000	4043.8\\
42000	4124.1\\
44000	4061.6\\
48000	4069.6\\
50000	4049.3\\
52000	3918.8\\
54000	4168.6\\
56000	3918.8\\
58000	4039.9\\
60000	3978.5\\
62000	4124.5\\
64000	4031.2\\
66000	3962.8\\
68000	3979.2\\
70000	4065.6\\
72000	3984.8\\
74000	4188.9\\
76000	3930.8\\
78000	4043.9\\
80000	3988.6\\
82000	3962.7\\
84000	4124.3\\
86000	3924.4\\
88000	3996.8\\
90000	4026.1\\
92000	4008.3\\
94000	3909.7\\
96000	3979.9\\
1e+05	3997.5\\
};
\addplot [color=mycolor35]
  table[row sep=crcr]{%
0	4395.3\\
2000	9665.7\\
4000	10034\\
6000	10011\\
8000	10563\\
10000	10401\\
12000	10153\\
14000	10269\\
16000	10285\\
18000	10081\\
20000	10246\\
24000	10302\\
26000	10297\\
28000	10054\\
30000	9788.4\\
32000	10045\\
34000	10127\\
36000	10180\\
38000	10295\\
40000	10126\\
42000	10189\\
48000	10120\\
50000	10226\\
52000	10241\\
54000	10168\\
56000	10352\\
60000	10312\\
62000	10187\\
64000	9975.3\\
66000	10485\\
68000	10027\\
70000	10181\\
72000	10417\\
74000	9955\\
76000	10373\\
78000	10229\\
80000	10560\\
84000	10309\\
86000	10099\\
90000	10135\\
92000	10318\\
94000	10015\\
96000	10595\\
98000	10458\\
1e+05	10148\\
};
\addlegendentry{Classe k=83}

\addplot [color=mycolor36, forget plot]
  table[row sep=crcr]{%
0	4395.3\\
2000	4380.1\\
4000	4177.7\\
6000	4304.8\\
8000	4229.4\\
10000	4222.6\\
12000	4177.2\\
14000	4252.9\\
16000	4250\\
18000	4407.7\\
20000	4297.5\\
22000	4233.7\\
26000	4218.4\\
28000	4488.5\\
30000	4253.1\\
32000	4275.5\\
34000	4199.9\\
36000	4082.4\\
38000	4202.8\\
40000	4268.6\\
42000	4298\\
44000	4137.3\\
46000	4310.1\\
48000	4228.5\\
50000	4309.7\\
52000	4370.8\\
54000	4387.3\\
56000	4152.2\\
58000	4251.7\\
60000	4369.2\\
62000	4221.8\\
66000	4349.8\\
68000	4338.9\\
70000	4237.6\\
72000	4297.1\\
74000	4316.1\\
76000	4248.8\\
78000	4403.3\\
80000	4210.2\\
82000	4367.3\\
84000	4257.2\\
86000	4001.5\\
88000	4116.3\\
90000	4491.5\\
92000	4356.8\\
94000	4313.5\\
96000	4176.9\\
98000	4296.9\\
1e+05	4192.4\\
};
\addplot [color=mycolor37, forget plot]
  table[row sep=crcr]{%
0	4395.3\\
2000	3799.5\\
4000	3661.9\\
6000	3609.9\\
8000	3665.3\\
10000	3576.3\\
12000	3602.1\\
14000	3584\\
16000	3728.5\\
18000	3579.6\\
20000	3776.7\\
22000	3611.2\\
24000	3593.9\\
26000	3601.4\\
28000	3571.2\\
30000	3645.2\\
32000	3567.4\\
34000	3645.9\\
36000	3565.5\\
38000	3671.3\\
40000	3647\\
42000	3530.4\\
44000	3636.4\\
46000	3611\\
48000	3570.3\\
50000	3609.1\\
52000	3523\\
54000	3592.9\\
56000	3624.4\\
58000	3523.2\\
60000	3598.7\\
62000	3565.6\\
64000	3562.4\\
66000	3682.5\\
68000	3582.2\\
70000	3624.2\\
72000	3543.4\\
74000	3565.5\\
76000	3628.1\\
78000	3556.6\\
80000	3609.3\\
82000	3592.3\\
84000	3687.8\\
86000	3640.4\\
88000	3489.6\\
90000	3625\\
92000	3642.7\\
94000	3584.1\\
96000	3546.8\\
1e+05	3508.2\\
};
\addplot [color=mycolor38, forget plot]
  table[row sep=crcr]{%
0	4395.3\\
2000	3176.3\\
4000	3014.5\\
8000	3003.2\\
10000	2947.7\\
12000	3001.8\\
14000	3047.7\\
16000	2941.2\\
18000	3017.3\\
20000	2981.5\\
22000	2869.3\\
24000	2881.2\\
26000	2955.1\\
28000	2953.8\\
30000	3019.9\\
32000	3052.8\\
34000	2996.6\\
38000	3090.6\\
40000	3011.9\\
42000	2871.6\\
44000	3008.2\\
46000	2919.7\\
48000	3021.8\\
50000	2988.3\\
54000	2944.4\\
56000	3061.8\\
58000	3033.5\\
60000	2962.7\\
62000	3034.7\\
64000	3012.7\\
66000	3128\\
68000	2963.4\\
70000	3025\\
72000	2960.3\\
74000	2968.8\\
76000	2991.9\\
78000	2989.9\\
80000	2962.9\\
82000	2964.8\\
84000	3076.6\\
86000	2918.1\\
88000	3024.5\\
90000	3065\\
92000	2938.4\\
94000	2954.1\\
96000	3011.2\\
98000	2960.1\\
1e+05	2802.6\\
};
\addplot [color=mycolor39, forget plot]
  table[row sep=crcr]{%
0	4395.3\\
2000	11050\\
4000	11810\\
6000	11532\\
8000	12254\\
10000	11821\\
12000	12211\\
14000	11932\\
16000	11619\\
18000	11990\\
20000	12175\\
22000	11406\\
24000	11852\\
26000	11994\\
30000	11897\\
32000	11600\\
34000	11743\\
36000	12066\\
38000	11915\\
40000	11711\\
42000	11793\\
44000	12490\\
46000	11664\\
48000	11604\\
50000	11656\\
52000	12090\\
54000	12222\\
56000	11875\\
58000	11897\\
60000	11726\\
62000	11657\\
64000	11758\\
68000	11656\\
70000	11673\\
72000	12015\\
74000	12227\\
76000	11863\\
78000	11982\\
80000	11852\\
82000	11629\\
84000	11814\\
86000	11694\\
90000	11920\\
92000	11481\\
94000	11677\\
96000	12011\\
98000	11622\\
1e+05	11833\\
};
\addplot [color=mycolor40, forget plot]
  table[row sep=crcr]{%
0	4395.3\\
2000	5835.3\\
4000	5791.6\\
6000	5704.3\\
8000	5720.5\\
10000	5784.2\\
14000	5729.9\\
16000	5707.5\\
18000	5873\\
20000	5683.8\\
22000	5836.9\\
24000	5665.2\\
26000	5743.1\\
28000	5794.4\\
30000	5781.4\\
32000	5703.9\\
34000	5759\\
36000	5848.3\\
38000	5893.7\\
40000	5649.9\\
42000	5817.9\\
44000	5734.1\\
46000	5816.5\\
48000	5784.2\\
50000	5673.8\\
52000	5847\\
54000	5720.6\\
56000	5648.2\\
58000	5755.8\\
60000	5625.7\\
62000	5783.8\\
64000	5737.2\\
66000	5612.2\\
68000	5952.4\\
70000	5817.6\\
72000	5647.5\\
74000	5836.9\\
76000	5733.8\\
78000	5864.6\\
80000	5891.4\\
82000	5781.2\\
84000	5722\\
86000	5758.2\\
88000	5592.2\\
90000	5789.9\\
92000	5678.3\\
94000	5732.4\\
96000	5853.5\\
98000	5735.2\\
1e+05	5641\\
};
\addplot [color=mycolor41, forget plot]
  table[row sep=crcr]{%
0	4395.3\\
2000	5538.3\\
4000	5315.2\\
6000	5331.5\\
8000	5366.6\\
12000	5337.4\\
14000	5474.5\\
16000	5376.9\\
18000	5479.7\\
20000	5189.4\\
22000	5223.4\\
24000	5346\\
26000	5436.2\\
28000	5371\\
30000	5382\\
32000	5363.5\\
34000	5286\\
36000	5411.7\\
38000	5455.3\\
40000	5307.9\\
42000	5282.6\\
44000	5342.2\\
46000	5505.1\\
48000	5443.4\\
52000	5460\\
54000	5388.2\\
56000	5228.7\\
58000	5254.3\\
60000	5520.5\\
62000	5474.8\\
64000	5361.7\\
66000	5526\\
68000	5431.8\\
70000	5510\\
74000	5334.9\\
76000	5333.7\\
78000	5382.3\\
80000	5367.8\\
82000	5399.4\\
84000	5344.7\\
86000	5151.9\\
88000	5339.6\\
90000	5491.8\\
92000	5411.3\\
94000	5260\\
96000	5293.6\\
98000	5299.8\\
1e+05	5391.8\\
};
\addplot [color=mycolor42, forget plot]
  table[row sep=crcr]{%
0	4395.3\\
2000	17723\\
4000	20261\\
6000	20696\\
8000	20732\\
10000	21199\\
12000	20856\\
16000	20785\\
18000	20297\\
20000	19485\\
22000	20546\\
24000	19738\\
26000	19958\\
28000	20387\\
30000	20478\\
32000	21066\\
34000	20743\\
36000	20587\\
38000	19661\\
40000	20521\\
42000	20027\\
44000	20418\\
46000	20192\\
48000	21614\\
50000	21734\\
52000	21075\\
54000	20776\\
56000	20083\\
58000	20202\\
60000	21107\\
62000	20510\\
64000	20679\\
66000	21091\\
68000	20721\\
70000	21134\\
72000	20270\\
74000	21938\\
76000	21301\\
80000	21530\\
82000	19885\\
84000	20803\\
86000	21123\\
88000	20512\\
90000	20889\\
92000	20275\\
94000	21461\\
96000	20757\\
98000	21023\\
1e+05	20657\\
};
\addplot [color=mycolor43, forget plot]
  table[row sep=crcr]{%
0	4395.3\\
2000	5014.4\\
4000	4906.9\\
6000	4876.1\\
8000	4937.2\\
10000	4782.9\\
12000	4815.6\\
14000	5005.6\\
16000	4896.3\\
18000	4856.2\\
20000	4853.2\\
22000	4802.4\\
24000	5011.6\\
26000	4846\\
28000	4874.6\\
30000	5072.9\\
32000	4912.5\\
34000	4898.7\\
36000	4903.5\\
40000	4852\\
42000	4937\\
44000	4996\\
50000	4803.2\\
52000	4912.1\\
54000	4927.9\\
56000	4865\\
58000	4874.5\\
62000	4956\\
64000	4918.4\\
66000	4911.1\\
68000	5000.2\\
70000	4896.8\\
72000	5005.4\\
74000	4944\\
76000	4943\\
78000	4919.9\\
80000	4867.6\\
82000	4764.7\\
86000	4900.3\\
88000	4924.1\\
90000	4977.2\\
94000	4865.5\\
96000	4968.5\\
98000	4897.7\\
1e+05	4866.7\\
};
\addplot [color=mycolor44, forget plot]
  table[row sep=crcr]{%
0	4395.3\\
2000	3579.3\\
4000	3258.6\\
6000	3219.6\\
8000	3355.5\\
10000	3349.2\\
12000	3490\\
14000	3292.7\\
16000	3368.2\\
18000	3465.5\\
20000	3269.9\\
22000	3237.3\\
24000	3492\\
26000	3343.2\\
28000	3333.5\\
30000	3362.7\\
32000	3319.9\\
34000	3263.1\\
36000	3405.8\\
38000	3359.3\\
40000	3370.7\\
42000	3318.1\\
44000	3204.4\\
46000	3271.8\\
50000	3328.9\\
52000	3363.3\\
54000	3274.1\\
56000	3285.3\\
58000	3432.3\\
60000	3267.2\\
62000	3273.2\\
64000	3175.1\\
66000	3291.7\\
68000	3373.1\\
70000	3391.1\\
72000	3396.7\\
74000	3275.9\\
76000	3252.2\\
78000	3392\\
80000	3367.2\\
82000	3384.7\\
84000	3414.2\\
86000	3271\\
90000	3360.4\\
92000	3281.6\\
94000	3168.2\\
96000	3356.5\\
98000	3288.4\\
1e+05	3448.9\\
};
\addplot [color=mycolor45, forget plot]
  table[row sep=crcr]{%
0	4395.3\\
2000	6861.1\\
4000	7126.2\\
6000	6900\\
8000	6846\\
10000	6984.1\\
12000	6866\\
14000	6841.3\\
16000	6906.8\\
18000	6892\\
20000	6942.5\\
22000	7040.2\\
24000	6959.6\\
28000	6953.3\\
30000	6865.4\\
32000	6956.5\\
34000	7016.1\\
36000	6911.2\\
38000	6998.7\\
40000	6809.9\\
42000	7046.7\\
44000	7004.4\\
46000	6776.6\\
48000	6951.8\\
50000	6881.3\\
54000	7180.8\\
56000	6972.3\\
58000	7130.9\\
60000	7002.8\\
62000	6848.9\\
64000	7039.4\\
66000	7011\\
68000	6909.9\\
70000	6705.3\\
72000	6938.8\\
74000	7044.7\\
76000	7063.7\\
78000	6897.6\\
80000	7044\\
82000	6914.3\\
84000	6906.1\\
86000	6777.4\\
88000	6910.2\\
90000	6962.2\\
92000	6843.4\\
94000	6899.1\\
96000	6979.7\\
98000	6778.6\\
1e+05	6932.2\\
};
\addplot [color=mycolor46, forget plot]
  table[row sep=crcr]{%
0	4395.3\\
2000	4154.5\\
4000	4060.3\\
6000	3995.8\\
10000	4122.2\\
12000	4124.4\\
14000	4046.9\\
16000	4147.1\\
18000	3991\\
20000	4022.1\\
22000	4076.2\\
24000	4071.7\\
26000	4027.3\\
28000	4143.8\\
30000	3964.2\\
32000	4060\\
34000	4071.6\\
36000	4131.2\\
38000	3981.6\\
40000	3946.4\\
42000	4085.3\\
44000	4153.9\\
46000	4000.3\\
50000	4062.1\\
52000	4177.3\\
54000	4069.6\\
56000	4062.1\\
58000	4029\\
60000	4084\\
62000	4064.9\\
64000	3968\\
66000	4110.7\\
68000	4144.4\\
70000	4100.7\\
72000	3920.6\\
74000	4166.8\\
76000	4186.2\\
78000	4059.1\\
80000	4089.1\\
82000	4149.2\\
84000	4026\\
86000	4071.9\\
88000	3991.5\\
90000	4088.9\\
92000	4169.7\\
96000	3985.6\\
98000	4161.7\\
1e+05	3950.1\\
};
\addplot [color=mycolor47, forget plot]
  table[row sep=crcr]{%
0	4395.3\\
2000	5242.6\\
4000	5354.9\\
6000	5120.4\\
8000	5432.4\\
10000	5206.6\\
12000	5290.1\\
16000	5222.6\\
18000	5222.6\\
22000	5278\\
26000	5282\\
28000	5289.1\\
30000	5186.3\\
32000	5334.5\\
34000	5188.9\\
36000	5156.4\\
38000	5076.7\\
40000	5343.5\\
42000	5452.7\\
44000	5216.6\\
46000	5347\\
48000	5429.1\\
50000	5360.3\\
52000	5174\\
54000	5146.3\\
58000	5236.1\\
60000	5246.1\\
62000	5353.2\\
64000	5212.4\\
66000	5353.7\\
68000	5308.9\\
70000	5211.9\\
72000	5251.2\\
74000	5191.3\\
76000	5363.3\\
78000	5362.4\\
80000	5144.4\\
84000	5390.7\\
86000	5334.6\\
88000	5430.5\\
90000	5415.4\\
92000	5347.4\\
94000	5325.9\\
96000	5237.1\\
1e+05	5333.5\\
};
\addplot [color=mycolor48, forget plot]
  table[row sep=crcr]{%
0	4395.3\\
2000	3085\\
4000	2866.4\\
6000	2881.4\\
10000	2861.6\\
12000	2938\\
14000	2819.1\\
16000	2868.1\\
18000	2882.1\\
20000	2831\\
22000	2911.5\\
24000	2859.6\\
26000	2822.3\\
28000	2861.6\\
30000	2839.7\\
32000	2897.5\\
36000	2885.7\\
38000	2947.2\\
40000	2884.7\\
42000	2904.7\\
44000	2914.1\\
46000	2891\\
48000	2932.2\\
50000	2935.4\\
52000	2804.9\\
54000	2873.1\\
56000	2857.7\\
58000	2876.2\\
60000	2919.7\\
62000	2872.3\\
64000	2962.4\\
66000	2944\\
68000	2899.4\\
70000	2846.6\\
72000	2919.5\\
74000	2900.9\\
76000	2925.6\\
78000	2938.7\\
80000	2887.9\\
82000	2906.6\\
84000	2908.1\\
86000	2865.9\\
88000	2931\\
90000	2847.1\\
92000	2886.8\\
96000	2897\\
98000	2923.7\\
1e+05	2884.9\\
};
\addplot [color=mycolor49, forget plot]
  table[row sep=crcr]{%
0	4395.3\\
2000	4315.9\\
4000	4105.3\\
6000	4288.6\\
8000	4360\\
10000	4373.3\\
14000	4340.7\\
16000	4309.4\\
18000	4155\\
20000	4282.7\\
22000	4259.9\\
24000	4490.5\\
26000	4214\\
28000	4366.5\\
30000	4371.7\\
32000	4277.7\\
34000	4384\\
36000	4260.8\\
38000	4374\\
40000	4260.5\\
42000	4336.7\\
44000	4233.1\\
46000	4259.3\\
48000	4413.6\\
50000	4301.7\\
52000	4469.2\\
54000	4306.8\\
56000	4379.6\\
58000	4527.9\\
60000	4384.1\\
62000	4390.1\\
64000	4274.7\\
66000	4354.2\\
70000	4451.4\\
72000	4372.2\\
74000	4274.2\\
76000	4417\\
78000	4473.2\\
80000	4327.8\\
82000	4344.4\\
84000	4291.8\\
86000	4405.3\\
90000	4235.5\\
92000	4384.2\\
94000	4285.1\\
96000	4280.3\\
98000	4261.2\\
1e+05	4446.4\\
};
\addplot [color=mycolor50, forget plot]
  table[row sep=crcr]{%
0	4395.3\\
2000	2702.4\\
4000	2518.5\\
6000	2416.3\\
8000	2488.4\\
10000	2488.8\\
12000	2444.6\\
14000	2495.3\\
16000	2499.7\\
18000	2462.7\\
20000	2486.1\\
22000	2448.1\\
24000	2518.7\\
26000	2345.1\\
28000	2411.1\\
30000	2424.7\\
32000	2524.4\\
34000	2454.6\\
36000	2448.8\\
38000	2434.7\\
40000	2467.7\\
42000	2439.1\\
44000	2399.9\\
46000	2456\\
48000	2414.3\\
50000	2440.2\\
52000	2529.2\\
54000	2511\\
56000	2433.3\\
58000	2464.7\\
60000	2412.9\\
62000	2426.8\\
64000	2472.1\\
66000	2479.6\\
68000	2476.7\\
70000	2463.1\\
72000	2410.8\\
74000	2456.9\\
76000	2490.1\\
78000	2481.3\\
80000	2393.2\\
82000	2405.5\\
84000	2460.9\\
86000	2510.6\\
88000	2476.8\\
90000	2435.9\\
92000	2472.4\\
94000	2410.7\\
96000	2441.9\\
98000	2359.4\\
1e+05	2461.6\\
};
\addplot [color=mycolor51, forget plot]
  table[row sep=crcr]{%
0	4395.3\\
2000	5780.1\\
4000	5938.9\\
6000	5981.6\\
8000	6052.8\\
10000	5834.2\\
12000	5753.5\\
14000	5897.1\\
16000	5871.3\\
18000	6168.7\\
20000	5951.8\\
22000	5912.2\\
26000	5928.3\\
28000	5936\\
30000	6033.7\\
32000	5922\\
34000	5992.5\\
36000	5999.2\\
38000	5682.1\\
40000	6163\\
42000	6172.3\\
46000	5905.4\\
48000	6346.2\\
50000	5922.6\\
52000	6201.5\\
54000	6068\\
56000	5703.5\\
58000	5750.2\\
60000	5956\\
64000	5880.3\\
66000	5895.2\\
68000	6068.5\\
70000	5893.1\\
72000	5836.4\\
74000	5917.9\\
76000	5949.7\\
78000	5926.9\\
80000	6014.1\\
82000	6268.8\\
84000	5898\\
86000	5634.4\\
88000	5797.5\\
90000	5908.4\\
96000	5721.2\\
98000	5781.4\\
1e+05	5704\\
};
\addplot [color=mycolor52, forget plot]
  table[row sep=crcr]{%
0	4395.3\\
2000	3615.9\\
4000	3519.6\\
6000	3463\\
8000	3560.8\\
10000	3463.5\\
12000	3593.2\\
14000	3324.2\\
16000	3631.5\\
18000	3497.5\\
20000	3495.4\\
22000	3309.5\\
24000	3406.6\\
26000	3432.5\\
28000	3391.9\\
30000	3523.1\\
32000	3456.3\\
34000	3650.2\\
36000	3406.5\\
38000	3551.1\\
40000	3517.6\\
42000	3440.9\\
44000	3492.4\\
46000	3581.3\\
48000	3595.6\\
50000	3500.6\\
52000	3450.3\\
54000	3568.5\\
56000	3359\\
58000	3453.3\\
60000	3561.7\\
62000	3350.3\\
64000	3310.4\\
66000	3549.3\\
70000	3486.8\\
72000	3413.4\\
74000	3453\\
76000	3505\\
78000	3477.5\\
80000	3362.8\\
82000	3353.4\\
84000	3405.8\\
86000	3478.4\\
88000	3365.8\\
90000	3445.6\\
92000	3398\\
94000	3475.5\\
96000	3384.2\\
98000	3480.4\\
1e+05	3635\\
};
\addplot [color=mycolor53, forget plot]
  table[row sep=crcr]{%
0	4395.3\\
2000	5574.9\\
4000	5794.1\\
6000	5620.9\\
8000	5714.7\\
10000	5472\\
12000	5788.9\\
14000	5912\\
16000	5649.1\\
18000	5603.3\\
24000	5794.1\\
26000	5475.6\\
28000	5456.6\\
30000	5600.8\\
32000	5571.3\\
34000	5643.8\\
36000	5656.5\\
38000	5978.9\\
40000	5656.3\\
42000	5667.4\\
44000	5511.2\\
46000	5689.8\\
48000	5615.3\\
50000	5703.2\\
52000	5548.3\\
54000	5535.6\\
56000	5625.5\\
58000	5757.2\\
60000	5593.4\\
64000	5647.6\\
66000	5643.3\\
68000	5615.2\\
70000	5451.1\\
72000	5823.5\\
74000	5439.6\\
76000	5599.8\\
78000	5686\\
80000	5669\\
82000	5635\\
84000	5513.5\\
86000	5815.7\\
88000	5963.5\\
90000	5624.2\\
92000	5654.9\\
94000	5557.3\\
96000	5944.1\\
1e+05	5715.9\\
};
\addplot [color=mycolor54, forget plot]
  table[row sep=crcr]{%
0	4395.3\\
2000	4677.1\\
4000	4631.6\\
6000	4560.8\\
8000	4437.7\\
10000	4688.3\\
12000	4607.5\\
14000	4753\\
16000	4584.9\\
18000	4605.1\\
20000	4564.9\\
22000	4555.4\\
24000	4523.1\\
26000	4620.5\\
28000	4480.8\\
30000	4590.7\\
32000	4797.8\\
34000	4543.2\\
36000	4547\\
38000	4605.9\\
40000	4454.3\\
46000	4560.9\\
48000	4604.2\\
50000	4526.9\\
52000	4646.8\\
56000	4630.4\\
58000	4517\\
60000	4449.4\\
62000	4439.1\\
64000	4654\\
66000	4636.1\\
68000	4605\\
70000	4625.8\\
72000	4683.4\\
76000	4475.4\\
78000	4607.5\\
80000	4691.1\\
82000	4582.1\\
84000	4585.8\\
86000	4475.9\\
88000	4561.1\\
90000	4535.7\\
92000	4559.4\\
96000	4478.1\\
98000	4501.4\\
1e+05	4476.1\\
};
\addplot [color=mycolor55, forget plot]
  table[row sep=crcr]{%
0	4395.3\\
2000	5827.7\\
4000	5829.5\\
6000	5988.7\\
8000	5823.5\\
10000	5916.7\\
12000	5791.3\\
16000	5625.7\\
22000	5795\\
24000	5774.6\\
26000	5851.6\\
28000	5783.1\\
30000	5755.2\\
32000	5675.2\\
34000	5846.2\\
36000	5689.9\\
38000	5668.6\\
40000	5820.5\\
42000	5684.7\\
44000	5751.6\\
46000	5886.3\\
48000	5634.3\\
50000	5698.5\\
52000	5790.9\\
54000	5862.9\\
56000	5748.3\\
58000	5708.3\\
60000	5536\\
62000	5831.7\\
64000	5798.7\\
66000	5650.7\\
68000	5810.7\\
70000	5751.3\\
72000	5803.8\\
74000	5897.1\\
76000	5693.9\\
80000	6108.6\\
84000	5713.5\\
86000	5756.8\\
88000	5676\\
90000	5829.7\\
92000	5772.6\\
94000	5876.4\\
96000	5753.1\\
98000	5831.5\\
1e+05	5787.8\\
};
\addplot [color=mycolor56, forget plot]
  table[row sep=crcr]{%
0	4395.3\\
2000	4119.4\\
4000	3923.9\\
6000	4028\\
8000	3968.6\\
10000	3998.7\\
12000	3985.2\\
14000	3912.9\\
16000	4002\\
18000	3997.2\\
20000	4052.4\\
22000	3899.4\\
24000	4016\\
26000	3909.1\\
28000	3929\\
30000	3995.6\\
32000	4030\\
36000	4001.3\\
38000	4031.8\\
40000	3932.6\\
42000	3889.9\\
44000	3941.7\\
46000	3873.4\\
48000	3992.8\\
50000	4098.7\\
52000	3824.5\\
54000	3931.9\\
56000	3974.2\\
58000	4056.3\\
60000	3949.8\\
62000	3982.5\\
64000	3917.4\\
66000	4120.2\\
68000	3936.5\\
70000	4038.7\\
74000	4049.6\\
76000	3918\\
78000	4066.5\\
80000	3932.5\\
82000	4004.6\\
84000	4014.6\\
86000	4003.2\\
88000	3960.8\\
90000	4003.3\\
92000	3958.2\\
94000	4058.5\\
98000	3741.6\\
1e+05	3930.8\\
};
\addplot [color=mycolor57, forget plot]
  table[row sep=crcr]{%
0	4395.3\\
2000	6392.9\\
4000	6660\\
6000	6719.4\\
8000	6730.2\\
10000	6636.3\\
12000	6796.9\\
14000	6758.2\\
16000	6765.7\\
18000	6510.5\\
20000	6578\\
22000	6814.8\\
24000	6613\\
26000	6849.1\\
28000	6581.8\\
30000	6758.5\\
32000	6628\\
34000	6723\\
36000	6453.5\\
38000	6755.3\\
40000	6763.5\\
42000	6660.7\\
44000	6832.2\\
46000	6689.5\\
48000	6740.2\\
50000	6837.2\\
52000	6704.8\\
54000	6715\\
56000	6868.2\\
58000	6736.6\\
60000	6770.2\\
62000	6630.5\\
64000	6726.3\\
66000	6670.2\\
68000	6760.1\\
70000	6891.7\\
72000	6793.3\\
74000	6631.6\\
76000	6550.9\\
78000	6801.2\\
80000	6584.8\\
84000	6802.2\\
86000	6800\\
88000	6890.2\\
90000	6790.9\\
92000	6798.7\\
94000	6917.4\\
96000	6639\\
98000	6507.7\\
1e+05	6751.8\\
};
\addplot [color=mycolor58, forget plot]
  table[row sep=crcr]{%
0	4395.3\\
2000	3709\\
4000	3573.4\\
8000	3488.5\\
10000	3525.4\\
12000	3538.5\\
14000	3452.1\\
16000	3698.6\\
18000	3527.6\\
22000	3652.8\\
24000	3559.9\\
26000	3537.9\\
28000	3557.8\\
30000	3526.1\\
32000	3422.7\\
34000	3569.5\\
36000	3502\\
38000	3547.7\\
40000	3550.2\\
42000	3609.7\\
44000	3477\\
46000	3520.2\\
48000	3476.6\\
50000	3514.5\\
52000	3568.4\\
54000	3607\\
56000	3530.8\\
58000	3512.5\\
60000	3638.5\\
62000	3528.7\\
64000	3555.2\\
66000	3674.5\\
68000	3566.1\\
70000	3577.8\\
72000	3571.7\\
74000	3520.8\\
76000	3568.7\\
78000	3530.6\\
80000	3541.4\\
82000	3397.2\\
84000	3611.5\\
86000	3609.1\\
88000	3483.2\\
90000	3594.6\\
92000	3589.6\\
94000	3531.8\\
96000	3617.5\\
98000	3529\\
1e+05	3498.2\\
};
\addplot [color=mycolor59, forget plot]
  table[row sep=crcr]{%
0	4395.3\\
2000	5019.6\\
6000	5126\\
8000	5118.9\\
10000	4998.2\\
12000	4906.1\\
14000	5154\\
18000	5092.3\\
20000	5017.3\\
24000	5134.3\\
26000	5171.1\\
28000	5104.4\\
30000	4932.5\\
32000	4935.4\\
34000	5124.4\\
36000	4971.1\\
38000	4915\\
40000	5102.8\\
42000	5071\\
44000	4949.5\\
46000	5241.7\\
48000	5061.5\\
50000	4957.9\\
52000	5064.8\\
54000	5086.3\\
56000	5151.7\\
58000	5099.9\\
60000	5271.4\\
62000	5085.7\\
64000	4920.8\\
68000	5028.8\\
70000	4993.7\\
72000	5075.5\\
74000	5190.3\\
76000	4978.2\\
78000	5088.2\\
80000	5049.7\\
82000	4951.4\\
84000	5045.1\\
88000	5027.3\\
90000	4873\\
92000	5063.3\\
94000	4889.2\\
96000	5077.2\\
98000	5002.2\\
1e+05	5101.4\\
};
\addplot [color=mycolor60, forget plot]
  table[row sep=crcr]{%
0	4395.3\\
2000	3613.2\\
4000	3501.9\\
6000	3587.7\\
8000	3531.1\\
12000	3493\\
14000	3432.8\\
16000	3502.9\\
18000	3549.7\\
22000	3538.9\\
24000	3595.4\\
26000	3540\\
28000	3630\\
30000	3505.3\\
32000	3528.6\\
34000	3577.5\\
36000	3477.4\\
38000	3486.8\\
40000	3471.9\\
42000	3504.4\\
44000	3549.5\\
46000	3532.9\\
50000	3540.5\\
52000	3602.6\\
54000	3500.6\\
58000	3604\\
60000	3447.6\\
64000	3519\\
66000	3482.1\\
68000	3515.6\\
70000	3537\\
72000	3527.3\\
74000	3582.2\\
76000	3422.2\\
78000	3596.7\\
80000	3517.9\\
82000	3581.9\\
84000	3468\\
86000	3532.4\\
90000	3575.7\\
94000	3581.2\\
96000	3555.8\\
98000	3583.5\\
1e+05	3520.5\\
};
\addplot [color=mycolor61, forget plot]
  table[row sep=crcr]{%
0	4395.3\\
2000	3243.6\\
4000	3092.8\\
6000	3123.9\\
8000	3131.6\\
12000	3060.1\\
14000	3152.7\\
16000	3167.1\\
18000	3109.1\\
20000	3082.6\\
22000	3132.4\\
24000	3100.9\\
26000	3153.1\\
30000	3124.1\\
32000	3120.2\\
34000	3099.5\\
36000	3116.6\\
38000	3078.7\\
40000	3089.2\\
42000	3118.8\\
44000	3072.3\\
46000	3079.4\\
48000	3061.6\\
52000	3134.3\\
54000	3082.9\\
58000	3045\\
64000	3114.7\\
66000	3103.1\\
68000	3123.9\\
70000	3111.3\\
72000	3055.5\\
74000	3180.1\\
76000	3098\\
78000	3074.1\\
80000	3137.7\\
82000	3152.5\\
84000	3090.3\\
86000	3129.6\\
88000	3093.2\\
90000	3029.6\\
92000	3081.9\\
94000	3094.9\\
96000	3189.7\\
98000	3165.9\\
1e+05	3037.7\\
};
\addplot [color=mycolor62, forget plot]
  table[row sep=crcr]{%
0	4395.3\\
2000	4903.7\\
4000	5066.5\\
6000	4915.5\\
8000	5015.9\\
12000	4884.7\\
14000	4929.2\\
16000	4929.8\\
18000	4968.6\\
20000	5113.6\\
22000	4924.2\\
24000	4842.7\\
26000	4938.9\\
28000	4864.6\\
30000	4973\\
32000	4998.1\\
34000	4945\\
36000	5089.3\\
38000	4964.4\\
40000	4996.7\\
42000	4936.5\\
44000	4983\\
46000	4917\\
48000	4984.9\\
50000	5009.8\\
52000	4981.9\\
54000	4936.9\\
56000	4949.2\\
58000	4912.5\\
60000	4956.5\\
62000	4957.1\\
64000	5076.8\\
66000	4926.1\\
68000	4894.5\\
70000	4913.1\\
72000	5083.4\\
74000	4975.8\\
76000	4920.5\\
78000	4997.7\\
80000	4971.6\\
82000	4980.6\\
84000	5013.4\\
86000	4898.1\\
88000	4994.5\\
90000	5021.8\\
92000	4994.3\\
94000	4890.9\\
96000	4917\\
1e+05	5141.5\\
};
\addplot [color=mycolor63, forget plot]
  table[row sep=crcr]{%
0	4395.3\\
2000	2624.2\\
4000	2398.3\\
6000	2388.2\\
12000	2423.1\\
14000	2381.3\\
16000	2375.1\\
18000	2409.8\\
20000	2354.9\\
24000	2364.6\\
26000	2385\\
28000	2298.4\\
30000	2376.5\\
32000	2408.3\\
34000	2366.2\\
36000	2380.5\\
38000	2423.6\\
40000	2418.7\\
42000	2402.1\\
44000	2369.7\\
46000	2344.3\\
48000	2458.8\\
50000	2349.3\\
52000	2323.7\\
54000	2366.6\\
56000	2388.5\\
58000	2386.2\\
60000	2400.9\\
62000	2389.3\\
64000	2399.4\\
66000	2346.8\\
68000	2359.1\\
70000	2359.6\\
72000	2389\\
74000	2366.6\\
76000	2387.8\\
78000	2396.3\\
80000	2432.6\\
82000	2430\\
84000	2420.3\\
86000	2458.5\\
88000	2391.9\\
90000	2365.4\\
92000	2381\\
94000	2379.2\\
98000	2412.6\\
1e+05	2391.1\\
};
\addplot [color=mycolor64, forget plot]
  table[row sep=crcr]{%
0	4395.3\\
2000	3218.9\\
4000	3057.2\\
6000	3093.4\\
8000	3037.6\\
10000	3008.8\\
12000	3085.3\\
14000	3078.3\\
16000	3041.7\\
18000	3066.8\\
20000	3045.4\\
22000	3063.8\\
24000	3031.2\\
26000	3086.8\\
28000	3076.6\\
30000	3027.2\\
32000	3040.6\\
34000	3009.8\\
38000	2982.3\\
40000	3030.8\\
42000	3094.1\\
44000	2989.4\\
46000	3061.5\\
48000	3049.3\\
50000	2978.9\\
52000	2992.8\\
54000	3027.3\\
56000	3030.8\\
58000	2978\\
60000	3025.3\\
64000	3015.4\\
68000	3015.3\\
70000	2989.8\\
72000	3107.7\\
74000	3094\\
76000	2995.7\\
78000	3076.9\\
80000	3021.6\\
82000	2991.2\\
84000	2987.8\\
86000	3071.4\\
88000	3055.3\\
90000	2957\\
92000	3044.9\\
94000	3059.4\\
96000	3002.3\\
98000	3006.5\\
1e+05	3073.1\\
};
\addplot [color=mycolor65, forget plot]
  table[row sep=crcr]{%
0	4395.3\\
2000	2799.9\\
4000	2561.9\\
6000	2551.5\\
8000	2557.4\\
12000	2507.3\\
14000	2499.8\\
16000	2459.1\\
18000	2552.6\\
20000	2543.3\\
22000	2552.8\\
24000	2546.4\\
26000	2497\\
30000	2470.6\\
32000	2520.3\\
34000	2534.3\\
36000	2558.5\\
38000	2470.6\\
40000	2504.6\\
44000	2515.5\\
46000	2546\\
48000	2524.9\\
50000	2449.7\\
52000	2531.1\\
54000	2557.9\\
56000	2494.9\\
58000	2495.2\\
60000	2530\\
62000	2532.9\\
64000	2556\\
66000	2504.1\\
68000	2519.4\\
70000	2551.1\\
72000	2519.7\\
74000	2506.5\\
76000	2513.4\\
78000	2482.8\\
80000	2508.9\\
82000	2553\\
84000	2580\\
86000	2518.8\\
88000	2522.4\\
90000	2445.7\\
92000	2542.7\\
94000	2547\\
96000	2472.4\\
98000	2562.5\\
1e+05	2527.9\\
};
\addplot [color=mycolor66, forget plot]
  table[row sep=crcr]{%
0	4395.3\\
2000	3523.2\\
4000	3406.3\\
6000	3393.9\\
8000	3369.9\\
12000	3380\\
14000	3407.8\\
16000	3485.2\\
18000	3453\\
22000	3349.3\\
24000	3354.6\\
26000	3374.7\\
28000	3345.5\\
30000	3362.8\\
32000	3416.2\\
34000	3373.8\\
36000	3464.7\\
38000	3386.3\\
40000	3439.8\\
42000	3392\\
44000	3374.5\\
46000	3403.4\\
48000	3367.1\\
50000	3396.3\\
52000	3377.1\\
54000	3381.5\\
58000	3456.2\\
60000	3373.9\\
62000	3404.5\\
64000	3423.1\\
68000	3407.6\\
70000	3347.1\\
72000	3444.6\\
74000	3339.1\\
76000	3436.6\\
78000	3433.2\\
80000	3415\\
82000	3416.4\\
84000	3361.2\\
86000	3417.2\\
88000	3388.1\\
90000	3404.6\\
92000	3405.7\\
94000	3367.5\\
96000	3352.9\\
98000	3429\\
1e+05	3380.7\\
};
\addplot [color=mycolor67, forget plot]
  table[row sep=crcr]{%
0	4395.3\\
2000	3406.7\\
4000	3203.6\\
6000	3206.9\\
8000	3181.3\\
10000	3179\\
12000	3070\\
14000	3126.2\\
16000	3160.6\\
18000	3168.5\\
20000	3034\\
22000	3110.2\\
26000	3155\\
28000	3117.9\\
30000	3155\\
32000	3170.2\\
36000	3134.9\\
38000	3183.9\\
40000	3116.6\\
42000	3164.4\\
44000	3167.2\\
46000	3099.9\\
48000	3074.7\\
50000	3131.4\\
52000	3103.6\\
54000	3227.4\\
56000	3166.7\\
58000	3166.9\\
60000	3195.1\\
62000	3194.4\\
64000	3158.4\\
66000	3204.8\\
68000	3139.8\\
70000	3138.2\\
72000	3155.2\\
74000	3187\\
78000	3112.6\\
80000	3173\\
82000	3116.2\\
86000	3149.5\\
90000	3172.8\\
92000	3085.5\\
94000	3098.6\\
96000	3176.5\\
1e+05	3120.8\\
};
\addplot [color=mycolor67]
  table[row sep=crcr]{%
0	4395.3\\
2000	2863.8\\
4000	2594.6\\
6000	2651.9\\
10000	2587.6\\
12000	2633.9\\
14000	2586.9\\
16000	2617.7\\
18000	2626.2\\
20000	2650.8\\
22000	2630.7\\
24000	2583.3\\
26000	2621.8\\
28000	2579.3\\
30000	2624.6\\
32000	2599.6\\
36000	2638.3\\
40000	2641.6\\
42000	2597\\
44000	2632.2\\
46000	2605\\
48000	2626.1\\
50000	2601.2\\
52000	2605.5\\
54000	2564.4\\
56000	2623.5\\
60000	2597.9\\
62000	2628.2\\
64000	2628.3\\
66000	2583.6\\
68000	2598.6\\
70000	2642.2\\
72000	2630\\
74000	2606.3\\
76000	2675\\
78000	2658.4\\
80000	2606.2\\
82000	2618\\
84000	2585.9\\
86000	2613.7\\
90000	2583.8\\
92000	2597.1\\
94000	2596.8\\
96000	2660.4\\
98000	2593.1\\
1e+05	2561\\
};
\addlegendentry{Classe k=43}

\addplot [color=mycolor68, forget plot]
  table[row sep=crcr]{%
0	4395.3\\
2000	2737.9\\
4000	2437.3\\
6000	2410\\
8000	2432.3\\
12000	2417.9\\
14000	2453.7\\
16000	2441.2\\
18000	2457.3\\
20000	2445.7\\
22000	2443.9\\
24000	2430.5\\
26000	2443\\
28000	2413.6\\
30000	2422.8\\
32000	2447.8\\
36000	2409.8\\
38000	2439.7\\
40000	2394.3\\
42000	2403.5\\
44000	2435.2\\
46000	2443.6\\
48000	2442.6\\
50000	2482.6\\
52000	2411.3\\
54000	2414\\
56000	2436.9\\
58000	2471.7\\
60000	2434.5\\
62000	2427.3\\
64000	2452.7\\
68000	2482\\
70000	2427\\
72000	2465.2\\
74000	2408.2\\
76000	2417\\
78000	2441.7\\
82000	2465.9\\
84000	2446.1\\
86000	2412.7\\
88000	2432.9\\
90000	2424.1\\
94000	2449.8\\
96000	2444.2\\
98000	2447.9\\
1e+05	2436.7\\
};
\addplot [color=mycolor69, forget plot]
  table[row sep=crcr]{%
0	4395.3\\
2000	3103.8\\
4000	2871.4\\
6000	2931.2\\
8000	2874.8\\
10000	2875\\
12000	2883.7\\
16000	2836.2\\
18000	2886.3\\
22000	2853.5\\
24000	2919.6\\
26000	2842.3\\
28000	2869.5\\
30000	2883\\
32000	2884.6\\
34000	2865.2\\
36000	2859.9\\
38000	2876.6\\
40000	2851.3\\
42000	2920\\
44000	2864.1\\
48000	2887.1\\
50000	2851\\
52000	2905\\
54000	2920.8\\
56000	2906.8\\
58000	2869.2\\
60000	2879.9\\
62000	2877.7\\
64000	2899.3\\
66000	2901.1\\
70000	2851.9\\
72000	2898.8\\
74000	2881.9\\
76000	2924.8\\
78000	2857.4\\
80000	2862.3\\
82000	2824.4\\
86000	2872.8\\
88000	2937.2\\
90000	2892\\
92000	2888.4\\
96000	2856\\
98000	2881.3\\
1e+05	2866.8\\
};
\addplot [color=mycolor70, forget plot]
  table[row sep=crcr]{%
0	4395.3\\
2000	2380.4\\
4000	2063.8\\
6000	2063.4\\
8000	2074.1\\
10000	2095.6\\
12000	2036.1\\
14000	2122.2\\
18000	2050.5\\
20000	2058.8\\
22000	2073.9\\
24000	2071.7\\
26000	2125.9\\
28000	2059.2\\
30000	2046.9\\
34000	2103.7\\
36000	2098.2\\
38000	2074.8\\
40000	2113\\
42000	2059.2\\
46000	2068.2\\
48000	2064.8\\
50000	2089.3\\
52000	2122.2\\
54000	2078.4\\
58000	2114.9\\
60000	2119.1\\
62000	2102.6\\
64000	2069.8\\
66000	2078.9\\
68000	2107.4\\
70000	2070.5\\
74000	2099.3\\
76000	2040.8\\
78000	2070\\
80000	2080.3\\
82000	2126.7\\
84000	2103.8\\
86000	2033\\
88000	2043\\
90000	2103.5\\
92000	2112.3\\
94000	2090.5\\
96000	2094.4\\
98000	2049.2\\
1e+05	2083.5\\
};
\addplot [color=mycolor71, forget plot]
  table[row sep=crcr]{%
0	4395.3\\
2000	2694.8\\
4000	2466.1\\
6000	2465.3\\
8000	2477.2\\
10000	2442\\
14000	2477.7\\
16000	2457.6\\
18000	2458.2\\
20000	2496.4\\
22000	2433.4\\
24000	2427.9\\
26000	2439.8\\
28000	2482.2\\
30000	2447.9\\
32000	2483.8\\
34000	2460.1\\
36000	2460.4\\
40000	2440.5\\
42000	2492.6\\
44000	2486.9\\
46000	2454.1\\
48000	2469.5\\
50000	2495.5\\
52000	2466.2\\
54000	2478.9\\
56000	2428.3\\
58000	2467\\
60000	2413.7\\
64000	2457.4\\
66000	2471.2\\
68000	2464.9\\
70000	2476.5\\
72000	2473.9\\
78000	2436.7\\
80000	2463.4\\
82000	2479.7\\
84000	2461.9\\
86000	2453.2\\
88000	2462.4\\
90000	2491.7\\
92000	2491.3\\
94000	2440.6\\
96000	2459.1\\
1e+05	2466.4\\
};
\addplot [color=mycolor72, forget plot]
  table[row sep=crcr]{%
0	4395.3\\
2000	2475.2\\
4000	2207.8\\
6000	2206.9\\
8000	2186.3\\
10000	2208.5\\
12000	2223.3\\
14000	2175.9\\
16000	2205.2\\
18000	2168.8\\
20000	2142.6\\
22000	2162.4\\
24000	2126.1\\
26000	2238.6\\
28000	2203.7\\
30000	2130.3\\
32000	2171.5\\
34000	2182.1\\
36000	2158.5\\
38000	2156.2\\
40000	2170.3\\
42000	2228.9\\
44000	2128.8\\
48000	2166.5\\
50000	2142.4\\
52000	2209.9\\
54000	2156\\
56000	2199.8\\
58000	2159.1\\
60000	2083.2\\
62000	2152.2\\
64000	2174.3\\
66000	2147.1\\
68000	2145.6\\
70000	2187.1\\
72000	2156.9\\
74000	2172.5\\
76000	2151.6\\
78000	2146\\
80000	2192.9\\
82000	2172\\
84000	2188.1\\
86000	2174.5\\
88000	2151.6\\
92000	2174.1\\
94000	2197\\
96000	2164.2\\
98000	2173.2\\
1e+05	2143.6\\
};
\addplot [color=mycolor73, forget plot]
  table[row sep=crcr]{%
0	4395.3\\
2000	2832.2\\
4000	2591.8\\
6000	2603.2\\
10000	2572.5\\
12000	2584.5\\
14000	2539.4\\
16000	2603.1\\
18000	2630.7\\
20000	2618.5\\
22000	2555.5\\
24000	2563.2\\
26000	2563.3\\
28000	2611.3\\
30000	2582.5\\
32000	2581.5\\
34000	2566.1\\
36000	2587.3\\
38000	2587.7\\
40000	2563\\
42000	2610.4\\
44000	2568.2\\
46000	2576.1\\
48000	2592.1\\
50000	2595.5\\
52000	2575\\
54000	2579.1\\
56000	2598.5\\
58000	2603.8\\
60000	2596.6\\
62000	2641.5\\
64000	2569.4\\
66000	2566.8\\
68000	2602\\
70000	2601.8\\
72000	2583.2\\
74000	2531.5\\
76000	2551.1\\
78000	2623.4\\
80000	2634.1\\
82000	2585.1\\
84000	2638\\
86000	2586.2\\
88000	2565.1\\
90000	2587\\
92000	2617.1\\
94000	2568.3\\
98000	2591.4\\
1e+05	2561.5\\
};
\addplot [color=mycolor74, forget plot]
  table[row sep=crcr]{%
0	4395.3\\
2000	1863.8\\
4000	1522\\
6000	1509.7\\
8000	1528.9\\
10000	1510.6\\
12000	1503.7\\
18000	1512.6\\
20000	1505.3\\
22000	1520.6\\
24000	1495.6\\
26000	1508.5\\
28000	1499.9\\
30000	1520.6\\
32000	1515.8\\
34000	1517.6\\
36000	1527.4\\
38000	1520.7\\
40000	1536.4\\
42000	1562.1\\
44000	1501.7\\
46000	1536.7\\
48000	1502.3\\
50000	1511.5\\
52000	1511.9\\
54000	1496.6\\
56000	1513.7\\
62000	1499.2\\
64000	1519.2\\
66000	1514.6\\
68000	1505.1\\
70000	1522.8\\
72000	1510.4\\
74000	1543.1\\
76000	1490\\
78000	1511.2\\
80000	1520.9\\
82000	1510.2\\
84000	1522.4\\
86000	1514.6\\
88000	1528.5\\
90000	1515\\
92000	1513.5\\
94000	1495.9\\
96000	1531.2\\
98000	1474.4\\
1e+05	1532.9\\
};
\addplot [color=mycolor75, forget plot]
  table[row sep=crcr]{%
0	4395.3\\
2000	2525.1\\
4000	2183\\
6000	2153.8\\
8000	2164.3\\
10000	2144.5\\
12000	2150.7\\
14000	2192.1\\
16000	2177.4\\
18000	2211.5\\
20000	2170.3\\
22000	2187.7\\
24000	2172\\
26000	2177.6\\
28000	2213.9\\
30000	2206.4\\
32000	2181.1\\
34000	2163.1\\
38000	2199.3\\
40000	2182.5\\
42000	2195.6\\
44000	2168.9\\
46000	2207.4\\
48000	2156.4\\
50000	2160.3\\
52000	2182\\
54000	2211.4\\
56000	2179.8\\
58000	2136.8\\
60000	2227.6\\
64000	2178.6\\
68000	2164\\
70000	2189.6\\
72000	2156.6\\
76000	2179.2\\
78000	2180.3\\
80000	2217.6\\
82000	2191.9\\
84000	2192.7\\
86000	2166.4\\
88000	2161.1\\
90000	2189.6\\
92000	2192.4\\
94000	2208.2\\
96000	2181.3\\
98000	2189.7\\
1e+05	2182.4\\
};
\addplot [color=mycolor76, forget plot]
  table[row sep=crcr]{%
0	4395.3\\
2000	2323.8\\
4000	2019.2\\
6000	1953.8\\
8000	1962.1\\
10000	1949.4\\
12000	1993.1\\
14000	1965.3\\
16000	2017.6\\
18000	1992.4\\
20000	2002.2\\
22000	1966.7\\
24000	1996.9\\
26000	2015.3\\
30000	1978\\
34000	2024.5\\
36000	1973.5\\
38000	1983.9\\
40000	1975.7\\
42000	1973.6\\
44000	1994.2\\
46000	2025.2\\
48000	1999.5\\
52000	1981.4\\
54000	2007.8\\
56000	1987.4\\
58000	2000.1\\
60000	1997.8\\
62000	1968.9\\
64000	1996.7\\
66000	1995.7\\
68000	1963.8\\
70000	2008\\
72000	2016.4\\
74000	2008.5\\
76000	1995\\
78000	1997.2\\
80000	2013.1\\
82000	1979.6\\
84000	2025.8\\
86000	1993.4\\
88000	2010.8\\
90000	1976.2\\
92000	1967.1\\
94000	1969.6\\
96000	1986.7\\
98000	1989.7\\
1e+05	1955.2\\
};
\addplot [color=mycolor77, forget plot]
  table[row sep=crcr]{%
0	4395.3\\
2000	1937.9\\
4000	1539.4\\
6000	1524.9\\
8000	1531.9\\
10000	1525.6\\
12000	1532.8\\
16000	1533.7\\
18000	1549.7\\
20000	1530.8\\
22000	1525.7\\
28000	1538.3\\
30000	1533.5\\
32000	1544.1\\
34000	1540.9\\
36000	1525.5\\
40000	1551.5\\
42000	1544.6\\
44000	1522.7\\
46000	1517.4\\
48000	1541.4\\
50000	1517.3\\
54000	1545.4\\
56000	1531.1\\
58000	1551.4\\
60000	1527.1\\
62000	1541.4\\
64000	1528.6\\
66000	1569.5\\
68000	1572.8\\
70000	1548.9\\
72000	1558.5\\
74000	1517.2\\
76000	1525\\
78000	1541.2\\
80000	1553\\
82000	1527.7\\
84000	1509.4\\
86000	1524.1\\
88000	1518.3\\
90000	1524.4\\
92000	1500.3\\
94000	1525.1\\
96000	1534.2\\
98000	1551.7\\
1e+05	1538.4\\
};
\addplot [color=mycolor78, forget plot]
  table[row sep=crcr]{%
0	4395.3\\
2000	2608.3\\
4000	2226.3\\
6000	2288.8\\
8000	2275\\
10000	2273.7\\
14000	2216.6\\
16000	2272.2\\
18000	2288.4\\
20000	2216.8\\
22000	2195\\
24000	2289.4\\
26000	2281.8\\
28000	2300.6\\
30000	2278.9\\
32000	2301.5\\
34000	2228.6\\
36000	2305.9\\
38000	2293.3\\
40000	2293.8\\
42000	2279\\
44000	2285.7\\
46000	2330.1\\
48000	2289.1\\
50000	2292.2\\
52000	2254.6\\
54000	2247\\
56000	2265\\
60000	2265.7\\
62000	2293.4\\
64000	2251\\
66000	2292.8\\
68000	2252.3\\
72000	2272.8\\
74000	2246.9\\
76000	2230.2\\
78000	2263.8\\
80000	2271\\
84000	2270.2\\
86000	2247.7\\
88000	2211.3\\
90000	2293.7\\
92000	2300.7\\
94000	2240\\
96000	2301.9\\
98000	2265\\
1e+05	2279.5\\
};
\addplot [color=mycolor79, forget plot]
  table[row sep=crcr]{%
0	4395.3\\
2000	2517.8\\
4000	2147.1\\
6000	2123.2\\
8000	2129.4\\
10000	2162.8\\
12000	2125.7\\
14000	2135.8\\
16000	2156.1\\
20000	2102.3\\
22000	2109.1\\
24000	2144.9\\
26000	2115\\
30000	2204.5\\
32000	2138.8\\
34000	2116.9\\
36000	2116.9\\
38000	2187.8\\
40000	2131.9\\
42000	2122.3\\
44000	2100.4\\
46000	2092.1\\
48000	2127.2\\
50000	2084.6\\
52000	2190\\
54000	2144.3\\
56000	2156.8\\
58000	2090.1\\
60000	2133.2\\
62000	2081\\
64000	2112.5\\
66000	2162\\
68000	2129.5\\
70000	2145.1\\
72000	2122.6\\
74000	2109.7\\
76000	2126.9\\
78000	2103.8\\
80000	2132.9\\
82000	2091.5\\
84000	2154.1\\
86000	2125\\
88000	2151.9\\
90000	2141.6\\
92000	2116\\
94000	2132.1\\
96000	2131.7\\
98000	2181.4\\
1e+05	2159.4\\
};
\addplot [color=mycolor79, forget plot]
  table[row sep=crcr]{%
0	4395.3\\
2000	2352.6\\
4000	1980.8\\
6000	1984.1\\
8000	1979.5\\
10000	1932.2\\
12000	1947.9\\
14000	2013.5\\
16000	1976.4\\
18000	2019.3\\
20000	1976.3\\
22000	2023.5\\
24000	2000.3\\
26000	1944.5\\
28000	2033.9\\
32000	1962.5\\
34000	1966\\
36000	1950.3\\
38000	1979.4\\
40000	2034.5\\
42000	2034.4\\
44000	1947.3\\
46000	2018.4\\
48000	1981.4\\
50000	1989\\
52000	1974\\
54000	1985.8\\
56000	2026.4\\
58000	2030.5\\
60000	1969.7\\
62000	2014.3\\
64000	1977.5\\
66000	1964.1\\
68000	1999.4\\
70000	2014.1\\
72000	1948.3\\
74000	1958.3\\
76000	2025.4\\
78000	1973\\
80000	1983.5\\
82000	1978.8\\
84000	2012.1\\
86000	1992.7\\
88000	1983.6\\
92000	2004.1\\
94000	1979.6\\
96000	1933.7\\
98000	2015.5\\
1e+05	1941.1\\
};
\addplot [color=mycolor80, forget plot]
  table[row sep=crcr]{%
0	4395.3\\
2000	1938.6\\
4000	1549.8\\
6000	1552.8\\
8000	1531.6\\
10000	1499.8\\
12000	1541.9\\
14000	1525.1\\
16000	1537.8\\
18000	1528.7\\
20000	1525.5\\
22000	1502.8\\
24000	1535.5\\
26000	1503.2\\
28000	1526.6\\
30000	1517.7\\
32000	1519.8\\
34000	1497\\
38000	1517.1\\
42000	1519.3\\
44000	1499.7\\
46000	1490.4\\
48000	1531.7\\
50000	1522.6\\
52000	1538.2\\
54000	1527.5\\
56000	1521.9\\
58000	1509\\
60000	1540.5\\
62000	1523.6\\
64000	1524.3\\
66000	1520.5\\
68000	1534.8\\
70000	1531.6\\
72000	1524.2\\
74000	1503.7\\
76000	1569.6\\
78000	1554\\
80000	1529.3\\
82000	1531.5\\
84000	1541\\
86000	1530\\
88000	1555.9\\
90000	1553.4\\
92000	1514.8\\
94000	1530.4\\
96000	1528.7\\
98000	1542\\
1e+05	1496.4\\
};
\addplot [color=mycolor81, forget plot]
  table[row sep=crcr]{%
0	4395.3\\
2000	1870.4\\
4000	1426\\
6000	1432.5\\
10000	1436.9\\
12000	1412.3\\
14000	1438.9\\
16000	1416.8\\
18000	1417\\
22000	1438\\
26000	1423.7\\
28000	1448.5\\
30000	1451.8\\
32000	1448.1\\
34000	1431.7\\
36000	1450.6\\
38000	1412.1\\
40000	1442\\
42000	1439.5\\
44000	1422.3\\
46000	1450.8\\
48000	1419.3\\
50000	1392.6\\
54000	1436\\
56000	1396.6\\
58000	1445.3\\
60000	1410.3\\
62000	1445.1\\
64000	1431.7\\
66000	1446.7\\
68000	1455.3\\
70000	1416.1\\
72000	1397.6\\
74000	1430\\
76000	1409.6\\
78000	1429.2\\
80000	1420.5\\
84000	1435.1\\
86000	1418.5\\
88000	1441.7\\
90000	1427.7\\
92000	1433.2\\
94000	1418\\
96000	1427.1\\
98000	1411.5\\
1e+05	1413.1\\
};
\addplot [color=mycolor82, forget plot]
  table[row sep=crcr]{%
0	4395.3\\
2000	2091.9\\
4000	1658.7\\
6000	1674.3\\
10000	1672.5\\
12000	1688.9\\
16000	1662\\
18000	1645.6\\
20000	1670.7\\
22000	1651\\
24000	1683.1\\
26000	1666.8\\
28000	1687.7\\
30000	1700.4\\
32000	1692.2\\
34000	1690.7\\
36000	1665\\
38000	1683.7\\
40000	1690\\
42000	1665.6\\
44000	1654.8\\
46000	1680.5\\
50000	1687.3\\
52000	1641.4\\
54000	1649.7\\
56000	1672.3\\
64000	1663.4\\
66000	1648.1\\
68000	1682.3\\
70000	1662.9\\
72000	1663.4\\
78000	1718.1\\
80000	1705.8\\
86000	1644.1\\
88000	1662.9\\
90000	1657.1\\
92000	1686.9\\
94000	1650.6\\
96000	1643.3\\
98000	1658.7\\
1e+05	1682.2\\
};
\addplot [color=mycolor83, forget plot]
  table[row sep=crcr]{%
0	4395.3\\
2000	2142.2\\
4000	1674.3\\
6000	1667.3\\
8000	1708.1\\
10000	1679.8\\
12000	1665.1\\
14000	1696.8\\
16000	1684.2\\
18000	1691.7\\
22000	1699.2\\
24000	1687.1\\
26000	1685.6\\
28000	1710.5\\
30000	1711\\
34000	1691.1\\
36000	1656.6\\
38000	1696.6\\
40000	1704.6\\
42000	1694.5\\
44000	1705.2\\
46000	1722.6\\
48000	1699.3\\
52000	1674.6\\
54000	1704\\
56000	1684.1\\
58000	1697.7\\
60000	1697.5\\
64000	1722.5\\
66000	1684.5\\
68000	1683.7\\
70000	1695.6\\
72000	1719.3\\
74000	1715.3\\
76000	1695.4\\
78000	1698.6\\
80000	1725.7\\
82000	1705.4\\
84000	1649.8\\
88000	1686.5\\
90000	1697.1\\
92000	1696.6\\
94000	1703.3\\
96000	1716.2\\
98000	1667.7\\
1e+05	1697.7\\
};
\addplot [color=mycolor83, forget plot]
  table[row sep=crcr]{%
0	4395.3\\
2000	1723.2\\
4000	1273.9\\
6000	1245.8\\
8000	1230.2\\
10000	1254.7\\
12000	1237.6\\
14000	1244.7\\
16000	1243.6\\
18000	1229.6\\
20000	1227.6\\
22000	1217.7\\
24000	1231.7\\
26000	1236.9\\
28000	1252.8\\
30000	1239.1\\
32000	1262.4\\
34000	1213.1\\
36000	1241.8\\
38000	1230\\
40000	1253\\
42000	1227.5\\
44000	1223.3\\
46000	1226.8\\
48000	1243.2\\
50000	1244.6\\
52000	1227.2\\
54000	1233.2\\
58000	1233.7\\
60000	1238\\
62000	1224.1\\
64000	1235\\
66000	1221.4\\
68000	1265.5\\
70000	1229.5\\
72000	1238.5\\
74000	1239.4\\
76000	1268.7\\
78000	1245\\
80000	1236.6\\
82000	1247.7\\
84000	1236.5\\
86000	1249.5\\
88000	1223.9\\
90000	1239.9\\
92000	1225.8\\
94000	1224.9\\
96000	1217.3\\
98000	1249.2\\
1e+05	1259\\
};
\addplot [color=mycolor84, forget plot]
  table[row sep=crcr]{%
0	4395.3\\
2000	1794\\
4000	1276.1\\
6000	1246.2\\
8000	1252.2\\
10000	1238.1\\
12000	1206.3\\
14000	1275.5\\
16000	1290.7\\
18000	1234.8\\
20000	1240.8\\
22000	1219.3\\
26000	1269.1\\
28000	1245.5\\
30000	1253.4\\
32000	1232.7\\
34000	1277.3\\
36000	1259.5\\
40000	1253.4\\
42000	1268.1\\
44000	1249.4\\
46000	1239.8\\
50000	1236.6\\
52000	1260.8\\
54000	1237\\
56000	1260.9\\
60000	1231.9\\
62000	1236.7\\
64000	1247.7\\
66000	1296.6\\
68000	1256.5\\
70000	1279.7\\
72000	1246.3\\
74000	1268.4\\
76000	1282.4\\
78000	1232.3\\
82000	1264.2\\
84000	1269.9\\
86000	1248.3\\
88000	1254.6\\
90000	1209\\
92000	1249.3\\
94000	1265.5\\
96000	1239.4\\
98000	1243.9\\
1e+05	1277.2\\
};
\addplot [color=mycolor85, forget plot]
  table[row sep=crcr]{%
0	4395.3\\
2000	1900.1\\
4000	1370.8\\
6000	1347\\
8000	1358.2\\
10000	1382.1\\
12000	1357.6\\
14000	1357.4\\
16000	1365.2\\
18000	1381.5\\
20000	1383.7\\
22000	1368.1\\
24000	1373.6\\
26000	1367.3\\
28000	1352.5\\
30000	1360.2\\
32000	1359.8\\
34000	1405.7\\
36000	1361.9\\
38000	1326.3\\
40000	1374.9\\
42000	1365.8\\
44000	1353\\
46000	1367.7\\
48000	1349.4\\
50000	1377.4\\
52000	1363.7\\
54000	1345\\
58000	1353.7\\
60000	1411.5\\
62000	1370\\
64000	1334\\
66000	1343.9\\
68000	1380.1\\
70000	1366.6\\
72000	1394.4\\
74000	1384.8\\
76000	1347.7\\
78000	1356.1\\
80000	1426.2\\
82000	1404.4\\
84000	1350.7\\
86000	1407.4\\
88000	1392.9\\
90000	1348.5\\
92000	1346.6\\
94000	1382.1\\
96000	1360.9\\
98000	1354.8\\
1e+05	1354.3\\
};
\addplot [color=mycolor86, forget plot]
  table[row sep=crcr]{%
0	4395.3\\
2000	1740\\
4000	1152.9\\
6000	1138.8\\
8000	1165\\
10000	1177.1\\
12000	1139.5\\
14000	1156.6\\
16000	1127.3\\
18000	1153.4\\
20000	1131.1\\
22000	1137.1\\
24000	1116.8\\
26000	1153\\
28000	1182.6\\
30000	1163.2\\
32000	1175\\
34000	1135.9\\
36000	1170.6\\
38000	1155.9\\
40000	1149.7\\
42000	1176.1\\
44000	1120\\
46000	1146.6\\
48000	1139.4\\
50000	1166.2\\
52000	1155.8\\
56000	1147.7\\
58000	1123.5\\
60000	1125.5\\
62000	1136.5\\
64000	1141.1\\
66000	1159.3\\
68000	1164.1\\
70000	1127.5\\
72000	1148.3\\
74000	1154.6\\
76000	1152.1\\
78000	1161.6\\
80000	1151\\
82000	1130.8\\
84000	1153.6\\
86000	1112.1\\
88000	1172.4\\
90000	1150.5\\
92000	1118.2\\
94000	1128.5\\
96000	1154.5\\
98000	1163.2\\
1e+05	1154.4\\
};
\addplot [color=mycolor87, forget plot]
  table[row sep=crcr]{%
0	4395.3\\
2000	1566.2\\
4000	969.62\\
6000	940.46\\
8000	941.61\\
10000	992.86\\
12000	994.52\\
14000	975.36\\
16000	983.34\\
18000	1002.7\\
20000	1016.9\\
22000	1005.1\\
24000	965.86\\
26000	962.39\\
28000	971.71\\
30000	994.9\\
32000	928.62\\
36000	958.06\\
38000	1017\\
40000	957.94\\
42000	970.37\\
44000	963.01\\
46000	1019.5\\
48000	1010.8\\
50000	1005.1\\
52000	991.03\\
54000	962.72\\
56000	924.11\\
58000	1022.7\\
60000	922.61\\
62000	958.58\\
64000	1008\\
66000	997.33\\
68000	967.54\\
70000	924.5\\
72000	985.69\\
74000	966.98\\
76000	970.06\\
78000	930.16\\
80000	961.64\\
82000	933.77\\
86000	930.18\\
88000	952.2\\
90000	959.24\\
92000	981.61\\
94000	983.74\\
96000	942.96\\
1e+05	994.29\\
};
\addplot [color=mycolor88, forget plot]
  table[row sep=crcr]{%
0	4395.3\\
2000	1810.9\\
4000	1159.1\\
6000	1180.3\\
8000	1180\\
10000	1151.2\\
12000	1166.9\\
14000	1155.6\\
16000	1166.8\\
18000	1168.7\\
20000	1127.5\\
22000	1182.6\\
24000	1139.4\\
26000	1139\\
28000	1170\\
30000	1124.5\\
32000	1123.8\\
34000	1158.1\\
36000	1160.4\\
38000	1138.1\\
40000	1138.6\\
42000	1144.7\\
44000	1133.6\\
46000	1163.4\\
48000	1168\\
50000	1159\\
52000	1126\\
54000	1150.5\\
56000	1142.1\\
58000	1179.7\\
62000	1135.8\\
64000	1117.4\\
68000	1162.2\\
70000	1171.6\\
72000	1160\\
74000	1184.7\\
76000	1160.8\\
78000	1156.8\\
80000	1172.5\\
82000	1158.2\\
84000	1169.4\\
86000	1177.3\\
88000	1143.5\\
90000	1129.5\\
92000	1149.9\\
94000	1174.4\\
96000	1141.3\\
98000	1157.3\\
1e+05	1157.4\\
};
\addplot [color=mycolor89, forget plot]
  table[row sep=crcr]{%
0	4395.3\\
2000	1501\\
4000	860.13\\
6000	851.26\\
8000	825.42\\
10000	816.87\\
12000	822.09\\
14000	817.03\\
16000	825.27\\
18000	846.21\\
20000	820.95\\
22000	832.36\\
24000	817.44\\
26000	807.4\\
28000	812.53\\
30000	827.99\\
32000	816.42\\
34000	840.61\\
36000	815.9\\
38000	828.14\\
40000	830.99\\
42000	815.72\\
46000	830.15\\
48000	812.41\\
50000	848.25\\
52000	829.31\\
54000	825.96\\
56000	853.71\\
58000	826.99\\
60000	846.07\\
62000	838.4\\
64000	836.9\\
66000	847.71\\
68000	835.81\\
70000	814.03\\
72000	831.16\\
74000	822.05\\
76000	825.03\\
78000	806.68\\
80000	831.61\\
82000	819.31\\
84000	803.61\\
86000	832.65\\
88000	826.1\\
90000	837.35\\
92000	851.38\\
94000	845.27\\
96000	817.13\\
98000	810.44\\
1e+05	801.49\\
};
\addplot [color=mycolor90, forget plot]
  table[row sep=crcr]{%
0	4395.3\\
2000	1452.5\\
4000	815.08\\
6000	790.26\\
8000	802.13\\
10000	811.48\\
12000	828.4\\
14000	782.91\\
16000	802.33\\
18000	784.87\\
20000	798.93\\
22000	791.4\\
24000	790.88\\
26000	794.28\\
28000	778.12\\
30000	798.42\\
32000	814.4\\
36000	810.07\\
38000	790.14\\
40000	821.5\\
42000	817.19\\
46000	792.2\\
48000	823.24\\
50000	795.27\\
52000	797.13\\
54000	802.15\\
56000	797.44\\
58000	784.53\\
60000	822.49\\
62000	799.78\\
64000	792.26\\
66000	811.56\\
68000	798.25\\
70000	788.06\\
72000	797.44\\
74000	795.56\\
78000	819.97\\
80000	807.05\\
82000	810.73\\
84000	819.46\\
86000	800.88\\
88000	811.65\\
90000	801.11\\
92000	812.26\\
94000	779.15\\
96000	801.51\\
98000	837.13\\
1e+05	787.51\\
};
\addplot [color=mycolor91, forget plot]
  table[row sep=crcr]{%
0	4395.3\\
2000	1349.4\\
4000	718.14\\
6000	676.73\\
8000	703.35\\
10000	708.54\\
12000	685.79\\
16000	684.13\\
18000	680.81\\
20000	698.4\\
22000	660.78\\
24000	681.92\\
26000	695.38\\
28000	695.32\\
30000	706.27\\
32000	703.55\\
34000	697.9\\
36000	709.64\\
38000	674.04\\
42000	685.54\\
44000	669.61\\
46000	687.68\\
48000	697\\
50000	679.94\\
52000	715.04\\
54000	684.79\\
56000	674.37\\
58000	701.15\\
60000	700.8\\
62000	682.85\\
64000	702.87\\
66000	720.52\\
70000	677.86\\
72000	706.54\\
74000	708.19\\
76000	700.96\\
78000	685.15\\
80000	689.67\\
82000	673.28\\
84000	677.79\\
86000	687.78\\
88000	682.25\\
90000	708.08\\
92000	679.61\\
94000	689.27\\
96000	682.28\\
98000	712.43\\
1e+05	681.9\\
};
\addplot [color=mycolor92, forget plot]
  table[row sep=crcr]{%
0	4395.3\\
2000	1618.9\\
4000	880.44\\
6000	889.63\\
8000	901.96\\
10000	901.78\\
12000	857.72\\
14000	860.52\\
16000	849.94\\
18000	852.5\\
20000	889.42\\
22000	883.91\\
24000	865.39\\
26000	865.22\\
28000	849.46\\
30000	851.26\\
32000	838.51\\
34000	890.64\\
36000	871.87\\
38000	883.06\\
40000	864.02\\
42000	861.42\\
44000	877.69\\
46000	864.17\\
48000	837.85\\
50000	834.58\\
52000	860.52\\
54000	894.48\\
56000	873.39\\
58000	866.24\\
60000	893.46\\
62000	865.93\\
64000	873.71\\
66000	878.83\\
70000	834.28\\
72000	899.57\\
74000	875.95\\
76000	864.08\\
78000	878.44\\
80000	883.25\\
82000	876.05\\
84000	857.91\\
86000	884.94\\
88000	851.9\\
90000	874.68\\
92000	876.96\\
94000	870.77\\
96000	884.28\\
98000	872.09\\
1e+05	857\\
};
\addplot [color=mycolor93, forget plot]
  table[row sep=crcr]{%
0	4395.3\\
2000	1530.8\\
4000	757.52\\
6000	745.65\\
8000	763.36\\
10000	701.95\\
12000	694.67\\
14000	708.96\\
16000	741.78\\
18000	735.68\\
20000	708.9\\
22000	705.09\\
24000	708.54\\
26000	740.45\\
28000	760.84\\
30000	718.3\\
32000	735.14\\
34000	720.43\\
36000	743.39\\
38000	716.29\\
40000	705.01\\
42000	710.96\\
44000	711.38\\
46000	682.26\\
48000	678.33\\
50000	713.87\\
52000	723.01\\
54000	698.32\\
56000	740\\
58000	725.27\\
60000	712.97\\
62000	731.08\\
64000	738.9\\
66000	722.87\\
68000	719.52\\
72000	705.05\\
74000	709.98\\
76000	737.03\\
78000	748.82\\
80000	721.21\\
82000	793.85\\
84000	723.59\\
86000	725.37\\
88000	710.59\\
90000	705.77\\
92000	755.64\\
94000	725.27\\
96000	729.19\\
98000	717.79\\
1e+05	697.03\\
};
\addplot [color=mycolor94, forget plot]
  table[row sep=crcr]{%
0	4395.3\\
2000	1518.5\\
4000	708.58\\
6000	686\\
10000	690.77\\
12000	685.53\\
16000	680.82\\
18000	717.85\\
22000	660.58\\
24000	698.97\\
26000	714.58\\
28000	702.76\\
30000	664.26\\
32000	671.41\\
34000	671.86\\
36000	691.64\\
38000	714.69\\
40000	690.01\\
42000	696.92\\
44000	678.85\\
46000	685.32\\
48000	698.21\\
50000	668.58\\
52000	676.72\\
54000	705.61\\
56000	702.62\\
58000	667.95\\
60000	691.8\\
62000	653.57\\
68000	690.08\\
70000	660.64\\
72000	670.11\\
74000	684.39\\
76000	690.15\\
78000	679.18\\
80000	680.41\\
82000	684.98\\
84000	666.47\\
86000	670.5\\
88000	692.04\\
90000	698.28\\
92000	692.61\\
94000	701.34\\
96000	688.99\\
98000	717.84\\
1e+05	709.96\\
};
\addplot [color=mycolor95, forget plot]
  table[row sep=crcr]{%
0	4395.3\\
2000	1241.4\\
4000	420.41\\
6000	403.6\\
8000	397.29\\
10000	421.05\\
12000	408.74\\
14000	412.61\\
16000	402.1\\
18000	400.35\\
20000	423.33\\
22000	416.9\\
24000	425.8\\
26000	418.07\\
28000	390.89\\
30000	390.66\\
32000	396.37\\
34000	395.52\\
36000	393.02\\
38000	392.76\\
40000	412.77\\
42000	407.68\\
44000	385.01\\
46000	387.04\\
48000	404.78\\
50000	404.62\\
52000	410.27\\
54000	412.52\\
56000	410.83\\
58000	393.07\\
60000	421.82\\
64000	385.92\\
66000	405.78\\
68000	414.76\\
70000	397.36\\
72000	409.36\\
74000	433.83\\
76000	406.43\\
78000	401.79\\
80000	384.78\\
82000	410.52\\
84000	397.63\\
86000	396.73\\
88000	407.24\\
90000	383.04\\
92000	400.62\\
94000	417.66\\
96000	399.45\\
98000	402.29\\
1e+05	402.27\\
};
\addplot [color=mycolor96, forget plot]
  table[row sep=crcr]{%
0	4395.3\\
2000	1331.6\\
4000	472.23\\
6000	422.36\\
8000	415.38\\
10000	405.61\\
12000	434.31\\
14000	448.46\\
16000	446.07\\
18000	442.2\\
20000	465.23\\
22000	431.18\\
24000	439.87\\
26000	443.72\\
28000	419.02\\
30000	412.28\\
32000	444.51\\
34000	433.36\\
36000	459.71\\
38000	407.54\\
40000	399.32\\
42000	432.8\\
44000	438.13\\
46000	416.25\\
48000	429.02\\
50000	428.46\\
52000	442.34\\
54000	443.79\\
56000	449.98\\
58000	485.5\\
60000	468.27\\
62000	417.96\\
64000	428.51\\
66000	441.74\\
68000	468.51\\
70000	443.34\\
74000	420.87\\
76000	417.6\\
78000	417.49\\
80000	459.29\\
82000	417.47\\
84000	422.43\\
86000	432.89\\
88000	419.45\\
90000	450.63\\
92000	440.19\\
94000	436.26\\
96000	429.71\\
98000	447.55\\
1e+05	415.85\\
};
\addplot [color=mycolor97, forget plot]
  table[row sep=crcr]{%
0	4395.3\\
2000	1223.6\\
4000	326.56\\
6000	302.75\\
8000	311.15\\
10000	298.85\\
12000	320.56\\
14000	314.83\\
16000	308.28\\
18000	297.05\\
20000	297.81\\
22000	285.55\\
24000	305.41\\
26000	300.1\\
28000	310.43\\
30000	310.72\\
32000	313.92\\
34000	296.4\\
36000	287.52\\
38000	295.05\\
40000	301.31\\
42000	283.57\\
44000	279.26\\
46000	293.37\\
48000	294.49\\
50000	294.07\\
52000	311.33\\
54000	308.24\\
56000	291.71\\
58000	291.22\\
60000	303.96\\
62000	312.86\\
64000	289\\
66000	269.92\\
68000	307.88\\
70000	289.26\\
72000	310.46\\
74000	267.24\\
76000	302.06\\
78000	330.31\\
80000	337.27\\
82000	311.07\\
84000	307.69\\
86000	295.45\\
88000	307.82\\
90000	321.8\\
92000	282.4\\
94000	312.79\\
96000	289.77\\
98000	276.99\\
1e+05	275.47\\
};
\addplot [color=mycolor98]
  table[row sep=crcr]{%
0	4395.3\\
2000	1466.9\\
4000	467.95\\
6000	411.05\\
8000	408\\
10000	424.4\\
12000	408.31\\
14000	418.31\\
16000	437.28\\
18000	406.48\\
22000	391.43\\
24000	386.18\\
26000	401.62\\
28000	409.45\\
30000	395.45\\
32000	377.78\\
34000	397.27\\
36000	419.94\\
38000	431\\
40000	391.34\\
42000	397.43\\
48000	389.6\\
50000	432.37\\
52000	429.13\\
54000	432.97\\
56000	403.11\\
58000	404.58\\
60000	419.02\\
62000	412.27\\
64000	418.7\\
66000	419.37\\
68000	391.73\\
70000	410.13\\
72000	422.03\\
74000	386.18\\
76000	394.76\\
78000	385.21\\
80000	414.34\\
82000	416.76\\
84000	396.27\\
86000	387.23\\
88000	382.02\\
90000	384.89\\
92000	381.06\\
94000	407.81\\
96000	417.93\\
98000	416.57\\
1e+05	411.85\\
};
\addlegendentry{Classe k=10}

            % \input{../../../../.tex/20/KS/SizeEvolutionsfig8.tex}
            \legend{}
        \end{groupplot}

        \begin{groupplot}[
            group style={
                group name=Row5,
                plotLastGroup,
                group size=2 by 1,
            },
            plotRow,
            row5Keys,
            plotLegendSE,
            %
            xmin=0,xmax=1e+05,
            ymin=100,ymax=1.6196e+05,
        ]
            \nextgroupplot[%
                at={($(Row4 c1r1.south west)-(0,\yPlotSep em)-(0,\plotHeight cm)$)},
                % xmin=0,xmax=1e+05,
                % ymin=100,ymax=1.5829e+05,
            ] % This file was created by matlab2tikz.
%
\definecolor{mycolor1}{rgb}{0.00146,0.00047,0.01387}%
\definecolor{mycolor2}{rgb}{0.01166,0.00942,0.06346}%
\definecolor{mycolor3}{rgb}{0.02943,0.02150,0.11462}%
\definecolor{mycolor4}{rgb}{0.06134,0.03659,0.17764}%
\definecolor{mycolor5}{rgb}{0.08741,0.04456,0.22481}%
\definecolor{mycolor6}{rgb}{0.12291,0.04754,0.28162}%
\definecolor{mycolor7}{rgb}{0.14907,0.04547,0.31709}%
\definecolor{mycolor8}{rgb}{0.18343,0.04033,0.35497}%
\definecolor{mycolor9}{rgb}{0.21795,0.03662,0.38352}%
\definecolor{mycolor10}{rgb}{0.24497,0.03705,0.40001}%
\definecolor{mycolor11}{rgb}{0.27135,0.04092,0.41198}%
\definecolor{mycolor12}{rgb}{0.29076,0.04564,0.41864}%
\definecolor{mycolor13}{rgb}{0.31628,0.05349,0.42512}%
\definecolor{mycolor14}{rgb}{0.33522,0.06006,0.42852}%
\definecolor{mycolor15}{rgb}{0.36028,0.06925,0.43150}%
\definecolor{mycolor16}{rgb}{0.37900,0.07625,0.43272}%
\definecolor{mycolor17}{rgb}{0.39767,0.08326,0.43318}%
\definecolor{mycolor18}{rgb}{0.41633,0.09020,0.43294}%
\definecolor{mycolor19}{rgb}{0.42877,0.09479,0.43241}%
\definecolor{mycolor20}{rgb}{0.44743,0.10160,0.43108}%
\definecolor{mycolor21}{rgb}{0.46610,0.10832,0.42912}%
\definecolor{mycolor22}{rgb}{0.47856,0.11276,0.42747}%
\definecolor{mycolor23}{rgb}{0.49726,0.11938,0.42449}%
\definecolor{mycolor24}{rgb}{0.50973,0.12377,0.42216}%
\definecolor{mycolor25}{rgb}{0.52221,0.12815,0.41955}%
\definecolor{mycolor26}{rgb}{0.53468,0.13253,0.41667}%
\definecolor{mycolor27}{rgb}{0.55339,0.13913,0.41183}%
\definecolor{mycolor28}{rgb}{0.56585,0.14357,0.40826}%
\definecolor{mycolor29}{rgb}{0.57830,0.14804,0.40441}%
\definecolor{mycolor30}{rgb}{0.59073,0.15256,0.40029}%
\definecolor{mycolor31}{rgb}{0.60314,0.15715,0.39589}%
\definecolor{mycolor32}{rgb}{0.61551,0.16182,0.39122}%
\definecolor{mycolor33}{rgb}{0.62169,0.16418,0.38878}%
\definecolor{mycolor34}{rgb}{0.63400,0.16899,0.38370}%
\definecolor{mycolor35}{rgb}{0.64626,0.17391,0.37836}%
\definecolor{mycolor36}{rgb}{0.65846,0.17896,0.37275}%
\definecolor{mycolor37}{rgb}{0.66454,0.18154,0.36985}%
\definecolor{mycolor38}{rgb}{0.67664,0.18681,0.36385}%
\definecolor{mycolor39}{rgb}{0.68865,0.19224,0.35760}%
\definecolor{mycolor40}{rgb}{0.69463,0.19502,0.35439}%
\definecolor{mycolor41}{rgb}{0.70650,0.20073,0.34778}%
\definecolor{mycolor42}{rgb}{0.71240,0.20366,0.34438}%
\definecolor{mycolor43}{rgb}{0.72410,0.20967,0.33742}%
\definecolor{mycolor44}{rgb}{0.72991,0.21276,0.33386}%
\definecolor{mycolor45}{rgb}{0.74142,0.21911,0.32658}%
\definecolor{mycolor46}{rgb}{0.74713,0.22238,0.32286}%
\definecolor{mycolor47}{rgb}{0.75279,0.22571,0.31909}%
\definecolor{mycolor48}{rgb}{0.76401,0.23255,0.31140}%
\definecolor{mycolor49}{rgb}{0.76956,0.23608,0.30749}%
\definecolor{mycolor50}{rgb}{0.77506,0.23967,0.30353}%
\definecolor{mycolor51}{rgb}{0.79129,0.25086,0.29139}%
\definecolor{mycolor52}{rgb}{0.79661,0.25473,0.28726}%
\definecolor{mycolor53}{rgb}{0.80187,0.25867,0.28310}%
\definecolor{mycolor54}{rgb}{0.81224,0.26679,0.27466}%
\definecolor{mycolor55}{rgb}{0.81734,0.27095,0.27039}%
\definecolor{mycolor56}{rgb}{0.82737,0.27952,0.26175}%
\definecolor{mycolor57}{rgb}{0.83230,0.28391,0.25738}%
\definecolor{mycolor58}{rgb}{0.83717,0.28839,0.25299}%
\definecolor{mycolor59}{rgb}{0.84671,0.29756,0.24411}%
\definecolor{mycolor60}{rgb}{0.85138,0.30226,0.23964}%
\definecolor{mycolor61}{rgb}{0.85599,0.30704,0.23513}%
\definecolor{mycolor62}{rgb}{0.86053,0.31189,0.23061}%
\definecolor{mycolor63}{rgb}{0.87374,0.32691,0.21689}%
\definecolor{mycolor64}{rgb}{0.87800,0.33206,0.21227}%
\definecolor{mycolor65}{rgb}{0.89034,0.34796,0.19829}%
\definecolor{mycolor66}{rgb}{0.89819,0.35891,0.18886}%
\definecolor{mycolor67}{rgb}{0.90200,0.36449,0.18412}%
\definecolor{mycolor68}{rgb}{0.90573,0.37014,0.17935}%
\definecolor{mycolor69}{rgb}{0.90939,0.37586,0.17456}%
\definecolor{mycolor70}{rgb}{0.91297,0.38164,0.16975}%
\definecolor{mycolor71}{rgb}{0.91646,0.38748,0.16492}%
\definecolor{mycolor72}{rgb}{0.91988,0.39339,0.16007}%
\definecolor{mycolor73}{rgb}{0.92322,0.39936,0.15519}%
\definecolor{mycolor74}{rgb}{0.92647,0.40539,0.15029}%
\definecolor{mycolor75}{rgb}{0.92964,0.41148,0.14537}%
\definecolor{mycolor76}{rgb}{0.93274,0.41763,0.14042}%
\definecolor{mycolor77}{rgb}{0.93575,0.42383,0.13544}%
\definecolor{mycolor78}{rgb}{0.93868,0.43009,0.13044}%
\definecolor{mycolor79}{rgb}{0.94152,0.43640,0.12541}%
\definecolor{mycolor80}{rgb}{0.94429,0.44277,0.12035}%
\definecolor{mycolor81}{rgb}{0.94956,0.45566,0.11016}%
\definecolor{mycolor82}{rgb}{0.95208,0.46218,0.10503}%
\definecolor{mycolor83}{rgb}{0.96129,0.48872,0.08429}%
\definecolor{mycolor84}{rgb}{0.96540,0.50225,0.07386}%
\definecolor{mycolor85}{rgb}{0.96732,0.50908,0.06866}%
\definecolor{mycolor86}{rgb}{0.97259,0.52980,0.05332}%
\definecolor{mycolor87}{rgb}{0.97568,0.54380,0.04362}%
\definecolor{mycolor88}{rgb}{0.98082,0.57221,0.02851}%
\definecolor{mycolor89}{rgb}{0.98189,0.57939,0.02625}%
\definecolor{mycolor90}{rgb}{0.98378,0.59385,0.02377}%
\definecolor{mycolor91}{rgb}{0.98532,0.60842,0.02420}%
\definecolor{mycolor92}{rgb}{0.98696,0.63048,0.03091}%
\definecolor{mycolor93}{rgb}{0.98793,0.66025,0.05175}%
\definecolor{mycolor94}{rgb}{0.98771,0.68281,0.07249}%
\definecolor{mycolor95}{rgb}{0.97750,0.78226,0.18592}%
\definecolor{mycolor96}{rgb}{0.96624,0.83619,0.26153}%
\definecolor{mycolor97}{rgb}{0.96252,0.85148,0.28555}%
\definecolor{mycolor98}{rgb}{0.98836,0.99836,0.64492}%
%
\addplot [color=mycolor1]
  table[row sep=crcr]{%
0	4395.3\\
2000	1.2183e+05\\
4000	1.5449e+05\\
6000	1.5023e+05\\
10000	1.5729e+05\\
12000	1.5314e+05\\
14000	1.5712e+05\\
16000	1.5247e+05\\
18000	1.5347e+05\\
20000	1.5329e+05\\
22000	1.5248e+05\\
24000	1.5638e+05\\
26000	1.5682e+05\\
28000	1.5641e+05\\
30000	1.5516e+05\\
32000	1.5261e+05\\
34000	1.5454e+05\\
36000	1.5222e+05\\
38000	1.5168e+05\\
40000	1.5475e+05\\
42000	1.5471e+05\\
44000	1.5684e+05\\
48000	1.5365e+05\\
50000	1.5664e+05\\
52000	1.5741e+05\\
54000	1.5683e+05\\
56000	1.5504e+05\\
58000	1.5513e+05\\
60000	1.5727e+05\\
62000	1.5718e+05\\
64000	1.55e+05\\
74000	1.5312e+05\\
76000	1.5087e+05\\
78000	1.5441e+05\\
80000	1.5425e+05\\
82000	1.5502e+05\\
84000	1.5829e+05\\
88000	1.5485e+05\\
92000	1.5437e+05\\
94000	1.4949e+05\\
96000	1.543e+05\\
98000	1.5647e+05\\
1e+05	1.5727e+05\\
};
\addlegendentry{Classe k=283}

\addplot [color=mycolor2, forget plot]
  table[row sep=crcr]{%
0	4395.3\\
2000	25558\\
4000	27088\\
6000	25935\\
8000	26502\\
10000	26862\\
12000	26609\\
14000	26484\\
16000	27025\\
18000	26263\\
20000	26401\\
22000	26095\\
24000	26875\\
28000	26221\\
30000	25953\\
32000	26579\\
34000	26556\\
36000	26789\\
38000	27594\\
40000	26231\\
42000	25623\\
44000	25147\\
46000	26311\\
48000	26052\\
50000	26331\\
52000	25860\\
56000	26510\\
58000	26144\\
60000	26274\\
62000	26016\\
64000	26086\\
66000	25560\\
68000	27202\\
70000	26740\\
72000	26157\\
74000	26218\\
76000	25558\\
78000	25919\\
80000	25424\\
82000	26451\\
84000	26647\\
86000	26113\\
88000	26616\\
90000	26870\\
92000	25943\\
94000	26073\\
96000	26439\\
98000	25933\\
1e+05	26481\\
};
\addplot [color=mycolor3, forget plot]
  table[row sep=crcr]{%
0	4395.3\\
2000	38508\\
4000	39695\\
6000	40436\\
8000	39062\\
10000	39251\\
12000	41305\\
14000	40757\\
16000	40668\\
18000	41279\\
20000	41622\\
22000	40576\\
24000	39365\\
26000	40485\\
28000	40314\\
30000	39739\\
32000	40814\\
34000	41098\\
38000	39195\\
40000	39910\\
42000	39861\\
44000	40898\\
46000	39385\\
48000	39357\\
50000	38994\\
52000	40542\\
54000	39941\\
56000	41388\\
58000	40859\\
60000	40588\\
62000	40700\\
64000	41110\\
66000	40462\\
68000	40323\\
70000	39469\\
72000	40200\\
74000	41483\\
76000	39493\\
78000	40294\\
80000	40202\\
82000	40758\\
84000	41139\\
86000	41110\\
88000	39905\\
90000	39232\\
92000	41213\\
94000	39557\\
96000	40240\\
98000	40511\\
1e+05	40136\\
};
\addplot [color=mycolor4, forget plot]
  table[row sep=crcr]{%
0	4395.3\\
2000	14825\\
4000	14586\\
6000	14811\\
8000	14487\\
10000	14898\\
12000	14142\\
14000	14268\\
16000	14677\\
18000	14197\\
20000	14721\\
22000	14587\\
24000	14143\\
26000	14240\\
28000	14845\\
30000	14558\\
32000	14631\\
34000	14752\\
36000	14396\\
40000	14609\\
42000	14291\\
44000	14671\\
46000	14802\\
48000	14370\\
50000	14247\\
52000	14347\\
54000	14401\\
56000	14895\\
58000	14388\\
60000	14747\\
62000	14090\\
64000	14689\\
66000	14752\\
68000	14323\\
70000	15026\\
72000	13859\\
74000	14370\\
76000	14330\\
78000	14816\\
80000	14472\\
82000	14320\\
84000	14338\\
86000	14474\\
88000	14547\\
90000	14926\\
92000	14895\\
94000	14820\\
98000	14154\\
1e+05	14523\\
};
\addplot [color=mycolor5, forget plot]
  table[row sep=crcr]{%
0	4395.3\\
2000	56972\\
4000	66584\\
6000	64386\\
8000	62663\\
10000	63418\\
12000	62660\\
16000	64607\\
18000	65186\\
20000	63901\\
22000	65987\\
24000	64923\\
26000	63515\\
28000	63478\\
30000	65492\\
32000	65027\\
34000	60903\\
36000	63343\\
38000	64629\\
40000	64278\\
46000	64127\\
48000	64325\\
50000	65549\\
54000	63829\\
56000	64243\\
58000	62850\\
62000	63535\\
64000	67241\\
66000	64317\\
72000	65878\\
74000	64985\\
76000	66497\\
78000	64483\\
80000	65726\\
82000	66741\\
84000	63592\\
86000	64159\\
88000	65172\\
90000	65730\\
92000	63517\\
94000	66177\\
96000	64099\\
98000	65111\\
1e+05	66425\\
};
\addplot [color=mycolor6, forget plot]
  table[row sep=crcr]{%
0	4395.3\\
2000	47776\\
4000	55540\\
6000	57736\\
8000	59253\\
10000	55876\\
12000	57349\\
14000	57862\\
16000	57772\\
18000	58099\\
20000	58639\\
22000	55650\\
28000	56856\\
30000	56074\\
32000	58863\\
34000	57679\\
36000	58506\\
38000	57158\\
40000	56438\\
42000	56558\\
44000	57825\\
46000	58070\\
48000	56092\\
50000	55492\\
52000	56695\\
54000	58665\\
56000	58171\\
58000	59730\\
60000	57064\\
62000	58030\\
64000	58795\\
66000	56683\\
68000	56563\\
70000	59328\\
72000	58969\\
74000	56698\\
76000	58805\\
78000	58063\\
80000	58099\\
82000	55482\\
84000	57675\\
86000	56101\\
88000	57687\\
90000	57121\\
92000	59512\\
94000	58590\\
96000	58426\\
98000	59205\\
1e+05	56353\\
};
\addplot [color=mycolor7, forget plot]
  table[row sep=crcr]{%
0	4395.3\\
2000	25305\\
4000	28230\\
6000	30031\\
8000	29450\\
10000	29100\\
12000	28450\\
14000	28894\\
16000	28438\\
18000	29020\\
20000	29088\\
22000	29450\\
24000	28647\\
26000	27642\\
28000	28879\\
30000	29663\\
32000	28479\\
34000	28472\\
36000	29025\\
38000	28893\\
40000	28503\\
42000	28452\\
44000	29018\\
48000	28304\\
50000	28657\\
52000	28874\\
54000	28494\\
56000	29414\\
58000	28370\\
60000	28642\\
62000	27264\\
64000	28101\\
66000	28875\\
68000	28315\\
70000	28563\\
72000	28475\\
74000	28747\\
76000	29511\\
78000	28269\\
80000	29158\\
82000	29244\\
84000	27603\\
86000	28205\\
88000	28511\\
90000	28706\\
92000	28619\\
94000	28901\\
96000	29010\\
98000	28171\\
1e+05	28616\\
};
\addplot [color=mycolor8, forget plot]
  table[row sep=crcr]{%
0	4395.3\\
2000	21482\\
4000	21534\\
6000	20813\\
8000	21146\\
10000	20417\\
12000	20464\\
14000	20340\\
16000	21361\\
18000	20602\\
20000	20473\\
22000	20468\\
24000	20847\\
26000	20958\\
28000	20427\\
30000	20978\\
32000	20558\\
34000	20710\\
36000	20569\\
38000	20754\\
40000	21218\\
42000	20481\\
44000	20338\\
46000	19980\\
48000	20862\\
50000	20707\\
52000	20968\\
54000	20718\\
56000	21698\\
58000	20781\\
60000	20905\\
62000	20840\\
64000	20866\\
66000	21221\\
68000	20645\\
70000	20630\\
72000	20128\\
74000	21556\\
78000	20495\\
80000	20475\\
82000	20750\\
84000	21348\\
86000	20977\\
88000	20831\\
90000	21235\\
92000	20764\\
94000	20989\\
96000	20316\\
98000	20773\\
1e+05	20961\\
};
\addplot [color=mycolor9, forget plot]
  table[row sep=crcr]{%
0	4395.3\\
2000	13251\\
4000	13822\\
6000	14537\\
8000	14486\\
10000	14145\\
12000	14192\\
14000	13822\\
16000	13951\\
18000	14244\\
20000	14249\\
22000	14598\\
24000	14326\\
26000	14489\\
30000	14290\\
32000	13923\\
34000	14141\\
36000	14252\\
38000	14214\\
40000	14119\\
42000	14444\\
44000	14240\\
46000	13620\\
48000	14681\\
50000	14287\\
52000	14234\\
54000	14420\\
56000	13646\\
58000	13729\\
60000	14255\\
62000	14130\\
64000	14180\\
66000	14700\\
68000	14504\\
70000	14604\\
74000	13914\\
76000	14130\\
78000	14046\\
80000	14046\\
82000	14264\\
84000	13945\\
86000	14220\\
88000	14099\\
90000	14118\\
92000	13980\\
94000	14443\\
96000	14269\\
98000	14683\\
1e+05	14290\\
};
\addplot [color=mycolor10, forget plot]
  table[row sep=crcr]{%
0	4395.3\\
2000	19424\\
4000	22170\\
6000	22335\\
10000	23008\\
14000	22211\\
16000	22664\\
18000	22384\\
20000	22614\\
22000	22198\\
24000	22400\\
26000	22955\\
28000	22415\\
30000	22319\\
32000	22464\\
34000	22905\\
36000	21974\\
38000	22926\\
40000	22019\\
42000	22417\\
44000	22256\\
46000	22834\\
48000	22831\\
50000	22529\\
52000	22127\\
54000	22324\\
56000	22257\\
58000	22812\\
60000	22977\\
62000	22921\\
64000	22710\\
66000	22823\\
68000	22607\\
70000	21769\\
72000	22762\\
74000	22318\\
76000	22017\\
78000	22613\\
80000	22393\\
82000	22477\\
84000	22067\\
86000	21946\\
88000	21891\\
90000	21979\\
92000	21456\\
94000	22375\\
96000	22594\\
98000	22508\\
1e+05	22715\\
};
\addplot [color=mycolor11, forget plot]
  table[row sep=crcr]{%
0	4395.3\\
2000	9614.8\\
4000	8855.7\\
6000	8477.7\\
8000	8727.4\\
10000	8705.1\\
12000	8552.9\\
14000	8662.2\\
16000	8528.5\\
18000	8745.6\\
20000	8871.7\\
22000	8672.2\\
24000	8713.9\\
26000	8605.4\\
28000	8534.3\\
30000	8499.4\\
32000	8597.4\\
34000	8886.2\\
36000	8809.8\\
38000	8491.8\\
40000	8495.7\\
42000	8462.7\\
44000	8335\\
46000	8607.7\\
48000	8615.3\\
50000	8827.9\\
52000	8370.5\\
54000	8874.1\\
56000	8762.4\\
58000	8732.1\\
60000	8951\\
62000	8796.4\\
64000	8596.2\\
66000	8730.3\\
68000	8926\\
70000	8946.8\\
72000	8646.1\\
74000	8475.6\\
76000	8890.6\\
80000	8672.9\\
82000	8860.1\\
84000	8652.7\\
86000	8727.3\\
88000	8513.5\\
90000	8914.6\\
92000	8984.1\\
94000	8426.5\\
96000	8813.4\\
98000	8456.8\\
1e+05	8708.4\\
};
\addplot [color=mycolor12, forget plot]
  table[row sep=crcr]{%
0	4395.3\\
2000	13519\\
6000	14956\\
10000	15071\\
12000	14765\\
14000	14948\\
16000	14948\\
18000	15513\\
20000	14221\\
22000	14684\\
24000	14747\\
26000	14921\\
28000	14756\\
30000	14368\\
32000	14697\\
34000	15338\\
36000	14487\\
38000	15375\\
40000	14851\\
42000	15074\\
44000	14871\\
46000	15127\\
48000	15152\\
50000	14990\\
52000	14734\\
54000	14584\\
56000	14897\\
58000	14995\\
60000	14734\\
62000	14680\\
64000	14937\\
66000	14759\\
68000	14698\\
70000	15130\\
72000	14785\\
74000	14502\\
76000	14680\\
78000	15384\\
80000	14559\\
82000	14593\\
84000	14586\\
86000	14893\\
88000	14218\\
90000	14669\\
92000	14689\\
94000	15167\\
96000	14793\\
98000	14794\\
1e+05	14587\\
};
\addplot [color=mycolor13, forget plot]
  table[row sep=crcr]{%
0	4395.3\\
2000	10564\\
4000	11169\\
6000	10968\\
8000	11375\\
10000	11199\\
12000	11117\\
14000	10985\\
18000	11552\\
20000	11164\\
22000	11318\\
24000	10949\\
26000	11218\\
28000	11425\\
30000	10917\\
32000	11385\\
34000	11127\\
36000	11410\\
38000	11287\\
40000	11332\\
42000	11217\\
46000	11839\\
48000	11230\\
50000	10950\\
52000	11024\\
54000	10810\\
56000	11289\\
58000	11293\\
60000	10904\\
62000	11342\\
64000	11148\\
66000	11174\\
68000	11269\\
70000	11269\\
72000	10818\\
74000	11176\\
76000	11325\\
78000	11013\\
80000	11589\\
82000	11400\\
84000	10964\\
86000	11685\\
88000	11269\\
90000	11236\\
92000	11026\\
94000	11133\\
96000	11515\\
98000	11355\\
1e+05	10899\\
};
\addplot [color=mycolor14, forget plot]
  table[row sep=crcr]{%
0	4395.3\\
2000	21162\\
4000	24324\\
6000	25657\\
8000	25335\\
10000	25337\\
12000	24360\\
14000	24765\\
16000	24709\\
18000	24583\\
20000	25283\\
22000	25361\\
24000	26047\\
26000	24364\\
28000	24979\\
30000	24027\\
32000	25237\\
34000	24602\\
36000	25136\\
38000	24106\\
40000	24975\\
42000	26142\\
44000	24974\\
46000	25115\\
48000	25108\\
50000	25350\\
52000	24441\\
54000	24685\\
56000	24227\\
58000	24557\\
60000	24700\\
62000	24089\\
64000	24203\\
66000	23879\\
68000	25869\\
70000	24892\\
72000	24568\\
74000	25296\\
76000	25104\\
78000	24158\\
80000	24441\\
82000	23899\\
84000	25111\\
86000	24098\\
88000	23721\\
90000	24442\\
92000	24230\\
94000	25510\\
96000	24868\\
98000	25246\\
1e+05	25271\\
};
\addplot [color=mycolor15, forget plot]
  table[row sep=crcr]{%
0	4395.3\\
2000	15482\\
4000	17250\\
6000	17505\\
8000	17244\\
10000	17042\\
12000	17508\\
14000	16902\\
16000	17292\\
18000	17344\\
20000	18048\\
22000	17524\\
24000	16583\\
26000	17073\\
28000	16689\\
30000	17171\\
32000	16940\\
34000	17893\\
36000	17305\\
38000	17151\\
40000	17310\\
42000	16998\\
44000	17712\\
46000	17390\\
48000	17905\\
50000	17686\\
52000	17058\\
54000	17369\\
56000	17852\\
58000	17552\\
60000	17758\\
62000	16855\\
64000	17038\\
68000	16890\\
70000	17081\\
72000	17337\\
74000	16833\\
76000	16981\\
78000	17071\\
80000	17103\\
82000	17533\\
84000	17787\\
86000	17054\\
88000	16922\\
90000	16858\\
92000	17033\\
94000	17544\\
96000	16307\\
98000	17619\\
1e+05	16836\\
};
\addplot [color=mycolor16, forget plot]
  table[row sep=crcr]{%
0	4395.3\\
2000	7209.7\\
4000	6971.6\\
6000	6895.2\\
8000	7117.5\\
10000	6941.8\\
12000	7122.8\\
14000	6808.4\\
16000	7087.9\\
18000	6852.9\\
20000	7447.3\\
22000	7098.9\\
24000	7039.3\\
26000	7164.8\\
28000	7024\\
30000	7077.7\\
32000	7239.2\\
34000	7290.5\\
36000	7119.6\\
38000	7126.2\\
40000	6759.4\\
42000	6968.9\\
46000	6805.2\\
48000	7443\\
50000	7045.2\\
52000	7296.5\\
54000	6865.6\\
56000	7340.9\\
58000	6806.6\\
60000	7095.5\\
62000	7072.1\\
64000	7257\\
66000	7223.7\\
68000	7112.6\\
70000	6883.8\\
72000	7234.5\\
74000	7143.2\\
76000	7253.7\\
78000	6929.4\\
80000	6958.3\\
82000	7162\\
88000	6859.5\\
90000	7096\\
92000	7004.1\\
94000	7050.5\\
96000	6871.4\\
98000	6855.8\\
1e+05	6982.9\\
};
\addplot [color=mycolor17, forget plot]
  table[row sep=crcr]{%
0	4395.3\\
2000	14523\\
4000	16174\\
6000	16291\\
8000	15929\\
12000	16610\\
14000	15938\\
16000	15903\\
18000	16342\\
20000	16641\\
22000	15901\\
24000	16675\\
26000	16871\\
28000	15947\\
30000	16360\\
34000	15628\\
36000	16275\\
38000	16324\\
40000	16178\\
42000	17086\\
44000	16392\\
46000	16911\\
48000	16748\\
50000	15461\\
52000	16311\\
54000	16352\\
56000	16304\\
58000	16616\\
60000	15778\\
62000	15844\\
64000	16353\\
66000	16619\\
68000	16288\\
70000	16527\\
72000	15978\\
74000	16211\\
76000	15786\\
78000	16164\\
80000	15994\\
82000	16710\\
84000	15868\\
86000	16356\\
88000	16004\\
90000	16388\\
92000	16602\\
94000	16093\\
96000	16174\\
98000	16538\\
1e+05	15789\\
};
\addplot [color=mycolor18, forget plot]
  table[row sep=crcr]{%
0	4395.3\\
2000	21051\\
4000	24003\\
6000	25899\\
8000	25286\\
10000	25990\\
12000	24906\\
14000	25904\\
16000	25196\\
20000	24003\\
22000	24590\\
24000	24449\\
26000	24853\\
28000	24686\\
30000	26064\\
32000	26020\\
34000	24867\\
36000	25465\\
38000	24926\\
40000	24049\\
42000	25075\\
44000	25184\\
46000	25383\\
50000	25392\\
52000	25539\\
54000	25183\\
56000	25745\\
58000	24704\\
60000	24216\\
62000	26002\\
64000	25256\\
66000	25336\\
68000	24149\\
70000	24864\\
72000	24704\\
74000	26013\\
76000	25650\\
78000	24496\\
80000	25395\\
82000	24362\\
84000	25535\\
86000	24472\\
88000	24143\\
90000	24765\\
92000	24704\\
94000	25563\\
96000	24250\\
98000	24015\\
1e+05	24087\\
};
\addplot [color=mycolor19, forget plot]
  table[row sep=crcr]{%
0	4395.3\\
2000	13609\\
4000	14790\\
6000	15028\\
8000	15006\\
10000	15379\\
12000	15175\\
14000	14751\\
18000	15320\\
20000	14994\\
22000	15080\\
24000	14856\\
26000	15147\\
28000	15262\\
30000	15204\\
32000	14780\\
36000	15202\\
38000	15276\\
40000	14918\\
42000	15233\\
44000	14526\\
46000	15323\\
48000	15080\\
50000	15153\\
52000	14836\\
54000	15924\\
56000	14537\\
58000	15402\\
60000	15237\\
62000	15135\\
64000	14533\\
66000	14853\\
68000	14606\\
70000	15536\\
72000	14777\\
74000	15363\\
76000	15126\\
78000	14837\\
80000	15260\\
82000	14906\\
84000	15002\\
86000	15683\\
88000	15083\\
90000	15176\\
92000	15069\\
94000	15206\\
96000	15445\\
98000	14916\\
1e+05	15065\\
};
\addplot [color=mycolor20, forget plot]
  table[row sep=crcr]{%
0	4395.3\\
2000	17219\\
4000	19157\\
6000	19382\\
8000	19394\\
10000	19072\\
12000	19691\\
14000	18876\\
18000	19942\\
20000	20240\\
24000	19165\\
26000	18950\\
28000	19762\\
30000	20169\\
32000	19317\\
34000	19904\\
36000	20359\\
38000	19285\\
40000	19691\\
42000	19377\\
44000	20667\\
46000	19747\\
48000	20277\\
50000	20162\\
52000	19498\\
54000	19304\\
56000	19829\\
58000	19259\\
60000	19116\\
62000	19777\\
64000	19144\\
66000	20134\\
68000	20051\\
70000	20341\\
72000	19546\\
74000	19171\\
76000	18995\\
78000	19377\\
80000	18998\\
82000	19597\\
84000	20021\\
86000	19159\\
88000	19982\\
90000	19425\\
92000	19128\\
94000	19986\\
96000	19516\\
98000	19465\\
1e+05	19648\\
};
\addplot [color=mycolor21, forget plot]
  table[row sep=crcr]{%
0	4395.3\\
2000	8894\\
4000	9493.1\\
6000	9843.2\\
8000	9762.7\\
12000	9828.2\\
14000	10114\\
16000	9886.5\\
18000	9865.2\\
20000	9950.2\\
22000	9912.3\\
24000	9971\\
26000	9765.9\\
28000	10046\\
30000	9846.7\\
32000	9749.7\\
34000	9701.1\\
36000	9865.3\\
40000	9506.5\\
42000	9941.9\\
44000	9954.8\\
46000	9557.1\\
48000	9913.2\\
50000	9680.3\\
52000	9802.8\\
54000	9685.8\\
56000	9810.3\\
58000	10080\\
60000	9610.4\\
62000	9846.7\\
64000	9668.8\\
66000	9532.4\\
68000	9648.9\\
70000	9639.3\\
72000	9786.1\\
74000	9726.3\\
76000	9756.2\\
78000	9940.8\\
80000	9835.4\\
82000	9571.5\\
84000	9874.9\\
88000	9669.2\\
90000	9871.5\\
92000	9729.6\\
94000	10112\\
96000	10001\\
98000	9578.1\\
1e+05	9956.4\\
};
\addplot [color=mycolor22, forget plot]
  table[row sep=crcr]{%
0	4395.3\\
2000	25849\\
4000	26794\\
6000	26637\\
8000	24590\\
10000	25054\\
12000	25020\\
14000	25583\\
16000	25492\\
18000	25562\\
20000	24624\\
22000	26196\\
24000	26105\\
26000	26543\\
28000	25739\\
30000	25920\\
32000	26210\\
34000	26120\\
36000	25148\\
38000	26247\\
40000	26422\\
42000	25733\\
44000	24677\\
46000	25322\\
48000	25439\\
50000	24683\\
52000	25723\\
54000	26024\\
56000	24916\\
58000	25636\\
60000	26077\\
62000	25195\\
64000	24786\\
66000	25512\\
68000	25334\\
70000	24872\\
72000	25628\\
74000	25350\\
76000	26669\\
78000	25765\\
80000	26767\\
82000	26460\\
84000	24828\\
88000	24867\\
90000	25196\\
92000	25716\\
94000	25790\\
96000	25387\\
98000	24526\\
1e+05	25520\\
};
\addplot [color=mycolor23, forget plot]
  table[row sep=crcr]{%
0	4395.3\\
2000	12002\\
4000	12461\\
6000	12718\\
8000	12709\\
10000	12551\\
12000	12328\\
14000	12563\\
16000	12487\\
18000	12297\\
20000	12340\\
22000	12489\\
24000	12355\\
26000	12474\\
28000	12141\\
30000	12496\\
36000	12340\\
40000	12483\\
42000	12361\\
44000	12432\\
46000	12292\\
48000	12441\\
50000	12756\\
52000	12595\\
54000	12598\\
56000	12512\\
58000	12499\\
60000	12763\\
62000	12489\\
64000	12517\\
66000	12221\\
68000	12253\\
70000	12510\\
72000	12129\\
74000	12540\\
76000	12571\\
78000	12429\\
80000	12341\\
82000	12187\\
84000	12291\\
86000	12848\\
88000	12614\\
90000	12590\\
92000	12665\\
94000	12440\\
96000	12625\\
98000	12289\\
1e+05	12100\\
};
\addplot [color=mycolor24, forget plot]
  table[row sep=crcr]{%
0	4395.3\\
2000	5189.6\\
4000	4677.3\\
6000	4859.7\\
8000	4630\\
10000	4663.4\\
14000	4803.4\\
16000	4740.2\\
18000	4743.6\\
20000	4931.4\\
22000	4750.1\\
24000	4866.5\\
26000	4802.2\\
28000	4673.8\\
30000	4829.2\\
32000	4784.2\\
34000	4612.4\\
36000	4719.3\\
42000	4797.2\\
44000	4754.9\\
46000	4859\\
48000	4786.9\\
50000	4733.5\\
52000	4771.8\\
54000	4924.7\\
56000	4730.1\\
58000	5016.4\\
60000	4989.9\\
62000	4562.6\\
64000	4858.5\\
68000	4748.4\\
70000	5003.6\\
72000	4934.6\\
74000	4772.9\\
76000	4644.3\\
78000	4863.7\\
80000	4740.2\\
82000	4809.1\\
84000	4792.7\\
86000	4801.1\\
88000	4754.8\\
90000	4659.7\\
92000	4675.2\\
94000	4738.6\\
96000	4638.1\\
98000	4621.9\\
1e+05	4756.9\\
};
\addplot [color=mycolor25, forget plot]
  table[row sep=crcr]{%
0	4395.3\\
2000	11077\\
4000	12146\\
6000	12488\\
8000	12369\\
10000	13150\\
12000	12461\\
14000	12701\\
16000	12409\\
18000	12681\\
20000	12466\\
22000	12784\\
24000	12436\\
26000	12647\\
28000	13018\\
30000	12230\\
32000	12863\\
34000	12541\\
36000	12629\\
38000	12477\\
40000	12782\\
42000	12573\\
44000	12053\\
46000	12231\\
48000	12616\\
50000	12542\\
52000	12156\\
54000	12367\\
56000	12660\\
58000	12160\\
60000	12791\\
62000	12524\\
64000	12680\\
66000	12708\\
68000	12514\\
70000	12513\\
72000	12595\\
74000	12291\\
76000	12635\\
78000	12873\\
80000	12533\\
82000	12367\\
84000	12254\\
86000	11949\\
88000	12441\\
90000	12372\\
92000	12370\\
94000	12276\\
96000	12398\\
98000	12395\\
1e+05	12434\\
};
\addplot [color=mycolor26, forget plot]
  table[row sep=crcr]{%
0	4395.3\\
2000	7700.1\\
4000	8357.7\\
6000	8215.1\\
8000	8404.1\\
10000	8250.9\\
12000	8506\\
14000	8584.5\\
16000	8571.1\\
18000	8490.9\\
20000	8258.9\\
22000	8563.8\\
24000	8344.9\\
26000	8406.9\\
28000	8420.7\\
30000	8292.7\\
32000	8546.4\\
34000	7923.5\\
36000	8354\\
38000	8259.4\\
40000	8280.1\\
42000	8200.1\\
44000	8602.1\\
46000	8540.1\\
48000	8263.5\\
50000	8462.7\\
52000	8327.4\\
54000	8587.4\\
56000	8354.1\\
58000	8208\\
60000	7958.9\\
62000	8505.5\\
64000	8427.6\\
66000	8380.7\\
68000	8151.3\\
70000	8386.8\\
72000	8407.4\\
74000	8041.3\\
76000	7783.1\\
78000	8157.7\\
80000	8211.7\\
84000	8505.6\\
86000	8084.8\\
88000	8212.4\\
90000	8015.1\\
92000	8011.8\\
94000	8554.2\\
96000	8186.1\\
98000	8444.4\\
1e+05	8322.5\\
};
\addplot [color=mycolor27, forget plot]
  table[row sep=crcr]{%
0	4395.3\\
2000	10120\\
4000	9631.1\\
6000	9533\\
8000	9728.3\\
10000	9825.2\\
12000	9771.7\\
14000	9526.5\\
16000	9973\\
18000	9695.6\\
20000	9905.4\\
22000	9990.1\\
24000	9996.3\\
26000	9857.4\\
28000	9950.6\\
30000	9395.4\\
32000	9725.1\\
34000	10475\\
36000	9551.7\\
38000	9532.4\\
40000	9957.5\\
42000	10250\\
44000	9729.2\\
46000	9638.1\\
48000	9645.2\\
50000	9768\\
54000	9567.6\\
56000	9777.1\\
58000	9671.3\\
60000	10320\\
62000	10081\\
64000	10022\\
66000	10008\\
68000	9952.9\\
70000	9558.6\\
72000	9744.6\\
74000	9455.7\\
76000	9515.3\\
78000	9970.3\\
80000	9609.3\\
82000	9841.7\\
84000	9558\\
86000	9632.3\\
88000	9549.2\\
90000	9670\\
92000	10062\\
98000	9741.5\\
1e+05	9835.4\\
};
\addplot [color=mycolor28, forget plot]
  table[row sep=crcr]{%
0	4395.3\\
2000	5907\\
4000	5303.6\\
6000	5281.8\\
8000	5416.9\\
10000	5314\\
12000	5312\\
14000	5425.5\\
16000	5305.3\\
18000	5424.4\\
20000	5671.8\\
22000	5563.2\\
24000	5624.3\\
26000	5496.1\\
28000	5667.7\\
30000	5622.6\\
32000	5440.2\\
34000	5421.1\\
38000	5605.1\\
40000	5511.2\\
42000	5344.8\\
44000	5302.2\\
46000	5450.8\\
48000	5443.3\\
50000	5453.5\\
52000	5437.6\\
54000	5300.5\\
56000	5683.7\\
58000	5300.7\\
60000	5286.9\\
62000	5322.4\\
64000	5533.3\\
66000	5506\\
68000	5559\\
70000	5420.4\\
74000	5580.6\\
76000	5242.6\\
78000	5246.3\\
80000	5539.1\\
82000	5582.4\\
84000	5693.9\\
86000	5509.2\\
88000	5438\\
90000	5589\\
92000	5656.1\\
94000	5423.5\\
96000	5315.4\\
98000	5632.3\\
1e+05	5444.9\\
};
\addplot [color=mycolor29, forget plot]
  table[row sep=crcr]{%
0	4395.3\\
2000	11518\\
4000	11724\\
6000	11448\\
8000	11653\\
10000	12016\\
14000	11730\\
16000	12149\\
18000	11869\\
20000	12056\\
22000	11610\\
24000	11671\\
26000	11903\\
28000	11569\\
30000	11750\\
32000	11845\\
34000	11315\\
36000	12140\\
38000	11499\\
40000	12232\\
42000	11790\\
44000	11759\\
46000	11488\\
48000	11351\\
50000	11509\\
52000	11461\\
54000	12076\\
56000	11292\\
58000	11750\\
60000	11532\\
62000	11703\\
64000	11213\\
66000	11693\\
68000	12062\\
70000	11666\\
72000	12068\\
74000	12119\\
78000	11623\\
80000	10991\\
82000	11741\\
86000	12086\\
88000	11557\\
90000	11476\\
92000	11853\\
94000	12118\\
96000	11595\\
98000	11627\\
1e+05	12212\\
};
\addplot [color=mycolor30, forget plot]
  table[row sep=crcr]{%
0	4395.3\\
2000	7314.5\\
4000	7630.5\\
6000	7858\\
10000	8013.1\\
12000	8007.5\\
14000	7826.7\\
16000	7775.6\\
18000	7816\\
20000	8056.1\\
22000	7969.7\\
24000	7657.4\\
26000	7946.7\\
28000	7625.4\\
30000	8126.2\\
32000	8116.5\\
34000	7825.7\\
36000	8049.6\\
38000	8143.9\\
40000	7720\\
44000	7865.1\\
46000	7875.4\\
48000	8004.1\\
50000	7773.9\\
52000	7947.3\\
54000	8227.3\\
56000	7941.3\\
58000	7811.8\\
60000	8060.9\\
62000	7721.9\\
64000	7955.6\\
66000	7888.6\\
68000	8054\\
70000	7832.7\\
72000	7806.5\\
74000	7675.7\\
76000	7838.7\\
78000	7978.5\\
80000	7981.9\\
82000	7879.3\\
84000	7962.3\\
86000	7898.4\\
88000	7945.2\\
90000	7921.2\\
92000	7666.7\\
94000	8019.9\\
96000	7619.6\\
98000	8121.9\\
1e+05	7976.1\\
};
\addplot [color=mycolor31, forget plot]
  table[row sep=crcr]{%
0	4395.3\\
2000	6719.6\\
4000	6595.4\\
6000	6257\\
8000	6532.8\\
10000	6484.9\\
12000	6478.8\\
14000	6846.9\\
16000	6469.9\\
18000	6534.2\\
20000	6453.9\\
22000	6212.6\\
24000	6544.8\\
26000	6490.7\\
28000	6703\\
30000	6472.1\\
32000	6605.1\\
34000	6273.7\\
38000	6654.3\\
40000	6233.7\\
42000	6580.7\\
44000	6609.4\\
46000	6289.7\\
48000	6672.3\\
50000	6482.3\\
52000	6640.6\\
54000	6597.3\\
56000	6702.4\\
60000	6337.2\\
62000	6476.9\\
64000	6463.6\\
66000	6541.7\\
68000	6410.8\\
70000	6620.6\\
72000	6596.9\\
74000	6666\\
76000	6310.3\\
78000	6743.8\\
80000	6447.4\\
82000	6666.2\\
84000	6612.3\\
86000	6499.5\\
90000	6741.8\\
92000	6719.2\\
94000	6397.5\\
96000	6528\\
98000	6613.7\\
1e+05	6801.3\\
};
\addplot [color=mycolor32, forget plot]
  table[row sep=crcr]{%
0	4395.3\\
2000	11036\\
4000	11096\\
6000	10742\\
8000	10433\\
10000	10383\\
12000	11268\\
14000	10241\\
16000	11094\\
18000	10639\\
20000	10408\\
22000	10674\\
24000	10697\\
26000	10818\\
28000	10759\\
30000	10356\\
32000	10514\\
34000	10362\\
36000	10556\\
38000	10201\\
40000	10778\\
42000	10918\\
44000	10714\\
46000	10393\\
48000	10797\\
50000	11249\\
52000	10454\\
54000	10271\\
58000	10602\\
60000	10392\\
62000	10303\\
64000	10782\\
66000	10655\\
68000	11192\\
70000	10046\\
72000	10455\\
74000	10958\\
76000	10647\\
78000	10761\\
80000	10495\\
82000	10879\\
84000	10664\\
86000	10772\\
88000	11003\\
90000	10725\\
94000	10713\\
96000	10455\\
98000	10652\\
1e+05	10357\\
};
\addplot [color=mycolor33, forget plot]
  table[row sep=crcr]{%
0	4395.3\\
2000	7678.4\\
4000	7819.1\\
6000	8034.7\\
8000	8371.9\\
10000	8449.2\\
12000	8295.6\\
14000	8339.7\\
16000	8013.4\\
18000	8284.4\\
20000	8164.9\\
22000	8108.1\\
24000	8367.9\\
26000	8341.5\\
28000	8377.3\\
30000	8294.3\\
32000	7968.4\\
34000	8199.5\\
36000	8589.2\\
38000	8186.2\\
40000	8349.9\\
42000	8231.6\\
44000	8277.7\\
46000	8369.4\\
48000	8167.9\\
50000	8265.8\\
52000	8561.9\\
54000	8331.4\\
56000	8473.9\\
58000	8532.2\\
60000	8054.9\\
62000	8527.1\\
64000	8196.1\\
66000	8439.3\\
68000	8340\\
70000	8189.9\\
72000	8236\\
74000	8430.8\\
76000	7969.3\\
78000	8248.9\\
80000	8170.8\\
82000	8198.7\\
84000	8058.2\\
86000	8132.3\\
88000	8050\\
90000	8205.5\\
92000	8418\\
94000	8128\\
98000	8348.2\\
1e+05	8101.5\\
};
\addplot [color=mycolor34, forget plot]
  table[row sep=crcr]{%
0	4395.3\\
2000	4407.4\\
4000	3967.5\\
6000	4112\\
8000	3974.9\\
10000	3922.5\\
14000	4024.1\\
16000	4121.2\\
18000	4073.8\\
20000	4083.2\\
22000	3989.5\\
24000	4070.5\\
28000	4002.9\\
30000	4110.9\\
32000	3954.4\\
34000	3906.9\\
36000	4018.1\\
38000	4083.9\\
40000	3913.3\\
42000	4097.8\\
44000	4051.4\\
46000	4093.4\\
48000	4093.7\\
50000	3948.8\\
52000	3943.6\\
54000	3999\\
56000	3951.7\\
58000	4140.5\\
60000	4128.8\\
62000	4062.8\\
64000	4058.4\\
66000	3965.9\\
68000	4097.5\\
70000	4075.5\\
72000	4019.9\\
74000	4060\\
76000	3964.3\\
78000	3894.9\\
80000	4082.9\\
82000	3981.3\\
84000	4038.3\\
86000	4004.6\\
88000	4008\\
90000	3922\\
92000	4052.3\\
94000	4022.5\\
96000	3937.2\\
98000	4114.9\\
1e+05	3832.7\\
};
\addplot [color=mycolor35]
  table[row sep=crcr]{%
0	4395.3\\
2000	9749.6\\
4000	10224\\
6000	10406\\
8000	10514\\
10000	10307\\
12000	10359\\
14000	10171\\
16000	10195\\
18000	10567\\
20000	10439\\
22000	10375\\
24000	10144\\
26000	10452\\
28000	10227\\
30000	10528\\
32000	10247\\
34000	10014\\
36000	10413\\
38000	10132\\
40000	10500\\
42000	10330\\
44000	10497\\
46000	10194\\
48000	10375\\
50000	10234\\
52000	10194\\
54000	10481\\
56000	10225\\
58000	10372\\
60000	10168\\
62000	10351\\
64000	10207\\
66000	10415\\
68000	10236\\
70000	10133\\
72000	10390\\
74000	10344\\
76000	10350\\
78000	10470\\
80000	10127\\
82000	10582\\
84000	10262\\
86000	10490\\
88000	10569\\
92000	10044\\
94000	10389\\
96000	10347\\
98000	10001\\
1e+05	10269\\
};
\addlegendentry{Classe k=83}

\addplot [color=mycolor36, forget plot]
  table[row sep=crcr]{%
0	4395.3\\
2000	4952.7\\
4000	4331.2\\
6000	4113.1\\
8000	4314.7\\
10000	4341.1\\
12000	4248.6\\
14000	4220.1\\
16000	4115.1\\
18000	4339.5\\
20000	4378\\
22000	4257.1\\
24000	4123.5\\
26000	4454.2\\
28000	4269.6\\
30000	4257.1\\
32000	4143.6\\
34000	4254.9\\
36000	4299.9\\
38000	4391.8\\
40000	4227.4\\
42000	4281.9\\
44000	4140.2\\
46000	4233.6\\
48000	4082.1\\
50000	4260.8\\
52000	4390.6\\
54000	4239.4\\
56000	4488.6\\
58000	4136\\
60000	4172.9\\
62000	4334.5\\
64000	4208.5\\
66000	4126\\
68000	4144.3\\
70000	4349.5\\
72000	4162.2\\
74000	4331.2\\
76000	4250.4\\
78000	4292.9\\
80000	4199.5\\
82000	4272.1\\
84000	4159.1\\
86000	4253.1\\
90000	4284.2\\
92000	4314.2\\
94000	4239\\
96000	4310.8\\
98000	4260.1\\
1e+05	4302.9\\
};
\addplot [color=mycolor37, forget plot]
  table[row sep=crcr]{%
0	4395.3\\
2000	3897.9\\
4000	3577.3\\
6000	3573.9\\
8000	3635.1\\
10000	3631.9\\
12000	3654.5\\
14000	3547.6\\
16000	3616.4\\
18000	3495\\
20000	3684.9\\
22000	3576.3\\
24000	3617.2\\
26000	3510.8\\
28000	3607.7\\
30000	3655\\
32000	3568.2\\
34000	3549\\
36000	3641.5\\
38000	3623.1\\
40000	3511.6\\
42000	3567.1\\
44000	3778.9\\
46000	3638\\
48000	3670.1\\
50000	3670.1\\
52000	3654.3\\
54000	3583.9\\
56000	3571.9\\
58000	3543.2\\
60000	3618.8\\
62000	3616.3\\
64000	3507.1\\
66000	3636.3\\
68000	3618\\
70000	3571.9\\
72000	3683.6\\
74000	3701.5\\
76000	3544.7\\
78000	3587.3\\
80000	3545.1\\
82000	3645.2\\
84000	3561.8\\
86000	3548.1\\
88000	3629\\
90000	3620.8\\
92000	3511\\
94000	3664.1\\
96000	3586.9\\
98000	3626.8\\
1e+05	3628.6\\
};
\addplot [color=mycolor38, forget plot]
  table[row sep=crcr]{%
0	4395.3\\
2000	3176.3\\
4000	2969.9\\
6000	2960.2\\
8000	2920.5\\
10000	2972.3\\
12000	2988.8\\
14000	2992.4\\
16000	2882.8\\
18000	2825.1\\
20000	2834.8\\
22000	2955.9\\
24000	2949.9\\
26000	2898.3\\
30000	3044.6\\
32000	3008.6\\
34000	2898.2\\
36000	2969.9\\
38000	2921.1\\
40000	3172.1\\
42000	2915.4\\
44000	2850.9\\
46000	2884.9\\
48000	2929\\
50000	2938\\
54000	2868.8\\
56000	2806.8\\
58000	2853.6\\
60000	2886\\
62000	2886\\
64000	2926.3\\
66000	2923.6\\
68000	3016\\
70000	2848.9\\
72000	2933.8\\
74000	2931.8\\
76000	3008.5\\
78000	2951.3\\
80000	3027.5\\
82000	2982.8\\
84000	2854.9\\
86000	3045.6\\
88000	2864.6\\
94000	2933.9\\
96000	2827\\
98000	2809.8\\
1e+05	2896.5\\
};
\addplot [color=mycolor39, forget plot]
  table[row sep=crcr]{%
0	4395.3\\
2000	12354\\
4000	12324\\
6000	12441\\
8000	11691\\
10000	11711\\
12000	11669\\
14000	11523\\
16000	12015\\
18000	12028\\
20000	12108\\
22000	11880\\
24000	12142\\
26000	11566\\
28000	12176\\
30000	12245\\
32000	12690\\
34000	12756\\
36000	11774\\
38000	11883\\
40000	11820\\
42000	12120\\
44000	11734\\
46000	12421\\
48000	11662\\
50000	12138\\
52000	11542\\
54000	11786\\
56000	11958\\
58000	11715\\
60000	11820\\
62000	12142\\
64000	11913\\
66000	11632\\
68000	11841\\
70000	11718\\
72000	11917\\
74000	11164\\
76000	12118\\
78000	12143\\
80000	12044\\
82000	11880\\
84000	11625\\
86000	12230\\
88000	12454\\
90000	11878\\
92000	12088\\
94000	12340\\
96000	11674\\
1e+05	12246\\
};
\addplot [color=mycolor40, forget plot]
  table[row sep=crcr]{%
0	4395.3\\
2000	5283\\
4000	5695.9\\
12000	5680.5\\
16000	5722.1\\
18000	5673.3\\
20000	5715\\
22000	5616.3\\
24000	5766.6\\
26000	5578.1\\
28000	5844.8\\
30000	5665.1\\
32000	5598.4\\
36000	5880\\
38000	5698.1\\
40000	5861.3\\
42000	5822.2\\
44000	5748.1\\
46000	5724.2\\
48000	5840.7\\
50000	5714.1\\
52000	5732.6\\
54000	5694.3\\
56000	5752.4\\
58000	5884.2\\
60000	5667.6\\
62000	5836.1\\
64000	5829\\
66000	5771.6\\
68000	5906.1\\
70000	5769\\
72000	5707.8\\
74000	5714.9\\
76000	5777.2\\
78000	5808.6\\
80000	5777.1\\
82000	5481.7\\
84000	5761.2\\
86000	5663.5\\
88000	5648.6\\
90000	5677.2\\
92000	5657.1\\
96000	5710.5\\
98000	5593.7\\
1e+05	5618.1\\
};
\addplot [color=mycolor41, forget plot]
  table[row sep=crcr]{%
0	4395.3\\
2000	4811.1\\
4000	5113.3\\
6000	5379.3\\
8000	5381.1\\
10000	5274.5\\
12000	5447.4\\
14000	5329.2\\
18000	5433.8\\
20000	5329.8\\
22000	5253.8\\
26000	5453.4\\
28000	5392.1\\
30000	5479.5\\
32000	5426.5\\
34000	5446.3\\
36000	5244.9\\
38000	5476.4\\
40000	5294.6\\
42000	5148.7\\
46000	5432.9\\
48000	5449.3\\
50000	5258.2\\
52000	5160.6\\
54000	5380\\
56000	5442.7\\
58000	5443.3\\
60000	5290.1\\
66000	5217.3\\
68000	5473.8\\
70000	5522.9\\
72000	5255.4\\
74000	5310.8\\
78000	5368.4\\
80000	5550.2\\
82000	5354.9\\
84000	5436.4\\
86000	5362.3\\
88000	5258.1\\
90000	5524\\
92000	5367\\
94000	5555.2\\
96000	5268.8\\
98000	5250.9\\
1e+05	5275.7\\
};
\addplot [color=mycolor42, forget plot]
  table[row sep=crcr]{%
0	4395.3\\
2000	17049\\
4000	20563\\
6000	20972\\
8000	20545\\
10000	21440\\
12000	20560\\
14000	21934\\
16000	21187\\
18000	20254\\
20000	20952\\
22000	20558\\
24000	20681\\
26000	20408\\
28000	21528\\
30000	21511\\
32000	20128\\
34000	21148\\
36000	22066\\
38000	21737\\
40000	20734\\
42000	21596\\
44000	21587\\
46000	21405\\
48000	20792\\
50000	20321\\
52000	21373\\
54000	20950\\
56000	20125\\
58000	20760\\
60000	21924\\
62000	21331\\
64000	20146\\
66000	21091\\
68000	21263\\
70000	20671\\
72000	21181\\
74000	21898\\
76000	21780\\
78000	20840\\
80000	20770\\
82000	20569\\
84000	20744\\
86000	21080\\
88000	20747\\
90000	21324\\
92000	20178\\
94000	21271\\
96000	21719\\
98000	20922\\
1e+05	21086\\
};
\addplot [color=mycolor43, forget plot]
  table[row sep=crcr]{%
0	4395.3\\
2000	5329.4\\
4000	4968.4\\
6000	4868.6\\
8000	4924.8\\
10000	4900.4\\
12000	4947.3\\
14000	5023.4\\
16000	4891.5\\
18000	4909.9\\
20000	4819.8\\
24000	4967.7\\
26000	4817.9\\
28000	4885.2\\
30000	4911.9\\
32000	4991.7\\
34000	4999.2\\
36000	4899.8\\
38000	5007.2\\
40000	4966.5\\
42000	5043.4\\
44000	4731.1\\
46000	4804\\
48000	4789.9\\
50000	4866.1\\
52000	5000.6\\
54000	4844.3\\
56000	4954.5\\
58000	4874.2\\
60000	4866.6\\
62000	4876.7\\
66000	4847\\
68000	4903.1\\
70000	4866.3\\
72000	4857.1\\
74000	4949.4\\
76000	4936\\
78000	5072.4\\
80000	4822.6\\
82000	4926.1\\
84000	4966.8\\
86000	5070.7\\
88000	5011.9\\
90000	5040.1\\
92000	5037\\
94000	4785.6\\
96000	4738\\
98000	4948.5\\
1e+05	4817.8\\
};
\addplot [color=mycolor44, forget plot]
  table[row sep=crcr]{%
0	4395.3\\
2000	3774.7\\
4000	3445.1\\
6000	3243\\
8000	3245.9\\
10000	3299.3\\
12000	3197\\
14000	3227.6\\
16000	3350.4\\
18000	3260.8\\
20000	3228.7\\
22000	3385\\
24000	3185.6\\
26000	3457.8\\
30000	3335.9\\
32000	3336.1\\
34000	3362.1\\
36000	3133.4\\
38000	3424.8\\
40000	3394.8\\
42000	3384.7\\
44000	3357.7\\
46000	3368.6\\
48000	3129.9\\
50000	3357.4\\
52000	3404\\
54000	3318.8\\
56000	3331.9\\
58000	3273\\
60000	3139.7\\
62000	3152\\
64000	3341\\
66000	3469.5\\
68000	3302.1\\
70000	3388.8\\
72000	3273.5\\
74000	3404.2\\
76000	3318.6\\
78000	3340.1\\
80000	3208.9\\
82000	3324.8\\
84000	3542.8\\
86000	3378.7\\
88000	3281.4\\
90000	3289.7\\
92000	3419.9\\
94000	3293.9\\
96000	3402.3\\
98000	3240.5\\
1e+05	3205.5\\
};
\addplot [color=mycolor45, forget plot]
  table[row sep=crcr]{%
0	4395.3\\
2000	6161.6\\
4000	6597.3\\
6000	6787.6\\
8000	6798.7\\
10000	6997.3\\
12000	6946.1\\
14000	6871.6\\
16000	7035.3\\
18000	6851.9\\
20000	6723.5\\
22000	6975.1\\
24000	6851.1\\
26000	7027.4\\
30000	6939.1\\
32000	6865.5\\
34000	6927.9\\
36000	6943.1\\
38000	7017.3\\
40000	6984.6\\
42000	6907.8\\
44000	7009.8\\
46000	6876.1\\
48000	6884.8\\
50000	6924.1\\
52000	6682\\
54000	6876\\
56000	6871.3\\
58000	6816.7\\
60000	6945.2\\
62000	7120.6\\
64000	6746.4\\
66000	7063.1\\
68000	7036.8\\
70000	6882.5\\
72000	6888.8\\
74000	6602.9\\
76000	6903.5\\
78000	7042.9\\
80000	7128.5\\
82000	6889.6\\
84000	6732.8\\
86000	7120.5\\
88000	6972.8\\
90000	6932.5\\
92000	7098.9\\
94000	7100.8\\
98000	7190.6\\
1e+05	6671.8\\
};
\addplot [color=mycolor46, forget plot]
  table[row sep=crcr]{%
0	4395.3\\
2000	3838.2\\
4000	3959.2\\
6000	3925.4\\
8000	4120.6\\
10000	4004.5\\
12000	4090.5\\
14000	4085.2\\
16000	4067.3\\
18000	3936.7\\
20000	4063.9\\
22000	4073.3\\
24000	3913.5\\
26000	4048.1\\
28000	4070.2\\
30000	3932\\
32000	3968.1\\
34000	4086.6\\
36000	4116.8\\
38000	3996.4\\
40000	3978\\
42000	3985.9\\
44000	4044.9\\
46000	4063.4\\
48000	4095.2\\
50000	4175.4\\
52000	4054\\
54000	4145.3\\
58000	4028.8\\
60000	4091.6\\
64000	3949.6\\
66000	3964.7\\
68000	4089.7\\
70000	3937.7\\
72000	4013.7\\
76000	4071.6\\
78000	4071.3\\
80000	4206.9\\
82000	4068.6\\
84000	4102.8\\
86000	4005.2\\
88000	3968.3\\
90000	4006.4\\
92000	4081.8\\
94000	4057.9\\
96000	3961.7\\
98000	3893.5\\
1e+05	3989\\
};
\addplot [color=mycolor47, forget plot]
  table[row sep=crcr]{%
0	4395.3\\
2000	5433.2\\
4000	5144.6\\
6000	5402.6\\
8000	5333.4\\
10000	5373.3\\
12000	5333.4\\
14000	5321.5\\
16000	5239\\
18000	5348.7\\
20000	5347.3\\
22000	5373.1\\
24000	5259.5\\
26000	5227.2\\
28000	5524.3\\
30000	5039.6\\
32000	5310\\
36000	5320.4\\
38000	5277.8\\
40000	5286.6\\
42000	5216\\
44000	5276.7\\
46000	5253.4\\
48000	5407.8\\
50000	5185.4\\
52000	5333.4\\
54000	5239\\
56000	5250.5\\
58000	5100.1\\
60000	5116\\
62000	5266.5\\
64000	5312.5\\
66000	5184.4\\
68000	5198\\
70000	5231.8\\
72000	5304.1\\
74000	5340.2\\
76000	5105\\
78000	5319.3\\
80000	5201.7\\
82000	5323.8\\
84000	5354.8\\
86000	5068.7\\
88000	5215.1\\
90000	5281\\
92000	5105.9\\
94000	5109.1\\
96000	5320.2\\
98000	5369.3\\
1e+05	5319.8\\
};
\addplot [color=mycolor48, forget plot]
  table[row sep=crcr]{%
0	4395.3\\
2000	3196.1\\
4000	2856.7\\
6000	2819.9\\
10000	2886.3\\
12000	2808.8\\
16000	2884.6\\
18000	2917.7\\
20000	2821.3\\
22000	2883.9\\
24000	2852.6\\
26000	2864\\
28000	2817.4\\
30000	2826.1\\
32000	2889.4\\
34000	2945.3\\
36000	2883.2\\
44000	2838.7\\
46000	2892.3\\
48000	2871.8\\
50000	2888\\
52000	2967.9\\
54000	2884\\
56000	2824.3\\
58000	2890\\
60000	2917.2\\
62000	2845.3\\
64000	2980.5\\
66000	2888.8\\
68000	2908.8\\
72000	2911\\
76000	2878.5\\
78000	2929.2\\
80000	2863.8\\
82000	2859.9\\
84000	2870.4\\
88000	2913.5\\
92000	2839.5\\
96000	2857.1\\
98000	2850.9\\
1e+05	2879.9\\
};
\addplot [color=mycolor49, forget plot]
  table[row sep=crcr]{%
0	4395.3\\
2000	4550.7\\
4000	4306.9\\
6000	4308.1\\
8000	4421.2\\
10000	4254.4\\
14000	4278.1\\
16000	4183.8\\
18000	4255\\
20000	4498.8\\
22000	4440.5\\
24000	4359.4\\
26000	4137.1\\
28000	4401.2\\
30000	4407.6\\
32000	4338.8\\
34000	4377.5\\
36000	4484.1\\
38000	4322.7\\
40000	4464.7\\
42000	4259.9\\
44000	4212.1\\
46000	4266.5\\
48000	4120.2\\
50000	4449.9\\
52000	4469.7\\
54000	4210.1\\
56000	4352.3\\
58000	4220.4\\
60000	4413.2\\
62000	4421\\
66000	4320.2\\
68000	4539.6\\
70000	4361.9\\
74000	4053.2\\
76000	4241.8\\
78000	4286.9\\
80000	4212.9\\
82000	4122\\
84000	4230\\
86000	4421.8\\
88000	4412.8\\
90000	4515.7\\
92000	4381.6\\
94000	4404.3\\
96000	4370.6\\
98000	4292.6\\
1e+05	4377\\
};
\addplot [color=mycolor50, forget plot]
  table[row sep=crcr]{%
0	4395.3\\
2000	2405.2\\
4000	2346.3\\
6000	2355.9\\
8000	2448.5\\
10000	2348.5\\
12000	2436.3\\
14000	2346.2\\
16000	2353.5\\
18000	2384.5\\
20000	2366.1\\
22000	2528.1\\
24000	2361.8\\
26000	2360.8\\
28000	2400.6\\
30000	2407.4\\
32000	2311.6\\
34000	2399.8\\
36000	2351.9\\
38000	2362.1\\
40000	2472.5\\
44000	2397.8\\
46000	2394.1\\
48000	2420.2\\
50000	2409.3\\
52000	2417.4\\
54000	2440.8\\
56000	2411.1\\
58000	2401.4\\
62000	2401.6\\
64000	2385.2\\
66000	2401.3\\
68000	2374.2\\
70000	2419.4\\
72000	2456.7\\
74000	2362.6\\
76000	2496.3\\
78000	2471\\
80000	2463.7\\
82000	2412.2\\
84000	2386.9\\
86000	2397.2\\
88000	2392\\
90000	2405.6\\
96000	2368\\
98000	2338.7\\
1e+05	2433.3\\
};
\addplot [color=mycolor51, forget plot]
  table[row sep=crcr]{%
0	4395.3\\
2000	5687.9\\
4000	5770.1\\
6000	5823.4\\
8000	6185.4\\
10000	5732.5\\
12000	6003.1\\
14000	5934.9\\
16000	6146.9\\
18000	5916.1\\
20000	5985.8\\
22000	5873.1\\
24000	6116.5\\
26000	6188.7\\
28000	5883.3\\
30000	5959.2\\
32000	5877.8\\
34000	6011.8\\
36000	5959.4\\
38000	5969.3\\
40000	5825.8\\
42000	5710.3\\
44000	6183.9\\
46000	6047.2\\
48000	6094.4\\
50000	5895.1\\
54000	5916.1\\
56000	5992.1\\
58000	5741.4\\
60000	6000.2\\
62000	6381.3\\
64000	5975.7\\
66000	5852\\
68000	6057\\
70000	6114.9\\
72000	5958\\
74000	6350.7\\
76000	6008.9\\
78000	5609.9\\
80000	5989.5\\
82000	6154.2\\
84000	5720.1\\
86000	5925.9\\
88000	5806.1\\
90000	5913\\
92000	5789\\
94000	5550.3\\
96000	5773.8\\
98000	5934.6\\
1e+05	5981.8\\
};
\addplot [color=mycolor52, forget plot]
  table[row sep=crcr]{%
0	4395.3\\
2000	4009.6\\
4000	3477.9\\
6000	3490.3\\
8000	3485.2\\
14000	3513.1\\
16000	3432.6\\
18000	3569.3\\
20000	3319.2\\
22000	3331.9\\
24000	3497.6\\
26000	3466.2\\
28000	3487.5\\
30000	3456.1\\
32000	3476.9\\
34000	3378.8\\
36000	3450.6\\
38000	3493.3\\
40000	3650.9\\
42000	3491.2\\
44000	3427.7\\
46000	3380.1\\
48000	3354.3\\
50000	3425.2\\
54000	3431.9\\
56000	3228.9\\
58000	3335\\
60000	3300.9\\
62000	3520.6\\
66000	3358.4\\
68000	3706.2\\
70000	3495.7\\
72000	3352.6\\
74000	3502.5\\
76000	3404.6\\
78000	3399.9\\
80000	3573.2\\
82000	3490.2\\
84000	3583.1\\
86000	3539.7\\
88000	3785.8\\
90000	3577.3\\
92000	3503.1\\
94000	3533.8\\
96000	3433.7\\
98000	3352.2\\
1e+05	3396.8\\
};
\addplot [color=mycolor53, forget plot]
  table[row sep=crcr]{%
0	4395.3\\
2000	6340.7\\
4000	5916.7\\
6000	5734.1\\
8000	5657.8\\
10000	5411.7\\
12000	5657.2\\
14000	5509.9\\
16000	5592.5\\
18000	5737\\
20000	5464.8\\
22000	5655.8\\
24000	5725.5\\
26000	5641.6\\
28000	5486.9\\
30000	5772\\
34000	5727.5\\
36000	5701.3\\
38000	5974.1\\
40000	5697.4\\
42000	5699.1\\
44000	5583.5\\
46000	5640.9\\
50000	5676.5\\
52000	5582.5\\
54000	5613.5\\
56000	5698.9\\
58000	5844.5\\
60000	5660.1\\
62000	5659.1\\
64000	5596.6\\
66000	5806\\
68000	5680.7\\
70000	5712.4\\
72000	5564.3\\
74000	5545.3\\
76000	5897.4\\
78000	5642.4\\
80000	5894\\
82000	5496.3\\
84000	5521.3\\
86000	5567.9\\
88000	5749.5\\
90000	5600.5\\
92000	5738.2\\
94000	5623.4\\
96000	5889.9\\
98000	5684.3\\
1e+05	5641.2\\
};
\addplot [color=mycolor54, forget plot]
  table[row sep=crcr]{%
0	4395.3\\
2000	4755.8\\
4000	4641.9\\
6000	4505.5\\
8000	4544.9\\
10000	4542.7\\
12000	4603.4\\
14000	4411.4\\
16000	4625.2\\
18000	4649\\
20000	4608.8\\
22000	4638.7\\
24000	4548.2\\
26000	4567.6\\
28000	4490\\
30000	4565\\
32000	4565.8\\
34000	4530.3\\
36000	4372.4\\
38000	4464.1\\
40000	4599.4\\
42000	4536\\
44000	4658.7\\
46000	4497.9\\
48000	4598.4\\
50000	4493\\
52000	4628.6\\
54000	4572.1\\
56000	4616.7\\
58000	4430.2\\
60000	4616.2\\
62000	4603.5\\
64000	4798.3\\
66000	4510.3\\
68000	4506.2\\
70000	4469.5\\
72000	4675.3\\
74000	4489.5\\
76000	4354.5\\
78000	4514.2\\
80000	4580.3\\
82000	4767.6\\
84000	4502.7\\
86000	4608.7\\
88000	4632.9\\
90000	4548.6\\
92000	4632.8\\
94000	4632.1\\
96000	4707.8\\
98000	4480.3\\
1e+05	4571.6\\
};
\addplot [color=mycolor55, forget plot]
  table[row sep=crcr]{%
0	4395.3\\
2000	5222.7\\
4000	5792.9\\
6000	5619.1\\
8000	5687.4\\
10000	5842.4\\
14000	5785.6\\
16000	5780.7\\
18000	5929.9\\
20000	5564.6\\
22000	5699.7\\
24000	5697.7\\
26000	5675.3\\
28000	5765.5\\
30000	5804.6\\
32000	5806.1\\
34000	5765\\
36000	5641.2\\
38000	5732\\
40000	5680.5\\
42000	5835.7\\
44000	5735.8\\
46000	5914.6\\
48000	5721.3\\
50000	5764.9\\
52000	5825.9\\
54000	5745.6\\
56000	5618.9\\
58000	5700.9\\
60000	5702.9\\
62000	5861.8\\
64000	5651.4\\
66000	5745.9\\
68000	5727.6\\
70000	5764\\
72000	5606.1\\
74000	5746.6\\
76000	5789.9\\
78000	5631.9\\
80000	5832.1\\
82000	5655.7\\
84000	5751.5\\
90000	5744.3\\
92000	5876.1\\
94000	5895.9\\
96000	5680.9\\
98000	5814.3\\
1e+05	5648.6\\
};
\addplot [color=mycolor56, forget plot]
  table[row sep=crcr]{%
0	4395.3\\
2000	4192.4\\
4000	4093.7\\
6000	3967.1\\
8000	3966.6\\
10000	4053.2\\
12000	4033.8\\
14000	3977\\
16000	4004.4\\
18000	3914.1\\
20000	3950.6\\
22000	3923.5\\
24000	3911.2\\
26000	3945.5\\
28000	3857.5\\
30000	3967\\
32000	3864.1\\
34000	3988.5\\
36000	3983.6\\
38000	3938.3\\
40000	4037\\
42000	3971.2\\
44000	3937.5\\
46000	3987.4\\
48000	4003\\
50000	3910.9\\
52000	4041\\
54000	3923\\
58000	3980.2\\
60000	3912.4\\
62000	3974.3\\
64000	3885.2\\
66000	3994.7\\
68000	3930.6\\
70000	3930.8\\
72000	3980.1\\
74000	3891.9\\
76000	3929.5\\
82000	4121.4\\
84000	3981.5\\
86000	3926.8\\
88000	4082\\
90000	3962.6\\
92000	4125.7\\
94000	3877.8\\
96000	3933.2\\
98000	3896.7\\
1e+05	3973.1\\
};
\addplot [color=mycolor57, forget plot]
  table[row sep=crcr]{%
0	4395.3\\
2000	6507.1\\
4000	6801.9\\
6000	7012.1\\
8000	6693.8\\
10000	6540.8\\
12000	6859.7\\
14000	6576.1\\
16000	6808.7\\
20000	6629.6\\
22000	6726.1\\
24000	6680\\
26000	6814.8\\
28000	6657.5\\
32000	6681.8\\
34000	6917.4\\
36000	6780.5\\
38000	6596\\
40000	6638\\
42000	6734.2\\
46000	6728.8\\
48000	6594.6\\
50000	6784.8\\
52000	6587.3\\
54000	6664.2\\
56000	6540.1\\
58000	6862.1\\
60000	6622.9\\
62000	6631.9\\
64000	6801.2\\
66000	6600.3\\
68000	6757.6\\
70000	6743.4\\
72000	6846.4\\
74000	6922.9\\
76000	6896\\
78000	6848.5\\
80000	6929.5\\
82000	6819.1\\
84000	6674.6\\
86000	6425.5\\
88000	6838.9\\
90000	6779.3\\
92000	6673.1\\
94000	6588.1\\
96000	6863.6\\
98000	6719.7\\
1e+05	6691.3\\
};
\addplot [color=mycolor58, forget plot]
  table[row sep=crcr]{%
0	4395.3\\
2000	3912.7\\
4000	3504.9\\
6000	3483.4\\
8000	3505.4\\
10000	3470\\
12000	3499.2\\
14000	3451.2\\
16000	3524.7\\
18000	3513.7\\
20000	3485.7\\
22000	3569.4\\
24000	3467.2\\
26000	3590.2\\
28000	3518.6\\
30000	3460.9\\
32000	3534.3\\
34000	3489.8\\
36000	3523.3\\
38000	3495.9\\
42000	3582.6\\
44000	3416.4\\
46000	3464.9\\
48000	3413.5\\
50000	3493.1\\
52000	3522.8\\
54000	3539.1\\
58000	3421.9\\
62000	3505.7\\
64000	3418.1\\
66000	3482.8\\
68000	3521.9\\
70000	3474.5\\
72000	3502.4\\
74000	3616.5\\
76000	3519.5\\
78000	3472.7\\
80000	3542.1\\
82000	3474.1\\
84000	3490.4\\
86000	3491.5\\
88000	3508.4\\
90000	3499.3\\
92000	3601\\
94000	3482.9\\
96000	3558\\
98000	3528.1\\
1e+05	3596.9\\
};
\addplot [color=mycolor59, forget plot]
  table[row sep=crcr]{%
0	4395.3\\
2000	4818.5\\
4000	4973.5\\
6000	4960.8\\
8000	5077.8\\
10000	4928\\
12000	5112.5\\
14000	4953.1\\
16000	4979.3\\
18000	5111.7\\
20000	5081.6\\
22000	4970.9\\
24000	5025.8\\
28000	5052.9\\
30000	4999.5\\
32000	4987.3\\
34000	5019.2\\
36000	4954.8\\
38000	5005.6\\
40000	4955.9\\
42000	4988.4\\
44000	4896.1\\
46000	5026\\
48000	4936.2\\
50000	4985.8\\
52000	4917.9\\
54000	4988.9\\
56000	5109.7\\
58000	5082.3\\
60000	5112.8\\
62000	5045.5\\
64000	5008.9\\
66000	4862.3\\
68000	5041.3\\
70000	5024.2\\
72000	4965.1\\
74000	4999.9\\
76000	5168.7\\
78000	5058.8\\
80000	4880.7\\
82000	4997\\
86000	5038.7\\
88000	4943.4\\
90000	4889.4\\
92000	5066.5\\
94000	4921.6\\
96000	4928.3\\
98000	5090.4\\
1e+05	5058.9\\
};
\addplot [color=mycolor60, forget plot]
  table[row sep=crcr]{%
0	4395.3\\
2000	3652.1\\
4000	3587.7\\
10000	3525\\
12000	3578.9\\
14000	3509.2\\
16000	3525.7\\
18000	3520.5\\
20000	3525.9\\
22000	3480.6\\
24000	3479.4\\
26000	3507\\
28000	3432.6\\
30000	3435.5\\
32000	3453\\
34000	3510.9\\
36000	3491.3\\
38000	3504.1\\
40000	3535.4\\
42000	3449.8\\
44000	3543\\
46000	3466.9\\
48000	3530.3\\
50000	3473.4\\
54000	3456\\
56000	3493.5\\
58000	3556.8\\
60000	3502.4\\
62000	3631.8\\
64000	3507\\
68000	3466.9\\
70000	3514.1\\
72000	3515.2\\
74000	3547.1\\
76000	3556.4\\
78000	3490.8\\
80000	3481.8\\
82000	3457.1\\
84000	3503\\
86000	3527.9\\
88000	3506.2\\
90000	3550.5\\
92000	3478.8\\
94000	3465\\
96000	3585.5\\
1e+05	3516\\
};
\addplot [color=mycolor61, forget plot]
  table[row sep=crcr]{%
0	4395.3\\
2000	3119.4\\
4000	3116\\
6000	3031.8\\
8000	3062.8\\
10000	3082.1\\
12000	3055.5\\
14000	3019.4\\
16000	3071.1\\
18000	3094.8\\
22000	3042.1\\
24000	3042\\
28000	3069.8\\
30000	2976\\
32000	3045.4\\
38000	3098.4\\
40000	3060.6\\
44000	3104.6\\
46000	3108\\
48000	3082.3\\
50000	3047.6\\
52000	3079.3\\
54000	3036.8\\
56000	3032.5\\
58000	3008.7\\
60000	3086.9\\
62000	3027\\
64000	3068.6\\
66000	3095.1\\
70000	3067.2\\
72000	3091.6\\
74000	3099.6\\
76000	3024.5\\
78000	3072.8\\
80000	3045.1\\
82000	3043.9\\
84000	3118.5\\
86000	3061.2\\
88000	2993.2\\
90000	3014.9\\
92000	2968\\
94000	3073.8\\
96000	3041\\
98000	3059.5\\
1e+05	3020.6\\
};
\addplot [color=mycolor62, forget plot]
  table[row sep=crcr]{%
0	4395.3\\
2000	4845.5\\
4000	5010.7\\
6000	4937.9\\
8000	4906.2\\
10000	4959.3\\
12000	4983.6\\
14000	5053.8\\
16000	4971.8\\
18000	4985.7\\
20000	5075.9\\
22000	5015.8\\
24000	5038\\
26000	4931.7\\
28000	4891.3\\
30000	4949.9\\
32000	5034.1\\
34000	4975.2\\
36000	4977.9\\
38000	5101.2\\
42000	4917.2\\
44000	5021.9\\
46000	4982.6\\
48000	4892\\
50000	4927.1\\
52000	5063.3\\
54000	4957.8\\
56000	5055.7\\
58000	4924.9\\
60000	4938.4\\
62000	4988\\
64000	4932.2\\
66000	5034.4\\
68000	4887.2\\
70000	4958\\
72000	4970.3\\
74000	4883.2\\
76000	4923.6\\
80000	4940.4\\
82000	4786.7\\
84000	4905.7\\
86000	4954\\
88000	4828.7\\
90000	4862.8\\
92000	4984.9\\
94000	4912.5\\
96000	5011.4\\
98000	5046.2\\
1e+05	4911.7\\
};
\addplot [color=mycolor63, forget plot]
  table[row sep=crcr]{%
0	4395.3\\
2000	2736.4\\
4000	2375.1\\
6000	2303\\
8000	2370.7\\
10000	2335.9\\
12000	2327.9\\
14000	2352.8\\
16000	2330.2\\
20000	2345.2\\
22000	2369.1\\
24000	2357.4\\
26000	2370.5\\
28000	2374.1\\
30000	2347.8\\
32000	2340.7\\
34000	2293.1\\
36000	2351.8\\
40000	2382.5\\
42000	2355.2\\
44000	2364.7\\
48000	2350.4\\
50000	2393.1\\
52000	2443.5\\
54000	2377.1\\
56000	2337.2\\
58000	2345.9\\
62000	2378\\
64000	2329.2\\
66000	2413.1\\
68000	2363.4\\
70000	2339.4\\
72000	2334.4\\
74000	2348.2\\
76000	2335.2\\
80000	2334.7\\
82000	2345.6\\
84000	2296.4\\
86000	2424.3\\
88000	2421.4\\
90000	2408.4\\
92000	2378.3\\
94000	2319.7\\
96000	2357.2\\
98000	2360.7\\
1e+05	2377.3\\
};
\addplot [color=mycolor64, forget plot]
  table[row sep=crcr]{%
0	4395.3\\
2000	2996\\
4000	2993.3\\
6000	2950.8\\
8000	3056.2\\
10000	2930.9\\
12000	3011.6\\
14000	3008.3\\
16000	2907.6\\
18000	2928.1\\
22000	3028.4\\
24000	2952.1\\
26000	2912.5\\
28000	3021\\
30000	2955.2\\
34000	3039\\
36000	3007.8\\
38000	3019.1\\
40000	2958\\
42000	2983.9\\
44000	2963.2\\
46000	2993.7\\
48000	2978.3\\
50000	2989.7\\
52000	2961.2\\
54000	3032.3\\
56000	3023.5\\
58000	2995.1\\
60000	3016.6\\
62000	2969.8\\
64000	2982.1\\
66000	3012.8\\
68000	2980.3\\
70000	2979.8\\
72000	3045\\
74000	2941\\
76000	2957.7\\
78000	2938.6\\
80000	2964.3\\
82000	2975.2\\
84000	2962.8\\
86000	3018.5\\
88000	3044.3\\
90000	3048.5\\
92000	3031.2\\
94000	2942.3\\
96000	3009\\
98000	2945.3\\
1e+05	2995.1\\
};
\addplot [color=mycolor65, forget plot]
  table[row sep=crcr]{%
0	4395.3\\
2000	2634.1\\
4000	2542.5\\
6000	2471\\
8000	2431.6\\
10000	2449.2\\
12000	2500.4\\
14000	2443.3\\
16000	2490.3\\
18000	2454.5\\
22000	2500.9\\
26000	2470.1\\
28000	2474\\
30000	2516.6\\
32000	2480.4\\
34000	2430.3\\
38000	2517.1\\
40000	2471.3\\
42000	2449\\
44000	2513.3\\
46000	2467.3\\
48000	2512.9\\
52000	2475.8\\
54000	2529.1\\
56000	2478.2\\
58000	2585.7\\
60000	2509.7\\
62000	2487\\
64000	2446.4\\
66000	2428\\
68000	2543.9\\
70000	2501.9\\
72000	2487.3\\
74000	2500.9\\
76000	2464.7\\
78000	2451\\
80000	2476.8\\
82000	2455.1\\
84000	2510.1\\
86000	2543.5\\
88000	2506.4\\
90000	2431.8\\
92000	2474.3\\
94000	2578.4\\
96000	2477.9\\
98000	2430.4\\
1e+05	2489.1\\
};
\addplot [color=mycolor66, forget plot]
  table[row sep=crcr]{%
0	4395.3\\
2000	3391.4\\
4000	3382.5\\
6000	3404.6\\
8000	3362.4\\
10000	3406.5\\
12000	3316.4\\
14000	3409.6\\
16000	3410.4\\
18000	3338.5\\
20000	3381.1\\
22000	3390.1\\
24000	3376.1\\
26000	3412.4\\
28000	3326.7\\
30000	3364.2\\
32000	3317.8\\
34000	3348.5\\
36000	3390.2\\
38000	3408.2\\
40000	3465.2\\
42000	3389.2\\
44000	3401.3\\
46000	3338.8\\
48000	3371.2\\
50000	3352\\
52000	3397\\
54000	3357.8\\
56000	3377.5\\
58000	3317.2\\
60000	3376.2\\
64000	3372.2\\
66000	3415.7\\
70000	3344.2\\
72000	3350.1\\
74000	3398.4\\
76000	3367.7\\
78000	3381.4\\
80000	3369.6\\
82000	3394.7\\
84000	3369.8\\
86000	3453.5\\
88000	3332\\
90000	3386\\
92000	3404.2\\
94000	3385\\
96000	3413.2\\
1e+05	3399.3\\
};
\addplot [color=mycolor67, forget plot]
  table[row sep=crcr]{%
0	4395.3\\
2000	3370.6\\
4000	3129.9\\
6000	3154.4\\
8000	3104.3\\
10000	3135.8\\
12000	3079.3\\
14000	3100.6\\
16000	3209.7\\
18000	3139.5\\
20000	3131.9\\
22000	3068.6\\
24000	3145.9\\
26000	3180.4\\
28000	3118.5\\
30000	3111.3\\
32000	3141.6\\
34000	3099.7\\
36000	3142.4\\
38000	3090.8\\
40000	3156.8\\
42000	3085.8\\
44000	3039.1\\
46000	3099.6\\
48000	3177.8\\
50000	3150.1\\
52000	3087.2\\
54000	3168.1\\
56000	3100.1\\
58000	3195\\
60000	3120.6\\
64000	3147.8\\
66000	3093.1\\
68000	3120.4\\
70000	3159.5\\
72000	3147.8\\
74000	3057.6\\
76000	3150.9\\
78000	3167\\
80000	3120.2\\
82000	3165.5\\
84000	3122.3\\
86000	3101.2\\
88000	3202\\
90000	3148.7\\
92000	3134.4\\
94000	3077.7\\
96000	3184.6\\
98000	3199.7\\
1e+05	3141.7\\
};
\addplot [color=mycolor67]
  table[row sep=crcr]{%
0	4395.3\\
2000	2660.7\\
4000	2547.7\\
6000	2575.2\\
8000	2534.2\\
12000	2590.7\\
14000	2580.9\\
16000	2579.5\\
18000	2614.3\\
20000	2543.1\\
22000	2536.4\\
24000	2594.6\\
26000	2515.7\\
28000	2560.5\\
32000	2552.1\\
34000	2585.8\\
36000	2569.3\\
38000	2519.6\\
40000	2556.9\\
42000	2564.8\\
44000	2527.8\\
46000	2559.7\\
48000	2581.4\\
52000	2593.3\\
54000	2592.3\\
56000	2579.1\\
58000	2533.8\\
60000	2539.8\\
62000	2600.4\\
64000	2501.3\\
66000	2589.5\\
68000	2555.8\\
70000	2582.6\\
72000	2579.4\\
74000	2531.8\\
76000	2598.2\\
78000	2566.6\\
80000	2583.8\\
82000	2574.1\\
84000	2593\\
86000	2593.7\\
88000	2515.2\\
90000	2520.3\\
92000	2551.5\\
94000	2613.8\\
96000	2592.7\\
98000	2553.4\\
1e+05	2584\\
};
\addlegendentry{Classe k=43}

\addplot [color=mycolor68, forget plot]
  table[row sep=crcr]{%
0	4395.3\\
2000	2676.4\\
4000	2453\\
6000	2444.9\\
8000	2402.4\\
10000	2400\\
12000	2412.3\\
16000	2396.7\\
20000	2411.7\\
22000	2437.8\\
24000	2391.2\\
26000	2368.2\\
28000	2408.4\\
30000	2422\\
34000	2398.3\\
36000	2411.9\\
38000	2390.7\\
40000	2434.3\\
42000	2389.8\\
46000	2378.2\\
48000	2403.1\\
50000	2420.7\\
52000	2453\\
54000	2426.8\\
56000	2410.2\\
58000	2385.2\\
60000	2386.6\\
62000	2354.8\\
64000	2394.8\\
68000	2429.7\\
70000	2382\\
72000	2398.3\\
74000	2427\\
76000	2431.4\\
78000	2421\\
80000	2425\\
82000	2413.9\\
84000	2435.5\\
86000	2422.9\\
90000	2433.1\\
92000	2444.6\\
94000	2378.1\\
96000	2363\\
98000	2404.5\\
1e+05	2406.8\\
};
\addplot [color=mycolor69, forget plot]
  table[row sep=crcr]{%
0	4395.3\\
2000	2943.2\\
4000	2825.1\\
6000	2796.8\\
8000	2816\\
10000	2872\\
12000	2845.3\\
14000	2845.6\\
16000	2856.8\\
18000	2806.8\\
20000	2866.9\\
22000	2848.7\\
24000	2866.6\\
26000	2872.7\\
28000	2864\\
30000	2865.5\\
32000	2832.7\\
34000	2898.7\\
36000	2855.1\\
42000	2838.4\\
44000	2868.3\\
46000	2850.9\\
48000	2879.7\\
50000	2836.2\\
52000	2813.5\\
56000	2827.9\\
58000	2845.1\\
60000	2840.7\\
62000	2818.7\\
64000	2806.4\\
66000	2841.8\\
68000	2801\\
70000	2805.2\\
72000	2861\\
74000	2832.1\\
76000	2894.3\\
78000	2854.8\\
80000	2834.7\\
82000	2874.3\\
84000	2841.2\\
86000	2847\\
88000	2781.1\\
90000	2810.7\\
92000	2848.9\\
94000	2819.5\\
96000	2840.1\\
98000	2871.1\\
1e+05	2867.9\\
};
\addplot [color=mycolor70, forget plot]
  table[row sep=crcr]{%
0	4395.3\\
2000	2150.4\\
4000	2011\\
6000	2044\\
8000	2062.8\\
10000	2021.8\\
12000	2066.5\\
14000	2094.7\\
16000	2044.3\\
18000	2071.9\\
20000	2055.2\\
22000	2065.2\\
24000	2022.9\\
26000	2062.6\\
28000	2001\\
30000	1999.8\\
32000	2033.4\\
36000	2053.7\\
38000	2030.3\\
40000	2082.9\\
42000	2048.1\\
44000	2037.8\\
46000	2021.7\\
48000	2082.3\\
50000	2061.8\\
52000	2050.3\\
54000	2033.1\\
56000	2066.7\\
58000	2077.3\\
60000	2047.3\\
62000	2031.2\\
64000	2006.1\\
66000	2025\\
68000	2063.6\\
70000	2084.4\\
72000	2002.8\\
76000	2031.1\\
78000	2036.7\\
80000	2033.5\\
82000	2016.5\\
84000	2011.5\\
86000	2052\\
88000	2059.7\\
90000	2039.4\\
92000	2040.1\\
94000	2014.2\\
96000	2062.5\\
98000	2063.3\\
1e+05	2043.4\\
};
\addplot [color=mycolor71, forget plot]
  table[row sep=crcr]{%
0	4395.3\\
2000	2661.5\\
4000	2450.6\\
6000	2444.6\\
8000	2449\\
10000	2409.3\\
12000	2461.8\\
16000	2400.6\\
18000	2439.1\\
20000	2387.3\\
24000	2474.6\\
26000	2412.4\\
28000	2406.6\\
30000	2427.5\\
32000	2439.5\\
34000	2436.5\\
36000	2459.3\\
38000	2475.4\\
40000	2425.6\\
42000	2444.7\\
44000	2411.4\\
46000	2441\\
48000	2420.1\\
50000	2440.2\\
52000	2439\\
54000	2377.9\\
56000	2398.2\\
58000	2438.1\\
60000	2403.2\\
62000	2405.7\\
64000	2431.6\\
66000	2422.4\\
68000	2405.3\\
70000	2435.3\\
72000	2441.8\\
74000	2436.5\\
76000	2452\\
78000	2425.7\\
80000	2426.1\\
82000	2443.1\\
84000	2367.4\\
86000	2389.7\\
88000	2444.6\\
92000	2420.9\\
96000	2431.4\\
98000	2449.1\\
1e+05	2422.9\\
};
\addplot [color=mycolor72, forget plot]
  table[row sep=crcr]{%
0	4395.3\\
2000	2411.9\\
4000	2150.5\\
6000	2196.3\\
8000	2194.9\\
10000	2130.9\\
12000	2146.9\\
14000	2107.1\\
16000	2173.5\\
18000	2121.1\\
20000	2173.8\\
22000	2145.6\\
24000	2074.2\\
26000	2121.5\\
28000	2131.8\\
30000	2154.6\\
32000	2135.2\\
34000	2174.2\\
36000	2117.9\\
38000	2155.6\\
44000	2121.9\\
46000	2145.6\\
48000	2150.3\\
50000	2181.5\\
52000	2148.9\\
54000	2050.1\\
56000	2092.1\\
58000	2082.9\\
60000	2094\\
62000	2113.3\\
64000	2190.1\\
66000	2139.4\\
68000	2155.9\\
72000	2106.3\\
74000	2172.6\\
76000	2105.2\\
78000	2073.3\\
80000	2147.6\\
82000	2111.5\\
84000	2154.5\\
86000	2157.3\\
88000	2193.1\\
90000	2158.3\\
92000	2175.7\\
94000	2166.3\\
96000	2111.2\\
98000	2082.5\\
1e+05	2121.5\\
};
\addplot [color=mycolor73, forget plot]
  table[row sep=crcr]{%
0	4395.3\\
2000	2769.6\\
4000	2591.7\\
6000	2585.9\\
8000	2629.1\\
10000	2528.6\\
12000	2561.1\\
14000	2500.5\\
18000	2563.3\\
20000	2575.5\\
22000	2546\\
24000	2574.1\\
26000	2527.6\\
28000	2537.4\\
30000	2535.9\\
32000	2554.4\\
34000	2583\\
36000	2571.8\\
38000	2612.3\\
40000	2554.6\\
42000	2602.4\\
44000	2545.2\\
46000	2544.4\\
48000	2596.8\\
50000	2489.8\\
52000	2521.5\\
56000	2504.3\\
58000	2539.5\\
60000	2590.1\\
62000	2519.6\\
64000	2578.7\\
66000	2574.7\\
70000	2530.9\\
72000	2548.7\\
74000	2624.2\\
78000	2516.7\\
80000	2587\\
82000	2571\\
84000	2581.8\\
88000	2559.3\\
90000	2533.1\\
92000	2571.7\\
94000	2570.9\\
96000	2584.7\\
1e+05	2513.9\\
};
\addplot [color=mycolor74, forget plot]
  table[row sep=crcr]{%
0	4395.3\\
2000	1847.2\\
4000	1504.5\\
6000	1491.2\\
8000	1509.8\\
10000	1462.9\\
12000	1490\\
14000	1509.4\\
16000	1456.6\\
18000	1480.2\\
20000	1489.6\\
22000	1507.4\\
26000	1493.1\\
28000	1484.9\\
30000	1467.7\\
32000	1482.7\\
34000	1462.5\\
38000	1464.9\\
40000	1482.1\\
42000	1468.4\\
44000	1493\\
46000	1500.9\\
48000	1452.4\\
50000	1505.8\\
52000	1477.4\\
54000	1491\\
56000	1477.4\\
58000	1482.2\\
60000	1498.3\\
62000	1504.7\\
66000	1486.1\\
68000	1484\\
70000	1456.8\\
72000	1489.6\\
74000	1484.9\\
76000	1466.4\\
78000	1468.6\\
80000	1496.3\\
82000	1461.3\\
84000	1489\\
88000	1479.4\\
90000	1500.7\\
92000	1443.2\\
94000	1500.9\\
96000	1489.6\\
98000	1452.4\\
1e+05	1490.2\\
};
\addplot [color=mycolor75, forget plot]
  table[row sep=crcr]{%
0	4395.3\\
2000	2416.8\\
4000	2162.5\\
6000	2127.9\\
8000	2165.4\\
10000	2151.8\\
12000	2120\\
14000	2102.7\\
16000	2075.9\\
18000	2146.8\\
20000	2131.9\\
22000	2188.2\\
24000	2152.6\\
26000	2187\\
30000	2139.3\\
32000	2124.1\\
34000	2177.5\\
38000	2121.9\\
40000	2152.5\\
42000	2153.4\\
44000	2111.7\\
46000	2162.5\\
48000	2154.2\\
50000	2152\\
52000	2089.5\\
54000	2152\\
56000	2169.8\\
58000	2162.1\\
60000	2110.9\\
62000	2135.9\\
64000	2139.2\\
66000	2185.8\\
68000	2141.4\\
70000	2137.2\\
74000	2162.8\\
76000	2135.2\\
78000	2171.3\\
80000	2131.6\\
82000	2106.6\\
84000	2128.7\\
86000	2143.7\\
88000	2186.6\\
90000	2147.1\\
92000	2177.2\\
94000	2135.3\\
98000	2146.7\\
1e+05	2152.4\\
};
\addplot [color=mycolor76, forget plot]
  table[row sep=crcr]{%
0	4395.3\\
2000	2238.5\\
4000	1978.9\\
6000	1998.8\\
8000	1963.2\\
10000	1969.2\\
12000	1956.3\\
14000	1977.6\\
16000	1967.9\\
18000	1933.3\\
22000	1943.7\\
24000	1957.3\\
26000	1962.5\\
28000	1948.3\\
30000	1967.6\\
32000	1993.5\\
34000	1919.5\\
36000	1965.6\\
38000	1945.8\\
40000	1946.8\\
42000	1957.1\\
44000	1934\\
48000	2011.4\\
50000	1920.2\\
52000	1929.4\\
54000	1912.3\\
56000	1964.4\\
58000	1947.9\\
60000	1951.4\\
62000	1922.6\\
64000	1956.2\\
66000	1962.1\\
68000	1955.6\\
70000	1957.9\\
72000	1995\\
74000	1944.2\\
76000	1938.2\\
78000	1973.8\\
80000	1953.3\\
82000	1977.1\\
84000	1961.6\\
86000	1936.5\\
88000	1964.8\\
90000	1957.1\\
92000	1936\\
94000	1972.7\\
96000	1959.5\\
98000	1964\\
1e+05	1959.9\\
};
\addplot [color=mycolor77, forget plot]
  table[row sep=crcr]{%
0	4395.3\\
2000	1875.3\\
4000	1553.5\\
6000	1490.1\\
8000	1525.6\\
10000	1496.7\\
12000	1513\\
14000	1489.4\\
16000	1495.6\\
18000	1510.5\\
20000	1484.9\\
22000	1511.2\\
26000	1510.6\\
28000	1487.9\\
30000	1501.5\\
34000	1510.7\\
38000	1494.6\\
40000	1515.4\\
42000	1514.8\\
44000	1493\\
46000	1511.3\\
48000	1499.8\\
50000	1523.9\\
52000	1507.8\\
54000	1485.4\\
56000	1486.8\\
58000	1504\\
60000	1476.4\\
62000	1511.3\\
66000	1504.8\\
68000	1491.3\\
70000	1519.1\\
72000	1501\\
74000	1490.2\\
78000	1536.8\\
80000	1533.2\\
82000	1491.6\\
84000	1537.1\\
86000	1537.1\\
88000	1551.8\\
90000	1526.1\\
92000	1529.9\\
94000	1489.2\\
98000	1495.9\\
1e+05	1491.4\\
};
\addplot [color=mycolor78, forget plot]
  table[row sep=crcr]{%
0	4395.3\\
2000	2553.5\\
4000	2288.6\\
6000	2206\\
8000	2183.9\\
10000	2235\\
12000	2260.1\\
14000	2266\\
16000	2217.6\\
18000	2211.7\\
20000	2233\\
22000	2205.3\\
24000	2233.5\\
26000	2238.1\\
28000	2249.4\\
30000	2285\\
32000	2264.2\\
34000	2210.5\\
36000	2191.6\\
38000	2227.4\\
40000	2203.5\\
42000	2226\\
44000	2255.3\\
46000	2242.2\\
48000	2260.1\\
50000	2246.7\\
52000	2224.4\\
54000	2244.2\\
56000	2190.1\\
58000	2246.3\\
60000	2221.3\\
62000	2258.5\\
64000	2238.5\\
66000	2250.2\\
68000	2233\\
70000	2253.8\\
72000	2260.1\\
74000	2291.3\\
76000	2247.2\\
78000	2254.1\\
82000	2188\\
84000	2226.9\\
86000	2226.3\\
88000	2286.8\\
92000	2241.8\\
94000	2266\\
98000	2226.4\\
1e+05	2251.4\\
};
\addplot [color=mycolor79, forget plot]
  table[row sep=crcr]{%
0	4395.3\\
2000	2382.6\\
4000	2184.9\\
6000	2105\\
8000	2081.9\\
10000	2106.7\\
14000	2101.9\\
16000	2122.8\\
18000	2047.5\\
20000	2042.5\\
22000	2089.9\\
24000	2058.1\\
26000	2074.6\\
28000	2105.4\\
30000	2080.4\\
32000	2096.4\\
34000	2092.4\\
36000	2098.5\\
40000	2122.8\\
42000	2121.2\\
44000	2111.1\\
46000	2124.1\\
48000	2095.1\\
50000	2099.1\\
52000	2146.1\\
56000	2054.8\\
58000	2085\\
60000	2063\\
62000	2059.1\\
64000	2083\\
66000	2044.2\\
68000	2111.5\\
70000	2128.7\\
72000	2104.9\\
74000	2106.8\\
76000	2140.6\\
78000	2098\\
80000	2071.2\\
82000	2136.3\\
84000	2120.6\\
86000	2130.8\\
88000	2134.7\\
90000	2148.9\\
92000	2086.7\\
94000	2095.5\\
96000	2078.6\\
1e+05	2093.5\\
};
\addplot [color=mycolor79, forget plot]
  table[row sep=crcr]{%
0	4395.3\\
2000	2347.4\\
4000	1948.8\\
6000	1955.7\\
8000	1935.6\\
10000	1963.8\\
12000	2041.3\\
14000	1945.3\\
16000	1965.5\\
18000	1942.6\\
20000	1954.1\\
24000	2014.9\\
26000	1975\\
28000	1979.8\\
30000	2005.7\\
32000	1971.2\\
34000	1966.3\\
36000	1983.1\\
38000	1938.9\\
40000	1993.6\\
42000	1927.5\\
44000	1973.2\\
46000	1944.1\\
48000	1940.5\\
50000	1966.2\\
54000	1907.1\\
58000	1932.7\\
62000	1948.9\\
64000	1909.8\\
66000	1949.8\\
68000	1917.2\\
70000	1973.3\\
72000	1957.8\\
74000	1901.7\\
76000	1930.2\\
78000	1950.9\\
80000	1930.4\\
82000	1940.3\\
84000	1968.8\\
86000	1962.7\\
88000	1994.1\\
90000	1948.9\\
92000	1953.7\\
94000	1928.2\\
96000	1969.8\\
98000	1944\\
1e+05	1973.2\\
};
\addplot [color=mycolor80, forget plot]
  table[row sep=crcr]{%
0	4395.3\\
2000	1828.4\\
4000	1491.7\\
6000	1501.7\\
8000	1475.8\\
10000	1500.2\\
12000	1473\\
14000	1516.8\\
16000	1492.3\\
18000	1483.6\\
20000	1529.1\\
22000	1497.7\\
24000	1480.5\\
26000	1490.2\\
28000	1518.3\\
30000	1509.2\\
32000	1485.2\\
36000	1503.1\\
38000	1476.3\\
40000	1484.3\\
42000	1507.7\\
46000	1478\\
48000	1477.5\\
50000	1486.9\\
52000	1470.9\\
54000	1498.3\\
56000	1501.1\\
58000	1499.2\\
60000	1508.5\\
62000	1472.6\\
66000	1522.9\\
68000	1478.7\\
70000	1486.8\\
72000	1490.6\\
74000	1505.1\\
76000	1508.2\\
78000	1477.2\\
80000	1501.7\\
82000	1513.1\\
84000	1507.9\\
86000	1489.4\\
88000	1519.7\\
90000	1514.3\\
94000	1472.9\\
96000	1508.7\\
98000	1502.2\\
1e+05	1500.8\\
};
\addplot [color=mycolor81, forget plot]
  table[row sep=crcr]{%
0	4395.3\\
2000	1780.9\\
4000	1404.8\\
6000	1390.3\\
8000	1402.3\\
10000	1384.7\\
12000	1419.3\\
14000	1374.4\\
16000	1386.4\\
18000	1405.9\\
20000	1368.6\\
22000	1410.2\\
24000	1399.8\\
26000	1395.2\\
28000	1395.5\\
30000	1388.1\\
34000	1388\\
36000	1380.9\\
38000	1382.4\\
40000	1417.5\\
42000	1396\\
44000	1406.6\\
46000	1402.3\\
48000	1378.5\\
52000	1395.8\\
54000	1392.1\\
56000	1407.6\\
58000	1399.1\\
60000	1410.3\\
62000	1397.6\\
64000	1420.3\\
68000	1368.2\\
70000	1389.1\\
72000	1392.7\\
76000	1376.9\\
78000	1389.9\\
80000	1397.7\\
82000	1371.5\\
84000	1381.8\\
86000	1423.5\\
88000	1406.2\\
90000	1393.7\\
92000	1397.8\\
94000	1388.1\\
96000	1389.2\\
98000	1396.7\\
1e+05	1397.6\\
};
\addplot [color=mycolor82, forget plot]
  table[row sep=crcr]{%
0	4395.3\\
2000	2024.7\\
4000	1687.7\\
6000	1635\\
8000	1664.3\\
10000	1629\\
12000	1675.9\\
14000	1679.5\\
18000	1613.7\\
20000	1657.7\\
22000	1676.9\\
24000	1673.5\\
26000	1631.3\\
28000	1670.3\\
30000	1618.5\\
32000	1636.9\\
34000	1630.1\\
36000	1648.4\\
38000	1692.2\\
40000	1681.7\\
42000	1622.4\\
44000	1616\\
46000	1650.3\\
48000	1666\\
50000	1674.6\\
52000	1645.8\\
54000	1664.3\\
56000	1628.5\\
58000	1606\\
60000	1622.3\\
62000	1645\\
64000	1650.6\\
66000	1639.8\\
68000	1645.2\\
70000	1609.9\\
72000	1628.4\\
74000	1641.1\\
76000	1634.9\\
78000	1647.4\\
80000	1627.9\\
82000	1630\\
84000	1597.9\\
86000	1647\\
88000	1669.6\\
90000	1631.8\\
92000	1635.4\\
94000	1658.8\\
96000	1620.7\\
98000	1617.5\\
1e+05	1605.1\\
};
\addplot [color=mycolor83, forget plot]
  table[row sep=crcr]{%
0	4395.3\\
2000	2130.6\\
4000	1691.8\\
6000	1676.1\\
8000	1648.6\\
10000	1636.5\\
12000	1668.6\\
14000	1668\\
16000	1685.5\\
18000	1637\\
20000	1639.1\\
22000	1655.2\\
24000	1687.2\\
26000	1696.1\\
28000	1667\\
30000	1657.4\\
32000	1682.2\\
34000	1691.9\\
36000	1635.3\\
38000	1673.3\\
40000	1675\\
42000	1684\\
44000	1653.6\\
46000	1661.5\\
48000	1640.7\\
50000	1654.6\\
52000	1659.5\\
54000	1669.6\\
56000	1639.2\\
58000	1662\\
60000	1656.3\\
62000	1669.4\\
64000	1661.6\\
66000	1662.3\\
68000	1656.8\\
70000	1697.5\\
72000	1649.1\\
74000	1676.3\\
76000	1659.1\\
78000	1661\\
80000	1684.9\\
82000	1675.8\\
84000	1647.9\\
86000	1677.1\\
90000	1688.6\\
92000	1705.9\\
94000	1675.8\\
96000	1680.9\\
98000	1656.3\\
1e+05	1652\\
};
\addplot [color=mycolor83, forget plot]
  table[row sep=crcr]{%
0	4395.3\\
2000	1535.8\\
4000	1210.6\\
6000	1206.6\\
10000	1227.2\\
12000	1225.4\\
14000	1216.4\\
22000	1200.7\\
24000	1187.2\\
26000	1186.1\\
30000	1217.6\\
34000	1211.7\\
36000	1198.2\\
44000	1218.3\\
48000	1201.7\\
50000	1229.3\\
52000	1218\\
54000	1203\\
56000	1221\\
58000	1213.6\\
60000	1211.6\\
62000	1199.2\\
64000	1207.6\\
66000	1232.9\\
68000	1221.6\\
70000	1237\\
72000	1242.9\\
74000	1220.7\\
76000	1218.7\\
78000	1227\\
80000	1178.2\\
82000	1208.4\\
84000	1195.1\\
86000	1192.7\\
88000	1229.8\\
92000	1230.5\\
94000	1210.4\\
96000	1204.6\\
98000	1207.3\\
1e+05	1198.1\\
};
\addplot [color=mycolor84, forget plot]
  table[row sep=crcr]{%
0	4395.3\\
2000	1597.3\\
4000	1230.3\\
6000	1243.4\\
8000	1200.6\\
10000	1205.4\\
12000	1238.2\\
14000	1211.5\\
16000	1227.4\\
18000	1197.2\\
20000	1189.4\\
24000	1211.7\\
26000	1207\\
28000	1232.2\\
30000	1220.4\\
32000	1221.2\\
34000	1176.8\\
36000	1218.5\\
38000	1220.9\\
40000	1194\\
42000	1193.6\\
44000	1208.8\\
46000	1189.9\\
48000	1202\\
50000	1194.8\\
52000	1231.7\\
54000	1218.6\\
56000	1241.4\\
58000	1236.1\\
60000	1206.2\\
62000	1181.1\\
64000	1188.5\\
66000	1247.3\\
68000	1221.3\\
70000	1181.9\\
72000	1234.4\\
74000	1232.3\\
76000	1225.2\\
78000	1167.7\\
80000	1216.1\\
82000	1229.8\\
84000	1227.3\\
86000	1184.1\\
88000	1204.5\\
90000	1190.3\\
92000	1238.1\\
94000	1219.1\\
96000	1247.6\\
98000	1224.1\\
1e+05	1214.9\\
};
\addplot [color=mycolor85, forget plot]
  table[row sep=crcr]{%
0	4395.3\\
2000	1781.8\\
4000	1313.9\\
6000	1358.4\\
8000	1376.6\\
10000	1351.9\\
12000	1360.7\\
14000	1336.5\\
16000	1367.2\\
18000	1310.2\\
20000	1353.4\\
24000	1315.7\\
26000	1363.2\\
28000	1397.1\\
32000	1330\\
34000	1366.7\\
36000	1341.2\\
40000	1340.2\\
42000	1348.8\\
44000	1379.6\\
48000	1337.5\\
50000	1397.8\\
52000	1338.8\\
54000	1315.3\\
56000	1324.6\\
58000	1392.6\\
60000	1339.6\\
62000	1388\\
64000	1351.6\\
66000	1328.7\\
70000	1341\\
72000	1368.7\\
74000	1338.9\\
76000	1350.2\\
78000	1367.8\\
80000	1365.7\\
82000	1345.1\\
84000	1417.5\\
86000	1356.8\\
90000	1355.6\\
92000	1363.4\\
94000	1394.1\\
96000	1327.6\\
98000	1306.4\\
1e+05	1403.4\\
};
\addplot [color=mycolor86, forget plot]
  table[row sep=crcr]{%
0	4395.3\\
2000	1597.7\\
4000	1131.4\\
6000	1140.3\\
8000	1095.7\\
10000	1102.2\\
12000	1143.5\\
14000	1087\\
16000	1140.6\\
18000	1162.6\\
20000	1124.6\\
22000	1113.4\\
24000	1116.2\\
26000	1111.1\\
28000	1123.6\\
30000	1145.4\\
32000	1133.7\\
34000	1113.7\\
36000	1146.2\\
42000	1128.9\\
44000	1131.9\\
46000	1123.7\\
48000	1108.1\\
50000	1103.4\\
52000	1091.6\\
54000	1087.9\\
56000	1123.6\\
60000	1119\\
62000	1118.6\\
64000	1149\\
68000	1133\\
70000	1127.3\\
72000	1127.7\\
74000	1095.4\\
76000	1125.1\\
78000	1144.2\\
80000	1124.9\\
82000	1121.4\\
84000	1126.9\\
86000	1079.3\\
88000	1137\\
90000	1141.2\\
92000	1114.7\\
94000	1117.4\\
96000	1140.3\\
98000	1175\\
1e+05	1127.9\\
};
\addplot [color=mycolor87, forget plot]
  table[row sep=crcr]{%
0	4395.3\\
2000	1364.3\\
4000	891.59\\
6000	941.55\\
8000	912.19\\
10000	985.39\\
12000	939.15\\
14000	943.44\\
18000	906.11\\
20000	981.59\\
22000	951.76\\
24000	944.51\\
28000	944.32\\
30000	911.26\\
32000	932.32\\
34000	930.48\\
36000	924.85\\
38000	933.8\\
40000	969.1\\
42000	979.61\\
44000	942.13\\
46000	965.59\\
48000	946.4\\
50000	921.73\\
52000	942.41\\
54000	887.08\\
56000	938.01\\
58000	962.1\\
60000	967.22\\
62000	920.14\\
64000	904.99\\
66000	954.63\\
68000	900.19\\
70000	897.66\\
72000	911.97\\
74000	891.37\\
76000	952.24\\
80000	921.91\\
82000	870.95\\
84000	944.92\\
86000	893.98\\
88000	871.16\\
90000	918.21\\
92000	939.8\\
94000	897.81\\
96000	978.06\\
98000	925.69\\
1e+05	919.56\\
};
\addplot [color=mycolor88, forget plot]
  table[row sep=crcr]{%
0	4395.3\\
2000	1573.7\\
4000	1183.2\\
6000	1128.7\\
8000	1127.5\\
10000	1113.3\\
12000	1149.9\\
14000	1115.3\\
16000	1138.3\\
18000	1123.9\\
20000	1134.1\\
24000	1137.7\\
26000	1119.7\\
28000	1132.3\\
30000	1116.3\\
32000	1103.8\\
34000	1164\\
36000	1092.6\\
38000	1132.4\\
40000	1124\\
42000	1120.7\\
44000	1099.8\\
46000	1085\\
48000	1136.1\\
50000	1123\\
52000	1095\\
54000	1125.9\\
56000	1137.9\\
58000	1145.6\\
60000	1140.4\\
62000	1121.1\\
66000	1102.8\\
70000	1137.4\\
72000	1112.9\\
74000	1116.6\\
76000	1125\\
78000	1097.4\\
80000	1134.1\\
82000	1135.6\\
86000	1128.5\\
88000	1151.6\\
90000	1162.6\\
92000	1185.5\\
94000	1112.7\\
96000	1102.3\\
98000	1112.9\\
1e+05	1088.3\\
};
\addplot [color=mycolor89, forget plot]
  table[row sep=crcr]{%
0	4395.3\\
2000	1428.8\\
4000	807.2\\
8000	785.67\\
10000	798.74\\
12000	791.69\\
14000	812.89\\
16000	801.73\\
18000	777.97\\
20000	812.25\\
22000	780.62\\
24000	822.29\\
26000	778.62\\
28000	795.76\\
30000	784.61\\
32000	789.02\\
34000	774.52\\
36000	784.59\\
38000	822.33\\
40000	820.06\\
42000	805.82\\
44000	794.09\\
46000	803.36\\
48000	804.24\\
50000	813.38\\
52000	815.45\\
54000	780.4\\
56000	808.14\\
58000	831.74\\
60000	779.86\\
62000	794.4\\
64000	795.75\\
66000	810.11\\
68000	815.18\\
70000	829.01\\
72000	795.96\\
74000	790.55\\
78000	812.57\\
80000	814.08\\
84000	805.62\\
88000	808.14\\
90000	821.92\\
92000	829.36\\
94000	801.89\\
98000	801.19\\
1e+05	817.78\\
};
\addplot [color=mycolor90, forget plot]
  table[row sep=crcr]{%
0	4395.3\\
2000	1229.5\\
4000	777.36\\
6000	751.1\\
8000	775.47\\
10000	790.03\\
12000	761.92\\
14000	770.28\\
16000	764.6\\
18000	785.18\\
20000	750.43\\
22000	773.02\\
24000	808.18\\
26000	748.26\\
30000	783.49\\
32000	762.82\\
34000	785.68\\
36000	803.1\\
38000	800.92\\
40000	750.54\\
42000	748.25\\
44000	773.75\\
46000	784.28\\
48000	761.79\\
50000	778.25\\
52000	786.63\\
54000	768.35\\
56000	775.51\\
58000	790.73\\
60000	780.92\\
62000	763.83\\
64000	778.7\\
66000	767.56\\
68000	783.19\\
70000	787.61\\
74000	774.12\\
76000	787.6\\
78000	774.4\\
80000	771.66\\
82000	771.7\\
84000	762.62\\
86000	743.4\\
88000	776.43\\
90000	763.77\\
92000	765.33\\
94000	775.77\\
96000	761.94\\
98000	782.36\\
1e+05	771.2\\
};
\addplot [color=mycolor91, forget plot]
  table[row sep=crcr]{%
0	4395.3\\
2000	1274.2\\
4000	667.97\\
6000	657.54\\
8000	665.59\\
10000	675.75\\
12000	662.82\\
14000	652.47\\
16000	664.64\\
18000	655.55\\
20000	689.06\\
22000	701.75\\
24000	657.56\\
26000	662.41\\
28000	716.86\\
30000	649.18\\
32000	693.47\\
34000	688.26\\
36000	647.41\\
38000	678.89\\
40000	666.21\\
42000	691.74\\
44000	648.47\\
46000	628.74\\
48000	696.08\\
50000	675.96\\
52000	653.11\\
54000	678.57\\
56000	689.69\\
58000	660.29\\
60000	669.91\\
62000	692.13\\
64000	672.77\\
70000	676.77\\
72000	667.24\\
74000	668.97\\
76000	689.37\\
78000	666.06\\
80000	665.37\\
82000	682.88\\
84000	672.81\\
86000	656.41\\
88000	660.29\\
90000	655.86\\
92000	688.87\\
94000	660.95\\
96000	666.01\\
98000	651.66\\
1e+05	701.61\\
};
\addplot [color=mycolor92, forget plot]
  table[row sep=crcr]{%
0	4395.3\\
2000	1409\\
4000	844.17\\
6000	839.34\\
8000	851.23\\
10000	848.1\\
12000	818.49\\
14000	850.61\\
16000	841.19\\
18000	850.3\\
20000	828.87\\
22000	883.08\\
24000	870.49\\
26000	838.02\\
28000	828.64\\
30000	846.18\\
32000	847.5\\
34000	856.12\\
36000	857.17\\
38000	842.01\\
40000	820.32\\
42000	833.44\\
44000	810.94\\
46000	821.52\\
48000	841.33\\
50000	847.48\\
52000	834.94\\
54000	846.47\\
56000	820.5\\
58000	792.59\\
60000	854.41\\
62000	813.43\\
64000	864.52\\
66000	851.95\\
68000	842.13\\
70000	882.47\\
72000	847.48\\
74000	852.3\\
76000	828.93\\
78000	852.46\\
80000	840.42\\
82000	850.52\\
84000	826.13\\
86000	833.13\\
88000	853.73\\
90000	833.08\\
92000	856.02\\
94000	854.49\\
96000	858.64\\
98000	855.43\\
1e+05	874.01\\
};
\addplot [color=mycolor93, forget plot]
  table[row sep=crcr]{%
0	4395.3\\
2000	1359.9\\
4000	706.84\\
6000	754.51\\
8000	685.68\\
10000	652.13\\
12000	738.33\\
14000	727.53\\
16000	699.21\\
20000	699.43\\
22000	696.84\\
24000	690.99\\
26000	735.55\\
28000	695.38\\
30000	677.09\\
32000	708.39\\
36000	696.47\\
38000	744.81\\
40000	733.37\\
42000	661.71\\
44000	696.09\\
46000	702.26\\
48000	720.62\\
50000	712.74\\
52000	698.04\\
54000	708.23\\
56000	738.23\\
58000	738.65\\
60000	695.5\\
62000	705.8\\
64000	722.18\\
66000	722.01\\
68000	708.72\\
70000	683.34\\
72000	697.71\\
74000	695.08\\
76000	700.36\\
78000	727.25\\
80000	669.05\\
82000	700.67\\
84000	692.13\\
86000	668.62\\
88000	717.95\\
90000	705.1\\
92000	740.55\\
94000	689.02\\
96000	710.32\\
98000	702.1\\
1e+05	707.91\\
};
\addplot [color=mycolor94, forget plot]
  table[row sep=crcr]{%
0	4395.3\\
2000	1362.7\\
4000	640.44\\
6000	658.67\\
8000	649.92\\
10000	643.14\\
12000	661.63\\
14000	662.13\\
16000	658.92\\
18000	675.05\\
20000	666.44\\
22000	679.99\\
24000	656.92\\
26000	650.73\\
28000	667.2\\
30000	643.47\\
32000	664.11\\
34000	669.5\\
36000	652.84\\
38000	680.94\\
40000	660.36\\
42000	671\\
44000	646.03\\
46000	646.94\\
48000	659.4\\
52000	672.26\\
54000	633.46\\
56000	658.58\\
58000	641.24\\
60000	673.5\\
62000	661.97\\
64000	642.02\\
66000	660.46\\
68000	651.27\\
70000	669.71\\
72000	665.45\\
74000	670.34\\
78000	667.99\\
80000	658.25\\
82000	696.09\\
84000	641.33\\
86000	646.22\\
88000	682.52\\
90000	685.94\\
92000	661.67\\
94000	663.76\\
96000	679.45\\
98000	687.97\\
1e+05	654.85\\
};
\addplot [color=mycolor95, forget plot]
  table[row sep=crcr]{%
0	4395.3\\
2000	1059.6\\
4000	401.68\\
6000	394.07\\
8000	374.4\\
10000	374.5\\
12000	388.3\\
14000	383.87\\
16000	387.27\\
18000	407.93\\
20000	384.88\\
22000	408.27\\
24000	392.12\\
26000	397.39\\
28000	406.92\\
30000	402.42\\
32000	381.27\\
34000	390.59\\
36000	403.17\\
38000	383.79\\
40000	401.08\\
42000	372.78\\
44000	378.11\\
46000	406.01\\
48000	401.87\\
50000	369.46\\
52000	375.81\\
54000	392.65\\
56000	421.74\\
58000	386.29\\
60000	406.51\\
62000	387.39\\
64000	375.57\\
66000	394.51\\
68000	380.13\\
70000	394.02\\
72000	391.02\\
74000	395.61\\
76000	383.92\\
78000	397.46\\
80000	403.62\\
82000	390.87\\
84000	380.55\\
86000	405.48\\
88000	400.64\\
90000	401.01\\
92000	396.93\\
94000	403.8\\
96000	379.61\\
98000	370.52\\
1e+05	401.5\\
};
\addplot [color=mycolor96, forget plot]
  table[row sep=crcr]{%
0	4395.3\\
2000	1126.4\\
4000	392.13\\
6000	391.44\\
8000	451.11\\
10000	423.85\\
12000	396.9\\
14000	390.84\\
16000	398.62\\
18000	402.96\\
20000	395.77\\
22000	401.02\\
24000	408.87\\
26000	408.98\\
28000	417.9\\
30000	392.51\\
32000	391.9\\
34000	418.46\\
36000	404.94\\
38000	412.23\\
40000	411.57\\
42000	433.47\\
44000	420.9\\
46000	439.98\\
48000	426.94\\
50000	421.91\\
52000	431\\
54000	414.33\\
56000	409\\
58000	431.94\\
60000	457.99\\
62000	448.15\\
64000	405.72\\
66000	396.84\\
68000	446.58\\
70000	436.21\\
72000	431.59\\
74000	388.06\\
76000	389.4\\
78000	407.28\\
80000	411.93\\
82000	399.84\\
84000	410.83\\
86000	442.37\\
88000	421.77\\
90000	424.8\\
92000	380.5\\
94000	399.98\\
96000	403.71\\
98000	388.15\\
1e+05	407.55\\
};
\addplot [color=mycolor97, forget plot]
  table[row sep=crcr]{%
0	4395.3\\
2000	913.59\\
4000	303.92\\
6000	321.63\\
8000	279.23\\
10000	279.13\\
12000	298.28\\
14000	312.68\\
16000	311.7\\
20000	279.67\\
24000	306.16\\
26000	325.69\\
28000	285.02\\
30000	284.79\\
32000	315.28\\
34000	288.81\\
36000	287.17\\
38000	303.74\\
40000	290.27\\
42000	311.41\\
44000	303.39\\
46000	293.59\\
48000	299.13\\
50000	300.85\\
52000	297.04\\
54000	282.83\\
56000	290.3\\
58000	276.59\\
60000	333.34\\
62000	277.17\\
64000	303.89\\
66000	308.06\\
68000	288.13\\
70000	282.17\\
72000	340.87\\
74000	284.3\\
76000	272.28\\
78000	316.13\\
80000	330.58\\
82000	320.57\\
84000	295.31\\
86000	331.41\\
88000	321.72\\
90000	294.53\\
92000	295.68\\
94000	304.06\\
96000	311.16\\
98000	276.22\\
1e+05	282.83\\
};
\addplot [color=mycolor98]
  table[row sep=crcr]{%
0	4395.3\\
2000	1214.1\\
4000	443.92\\
6000	397.83\\
8000	392.66\\
10000	383.23\\
12000	405.92\\
14000	425.59\\
16000	409.17\\
18000	383.18\\
20000	372.84\\
22000	403.63\\
24000	418.06\\
26000	397.49\\
28000	389.48\\
30000	406.85\\
32000	411.31\\
34000	388.27\\
36000	373.2\\
38000	404.11\\
40000	420.49\\
42000	383.29\\
44000	395.31\\
46000	396.76\\
48000	391.96\\
50000	412.12\\
52000	383.84\\
54000	402.05\\
56000	407.22\\
58000	386.7\\
60000	383.51\\
62000	382.2\\
64000	393.54\\
66000	369.08\\
68000	384.55\\
70000	408.57\\
72000	388.52\\
76000	395.33\\
78000	387.28\\
80000	407.59\\
82000	399.92\\
84000	389.84\\
86000	390.26\\
88000	400.68\\
90000	376.45\\
92000	423.67\\
94000	379.63\\
96000	390.73\\
98000	397.66\\
1e+05	408.74\\
};
\addlegendentry{Classe k=10}

            % \input{../../../../.tex/20/KS/SizeEvolutionsfig9.tex}
            \legend{}

            \nextgroupplot[%
                % xmin=0,xmax=1e+05,
                % ymin=100,ymax=1.6196e+05,
            ] % This file was created by matlab2tikz.
%
\definecolor{mycolor1}{rgb}{0.00146,0.00047,0.01387}%
\definecolor{mycolor2}{rgb}{0.01166,0.00942,0.06346}%
\definecolor{mycolor3}{rgb}{0.02943,0.02150,0.11462}%
\definecolor{mycolor4}{rgb}{0.06134,0.03659,0.17764}%
\definecolor{mycolor5}{rgb}{0.08741,0.04456,0.22481}%
\definecolor{mycolor6}{rgb}{0.12291,0.04754,0.28162}%
\definecolor{mycolor7}{rgb}{0.14907,0.04547,0.31709}%
\definecolor{mycolor8}{rgb}{0.18343,0.04033,0.35497}%
\definecolor{mycolor9}{rgb}{0.21795,0.03662,0.38352}%
\definecolor{mycolor10}{rgb}{0.24497,0.03705,0.40001}%
\definecolor{mycolor11}{rgb}{0.27135,0.04092,0.41198}%
\definecolor{mycolor12}{rgb}{0.29076,0.04564,0.41864}%
\definecolor{mycolor13}{rgb}{0.31628,0.05349,0.42512}%
\definecolor{mycolor14}{rgb}{0.33522,0.06006,0.42852}%
\definecolor{mycolor15}{rgb}{0.36028,0.06925,0.43150}%
\definecolor{mycolor16}{rgb}{0.37900,0.07625,0.43272}%
\definecolor{mycolor17}{rgb}{0.39767,0.08326,0.43318}%
\definecolor{mycolor18}{rgb}{0.41633,0.09020,0.43294}%
\definecolor{mycolor19}{rgb}{0.42877,0.09479,0.43241}%
\definecolor{mycolor20}{rgb}{0.44743,0.10160,0.43108}%
\definecolor{mycolor21}{rgb}{0.46610,0.10832,0.42912}%
\definecolor{mycolor22}{rgb}{0.47856,0.11276,0.42747}%
\definecolor{mycolor23}{rgb}{0.49726,0.11938,0.42449}%
\definecolor{mycolor24}{rgb}{0.50973,0.12377,0.42216}%
\definecolor{mycolor25}{rgb}{0.52221,0.12815,0.41955}%
\definecolor{mycolor26}{rgb}{0.53468,0.13253,0.41667}%
\definecolor{mycolor27}{rgb}{0.55339,0.13913,0.41183}%
\definecolor{mycolor28}{rgb}{0.56585,0.14357,0.40826}%
\definecolor{mycolor29}{rgb}{0.57830,0.14804,0.40441}%
\definecolor{mycolor30}{rgb}{0.59073,0.15256,0.40029}%
\definecolor{mycolor31}{rgb}{0.60314,0.15715,0.39589}%
\definecolor{mycolor32}{rgb}{0.61551,0.16182,0.39122}%
\definecolor{mycolor33}{rgb}{0.62169,0.16418,0.38878}%
\definecolor{mycolor34}{rgb}{0.63400,0.16899,0.38370}%
\definecolor{mycolor35}{rgb}{0.64626,0.17391,0.37836}%
\definecolor{mycolor36}{rgb}{0.65846,0.17896,0.37275}%
\definecolor{mycolor37}{rgb}{0.66454,0.18154,0.36985}%
\definecolor{mycolor38}{rgb}{0.67664,0.18681,0.36385}%
\definecolor{mycolor39}{rgb}{0.68865,0.19224,0.35760}%
\definecolor{mycolor40}{rgb}{0.69463,0.19502,0.35439}%
\definecolor{mycolor41}{rgb}{0.70650,0.20073,0.34778}%
\definecolor{mycolor42}{rgb}{0.71240,0.20366,0.34438}%
\definecolor{mycolor43}{rgb}{0.72410,0.20967,0.33742}%
\definecolor{mycolor44}{rgb}{0.72991,0.21276,0.33386}%
\definecolor{mycolor45}{rgb}{0.74142,0.21911,0.32658}%
\definecolor{mycolor46}{rgb}{0.74713,0.22238,0.32286}%
\definecolor{mycolor47}{rgb}{0.75279,0.22571,0.31909}%
\definecolor{mycolor48}{rgb}{0.76401,0.23255,0.31140}%
\definecolor{mycolor49}{rgb}{0.76956,0.23608,0.30749}%
\definecolor{mycolor50}{rgb}{0.77506,0.23967,0.30353}%
\definecolor{mycolor51}{rgb}{0.79129,0.25086,0.29139}%
\definecolor{mycolor52}{rgb}{0.79661,0.25473,0.28726}%
\definecolor{mycolor53}{rgb}{0.80187,0.25867,0.28310}%
\definecolor{mycolor54}{rgb}{0.81224,0.26679,0.27466}%
\definecolor{mycolor55}{rgb}{0.81734,0.27095,0.27039}%
\definecolor{mycolor56}{rgb}{0.82737,0.27952,0.26175}%
\definecolor{mycolor57}{rgb}{0.83230,0.28391,0.25738}%
\definecolor{mycolor58}{rgb}{0.83717,0.28839,0.25299}%
\definecolor{mycolor59}{rgb}{0.84671,0.29756,0.24411}%
\definecolor{mycolor60}{rgb}{0.85138,0.30226,0.23964}%
\definecolor{mycolor61}{rgb}{0.85599,0.30704,0.23513}%
\definecolor{mycolor62}{rgb}{0.86053,0.31189,0.23061}%
\definecolor{mycolor63}{rgb}{0.87374,0.32691,0.21689}%
\definecolor{mycolor64}{rgb}{0.87800,0.33206,0.21227}%
\definecolor{mycolor65}{rgb}{0.89034,0.34796,0.19829}%
\definecolor{mycolor66}{rgb}{0.89819,0.35891,0.18886}%
\definecolor{mycolor67}{rgb}{0.90200,0.36449,0.18412}%
\definecolor{mycolor68}{rgb}{0.90573,0.37014,0.17935}%
\definecolor{mycolor69}{rgb}{0.90939,0.37586,0.17456}%
\definecolor{mycolor70}{rgb}{0.91297,0.38164,0.16975}%
\definecolor{mycolor71}{rgb}{0.91646,0.38748,0.16492}%
\definecolor{mycolor72}{rgb}{0.91988,0.39339,0.16007}%
\definecolor{mycolor73}{rgb}{0.92322,0.39936,0.15519}%
\definecolor{mycolor74}{rgb}{0.92647,0.40539,0.15029}%
\definecolor{mycolor75}{rgb}{0.92964,0.41148,0.14537}%
\definecolor{mycolor76}{rgb}{0.93274,0.41763,0.14042}%
\definecolor{mycolor77}{rgb}{0.93575,0.42383,0.13544}%
\definecolor{mycolor78}{rgb}{0.93868,0.43009,0.13044}%
\definecolor{mycolor79}{rgb}{0.94152,0.43640,0.12541}%
\definecolor{mycolor80}{rgb}{0.94429,0.44277,0.12035}%
\definecolor{mycolor81}{rgb}{0.94956,0.45566,0.11016}%
\definecolor{mycolor82}{rgb}{0.95208,0.46218,0.10503}%
\definecolor{mycolor83}{rgb}{0.96129,0.48872,0.08429}%
\definecolor{mycolor84}{rgb}{0.96540,0.50225,0.07386}%
\definecolor{mycolor85}{rgb}{0.96732,0.50908,0.06866}%
\definecolor{mycolor86}{rgb}{0.97259,0.52980,0.05332}%
\definecolor{mycolor87}{rgb}{0.97568,0.54380,0.04362}%
\definecolor{mycolor88}{rgb}{0.98082,0.57221,0.02851}%
\definecolor{mycolor89}{rgb}{0.98189,0.57939,0.02625}%
\definecolor{mycolor90}{rgb}{0.98378,0.59385,0.02377}%
\definecolor{mycolor91}{rgb}{0.98532,0.60842,0.02420}%
\definecolor{mycolor92}{rgb}{0.98696,0.63048,0.03091}%
\definecolor{mycolor93}{rgb}{0.98793,0.66025,0.05175}%
\definecolor{mycolor94}{rgb}{0.98771,0.68281,0.07249}%
\definecolor{mycolor95}{rgb}{0.97750,0.78226,0.18592}%
\definecolor{mycolor96}{rgb}{0.96624,0.83619,0.26153}%
\definecolor{mycolor97}{rgb}{0.96252,0.85148,0.28555}%
\definecolor{mycolor98}{rgb}{0.98836,0.99836,0.64492}%
%
\addplot [color=mycolor1]
  table[row sep=crcr]{%
0	4395.3\\
2000	1.2822e+05\\
4000	1.5524e+05\\
6000	1.5658e+05\\
8000	1.5852e+05\\
10000	1.5536e+05\\
12000	1.5403e+05\\
14000	1.5778e+05\\
16000	1.5712e+05\\
18000	1.5878e+05\\
20000	1.5856e+05\\
22000	1.5609e+05\\
24000	1.5222e+05\\
28000	1.5978e+05\\
30000	1.5592e+05\\
32000	1.6077e+05\\
34000	1.5667e+05\\
36000	1.5523e+05\\
38000	1.5873e+05\\
40000	1.5621e+05\\
44000	1.6196e+05\\
46000	1.5907e+05\\
48000	1.5831e+05\\
50000	1.5527e+05\\
52000	1.5494e+05\\
54000	1.5676e+05\\
56000	1.5618e+05\\
58000	1.5706e+05\\
60000	1.6162e+05\\
62000	1.563e+05\\
64000	1.5664e+05\\
66000	1.5568e+05\\
70000	1.5597e+05\\
72000	1.5696e+05\\
74000	1.5184e+05\\
76000	1.5591e+05\\
78000	1.5474e+05\\
80000	1.5663e+05\\
82000	1.5671e+05\\
84000	1.5361e+05\\
86000	1.5699e+05\\
88000	1.5885e+05\\
90000	1.5769e+05\\
92000	1.5838e+05\\
94000	1.549e+05\\
96000	1.5254e+05\\
98000	1.5556e+05\\
1e+05	1.559e+05\\
};
\addlegendentry{Classe k=283}

\addplot [color=mycolor2, forget plot]
  table[row sep=crcr]{%
0	4395.3\\
2000	24024\\
4000	26831\\
6000	26796\\
8000	26159\\
10000	26541\\
12000	26562\\
14000	26850\\
16000	26014\\
18000	26185\\
20000	25803\\
22000	27149\\
24000	26907\\
26000	26052\\
28000	26782\\
30000	26299\\
32000	26199\\
34000	26373\\
36000	26211\\
38000	25854\\
40000	25820\\
42000	26744\\
44000	26090\\
46000	25734\\
48000	27018\\
50000	26284\\
52000	26255\\
54000	26916\\
56000	27028\\
58000	26105\\
60000	26297\\
62000	25914\\
64000	25653\\
66000	26777\\
68000	26126\\
70000	26181\\
72000	25855\\
76000	26356\\
80000	26290\\
82000	26544\\
84000	25952\\
86000	25692\\
92000	26431\\
94000	25962\\
96000	26471\\
98000	25682\\
1e+05	26013\\
};
\addplot [color=mycolor3, forget plot]
  table[row sep=crcr]{%
0	4395.3\\
2000	36275\\
4000	39677\\
6000	39686\\
8000	40600\\
10000	42125\\
12000	40432\\
14000	41124\\
16000	41079\\
18000	39899\\
20000	40362\\
22000	41422\\
24000	40494\\
26000	40629\\
28000	40333\\
30000	39733\\
32000	41031\\
34000	40222\\
36000	40321\\
40000	41257\\
42000	41002\\
44000	41205\\
46000	40573\\
50000	39830\\
52000	40565\\
54000	42970\\
56000	40885\\
58000	39983\\
60000	40860\\
62000	42038\\
64000	40810\\
66000	41999\\
68000	40569\\
70000	40529\\
72000	40326\\
74000	40277\\
76000	39230\\
78000	40675\\
80000	40815\\
82000	39254\\
84000	41319\\
86000	40677\\
88000	38931\\
90000	40660\\
92000	39742\\
94000	39935\\
96000	40343\\
98000	41190\\
1e+05	40387\\
};
\addplot [color=mycolor4, forget plot]
  table[row sep=crcr]{%
0	4395.3\\
2000	13701\\
4000	14015\\
6000	14672\\
8000	14420\\
10000	14607\\
12000	14272\\
14000	14194\\
16000	14617\\
18000	14274\\
20000	14223\\
24000	14758\\
26000	14470\\
28000	14478\\
30000	14342\\
32000	14471\\
34000	13871\\
36000	14330\\
38000	14149\\
40000	14366\\
42000	14267\\
44000	14597\\
46000	14126\\
48000	14553\\
50000	14474\\
52000	14807\\
54000	14109\\
56000	14708\\
58000	14082\\
60000	14338\\
62000	14332\\
64000	15264\\
66000	14305\\
68000	14546\\
70000	14341\\
72000	14411\\
74000	14434\\
76000	14330\\
78000	14270\\
80000	14332\\
82000	14518\\
84000	14941\\
86000	14489\\
88000	14252\\
90000	14347\\
92000	14243\\
94000	14526\\
96000	14582\\
98000	14346\\
1e+05	14683\\
};
\addplot [color=mycolor5, forget plot]
  table[row sep=crcr]{%
0	4395.3\\
2000	54731\\
4000	66411\\
6000	63679\\
8000	64086\\
12000	66385\\
14000	65290\\
16000	65852\\
18000	65711\\
20000	66253\\
22000	62872\\
24000	67856\\
26000	64730\\
28000	65095\\
30000	67949\\
32000	64009\\
34000	63224\\
36000	66589\\
38000	63775\\
40000	63073\\
42000	64215\\
44000	63495\\
48000	65051\\
50000	66039\\
52000	64189\\
54000	63204\\
56000	64345\\
58000	66506\\
60000	64991\\
62000	65590\\
64000	66000\\
66000	63709\\
68000	64651\\
70000	64679\\
72000	65099\\
74000	64984\\
76000	69375\\
78000	65809\\
80000	65927\\
82000	66813\\
84000	64399\\
86000	66682\\
88000	66930\\
90000	65331\\
92000	64051\\
94000	66644\\
96000	65609\\
98000	65706\\
1e+05	66850\\
};
\addplot [color=mycolor6, forget plot]
  table[row sep=crcr]{%
0	4395.3\\
2000	48458\\
4000	58168\\
6000	57657\\
8000	57373\\
10000	55800\\
12000	59132\\
14000	55942\\
16000	58906\\
18000	55862\\
20000	56756\\
22000	59595\\
24000	57524\\
26000	56350\\
28000	57906\\
30000	57215\\
34000	56913\\
36000	56246\\
38000	56706\\
40000	57629\\
42000	56777\\
44000	57294\\
46000	58279\\
48000	58357\\
50000	57133\\
52000	58970\\
56000	56392\\
58000	56428\\
60000	57234\\
62000	57329\\
64000	56704\\
66000	58331\\
68000	58687\\
70000	57855\\
74000	58700\\
76000	57717\\
78000	59221\\
80000	56467\\
82000	55056\\
84000	56172\\
86000	55631\\
88000	57222\\
90000	56566\\
92000	59349\\
94000	56518\\
96000	56529\\
98000	59008\\
1e+05	57869\\
};
\addplot [color=mycolor7, forget plot]
  table[row sep=crcr]{%
0	4395.3\\
2000	26194\\
4000	28904\\
6000	28941\\
8000	29265\\
10000	27799\\
12000	28272\\
14000	29663\\
16000	28241\\
18000	29173\\
20000	27813\\
22000	28888\\
24000	29258\\
26000	28803\\
28000	29739\\
30000	27993\\
32000	28899\\
34000	28746\\
36000	28954\\
38000	28539\\
40000	29602\\
42000	28400\\
44000	28190\\
46000	27575\\
48000	28524\\
50000	29030\\
52000	28197\\
54000	29364\\
56000	28283\\
58000	29264\\
60000	27682\\
62000	27256\\
64000	28716\\
66000	28408\\
68000	27757\\
70000	28194\\
72000	28129\\
74000	29609\\
76000	27802\\
78000	28803\\
80000	28762\\
82000	28963\\
84000	27801\\
86000	28518\\
88000	28517\\
90000	28836\\
92000	28389\\
94000	28938\\
96000	28081\\
98000	28261\\
1e+05	27615\\
};
\addplot [color=mycolor8, forget plot]
  table[row sep=crcr]{%
0	4395.3\\
2000	19247\\
4000	21398\\
6000	20622\\
8000	20731\\
10000	21329\\
12000	20959\\
14000	20293\\
16000	20305\\
18000	20613\\
20000	20619\\
22000	21186\\
24000	21029\\
26000	21015\\
28000	21180\\
30000	20703\\
32000	21162\\
34000	20243\\
36000	20653\\
40000	20735\\
42000	20743\\
44000	20602\\
46000	20642\\
48000	21180\\
50000	20030\\
52000	20830\\
54000	20864\\
56000	21120\\
58000	20529\\
60000	21041\\
62000	20917\\
64000	20541\\
68000	20495\\
70000	20863\\
72000	20870\\
76000	21201\\
78000	21053\\
80000	21177\\
82000	21616\\
84000	20414\\
86000	21375\\
88000	20536\\
90000	21019\\
92000	20469\\
94000	20524\\
96000	21282\\
98000	20828\\
1e+05	20895\\
};
\addplot [color=mycolor9, forget plot]
  table[row sep=crcr]{%
0	4395.3\\
2000	14088\\
4000	14255\\
6000	13762\\
8000	14328\\
10000	14124\\
12000	13858\\
14000	14211\\
16000	14464\\
18000	14067\\
20000	14330\\
22000	14337\\
24000	14476\\
26000	13921\\
28000	13722\\
30000	14149\\
32000	13938\\
34000	14315\\
38000	14202\\
40000	13961\\
42000	14253\\
44000	14099\\
46000	14493\\
48000	14711\\
50000	14038\\
52000	14753\\
54000	14104\\
56000	14133\\
60000	14566\\
62000	14246\\
64000	14364\\
66000	14203\\
68000	13870\\
70000	14059\\
72000	14036\\
74000	14369\\
76000	13990\\
78000	14094\\
80000	14345\\
82000	14031\\
84000	14559\\
86000	14347\\
88000	14178\\
90000	14254\\
92000	14391\\
94000	14151\\
96000	14111\\
98000	14453\\
1e+05	14103\\
};
\addplot [color=mycolor10, forget plot]
  table[row sep=crcr]{%
0	4395.3\\
2000	20239\\
4000	23077\\
6000	22334\\
8000	22026\\
10000	22319\\
12000	23219\\
14000	22394\\
16000	22461\\
18000	22328\\
20000	22462\\
22000	22692\\
24000	22082\\
26000	22597\\
28000	22512\\
30000	22307\\
32000	22674\\
34000	22633\\
38000	22341\\
40000	21788\\
42000	22590\\
44000	22637\\
46000	22524\\
48000	22627\\
52000	22439\\
54000	22659\\
56000	22375\\
58000	22736\\
60000	22123\\
62000	22593\\
64000	22229\\
66000	22683\\
68000	22310\\
70000	22364\\
72000	21786\\
74000	23081\\
76000	22619\\
78000	22368\\
80000	22442\\
82000	22344\\
84000	22826\\
86000	22130\\
88000	23328\\
90000	22672\\
92000	22646\\
94000	23015\\
96000	22618\\
98000	22727\\
1e+05	21827\\
};
\addplot [color=mycolor11, forget plot]
  table[row sep=crcr]{%
0	4395.3\\
2000	8682.5\\
6000	8452.6\\
8000	8663.2\\
10000	8505.1\\
12000	8649\\
14000	8501.3\\
16000	8645.1\\
18000	8542.5\\
20000	8412.1\\
22000	8514.8\\
24000	8500.3\\
26000	8953.6\\
28000	8588.4\\
30000	8592.2\\
32000	8549.6\\
34000	8970.7\\
36000	8759.3\\
40000	8488.7\\
42000	8499.7\\
44000	8739.8\\
46000	8631.4\\
48000	8589.5\\
50000	8429.6\\
52000	8894.7\\
54000	8570.9\\
58000	8401.9\\
60000	8721.3\\
62000	8662.2\\
64000	8792.7\\
66000	8571.3\\
68000	8536.1\\
70000	8965.6\\
72000	8573.6\\
76000	8651.1\\
78000	8924.4\\
80000	8751.9\\
82000	8359.3\\
84000	8638.6\\
86000	8671\\
88000	8604.2\\
90000	8769.6\\
92000	8656.9\\
94000	8650.5\\
96000	8825\\
98000	8523.7\\
1e+05	8506\\
};
\addplot [color=mycolor12, forget plot]
  table[row sep=crcr]{%
0	4395.3\\
2000	13982\\
4000	14905\\
6000	14708\\
8000	14752\\
12000	15285\\
14000	14744\\
16000	15118\\
18000	14440\\
20000	15211\\
22000	14844\\
24000	14667\\
26000	14739\\
28000	14768\\
30000	14639\\
32000	14632\\
34000	14816\\
36000	15199\\
38000	15117\\
40000	14654\\
42000	14482\\
44000	15219\\
46000	15156\\
48000	14468\\
50000	14727\\
52000	15090\\
54000	14468\\
56000	14626\\
58000	14467\\
60000	15062\\
62000	14384\\
64000	15332\\
66000	14192\\
68000	15074\\
70000	15215\\
72000	14504\\
76000	14833\\
80000	15263\\
82000	14662\\
84000	14762\\
86000	14521\\
88000	15207\\
90000	15310\\
92000	14549\\
94000	14944\\
98000	15319\\
1e+05	15252\\
};
\addplot [color=mycolor13, forget plot]
  table[row sep=crcr]{%
0	4395.3\\
2000	10867\\
4000	10853\\
6000	11112\\
8000	11335\\
10000	10997\\
12000	11275\\
14000	11003\\
16000	11077\\
18000	11009\\
20000	11552\\
22000	11143\\
24000	11339\\
28000	11201\\
30000	11605\\
32000	11384\\
36000	10856\\
38000	10854\\
40000	11367\\
42000	11353\\
44000	11439\\
46000	11274\\
48000	11404\\
50000	11220\\
52000	11544\\
54000	11331\\
56000	11016\\
58000	10904\\
60000	10933\\
62000	11063\\
64000	11316\\
66000	11323\\
68000	10974\\
70000	11106\\
72000	11442\\
74000	11379\\
76000	11028\\
78000	10753\\
82000	11128\\
84000	11513\\
86000	11118\\
88000	10884\\
90000	10790\\
92000	11088\\
94000	11156\\
96000	11622\\
98000	11381\\
1e+05	11189\\
};
\addplot [color=mycolor14, forget plot]
  table[row sep=crcr]{%
0	4395.3\\
2000	22236\\
4000	24222\\
6000	24794\\
8000	24807\\
10000	23553\\
12000	24113\\
14000	23776\\
18000	24611\\
20000	24631\\
22000	24182\\
24000	24236\\
26000	25146\\
28000	24555\\
30000	24328\\
32000	23910\\
34000	25329\\
36000	24643\\
38000	25464\\
40000	25028\\
42000	24270\\
44000	24815\\
46000	24171\\
48000	24422\\
50000	24456\\
52000	25534\\
54000	24109\\
56000	25037\\
58000	24644\\
60000	24624\\
62000	24933\\
66000	24584\\
68000	25065\\
70000	24221\\
72000	24648\\
74000	24677\\
76000	24583\\
78000	25278\\
80000	24438\\
84000	25189\\
86000	25115\\
88000	24293\\
90000	23851\\
92000	25190\\
94000	25587\\
96000	25164\\
98000	24986\\
1e+05	25083\\
};
\addplot [color=mycolor15, forget plot]
  table[row sep=crcr]{%
0	4395.3\\
2000	16508\\
4000	17140\\
6000	16938\\
8000	17295\\
10000	17564\\
12000	17217\\
14000	17244\\
16000	17191\\
20000	16851\\
22000	17092\\
24000	16914\\
26000	16480\\
28000	16839\\
30000	16998\\
32000	17358\\
34000	16999\\
36000	16783\\
38000	17670\\
40000	17121\\
44000	16694\\
46000	16817\\
48000	17298\\
50000	17454\\
52000	17958\\
54000	17433\\
56000	17184\\
58000	16854\\
60000	16836\\
62000	18504\\
64000	17614\\
66000	17895\\
68000	17248\\
70000	17946\\
72000	17537\\
74000	18064\\
76000	16666\\
78000	17292\\
80000	16823\\
82000	16708\\
84000	17533\\
86000	17538\\
88000	17885\\
90000	16956\\
92000	17076\\
94000	16559\\
96000	18053\\
98000	16945\\
1e+05	17659\\
};
\addplot [color=mycolor16, forget plot]
  table[row sep=crcr]{%
0	4395.3\\
2000	7359.4\\
6000	6941.4\\
8000	7139.7\\
10000	7152.6\\
12000	6997.8\\
14000	7057.8\\
16000	6968.2\\
18000	6767.3\\
20000	7209.4\\
22000	6739.6\\
24000	7161.4\\
26000	6896.2\\
28000	7030.5\\
30000	6903.8\\
32000	7037.2\\
34000	6924.3\\
36000	7000.9\\
38000	6808.2\\
40000	7069.4\\
42000	7061.2\\
44000	6775\\
46000	6941.4\\
48000	7217.9\\
50000	6940.8\\
52000	6780.3\\
54000	6785.5\\
56000	7148.4\\
58000	7097.7\\
60000	7072.8\\
62000	6975.4\\
64000	6708\\
66000	7093.4\\
68000	7016.3\\
70000	7269\\
72000	7086.1\\
74000	7212.7\\
76000	7041.5\\
78000	7265.6\\
80000	7041.3\\
82000	7019.7\\
84000	7077.2\\
88000	7292.8\\
90000	6966.9\\
92000	7215.4\\
94000	7004.5\\
96000	6885.3\\
98000	7050.7\\
1e+05	6687\\
};
\addplot [color=mycolor17, forget plot]
  table[row sep=crcr]{%
0	4395.3\\
2000	15505\\
4000	15896\\
6000	16481\\
8000	15987\\
10000	16439\\
12000	16243\\
14000	15715\\
18000	16219\\
20000	16208\\
22000	15681\\
24000	16350\\
26000	16053\\
28000	16008\\
30000	16075\\
32000	15845\\
34000	16207\\
36000	17029\\
38000	16668\\
40000	16038\\
42000	15953\\
44000	16201\\
46000	17009\\
48000	15763\\
50000	16370\\
52000	16240\\
54000	16557\\
56000	16331\\
58000	16233\\
60000	15706\\
64000	16245\\
66000	16398\\
68000	16238\\
70000	15345\\
72000	16579\\
74000	16619\\
76000	16528\\
78000	15865\\
80000	16383\\
82000	16174\\
84000	16125\\
86000	16198\\
88000	16345\\
90000	15855\\
92000	16956\\
94000	16525\\
96000	16531\\
98000	16715\\
1e+05	16307\\
};
\addplot [color=mycolor18, forget plot]
  table[row sep=crcr]{%
0	4395.3\\
2000	23041\\
4000	24396\\
8000	25472\\
10000	25812\\
12000	24840\\
14000	25954\\
16000	25481\\
18000	24463\\
20000	24818\\
22000	25432\\
24000	25295\\
26000	24180\\
28000	25290\\
30000	24967\\
32000	24384\\
34000	25484\\
36000	25587\\
38000	25172\\
40000	24394\\
42000	24730\\
44000	25259\\
46000	24825\\
48000	25602\\
50000	25540\\
52000	25253\\
54000	25376\\
56000	24422\\
58000	25041\\
60000	23331\\
62000	25122\\
64000	24905\\
66000	25476\\
68000	24987\\
70000	25532\\
72000	24936\\
74000	25264\\
76000	24871\\
80000	24949\\
82000	24609\\
84000	24662\\
86000	25590\\
88000	25227\\
90000	25762\\
92000	25664\\
94000	25484\\
96000	24400\\
1e+05	25544\\
};
\addplot [color=mycolor19, forget plot]
  table[row sep=crcr]{%
0	4395.3\\
2000	14088\\
4000	14917\\
6000	15322\\
8000	14792\\
10000	15483\\
12000	15141\\
14000	14664\\
16000	15189\\
18000	15099\\
20000	15063\\
22000	15182\\
24000	15058\\
26000	14780\\
28000	14844\\
30000	15245\\
32000	14776\\
34000	15801\\
36000	15293\\
38000	14679\\
40000	15601\\
44000	14981\\
46000	15348\\
48000	15564\\
50000	14603\\
52000	15393\\
54000	15140\\
56000	14733\\
60000	15386\\
62000	15307\\
64000	15294\\
66000	15347\\
68000	15276\\
70000	15529\\
72000	15606\\
74000	14889\\
76000	15328\\
78000	15021\\
80000	15005\\
82000	14608\\
84000	15822\\
86000	15595\\
88000	15108\\
90000	15550\\
92000	14743\\
94000	15355\\
96000	14536\\
98000	14734\\
1e+05	14314\\
};
\addplot [color=mycolor20, forget plot]
  table[row sep=crcr]{%
0	4395.3\\
2000	17463\\
4000	19496\\
6000	19834\\
8000	19782\\
10000	19643\\
12000	19371\\
14000	19554\\
16000	19563\\
18000	19756\\
20000	19181\\
22000	19205\\
24000	18832\\
26000	19418\\
28000	19020\\
30000	19337\\
32000	19263\\
34000	20405\\
36000	19989\\
38000	19518\\
40000	19235\\
42000	19121\\
44000	19620\\
48000	19265\\
50000	19964\\
52000	19564\\
54000	19852\\
56000	19048\\
58000	18684\\
60000	19235\\
62000	19914\\
64000	19323\\
66000	19683\\
68000	19160\\
72000	19246\\
74000	19628\\
76000	19649\\
78000	19151\\
80000	19382\\
82000	20412\\
84000	19983\\
86000	19096\\
88000	19038\\
90000	18755\\
92000	19278\\
94000	18746\\
96000	19862\\
98000	19871\\
1e+05	19697\\
};
\addplot [color=mycolor21, forget plot]
  table[row sep=crcr]{%
0	4395.3\\
2000	9617.5\\
4000	9918\\
6000	9747\\
8000	9622.7\\
10000	10003\\
12000	9780.7\\
14000	9669.4\\
16000	9873.3\\
18000	9797.6\\
20000	10048\\
22000	9700.3\\
24000	9682.9\\
26000	9991.8\\
30000	9637.2\\
32000	9662.3\\
34000	9714.9\\
36000	9652.3\\
38000	9799.5\\
42000	10021\\
44000	9700.6\\
46000	9720.5\\
48000	9575.5\\
50000	9672.5\\
52000	9657.5\\
54000	9391.2\\
56000	9993.1\\
58000	9685.3\\
60000	9860.7\\
62000	9655\\
64000	9673.7\\
66000	9366.2\\
68000	9891.3\\
70000	9705.5\\
72000	10033\\
74000	9561.5\\
76000	9925.7\\
78000	9619.9\\
80000	9608\\
82000	9689.2\\
84000	10023\\
86000	9856.9\\
88000	9541.4\\
90000	9727.2\\
94000	9675.9\\
96000	9949.9\\
98000	9945.4\\
1e+05	9693\\
};
\addplot [color=mycolor22, forget plot]
  table[row sep=crcr]{%
0	4395.3\\
2000	23026\\
4000	24748\\
6000	25349\\
8000	26179\\
10000	25568\\
12000	25935\\
14000	25766\\
16000	25709\\
18000	25765\\
20000	25042\\
22000	25546\\
24000	26137\\
26000	25598\\
28000	25262\\
30000	26446\\
32000	25668\\
34000	25411\\
36000	25463\\
38000	25776\\
40000	27136\\
42000	25228\\
44000	26449\\
46000	26105\\
48000	25386\\
50000	26628\\
52000	26008\\
54000	25668\\
56000	25945\\
58000	25487\\
60000	26160\\
62000	26513\\
64000	26188\\
66000	25071\\
68000	26330\\
70000	25147\\
72000	25117\\
74000	26002\\
76000	25604\\
78000	26123\\
80000	26403\\
82000	25664\\
84000	25212\\
86000	25603\\
88000	25675\\
90000	26216\\
92000	25626\\
94000	25945\\
96000	25988\\
98000	26971\\
1e+05	26475\\
};
\addplot [color=mycolor23, forget plot]
  table[row sep=crcr]{%
0	4395.3\\
2000	11956\\
4000	12409\\
6000	12569\\
8000	12547\\
10000	12488\\
12000	12359\\
14000	12660\\
16000	12524\\
18000	12520\\
20000	12719\\
22000	12172\\
24000	12308\\
26000	12324\\
28000	12502\\
30000	12602\\
32000	12212\\
34000	12455\\
36000	12448\\
38000	12346\\
40000	12359\\
42000	12476\\
44000	12492\\
46000	12167\\
48000	12507\\
50000	12257\\
52000	12286\\
56000	12618\\
58000	12454\\
60000	12550\\
62000	12509\\
64000	12363\\
66000	12479\\
68000	12232\\
70000	12594\\
72000	12059\\
74000	12666\\
76000	12495\\
78000	12655\\
80000	12559\\
82000	12594\\
84000	12370\\
86000	12407\\
88000	12327\\
90000	12313\\
92000	12561\\
94000	12223\\
96000	12902\\
98000	12425\\
1e+05	12512\\
};
\addplot [color=mycolor24, forget plot]
  table[row sep=crcr]{%
0	4395.3\\
2000	4732.1\\
4000	4636\\
6000	4908\\
8000	4680.7\\
10000	4885.2\\
12000	4790.5\\
14000	4640.5\\
16000	4677.5\\
18000	4629.8\\
20000	4705.4\\
22000	4796.4\\
24000	4646.8\\
26000	4642.8\\
28000	4851.2\\
30000	4802.8\\
32000	4640.3\\
34000	4647.3\\
38000	4713.9\\
40000	4802\\
42000	4605\\
44000	4915.5\\
46000	4594.7\\
48000	4789\\
50000	4704.7\\
52000	4725.8\\
54000	4783.7\\
56000	4752.5\\
58000	4696.1\\
60000	4788.9\\
62000	4684.2\\
64000	4648.4\\
66000	4736.6\\
68000	4803.4\\
70000	4849.6\\
74000	4667.9\\
76000	4831.8\\
78000	4799.4\\
80000	4853.5\\
82000	4723.1\\
84000	4669.8\\
86000	4856.1\\
90000	4735.9\\
92000	4860.6\\
94000	4866.3\\
96000	4702.6\\
98000	4628\\
1e+05	4694.1\\
};
\addplot [color=mycolor25, forget plot]
  table[row sep=crcr]{%
0	4395.3\\
2000	11762\\
4000	12557\\
8000	12687\\
10000	12518\\
12000	12163\\
14000	12815\\
16000	12771\\
18000	11965\\
22000	13003\\
24000	12491\\
26000	12095\\
28000	12527\\
30000	12831\\
32000	12904\\
34000	13187\\
36000	12112\\
38000	11944\\
40000	12789\\
42000	12757\\
44000	12407\\
46000	12530\\
48000	12012\\
50000	12942\\
52000	12590\\
54000	12838\\
56000	12894\\
58000	12147\\
60000	12104\\
62000	12217\\
64000	12558\\
66000	12691\\
68000	12632\\
70000	12501\\
72000	13185\\
74000	12483\\
76000	12215\\
78000	12898\\
80000	12245\\
82000	12986\\
84000	12752\\
86000	12326\\
90000	12499\\
92000	12435\\
94000	12840\\
96000	12662\\
98000	12551\\
1e+05	12827\\
};
\addplot [color=mycolor26, forget plot]
  table[row sep=crcr]{%
0	4395.3\\
2000	7965.5\\
4000	8284.2\\
6000	7977.2\\
8000	8519.3\\
10000	8279\\
12000	8615.6\\
14000	8614.8\\
16000	8266.6\\
18000	8428.9\\
20000	8194.8\\
22000	8352.8\\
24000	8378.5\\
26000	8085.2\\
28000	8270.4\\
30000	7999.4\\
32000	8574.4\\
34000	8328.5\\
36000	8230.3\\
38000	8597.6\\
40000	8191.9\\
42000	8183.4\\
44000	8251\\
46000	8264\\
50000	8210.6\\
52000	8629.5\\
56000	8348.7\\
58000	8563.9\\
60000	8012\\
62000	8063.9\\
64000	8348.1\\
66000	8342.2\\
68000	8130.3\\
70000	8443.3\\
72000	8209.9\\
74000	8462.6\\
76000	8246\\
78000	8244.3\\
80000	8491.3\\
82000	8523.9\\
84000	8190.4\\
86000	8281.3\\
88000	8216.4\\
90000	8447\\
92000	8401.3\\
94000	8391.1\\
96000	8326.9\\
98000	8435\\
1e+05	8574.2\\
};
\addplot [color=mycolor27, forget plot]
  table[row sep=crcr]{%
0	4395.3\\
2000	9480.7\\
4000	9870.9\\
6000	9780.8\\
8000	9739\\
10000	9899.6\\
12000	9817.2\\
14000	9592.9\\
16000	9844.5\\
18000	9845.4\\
20000	9936.7\\
22000	10081\\
24000	10086\\
26000	9744.1\\
28000	9880.7\\
30000	9833.5\\
32000	10034\\
34000	9742.2\\
36000	9867.6\\
38000	9907.4\\
40000	9675.9\\
42000	9747.2\\
44000	9889.6\\
46000	9796\\
48000	10081\\
50000	10026\\
52000	9151.2\\
54000	10014\\
56000	9516.4\\
58000	9872.8\\
60000	9666.8\\
62000	9429.3\\
64000	9987.2\\
66000	9456.8\\
68000	10265\\
70000	9737\\
72000	9650.2\\
74000	9332.2\\
76000	9745.7\\
78000	9882.5\\
80000	9922.3\\
82000	10187\\
84000	9930.7\\
86000	9583.1\\
88000	10061\\
90000	9892.9\\
92000	10040\\
94000	9955.6\\
96000	9900.1\\
98000	9526.2\\
1e+05	9833.9\\
};
\addplot [color=mycolor28, forget plot]
  table[row sep=crcr]{%
0	4395.3\\
2000	5530.7\\
4000	5586.4\\
6000	5722.5\\
8000	5353.2\\
10000	5435.3\\
12000	5536.5\\
14000	5347.3\\
16000	5262.3\\
18000	5329.9\\
20000	5170.8\\
22000	5409.7\\
26000	5270\\
28000	5253.4\\
30000	5469.6\\
32000	5296.1\\
34000	5217.1\\
36000	5567.8\\
38000	5715.1\\
40000	5277.1\\
42000	5511.8\\
44000	5548.6\\
46000	5507.6\\
48000	5269.2\\
50000	5426\\
52000	5484.5\\
54000	5329.3\\
56000	5694.8\\
58000	5513.6\\
60000	5354.2\\
62000	5376.6\\
64000	5443.6\\
66000	5477.2\\
68000	5238.9\\
72000	5220.6\\
74000	5299.8\\
76000	5411.9\\
78000	5348.8\\
80000	5496.9\\
82000	5421.8\\
84000	5239.8\\
86000	5546.3\\
88000	5554.9\\
90000	5386.3\\
92000	5343\\
94000	5459.8\\
96000	5427\\
98000	5412.6\\
1e+05	5565\\
};
\addplot [color=mycolor29, forget plot]
  table[row sep=crcr]{%
0	4395.3\\
2000	10959\\
4000	11400\\
6000	11502\\
8000	11641\\
10000	12202\\
12000	11412\\
14000	11448\\
20000	11732\\
22000	11756\\
24000	11289\\
26000	11981\\
28000	11832\\
32000	11416\\
34000	11779\\
36000	11781\\
38000	11318\\
40000	11914\\
44000	11310\\
46000	11827\\
48000	12109\\
52000	11627\\
54000	11271\\
56000	11831\\
58000	11603\\
60000	11990\\
62000	12214\\
64000	11694\\
66000	12090\\
68000	11787\\
70000	11108\\
72000	11562\\
74000	11922\\
76000	12062\\
78000	12154\\
80000	11661\\
82000	11957\\
84000	12148\\
86000	11655\\
88000	11682\\
90000	11964\\
92000	11742\\
94000	11835\\
96000	11522\\
98000	11742\\
1e+05	11920\\
};
\addplot [color=mycolor30, forget plot]
  table[row sep=crcr]{%
0	4395.3\\
2000	7975.7\\
4000	7894.9\\
6000	8134.5\\
8000	7660.4\\
10000	7730.9\\
12000	7852.6\\
14000	7839.2\\
16000	7977.8\\
18000	8009.4\\
20000	8070.1\\
22000	7622.8\\
24000	7889.9\\
26000	7927.3\\
28000	8046.3\\
30000	8026.4\\
32000	7803.8\\
34000	7993.7\\
36000	7756.5\\
38000	7998.2\\
40000	7945.4\\
42000	8082.6\\
44000	8059.2\\
46000	7883.8\\
48000	7880.5\\
50000	7945.2\\
52000	7795\\
54000	7733.2\\
56000	8014.8\\
58000	7784\\
60000	7918\\
62000	7959.1\\
64000	7786.6\\
66000	7862.5\\
68000	8074.4\\
70000	8099.3\\
72000	7786.3\\
74000	7730.5\\
76000	8078.9\\
78000	7937.6\\
82000	8083\\
84000	7967.8\\
86000	7943.6\\
88000	7844.9\\
90000	7959.1\\
92000	7988.2\\
94000	7896.3\\
96000	7930.4\\
98000	7849.8\\
1e+05	7878.1\\
};
\addplot [color=mycolor31, forget plot]
  table[row sep=crcr]{%
0	4395.3\\
2000	6153.8\\
4000	6760.4\\
6000	6200.7\\
8000	6451.3\\
10000	6337.9\\
12000	6295.4\\
14000	6369.3\\
16000	6395.3\\
18000	6674.6\\
20000	6695.4\\
22000	6468.3\\
26000	6601.7\\
30000	6619.6\\
32000	6532.1\\
34000	6565.8\\
36000	6575.6\\
38000	6782.6\\
40000	6558.3\\
42000	6524.7\\
44000	6462.7\\
46000	6578.7\\
48000	6405.5\\
50000	6649.2\\
52000	6175.9\\
54000	6448.2\\
56000	6464\\
58000	6607.1\\
60000	6637.2\\
64000	6497.6\\
68000	6390.4\\
70000	6671.8\\
72000	6577.1\\
74000	6581.1\\
76000	6351.6\\
78000	6751.5\\
80000	6551\\
82000	6495\\
84000	6188.5\\
86000	6337.4\\
88000	6459.8\\
90000	6473.1\\
92000	6618.1\\
94000	6487.7\\
96000	6521.6\\
98000	6414.5\\
1e+05	6528\\
};
\addplot [color=mycolor32, forget plot]
  table[row sep=crcr]{%
0	4395.3\\
2000	10202\\
4000	10560\\
6000	10401\\
8000	10537\\
10000	10778\\
12000	10120\\
14000	10738\\
16000	9959.6\\
18000	10951\\
20000	10395\\
22000	10272\\
24000	10680\\
26000	10860\\
28000	10375\\
30000	11196\\
32000	10666\\
34000	10728\\
36000	10824\\
38000	10550\\
40000	10663\\
42000	10623\\
44000	10954\\
46000	10076\\
48000	10521\\
50000	10861\\
52000	10293\\
54000	10478\\
56000	10393\\
58000	10491\\
60000	10650\\
62000	10542\\
64000	10805\\
66000	10687\\
68000	10934\\
70000	10407\\
72000	10391\\
74000	10662\\
76000	10489\\
78000	10500\\
80000	10752\\
82000	10565\\
84000	10326\\
86000	10839\\
88000	10509\\
90000	10656\\
92000	10423\\
94000	10909\\
96000	10803\\
98000	10493\\
1e+05	10358\\
};
\addplot [color=mycolor33, forget plot]
  table[row sep=crcr]{%
0	4395.3\\
2000	8048\\
4000	8161\\
6000	8245.3\\
8000	8536.2\\
10000	8394.4\\
12000	8441.7\\
14000	8186.8\\
16000	8421.5\\
18000	8282.9\\
20000	8419.1\\
22000	8260.3\\
24000	8357.7\\
26000	7993.3\\
28000	8285.7\\
30000	8323\\
32000	8086.3\\
34000	8135.9\\
36000	8014.9\\
38000	8385\\
40000	8496.2\\
42000	8582.7\\
44000	7883.8\\
46000	8387.4\\
48000	8632.4\\
50000	8555.1\\
52000	8297.6\\
54000	8089.8\\
56000	8365.2\\
58000	8107.2\\
60000	8075.8\\
62000	8224.5\\
64000	8129\\
66000	8208.2\\
68000	8148.3\\
72000	8263.3\\
74000	8240.9\\
76000	7875.8\\
78000	8323\\
80000	8197.7\\
82000	8139.5\\
86000	8255.1\\
88000	8420.3\\
90000	8323.7\\
92000	8348.4\\
94000	8116.1\\
96000	8532\\
98000	8208.8\\
1e+05	8503.9\\
};
\addplot [color=mycolor34, forget plot]
  table[row sep=crcr]{%
0	4395.3\\
2000	4035.7\\
4000	3972\\
6000	3954.1\\
8000	4017.1\\
10000	3924\\
14000	3986.6\\
16000	3961.1\\
18000	3837.9\\
20000	3956.3\\
22000	4010\\
24000	3848.2\\
26000	3965\\
28000	3840.2\\
30000	3842.2\\
32000	3959.7\\
34000	4095.5\\
36000	3962.2\\
38000	3847.2\\
40000	3987\\
42000	4038.6\\
44000	3968.1\\
46000	4003.2\\
48000	4102.5\\
50000	3964.3\\
52000	3953.7\\
54000	3994.6\\
56000	3981.9\\
60000	4070.6\\
62000	4094.1\\
64000	4006.3\\
66000	3981.5\\
68000	3977.6\\
70000	4017.4\\
72000	4010\\
74000	3963.6\\
76000	3944.5\\
78000	3878.1\\
80000	4071.5\\
82000	3864.8\\
84000	3849.4\\
86000	3854.8\\
92000	3944.4\\
96000	3967.2\\
98000	3832\\
1e+05	3831.1\\
};
\addplot [color=mycolor35]
  table[row sep=crcr]{%
0	4395.3\\
2000	9656.4\\
4000	10290\\
6000	10478\\
8000	10240\\
10000	10295\\
12000	10291\\
14000	10105\\
18000	10360\\
20000	10077\\
22000	10256\\
24000	10126\\
26000	10462\\
28000	10445\\
30000	10099\\
32000	10363\\
34000	10323\\
36000	10511\\
38000	10219\\
40000	10442\\
42000	10222\\
46000	10266\\
48000	10343\\
52000	10204\\
54000	10276\\
56000	10047\\
58000	10226\\
60000	10295\\
62000	10199\\
64000	10331\\
66000	10318\\
68000	10238\\
70000	9946.4\\
72000	10558\\
74000	10214\\
76000	10347\\
78000	10173\\
80000	10229\\
82000	10386\\
84000	10470\\
86000	10202\\
88000	10179\\
90000	10294\\
92000	10160\\
94000	10216\\
96000	10134\\
98000	10207\\
1e+05	10237\\
};
\addlegendentry{Classe k=83}

\addplot [color=mycolor36, forget plot]
  table[row sep=crcr]{%
0	4395.3\\
2000	4314.5\\
4000	4183.1\\
6000	4280.6\\
8000	4198.7\\
10000	4163\\
12000	4329.7\\
14000	4160.6\\
16000	4256.4\\
18000	4225.4\\
20000	4343.4\\
22000	4242.9\\
24000	4221.4\\
26000	4102\\
28000	4269\\
32000	4234.6\\
34000	4427.1\\
36000	4227.2\\
38000	4221.9\\
40000	4195.9\\
42000	4232.6\\
44000	4287.8\\
46000	4391.6\\
48000	4163.7\\
50000	4198\\
52000	4258.3\\
54000	4279.7\\
56000	4162.1\\
58000	4212.1\\
60000	4169.1\\
62000	4188.1\\
64000	4181.2\\
66000	4274.8\\
68000	4229.6\\
70000	4323.9\\
72000	4087.3\\
74000	4194.3\\
76000	4191.3\\
78000	4289.9\\
80000	4228.4\\
82000	4273.7\\
84000	4339.6\\
86000	4233.4\\
88000	4338.1\\
90000	4311.1\\
92000	4184.5\\
94000	4136.4\\
96000	4220\\
98000	4175.1\\
1e+05	4162.8\\
};
\addplot [color=mycolor37, forget plot]
  table[row sep=crcr]{%
0	4395.3\\
2000	3777.1\\
4000	3584\\
6000	3586.6\\
8000	3601.2\\
12000	3562.3\\
14000	3605.1\\
16000	3509.2\\
18000	3604.4\\
20000	3585.5\\
22000	3540.2\\
24000	3660.8\\
26000	3624.3\\
28000	3633.8\\
30000	3502\\
32000	3538.5\\
34000	3615.2\\
36000	3555.6\\
38000	3613.1\\
40000	3559.9\\
42000	3546.9\\
44000	3598.4\\
46000	3569.9\\
48000	3605\\
50000	3671.4\\
52000	3557.9\\
54000	3597.2\\
56000	3545.6\\
58000	3550.6\\
60000	3573.8\\
62000	3587\\
64000	3642.9\\
66000	3584.4\\
68000	3441.5\\
70000	3595.2\\
72000	3601.9\\
74000	3530.7\\
76000	3527.7\\
78000	3698.8\\
82000	3486.1\\
84000	3730\\
86000	3481.9\\
90000	3597.8\\
92000	3513.8\\
96000	3629.1\\
1e+05	3603.4\\
};
\addplot [color=mycolor38, forget plot]
  table[row sep=crcr]{%
0	4395.3\\
2000	3083.8\\
8000	2862.1\\
10000	2993.6\\
12000	2969.5\\
14000	2976.2\\
16000	2902.8\\
18000	2989\\
20000	2895.1\\
22000	2975.4\\
24000	2955.6\\
26000	2858.4\\
28000	2971.5\\
30000	2901.7\\
32000	2977.5\\
34000	2855.9\\
36000	2908.4\\
40000	2968.5\\
42000	2970.2\\
44000	2912.1\\
46000	3002.3\\
48000	2831.3\\
52000	2904.5\\
54000	2977.8\\
56000	2918.1\\
58000	2922.1\\
60000	2981.5\\
62000	2856.2\\
64000	2916.8\\
66000	2923.3\\
68000	2921.2\\
72000	2983.3\\
74000	2912.1\\
76000	3001.3\\
78000	2983.7\\
80000	2922.2\\
82000	2979.8\\
84000	2965.6\\
88000	2894.1\\
90000	2983.3\\
92000	2946.4\\
94000	2998.9\\
96000	2933.2\\
98000	3043\\
1e+05	2827.5\\
};
\addplot [color=mycolor39, forget plot]
  table[row sep=crcr]{%
0	4395.3\\
2000	11344\\
4000	11832\\
6000	11781\\
8000	11853\\
10000	11817\\
12000	11949\\
14000	11638\\
16000	11615\\
18000	11993\\
20000	11740\\
22000	11884\\
24000	11973\\
26000	11988\\
28000	12047\\
32000	12336\\
34000	11650\\
36000	12053\\
38000	11964\\
40000	11801\\
42000	12281\\
44000	11708\\
46000	11716\\
48000	11195\\
50000	11853\\
52000	11742\\
54000	11749\\
56000	11709\\
58000	11872\\
60000	11979\\
62000	12199\\
64000	11748\\
66000	11659\\
68000	11723\\
70000	11992\\
72000	11932\\
74000	11567\\
76000	12374\\
78000	12075\\
80000	12166\\
82000	11973\\
84000	11867\\
86000	11820\\
88000	12117\\
90000	12325\\
92000	11808\\
94000	12022\\
96000	11884\\
98000	11948\\
1e+05	11880\\
};
\addplot [color=mycolor40, forget plot]
  table[row sep=crcr]{%
0	4395.3\\
2000	5844.8\\
4000	5831.5\\
8000	5707\\
10000	5688.3\\
12000	5862\\
14000	5666.4\\
16000	5721.6\\
20000	5709.1\\
22000	5643.9\\
24000	5717.5\\
26000	5680.7\\
28000	5774.1\\
30000	5606.3\\
32000	5850.2\\
34000	5775.9\\
36000	5721.8\\
38000	5645.4\\
42000	5727.8\\
44000	5727.1\\
46000	5784.9\\
48000	5674.4\\
50000	5757\\
52000	5820.1\\
54000	5842\\
56000	5677.3\\
58000	5789.5\\
60000	5808\\
62000	5730.9\\
64000	5764.3\\
66000	5820.6\\
68000	5949.7\\
70000	5737.7\\
72000	5850.6\\
74000	5843\\
76000	5708.6\\
78000	5643.4\\
80000	5708.1\\
82000	5592.5\\
84000	5786.9\\
86000	5694.5\\
88000	5675\\
90000	5752.7\\
92000	5727\\
96000	5724.5\\
98000	5587.1\\
1e+05	5727.7\\
};
\addplot [color=mycolor41, forget plot]
  table[row sep=crcr]{%
0	4395.3\\
2000	5259\\
4000	5373.7\\
6000	5660\\
8000	5492.4\\
10000	5365.4\\
12000	5322.5\\
14000	5565.4\\
16000	5598.1\\
18000	5565.4\\
20000	5409.7\\
22000	5483.3\\
26000	5177.5\\
28000	5333.1\\
30000	5291.6\\
32000	5273.5\\
34000	5428.8\\
38000	5273\\
40000	5560.1\\
42000	5301.7\\
44000	5328.2\\
46000	5406.2\\
48000	5376.9\\
50000	5263.6\\
52000	5365.6\\
54000	5633.1\\
56000	5327.8\\
58000	5473.6\\
60000	5375.2\\
62000	5357.4\\
64000	5373.9\\
66000	5477.7\\
68000	5418.5\\
70000	5414.1\\
72000	5548.5\\
74000	5217.2\\
76000	5471.3\\
78000	5172.5\\
80000	5397.1\\
82000	5423.8\\
84000	5272.9\\
86000	5384.1\\
88000	5513.3\\
90000	5363.3\\
92000	5273.4\\
96000	5541.5\\
98000	5537.2\\
1e+05	5448.1\\
};
\addplot [color=mycolor42, forget plot]
  table[row sep=crcr]{%
0	4395.3\\
2000	19689\\
4000	21281\\
6000	21641\\
8000	21086\\
10000	20460\\
12000	20826\\
14000	21543\\
16000	20964\\
18000	21204\\
20000	21540\\
22000	20922\\
24000	21117\\
26000	20866\\
28000	20501\\
30000	20618\\
32000	20000\\
34000	21208\\
36000	21119\\
38000	20857\\
40000	21066\\
42000	20671\\
44000	20665\\
46000	21403\\
48000	20061\\
50000	21471\\
52000	21006\\
54000	20912\\
56000	20902\\
58000	20718\\
60000	21149\\
62000	20311\\
64000	20424\\
66000	21683\\
68000	21037\\
70000	21101\\
72000	21439\\
74000	21509\\
76000	21084\\
78000	21425\\
80000	21276\\
82000	20947\\
84000	21738\\
86000	21946\\
88000	21091\\
90000	20566\\
92000	20472\\
94000	21529\\
96000	20395\\
98000	19547\\
1e+05	20704\\
};
\addplot [color=mycolor43, forget plot]
  table[row sep=crcr]{%
0	4395.3\\
2000	4925.8\\
4000	4920\\
6000	4843.8\\
8000	5015.4\\
10000	4932.3\\
14000	4816.8\\
16000	4922.7\\
18000	4982.9\\
20000	4903.9\\
22000	4898.6\\
24000	4760.4\\
26000	4984.6\\
28000	4885.7\\
30000	4980.2\\
32000	4865.9\\
36000	4885.3\\
38000	4945.7\\
40000	4925.2\\
42000	4889.2\\
44000	4895\\
46000	4984.5\\
48000	5013.7\\
50000	4875.4\\
52000	5074.4\\
54000	4886.9\\
56000	4946.5\\
60000	4929.6\\
64000	4835.3\\
66000	4838.6\\
68000	4921.5\\
70000	4842.3\\
72000	4803.2\\
74000	4919.4\\
76000	4961.1\\
80000	4804.7\\
82000	4924.7\\
84000	5063.4\\
86000	4888.4\\
88000	4944.6\\
90000	4857.1\\
92000	4908.1\\
94000	4933.2\\
96000	4812.5\\
98000	4820\\
1e+05	4952.2\\
};
\addplot [color=mycolor44, forget plot]
  table[row sep=crcr]{%
0	4395.3\\
2000	3467.2\\
4000	3328\\
6000	3169.5\\
8000	3372.5\\
10000	3244.2\\
12000	3247.9\\
14000	3317.7\\
16000	3302.5\\
18000	3205.7\\
20000	3208.3\\
22000	3400.4\\
24000	3224.3\\
26000	3359.3\\
28000	3431.1\\
30000	3272.7\\
32000	3292.4\\
34000	3459.8\\
36000	3359.7\\
38000	3194\\
40000	3253.1\\
42000	3331.2\\
44000	3372.4\\
46000	3263.4\\
48000	3245.3\\
50000	3209\\
52000	3269.8\\
54000	3179.9\\
56000	3330.3\\
58000	3426.2\\
60000	3397.1\\
62000	3423.5\\
64000	3307.5\\
66000	3365.1\\
68000	3252.4\\
70000	3242.7\\
72000	3369.1\\
74000	3266.7\\
76000	3183.2\\
78000	3288.9\\
80000	3308.2\\
82000	3380.5\\
84000	3239.7\\
86000	3087.1\\
88000	3176.5\\
90000	3217.8\\
92000	3451.8\\
94000	3284.4\\
96000	3145\\
98000	3322.8\\
1e+05	3322.7\\
};
\addplot [color=mycolor45, forget plot]
  table[row sep=crcr]{%
0	4395.3\\
2000	6747.7\\
4000	7031.3\\
6000	7012.2\\
8000	6898.2\\
10000	7026.4\\
12000	7046.9\\
14000	6720.3\\
16000	6864.9\\
18000	7159.7\\
20000	6952.5\\
22000	6920.4\\
24000	7158.1\\
26000	7086.5\\
28000	6972.6\\
30000	7082.8\\
32000	7079.2\\
34000	6879.8\\
36000	7127.3\\
38000	6841\\
40000	7057.2\\
42000	6715.3\\
44000	6874.6\\
46000	6853.3\\
48000	6948.3\\
50000	6944.6\\
52000	6779.9\\
54000	6788.2\\
56000	7058.3\\
58000	7107.1\\
60000	6831.4\\
62000	6890.5\\
66000	6781.2\\
68000	7144\\
70000	6968.4\\
72000	6999.8\\
74000	6939.3\\
76000	6941.5\\
78000	6848.3\\
80000	6872.4\\
82000	7039.8\\
84000	6667.8\\
86000	6834.5\\
88000	6851\\
90000	6795.9\\
92000	7043\\
94000	6926.9\\
98000	6970\\
1e+05	6835.7\\
};
\addplot [color=mycolor46, forget plot]
  table[row sep=crcr]{%
0	4395.3\\
2000	4020.5\\
10000	4054.8\\
12000	4004\\
14000	3972.4\\
16000	4124.8\\
18000	4046\\
20000	4184.5\\
22000	4023.9\\
24000	3996.3\\
26000	4087.1\\
28000	4092.1\\
30000	3876.2\\
32000	4039.3\\
34000	3979.9\\
36000	3972.7\\
38000	4029.1\\
40000	4018.1\\
42000	4109.6\\
44000	4055.1\\
46000	3869.8\\
50000	3989.8\\
52000	4090.1\\
54000	4037.9\\
56000	3961.4\\
58000	3930.2\\
60000	3878.7\\
62000	3877.8\\
64000	3905.3\\
66000	4165.2\\
68000	3939.4\\
70000	4097.5\\
72000	3962\\
74000	4047.7\\
76000	3970.3\\
78000	4078.9\\
80000	4001.2\\
82000	4212.4\\
84000	4028\\
86000	3921.7\\
88000	4065.8\\
90000	4075\\
92000	3978.1\\
94000	4163.6\\
96000	3973.5\\
98000	3855.7\\
1e+05	4227.5\\
};
\addplot [color=mycolor47, forget plot]
  table[row sep=crcr]{%
0	4395.3\\
2000	5255.8\\
4000	5269\\
6000	5209.6\\
8000	5345\\
10000	5149\\
12000	5265.9\\
14000	5218.7\\
16000	5326.7\\
18000	5250.8\\
20000	5332.1\\
22000	5438.8\\
24000	5400.5\\
26000	5165.4\\
28000	5164.2\\
30000	5335.6\\
32000	5086.2\\
36000	5352.9\\
38000	5156.5\\
40000	5232.9\\
42000	5265.6\\
44000	5149.2\\
46000	5178.8\\
48000	5150.2\\
50000	5359.9\\
52000	5058.1\\
54000	5341.5\\
58000	5276.2\\
60000	5142.5\\
62000	5302.8\\
64000	5250.7\\
66000	5328.6\\
68000	5346.9\\
70000	5386.7\\
72000	5350.6\\
74000	5256.7\\
76000	5424.6\\
78000	5098.3\\
80000	5298.9\\
82000	5231.4\\
84000	5210.5\\
86000	5260.5\\
88000	5405.9\\
90000	5273.4\\
92000	5186.6\\
94000	5390.1\\
96000	5321.5\\
98000	5283.7\\
1e+05	5381.6\\
};
\addplot [color=mycolor48, forget plot]
  table[row sep=crcr]{%
0	4395.3\\
2000	3057\\
4000	2902.9\\
6000	2842.4\\
8000	2893.4\\
10000	2838.7\\
14000	2860.3\\
16000	2818.8\\
18000	2840.6\\
20000	2901\\
22000	2805.8\\
24000	2822.9\\
26000	2864.9\\
28000	2803\\
30000	2823.9\\
32000	2878.5\\
34000	2851.1\\
36000	2778.1\\
38000	2880.4\\
40000	2827.9\\
42000	2837.9\\
44000	2866.1\\
46000	2823.5\\
48000	2822.3\\
52000	2903.6\\
56000	2784.3\\
58000	2807.5\\
60000	2867\\
62000	2868.6\\
64000	2827.1\\
66000	2802.6\\
68000	2826.4\\
70000	2835.6\\
72000	2804\\
74000	2837.8\\
76000	2800.5\\
80000	2907.7\\
82000	2884.2\\
84000	2842.2\\
86000	2884.5\\
88000	2862.8\\
90000	2864.7\\
92000	2891.3\\
94000	2891.7\\
96000	2905.5\\
98000	2795.1\\
1e+05	2846.7\\
};
\addplot [color=mycolor49, forget plot]
  table[row sep=crcr]{%
0	4395.3\\
2000	4424.2\\
4000	4164.9\\
6000	4393.1\\
8000	4306.1\\
10000	4197.9\\
12000	4526\\
14000	4312.8\\
16000	4066.8\\
18000	4231.2\\
20000	4482.7\\
22000	4335.4\\
24000	4383.8\\
26000	4273.2\\
28000	4287.9\\
30000	4487\\
32000	4545.6\\
34000	4125.1\\
36000	4351.7\\
38000	4315.6\\
40000	4332.5\\
42000	4284.4\\
46000	4520.5\\
48000	4240.8\\
50000	4250\\
52000	4336.1\\
54000	4395\\
56000	4317.7\\
58000	4194.1\\
60000	4391.6\\
62000	4342.2\\
64000	4149.2\\
68000	4254.7\\
70000	4211.6\\
72000	4404.8\\
74000	4260.3\\
76000	4319.7\\
80000	4175\\
82000	4398.8\\
84000	4130.2\\
86000	4171.5\\
88000	4038.9\\
90000	4364.2\\
92000	4211.3\\
94000	4235.4\\
96000	4420.4\\
98000	4256.5\\
1e+05	4512.1\\
};
\addplot [color=mycolor50, forget plot]
  table[row sep=crcr]{%
0	4395.3\\
2000	2554\\
4000	2430.8\\
6000	2358.6\\
8000	2343\\
12000	2385.4\\
14000	2417.5\\
16000	2402.5\\
18000	2358.5\\
24000	2438.8\\
26000	2354.7\\
28000	2432.8\\
30000	2371.2\\
32000	2415.5\\
34000	2364.5\\
36000	2434.9\\
38000	2474.4\\
40000	2462.5\\
42000	2357.9\\
44000	2443.1\\
46000	2455\\
48000	2400.3\\
50000	2420.5\\
52000	2367\\
54000	2435.4\\
56000	2454.6\\
58000	2372.8\\
60000	2447.8\\
62000	2326.8\\
64000	2422.8\\
66000	2471.7\\
68000	2438.5\\
70000	2455.3\\
72000	2385.3\\
74000	2452.9\\
76000	2441.5\\
80000	2382.4\\
82000	2314.7\\
84000	2490.5\\
86000	2511.1\\
88000	2319.6\\
90000	2415\\
92000	2424.8\\
94000	2422\\
96000	2437.4\\
98000	2397.6\\
1e+05	2368.3\\
};
\addplot [color=mycolor51, forget plot]
  table[row sep=crcr]{%
0	4395.3\\
2000	5635.1\\
4000	5909.6\\
6000	6239\\
8000	5717.9\\
10000	5923.2\\
12000	5838.9\\
14000	5948.1\\
16000	6092.7\\
18000	6101.2\\
20000	6427.4\\
22000	6326.8\\
24000	6045.7\\
26000	5942.2\\
28000	5950.9\\
30000	5921.1\\
32000	5800.5\\
34000	6116.9\\
36000	6003.3\\
38000	6188.6\\
40000	5948.3\\
42000	5600.6\\
44000	5838\\
46000	5991.5\\
48000	5919.1\\
50000	6026.3\\
52000	6090.8\\
54000	5718.8\\
56000	6000.7\\
58000	5907.8\\
60000	5675.7\\
62000	5858.8\\
64000	5826.4\\
66000	6001\\
68000	5825.8\\
70000	5771.7\\
72000	6081.6\\
74000	5796.8\\
76000	5644.5\\
78000	5798.4\\
80000	5611.5\\
82000	5805.4\\
84000	5841.6\\
86000	5735.2\\
88000	5771.9\\
90000	5617.3\\
92000	5817.5\\
94000	5914.6\\
96000	5780.8\\
98000	6113.9\\
1e+05	6168.5\\
};
\addplot [color=mycolor52, forget plot]
  table[row sep=crcr]{%
0	4395.3\\
2000	3762.5\\
4000	3494.3\\
6000	3289.7\\
8000	3332\\
10000	3503.8\\
12000	3537\\
14000	3471.2\\
18000	3316.7\\
20000	3480.4\\
22000	3333\\
24000	3626\\
26000	3427.2\\
28000	3611.5\\
30000	3434.3\\
32000	3421\\
34000	3442\\
36000	3358.6\\
38000	3462.1\\
40000	3541.5\\
42000	3427.9\\
44000	3512.7\\
46000	3181.4\\
48000	3495.3\\
50000	3372.7\\
52000	3518.2\\
54000	3464\\
56000	3451.5\\
58000	3365.8\\
60000	3521.3\\
62000	3525.3\\
64000	3462.4\\
66000	3373\\
68000	3394.1\\
70000	3320.6\\
72000	3336.8\\
74000	3297\\
76000	3452\\
78000	3399.4\\
80000	3564.3\\
82000	3501.3\\
84000	3475.7\\
88000	3364\\
90000	3328.2\\
92000	3374.1\\
94000	3547.8\\
96000	3586.5\\
98000	3472.8\\
1e+05	3545.9\\
};
\addplot [color=mycolor53, forget plot]
  table[row sep=crcr]{%
0	4395.3\\
2000	5560.4\\
4000	5591.6\\
6000	5858.2\\
8000	5582.2\\
10000	5600.2\\
12000	5583.5\\
14000	5676.3\\
16000	5742.5\\
18000	5725.4\\
20000	5470.7\\
22000	5631.3\\
24000	5713\\
26000	5487.2\\
28000	5753.8\\
30000	5753.2\\
32000	5567.2\\
34000	5625.3\\
36000	5518.1\\
38000	5553.7\\
40000	5324.2\\
42000	5556.7\\
44000	5459.1\\
46000	5684.1\\
48000	5478.6\\
50000	5490.1\\
52000	5622.2\\
54000	5560.7\\
56000	5602.4\\
58000	5770.2\\
60000	5543.9\\
62000	5674.6\\
64000	5477.8\\
66000	5377.1\\
68000	5662.2\\
70000	5436.4\\
72000	5755.9\\
74000	5448.6\\
76000	5706.4\\
78000	5643.8\\
80000	5811.5\\
82000	5767.5\\
84000	5951.1\\
86000	5540.9\\
88000	5834.3\\
90000	5647.7\\
92000	5497.6\\
94000	5618.1\\
96000	5664.3\\
98000	5820.9\\
1e+05	5625.7\\
};
\addplot [color=mycolor54, forget plot]
  table[row sep=crcr]{%
0	4395.3\\
2000	4588.7\\
4000	4304.9\\
6000	4445\\
8000	4691.7\\
10000	4462.8\\
12000	4524.5\\
14000	4354.9\\
16000	4382.9\\
18000	4579.1\\
20000	4648.3\\
22000	4465.5\\
24000	4488.7\\
26000	4640.5\\
28000	4662.8\\
30000	4529.5\\
32000	4516.6\\
34000	4645.8\\
36000	4526.8\\
38000	4702.6\\
40000	4646\\
42000	4576.4\\
44000	4436.4\\
46000	4590.1\\
48000	4459\\
50000	4632.4\\
52000	4456.9\\
54000	4522.9\\
58000	4519.7\\
60000	4582.6\\
62000	4434.2\\
64000	4414\\
66000	4295.8\\
68000	4505.5\\
70000	4516.9\\
72000	4640.7\\
74000	4627.1\\
76000	4600\\
78000	4366.2\\
80000	4530.4\\
82000	4632.2\\
84000	4443.9\\
86000	4569.3\\
88000	4571.8\\
90000	4555.1\\
94000	4542.9\\
96000	4609.5\\
98000	4572.7\\
1e+05	4579.7\\
};
\addplot [color=mycolor55, forget plot]
  table[row sep=crcr]{%
0	4395.3\\
2000	5661.3\\
4000	5664\\
6000	5852.5\\
8000	5738.3\\
10000	5831.5\\
12000	5836.6\\
14000	5902.2\\
16000	5653.7\\
18000	5748\\
22000	5774.7\\
24000	5837\\
26000	5713.8\\
28000	5652.6\\
30000	5773.5\\
32000	5731.5\\
34000	5798.8\\
36000	5729.2\\
38000	5830.3\\
40000	5790.2\\
42000	5831.9\\
44000	5808.1\\
46000	5699.8\\
48000	5803\\
50000	5713.9\\
52000	5704.8\\
54000	5818.6\\
56000	5731.6\\
58000	5819\\
60000	5681.2\\
62000	5800.2\\
64000	5958.3\\
66000	5910.7\\
68000	5707.8\\
70000	5850.1\\
72000	5649.5\\
74000	5897.5\\
76000	5686.6\\
78000	5749.5\\
80000	5729.8\\
82000	5840.7\\
84000	5906.9\\
86000	5785.8\\
90000	5814.4\\
92000	5889\\
94000	5772.4\\
96000	5812.8\\
98000	5691.1\\
1e+05	5782.6\\
};
\addplot [color=mycolor56, forget plot]
  table[row sep=crcr]{%
0	4395.3\\
2000	3941.6\\
4000	3843.3\\
6000	4018.3\\
8000	3900\\
10000	3944\\
12000	4014.4\\
14000	3861.6\\
16000	4042.7\\
18000	3881.1\\
20000	4016.5\\
22000	3890.5\\
24000	3954.3\\
26000	4005.8\\
28000	3824.9\\
30000	3842\\
32000	3960.7\\
34000	3852.2\\
36000	3885.2\\
38000	3979.7\\
40000	3951.7\\
42000	3895.1\\
44000	3892.7\\
46000	3917.1\\
48000	3902.7\\
50000	3945.1\\
54000	3926.7\\
56000	3940.4\\
58000	4036.9\\
60000	3895.5\\
62000	3891.9\\
64000	3955.4\\
66000	4008.6\\
68000	3936\\
70000	3983.6\\
72000	3980.7\\
74000	3889.4\\
76000	3894.2\\
78000	3951.5\\
80000	3922\\
82000	3914.6\\
84000	3932.9\\
86000	3918.6\\
88000	3829\\
90000	3898\\
92000	3855.2\\
96000	4114.3\\
98000	4010.4\\
1e+05	3941.5\\
};
\addplot [color=mycolor57, forget plot]
  table[row sep=crcr]{%
0	4395.3\\
2000	6508.3\\
4000	6713.7\\
6000	6743.9\\
8000	6541.1\\
10000	6714.2\\
12000	6762.7\\
14000	6753.3\\
16000	6621.8\\
18000	6601.4\\
20000	6548.3\\
22000	6585.7\\
24000	6668.9\\
26000	6842.2\\
28000	6585.7\\
30000	6734.7\\
32000	6788.8\\
34000	6817.4\\
36000	6883.2\\
38000	6641.1\\
40000	6924.6\\
44000	6522.6\\
46000	6703.8\\
48000	6616.8\\
50000	6918\\
52000	6773.9\\
54000	6572.7\\
56000	6809\\
58000	6827.1\\
60000	6657.4\\
62000	6698.9\\
64000	6702.7\\
66000	6761.5\\
68000	6774.4\\
72000	6690.5\\
74000	6693.9\\
76000	6417.7\\
78000	6634.6\\
80000	6740.5\\
82000	6804\\
84000	6562.6\\
86000	6667.6\\
90000	6750.7\\
92000	6620.5\\
94000	6740.4\\
96000	6633.4\\
98000	6744.7\\
1e+05	6734.2\\
};
\addplot [color=mycolor58, forget plot]
  table[row sep=crcr]{%
0	4395.3\\
2000	3681\\
4000	3688\\
6000	3463.4\\
8000	3436.1\\
12000	3524.1\\
14000	3500.9\\
16000	3453.5\\
18000	3513.8\\
22000	3499.9\\
24000	3398\\
26000	3503.3\\
28000	3499\\
30000	3529.8\\
32000	3463.1\\
34000	3566.4\\
36000	3515.3\\
38000	3423.9\\
40000	3531.1\\
42000	3490.1\\
44000	3558.6\\
46000	3527\\
48000	3565.6\\
50000	3554.7\\
52000	3603.2\\
54000	3440.8\\
58000	3611.2\\
60000	3589.1\\
62000	3493.3\\
64000	3615\\
66000	3562.8\\
68000	3560.2\\
72000	3405.6\\
74000	3604.4\\
76000	3324.3\\
78000	3445.5\\
80000	3512.9\\
84000	3493\\
86000	3456.2\\
88000	3611.8\\
90000	3527.6\\
92000	3490.9\\
94000	3407.5\\
96000	3541.9\\
98000	3508.8\\
1e+05	3424.1\\
};
\addplot [color=mycolor59, forget plot]
  table[row sep=crcr]{%
0	4395.3\\
2000	4947.7\\
4000	5094.6\\
6000	4913.9\\
8000	5081.5\\
10000	5063.9\\
12000	4982.6\\
14000	4976.2\\
16000	5031.7\\
18000	4975.4\\
20000	5048\\
22000	5044\\
24000	5126.8\\
26000	5026.2\\
28000	4992.1\\
30000	5074.6\\
32000	5062.1\\
34000	5069.3\\
36000	5024.7\\
38000	5091.9\\
40000	5116.4\\
42000	4950\\
44000	4944.9\\
46000	5154\\
48000	5054.2\\
50000	5120.3\\
52000	5165.5\\
54000	5052.7\\
56000	5063.3\\
58000	5172.9\\
60000	5000.6\\
62000	5137.8\\
64000	5145.9\\
66000	5010.6\\
68000	4781.8\\
70000	4993.6\\
72000	4974.1\\
74000	5049.1\\
76000	5020.6\\
78000	5119.3\\
80000	4940\\
82000	5118.1\\
84000	5070.6\\
86000	4996\\
88000	5032.9\\
90000	5039.8\\
92000	4921\\
96000	5034.7\\
98000	4887.1\\
1e+05	4995.5\\
};
\addplot [color=mycolor60, forget plot]
  table[row sep=crcr]{%
0	4395.3\\
2000	3644.3\\
4000	3409.4\\
6000	3585.3\\
8000	3464\\
10000	3637.6\\
12000	3528.3\\
14000	3438\\
16000	3528\\
18000	3456\\
20000	3516.5\\
22000	3477.4\\
24000	3574.3\\
26000	3612.6\\
28000	3418.7\\
30000	3502.3\\
32000	3425.4\\
34000	3532.9\\
38000	3534.6\\
40000	3348.1\\
42000	3423.5\\
44000	3454.1\\
46000	3538.2\\
48000	3398.6\\
50000	3359.9\\
52000	3411.5\\
54000	3566.9\\
58000	3489\\
60000	3466.7\\
62000	3526.1\\
64000	3519\\
66000	3604\\
68000	3562.1\\
70000	3534.8\\
72000	3477.5\\
74000	3439.2\\
76000	3483.2\\
80000	3480.5\\
82000	3455.3\\
84000	3487.7\\
86000	3460.8\\
88000	3558.4\\
90000	3524.4\\
92000	3501\\
94000	3388.2\\
96000	3444.6\\
98000	3554.5\\
1e+05	3450\\
};
\addplot [color=mycolor61, forget plot]
  table[row sep=crcr]{%
0	4395.3\\
2000	3260.1\\
4000	3059.5\\
6000	3070.7\\
8000	3090.6\\
10000	3094.4\\
12000	3078.3\\
14000	3114.8\\
16000	3101.5\\
18000	3039.7\\
20000	3048.2\\
22000	3044.3\\
24000	3062.8\\
26000	3108\\
28000	3013.4\\
32000	3103.7\\
34000	3072\\
38000	3039.1\\
40000	3049.8\\
42000	3107.2\\
44000	3087.4\\
46000	3037.6\\
48000	3064.1\\
50000	3065.6\\
52000	3044.7\\
54000	3079.3\\
56000	3050.5\\
58000	3091.1\\
60000	3056\\
62000	3062.2\\
64000	3039.3\\
68000	3089\\
70000	3010.9\\
72000	3045.9\\
74000	3046.7\\
76000	3006.6\\
78000	3056.7\\
80000	3010.9\\
82000	3104.3\\
84000	3033.4\\
86000	3027.9\\
88000	2984.6\\
90000	3052.1\\
92000	3085.8\\
94000	3001.3\\
96000	3076.2\\
98000	3093.7\\
1e+05	3018.3\\
};
\addplot [color=mycolor62, forget plot]
  table[row sep=crcr]{%
0	4395.3\\
2000	4810.9\\
4000	4986.9\\
6000	4905.4\\
8000	5047.9\\
12000	4968.9\\
14000	4932.3\\
16000	4876.6\\
18000	4989.9\\
20000	4961.3\\
22000	4858.5\\
24000	4912.9\\
26000	5028\\
28000	4882.8\\
30000	5031.5\\
32000	4914.9\\
34000	4869.5\\
36000	4970.5\\
38000	4867.2\\
40000	4938.9\\
42000	4882.2\\
44000	4884.2\\
46000	5000.1\\
48000	5020.1\\
50000	4835.2\\
52000	4974.2\\
54000	5035.2\\
56000	4969.7\\
58000	5039.9\\
60000	5007.5\\
62000	4868.1\\
64000	4799.3\\
66000	4906.6\\
68000	4898.6\\
70000	4870.6\\
72000	4959.8\\
74000	4866.7\\
76000	4994.4\\
78000	4776.8\\
80000	4886.9\\
82000	4894.9\\
84000	5012.3\\
86000	5016.7\\
88000	4998.7\\
90000	5016.1\\
92000	4884.4\\
94000	5011.9\\
96000	4990\\
1e+05	4891.1\\
};
\addplot [color=mycolor63, forget plot]
  table[row sep=crcr]{%
0	4395.3\\
2000	2553.3\\
4000	2374.7\\
6000	2301.2\\
8000	2345\\
10000	2341.3\\
12000	2299.1\\
14000	2335.2\\
16000	2291\\
18000	2378.2\\
20000	2261.3\\
22000	2415\\
24000	2325.6\\
28000	2321.5\\
30000	2330.4\\
32000	2368.2\\
34000	2313.5\\
36000	2328.3\\
40000	2397.5\\
42000	2344.8\\
44000	2300.2\\
46000	2379.8\\
48000	2327.6\\
50000	2313.7\\
52000	2308.2\\
54000	2409.7\\
58000	2337.5\\
60000	2312\\
62000	2371.5\\
64000	2351\\
66000	2346.9\\
68000	2356.7\\
74000	2365.2\\
76000	2351.1\\
78000	2408.5\\
80000	2333.1\\
82000	2303.7\\
84000	2336.9\\
86000	2315.1\\
88000	2332.9\\
90000	2326.3\\
92000	2349.4\\
94000	2328\\
96000	2406.5\\
98000	2332.4\\
1e+05	2377.6\\
};
\addplot [color=mycolor64, forget plot]
  table[row sep=crcr]{%
0	4395.3\\
2000	3125.7\\
4000	2903.8\\
6000	3002.6\\
8000	3040.3\\
10000	3037.8\\
12000	2943.2\\
14000	2973.1\\
16000	2955.7\\
18000	3042\\
20000	2972.1\\
22000	2982.5\\
24000	3010.7\\
26000	2994.1\\
28000	2994.9\\
32000	3035.1\\
34000	2959.8\\
36000	3046.2\\
38000	3020.7\\
40000	3027.8\\
44000	3005.1\\
46000	3031.1\\
48000	2969.9\\
50000	3022.8\\
52000	3048.6\\
54000	2997.2\\
56000	3032.5\\
58000	3003.9\\
60000	3036.8\\
62000	3033.4\\
64000	2988.3\\
66000	2973.3\\
68000	2969.6\\
70000	2990\\
72000	2994.4\\
74000	3028.5\\
76000	3027.1\\
78000	2983.3\\
80000	2907.8\\
84000	3003.4\\
86000	2951.6\\
88000	2939\\
90000	2994.4\\
92000	2975.6\\
94000	2992.6\\
96000	3041.7\\
98000	2952.5\\
1e+05	2954.7\\
};
\addplot [color=mycolor65, forget plot]
  table[row sep=crcr]{%
0	4395.3\\
2000	2741.9\\
4000	2504.6\\
8000	2454.8\\
10000	2478\\
12000	2511.8\\
14000	2478.9\\
16000	2487.6\\
18000	2432.2\\
20000	2421.7\\
22000	2533.5\\
24000	2463.1\\
26000	2416.3\\
28000	2460.3\\
30000	2522.5\\
32000	2429.7\\
34000	2526.7\\
36000	2474.1\\
38000	2529.8\\
40000	2489.3\\
42000	2464.6\\
44000	2506.7\\
46000	2485.3\\
48000	2441.2\\
50000	2437.4\\
52000	2440.9\\
56000	2534.9\\
58000	2474.9\\
60000	2470\\
64000	2481.8\\
66000	2507.6\\
68000	2510.4\\
70000	2552\\
72000	2615.4\\
74000	2480.1\\
76000	2439.2\\
78000	2449.6\\
80000	2492\\
82000	2506.5\\
84000	2479.8\\
86000	2438.8\\
88000	2483.4\\
90000	2485.1\\
92000	2435.6\\
94000	2488.4\\
96000	2461\\
98000	2517.6\\
1e+05	2503.5\\
};
\addplot [color=mycolor66, forget plot]
  table[row sep=crcr]{%
0	4395.3\\
2000	3466.4\\
4000	3410.9\\
6000	3335.4\\
10000	3329.3\\
12000	3376.6\\
14000	3452.7\\
16000	3432.3\\
18000	3350.4\\
20000	3329.6\\
22000	3372.2\\
24000	3402.5\\
26000	3366.3\\
32000	3436.3\\
34000	3318\\
36000	3367\\
38000	3399.5\\
40000	3377.5\\
42000	3395.9\\
44000	3352.2\\
46000	3361.5\\
48000	3357.4\\
50000	3317.4\\
52000	3335.5\\
54000	3399.2\\
56000	3364.2\\
58000	3372.1\\
60000	3296.8\\
62000	3366.5\\
64000	3420.7\\
66000	3325.1\\
68000	3414.8\\
72000	3305\\
74000	3378.3\\
76000	3384\\
78000	3361.1\\
80000	3352.5\\
82000	3414.5\\
84000	3461.2\\
86000	3345.4\\
88000	3391.1\\
90000	3333.9\\
92000	3346.8\\
94000	3431.9\\
96000	3339.3\\
98000	3364.6\\
1e+05	3319.7\\
};
\addplot [color=mycolor67, forget plot]
  table[row sep=crcr]{%
0	4395.3\\
2000	3248.4\\
4000	3140.6\\
8000	3114.5\\
10000	3081.9\\
12000	3129.9\\
14000	3120.7\\
16000	3097.6\\
18000	3087.7\\
20000	3126.3\\
22000	3151.3\\
24000	3123.2\\
26000	3081.3\\
28000	3077.2\\
30000	3115.2\\
32000	3186.9\\
34000	3110.5\\
36000	3135.4\\
38000	3127.2\\
40000	3095.3\\
42000	3172.8\\
44000	3073.4\\
46000	3125\\
48000	3190.9\\
50000	3141.4\\
52000	3134.9\\
54000	3066.5\\
56000	3088\\
58000	3069\\
60000	3131.1\\
62000	3109.1\\
64000	3127\\
66000	3122.2\\
68000	3076.5\\
70000	3133.1\\
72000	3161.8\\
74000	3161.5\\
76000	3066.9\\
78000	3136.5\\
80000	3067.4\\
82000	3080.6\\
84000	3148.8\\
86000	3139.2\\
88000	3069.4\\
90000	3081.5\\
92000	3051.9\\
94000	3061.3\\
96000	3079.8\\
98000	3028.5\\
1e+05	3082.3\\
};
\addplot [color=mycolor67]
  table[row sep=crcr]{%
0	4395.3\\
2000	2724\\
4000	2562.1\\
6000	2582.8\\
8000	2546.8\\
10000	2557.6\\
12000	2550.8\\
14000	2513.8\\
18000	2610.5\\
20000	2537\\
22000	2621.8\\
24000	2615.9\\
26000	2621.6\\
28000	2522.4\\
30000	2576.9\\
32000	2567.7\\
34000	2528.1\\
36000	2616.9\\
38000	2513\\
40000	2589.3\\
42000	2572.2\\
48000	2541.7\\
50000	2557.2\\
52000	2555.1\\
54000	2596.7\\
56000	2573.5\\
58000	2566.6\\
60000	2547.5\\
62000	2565.7\\
64000	2531.3\\
66000	2550.2\\
68000	2613.2\\
70000	2597.8\\
72000	2590.5\\
74000	2521.4\\
76000	2581.9\\
78000	2537.3\\
80000	2553.8\\
82000	2558.3\\
84000	2550.2\\
86000	2572.7\\
88000	2564.6\\
90000	2530.1\\
92000	2558.6\\
94000	2568.4\\
96000	2548.5\\
98000	2540.3\\
1e+05	2581.8\\
};
\addlegendentry{Classe k=43}

\addplot [color=mycolor68, forget plot]
  table[row sep=crcr]{%
0	4395.3\\
2000	2571\\
4000	2415.9\\
6000	2363.6\\
8000	2411.1\\
10000	2389.3\\
12000	2418.8\\
14000	2405.1\\
16000	2366.2\\
18000	2404\\
20000	2375.8\\
22000	2417\\
24000	2411.9\\
26000	2418.7\\
32000	2375\\
34000	2431.1\\
36000	2375.3\\
38000	2408.4\\
40000	2346.3\\
42000	2380.6\\
44000	2358.3\\
46000	2384.5\\
48000	2387.5\\
50000	2414.5\\
52000	2367.2\\
54000	2351.6\\
56000	2463.3\\
58000	2360.7\\
60000	2399.3\\
62000	2402.9\\
64000	2429.7\\
66000	2365.4\\
68000	2370.7\\
70000	2394.4\\
72000	2405.5\\
76000	2389.4\\
78000	2391.5\\
80000	2367.4\\
82000	2356.3\\
84000	2384.4\\
86000	2399.3\\
88000	2347.5\\
90000	2388.7\\
92000	2398\\
94000	2415.5\\
96000	2400.5\\
98000	2370.4\\
1e+05	2389.1\\
};
\addplot [color=mycolor69, forget plot]
  table[row sep=crcr]{%
0	4395.3\\
2000	3037\\
4000	2875.2\\
8000	2849.7\\
10000	2919.3\\
12000	2839.7\\
14000	2845.9\\
18000	2799\\
20000	2853.4\\
22000	2863.2\\
24000	2861.9\\
26000	2816.4\\
28000	2818.3\\
30000	2849.2\\
32000	2844.5\\
34000	2848.4\\
36000	2797.2\\
38000	2902.8\\
40000	2794.2\\
42000	2839.1\\
44000	2820.6\\
46000	2824.7\\
48000	2865.2\\
50000	2838.5\\
52000	2842.4\\
54000	2818.5\\
56000	2832.1\\
58000	2888\\
60000	2839.9\\
62000	2893.8\\
64000	2827.7\\
66000	2874.3\\
68000	2831.3\\
70000	2886.7\\
72000	2855.6\\
74000	2861.6\\
76000	2917.2\\
78000	2899\\
80000	2859.6\\
82000	2836.1\\
84000	2855.2\\
86000	2889\\
88000	2815.9\\
92000	2853.3\\
94000	2838.9\\
96000	2837.5\\
98000	2847.4\\
1e+05	2798.1\\
};
\addplot [color=mycolor70, forget plot]
  table[row sep=crcr]{%
0	4395.3\\
2000	2313.2\\
4000	2052.7\\
6000	2088.7\\
8000	1986.4\\
10000	2038.1\\
12000	2000.4\\
14000	2039.5\\
16000	2046.2\\
18000	2067.1\\
20000	2032.9\\
22000	2031.2\\
24000	2059\\
26000	2056.4\\
28000	2037.6\\
30000	2034.5\\
32000	2019.4\\
34000	2076.5\\
36000	2003.8\\
40000	2031.9\\
44000	2042.3\\
46000	2049.2\\
48000	2036.9\\
50000	2050.2\\
52000	2056\\
54000	2027.2\\
56000	2047.7\\
58000	1976.9\\
60000	2069.6\\
62000	1975.8\\
64000	2064\\
66000	2017.3\\
68000	2012.6\\
70000	2052.8\\
72000	2026.6\\
74000	2053.6\\
76000	2028.7\\
78000	2028.6\\
80000	2002.2\\
82000	2014.8\\
84000	2006.3\\
86000	2041.1\\
88000	2060.5\\
92000	2028.1\\
94000	2041.2\\
96000	2060.4\\
98000	2050.9\\
1e+05	2053\\
};
\addplot [color=mycolor71, forget plot]
  table[row sep=crcr]{%
0	4395.3\\
2000	2647.6\\
4000	2446.7\\
6000	2432\\
8000	2407.9\\
10000	2375\\
12000	2405.6\\
14000	2448.3\\
16000	2401\\
18000	2390.5\\
22000	2410.1\\
24000	2407.7\\
26000	2475\\
28000	2426.3\\
30000	2417.3\\
32000	2437.3\\
34000	2413.3\\
36000	2412\\
38000	2431.8\\
40000	2463.1\\
42000	2439.9\\
44000	2407.4\\
46000	2397.7\\
48000	2412.9\\
50000	2415.2\\
52000	2430.6\\
54000	2461.2\\
58000	2407\\
60000	2443.3\\
62000	2429.2\\
64000	2398.8\\
66000	2409.9\\
68000	2449.1\\
70000	2446.8\\
74000	2399.4\\
76000	2399.1\\
78000	2422\\
80000	2420.7\\
82000	2445.7\\
84000	2425.9\\
86000	2434\\
88000	2423.5\\
90000	2436.9\\
92000	2395.2\\
94000	2422.1\\
96000	2434.1\\
98000	2424.7\\
1e+05	2427.7\\
};
\addplot [color=mycolor72, forget plot]
  table[row sep=crcr]{%
0	4395.3\\
2000	2357.1\\
4000	2123.3\\
6000	2114.7\\
8000	2125.3\\
10000	2149.5\\
12000	2091\\
14000	2157\\
16000	2157.1\\
18000	2146.6\\
22000	2148.9\\
24000	2125.6\\
28000	2117.2\\
30000	2133.9\\
32000	2077.6\\
34000	2103.4\\
36000	2118.6\\
38000	2184.9\\
40000	2100.8\\
42000	2146.9\\
44000	2163.7\\
46000	2088.2\\
48000	2117.7\\
50000	2087\\
52000	2145.4\\
54000	2163.6\\
56000	2174.1\\
58000	2159.9\\
60000	2110.5\\
62000	2123.8\\
64000	2166.1\\
66000	2120.9\\
68000	2110.1\\
70000	2147.9\\
72000	2108.9\\
74000	2100.2\\
76000	2107.6\\
78000	2075\\
80000	2164.1\\
82000	2138.7\\
84000	2086\\
86000	2178.4\\
88000	2100.9\\
90000	2159.3\\
92000	2108.8\\
94000	2122.6\\
96000	2112.6\\
98000	2134.3\\
1e+05	2141.3\\
};
\addplot [color=mycolor73, forget plot]
  table[row sep=crcr]{%
0	4395.3\\
2000	2722.1\\
4000	2542\\
6000	2567.7\\
8000	2501.6\\
10000	2579.7\\
12000	2538.6\\
14000	2541.6\\
16000	2589.1\\
18000	2611\\
20000	2550.6\\
22000	2515\\
24000	2603.8\\
28000	2559.5\\
30000	2535.5\\
32000	2576\\
34000	2580.5\\
36000	2546\\
38000	2536.3\\
40000	2558.5\\
42000	2549.3\\
44000	2518.1\\
46000	2584.1\\
48000	2557.7\\
50000	2585.7\\
52000	2539.8\\
54000	2486.6\\
56000	2559\\
58000	2616.6\\
60000	2537.9\\
62000	2529.6\\
66000	2561\\
68000	2599.9\\
70000	2581\\
72000	2545.4\\
74000	2624\\
76000	2514.6\\
78000	2557.4\\
80000	2568.6\\
84000	2551\\
86000	2585.6\\
88000	2489.5\\
90000	2525\\
92000	2530.7\\
94000	2522\\
96000	2551.9\\
98000	2558.8\\
1e+05	2542.8\\
};
\addplot [color=mycolor74, forget plot]
  table[row sep=crcr]{%
0	4395.3\\
2000	1743.6\\
4000	1487\\
6000	1461.1\\
8000	1474.5\\
10000	1492.9\\
12000	1457.5\\
16000	1475.3\\
18000	1474.5\\
20000	1459.7\\
22000	1474.3\\
24000	1452.3\\
26000	1501\\
28000	1453.5\\
30000	1478.1\\
32000	1473.5\\
34000	1474.9\\
36000	1441.8\\
38000	1432.8\\
40000	1465.2\\
42000	1482.4\\
46000	1460.4\\
48000	1466.4\\
50000	1464.8\\
52000	1489.4\\
54000	1456.7\\
56000	1478.3\\
58000	1460.8\\
60000	1454.3\\
62000	1456.1\\
64000	1494.4\\
66000	1475.1\\
68000	1467.7\\
70000	1436.5\\
72000	1469.1\\
74000	1468.5\\
76000	1454.9\\
78000	1459.1\\
80000	1458.9\\
82000	1486.4\\
84000	1456.3\\
86000	1484.5\\
88000	1467.3\\
92000	1477.2\\
94000	1492.7\\
96000	1467.8\\
98000	1485.6\\
1e+05	1464.7\\
};
\addplot [color=mycolor75, forget plot]
  table[row sep=crcr]{%
0	4395.3\\
2000	2407.7\\
4000	2154.8\\
6000	2149.4\\
8000	2131.5\\
10000	2145.2\\
12000	2134.7\\
14000	2155.8\\
16000	2140.2\\
20000	2084\\
22000	2115.5\\
24000	2132.1\\
26000	2161.3\\
28000	2140.9\\
30000	2110.8\\
32000	2143.5\\
34000	2160.8\\
36000	2141.6\\
38000	2162.1\\
40000	2149.4\\
42000	2143.1\\
44000	2122.6\\
46000	2141.4\\
48000	2143.8\\
50000	2179.1\\
52000	2154.2\\
56000	2170.8\\
58000	2121.7\\
60000	2124.4\\
62000	2138.9\\
64000	2170\\
66000	2173.6\\
68000	2120.7\\
70000	2113.3\\
72000	2130.3\\
74000	2141.2\\
76000	2214.4\\
80000	2099.6\\
84000	2145.9\\
86000	2125.2\\
88000	2133\\
90000	2149.2\\
92000	2153.9\\
94000	2119\\
96000	2160.8\\
98000	2104.3\\
1e+05	2146.7\\
};
\addplot [color=mycolor76, forget plot]
  table[row sep=crcr]{%
0	4395.3\\
2000	2213.3\\
4000	1964\\
6000	1954\\
8000	1955.7\\
10000	1964.1\\
12000	1947.3\\
14000	1965.6\\
16000	1939.6\\
18000	1980.3\\
20000	1970.9\\
22000	1954.5\\
24000	1953.2\\
26000	1959.2\\
28000	1959.5\\
30000	1967.2\\
32000	1957.3\\
34000	1962.2\\
36000	1940.3\\
38000	1967.8\\
40000	1940.6\\
42000	1932\\
44000	1934.8\\
46000	1931.7\\
48000	1963.3\\
50000	1909.8\\
52000	1949.6\\
54000	1925.3\\
56000	1945.5\\
58000	1948.7\\
60000	1935.9\\
62000	1941.1\\
64000	1963.2\\
66000	1912.5\\
68000	1979.3\\
70000	1964\\
72000	1964.6\\
74000	1958.1\\
76000	1974.1\\
78000	1972.8\\
80000	1980.5\\
82000	1948.2\\
84000	1951.3\\
86000	1930\\
88000	1971.4\\
90000	1957.5\\
92000	1929.9\\
98000	1986.4\\
1e+05	1953.6\\
};
\addplot [color=mycolor77, forget plot]
  table[row sep=crcr]{%
0	4395.3\\
2000	1784.8\\
4000	1475.3\\
6000	1509\\
8000	1506\\
10000	1495.6\\
12000	1495.3\\
14000	1499.4\\
16000	1492\\
20000	1514.6\\
22000	1501\\
24000	1462\\
28000	1513.3\\
30000	1487.6\\
32000	1517.3\\
34000	1464.8\\
36000	1507.6\\
38000	1476\\
40000	1501.4\\
42000	1494.1\\
44000	1502\\
46000	1486\\
48000	1501.5\\
50000	1523.5\\
52000	1510.5\\
54000	1490.3\\
56000	1474.7\\
58000	1533.9\\
60000	1507.1\\
62000	1476.3\\
64000	1486.2\\
68000	1494.1\\
70000	1508.6\\
72000	1504.6\\
74000	1526.9\\
76000	1480.4\\
78000	1499.8\\
82000	1511.7\\
84000	1487.8\\
86000	1494.5\\
90000	1481.8\\
92000	1507.6\\
94000	1499.7\\
96000	1549.7\\
98000	1476.7\\
1e+05	1487.1\\
};
\addplot [color=mycolor78, forget plot]
  table[row sep=crcr]{%
0	4395.3\\
2000	2502.1\\
4000	2213.8\\
6000	2249.5\\
8000	2188.4\\
10000	2233\\
12000	2229.5\\
14000	2258.8\\
16000	2203.8\\
18000	2169.1\\
20000	2221.2\\
22000	2248.2\\
26000	2182.2\\
28000	2238.1\\
30000	2238.2\\
32000	2195.6\\
34000	2212.6\\
36000	2270.2\\
38000	2235.1\\
40000	2253.8\\
42000	2220.5\\
44000	2217.3\\
48000	2186.9\\
50000	2259.6\\
54000	2208.3\\
56000	2206.2\\
58000	2238.6\\
60000	2240.2\\
64000	2221\\
66000	2290.9\\
68000	2218.8\\
70000	2270.2\\
74000	2221.4\\
76000	2229.9\\
78000	2180.9\\
80000	2207.1\\
82000	2249.3\\
84000	2275.2\\
86000	2243.9\\
88000	2205.5\\
90000	2205.7\\
94000	2184.6\\
96000	2193.5\\
98000	2248.8\\
1e+05	2251.3\\
};
\addplot [color=mycolor79, forget plot]
  table[row sep=crcr]{%
0	4395.3\\
2000	2401.2\\
4000	2063.6\\
6000	2109.3\\
8000	2060\\
10000	2068.2\\
14000	2126.8\\
16000	2092.6\\
18000	2070.8\\
20000	2094.8\\
24000	2092.4\\
26000	2118.5\\
28000	2054.5\\
30000	2104.4\\
32000	2067.6\\
34000	2113\\
36000	2099.4\\
38000	2065.6\\
40000	2093.6\\
42000	2086\\
44000	2135.6\\
46000	2118.4\\
48000	2083.5\\
50000	2061.3\\
54000	2052.1\\
56000	2097.6\\
58000	2103\\
60000	2142.9\\
62000	2062.1\\
64000	2111.3\\
66000	2066.5\\
68000	2086.7\\
70000	2095.1\\
72000	2111.8\\
76000	2082.3\\
78000	2131.3\\
80000	2114.5\\
82000	2038.9\\
84000	2129.2\\
86000	2090.3\\
90000	2075.2\\
92000	2172.5\\
94000	2100.2\\
96000	2071\\
98000	2077.7\\
1e+05	2141.8\\
};
\addplot [color=mycolor79, forget plot]
  table[row sep=crcr]{%
0	4395.3\\
2000	2315.7\\
4000	1937.7\\
6000	1908.8\\
8000	1939\\
10000	1959.9\\
14000	1953.3\\
16000	1877.3\\
18000	1914.2\\
20000	1958\\
22000	1971.8\\
24000	2003\\
26000	1934.8\\
28000	1926.3\\
30000	1946\\
32000	1878.3\\
34000	2005\\
36000	1910.9\\
38000	1902.7\\
40000	2030.9\\
42000	1945.3\\
44000	1926.8\\
46000	2000.4\\
48000	1909.8\\
50000	1970.6\\
52000	1928.3\\
54000	1910.4\\
56000	1998.5\\
58000	1954.2\\
60000	1951.9\\
62000	1919.7\\
64000	1910.5\\
66000	1885.5\\
68000	1948.1\\
70000	1948.5\\
72000	1962.1\\
74000	1944.9\\
76000	1942.3\\
78000	1926\\
80000	1941.8\\
82000	1936.6\\
84000	1978.3\\
86000	1956.1\\
88000	1951.2\\
90000	1968.4\\
92000	1972.5\\
94000	1922.1\\
96000	1938\\
98000	1937.2\\
1e+05	1902.9\\
};
\addplot [color=mycolor80, forget plot]
  table[row sep=crcr]{%
0	4395.3\\
2000	1816.3\\
4000	1474\\
6000	1480.3\\
8000	1508.2\\
10000	1482.3\\
12000	1480.7\\
14000	1484.9\\
16000	1466.5\\
18000	1498.1\\
20000	1483.1\\
22000	1489.4\\
24000	1484.7\\
26000	1491.7\\
28000	1438.9\\
30000	1488.9\\
32000	1482\\
34000	1492.1\\
36000	1479\\
38000	1472.3\\
40000	1472.9\\
42000	1479.9\\
44000	1474.9\\
46000	1481\\
48000	1478.4\\
50000	1489.3\\
52000	1486.4\\
54000	1501.2\\
58000	1494.1\\
60000	1493\\
62000	1459.3\\
64000	1478.1\\
66000	1480.3\\
68000	1509.1\\
70000	1492\\
72000	1491.3\\
74000	1479.5\\
76000	1488.4\\
78000	1480.5\\
80000	1526.6\\
82000	1498.5\\
84000	1485\\
86000	1505.5\\
88000	1468.7\\
90000	1491.1\\
92000	1494.9\\
94000	1506.4\\
96000	1491.7\\
98000	1484.3\\
1e+05	1499.4\\
};
\addplot [color=mycolor81, forget plot]
  table[row sep=crcr]{%
0	4395.3\\
2000	1728.1\\
4000	1397.1\\
6000	1387.6\\
8000	1369.7\\
10000	1369.6\\
12000	1392.3\\
14000	1401.8\\
16000	1404.9\\
18000	1392\\
20000	1407.7\\
22000	1380.1\\
24000	1412.3\\
26000	1394.4\\
28000	1382.6\\
30000	1349.3\\
32000	1397.7\\
36000	1386.7\\
38000	1400.5\\
40000	1383.6\\
42000	1402.5\\
44000	1398.4\\
46000	1381.4\\
48000	1395\\
50000	1420.2\\
52000	1412.1\\
54000	1374.8\\
56000	1407.7\\
58000	1368\\
62000	1413.3\\
64000	1411.5\\
66000	1370.9\\
68000	1375.4\\
70000	1407\\
74000	1394.8\\
76000	1382.2\\
78000	1386.7\\
80000	1402.3\\
82000	1390.6\\
84000	1402.3\\
86000	1385.3\\
88000	1403.9\\
90000	1396.9\\
92000	1361.6\\
94000	1406.7\\
96000	1386.8\\
98000	1400.4\\
1e+05	1376.6\\
};
\addplot [color=mycolor82, forget plot]
  table[row sep=crcr]{%
0	4395.3\\
2000	1933.5\\
4000	1626.9\\
6000	1626.1\\
8000	1619.7\\
14000	1625.2\\
18000	1654.9\\
20000	1658.3\\
22000	1626\\
24000	1639.2\\
26000	1675.4\\
28000	1621.4\\
30000	1622\\
32000	1579.4\\
34000	1611.2\\
36000	1637.9\\
38000	1655\\
40000	1614.5\\
42000	1652.9\\
44000	1623.4\\
46000	1680.3\\
48000	1631.5\\
50000	1645.2\\
52000	1630.3\\
54000	1650.3\\
56000	1618\\
58000	1641.2\\
60000	1596\\
62000	1650.8\\
66000	1622.8\\
68000	1639.7\\
72000	1609.7\\
74000	1622.3\\
78000	1622.3\\
80000	1633.8\\
82000	1627.6\\
84000	1683.9\\
86000	1638.8\\
88000	1651\\
90000	1674.7\\
92000	1648.4\\
94000	1628\\
96000	1626.1\\
98000	1637.3\\
1e+05	1622.1\\
};
\addplot [color=mycolor83, forget plot]
  table[row sep=crcr]{%
0	4395.3\\
2000	2011.4\\
4000	1660.9\\
6000	1645.1\\
8000	1659.8\\
10000	1643.4\\
12000	1681.1\\
14000	1636.2\\
16000	1656.5\\
18000	1683.3\\
22000	1658.8\\
24000	1642.5\\
28000	1657.9\\
32000	1662\\
34000	1680.1\\
36000	1654\\
38000	1676.1\\
40000	1657.1\\
42000	1665.2\\
44000	1645.7\\
46000	1656.7\\
48000	1643.5\\
50000	1655.9\\
52000	1652.3\\
54000	1614.5\\
56000	1652.9\\
58000	1668.2\\
60000	1638.4\\
66000	1671.1\\
68000	1653.9\\
70000	1678\\
72000	1674.1\\
74000	1661.8\\
76000	1658.9\\
78000	1678.2\\
80000	1682.5\\
82000	1656\\
84000	1684.2\\
86000	1680.3\\
88000	1682.6\\
90000	1657.4\\
92000	1678.5\\
94000	1648.5\\
96000	1655.8\\
98000	1638\\
1e+05	1666.1\\
};
\addplot [color=mycolor83, forget plot]
  table[row sep=crcr]{%
0	4395.3\\
2000	1569.5\\
4000	1200.4\\
6000	1224.4\\
8000	1217.7\\
10000	1227.5\\
12000	1201\\
14000	1183.4\\
16000	1212.3\\
18000	1258.4\\
20000	1190.3\\
22000	1197.2\\
24000	1190.2\\
28000	1188.2\\
30000	1177.7\\
32000	1198\\
34000	1212.6\\
36000	1183.7\\
38000	1217.7\\
40000	1204.2\\
42000	1214.5\\
44000	1199.3\\
46000	1205.9\\
48000	1200.2\\
50000	1184.8\\
52000	1193.3\\
54000	1195.1\\
56000	1222.4\\
58000	1196.7\\
60000	1214.9\\
62000	1200.6\\
64000	1193\\
66000	1206.1\\
68000	1206.2\\
70000	1216.7\\
74000	1216.9\\
76000	1204.5\\
78000	1224.3\\
80000	1196.6\\
82000	1193.7\\
84000	1187.2\\
86000	1175.1\\
88000	1196.7\\
90000	1181.6\\
92000	1208.5\\
94000	1190.4\\
96000	1197.8\\
98000	1238\\
1e+05	1214.5\\
};
\addplot [color=mycolor84, forget plot]
  table[row sep=crcr]{%
0	4395.3\\
2000	1621.3\\
4000	1210.3\\
6000	1189.5\\
8000	1199.6\\
10000	1214.8\\
12000	1203.9\\
14000	1185.2\\
16000	1237.1\\
18000	1206.7\\
20000	1183.1\\
22000	1255.9\\
24000	1231.7\\
26000	1243.8\\
28000	1208.9\\
30000	1197.3\\
32000	1217.7\\
34000	1190.2\\
36000	1201.4\\
38000	1199.6\\
40000	1184.7\\
42000	1214\\
44000	1217.6\\
46000	1212.7\\
48000	1220.2\\
50000	1206.2\\
52000	1201.5\\
54000	1228.6\\
56000	1195.9\\
62000	1215.1\\
64000	1180.8\\
66000	1180.8\\
68000	1257.4\\
70000	1250.1\\
72000	1176.1\\
74000	1235.3\\
76000	1188.9\\
78000	1221.7\\
80000	1265.6\\
82000	1181.9\\
84000	1225.4\\
86000	1206.5\\
88000	1210.5\\
90000	1203.9\\
92000	1201\\
94000	1229.5\\
96000	1240.6\\
98000	1234.2\\
1e+05	1176.5\\
};
\addplot [color=mycolor85, forget plot]
  table[row sep=crcr]{%
0	4395.3\\
2000	1744.4\\
4000	1355.1\\
6000	1319.4\\
8000	1327.4\\
12000	1361.2\\
14000	1372.3\\
16000	1347.8\\
18000	1302.4\\
20000	1353.5\\
22000	1352.6\\
24000	1374.7\\
26000	1377.3\\
28000	1365.3\\
30000	1327.2\\
32000	1328.2\\
34000	1335\\
36000	1319.6\\
38000	1340.5\\
40000	1342.2\\
42000	1365.7\\
46000	1291.9\\
48000	1325.1\\
50000	1317.7\\
52000	1331.3\\
54000	1357.1\\
56000	1334.8\\
58000	1354.8\\
60000	1340.2\\
62000	1385.8\\
64000	1362.7\\
66000	1386.5\\
68000	1309\\
70000	1341.8\\
72000	1347.8\\
74000	1331.5\\
76000	1325.4\\
78000	1365\\
80000	1334.7\\
82000	1317.7\\
84000	1322.5\\
86000	1304.1\\
88000	1313.9\\
90000	1354.6\\
92000	1341.8\\
94000	1334.5\\
96000	1337.7\\
98000	1305.2\\
1e+05	1364.9\\
};
\addplot [color=mycolor86, forget plot]
  table[row sep=crcr]{%
0	4395.3\\
2000	1556\\
4000	1158.4\\
6000	1148\\
8000	1107.9\\
10000	1133.6\\
12000	1133.9\\
14000	1124\\
16000	1107.4\\
18000	1108.1\\
20000	1098.5\\
22000	1093.8\\
24000	1105.3\\
26000	1107.4\\
28000	1105\\
30000	1141.4\\
32000	1129.8\\
34000	1088.7\\
36000	1162\\
38000	1100.8\\
40000	1085.6\\
44000	1122.5\\
46000	1100.7\\
48000	1118.6\\
50000	1112.5\\
52000	1090.1\\
54000	1077.5\\
56000	1105.5\\
58000	1109.6\\
60000	1095.2\\
62000	1112.6\\
64000	1071.5\\
66000	1131\\
68000	1099.2\\
70000	1105.1\\
72000	1095.4\\
74000	1089.4\\
76000	1099.6\\
78000	1091.7\\
80000	1154.9\\
82000	1086.1\\
84000	1131.4\\
86000	1103.6\\
88000	1108.9\\
90000	1148.8\\
92000	1113.8\\
94000	1146.4\\
96000	1150.3\\
98000	1124.4\\
1e+05	1140.3\\
};
\addplot [color=mycolor87, forget plot]
  table[row sep=crcr]{%
0	4395.3\\
2000	1303\\
4000	921.51\\
6000	956.54\\
8000	943.1\\
10000	952.06\\
12000	924.99\\
14000	974.96\\
16000	923.8\\
18000	890.22\\
20000	933.58\\
22000	970.53\\
24000	954.85\\
26000	967.4\\
28000	950.98\\
30000	960.36\\
32000	965.9\\
36000	957.96\\
38000	884.37\\
40000	956.11\\
42000	972.59\\
44000	886.31\\
46000	916.27\\
48000	916.36\\
50000	937.55\\
52000	984.27\\
54000	891.59\\
56000	893.35\\
58000	939.13\\
60000	967.74\\
62000	915.29\\
64000	928.16\\
66000	921.08\\
68000	953.37\\
70000	977.8\\
72000	1045.8\\
74000	935.14\\
76000	878\\
78000	942.64\\
80000	923.36\\
82000	883.72\\
84000	947.21\\
86000	917.26\\
90000	930.91\\
92000	996.09\\
94000	944.03\\
96000	936.52\\
98000	908.64\\
1e+05	968.78\\
};
\addplot [color=mycolor88, forget plot]
  table[row sep=crcr]{%
0	4395.3\\
2000	1591.3\\
4000	1112\\
6000	1113.3\\
8000	1161.3\\
10000	1140.6\\
12000	1077.6\\
14000	1117.9\\
16000	1126.7\\
18000	1114.7\\
20000	1110.7\\
22000	1133.5\\
24000	1125\\
26000	1109.3\\
28000	1102.6\\
30000	1110.7\\
32000	1092.4\\
34000	1136.4\\
36000	1101\\
38000	1100.3\\
40000	1129\\
42000	1151.4\\
44000	1123.6\\
46000	1083.2\\
48000	1111\\
50000	1095.8\\
52000	1096.1\\
54000	1144.9\\
56000	1118.5\\
58000	1125.8\\
60000	1155.4\\
62000	1111\\
64000	1124.3\\
66000	1108.8\\
68000	1136.8\\
70000	1127.7\\
72000	1144.1\\
74000	1146.9\\
76000	1091.6\\
78000	1091.6\\
80000	1127.8\\
82000	1116.7\\
84000	1123.9\\
86000	1147.9\\
88000	1122.4\\
90000	1108.3\\
92000	1126\\
94000	1102.8\\
96000	1097.5\\
98000	1110.8\\
1e+05	1088\\
};
\addplot [color=mycolor89, forget plot]
  table[row sep=crcr]{%
0	4395.3\\
2000	1284.1\\
4000	793.7\\
6000	815.69\\
8000	813.62\\
10000	778.65\\
12000	802.36\\
14000	821.52\\
16000	813.41\\
18000	781.83\\
20000	799.63\\
22000	812.54\\
24000	800.13\\
26000	792.94\\
28000	815.16\\
30000	812.47\\
32000	802.12\\
34000	783.4\\
38000	778.85\\
40000	789.64\\
42000	790.84\\
44000	814.17\\
46000	787.03\\
48000	781.89\\
50000	805.6\\
52000	790.57\\
54000	790.41\\
56000	792.7\\
58000	791.63\\
60000	773.85\\
62000	789.9\\
64000	786.54\\
66000	797.06\\
68000	786.59\\
70000	798.29\\
72000	819.08\\
74000	785.15\\
76000	788.81\\
78000	800.86\\
80000	778.4\\
82000	802.64\\
84000	779.51\\
86000	802.69\\
88000	791.75\\
90000	769.52\\
92000	808.83\\
94000	780.82\\
96000	770.53\\
98000	777.21\\
1e+05	801.08\\
};
\addplot [color=mycolor90, forget plot]
  table[row sep=crcr]{%
0	4395.3\\
2000	1236.1\\
4000	780.04\\
6000	796.54\\
8000	782.03\\
10000	744.9\\
12000	771.09\\
14000	814.38\\
16000	778.23\\
18000	773.28\\
20000	775.6\\
22000	783.94\\
24000	779.4\\
26000	787.73\\
28000	762.27\\
30000	789.89\\
32000	793.26\\
34000	780.25\\
36000	770.23\\
38000	756.12\\
40000	768.09\\
42000	767.35\\
44000	777.79\\
46000	784.07\\
48000	776.03\\
50000	770.92\\
52000	785.97\\
54000	783.44\\
56000	799.83\\
58000	790.97\\
60000	775.43\\
62000	810.89\\
64000	786.22\\
66000	781.74\\
68000	766.89\\
70000	792.45\\
74000	756.36\\
76000	772.26\\
78000	749.19\\
80000	783.08\\
82000	792.83\\
84000	746.64\\
86000	779.98\\
88000	761.39\\
90000	760.03\\
92000	755.86\\
94000	744.54\\
96000	755.23\\
98000	787.31\\
1e+05	777.55\\
};
\addplot [color=mycolor91, forget plot]
  table[row sep=crcr]{%
0	4395.3\\
2000	1220.6\\
4000	644.79\\
6000	658.97\\
8000	666.55\\
10000	691.83\\
12000	668.88\\
14000	673.56\\
16000	638.52\\
18000	655.7\\
20000	644.61\\
22000	662.66\\
24000	660.23\\
26000	679.82\\
28000	645.91\\
30000	659.11\\
32000	675.64\\
34000	660.49\\
36000	662.44\\
38000	677.52\\
40000	660.93\\
42000	684.11\\
44000	633.72\\
46000	685.18\\
50000	665.64\\
52000	681.25\\
54000	659.27\\
56000	656.55\\
58000	673.39\\
60000	648.95\\
64000	660.27\\
66000	676.29\\
70000	663.02\\
72000	647.02\\
74000	674.35\\
76000	655.84\\
78000	642.64\\
80000	682.6\\
82000	657.3\\
84000	680.34\\
86000	688.15\\
88000	653.91\\
90000	710.13\\
92000	664.32\\
94000	663.9\\
96000	678.15\\
98000	653.24\\
1e+05	666.33\\
};
\addplot [color=mycolor92, forget plot]
  table[row sep=crcr]{%
0	4395.3\\
2000	1429.7\\
4000	848.45\\
6000	820.19\\
8000	865.67\\
10000	843.57\\
12000	844.38\\
14000	832.2\\
16000	835.35\\
18000	855.43\\
20000	849.26\\
22000	829.58\\
26000	847.56\\
28000	834.02\\
30000	813.59\\
32000	837.86\\
34000	819.56\\
36000	842.93\\
38000	823.03\\
42000	860.44\\
44000	864.75\\
46000	827.16\\
48000	876.94\\
54000	814.66\\
56000	812.54\\
58000	840.68\\
60000	838.3\\
62000	842.94\\
64000	855.61\\
66000	840.16\\
68000	842.17\\
70000	831.63\\
72000	825.36\\
74000	823.15\\
76000	804.25\\
78000	812.32\\
80000	824.11\\
82000	844.91\\
84000	816.92\\
86000	816.95\\
88000	825.74\\
90000	815.09\\
92000	845.79\\
94000	844.47\\
96000	836.32\\
98000	798.23\\
1e+05	816.33\\
};
\addplot [color=mycolor93, forget plot]
  table[row sep=crcr]{%
0	4395.3\\
2000	1273.7\\
4000	665.57\\
6000	697.82\\
8000	705.67\\
10000	693.6\\
12000	753.79\\
14000	671.74\\
16000	678.84\\
18000	714.71\\
20000	702.41\\
22000	710.32\\
24000	646.77\\
26000	686.04\\
28000	667.87\\
30000	704.54\\
32000	696.81\\
34000	716.99\\
36000	713.73\\
38000	749.44\\
40000	709.92\\
42000	665.95\\
44000	647.78\\
46000	690.18\\
48000	673.11\\
50000	663.7\\
52000	675.25\\
54000	696.65\\
56000	700.69\\
58000	688.34\\
60000	659.19\\
62000	693.15\\
64000	702.29\\
66000	673.24\\
68000	688.15\\
70000	671.84\\
72000	688.16\\
74000	700.94\\
76000	711.72\\
78000	671.56\\
80000	705.65\\
82000	721.75\\
84000	678.14\\
86000	657.05\\
88000	701.98\\
90000	711.06\\
92000	685\\
94000	720.86\\
96000	703.29\\
98000	702.42\\
1e+05	669.21\\
};
\addplot [color=mycolor94, forget plot]
  table[row sep=crcr]{%
0	4395.3\\
2000	1278.2\\
4000	672.8\\
6000	651.62\\
8000	671.66\\
10000	624.52\\
12000	646.02\\
14000	625.28\\
16000	646.83\\
20000	657.94\\
22000	650.39\\
24000	671.21\\
26000	664.17\\
28000	665.76\\
30000	654.22\\
32000	657.44\\
34000	641.51\\
36000	664.19\\
38000	665.8\\
40000	648.34\\
42000	674.34\\
44000	648.21\\
46000	652.58\\
48000	659.43\\
50000	647.03\\
52000	670.95\\
54000	669.18\\
56000	648.43\\
58000	641.8\\
60000	677.45\\
62000	653.13\\
64000	687.5\\
66000	662.13\\
68000	644.51\\
70000	649.55\\
72000	659.98\\
74000	662.99\\
76000	683.85\\
78000	675.97\\
80000	639.98\\
82000	639.48\\
84000	647.88\\
86000	684.5\\
88000	659.05\\
90000	672.45\\
92000	660.43\\
96000	642.4\\
98000	669.02\\
1e+05	644.31\\
};
\addplot [color=mycolor95, forget plot]
  table[row sep=crcr]{%
0	4395.3\\
2000	1005.4\\
4000	369.27\\
6000	392.64\\
8000	393.31\\
10000	398.19\\
12000	382.43\\
14000	392.9\\
16000	380.18\\
18000	393.98\\
20000	374.72\\
22000	372.1\\
24000	377.62\\
26000	368.84\\
28000	397.22\\
30000	383.28\\
32000	372.91\\
34000	378\\
36000	399.01\\
38000	391.46\\
40000	378.62\\
42000	402.24\\
44000	363.43\\
46000	373\\
48000	384.54\\
50000	394.03\\
52000	384.3\\
54000	392.62\\
56000	354.81\\
58000	378.38\\
60000	366.89\\
62000	374.14\\
64000	386.04\\
66000	382.19\\
68000	398.5\\
70000	403.05\\
72000	380.15\\
74000	405.48\\
76000	373.11\\
78000	382.05\\
80000	388.94\\
82000	400.42\\
84000	393.11\\
86000	389.51\\
88000	374.63\\
90000	382.85\\
92000	359.89\\
94000	386.05\\
96000	400.88\\
98000	364.18\\
1e+05	390.97\\
};
\addplot [color=mycolor96, forget plot]
  table[row sep=crcr]{%
0	4395.3\\
2000	1154.9\\
4000	403.97\\
6000	406.37\\
8000	391.96\\
10000	434.95\\
12000	437.64\\
14000	397.03\\
16000	396.86\\
18000	430.2\\
20000	410.22\\
22000	424.66\\
24000	436.49\\
26000	383.41\\
28000	410.48\\
30000	408.94\\
32000	433.96\\
34000	407.97\\
36000	426.58\\
38000	422.91\\
40000	388.77\\
42000	391.68\\
44000	408.97\\
48000	391.58\\
50000	431.04\\
52000	387.26\\
54000	414.47\\
56000	417.61\\
58000	402.07\\
60000	431.07\\
62000	391.24\\
64000	422.45\\
66000	407.29\\
68000	408.22\\
70000	430.27\\
72000	381.21\\
74000	415.45\\
76000	422.61\\
78000	393.86\\
80000	395.24\\
82000	417.73\\
84000	407.16\\
86000	430.32\\
90000	432.02\\
92000	386.24\\
94000	432.92\\
96000	446.31\\
98000	407.38\\
1e+05	403.57\\
};
\addplot [color=mycolor97, forget plot]
  table[row sep=crcr]{%
0	4395.3\\
2000	965.68\\
4000	283.26\\
6000	259.81\\
8000	292.1\\
10000	318.53\\
12000	277.48\\
14000	271.72\\
16000	297.03\\
18000	281.23\\
20000	276\\
22000	267.98\\
24000	286.45\\
26000	262.18\\
28000	316.21\\
30000	283.24\\
32000	272.49\\
34000	290.99\\
36000	271.78\\
38000	310.98\\
40000	280.48\\
42000	293.71\\
44000	283.27\\
46000	296.95\\
48000	277.02\\
50000	298.4\\
52000	302.97\\
56000	293.32\\
58000	272.07\\
60000	258.6\\
62000	252.05\\
64000	297.34\\
68000	279.69\\
70000	289.67\\
72000	274.04\\
74000	273.94\\
76000	283.44\\
78000	280.58\\
80000	287.68\\
82000	293.12\\
84000	305.91\\
86000	304.58\\
88000	282.02\\
90000	267.71\\
92000	293.93\\
94000	291.24\\
96000	293.02\\
98000	286.29\\
1e+05	273.16\\
};
\addplot [color=mycolor98]
  table[row sep=crcr]{%
0	4395.3\\
2000	1220.4\\
4000	396.06\\
6000	363.06\\
8000	384.35\\
10000	395.77\\
12000	377.56\\
14000	383.74\\
16000	381.23\\
18000	385.67\\
20000	405.71\\
22000	395.79\\
24000	404.5\\
26000	378.59\\
30000	392.44\\
32000	386.62\\
34000	395.19\\
36000	407.91\\
38000	382.52\\
40000	400.6\\
42000	382.33\\
44000	385.35\\
46000	414\\
48000	384.02\\
50000	399.9\\
52000	388.35\\
54000	360.45\\
56000	397.64\\
58000	392.9\\
60000	394.77\\
62000	370.8\\
64000	383.23\\
66000	369.41\\
68000	360.2\\
70000	375.55\\
72000	382.57\\
74000	396.74\\
76000	377.89\\
78000	379.08\\
80000	397.65\\
82000	396.81\\
84000	347.08\\
86000	367.61\\
88000	385.24\\
90000	372.47\\
92000	404.96\\
94000	388.46\\
96000	388.13\\
98000	379.84\\
1e+05	393.67\\
};
\addlegendentry{Classe k=10}

            % \input{../../../../.tex/20/KS/SizeEvolutionsfig10.tex}
            \legend{}
        \end{groupplot}
    %\endgroup

    %\begingroup Parametri
        
%\begingroup Parametro
\newcommand\varyingParameterSymbol\lambda
\newcommand\varyingParameterList[1]{
    \ifcase#1\relax%
    \or\num{1e-2}% Caso 1
    \or\num{8.2e-2}% Caso 2
    \or\num{1.55e-1}% Caso 3
    \or\num{2.27e-1}% Caso 4
    \else\num{3e-1}% Caso 5
    \fi
}
%\endgroup

%\begingroup beta
\newcommand\betaExactList[1]{
    \ifcase#1\relax%
    \or\num{1.267+-0.007}% Caso 1
    \or\num{1.285+-0.010}% Caso 2
    \or\num{1.264+-0.015}% Caso 3
    \or\num{1.246+-0.015}% Caso 4
    \else\num{1.210+-0.015}% Caso 5
    \fi
}
\newcommand\betaAprxList[1]{
    \ifcase#1\relax%
    \or\num{1.264+-0.006}% Caso 1
    \or\num{1.270+-0.011}% Caso 2
    \or\num{1.280+-0.013}% Caso 3
    \or\num{1.250+-0.015}% Caso 4
    \else\num{1.218+-0.014}% Caso 5
    \fi
}
%\endgroup

%\begingroup xi
\newcommand\xiExactList[1]{
    \ifcase#1\relax%
    \or\num{0.757+-0.004}% Caso 1
    \or\num{0.710+-0.026}% Caso 2
    \or\num{0.670+-0.035}% Caso 3
    \or\num{0.625+-0.052}% Caso 4
    \else\num{0.539+-0.059}% Caso 5
    \fi
}
\newcommand\xiAprxList[1]{
    \ifcase#1\relax%
    \or\num{0.749+-0.004}% Caso 1
    \or\num{0.721+-0.022}% Caso 2
    \or\num{0.645+-0.040}% Caso 3
    \or\num{0.622+-0.051}% Caso 4
    \else\num{0.567+-0.053}% Caso 5
    \fi
}
%\endgroup

%\begingroup mI
\newcommand\mIExactList[1]{
    \ifcase#1\relax%
    \or\num{3.265+-0.001}% Caso 1
    \or\num{3.238+-0.005}% Caso 2
    \or\num{3.220+-0.008}% Caso 3
    \or\num{3.213+-0.013}% Caso 4
    \else\num{3.174+-0.023}% Caso 5
    \fi
}
\newcommand\mIAprxList[1]{
    \ifcase#1\relax%
    \or\num{3.262+-0.001}% Caso 1
    \or\num{3.240+-0.005}% Caso 2
    \or\num{3.217+-0.010}% Caso 3
    \or\num{3.217+-0.012}% Caso 4
    \else\num{3.208+-0.009}% Caso 5
    \fi
}
%\endgroup

%\begingroup sI
\newcommand\sIExactList[1]{
    \ifcase#1\relax%
    \or\num{0.284+-0.001}% Caso 1
    \or\num{0.311+-0.008}% Caso 2
    \or\num{0.343+-0.010}% Caso 3
    \or\num{0.361+-0.020}% Caso 4
    \else\num{0.398+-0.029}% Caso 5
    \fi
}
\newcommand\sIAprxList[1]{
    \ifcase#1\relax%
    \or\num{0.281+-0.001}% Caso 1
    \or\num{0.315+-0.006}% Caso 2
    \or\num{0.334+-0.013}% Caso 3
    \or\num{0.366+-0.018}% Caso 4
    \else\num{0.411+-0.025}% Caso 5
    \fi
}
%\endgroup

%\begingroup mII
\newcommand\mIIExactList[1]{
    \ifcase#1\relax%
    \or\num{3.788+-0.005}% Caso 1
    \or\num{3.730+-0.029}% Caso 2
    \or\num{3.681+-0.052}% Caso 3
    \or\num{3.650+-0.072}% Caso 4
    \else\num{3.572+-0.095}% Caso 5
    \fi
}
\newcommand\mIIAprxList[1]{
    \ifcase#1\relax%
    \or\num{3.778+-0.006}% Caso 1
    \or\num{3.735+-0.029}% Caso 2
    \or\num{3.656+-0.048}% Caso 3
    \or\num{3.658+-0.084}% Caso 4
    \else\num{3.502+-0.078}% Caso 5
    \fi
}
%\endgroup

%\begingroup sII
\newcommand\sIIExactList[1]{
    \ifcase#1\relax%
    \or\num{0.455+-0.002}% Caso 1
    \or\num{0.486+-0.007}% Caso 2
    \or\num{0.515+-0.014}% Caso 3
    \or\num{0.533+-0.021}% Caso 4
    \else\num{0.570+-0.028}% Caso 5
    \fi
}
\newcommand\sIIAprxList[1]{
    \ifcase#1\relax%
    \or\num{0.457+-0.002}% Caso 1
    \or\num{0.490+-0.007}% Caso 2
    \or\num{0.521+-0.012}% Caso 3
    \or\num{0.545+-0.021}% Caso 4
    \else\num{0.586+-0.029}% Caso 5
    \fi
}
%\endgroup

        \ifpageisodd\writeVaryingParameter{1}{west}{east}{-.25em}{0em}
        \else\writeVaryingParameter{2}{east}{west}{+.25em}{0em}\fi
    %\endgroup

    %\begingroup Didascalie
        \writeSubCaptions[2]{south west}{south west}{0em}{0em}
        % https://tex.stackexchange.com/questions/632045/get-subcaption-for-every-plot-using-groupplot-pgf-tikz-after-2020-update
        \writeSizeEvolutionDescriptions[1]
    %\endgroup

    %\begingroup Titoli
        % \node[title,overlay] at (Row1 c1r1.north) {Esatto};
        % \node[title,overlay] at (Row1 c2r1.north) {Approssimato};
    %\endgroup

    \redefineTikZbounds{Row1 c1r1}{Row5 c2r1}
\end{tikzpicture}%
